\documentclass[11pt,letterpaper]{article}
\usepackage[T1]{fontenc}
\usepackage[utf8]{inputenc}
\usepackage{lmodern,microtype}
\usepackage[margin=1in,headheight=14pt]{geometry}
\usepackage{amsmath,amssymb,amsthm,mathtools}
\usepackage{booktabs,array,tabularx,longtable,enumitem}
\usepackage[dvipsnames]{xcolor}
\usepackage[colorlinks=true,linkcolor=MidnightBlue,citecolor=MidnightBlue,urlcolor=MidnightBlue]{hyperref}
\usepackage{fancyhdr}
\hypersetup{pdftitle={Observation and Recovery in Arithmetic and Probability: Prime Flows, Rare Events, and Infinite Mixtures},pdfauthor={KEIO},pdfsubject={Integrated research manuscript consolidating eleven source papers}}
\newtheorem{theorem}{Theorem}[section]
\newtheorem{lemma}[theorem]{Lemma}
\newtheorem{proposition}[theorem]{Proposition}
\newtheorem{corollary}[theorem]{Corollary}
\theoremstyle{definition}
\newtheorem{definition}[theorem]{Definition}
\newtheorem{assumption}[theorem]{Assumption}
\theoremstyle{remark}
\newtheorem{remark}[theorem]{Remark}
\newtheorem{example}[theorem]{Example}
\DeclareMathOperator{\Var}{Var}
\DeclareMathOperator{\Cov}{Cov}
\DeclareMathOperator{\supp}{supp}
\DeclareMathOperator{\diag}{diag}
\DeclareMathOperator{\rank}{rank}
\DeclareMathOperator{\tr}{tr}
\DeclareMathOperator{\dist}{dist}
\DeclareMathOperator{\sinc}{sinc}
\DeclareMathOperator{\TV}{TV}
\newcommand{\T}{\mathbb T}
\newcommand{\R}{\mathbb R}
\newcommand{\C}{\mathbb C}
\newcommand{\N}{\mathbb N}
\newcommand{\Z}{\mathbb Z}
\newcommand{\E}{\mathbb E}
\newcommand{\PP}{\mathbb P}
\newcommand{\Pp}{\mathbb P}
\newcommand{\eps}{\varepsilon}
\newcommand{\dd}{\,\mathrm d}
\newcommand{\norm}[1]{\left\lVert#1\right\rVert}
\newcommand{\Ree}{\operatorname{Re}}
\newcommand{\Imm}{\operatorname{Im}}
\renewcommand{\Re}{\operatorname{Re}}
\renewcommand{\Im}{\operatorname{Im}}
\newcommand{\1}{\mathbf 1}
\pagestyle{fancy}
\fancyhf{}
\fancyhead[L]{\small Observation and Recovery in Arithmetic and Probability}
\fancyhead[R]{\small KEIO}
\fancyfoot[C]{\thepage}
\setcounter{tocdepth}{1}
\title{Observation and Recovery\\in Arithmetic and Probability\\[5pt]
\large Prime Flows, Rare Events, and Infinite Mixtures}
\author{KEIO}
\date{14 October 2026\\\small Integrated research manuscript\\
\small Prepared with AI assistance; not peer reviewed}
\begin{document}
\maketitle
\begin{abstract}
We study what arithmetic and probabilistic observations reveal about
hidden infinite-dimensional structure. Quantitative recurrence of weighted
prime flows leads to finite-height boundary zeros and sharp Euler-edge
locking. Weighted Weil forms localize simple critical-line zeros in prime
phase regions, including shrinking targets, and connect selected phases
to conditional Gibbs profiles and joint edge laws. In weighted random
series, macroscopic joint rare events recover the microscopic law, while
microscopic fluctuations recover a discrete coefficient spectrum.
Controlled thinning and interior probes then identify unknown dimensions,
power ratios, and correlated laws, with explicit finite-data obstructions
and stability estimates. For infinite mixtures we give sparse
partial-sum identification criteria, construct analytic families with
arbitrary replication thresholds despite independence of all finite
parameter jets, and show that a finite-dimensional one-view defect
becomes supported on finitely many atoms after two observations.
A localized multiplicative-smoothing argument
proves the prescribed-moment assertion of B\"urgin, Epperlein, and Wirth
and yields an exact classification of integer moment coincidence sets
on positive compact intervals, realizable by strictly log-concave
Gevrey densities. The final results separate
exact identification from uniform statistical recovery through sharp
regularity boundaries, logarithmic rates, and arbitrarily slow examples.
This manuscript consolidates eleven source papers into five connected
parts, retaining their distinct proofs while replacing repeated arguments
by internal references. Conditional inputs, classical ingredients, and
remaining quantitative and priority questions are stated explicitly.
\end{abstract}
\clearpage
\section*{Introduction and organization}
\addcontentsline{toc}{section}{Introduction and organization}
\label{sec:integrated-introduction}

An observation can retain a surprisingly detailed description of an
infinite system, or conceal an entire family of alternatives. The
distinction depends on more than the smoothness of the observation
map. It depends on which coordinates are tested, how several
observations are coupled, and which limiting operations are allowed.
This manuscript develops that issue in arithmetic dynamics, rare-event
probability, and mixture identification.

The arithmetic starting point is a weighted metric on an infinite
prime torus. A small metric ball imposes simultaneous phase coherence
at many primes. Its probability controls an entropy scale, while the
actual logarithms of primes introduce an additional arithmetic cost.
Near the extremal edges of Euler products, analytic defects become
positive weighted phase energies, allowing recurrence estimates to
control the location and spacing of zeros and critical points.
Weil-form localization provides a different observation mechanism:
it selects critical-line zeros with prescribed phase constraints.

The probabilistic middle of the manuscript asks what is retained by
the resulting continuum and conditional limits. Joint rare-event
surfaces admit Mellin inversions that recover a local law, whereas
centered fluctuation series retain a discrete spectrum. Once the law
and the observation architecture are both unknown, extra channels,
controlled Bernoulli gates, and interior probes give exact recovery
criteria. Counterexamples show why finite marginal summaries,
resonances, and flat endpoint jets can fail.

The final part places these observations in the language of mixing
operators. If $P_\theta$ is a known component family and $\mu$ is an
unknown mixing law, then
\[
 A_q\mu=\int P_\theta^{\otimes q}\,\mu(d\theta)
\]
describes $q$ observations sharing a latent parameter. Products of
observable functions enlarge the collection of accessible moments.
This makes the number of repeated observations an algebraic property
of the observable function space. Sparse moments, finite derivative
tests, and closure of model classes then explain why finite mixtures,
arbitrary mixing laws, and uniformly estimable subclasses can have
different thresholds.

\paragraph{Five parts.}
Part~\ref{part:prime-flows} develops quantitative prime-flow recurrence,
boundary geometry, and the difference between entropy and arithmetic
time scales. Part~\ref{part:zeta-localization} establishes phase
localization at zeta zeros and develops shrinking targets, simultaneous
edge approximation, and conditional phase laws. Part~\ref{part:rare-events}
derives inverses and fluctuation laws for weighted series and analyzes
minimal multichannel observation. Part~\ref{part:controlled-observation}
uses thinning and interior probes for blind identification, including
finite noisy observation guarantees and exact obstructions.
Part~\ref{part:mixtures} treats partial-sum marginals, finite differential
barriers, replication thresholds, prescribed moment coincidences,
and statistical recovery.

\paragraph{Three kinds of conclusion.}
Throughout, exact recovery from an entire response function is
distinguished from stable recovery from noisy finite data. A theorem
about recurrence on a prime torus is also distinguished from its
arithmetic transfer to a particular Euler product. Finally, an explicit
construction with strong regularity is an existence theorem; it need
not supply an efficient numerical representation of the constructed
object. The hypotheses attached to each transfer and inverse theorem
are part of its statement.

\paragraph{Scope of the consolidation.}
The eleven source manuscripts are listed in
Section~\ref{sec:source-correspondence}. Their substantive results and
proof arguments are incorporated here. Where a later manuscript proves
a stronger version of an earlier result, the integrated exposition uses
the stronger proof and records the earlier specialization or its
reference. Questions resolved by a later part are no longer listed as
open. Independent results that use similar tools but concern different
objects remain separate.

The integration unifies notation, numbering, references, and dependencies.
It does not turn an explicitly conditional statement into an unconditional
one, and it does not establish historical priority or publication-level
significance. The source drafts contained several independent proof
checks; the present editorial pass adds coverage and cross-reference
checks rather than a new external peer review of every assertion.

\paragraph{Notation.}
The symbols $\R,\C,\N,\Z$ denote the usual number systems and $\T$
denotes a circle group, with its period specified when first used in a
part. The shift between period $1$ and period $2\pi$ is kept explicit.
For a signed measure, $\|\nu\|_{\rm TV}$ denotes its total variation
norm. For probabilities, $\operatorname{TV}(P,Q)=\sup_B|P(B)-Q(B)|$
is half the variation norm of $P-Q$. Wasserstein metrics refer to
the underlying parameter metric specified in the relevant section.
The symbols $A,E$ denote the two Euler-edge locations in the arithmetic
parts and positive power exponents in the weighted-series parts;
those sections state the intended parameters explicitly. Likewise
$N$ can denote an observation count or a spectral counting function
with an argument. No identification of these different uses is assumed.

\clearpage
\tableofcontents
\clearpage

% ===== Integrated Part 1 =====
\part{Prime flows and quantitative recurrence}
\label{part:prime-flows}

The first part develops a common quantitative framework for weighted
prime-phase flows and analytic boundary points. The geometry of a small
phase ball determines its Haar probability; the arithmetic of the prime
frequencies determines the additional cost of reaching it uniformly.
A quadratic real-part estimate then transports phase recurrence into
zeros or critical points near a vertical boundary. These three steps
will also supply the phase geometry used in the subsequent parts.

The qualitative prime-zeta boundary and almost-periodic transport of
zeros are classical \cite{SV,BCS,Farag,KR,Baroni}.
For Euler products, the right boundaries of $\zeta(s)=1$ and
$\zeta'(s)=0$, and their qualitative correspondence, are due to
Arias de Reyna and van de Lune \cite{AriasLune}. The quantitative
claims developed here concern shrinking targets, uniform maximal gaps,
and sharp local moduli between these two boundaries. The phase-hull
viewpoint motivating the prime-zeta construction is also present in
\cite{Keio004}; none of its density or correlation asymptotics is
assumed in the proofs below.

Finite-dimensional small-divisor covering estimates and
infinite-dimensional product-measure nonresonance arguments are
established methods \cite{Bounemoura,DumasFischler,Bourgain}.
Here their growing-dimensional costs are tracked explicitly. The
contrast between probability-one recurrence and residual slowness is
related to classical exceptional-rotation phenomena
\cite{GalatoloPeterlongo}. The Fej\'er coefficient inequality used in
the localization optimization is likewise classical
\cite{FejerMulti}; its weighted prime-budget consequences are proved
in Section~\ref{p02:sec:barrier}.

The results in this part concern the half-plane of absolute convergence.
They do not assert the Riemann hypothesis or a deterministic lower
first-waiting bound at the origin for ordinary prime zeta or ordinary
$\zeta'$. Maximal gaps, recurrence of long empty intervals, and
almost-sure first waiting times have different quantifiers throughout.

\section{Weighted prime flows and a finite Fourier certificate}
\label{p02:sec:foundations}

Let $\mathcal P$ be the set of primes, with finitely many primes possibly
removed, and let $\Omega=\T^{\mathcal P}$ carry normalized product Haar
measure $\mu$. Here $\T=\{z\in\C:|z|=1\}$, parametrized by
$z=e^{i\theta}$ with $\theta$ of period $2\pi$.
For positive summable weights $w=(w_p)$, put
\begin{equation}\label{p02:eq:energy}
 \mathcal E_w(\eta)=\sum_{p\in\mathcal P}w_p(1-\Re\eta_p),
 \qquad d_w(\eta,\zeta)^2=\sum_p w_p|\eta_p-\zeta_p|^2.
\end{equation}
Thus $d_w(\eta,1)^2=2\mathcal E_w(\eta)$. The metric $d_w$ induces
the product topology. For a real frequency vector $\lambda$, write
$\Phi_t\eta=(\eta_p e^{-it\lambda_p})_p$. This flow is an isometry.
We use the uniform covering time
\begin{equation}\label{p02:eq:covering}
 \mathcal R_{w,\lambda}(\eps)=\inf\{T>0:
 \text{for every }\eta\in\Omega\text{ there is }t\in[0,T]
 \text{ with }\mathcal E_w(\Phi_t\eta)<\eps\}.
\end{equation}
Translation in the group makes this the same as covering every target,
from every initial phase, with squared-distance tolerance $2\eps$.
All estimates below are also uniform in the initial time of an interval.
The fixed-origin waiting time is
$\tau_\eta(\eps)=\inf\{t\ge0:\mathcal E_w(\Phi_t\eta)<\eps\}$.

Our main weight class is
\begin{equation}\label{p02:eq:weights}
 c_0p^{-\sigma}(\log p)^\kappa\le w_p\le
 C_0p^{-\sigma}(\log p)^\kappa,
 \qquad \sigma>1,\quad\kappa\in\R.
\end{equation}
Throughout, implicit constants may depend on these parameters and on a
fixed finite set of omitted primes. Set
\begin{equation}\label{p02:eq:scales}
 \beta=\frac1{\sigma-1},\qquad \ell=\log(1/\eps),\qquad
 X_{\sigma,\kappa}(\eps)=\eps^{-\beta}\ell^{(\kappa-1)\beta},
 \qquad Y_{\sigma,\kappa}(\eps)=\frac{X_{\sigma,\kappa}(\eps)}{\ell}.
\end{equation}
We restrict to $0<\eps<e^{-1}$. Multiplying $\eps$ by a fixed positive
constant changes $X$ and $Y$ by bounded factors. If $x\asymp X(\eps)$,
then
\begin{equation}\label{p02:eq:cutoff-scale}
 x^{1-\sigma}(\log x)^{\kappa-1}\asymp\eps,
 \qquad x/\log x\asymp Y(\eps).
\end{equation}

\begin{lemma}[Prime sums and logarithmic budgets]\label{p01:lem:primes}
For $a<1$, $u>1$, and real $b$, as $x\to\infty$,
\begin{align}
 \sum_{p\le x}p^{-a}(\log p)^b
 &\asymp x^{1-a}(\log x)^{b-1},\label{p02:eq:prime-sum}\\
 \sum_{p>x}p^{-u}(\log p)^b
 &\asymp x^{1-u}(\log x)^{b-1}.\label{p02:eq:prime-tail}
\end{align}
In particular, for fixed $u>1$ and $0\le v<1$,
\begin{align}
 \sum_{p>x}p^{-u}&\asymp_u\frac{x^{1-u}}{\log x},
 &\sum_{3\le p\le x}p^{-v}&\asymp_v\frac{x^{1-v}}{\log x},
 \label{p01:eq:primesums}\\
 \sum_{3\le p\le x}p^{-v}\log p&\asymp_v x^{1-v},
 &\#\{p:x\le p\le2x\}&\asymp\frac{x}{\log x}.
 \label{p01:eq:primeblocks}
\end{align}
Moreover,
\begin{align}
 \sum_{p\le x}\log(x/p)&=O(x/\log x),
 \label{p02:eq:log-budget}\\
 \sum_{p\le x}\log(\log x/\log p)&=O(x/(\log x)^2).
 \label{p02:eq:loglog-budget}
\end{align}
Deleting finitely many primes changes none of these orders.
\end{lemma}
\begin{proof}
Partial summation from $\pi(x)\sim x/\log x$ and
$\sum_{p\le x}\log p\sim x$ gives the power-weight estimates and
$\pi(2x)-\pi(x)\asymp x/\log x$. For the logarithmic budgets,
\[
 \sum_{p\le x}\log(x/p)=\int_2^x\frac{\pi(t)}{t}\,dt,
 \qquad
 \sum_{p\le x}\log(\log x/\log p)
 =\int_2^x\frac{\pi(t)}{t\log t}\,dt.
\]
The prime number theorem gives the stated bounds. These are the only
asymptotic facts about prime distribution used in the recurrence
arguments; see \cite{MontgomeryVaughan}.
\end{proof}

For a finite trigonometric polynomial $P=\sum_k\widehat P(k)\eta^k$,
write $\|P\|_A=\sum_k|\widehat P(k)|$. Finite multi-indices are understood
throughout. Direct integration of each character gives
\begin{equation}\label{p02:eq:fourier-average}
 \left|\frac1T\int_0^T P(\Phi_t\eta)\,dt-\int P\,d\mu\right|
 \le\frac{2\|P\|_A}{T\Delta},
\end{equation}
provided every nonconstant frequency occurring in $P$ satisfies
$|k\cdot\lambda|\ge\Delta>0$. This bound is uniform in $\eta$.

\begin{lemma}[Positive Fourier certificate]\label{p02:lem:certificate}
Let $Q\subset\mathcal P$ be finite, and let $K\ge0$ be a trigonometric
polynomial depending only on $Q$,
with $\int K\,d\mu=1$. Write
\[
 \mathcal E_Q=\sum_{p\in Q}w_p(1-\Re\eta_p),\qquad
 m=\int\mathcal E_Q K\,d\mu,\qquad W=\sum_p w_p.
\]
Choose $a>m$, and suppose the nonzero frequencies of
$P=(a-\mathcal E_Q)K$ have absolute value at least $\Delta$.
If
\begin{equation}\label{p02:eq:certificate-time}
 T>\frac{2(a+2W)\|K\|_A}{\Delta(a-m)},
\end{equation}
then every orbit has a time in $[0,T]$ at which
$\mathcal E_w<a+2\sum_{p\notin Q}w_p$.
\end{lemma}
\begin{proof}
The mean of $P$ is $a-m>0$, and
$\|P\|_A\le(a+2W)\|K\|_A$. Equation
\eqref{p02:eq:fourier-average} makes its time average positive. At a time
where $P>0$, nonnegativity of $K$ forces $\mathcal E_Q<a$.
The remaining coordinates contribute at most twice their total weight.
\end{proof}

Haar integration on one circle is normalized to total mass one. For an
integer $m\ge1$, define
\begin{equation}\label{p01:eq:kernel}
 F_m(\theta)=\frac1m\left|\sum_{j=0}^{m-1}e^{ij\theta}\right|^2,
 \quad A_m=\frac{2m^2+1}{3m},\quad J_m=F_m^2/A_m.
\end{equation}
Equivalently writing $J_m(e^{i\theta})=J_m(\theta)$, we use the Fourier coefficient norm
$\norm{Q}_{\mathcal A}=\sum_k|\widehat Q(k)|$ for a trigonometric
polynomial; the analogous definition applies to several variables. Thus $\norm{Q}_{\mathcal A}=\|Q\|_A$.

\begin{lemma}[Squared Fej\'er identities]\label{p01:lem:jackson}
The polynomial $J_m$ is nonnegative, has degree $2m-2$, and has Haar
mean one. All its Fourier coefficients are nonnegative. Moreover,
\begin{align}
 \int(1-\cos\theta)J_m(\theta)\dd\theta
 &=\frac3{2m^2+1},\label{p01:eq:kernelenergy}\\
 \norm{J_m}_{\mathcal A}=J_m(0)
 &=\frac{3m^3}{2m^2+1}\le\frac32m.
 \label{p01:eq:kernelnorm}
\end{align}
Here and below the normalized Haar factor is included in
$\dd\theta$ when integration is over a circle.
\end{lemma}
\begin{proof}
Write
\[
 F_m(\theta)=\sum_{k\in\Z}a_ke^{ik\theta},\qquad
 a_k=(1-|k|/m)_+.
\]
Convolution shows positivity of every Fourier coefficient of $F_m^2$.
Its constant coefficient is
\[
 \sum_k a_k^2=1+2\sum_{k=1}^{m-1}(1-k/m)^2
 =\frac{2m^2+1}{3m}=A_m.
\]
If $b_0,b_1$ are the zeroth and first coefficients of $F_m^2$, then
\[
 b_0-b_1=\frac12\sum_{k\in\Z}(a_k-a_{k-1})^2=\frac1m.
\]
There are $2m$ nonzero differences, each of magnitude $1/m$.
Division by $A_m$ proves \eqref{p01:eq:kernelenergy}. Nonnegative Fourier
coefficients make the coefficient norm equal to the value at zero.
Since $F_m(0)=m$, this proves \eqref{p01:eq:kernelnorm}.
\end{proof}

The gain $m^{-2}$ in \eqref{p01:eq:kernelenergy} is the quantitative reason
for squaring the kernel. An unweighted, equal-precision treatment of all
prime phases loses the exponent needed below.

For $K=\prod_{p\in Q}J_{m_p}$, its localized mean is consequently
\begin{equation}\label{p02:eq:kernel-mean}
 m=\sum_{p\in Q}\frac{3w_p}{2m_p^2+1}.
\end{equation}

\section{Small balls, arithmetic recurrence, and recurrent holes}
\label{p02:sec:recurrence}

\begin{theorem}[Small-ball scale]\label{p02:thm:smallball}
Under \eqref{p02:eq:weights}, for every $\sigma>1$ and $\kappa\in\R$,
there are positive constants $c,C$ such that
\begin{equation}\label{p02:eq:smallball}
 e^{-CY(\eps)}\le\mu\{\mathcal E_w\le\eps\}
 \le e^{-cY(\eps)}
\end{equation}
for all sufficiently small $\eps$.
If $L=\sum_p w_p|\lambda_p|<\infty$, then also
\begin{equation}\label{p02:eq:time-smallball}
 \mu\{\min_{0\le t\le T}\mathcal E_w(\Phi_t\eta)\le\eps\}
 \le C(1+T/\eps)e^{-cY(\eps)}.
\end{equation}
In particular $\mathcal R_{w,\lambda}(\eps)\ge e^{cY(\eps)}$.
\end{theorem}
\begin{proof}
For the upper small-ball bound take $x=bX(\eps)$ with $b>0$ a
sufficiently small fixed constant. The block $x\le p\le2x$ contains
$N\asymp x/\log x$ primes, each of weight comparable to
$v=x^{-\sigma}(\log x)^\kappa$. As $Nv\asymp
b^{1-\sigma}\eps$, the event $\mathcal E_w\le\eps$ forces the
average of $N$ independent copies of $U=1-\cos\Theta$ to be at most
a fixed $q<1=\mathbb E U$. For small fixed $s>0$,
$e^{sq}\mathbb E e^{-sU}<1$. Markov's inequality thus bounds the
probability by $e^{-cN}$.

For the lower bound take $x=BX(\eps)$, with $B$ large enough that
$2\sum_{p>x}w_p\le\eps/2$. On the coordinates $p\le x$, require
\[
 |\arg\eta_p|\le
 \min\left\{1,b_0(p/x)^{\sigma/2}
                   (\log x/\log p)^{\kappa/2}\right\}.
\]
Each of their energy contributions is at most
$C b_0^2x^{-\sigma}(\log x)^\kappa$. Their sum is at most
$\eps/2$ after fixing $b_0$ small. The negative logarithm of the
probability of this box is bounded by a constant times
\[
 \pi(x)+\sum_{p\le x}\log(x/p)
             +\sum_{p\le x}\log(\log x/\log p)=O(x/\log x).
\]
This proves \eqref{p02:eq:smallball}.

The energy along any orbit is $L$-Lipschitz. A mesh of spacing
$\eps/(L+1)$ converts a hit below $\eps$ into a mesh value below
$2\eps$. Haar invariance, the union bound, and
\eqref{p02:eq:smallball} give \eqref{p02:eq:time-smallball}. Taking
$T=e^{c_1Y(\eps)}$ with $c_1$ small makes the last probability less
than one. Here $Y(\eps)/\log(1/\eps)\to\infty$, so the polynomial
mesh factor is harmless.
\end{proof}

There is also a metric interpretation of $Y$. If $N_w(r)$ denotes the
least number of $d_w$-balls of radius $r$ covering $\Omega$, then
\begin{equation}\label{p02:eq:metric-entropy}
 \log N_w(\sqrt\eps)\asymp Y(\eps).
\end{equation}
The lower estimate follows by measuring a cover. For the upper estimate,
take a maximal $\sqrt\eps$-separated set. Its half-radius balls are
disjoint, each of the same Haar measure, and the full-radius balls
cover. Compactness makes this a finite argument.

\begin{theorem}[Arithmetic prime frequencies]\label{p02:thm:arithmetic-flow}
Let $\lambda_p=\log p$. Under \eqref{p02:eq:weights} with $1<\sigma<3$,
\begin{equation}\label{p02:eq:arithmetic-flow}
 \exp(cY_{\sigma,\kappa}(\eps))
 \le\mathcal R_{w,\lambda}(\eps)
 \le\exp(CX_{\sigma,\kappa}(\eps)).
\end{equation}
The same conclusion holds for $\lambda_p=\log(p/p_0)$ with a fixed
prime $p_0$ omitted from $\mathcal P$.
\end{theorem}
\begin{proof}
Unique factorization shows that each nonzero finite integer vector $k$
has $\sum k_p\log p=\log(A/B)\ne0$, where $A,B$ are positive integers
at most $\exp(\sum|k_p|\log p)$. Integer separation gives
\begin{equation}\label{p02:eq:arithmetic-separation}
 \left|\sum k_p\log p\right|
 \ge \exp\left(-\sum|k_p|\log p\right).
\end{equation}
For relative frequencies the denominator cost is at most a fixed
multiple of this cost, since the additional prime is $p_0$.

Take $Q=\{p\le x\}$ and
\[
 m_p=\left\lceil\max\left\{1,
 (x/p)^{\sigma/3}(\log p/\log x)^{\kappa/3}\right\}\right\rceil.
\]
Equations \eqref{p02:eq:prime-sum}--\eqref{p02:eq:kernel-mean} give
\begin{align*}
 m&\le Cx^{-2\sigma/3}(\log x)^{2\kappa/3}
       \sum_{p\le x}p^{-\sigma/3}(\log p)^{\kappa/3}
       \le Cx^{1-\sigma}(\log x)^{\kappa-1},\\
 B_x:=\sum_{p\le x}(2m_p-1)\log p&\le Cx,\qquad
 \log\|K\|_A\le Cx/\log x.
\end{align*}
The estimate for $B_x$ uses $\sigma/3<1$; the norm estimate uses
\eqref{p02:eq:log-budget} and \eqref{p02:eq:loglog-budget}.
Choose $x$ a sufficiently large constant multiple of $X(\eps)$,
so that $m\le\eps/4$ and $2\sum_{p>x}w_p\le\eps/4$.
The polynomial $(\eps/2-\mathcal E_Q)K$ has coordinate degrees
at most $2m_p-1$. Equations \eqref{p02:eq:arithmetic-separation} and
\eqref{p02:eq:certificate-time} give a covering time at most $e^{Cx}$.
The lower estimate is Theorem~\ref{p02:thm:smallball}, since
$\sum w_p|\lambda_p|<\infty$.
\end{proof}

\begin{proposition}[Robust holes and first waiting times]\label{p02:prop:holes}
Suppose \eqref{p02:eq:weights} holds, $\sum w_p|\lambda_p|<\infty$, and
$\mathcal R_{w,\lambda}(\eps)\le U(\eps)<\infty$.
There are constants $c,b>0$ with the following property. For every
sufficiently small $\eps$, every orbit, and every $A\in\R$, a time
$t_0\in[A,A+2U(b\eps)]$ begins an interval of length
$e^{cY(\eps)}$ on which $\mathcal E_w(\Phi_t\eta)>\eps$.
Moreover, for Haar-almost every $\eta$,
\begin{equation}\label{p02:eq:waiting-lower}
 \tau_\eta(\eps)\ge e^{cY(\eps)}
 \quad\hbox{for all sufficiently small }\eps.
\end{equation}
If $U(\eps)=e^{CX(\eps)}$, the double-logarithmic waiting-time
exponent is $\beta$ almost surely. If $U(\eps)=e^{CY(\eps)}$,
then $\log\tau_\eta(\eps)\asymp Y(\eps)$ almost surely.
\end{proposition}
\begin{proof}
Apply \eqref{p02:eq:time-smallball} at threshold $16\eps$ and time
$H=e^{cY(\eps)}$ to obtain a phase $\eta^*$ satisfying
$\mathcal E_w(\Phi_t\eta^*)>16\eps$ throughout $[0,H]$.
Every phase within $d_w$-distance $\sqrt\eps$ of $\eta^*$ still
has energy greater than $\eps$ on that interval: use the triangle
inequality and the isometry property. Covering at tolerance $b\eps$
brings any prescribed orbit into that ball in time $2U(b\eps)$,
also after any initial time $A$. This proves the robust-hole assertion.

For the almost-sure assertion, take $\eps_n=e^{-n}$. With
$T_n=e^{cY(\eps_n)}$ and $c$ sufficiently small, the probabilities
in \eqref{p02:eq:time-smallball} are summable. Borel--Cantelli excludes
a hit before $T_n$ for all sufficiently large $n$, almost surely.
Monotonicity of $\tau$ and comparability of $Y(\eps_n)$ at consecutive
indices give \eqref{p02:eq:waiting-lower} for all small $\eps$ after
reducing $c$. The asserted upper bounds follow from uniform covering.
\end{proof}

The power-weight specialization and its explicit certificate are given next.
The endpoint $\sigma=3$ is resolved in Section~\ref{p02:sec:barrier};
the frequency-perturbation theorem in Section~\ref{p02:sec:generic}
has no endpoint restriction.

\section{Power weights and an explicit arithmetic certificate}\label{sec:power-certificate}

Specialize the preceding framework to the relative frequencies and write
\[
 \Omega=\prod_{p\ge3}\T,\qquad
 \T=\{z\in\C:|z|=1\},\qquad
 (\Phi_t\eta)_p=\eta_p e^{-it\lambda_p},\quad
 \lambda_p=\log(p/2).
\]
Equip $\Omega$ with product Haar probability $\mu$. Fix $\sigma>1$ and
positive weights $w=(w_p)_{p\ge3}$ satisfying
\begin{equation}\label{p01:eq:weights}
 c_0p^{-\sigma}\le w_p\le C_0p^{-\sigma}
 \qquad(p\ge3)
\end{equation}
for two fixed positive constants. Define
\begin{equation}\label{p01:eq:energy}
 E_w(\eta)=\sum_{p\ge3}w_p(1-\Ree\eta_p),\qquad
 d_w(\eta,\xi)^2=\sum_{p\ge3}w_p|\eta_p-\xi_p|^2.
\end{equation}
Here $E_w=\mathcal E_w$ in the notation of \eqref{p02:eq:energy}.

The covering time is
\begin{equation}\label{p01:eq:coverdef}
 \mathcal R_w(\eps)=\inf\left\{T>0:
  \text{for all }\eta,\xi\in\Omega\text{ there is }t\in[0,T]
  \text{ with }d_w(\Phi_t\eta,\xi)^2\le2\eps\right\}.
\end{equation}
Multiplication by $\bar\xi$ reduces the target to $1$. All logarithms
are natural. For $0<\eps<e^{-1}$ put
\begin{equation}\label{p01:eq:scales}
 \beta_\sigma=\frac1{\sigma-1},\qquad
 \ell_\eps=\log(1/\eps),\qquad
 X_\sigma(\eps)=\eps^{-\beta_\sigma}\ell_\eps^{-\beta_\sigma},\qquad
 Y_\sigma(\eps)=\frac{X_\sigma(\eps)}{\ell_\eps}.
\end{equation}
Constants may depend on the fixed weights through
$\sigma,c_0,C_0$, but never on the phases or the initial time.

\begin{theorem}[Weighted recurrence and small balls]\label{p01:thm:flow}
Assume \eqref{p01:eq:weights} and $1<\sigma<3$. There are constants
$c,C,\eps_0>0$ such that, for $0<\eps<\eps_0$,
\begin{align}
 e^{cY_\sigma(\eps)}
 &\le\mathcal R_w(\eps)\le e^{CX_\sigma(\eps)},
 \label{p01:eq:flowbounds}\\
 e^{-CY_\sigma(\eps)}
 &\le\mu\{E_w\le\eps\}\le e^{-cY_\sigma(\eps)}.
 \label{p01:eq:smallball}
\end{align}
The upper recurrence bound holds on every interval of the stated
length. In particular,
\begin{equation}\label{p01:eq:flowexponent}
 \lim_{\eps\downarrow0}
 \frac{\log\log\mathcal R_w(\eps)}{\log(1/\eps)}
 =\frac1{\sigma-1}.
\end{equation}
Moreover, if
$\tau_\eta(\eps)=\inf\{t\ge0:E_w(\Phi_t\eta)\le\eps\}$, then for
$\mu$-almost every $\eta$,
\begin{equation}\label{p01:eq:flowwaiting}
 \lim_{\eps\downarrow0}
 \frac{\log\log(e^e+\tau_\eta(\eps))}{\log(1/\eps)}
 =\frac1{\sigma-1}.
\end{equation}
\end{theorem}

The theorem is the $\kappa=0$ specialization of the preceding general
results. Its exponent is exact, while its upper and lower logarithms
differ by a factor $\log(1/\eps)$. Section~\ref{p02:sec:barrier}
extends the exact double-logarithmic exponent to $\sigma=3$.

\begin{lemma}[Arithmetic separation of frequencies]\label{p01:lem:separation}
Let $\mathcal Q$ be a finite set of primes at least $3$, let $d_p\ge1$
be integers, and put $B=\sum_{p\in\mathcal Q}d_p\log p$. If
$k\ne0$ and $|k_p|\le d_p$, then
\begin{equation}\label{p01:eq:separation}
 \left|\sum_{p\in\mathcal Q}k_p\log(p/2)\right|\ge e^{-2B}.
\end{equation}
Consequently, for any trigonometric polynomial $Q$ with these coordinate
degrees, any initial phase $\eta$, and any real interval $I$ of length $T$,
\begin{equation}\label{p01:eq:fourieraverage}
 \left|\frac1T\int_IQ(\Phi_t\eta)\dd t-\int Q\dd\mu\right|
 \le\frac{2e^{2B}}T\norm{Q}_{\mathcal A}.
\end{equation}
\end{lemma}
\begin{proof}
Exponentiating the frequency gives a positive rational $A/D$ in lowest
terms. Unique factorization at every prime $p\ge3$ implies $A\ne D$.
Both integers are at most
\[
 \exp\left(\sum|k_p|\log p+\left|\sum k_p\right|\log2\right)
 \le e^{2B}.
\]
For distinct positive integers,
$|\log A-\log D|\ge1/\max(A,D)$, proving \eqref{p01:eq:separation}.
The average of $e^{-it\nu}$ on an interval of length $T$ has magnitude
at most $2/(T|\nu|)$. Expand $Q$ and apply this to each nonconstant
coefficient. The phase and the initial endpoint have modulus one and
do not change the bound.
\end{proof}

\begin{proposition}[Finite certificate]\label{p01:prop:certificate}
Fix a nonempty finite set $\mathcal Q$, positive $w_p$, and integers $m_p\ge1$.
Put
\begin{align*}
 W_{\mathcal Q}&=\sum_{p\in\mathcal Q}w_p,&
 M&=\sum_{p\in\mathcal Q}\frac{3w_p}{2m_p^2+1},\\
 B&=\sum_{p\in\mathcal Q}(2m_p-1)\log p,&
 K_0&=\prod_{p\in\mathcal Q}\frac{3m_p^3}{2m_p^2+1}.
\end{align*}
Every interval of length
\begin{equation}\label{p01:eq:explicitT}
 T\ge 8e^{2B}\frac{(M+W_{\mathcal Q})K_0}{M}
\end{equation}
contains a $t$ for which
\begin{equation}\label{p01:eq:finitehit}
 \sum_{p\in\mathcal Q}w_p(1-\Ree(\Phi_t\eta)_p)<2M.
\end{equation}
The estimate is uniform in $\eta$ and remains true for any target phase
after translation.
\end{proposition}
\begin{proof}
Let $E_{\mathcal Q}$ denote the sum in \eqref{p01:eq:finitehit}, and set
$K=\prod_{p\in\mathcal Q}J_{m_p}$. Lemma~\ref{p01:lem:jackson} gives
$\int K=1$, $\int E_{\mathcal Q}K=M$, and
$\norm K_{\mathcal A}=K_0$. The real polynomial
\[
 Q=(2M-E_{\mathcal Q})K
\]
has mean $M>0$, coordinate degree at most $2m_p-1$, and
\[
 \norm Q_{\mathcal A}
 \le(2M+2W_{\mathcal Q})K_0.
\]
By \eqref{p01:eq:fourieraverage} and \eqref{p01:eq:explicitT}, its average on
the interval is at least $M/2$. Hence $Q$ is positive somewhere. Since
$K\ge0$, positivity implies \eqref{p01:eq:finitehit}.
\end{proof}

\begin{proof}[Proof of the upper bound in Theorem~\ref{p01:thm:flow}]
Take $\mathcal Q=\{3\le p\le x\}$ and allocate the degrees by
\begin{equation}\label{p01:eq:degrees}
 m_p=\left\lceil(x/p)^{\sigma/3}\right\rceil.
\end{equation}
Because $\sigma/3<1$, Lemma~\ref{p01:lem:primes} gives
\begin{align}
 M&\le\frac32C_0x^{-2\sigma/3}
                  \sum_{p\le x}p^{-\sigma/3}
       \ll\frac{x^{1-\sigma}}{\log x},\label{p01:eq:Mbound}\\
 B&\ll x^{\sigma/3}\sum_{p\le x}p^{-\sigma/3}\log p
       \ll x,\label{p01:eq:Bbound}\\
 \log K_0&\le\sum_{p\le x}\log(3m_p/2)
       \ll\pi(x)+\sum_{p\le x}\log(x/p)
       \ll\frac{x}{\log x}.\label{p01:eq:Kbound}
\end{align}
The single term $p=3$ also gives $M\gg x^{-2\sigma/3}$, so
$\log(1/M)=O(\log x)$. Thus \eqref{p01:eq:explicitT} is at most
$\exp(Cx)$ for large $x$. At the resulting time, the unprescribed tail
has the deterministic bound
\[
 E_w(\Phi_t\eta)
 \le 2M+2\sum_{p>x}w_p
 \ll\frac{x^{1-\sigma}}{\log x}.
\]
Choose $x=KX_\sigma(\eps)$ with a sufficiently large fixed $K$.
The scale identity
\begin{equation}\label{p01:eq:scaleid}
 \frac{X_\sigma(\eps)^{1-\sigma}}
      {\log X_\sigma(\eps)}\asymp_\sigma\eps
\end{equation}
shows that the last bound is at most $\eps$. This proves the upper
bound on every initial interval, for every phase and target.
\end{proof}

\begin{remark}[Why the degrees are anisotropic]
Ignoring integer rounding and slow logarithmic factors, minimizing
$\sum m_p$ subject to $\sum w_p/m_p^2\le\eps$ leads to
$m_p\propto w_p^{1/3}$. Formula~\eqref{p01:eq:degrees} implements that
allocation at the truncation scale. We do not claim exact optimization
of constants or logarithmic factors.
\end{remark}

\begin{lemma}[Power-weight small balls]\label{p01:lem:smallball}
For every $\sigma>1$ under \eqref{p01:eq:weights},
\[
 e^{-CY_\sigma(\eps)}\le\mu\{E_w\le\eps\}
 \le e^{-cY_\sigma(\eps)}.
\]
\end{lemma}
\begin{proof}
This is Theorem~\ref{p02:thm:smallball} with $\kappa=0$.
Its lower bound may in this specialization use the explicit event
\begin{equation}\label{p01:eq:boxevent}
 |\arg\eta_p|\le(p/x)^{\sigma/2}\quad(3\le p\le x),
 \qquad x=KX_\sigma(\eps).
\end{equation}
Its energy is at most
$C x^{-\sigma}\pi(x)+2\sum_{p>x}w_p\ll
x^{1-\sigma}/\log x$, and its negative log probability is
$(\pi(x)-1)\log\pi+(\sigma/2)\sum_{3\le p\le x}\log(x/p)
=O(x/\log x)$. Choosing $K$ large gives the assertion.
\end{proof}

\begin{lemma}[Finite-time avoidance estimate]\label{p01:lem:avoidance}
For weights satisfying \eqref{p01:eq:weights}, there are $c,C>0$ such that
for all small $\eps$ and all $T\ge0$,
\begin{equation}\label{p01:eq:avoidprob}
 \mu\left\{\eta:\min_{0\le t\le T}E_w(\Phi_t\eta)\le\eps\right\}
 \le C(1+T/\eps)e^{-cY_\sigma(\eps)}.
\end{equation}
In particular, for a sufficiently small fixed $c_1>0$, some $\eta^*$
satisfies
\begin{equation}\label{p01:eq:badphase}
 E_w(\Phi_t\eta^*)>\eps
 \quad\text{for every }0\le t\le e^{c_1Y_\sigma(\eps)}.
\end{equation}
\end{lemma}
\begin{proof}
The phase energy is uniformly Lipschitz:
\begin{equation}\label{p01:eq:lipschitz}
 \left|\frac{\dd}{\dd t}E_w(\Phi_t\eta)\right|
 \le L_w:=\sum_{p\ge3}w_p\lambda_p<\infty.
\end{equation}
Apply Theorem~\ref{p02:thm:smallball} with $\kappa=0$ and relative
prime frequencies. Its time-mesh bound is exactly
\eqref{p01:eq:avoidprob}; its nonempty complement at
$T=\exp(c_1Y_\sigma(\eps))$ gives \eqref{p01:eq:badphase}.
No independence between different times is used.
\end{proof}

\begin{proof}[Completion of Theorem~\ref{p01:thm:flow}]
The bad phase in Lemma~\ref{p01:lem:avoidance} proves the lower bound for
$\mathcal R_w$. The small-ball bounds are Theorem~\ref{p02:thm:smallball} with $\kappa=0$. Taking
two logarithms in \eqref{p01:eq:flowbounds} yields
\[
 \beta_\sigma\ell_\eps-(\beta_\sigma+1)\log\ell_\eps+O(1)
 \le\log\log\mathcal R_w(\eps)
 \le\beta_\sigma\ell_\eps-\beta_\sigma\log\ell_\eps+O(1),
\]
which proves \eqref{p01:eq:flowexponent}.

For the almost-sure assertion set $\eps_n=e^{-n}$ and
$T_n=\exp(c_1Y_\sigma(\eps_n))$. By \eqref{p01:eq:avoidprob}, decreasing
$c_1$ if needed,
\[
 \sum_{n=1}^\infty
 \mu\{\tau_\eta(\eps_n)\le T_n\}<\infty.
\]
The first Borel--Cantelli lemma implies
$\tau_\eta(\eps_n)>T_n$ eventually for almost every $\eta$.
The uniform upper bound applies to every $\eta$.
If $e^{-(n+1)}<\eps\le e^{-n}$, monotonicity sandwiches
$\tau_\eta(\eps)$ between its two adjacent values. Since
$(n+1)/n\to1$, the two logarithms give \eqref{p01:eq:flowwaiting}.
\end{proof}

\begin{remark}[Entropy and arithmetic cost]
The number of active primes is $\asymp X_\sigma(\eps)/\ell_\eps$.
This determines the small-ball probability. The separation bound for
prime logarithms pays the larger cost
$\sum m_p\log p=O(X_\sigma(\eps))$. The remaining logarithmic gap
between the upper and lower recurrence bounds occurs exactly here.
We have not identified the true leading order of $\log\mathcal R_w$.
\end{remark}

\section{Zeros of prime Dirichlet series in shrinking edge strips}\label{sec:prime-edge-statements}

In this section $\lambda_p=\log(p/2)$ and $\Omega=\prod_{p\ge3}\T$.
Let $a_p$ be real and satisfy
\begin{equation}\label{p01:eq:amplitudes}
 0<a_-\le a_p\le a_+<\infty,
 \qquad F_\omega(s)=\sum_p a_p\omega_pp^{-s},\quad|\omega_p|=1.
\end{equation}
The series and all its derivatives converge locally uniformly on
$\Ree s>1$. There is a unique $S>1$ satisfying
\begin{equation}\label{p01:eq:edge}
 a_2 2^{-S}=\sum_{p\ge3}a_pp^{-S}.
\end{equation}
Indeed, after division by $a_2 2^{-S}$ the right side is strictly
decreasing from infinity to zero. In the results that follow we assume
\begin{equation}\label{p01:eq:edgerange}
 1<S<3.
\end{equation}
No such restriction is needed for the elementary zero-free half-plane.

Write $\mathcal Z_\omega(\eps)$ for the set of ordinates of zeros in
\[
 \mathcal S_\eps=\{s:S-\eps<\Ree s\le S\}.
\]
Define the maximal vertical gap and the first waiting time by
\begin{align}
 G_\omega(\eps)&=\sup\{b-a:a<b,
              (a,b)\cap\mathcal Z_\omega(\eps)=\varnothing\},
 \label{p01:eq:gapdef}\\
 W_\omega(\eps)&=\inf\bigl(\mathcal Z_\omega(\eps)\cap[0,\infty)\bigr).
 \label{p01:eq:waitdef}
\end{align}
As usual the infimum of the empty set is infinity. The theorem shows it
is finite here. The definition of a gap ignores multiplicity.

\begin{theorem}[Finite-height edge zeros and recurrent long gaps]\label{p01:thm:zeros}
Assume \eqref{p01:eq:amplitudes}--\eqref{p01:eq:edgerange}. There are constants
$c,C,\eps_0>0$, independent of $\omega$, such that for
$0<\eps<\eps_0$ the following hold.
\begin{enumerate}[label=\textup{(\roman*)},leftmargin=2.1em]
\item Every real interval of length
$U(\eps)=\exp(CX_S(\eps))$ contains the ordinate of a simple zero in
$\mathcal S_\eps$. Every zero in this strip is simple.
\item For every $A\in\R$, some
$b\in[A,A+U(\eps)]$ has
\[
 [b,b+L(\eps)]\cap\mathcal Z_\omega(\eps)=\varnothing,
 \qquad L(\eps)=\exp(cY_S(\eps)).
\]
Thus the long empty intervals occur arbitrarily high for every
deterministic twist, with an explicit bound on the waiting time for
their starting points.
\item For every twist,
\begin{equation}\label{p01:eq:gapbounds}
 e^{cY_S(\eps)}\le G_\omega(\eps)\le e^{CX_S(\eps)},\qquad
 \lim_{\eps\downarrow0}
 \frac{\log\log G_\omega(\eps)}{\log(1/\eps)}
 =\frac1{S-1}.
\end{equation}
\item For product-Haar-almost every twist,
\begin{equation}\label{p01:eq:zerowaiting}
 \lim_{\eps\downarrow0}
 \frac{\log\log(e^e+W_\omega(\eps))}{\log(1/\eps)}
 =\frac1{S-1}.
\end{equation}
\end{enumerate}
\end{theorem}

The upper bound in (i) and the lower bound in (ii) concern different
lengths. They are compatible: $L(\eps)$ is much smaller than
$U(\eps)$. The long gaps in (ii) do not imply a corresponding lower
bound for the first zero above a fixed origin of a prescribed twist.

\begin{corollary}[Ordinary prime zeta]\label{p01:cor:primezeta}
Let $\mathcal P(s)=\sum_p p^{-s}$ and let $S$ be the solution of
$\mathcal P(S)=2^{1-S}$. Then $1<S<2$, the line $\Ree s=S$ is zero-free,
and the maximal gaps between ordinates of zeros in
$S-\eps<\Ree s<S$ satisfy \eqref{p01:eq:gapbounds}. Their exact
double-logarithmic exponent is $1/(S-1)$.
For orientation,
\[
 S=1.7795446535\ldots,\qquad \frac1{S-1}=1.2828001519\ldots .
\]
These decimals are illustrative; the theorem uses the exact edge equation.
\end{corollary}

\section{The edge germ and quadratic horizontal control}

Assume \eqref{p01:eq:amplitudes} and let $S$ satisfy \eqref{p01:eq:edge}.
Normalize the positive weights by
\begin{equation}\label{p01:eq:bweights}
 b_p=\frac{a_p}{a_2}(2/p)^S\quad(p\ge3),\qquad
 \sum_{p\ge3}b_p=1,\qquad
 h=\sum_{p\ge3}b_p\lambda_p>0.
\end{equation}
The weights $b_p$ satisfy \eqref{p01:eq:weights} with $\sigma=S$.
For an arbitrary phase $\eta\in\Omega$ put
\begin{equation}\label{p01:eq:germ}
 g_\eta(z)=1-\sum_{p\ge3}b_p\eta_pe^{-\lambda_p z},\qquad
 H(z)=g_1(z).
\end{equation}
They are analytic for $\Ree z>1-S$. We have
$H(0)=0$, $H'(0)=h$, and
\begin{equation}\label{p01:eq:realenergy}
 \Ree g_\eta(0)=E_b(\eta).
\end{equation}
For the original Dirichlet series define
\begin{equation}\label{p01:eq:phasequotient}
 \eta_p=-\frac{\omega_p}{\omega_2}.
\end{equation}
Then the exact identity
\begin{equation}\label{p01:eq:normalizedF}
 \frac{F_\omega(S+it+z)}
      {a_2\omega_2 2^{-(S+it+z)}}=g_{\Phi_t\eta}(z)
\end{equation}
holds throughout the domain.

\begin{lemma}[One-sided boundary and quadratic zero transfer]\label{p01:lem:transfer}
For all twists, $F_\omega$ has no zero in $\Ree s>S$.
There are $C,\eta_0>0$ such that whenever $E=E_b(\eta)\le\eta_0$,
$g_\eta$ has a simple zero $z_\eta$ satisfying
\begin{equation}\label{p01:eq:quadratictransfer}
 -CE\le\Ree z_\eta\le0,
 \qquad |\Imm z_\eta|\le C\sqrt E,
 \qquad
 z_\eta=-\frac{g_\eta(0)}h+O(E).
\end{equation}
All constants are uniform in the infinitely many phases.
\end{lemma}
\begin{proof}
For real $\sigma>S$,
$a_2 2^{-\sigma}>\sum_{p\ge3}a_pp^{-\sigma}$, so the first term
cannot be canceled. Equivalently, $g_\eta$ has no zero with positive
real part.

Choose a fixed small disk $|z|\le r$ with $S-2r>1$, on which $H$ has
only its simple zero at zero and $|H'(z)-h|<h/4$.
Cauchy--Schwarz gives, for $j=0,1,2$,
\begin{equation}\label{p01:eq:germcontrol}
 \sup_{|z|\le r}|(g_\eta-H)^{(j)}(z)|
 \le\sqrt{2E}
 \left(\sum_{p\ge3}b_p\lambda_p^{2j}e^{2r\lambda_p}\right)^{1/2}
 \le C_j\sqrt E.
\end{equation}
Each sum is finite by $S-2r>1$.
Rouch\'e's theorem on the boundary circle gives a unique zero in the
disk when $E$ is small, counted with multiplicity. The derivative
there is bounded below by $h/2$, so this zero is simple.
Shrinking the disk if necessary, $|H(z)|\ge c|z|$ throughout it.
At the zero, \eqref{p01:eq:germcontrol} therefore implies
$|z_\eta|\le C\sqrt E$.

Taylor's formula at zero, with a uniform bound for the second
derivative, gives
\[
 0=g_\eta(0)+h z_\eta
   +(g_\eta'(0)-h)z_\eta+O(|z_\eta|^2)
  =g_\eta(0)+h z_\eta+O(E).
\]
This proves the last estimate in \eqref{p01:eq:quadratictransfer}.
Its real part, together with \eqref{p01:eq:realenergy}, is $O(E)$ rather
than merely $O(\sqrt E)$. The zero-free half-plane supplies the upper
sign restriction. The other bounds follow.
\end{proof}

The distinction between horizontal and vertical displacement is
decisive: a direct uniform-error estimate without
\eqref{p01:eq:realenergy} would need phase energy of order $\eps^2$ and
would lose the principal exponent in Theorem~\ref{p01:thm:zeros}.

\begin{lemma}[Necessary energy at every edge zero]\label{p01:lem:necessary}
There are $B,\eps_0>0$ such that every zero
$S-\eps\le\sigma\le S$, $F_\omega(\sigma+it)=0$, satisfies
\begin{equation}\label{p01:eq:necessary}
 E_b(\Phi_t\eta)\le B\eps,
\end{equation}
where \eqref{p01:eq:phasequotient} defines $\eta$.
Every zero with $S-\eps_0\le\Ree s\le S$ is simple.
\end{lemma}
\begin{proof}
Put $\delta=S-\sigma\ge0$ and $q=\Phi_t\eta$. At the zero,
\begin{equation}\label{p01:eq:zeroequation}
 \sum_{p\ge3}b_pe^{\lambda_p\delta}q_p=1.
\end{equation}
Taking real parts gives the exact defect formula
\begin{equation}\label{p01:eq:defect}
 \sum_{p\ge3}b_pe^{\lambda_p\delta}(1-\Ree q_p)
 =A(\delta)-1,
 \quad A(\delta):=\sum_{p\ge3}b_pe^{\lambda_p\delta}.
\end{equation}
For a fixed sufficiently small $\delta_0<S-1$, this is an analytic
function with
$A(\delta)=1+h\delta+O(\delta^2)$, uniformly on $[0,\delta_0]$.
Since $e^{\lambda_p\delta}\ge1$ and each summand is nonnegative,
\eqref{p01:eq:defect} implies \eqref{p01:eq:necessary}.

Finally,
\begin{align*}
 |g_q'(-\delta)-h|
 &\le\sum b_p\lambda_p(e^{\lambda_p\delta}-1)
  +\left|\sum b_p\lambda_pe^{\lambda_p\delta}(q_p-1)\right|\\
 &\le O(\delta)
  +\sqrt{2(A(\delta)-1)}
       \left(\sum b_p\lambda_p^2e^{\lambda_p\delta}\right)^{1/2}
 =O(\sqrt\delta).
\end{align*}
For small $\delta_0$ this derivative is nonzero. By
\eqref{p01:eq:normalizedF}, the original zero is also simple. At
$\delta=0$, the defect vanishes and the derivative is exactly $h$.
\end{proof}

\subsection{Curvature transverse to the vertical flow}

The preceding lemmas require no coordinatewise small-angle assumption.
For smooth one-parameter phase perturbations, the geometry can be
identified more precisely.

\begin{proposition}[Exact quadratic form]\label{p01:prop:curvature}
Let $v_p$ be real and satisfy $|v_p|\le C_v\lambda_p$. Set
$\eta_p(u)=e^{iu v_p}$ for real $u$ near zero. The zero from
Lemma~\ref{p01:lem:transfer} is an analytic function $z(u)$ with $z(0)=0$.
If $V=\sum b_pv_p$, then
\begin{equation}\label{p01:eq:curvature}
 \Ree z(u)=-\frac{u^2}{2h}
    \sum_{p\ge3}b_p\left(v_p-\frac Vh\lambda_p\right)^2+O(u^4).
\end{equation}
The quadratic form vanishes exactly when $v_p=c\lambda_p$ for every
$p$. In that case $z(u)=icu$ identically near zero.
\end{proposition}
\begin{proof}
The growth condition ensures local analyticity of
\[
 G(z,u)=1-\sum b_p\exp(iuv_p-\lambda_p z)
\]
in both complex variables, by normal convergence. The implicit function
theorem applies since $G_z(0,0)=h>0$. Write
\[
 C_1=\sum b_p\lambda_pv_p,\quad V_2=\sum b_pv_p^2,
 \quad H_2=\sum b_p\lambda_p^2.
\]
Differentiation gives $z'(0)=iV/h$ and
\[
 z''(0)=-\frac1h\left(V_2-\frac{2VC_1}h
                                  +\frac{V^2H_2}{h^2}\right).
\]
The bracket is the sum of squares in \eqref{p01:eq:curvature}.
For real $u$, uniqueness gives $z(-u)=\overline{z(u)}$, so the real
part is even and its remainder after the quadratic term is $O(u^4)$.
All $b_p$ are positive, hence the sum vanishes precisely when
$v_p=(V/h)\lambda_p$. Direct substitution then gives the exact root
$z(u)=i(V/h)u$.
\end{proof}

Thus the only direction with vanishing quadratic inward displacement
is the actual vertical-translation direction. The $O(E)$ control in
Lemma~\ref{p01:lem:transfer} is consistent with this more precise local
geometry, while remaining valid for arbitrary tail phases.

\section{Proof of the finite-height and gap theorems}

\subsection{Uniform simple zeros in every long interval}

Apply the upper bound of Theorem~\ref{p01:thm:flow} with weights $b_p$ and
energy tolerance $c_2\eps$, where $c_2$ is a sufficiently small fixed
constant. In every interval of length
$\exp(CX_S(\eps))$ there is a time $t$ with
$E_b(\Phi_t\eta)\le c_2\eps$. By
Lemma~\ref{p01:lem:transfer}, the original function has a simple zero
\[
 \rho=S+it+z,\qquad -\eps/2\le\Ree z\le0,
 \quad |\Imm z|\le C\sqrt\eps.
\]
For small $\eps$, the last displacement is less than one. Given an
interval $[A,A+T+2]$, run recurrence on $[A+1,A+T+1]$; the zero lies
in the required height interval. Increasing the fixed constant in
the exponential absorbs the additive two. Lemma~\ref{p01:lem:necessary}
proves the simplicity of every zero in the strip. This proves
Theorem~\ref{p01:thm:zeros}(i), and also
\begin{equation}\label{p01:eq:waitupper}
 W_\omega(\eps)\le\exp(CX_S(\eps)).
\end{equation}

\subsection{Robust empty intervals and their recurrence}

Let $B$ be the constant in \eqref{p01:eq:necessary}. Apply
Lemma~\ref{p01:lem:avoidance} at tolerance $16B\eps$. It gives a phase
$\eta^*$ and a length
\[
 L=\exp(cY_S(\eps))
\]
such that
\begin{equation}\label{p01:eq:robustbad}
 E_b(\Phi_u\eta^*)>16B\eps\qquad(0\le u\le L).
\end{equation}
Changing the tolerance by a fixed factor changes only the constants
in this length. The strict inequality holds simultaneously on the
whole compact interval.

The metric triangle inequality and isometry of the flow imply
\begin{equation}\label{p01:eq:metrictransport}
 \left|\sqrt{E_b(\Phi_u\xi)}
             -\sqrt{E_b(\Phi_u\eta^*)}\right|
 \le\frac1{\sqrt2}d_b(\xi,\eta^*)
 \qquad(u\in\R).
\end{equation}
Therefore every phase $\xi$ satisfying
$d_b(\xi,\eta^*)^2\le2B\eps$ has
\[
 E_b(\Phi_u\xi)>9B\eps>B\eps\qquad(0\le u\le L).
\]
By \eqref{p01:eq:necessary}, its associated Dirichlet series has no zero
in the edge strip at any of these heights.

Now fix any deterministic twist $\omega$ and any initial height $A$.
The covering theorem, with target $\eta^*$ and tolerance $B\eps$,
provides a
\[
 b\in[A,A+\exp(CX_S(\eps))]
 \quad\text{such that}\quad
 d_b(\Phi_b\eta,\eta^*)^2\le2B\eps.
\]
The entire segment $[b,b+L]$ is consequently zero-free in the strip.
This proves (ii). It also shows why a probabilistically constructed
bad phase gives long holes for \emph{every} fixed twist: the same
weighted isometry preserves the exclusion for the entire segment,
and the covering theorem transports its starting point effectively.

Part (iii) follows at once. For the upper bound on a supremum of
open zero-free intervals, if an interval were longer than $U(\eps)$,
it would contain a closed subinterval of that length and hence a zero
by (i). Taking two logarithms as in Theorem~\ref{p01:thm:flow} gives the
limit in \eqref{p01:eq:gapbounds}.

\subsection{Almost-sure first waiting time}

For independent Haar phases $\omega_p$, the relative phases
$\eta_p=-\omega_p/\omega_2$, $p\ge3$, have product Haar law. One may
condition on $\omega_2$ to verify this. If an edge zero occurs at
some $t\in[0,T]$, Lemma~\ref{p01:lem:necessary} implies
$E_b(\Phi_t\eta)\le B\eps$. Hence
\begin{equation}\label{p01:eq:zeroprob}
 \mathbb P\{W_\omega(\eps)<T\}
 \le C(1+T/\eps)e^{-cY_S(\eps)}.
\end{equation}
The use of $<$ avoids any issue of an infimum at the open left boundary
of the strip; using $T+1$ yields the same bound for $\le T$ after
changing constants.

Apply this estimate with $\eps_n=e^{-n}$ and
$T_n=\exp(c_3Y_S(\eps_n))$, $c_3>0$ sufficiently small. The resulting
probabilities are summable. Borel--Cantelli, \eqref{p01:eq:waitupper},
and monotonicity in $\eps$ give part (iv). Independence between times
or between the events indexed by $n$ is not required.

\section{Consequences for prime zeta and boundary-attaining twists}

\begin{lemma}[Boundary classification]\label{p01:lem:boundary}
The function $F_\omega$ has a zero at $S+it$ if and only if
\begin{equation}\label{p01:eq:boundarytwist}
 -\frac{\omega_p}{\omega_2}e^{-it\log(p/2)}=1
 \qquad(p\ge3).
\end{equation}
Each twist has at most one such ordinate. The untwisted prime zeta
function has none.
\end{lemma}
\begin{proof}
At $\delta=0$, the defect identity \eqref{p01:eq:defect} is a sum of
nonnegative terms equal to zero, and every weight is positive. This
is equivalent to \eqref{p01:eq:boundarytwist}; the converse follows from
the edge equation.
If two distinct ordinates satisfied the condition, their difference
$t_0\ne0$ would have both $t_0\log(3/2)$ and $t_0\log(5/2)$ in
$2\pi\Z$. This would make
$\log(3/2)/\log(5/2)$ rational. Unique factorization rules that out:
an equality $(3/2)^m=(5/2)^n$ for integers forces $m=n=0$.
For the untwisted function, \eqref{p01:eq:boundarytwist} would instead
make these two products odd multiples of $\pi$, giving the same
contradiction. At $t=0$ the condition also fails.
\end{proof}

\begin{proof}[Proof of Corollary~\ref{p01:cor:primezeta}]
For $a_p=1$, divergence of $\sum p^{-\sigma}$ as $\sigma\downarrow1$
gives $S>1$. At $\sigma=2$,
\[
 \sum_{p\ge3}(2/p)^2
 \le4\sum_{\substack{n\ge3\\n\text{ odd}}}n^{-2}
 =\frac{\pi^2}{2}-4<1,
\]
so $S<2<3$. Theorem~\ref{p01:thm:zeros} applies, and
Lemma~\ref{p01:lem:boundary} removes the edge line from the strip without
changing its zero set.
\end{proof}

\begin{corollary}[Strict strips for every twist]\label{p01:cor:strict}
The maximal-gap exponent in Theorem~\ref{p01:thm:zeros} is unchanged if
one uses $S-\eps<\Ree s<S$ instead of $\mathcal S_\eps$, for any
twist. A uniform upper waiting bound of the same order also holds.
\end{corollary}
\begin{proof}
At most one ordinate is removed by Lemma~\ref{p01:lem:boundary}. Two
disjoint closed intervals of length $U(\eps)$, separated by one unit,
each contain an edge-strip zero. They have different ordinates, so at
least one zero is strictly to the left of $S$. This gives an upper
bound of length $2U(\eps)+1$, absorbed in the exponential constant.
The long holes already exclude the larger strip, so their lower
bound persists.
\end{proof}

\begin{corollary}[An explicit approach rate at finite height]\label{p01:cor:approach}
For every twist and all sufficiently large $T$, a simple zero has
$0\le\Imm\rho\le T$ and
\begin{equation}\label{p01:eq:approach}
 0\le S-\Ree\rho
 \le\frac{C}{(\log T)^{S-1}\log\log T}.
\end{equation}
The zero can be chosen strictly to the left of $S$, with a changed
constant.
\end{corollary}
\begin{proof}
Put $L=\log T$ and choose
$\eps=K L^{-(S-1)}(\log L)^{-1}$.
Then
\[
 \eps\log(1/\eps)
 =K(S-1)L^{-(S-1)}(1+o(1)),
\]
so $X_S(\eps)=(K(S-1))^{-1/(S-1)}L(1+o(1))$.
Choose $K$ large enough that the upper bound in
\eqref{p01:eq:waitupper} is less than $T$. Apply
Theorem~\ref{p01:thm:zeros}(i), or Corollary~\ref{p01:cor:strict} for the
strict version.
\end{proof}

\begin{remark}[Uniformity versus practical search]
Equation~\eqref{p01:eq:explicitT} is a finite, verifiable certificate once
the weights and degrees are supplied. Nevertheless, the heights in the
general bound are enormous. These results quantify existence and
recurrence; they do not provide a practical numerical search method
for extremely thin edge strips.
\end{remark}

\subsection{Quantifiers in the prime-zeta conclusions}

The same exponent appears in three logically different statements.
The following table records their quantifiers.
\begin{center}
\small
\begin{tabular}{@{}>{\raggedright\arraybackslash}p{.35\linewidth}
 >{\raggedright\arraybackslash}p{.32\linewidth}
 >{\raggedright\arraybackslash}p{.25\linewidth}@{}}
\toprule
Quantity & Phase quantifier & Conclusion\\
\midrule
Weighted torus covering time & All starts and targets & Exact exponent
$1/(\sigma-1)$\\[3pt]
Maximal shrinking-edge zero gap & Every fixed twist & Exact exponent
$1/(S-1)$\\[3pt]
Fixed-origin first zero & Haar-almost every twist & Exact exponent
$1/(S-1)$\\[3pt]
Fixed-origin first zero of $\mathcal P$ & One specified twist & Uniform
upper bound only\\
\bottomrule
\end{tabular}
\end{center}
Here ``exact exponent'' always means a limit of the form
$\log\log(\text{time})/\log(1/\eps)$, not a full asymptotic for the
time itself.

\subsection{Proof assumptions for the recurrence-to-zero transfer}

The recurrence-to-zero arguments above depend on the following distinctions.
\begin{enumerate}[leftmargin=2em]
\item \textbf{Growing dimension.} The number of retained primes tends
to infinity. All Fourier norms and frequency-separation constants
are tracked explicitly in \eqref{p01:eq:explicitT}; a fixed-dimensional
equidistribution theorem alone is insufficient.
\item \textbf{Relative frequencies.} Removing the common phase
replaces $\log p$ by $\log(p/2)$. The resulting powers of $2$ are
included in Lemma~\ref{p01:lem:separation}; unique factorization still
rules out nontrivial resonances.
\item \textbf{Energy versus amplitude.} Small energy gives a complex
function error of order $\sqrt E$. The exact real-part identity and
one Newton expansion are what improve the horizontal zero error to
$O(E)$.
\item \textbf{The unobserved tail.} It is bounded for all phases by
$2\sum_{p>x}w_p$. No finite phase box is treated as though it
controlled infinitely many coordinates.
\item \textbf{No mixing.} The avoidance estimate uses one-time Haar
laws and a finite union bound. The almost-sure step uses the first
Borel--Cantelli lemma, which requires no independence.
\item \textbf{Holes for a fixed function.} A positive-measure family
of bad initial phases alone gives only a worst-twist lower bound.
The isometry argument \eqref{p01:eq:metrictransport} and the uniform
covering theorem are explicitly used to transfer an entire bad
segment to every twist.
\item \textbf{Height endpoints.} The generated zero changes its
ordinate by $O(\sqrt\eps)$. Unit buffers are inserted before the
upper waiting statement is made.
\item \textbf{Sharpness.} The power of $\eps$ after two logarithms
is determined. The single-logarithmic asymptotic and its constant
are not determined.
\end{enumerate}

\section{Euler products: critical points and quantitative locking}
\label{p02:sec:euler}

For a sequence $\omega=(\omega_p)_p\in\prod_p\T$, define
\begin{equation}\label{p02:eul:product}
 Z_\omega(s)=\prod_p(1-\omega_pp^{-s})^{-1},\qquad
 Q_\omega(s)=-\frac{Z_\omega'(s)}{Z_\omega(s)}
 =\sum_p\frac{(\log p)\omega_pp^{-s}}{1-\omega_pp^{-s}}.
\end{equation}
These expressions and their derivatives converge locally uniformly in
$\Re s>1$, where $Z_\omega$ never vanishes. The phases of prime powers
in an Euler product are coupled; they cannot be replaced by independent
coordinates. The relevant flow here uses the full frequencies $\log p$,
including $p=2$.

Let $E_*>1$ be the unique solution of
\begin{equation}\label{p02:eul:critical-edge}
 \frac{\log2}{2^{E_*}+1}
 =\sum_{p\ge3}\frac{\log p}{p^{E_*}-1},
 \qquad \beta_*=(E_*-1)^{-1}.
\end{equation}
For $\omega_p=1$, this is the classical right boundary of the zeros
of $\zeta'$. Arias de Reyna and van de Lune
\cite[Theorem~5.1]{AriasLune} give
$E_*=2.813014020252898\ldots$, in particular $1<E_*<3$.
Their Theorem~6.1 also proves qualitative freezing of individual prime
phases along sequences approaching the boundary. The quantitative
energy estimates below strengthen that coordinatewise conclusion.

Write $\mathcal C_\omega(\eps)$ for the ordinates of critical points
of $Z_\omega$ in $E_*-\eps<\Re s\le E_*$, and define
\[
 G^{\rm crit}_\omega(\eps)=
 \sup\{b-a:(a,b)\cap\mathcal C_\omega(\eps)=\varnothing\}.
\]

\begin{theorem}[Critical-point gaps for every Euler-product twist]
\label{p02:eul:gap}
There are positive constants $c,C,\eps_0$, independent of $\omega$,
such that, for $0<\eps<\eps_0$, the following assertions hold. Put
\[
 U_*(\eps)=\exp(C\eps^{-\beta_*}),\qquad
 L_*(\eps)=\exp\!\left(
        \frac{c\eps^{-\beta_*}}{\log(1/\eps)}\right).
\]
Every real interval of length $U_*(\eps)$ contains an ordinate in
$\mathcal C_\omega(\eps)$, and all the critical points in this strip
are simple zeros of $Z_\omega'$. For every real $a$, there is
$b\in[a,a+U_*(\eps)]$ such that
\[
 [b,b+L_*(\eps)]\cap\mathcal C_\omega(\eps)=\varnothing.
\]
Consequently
\begin{equation}\label{p02:eul:gap-exponent}
 L_*(\eps)\le G^{\rm crit}_\omega(\eps)\le U_*(\eps),
 \qquad
 \lim_{\eps\downarrow0}
 \frac{\log\log G^{\rm crit}_\omega(\eps)}{\log(1/\eps)}
 =\frac1{E_*-1}.
\end{equation}
If $W^{\rm crit}_\omega(\eps)$ is the first nonnegative ordinate in
$\mathcal C_\omega(\eps)$, then for product-Haar-almost every $\omega$,
\begin{equation}\label{p02:eul:waiting}
 \lim_{\eps\downarrow0}
 \frac{\log\log(e^e+W^{\rm crit}_\omega(\eps))}{\log(1/\eps)}
 =\frac1{E_*-1}.
\end{equation}
\end{theorem}

\subsection{A circle identity and a positive boundary defect}

The frozen phase is
\[
 z^*_2=-1,\qquad z^*_p=1\quad(p\ge3).
\]
Define its weighted squared-distance defect by
\begin{equation}\label{p02:eul:critical-energy}
 D_E(z)=\sum_p(\log p)p^{-E_*}|z_p-z_p^*|^2.
\end{equation}
There is no common-rotation freedom in this definition.

\begin{lemma}[Quantitative rigidity of critical points]
\label{p02:eul:rigidity}
For every twist, $Q_\omega$ has no zeros in $\Re s>E_*$. If
$Q_\omega(\sigma+it)=0$ and $E_*-\eps<\sigma\le E_*$, then
\begin{equation}\label{p02:eul:necessary}
 D_E((\omega_pp^{-it})_p)\le C\eps
\end{equation}
for sufficiently small $\eps$, uniformly in the twist and the ordinate.
At $\sigma=E_*$ the phases equal $z^*$ exactly.
\end{lemma}

\begin{proof}
For real $\sigma>1$, put $r_p=p^{-\sigma}$ and
\[
 c_p=\frac{(\log p)r_p^2}{1-r_p^2},\qquad
 R_p=\frac{(\log p)r_p}{1-r_p^2},\qquad
 \xi_p=\frac{z_p-r_p}{1-r_pz_p}.
\]
For $|z_p|=1$, one has $|\xi_p|=1$ and the exact identity
\begin{equation}\label{p02:eul:circle}
 \frac{(\log p)r_pz_p}{1-r_pz_p}=c_p+R_p\xi_p.
\end{equation}
Set
\[
 C_0=\sum_pc_p,\qquad b_0=R_2,\qquad R_0=\sum_{p\ge3}R_p,
 \qquad \Delta(\sigma)=C_0+R_0-b_0.
\]
Then
\begin{equation}\label{p02:eul:Delta}
 \Delta(\sigma)=\sum_{p\ge3}\frac{\log p}{p^\sigma-1}
                   -\frac{\log2}{2^\sigma+1}.
\end{equation}
Divide the first sum by $h(\sigma)=(\log2)/(2^\sigma+1)$.
Each summand of the resulting ratio has logarithmic derivative
\[
 \frac{(\log2)2^\sigma}{2^\sigma+1}
 -\frac{(\log p)p^\sigma}{p^\sigma-1}<0.
\]
The ratio decreases from infinity to zero. Thus $\Delta$ has the
unique zero $E_*$ and is negative to its right.

At a zero of $Q$, take $z_p=\omega_pp^{-it}$. Equation
\eqref{p02:eul:circle} gives
\begin{equation}\label{p02:eul:vector-balance}
 C_0+b_0\xi_2=-\sum_{p\ge3}R_p\xi_p,
 \qquad |C_0+b_0\xi_2|\le R_0.
\end{equation}
If $\sigma>E_*$, then $b_0>C_0+R_0$, contradicting this inequality.
Near $E_*$ from the left, $b_0-C_0$ is positive and
$C_0,b_0,R_0$ stay bounded above and below by positive constants.
Squaring \eqref{p02:eul:vector-balance} yields
\begin{equation}\label{p02:eul:head-defect}
 C_0b_0|\xi_2+1|^2
 \le R_0^2-(b_0-C_0)^2
 =\Delta(\sigma)(R_0+b_0-C_0).
\end{equation}
Taking real parts in the same identity gives
\begin{equation}\label{p02:eul:tail-defect}
 \sum_{p\ge3}R_p|\xi_p-1|^2
 =2\Delta(\sigma)+b_0|\xi_2+1|^2.
\end{equation}
Since $\Delta(\sigma)=O(E_*-\sigma)$, both defects are $O(\eps)$.
Finally,
\[
 |\xi-1|=\frac{(1+r)|z-1|}{|1-rz|},\qquad
 |\xi+1|=\frac{(1-r)|z+1|}{|1-rz|}.
\]
These are uniform comparisons because all $r_p$ are bounded away
from $1$, and $R_p\ge(\log p)p^{-E_*}$ for $\sigma\le E_*$.
This proves \eqref{p02:eul:necessary}. At the boundary both right sides
in \eqref{p02:eul:head-defect}--\eqref{p02:eul:tail-defect} vanish, which
forces precisely $z=z^*$.
\end{proof}

\begin{lemma}[Realization with a quadratic real-part loss]
\label{p02:eul:local-critical}
There are constants $d_0,K>0$ such that $D=D_E(z)<d_0$ implies that
$Q_z$ has a simple zero $E_*+h_z$ with
\begin{equation}\label{p02:eul:local-bounds}
 -KD\le\Re h_z\le0,\qquad |\Im h_z|\le K\sqrt D.
\end{equation}
Furthermore every zero in a sufficiently thin right-edge strip is
simple.
\end{lemma}

\begin{proof}
The frozen function is $Q_{z^*}(s)=\Delta(s)$. Put
$a_E=\Delta'(E_*)$. If
$h=(\log2)/(2^{E_*}+1)$ and
$t_p=(\log p)/(p^{E_*}-1)$, then $\sum_{p\ge3}t_p=h$ and
\[
 a_E=h\frac{(\log2)2^{E_*}}{2^{E_*}+1}
 -\sum_{p\ge3}t_p\frac{(\log p)p^{E_*}}{p^{E_*}-1}
 <h(\log2-\log3)<0.
\]
For $F_z(u)=Q_z(E_*+u)$, the circle identity at $\sigma=E_*$,
followed by Cauchy--Schwarz, gives
\begin{equation}\label{p02:eul:taylor-data}
 |F_z(0)|\le K\sqrt D,\qquad
 |\Re F_z(0)|\le KD,\qquad
 |F_z'(0)-a_E|\le K\sqrt D.
\end{equation}
For the real estimate each tail summand has real-part change
$-R_p|\xi_p-1|^2/2$, and the $p=2$ summand has change
$R_2|\xi_2+1|^2/2$. For the derivative estimate the local difference
is at most $K(\log p)^2p^{-E_*}|z_p-z_p^*|$; its
Cauchy--Schwarz companion sum is
$\sum_p(\log p)^3p^{-E_*}<\infty$.

Choose a fixed $r_0>0$ with $E_*-2r_0>1$. On $|u|\le r_0$
the second derivatives are uniformly bounded, and hence
\begin{equation}\label{p02:eul:local-taylor}
 F_z(u)=F_z(0)+a_Eu+O(\sqrt D\,|u|+|u|^2).
\end{equation}
Rouch\'e's theorem on $|u|=K_0\sqrt D$, with $K_0$ sufficiently
large and then $D$ sufficiently small, gives one zero $h_z$ in
that disk. Its derivative is $a_E+O(\sqrt D)$, so it is simple.
Taking real parts in \eqref{p02:eul:local-taylor} proves
$|\Re h_z|\le KD$; Lemma~\ref{p02:eul:rigidity} supplies the sign.
The case $D=0$ is the frozen zero itself.

Conversely, translate any zero in the edge strip to its own ordinate.
Lemma~\ref{p02:eul:rigidity} gives $D=O(\eps)$, so its derivative is
$a_E+O(\sqrt\eps)\ne0$. At a zero of $Q_z$ one has
$Z_z''=-Z_zQ_z'$, proving the asserted simplicity for $Z_z'$.
\end{proof}

\subsection{The finite-height consequences}

\begin{proof}[Proof of Theorem~\ref{p02:eul:gap}]
Apply Theorem~\ref{p02:thm:arithmetic-flow} to full prime frequencies and
weights $(\log p)p^{-E_*}$. This is the logarithmic-weight case:
the truncation and entropy scales are respectively
$\eps^{-\beta_*}$ and $\eps^{-\beta_*}/\log(1/\eps)$.
The energy in \eqref{p02:eul:critical-energy} differs by a factor of
two from the energy used there, after translating the target to $z^*$.
Thus every orbit reaches $D_E\le c_0\eps$ on every interval of
length $\exp(C\eps^{-\beta_*})$. Lemma~\ref{p02:eul:local-critical}
produces a critical point in the required strip and moves its ordinate
by at most $K\sqrt\eps$. Padding the initial interval and enlarging
$C$ gives the stated upper bound.

By Lemma~\ref{p02:eul:rigidity}, the ordinate of every critical point in
the strip is a visit to $D_E\le K\eps$. The robust-hole assertion
of Proposition~\ref{p02:prop:holes}, applied to this enlarged target,
therefore gives the asserted empty intervals and their recurrence
windows. This proves \eqref{p02:eul:gap-exponent}.

For completeness, the almost-sure lower waiting estimate requires no
independence between times. The weighted energy is uniformly Lipschitz
along the flow, since $\sum_p(\log p)^2p^{-E_*}<\infty$.
A grid of mesh comparable to $\eps$ and the upper small-ball bound
in Theorem~\ref{p02:thm:smallball} give
\[
 \mathbb P\!\left\{\min_{0\le t\le T}D_E(\Phi_t\omega)\le K\eps\right\}
 \le C(1+T/\eps)
 \exp\!\left(-\frac{c\eps^{-\beta_*}}{\log(1/\eps)}\right).
\]
Take $\eps_n=e^{-n}$ and
$T_n=\exp(c_1\eps_n^{-\beta_*}/n)$ with $c_1$ sufficiently small.
The displayed probabilities are summable. Borel--Cantelli,
Lemma~\ref{p02:eul:rigidity}, monotone interpolation between consecutive
$\eps_n$, and the deterministic upper bound prove
\eqref{p02:eul:waiting}.
\end{proof}

\begin{corollary}[Characters of an odd modulus]\label{p02:eul:characters}
Let $q$ be odd, and let $E_q>1$ solve
\[
 \frac{\log2}{2^{E_q}+1}
 =\sum_{\substack{p\ge3\\p\nmid q}}
                       \frac{\log p}{p^{E_q}-1}.
\]
Then $1<E_q\le E_*<3$. For every Dirichlet character $\chi$ modulo
$q$, the deterministic assertions of Theorem~\ref{p02:eul:gap} hold for the critical points of
$L(s,\chi)$ with $E_q$ replacing $E_*$ and constants depending on
$q$. In particular their maximal-gap exponent is $(E_q-1)^{-1}$,
independently of the character.
\end{corollary}
\begin{proof}
Omit the finitely many primes dividing $q$ in every sum and product.
The smallest remaining prime is $2$, and the decreasing-ratio proof
gives the stated boundary. Omitting positive tail terms can only lower
it. All circle identities and local estimates remain valid.
The prime-flow and small-ball estimates are unchanged in scale by
finitely many omissions. Finally $\chi(p)$ supplies a fixed unimodular
twist on the remaining primes. This proves the deterministic assertions;
an almost-sure statement, if desired, concerns additional Haar twists,
not probability on the finite set of characters.
\end{proof}

\subsection{A sharp quantitative relation between 1-points and critical points}

Let $A_*>1$ be the unique solution of
\begin{equation}\label{p02:eul:one-edge}
 \sum_{\substack{n\ge3\\n\text{ odd}}}n^{-A_*}=2^{-A_*},
 \quad\text{equivalently}\quad
 \zeta(A_*)=\frac{2^{A_*}+1}{2^{A_*}-1}.
\end{equation}
Its classical value is $A_*=1.940101683743625\ldots$
\cite[Theorem~4.1]{AriasLune}. Put
\begin{equation}\label{p02:eul:locking-scale}
 \theta=\frac{A_*-1}{E_*-1}\in(0,1),\qquad
 \Psi(\delta)=\frac{\delta^\theta}{\log(1/\delta)},
 \qquad
 D_A(z)=\sum_pp^{-A_*}|z_p-z_p^*|^2.
\end{equation}
Theorem~6.2 of \cite{AriasLune} establishes a qualitative relation
between the two kinds of almost-extremal points. The next result gives
an explicit rate, including its logarithmic factor, and proves that
this rate cannot be uniformly improved.

\begin{theorem}[Quantitative locking and its sharpness]
\label{p02:eul:locking}
There are constants $K,\delta_0>0$, independent of the twist, with the
following properties.
\begin{enumerate}
\item If $Z_\omega'(E_*-\delta+i\gamma)=0$, where
$0<\delta<\delta_0$, then $Z_\omega$ has a $1$-point $s$ satisfying
\[
 0\le A_*-\Re s\le K\Psi(\delta),\qquad
 |\Im s-\gamma|\le K\sqrt\delta.
\]
\item If $Z_\omega(A_*-\eps+i\gamma)=1$, where
$0<\eps<\delta_0$, then it has a critical point $\rho$ satisfying
\[
 0\le E_*-\Re\rho\le K\eps,\qquad
 |\Im\rho-\gamma|\le K\sqrt\eps.
\]
\item The real-part rate $\Psi$ in the first assertion is sharp up
to constant factors, both uniformly among twists and along sequences
for the ordinary zeta function. More precisely, there are critical
points $\rho_n$ of $\zeta$ and $1$-points $s_n$ such that, with
$\delta_n=E_*-\Re\rho_n\downarrow0$ after passage to a subsequence,
\[
 A_*-\Re s_n\asymp\Psi(\delta_n),\qquad
 |\Im s_n-\Im\rho_n|\longrightarrow0.
\]
The $1$-point $s_n$ is the unique $1$-point in a fixed sufficiently
small disk centered at $A_*+i\Im\rho_n$. Thus the sharpness concerns
the local counterpart supplied by the first assertion.
\end{enumerate}
\end{theorem}

\begin{lemma}[Rigidity and realization for $1$-points]
\label{p02:eul:one-rigidity}
Every twist has no $1$-points in $\Re s>A_*$. A $1$-point
$\sigma+it$ with $A_*-\eps<\sigma\le A_*$ satisfies
\[
 D_A((\omega_pp^{-it})_p)\le K\eps.
\]
Conversely, if $D=D_A(z)$ is sufficiently small, there is a simple
$1$-point $A_*+u_z$ with
$-KD\le\Re u_z\le0$ and $|\Im u_z|\le K\sqrt D$.
\end{lemma}
\begin{proof}
For prime phases $z$, write $z(n)=\prod_pz_p^{v_p(n)}$, and set
$M(\sigma)=\sum_{n\ge3,\,n\text{ odd}}n^{-\sigma}$.
The equation $Z_z(\sigma)=1$ is equivalent to
\[
 \sum_{\substack{n\ge3\\n\text{ odd}}}z(n)n^{-\sigma}
             =-2^{-\sigma}z_2.
\]
The triangle inequality implies $2^{-\sigma}\le M(\sigma)$.
Since $M(\sigma)/2^{-\sigma}=\sum_{n\ge3,\,n\text{ odd}}(2/n)^\sigma$
is strictly decreasing, this proves the boundary assertion.
Put $\tau=-z_2$. The same equation gives the exact identity
\begin{equation}\label{p02:eul:odd-defect}
 \sum_{\substack{n\ge3\\n\text{ odd}}}
 n^{-\sigma}|z(n)-\tau|^2
       =2\bigl(M(\sigma)-2^{-\sigma}\bigr)=O(\eps).
\end{equation}
The terms $n=3,9$ show that
$|z_3-\tau|+|z_3^2-\tau|=O(\sqrt\eps)$; hence
$|\tau-1|=|\tau^2-\tau|=O(\sqrt\eps)$.
The prime terms in \eqref{p02:eul:odd-defect}, the inequality
$p^{-A_*}\le p^{-\sigma}$, and summability of $\sum_pp^{-A_*}$
now give $D_A(z)=O(\eps)$.

For the converse, the frozen product is
\[
 Z_{z^*}(s)=\frac{1-2^{-s}}{1+2^{-s}}\zeta(s),\qquad
 Z_{z^*}(s)-1=\frac{M(s)-2^{-s}}{1+2^{-s}}.
\]
Its derivative $a_A$ at $A_*$ is strictly negative, since
$M(s)/2^{-s}$ decreases strictly through $1$ there. Use the absolutely
convergent logarithm
$\mathcal L_z(s)=-\sum_p\log(1-z_pp^{-s})$, defined by its power
series. At $s=A_*$ the frozen logarithm is $0$. Each local real part
has a stationary extremum at its frozen phase, so
\[
 |\mathcal L_z(A_*)|=O(\sqrt D),\qquad
 |\Re\mathcal L_z(A_*)|=O(D).
\]
These estimates follow directly from
$|1-rz|^2=(1-r)^2+r|z-1|^2$ at $z^*=1$, and
$|1-rz|^2=(1+r)^2-r|z+1|^2$ at $z^*=-1$,
followed by Cauchy--Schwarz. Exponentiating preserves the bounds
$|Z_z(A_*)-1|=O(\sqrt D)$ and
$|\Re(Z_z(A_*)-1)|=O(D)$.
The derivative differs from $a_A$ by $O(\sqrt D)$ and the second
derivative is uniformly bounded in a fixed disk: the necessary sums
$\sum_p(\log p)^jp^{-A_*+2r_0}$ converge after choosing
$A_*-2r_0>1$. The Rouch\'e and real-part argument of
Lemma~\ref{p02:eul:local-critical} proves the assertion.
\end{proof}

\begin{proof}[Proof of Theorem~\ref{p02:eul:locking}]
First $D_E(z)\le KD_A(z)$, because
$(\log p)p^{-(E_*-A_*)}$ is bounded. In the reverse direction,
for sufficiently large $x$,
\begin{align}\label{p02:eul:interpolation}
 D_A(z)
 &\le K\frac{x^{E_*-A_*}}{\log x}D_E(z)
       +K\frac{x^{1-A_*}}{\log x}.
\end{align}
Indeed the ratio of the two weights for $p\le x$ is bounded by
$Kx^{E_*-A_*}/\log x$, and the tail uses
$|z_p-z_p^*|^2\le4$ and the prime-number estimate for
$\sum_{p>x}p^{-A_*}$.
For a critical point with real-part gap $\delta$,
Lemma~\ref{p02:eul:rigidity} gives $D_E\le K\delta$ at its own
ordinate. Taking $x=\delta^{-1/(E_*-1)}$ in
\eqref{p02:eul:interpolation} gives $D_A\le K\Psi(\delta)$.
The sharper ordinate estimate uses
\begin{equation}\label{p02:eul:crossweight}
 2A_*-E_*>1.
\end{equation}
The certified enclosures in Lemma~\ref{p02:lem:edge-enclosures} give
$2A_*-E_*>1.04$. Choose a fixed disk radius $r_0>0$ such that
$2A_*-E_*-2r_0>1$. Cauchy--Schwarz against the $D_E$ weights
bounds the logarithmic Euler germ and its first derivative in that
disk by $O(\sqrt{D_E})$: the companion series are bounded by
\[
 \sum_p\frac{p^{E_*-2A_*+2r_0}}{\log p},\qquad
 \sum_p(\log p)p^{E_*-2A_*+2r_0},
\]
both convergent. Thus $Z_z(A_*+u)-1$ differs in $C^1$ from the
frozen germ by $O(\sqrt\delta)$. Rouch\'e gives a zero with
$|u|=O(\sqrt\delta)$. At $u=0$ its real part remains $O(D_A)$,
by the proof of Lemma~\ref{p02:eul:one-rigidity}. Taylor's formula
therefore gives $|\Re u|=O(D_A+\delta)=O(\Psi(\delta))$.
The global $1$-point boundary supplies its sign. This proves the
first assertion, including the sharper ordinate bound.
The reverse assertion follows from that lemma, $D_E\le KD_A$,
and Lemma~\ref{p02:eul:local-critical}.

To prove sharpness, let $x\to\infty$ and take the real twist
\[
 z^{(x)}_2=-1,\quad z^{(x)}_p=1\ (3\le p\le x),\quad
 z^{(x)}_p=-1\ (p>x).
\]
These products converge locally uniformly, with all derivatives, to
$Z_{z^*}$ on $\Re s>1$. At the critical boundary,
\begin{equation}\label{p02:eul:sharp-critical}
 Q_{z^{(x)}}(E_*)
 =-2\sum_{p>x}\frac{(\log p)p^{-E_*}}{1-p^{-2E_*}}
 \sim-\frac{2}{E_*-1}x^{1-E_*}.
\end{equation}
At the $1$-point boundary,
\begin{equation}\label{p02:eul:sharp-one}
 Z_{z^{(x)}}(A_*)
 =\prod_{p>x}\frac{1-p^{-A_*}}{1+p^{-A_*}},\qquad
 Z_{z^{(x)}}(A_*)-1\sim-
              \frac{2}{A_*-1}\frac{x^{1-A_*}}{\log x}.
\end{equation}
The prime number theorem gives both comparisons. Because the two
frozen derivatives are real and strictly negative, the unique nearby
zeros are real, lie to the left, and have gaps
\begin{equation}\label{p02:eul:sharp-gaps}
 \delta_x\asymp x^{1-E_*},\qquad
 \eps_x\asymp\frac{x^{1-A_*}}{\log x}
            \asymp\Psi(\delta_x).
\end{equation}
More precisely, if
$c_E=2/((E_*-1)|a_E|)$ and
$c_A=2/((A_*-1)|a_A|)$, then
$\delta_x\sim c_Ex^{1-E_*}$ and
$\eps_x\sim c_Ax^{1-A_*}/\log x$, whence
$\eps_x/\Psi(\delta_x)\to c_A(E_*-1)c_E^{-\theta}>0$.
Reality follows alternatively from conjugation and local uniqueness.
Choose fixed small disks around $A_*$ and $E_*$ containing just the
respective frozen zeros. For all sufficiently large $x$, the nearby
zeros in these disks remain unique, by Rouch\'e's theorem.

Finally every fixed Euler-product twist is a locally uniform limit
of vertical translates of $\zeta$ in $\Re s>1$. To check this
directly, truncate the logarithmic Euler series so that its tail and
the required derivative tails are uniformly small on the two disks,
and use the rational independence of the finitely many $\log p$
to approximate the remaining prime phases. Such approximations occur
at arbitrarily large ordinates by Kronecker's theorem. For each member
of a sequence $x_n\to\infty$, approximate $Z_{z^{(x_n)}}$ and its
derivative closely enough that both local zeros move by $o(\delta_{x_n})$.
The derivatives at these zeros stay bounded away from zero, so this
follows from Rouch\'e's theorem, or the local inverse estimate.
The resulting ordinary-zeta points retain \eqref{p02:eul:sharp-gaps},
their ordinate difference tends to zero, and the local $1$-point is
unique in a fixed smaller disk. Passing to a subsequence makes the
critical gaps decrease strictly to zero. This proves the final assertion.
\end{proof}

\begin{remark}[What the sharpness assertion measures]
The construction above concerns the real-part transfer
$\delta\mapsto\delta^\theta/\log(1/\delta)$.
Proposition~\ref{p02:eul:vertical-sharpness} below treats the other rates.
The almost-sure assertion \eqref{p02:eul:waiting} does not provide a deterministic
first-occurrence lower bound for ordinary $\zeta'$.
\end{remark}


\subsection{Optimal ordinate estimates and the reverse horizontal rate}

The preceding tail construction concerns the forward horizontal
modulus. A different deformation proves sharpness of the remaining
three estimates.

\begin{proposition}[One-prime sharpness]\label{p02:eul:vertical-sharpness}
The square-root ordinate exponents in both directions of
Theorem~\ref{p02:eul:locking}, and the linear horizontal bound from a
$1$-point to a critical point, cannot be improved uniformly. Each
assertion is sharp along sequences for ordinary $\zeta$, with the
counterparts chosen as the unique roots in fixed local disks.
\end{proposition}
\begin{proof}
Fix a sufficiently large odd prime $P$. Starting from $z^*$, change
only $z_P$ to $e^{iu}$ for real $u$ near zero. The holomorphic
implicit-function theorem gives analytic functions $h_A(u),h_E(u)$
such that $A_*+h_A(u)$ is a $1$-point and $E_*+h_E(u)$ is a critical
point of the perturbed product, with $h_A(0)=h_E(0)=0$.
Let $a_A<0$ and $a_E<0$ be the frozen derivatives from the preceding
proofs. Differentiating in $u$ gives
\begin{equation}\label{p02:eul:phase-derivatives}
 h_A'(0)=\frac{-i}{a_A(P^{A_*}-1)},\qquad
 h_E'(0)=\frac{-i(\log P)P^{E_*}}
                   {a_E(P^{E_*}-1)^2}.
\end{equation}
Write their positive imaginary parts as $b_A(P)$ and $b_E(P)$.
Since $E_*>A_*$,
\[
 b_A(P)\sim P^{-A_*}/|a_A|,\qquad
 b_E(P)\sim(\log P)P^{-E_*}/|a_E|=o(b_A(P)).
\]
Thus $b=b_A(P)-b_E(P)>0$ after fixing $P$ large enough.
Conjugation gives $h_j(-u)=\overline{h_j(u)}$, so the real parts
are even and the imaginary parts odd.

Put $\delta_j(u)=-\Re h_j(u)$ and $t_j(u)=\Im h_j(u)$ for
$j=A,E$. The realization lemmas give $0\le\delta_j(u)=O(u^2)$.
For the reverse bound, recenter at $t_E(u)$. The prime-$2$ phase is
$-e^{-it_E(u)\log2}$, whose contribution to the critical energy is
\[
 (\log2)2^{-E_*}|1-e^{-it_E(u)\log2}|^2
 \sim (\log2)^3 2^{-E_*}b_E(P)^2u^2.
\]
The necessary-energy inequality bounds this by $C\delta_E(u)$.
The identical argument with the $1$-point energy at height $t_A(u)$
gives $\delta_A(u)\ge cu^2$. At zero deficit the necessary identities
still hold with zero right side, so there is no exception in this
argument. Analyticity and evenness now give positive constants
$k_A,k_E$ such that
\[
 \delta_A(u)=k_Au^2+O(u^4),\qquad
 \delta_E(u)=k_Eu^2+O(u^4),\qquad
 t_A(u)-t_E(u)=bu+O(u^3).
\]
Therefore the ordinate difference is comparable to the square root
of either deficit, and the two deficits are comparable to each other.

To realize these asymptotics for ordinary $\zeta$, take $u_n\to0$
and approximate each fixed twist by a vertical translate of $\zeta$
on slightly enlarged fixed disks around $A_*$ and $E_*$, with error
$o(u_n^2)$. Finite-coordinate Kronecker approximation and absolute
convergence permit this at arbitrarily large heights, exactly as in
the preceding sharpness proof. Cauchy's formula also controls the
derivatives. Uniform simple-root stability moves the two roots by
$o(u_n^2)$, preserving all three comparisons and local uniqueness.
\end{proof}


\subsection{A finite certificate for the cross-weight inequality}

\begin{lemma}\label{p02:lem:edge-enclosures}
The edges defined in \eqref{p02:eul:critical-edge} and
\eqref{p02:eul:one-edge} satisfy
\begin{equation}\label{p02:eq:edge-enclosures}
 1.93<A_*<2<E_*<2.82,
 \qquad 2A_*-E_*>1.04.
\end{equation}
\end{lemma}
\begin{proof}
For $a=193/100$, the finite positive sum gives
\[
 \sum_{\substack{3\le n\le499\\n\text{ odd}}}(2/n)^a
 >1.0054767>1.
\]
The decreasing-ratio characterization of $A_*$ implies $A_*>a$.
At $s=2$, the odd sum $M(2)=\pi^2/8-1<1/4$, since $\pi^2<10$;
hence $A_*<2$. The terms $p=3,5$ alone give
$\Delta(2)>0$ and therefore $E_*>2$.

For the remaining upper bound put $s=141/50$ and $u=s-1$. Direct
outward enclosures of the following finite expressions give
\begin{align*}
 P_{200}&:=\sum_{3\le p\le200}\frac{\log p}{p^s-1}
                       <0.085425810,\\
 T_{200}&:=\frac{200^{1-s}}{1-201^{-s}}
       \left(\frac{\log200}{u}+\frac1{u^2}\right)
                       <0.000208471,\\
 H_2&:=\frac{\log2}{2^s+1}>0.085981085.
\end{align*}
The positive decreasing function $(\log x)x^{-s}$ and the integral
test show
\[
 \sum_{p>200}\frac{\log p}{p^s-1}
 \le\frac1{1-201^{-s}}\int_{200}^\infty(\log x)x^{-s}\,dx
 =T_{200}.
\]
Thus $\Delta(s)<-0.000346804<0$, proving $E_*<2.82$.

The displayed numerical inequalities are only coarse finite
certificates, not assumed high-precision values of either edge.
They can be verified using rational interval arithmetic as follows.
Reduce logarithms by powers of $2$ to arguments in $[1,2]$ and use
\[
 \log y=2\sum_{j=0}^{N}\frac{t^{2j+1}}{2j+1}+R_N,
 \quad t=\frac{y-1}{y+1},\quad
 0\le R_N\le\frac{2t^{2N+3}}{(2N+3)(1-t^2)}.
\]
Taking $N=80$ suffices, since $0\le t\le1/3$. For exponentials of
nonnegative arguments $v\le18$, use the Taylor polynomial through
degree $256$ and the tail bound
\[
 0\le e^v-\sum_{j=0}^{256}\frac{v^j}{j!}
 \le\frac{v^{257}}{257!}\frac1{1-v/258}.
\]
Use reciprocals for negative arguments and propagate endpoints
outward. These errors are much smaller than the stated decimal
margins. Finally $2(1.93)-2.82=1.04$, proving the last inequality.
\end{proof}


\section{Small random perturbations and the matching recurrence scale}
\label{p02:sec:generic}

Let $\mathbb P$ denote the full set of primes. Fix $A>0$ and
$0<\rho<\log2/2$. Consider the compact product box
\[
 \mathcal B=\prod_{p\in\mathbb P}
 [\log p-\rho p^{-A},\log p+\rho p^{-A}],
 \qquad \lambda_p=\log p+\xi_p,
\]
with its product topology and the product probability law under which
the $\xi_p$ are independent and uniform on
$[-\rho p^{-A},\rho p^{-A}]$.
The radius $\rho$ may be arbitrarily small. We use either the full
frequencies $\nu_p=\lambda_p$, $p\in\mathbb P$, or the relative frequencies
$\nu_p=\lambda_p-\lambda_2$, $p\ge3$. The conclusions below hold
simultaneously for both models. Sums and products below range over
the coordinates of the chosen model. Write $\mathcal R_{w,\nu}$ for the
covering time associated with $\mathcal E_w$ and these frequencies.

\begin{theorem}[Matching scale on one full-measure frequency set]
\label{p02:thm:generic}
There is a set $\mathcal B_{\mathrm{good}}$ of full product probability
such that the following assertion holds for every
$\lambda\in\mathcal B_{\mathrm{good}}$, simultaneously for every
$\sigma>1$, every $\kappa\in\mathbb R$, and every positive weight
sequence satisfying
\[
 w_p\asymp p^{-\sigma}(\log p)^\kappa.
\]
There are $c,C>0$ such that, for all sufficiently small $\eps$,
\begin{equation}\label{p02:eq:generic-scale}
 \exp\bigl(cY_{\sigma,\kappa}(\eps)\bigr)
 \le\mathcal R_{w,\nu}(\eps)
 \le\exp\bigl(CY_{\sigma,\kappa}(\eps)\bigr).
\end{equation}
The upper bound is uniform in the starting phase, target phase, and
initial time. In particular, with $\beta=(\sigma-1)^{-1}$ and
$\ell_\eps=\log(1/\eps)$,
\begin{equation}\label{p02:eq:generic-second-order}
 \log\log\mathcal R_{w,\nu}(\eps)
 =\beta\ell_\eps+
   \bigl((\kappa-1)\beta-1\bigr)\log\ell_\eps+O(1).
\end{equation}
Constants and the small-$\eps$ threshold may depend on the fixed
weights and frequencies. No assertion of a limiting leading constant
for $\log\mathcal R_{w,\nu}$ is made.
The same estimates hold for the supremum of energy-avoidance interval
lengths along every individual orbit.
\end{theorem}

The phrase ``simultaneously'' is material: the exceptional frequency
set is chosen before the real parameters $\sigma,\kappa$ or the weight
sequence. This permits applications in which the weight exponent
itself depends on the realized frequencies.

\begin{lemma}[A countable family of frequency boxes]
\label{p02:lem:generic-small-divisors}
There is a full-probability set of frequencies with the following
property. For each integer $M\ge2$, there is $D_M>0$ such that, at
every sufficiently large dyadic $x$, every nonzero integer vector
supported on the relevant primes $p\le x$ and satisfying
\[
 |k_p|\le 2\left\lceil
 (x/p)^{M/2}(\log x/\log p)^{M/2}\right\rceil+1
\]
obeys
\begin{equation}\label{p02:eq:generic-separation}
 |k\cdot\nu|\ge\exp(-D_Mx/\log x).
\end{equation}
Both frequency models satisfy this property on the same set.
\end{lemma}
\begin{proof}
The elementary prime estimates needed for counting these vectors are
\begin{align}
 \sum_{p\le x}\log(x/p)&=O(x/\log x),\label{p02:eq:generic-logsum}\\
 \sum_{p\le x}\log(\log x/\log p)
 &=\int_2^x\frac{\pi(t)}{t\log t}\,\mathrm dt
 =O(x/\log^2x).\label{p02:eq:generic-loglogsum}
\end{align}
The equalities follow by summing the relevant integrals; the estimates
follow from the prime number theorem. Removing the prime $2$ changes
neither order estimate. Consequently the number of admitted vectors
is at most $\exp(C_Mx/\log x)$.

For any one nonzero vector, choose a coordinate $p$ with $k_p\ne0$.
Condition on every random coordinate except $\xi_p$. The linear form
$k\cdot\nu$ is then uniform on an interval of length
$2|k_p|\rho p^{-A}$. This remains true in the relative model:
$p\ge3$ there, and the common coordinate $\xi_2$ is among those
conditioned upon. Thus, for $u>0$,
\[
 \mathbb P\{|k\cdot\nu|<u\}
 \le\rho^{-1}x^Au.
\]
A union bound makes the probability of failure of
\eqref{p02:eq:generic-separation} at most
\[
 \rho^{-1}x^A\exp\bigl(-(D_M-C_M)x/\log x\bigr).
\]
Choose $D_M>C_M+1$. These probabilities are summable for $x=2^n$.
The first Borel--Cantelli lemma, followed by intersection over the
countably many integers $M$ and the two models, proves the assertion.
No independence between distinct linear forms is required.
\end{proof}

\begin{proof}[Proof of Theorem~\ref{p02:thm:generic}]
Fix frequencies furnished by Lemma~\ref{p02:lem:generic-small-divisors}.
At a dyadic truncation $x$ choose squared Fej\'er kernels with degrees
allocated through
\begin{equation}\label{p02:eq:generic-degrees}
 m_p=\left\lceil
 (x/p)^{\sigma/2}(\log p/\log x)^{\kappa/2}
 \right\rceil,
 \qquad K=\prod_{p\le x}J_{m_p}.
\end{equation}
The quantity inside the ceiling is allowed to be smaller than one.
The identities for $J_m$ give
\begin{align*}
 m:=\int \mathcal E_{w,\le x}K
 &=\sum_{p\le x}\frac{3w_p}{2m_p^2+1}\\
 &\ll x^{-\sigma}(\log x)^\kappa\pi(x)
 \ll x^{1-\sigma}(\log x)^{\kappa-1}.
\end{align*}
Also $\log\|K\|_A=O_{\sigma,\kappa}(x/\log x)$.
Indeed, if $h_p$ denotes the quantity rounded in
\eqref{p02:eq:generic-degrees}, then
\[
 \log\lceil h_p\rceil\le\log2+(\log h_p)_+
 \le\log2+\frac\sigma2\log(x/p)
       +\frac{|\kappa|}2\log(\log x/\log p),
\]
and \eqref{p02:eq:generic-logsum}--\eqref{p02:eq:generic-loglogsum} apply.

Choose an integer $M\ge\max\{2,\sigma,|\kappa|\}$. Every frequency
of $(2m-\mathcal E_{w,\le x})K$ lies in the box of
Lemma~\ref{p02:lem:generic-small-divisors}. Its nonzero frequencies
therefore have separation at least $\exp(-D_Mx/\log x)$.
The finite-frequency certificate, Lemma~\ref{p02:lem:certificate}, shows
that every interval of length $\exp(Cx/\log x)$ contains a time with
$\mathcal E_{w,\le x}<2m$. To see that the factor $1/m$ in the certificate
causes no loss, a single fixed prime gives
$m\gg x^{-\sigma}(\log x)^\kappa$ for large $x$, so
$\log(1/m)=O(\log x)$.

The uncontrolled tail contributes at most
\[
 2\sum_{p>x}w_p\ll x^{1-\sigma}(\log x)^{\kappa-1}.
\]
Take the least dyadic $x\ge K_1X_{\sigma,\kappa}(\eps)$, with a
sufficiently large fixed $K_1$. Since $x$ lies between $K_1X$ and
$2K_1X$, the total energy is less than $\eps$ and the elapsed time is
at most $\exp(CY_{\sigma,\kappa}(\eps))$. Target translation changes
Fourier coefficients only by unimodular factors. This proves the
upper bound with all the stated uniformities.

For the lower bound, every frequency vector in $\mathcal B$ satisfies
\[
 \sum_p w_p|\nu_p|<\infty.
\]
Thus $\mathcal E_w(\Phi_t\eta)$ is uniformly Lipschitz in $t$. A time mesh of
size proportional to $\eps$, Haar invariance, and the upper
small-ball estimate in Theorem~\ref{p02:thm:smallball} give
\[
 \mu\left\{\min_{0\le t\le T}\mathcal E_w(\Phi_t\eta)\le\eps\right\}
 \le C(1+T/\eps)\exp\bigl(-cY_{\sigma,\kappa}(\eps)\bigr).
\]
For $T=\exp(c_1Y_{\sigma,\kappa}(\eps))$ with $c_1$ small, this
probability is less than one. Some starting phase therefore avoids
the target for that long, proving the lower covering bound. Taking
two logarithms proves \eqref{p02:eq:generic-second-order}.
For completeness, the proved uniform covering bound makes every
orbit dense. The bad
segment just constructed has a strict energy margin on its compact
time interval. Continuity and the flow isometry preserve it in an
open neighborhood of its starting phase. Every orbit visits that
neighborhood, proving the same lower bound for each orbit's maximal
avoidance interval; the upper bound follows from the uniform covering
estimate.
\end{proof}

\begin{remark}
The deterministic prime-flow estimate paid the arithmetic cost
$\sum m_p\log p$. Here the almost-sure separation estimate pays
the logarithm of the number of Fourier modes, of order $x/\log x$.
This is why the bounds now match and why there is no restriction
$\sigma<3$. It is not a proof of the corresponding improvement for
the unperturbed prime logarithms.
\end{remark}

\section{Resonance obstructions and topologically typical slowness}
\label{p02:sec:residual}

Probability-one behavior and residual behavior in the same frequency
box are very different. The obstruction is already visible in one
integer character of a two-coordinate torus.

\begin{lemma}[Periodic exclusion intervals from one character]
\label{p02:lem:character-obstruction}
Let $w_p>0$ be summable, let $k\ne0$ be a finite integer vector, and
put
\[
 Q_w(k)^2=\sum_p\frac{k_p^2}{w_p},\qquad
 v=k\cdot\nu,\qquad \chi_k(\eta)=\prod_p\eta_p^{k_p}.
\]
Suppose $v\ne0$ and $\eps Q_w(k)^2\le1/8$. Every orbit has closed
intervals of length $\pi/|v|$ on which $\mathcal E_w>\eps$. These intervals
repeat with period $2\pi/|v|$. In particular,
\begin{equation}\label{p02:eq:dual-lower}
 \mathcal R_{w,\nu}(\eps)\ge\frac\pi{|k\cdot\nu|}.
\end{equation}
\end{lemma}
\begin{proof}
Choose principal arguments $\theta_p\in[-\pi,\pi]$. Since
$\theta_p^2\le(\pi^2/2)(1-\cos\theta_p)$, Cauchy--Schwarz gives
\[
 \left|\sum_p k_p\theta_p\right|
 \le\frac\pi{\sqrt2}\sqrt{\mathcal E_w(\eta)}\,Q_w(k).
\]
Hence $\mathcal E_w\le\eps$ forces $\chi_k(\eta)$ into the closed arc with
argument between $-\pi/4$ and $\pi/4$. Along an orbit this character
rotates with angular speed $v$:
$\chi_k(\Phi_t\eta)=\chi_k(\eta)e^{-itv}$.
Its passage through the opposite closed half-circle takes time
$\pi/|v|$ and excludes the target throughout. Choosing a starting
phase at the beginning of that passage proves \eqref{p02:eq:dual-lower}.
\end{proof}

\begin{theorem}[Residual slowness beyond any prescribed rate]
\label{p02:thm:residual}
Let $\Psi:(0,1)\longrightarrow[e,\infty)$ be any prescribed
finite-valued function. There is a dense $G_\delta$ subset
$\mathcal B_{\mathrm{slow}}(\Psi)$ of $\mathcal B$, consisting of
rationally independent frequency vectors, such that for every
$\lambda$ in this set and every positive summable weight sequence,
\begin{equation}\label{p02:eq:residual-slow}
 \limsup_{\eps\downarrow0}
 \frac{\log\mathcal R_{w,\nu}(\eps)}{\log\Psi(\eps)}=\infty.
\end{equation}
The same assertion holds when covering time is replaced by the
supremum of energy-avoidance interval lengths along any one orbit.
Both assertions hold for both the full and relative frequency models.
Each covering time at a fixed positive tolerance is nevertheless
finite.
\end{theorem}
\begin{proof}
First fix one of the two models and two of its coordinates: use
$\nu_2,\nu_3$ in the full model and $\nu_3,\nu_5$ in the relative
model. Denote them by $\alpha,\gamma$. For integers $q\ge1$ and
$p\in\mathbb Z$, put
\[
 N_{p,q}=p^2+q^2,\qquad e_{p,q}=N_{p,q}^{-2}.
\]
For each integer $n\ge2$, let $U_n$ be the union, over
$N_{p,q}\ge n$, of the relatively open slabs
\begin{equation}\label{p02:eq:resonance-slabs}
 |q\alpha-p\gamma|
 <\frac\pi{\Psi(e_{p,q})^n}.
\end{equation}
Every $U_n$ is dense in $\mathcal B$. Indeed, every nonempty basic
product neighborhood contains nondegenerate interior intervals for
the two chosen coordinates. In the relative model first hold
$\lambda_2$ fixed at an interior value and vary $\lambda_3,\lambda_5$.
This again supplies a rectangle of relative frequencies. Choose a
smaller rectangle on which $\gamma\ne0$. Rational slopes $p/q$ with
arbitrarily large denominators are dense, so a line
$q\alpha=p\gamma$ crosses the rectangle. Its open slab in
\eqref{p02:eq:resonance-slabs} intersects the neighborhood, however small
its width. This argument also covers neighborhoods touching endpoints
of the original closed intervals.

The compact metrizable product $\mathcal B$ is a Baire space. Thus
$\bigcap_{n\ge2}U_n$ is dense $G_\delta$. Intersect the sets obtained
for the two models and remove every nontrivial finite integer
resonance hyperplane for the full frequencies. The complements of
these countably many hyperplanes are open dense. The resulting set
is the required $\mathcal B_{\mathrm{slow}}(\Psi)$; full rational
independence implies relative rational independence as well.

For a frequency vector in this set, choose $(p_n,q_n)$ witnessing
membership in $U_n$. Then $N_n\ge n\to\infty$,
$e_n=N_n^{-2}\to0$, and $v_n=q_n\alpha-p_n\gamma\ne0$.
For any fixed positive weights on the two chosen coordinates,
\[
 e_n Q_w(k_n)^2
 \le\frac{\max\{w_{\alpha}^{-1},w_{\gamma}^{-1}\}}{N_n}
 \longrightarrow0,
\]
where $k_n=(q_n,-p_n)$ on those coordinates. Here the subscripts
$\alpha,\gamma$ refer to the chosen coordinate indices.
Lemma~\ref{p02:lem:character-obstruction} therefore applies eventually
and gives
\[
 \mathcal R_{w,\nu}(e_n)\ge\frac\pi{|v_n|}
 >\Psi(e_n)^n.
\]
It supplies intervals of the same length for every orbit. These
facts prove the limsup assertions. The frequency set was chosen
without reference to the weights; only the initial index from which
the dual condition holds depends on them.

Finally, rational independence makes every finite-frequency minimum
strictly positive. At any fixed tolerance choose a finite truncation
with sufficiently small tail, then choose sufficiently large kernel
degrees. Lemma~\ref{p02:lem:certificate} yields a finite covering time.
Thus the theorem describes arbitrarily slow dense motion, not a loss
of density.
\end{proof}

For each prescribed $\Psi$ the residual set can be chosen separately;
the theorem does not assert one set that outgrows every conceivable
function. Countably many prescribed rates can be imposed at once.
The underlying Liouville mechanism is classical in almost-periodic
dynamics. What is contrasted here is its residual strength with the
matching infinite-dimensional probability-one scale in
Theorem~\ref{p02:thm:generic}.


\section{Optimal positive kernels and the limitation of a height estimate}
\label{p02:sec:barrier}

The almost-sure small-divisor estimate removes a restriction that cannot be
removed merely by choosing a better positive localization kernel. We now
make that distinction precise. The optimization below concerns a Fourier
\emph{certificate cost}; it does not identify the covering time of the exact
prime flow when the weights decay faster than $p^{-3}$.

Throughout this section primes include $2$, and $w_p>0$ with
$\sum_p w_p<\infty$. For a finitely supported vector of nonnegative integers
$D=(D_p)$ put
\[
 B(D)=\sum_pD_p\log p,\qquad
 c(D_p)=1-\cos\frac{\pi}{D_p+2},\qquad
 M_w(D)=\sum_pw_pc(D_p).
\]
In particular, $c(0)=1$: a coordinate on which a kernel does not depend
has a uniform marginal, and contributes its full mean energy $w_p$.
Let $\mathcal K(D)$ be the nonnegative trigonometric polynomials of Haar
mean one whose degree in coordinate $p$ is at most $D_p$. Define
\begin{equation}\label{p02:eq:barrier-Psi}
 \Psi_w(B)=\inf_{B(D)\le B}M_w(D).
\end{equation}
The infimum can be taken over a finite set for each fixed $B$: a positive
$D_p$ requires $p\le e^B$, and every $D_p$ is bounded by $B/\log2$.

For clarity, $B(D)$ is a \emph{rectangular} Fourier budget. If $D_p$ are
the coordinate degrees of a polynomial $K$, its actual support height is
\[
 H(K)=\max_{k\in\operatorname{supp}\widehat K}
             \sum_p|k_p|\log p\le B(D).
\]
Equality need not hold for an arbitrary multivariate polynomial. It does
hold for the product extremizers constructed below. We will not identify
these two costs in a lower bound.

\subsection{Exact localization at prescribed coordinate degrees}

\begin{proposition}[Rectangular localization optimum]
\label{p02:prop:barrier-exact}
For every finitely supported bandwidth vector $D$,
\begin{equation}\label{p02:eq:barrier-exact}
 \inf_{K\in\mathcal K(D)}\int \mathcal E_w K\,d\mu
 =\sum_pw_p\left(1-\cos\frac{\pi}{D_p+2}\right)=M_w(D).
\end{equation}
A product of one-variable extremizers attains the infimum. Each factor
has nonnegative Fourier coefficients and Fourier coefficient norm at most
$D_p+1$.
\end{proposition}

\begin{proof}
The one-variable coefficient inequality used here is classical Fej\'er
extremality; see \cite{FejerMulti}. We include its short proof. A
nonnegative trigonometric polynomial $J$ of degree at most $d$ and mean one
has, by Fej\'er--Riesz factorization, a representation
\[
 J(\theta)=\left|\sum_{j=0}^d a_je^{ij\theta}\right|^2,
 \qquad \sum_{j=0}^d|a_j|^2=1.
\]
Consequently
\[
 \int\cos\theta\,J(\theta)\,d\theta
 =\operatorname{Re}\sum_{j=0}^{d-1}a_{j+1}\overline{a_j}
 \le\cos\frac{\pi}{d+2}.
\]
The last inequality is the Rayleigh bound for the $(d+1)$-vertex path
matrix with entries $1/2$ immediately above and below the diagonal. Its
largest eigenvalue is $\cos(\pi/(d+2))$, with normalized eigenvector
\[
 a_j=\sqrt{\frac2{d+2}}\sin\frac{(j+1)\pi}{d+2},
 \qquad 0\le j\le d.
\]
These positive coefficients give an extremizer $J_d$. All Fourier
coefficients of $J_d$ are nonnegative, and Cauchy--Schwarz gives
\[
 \|J_d\|_A=J_d(0)
 =\left(\sum_{j=0}^d a_j\right)^2\le d+1.
\]
For $d=0$, take $J_0=1$.

For a multivariate $K$, integrate out every coordinate except $p$.
The resulting marginal is nonnegative, has mean one, and has degree at
most $D_p$. Applying the preceding inequality to each marginal and
summing the nonnegative energy terms proves the lower bound in
\eqref{p02:eq:barrier-exact}. The product
$\prod_{D_p>0}J_{D_p}$ attains every marginal bound simultaneously.
\end{proof}

Thus allowing correlations among the coordinates cannot improve this
mean-energy objective when the rectangular degree bounds are fixed.
This observation concerns arbitrary nonnegative kernels, not just tensor
products or squared Fej\'er kernels.

\begin{lemma}[Budget inequality]\label{p02:lem:barrier-holder}
For a finite set of primes $Q$ set
\[
 L_Q=\sum_{p\in Q}w_p^{1/3}(\log p)^{2/3},\qquad
 \vartheta(Q)=\sum_{p\in Q}\log p.
\]
If $B(D)\le B$, then
\begin{equation}\label{p02:eq:barrier-holder}
 M_w(D)\ge\frac{2L_Q^3}{(B+2\vartheta(Q))^2}.
\end{equation}
\end{lemma}
\begin{proof}
For $d\ge0$, elementary cosine bounds give
\begin{equation}\label{p02:eq:barrier-cbounds}
 \frac2{(d+2)^2}\le c(d)\le\frac{\pi^2}{2(d+2)^2}.
\end{equation}
H\"older's inequality, applied to
\[
 w_p^{1/3}(\log p)^{2/3}
 =\left(\frac{w_p}{(D_p+2)^2}\right)^{1/3}
  \bigl((D_p+2)\log p\bigr)^{2/3},
\]
shows that
\[
 L_Q^3\le
 \left(\sum_{p\in Q}\frac{w_p}{(D_p+2)^2}\right)
 \left(\sum_{p\in Q}(D_p+2)\log p\right)^2.
\]
Now use \eqref{p02:eq:barrier-cbounds} and the budget constraint.
\end{proof}

\subsection{The optimal prime budget and its transition}

\begin{theorem}[Sharp rectangular cost]\label{p02:thm:barrier-cost}
Suppose $w_p\asymp p^{-\sigma}$ for a fixed $\sigma>1$. As $B\to\infty$,
\begin{equation}\label{p02:eq:barrier-orders}
 \Psi_w(B)\asymp
 \begin{cases}
 B^{1-\sigma}/\log B,&1<\sigma<3,\\
 B^{-2}(\log B)^2,&\sigma=3,\\
 B^{-2},&\sigma>3.
 \end{cases}
\end{equation}
More generally, if
\begin{equation}\label{p02:eq:barrier-L}
 L=\sum_pw_p^{1/3}(\log p)^{2/3}<\infty,
\end{equation}
then the exact limit is
\begin{equation}\label{p02:eq:barrier-constant}
 \lim_{B\to\infty}B^2\Psi_w(B)=\frac{\pi^2}{2}L^3.
\end{equation}
In particular, \eqref{p02:eq:barrier-constant} applies to every $\sigma>3$.
\end{theorem}

\begin{proof}
First consider the lower estimates in \eqref{p02:eq:barrier-orders}.
For $Q=\{p:p\le B\}$, prime partial summation gives
$\vartheta(Q)=O(B)$ and
\[
 L_Q\asymp
 \begin{cases}
 B^{1-\sigma/3}(\log B)^{-1/3},&1<\sigma<3,\\
 (\log B)^{2/3},&\sigma=3.
 \end{cases}
\]
Lemma~\ref{p02:lem:barrier-holder} gives the first two lower bounds.
For $\sigma>3$, its application to $Q=\{2\}$ already gives
$\Psi_w(B)\gg B^{-2}$.

For the upper bound when $1<\sigma<3$, choose $x=cB$ with a sufficiently
small fixed $c>0$, and set
\[
 D_p=\left\lceil(x/p)^{\sigma/3}\right\rceil\quad(p\le x),
 \qquad D_p=0\quad(p>x).
\]
Partial summation shows
\[
 B(D)\ll x^{\sigma/3}\sum_{p\le x}p^{-\sigma/3}\log p
          +\vartheta(x)\ll x.
\]
Thus $B(D)\le B$ after fixing $c$. Also,
\[
 M_w(D)\ll
 x^{-2\sigma/3}\sum_{p\le x}p^{-\sigma/3}
 +\sum_{p>x}p^{-\sigma}
 \ll\frac{x^{1-\sigma}}{\log x},
\]
which gives the required upper estimate.

At $\sigma=3$, put
\[
 A=\frac{B}{C(\log B)^{2/3}},\qquad
 D_p=\left\lfloor\frac{A}{p(\log p)^{1/3}}\right\rfloor,
\]
where $C$ is sufficiently large. The nonzero degrees occur for
$p\le x$, where $x(\log x)^{1/3}=A$; hence
$x\asymp B/\log B$. The budget satisfies
\[
 B(D)\le A\sum_{p\le x}\frac{(\log p)^{2/3}}p
 \ll A(\log x)^{2/3}\le B.
\]
Since $D_p+1\ge A/[p(\log p)^{1/3}]$ on the active coordinates,
\begin{align*}
 M_w(D)
 &\ll A^{-2}\sum_{p\le x}\frac{(\log p)^{2/3}}p
               +\sum_{p>x}p^{-3}\\
 &\ll A^{-2}(\log x)^{2/3}+\frac{x^{-2}}{\log x}
 \ll\frac{(\log B)^2}{B^2}.
\end{align*}
The inactive tail is only $O(B^{-2}\log B)$.

It remains to prove \eqref{p02:eq:barrier-constant}, which also supplies the
last upper bound. For its lower bound fix a finite set $Q$. Along any
sequence with $B^2M_w(D)$ bounded, every $D_p$ with $p\in Q$ tends to
infinity. Otherwise a fixed positive summand of $M_w(D)$ would remain.
Thus, uniformly on this fixed set,
\[
 c(D_p)=\frac{\pi^2/2+o(1)}{(D_p+2)^2}.
\]
The H\"older argument of Lemma~\ref{p02:lem:barrier-holder} gives
\[
 \liminf_{B\to\infty}B^2M_w(D)\ge\frac{\pi^2}{2}L_Q^3.
\]
Sequences with unbounded scaled energy do not weaken the lower bound.
Letting $Q$ increase to all primes proves the desired lower limit.

For the matching construction set
\[
 t_p=w_p^{1/3}(\log p)^{-1/3},\qquad
 D_p=\left\lfloor\frac{Bt_p}{L}\right\rfloor.
\]
Since $t_p\to0$, this vector has finite support, and $B(D)\le B$.
For each fixed prime,
\[
 B^2w_pc(D_p)\longrightarrow
 \frac{\pi^2}{2}L^2w_p^{1/3}(\log p)^{2/3}.
\]
Furthermore, $D_p+1\ge Bt_p/L$ and \eqref{p02:eq:barrier-cbounds} imply
\[
 B^2w_pc(D_p)\le C L^2w_p^{1/3}(\log p)^{2/3}.
\]
The right side is summable by \eqref{p02:eq:barrier-L}. Dominated convergence
therefore gives the upper limit $(\pi^2/2)L^3$.
\end{proof}

\subsection{Consequences for the unperturbed prime flow}

The exact localization optimum can be inserted into the same elementary
arithmetic averaging argument. Here $\mathcal R_{w,\log p}$ denotes the
covering time for frequencies $\lambda_p=\log p$.

\begin{corollary}[Endpoint and faster-decay upper bounds]
\label{p02:cor:barrier-recurrence}
For $w_p\asymp p^{-\sigma}$, $\sigma>1$, put
$\ell=\log(1/\varepsilon)$. For all sufficiently small $\varepsilon$,
\begin{equation}\label{p02:eq:barrier-Rupper}
 \log\mathcal R_{w,\log p}(\varepsilon)\ll
 \begin{cases}
 \varepsilon^{-1/(\sigma-1)}\ell^{-1/(\sigma-1)},&1<\sigma<3,\\
 \varepsilon^{-1/2}\ell,&\sigma=3,\\
 \varepsilon^{-1/2},&\sigma>3.
 \end{cases}
\end{equation}
In particular, at $\sigma=3$,
\begin{equation}\label{p02:eq:barrier-endpoint}
 e^{c\varepsilon^{-1/2}\ell^{-3/2}}
 \le\mathcal R_{w,\log p}(\varepsilon)
 \le e^{C\varepsilon^{-1/2}\ell},\qquad
 \lim_{\varepsilon\downarrow0}
 \frac{\log\log\mathcal R_{w,\log p}(\varepsilon)}{\ell}
 =\frac12.
\end{equation}
The same conclusions hold for the relative frequencies $\log(p/2)$,
with the prime $2$ omitted.
\end{corollary}

\begin{proof}
Choose $D$ with $M_w(D)\le\varepsilon/4$ and let
$Q=\{p:D_p>0\}$. For small $\varepsilon$, this set is nonempty.
Use the product extremizer $K$ from Proposition~\ref{p02:prop:barrier-exact}.
Its truncated energy mean is at most $\varepsilon/4$, while
\[
 \sum_{p\notin Q}w_p\le M_w(D)\le\varepsilon/4.
\]
The omitted coordinates therefore contribute at most $\varepsilon/2$
pointwise. The certificate polynomial
\[
 P=(\varepsilon/2-\mathcal E_Q)K
\]
has mean at least $\varepsilon/4$ and coordinate degrees at most
$D_p+1$. Its arithmetic degree budget is at most
\[
 \sum_{p\in Q}(D_p+1)\log p\le2B(D),
\]
since each active $D_p\ge1$. Moreover
\[
 \log\|K\|_A
 \le\sum_{p\in Q}\log(D_p+1)
 \le\sum_pD_p\le B(D)/\log2.
\]
A nonzero frequency $\sum k_p\log p$ with these coordinate bounds is
the logarithm of a nonunit positive rational whose numerator and
denominator are at most $\exp(2B(D))$. Its absolute value is therefore
at least $\exp(-2B(D))$. The averaging certificate of
Lemma~\ref{p02:lem:certificate}, or direct integration of each nonconstant
Fourier term, now gives a positive average of $P$ on every interval of
length $\exp(CB(D))$. Here $\log(1/\varepsilon)$ is absorbed into
$B(D)$; indeed its single-prime localization bound already forces
$B(D)\gg\varepsilon^{-1/2}$. Positivity of $P$ gives
$\mathcal E_Q<\varepsilon/2$, and hence $\mathcal E_w<\varepsilon$.

Inverting the three orders in Theorem~\ref{p02:thm:barrier-cost} proves
\eqref{p02:eq:barrier-Rupper}. The lower bound in
\eqref{p02:eq:barrier-endpoint} is the small-ball avoidance consequence of
Theorem~\ref{p02:thm:smallball} at $\sigma=3$; taking two logarithms gives
the limit. For relative frequencies, clearing the additional power of
$2$ changes the elementary frequency estimate by only a fixed factor
in its exponent. Deleting one coordinate changes none of the orders.
\end{proof}

The local analytic transfer of phase energy to an edge zero needs only
$S>1$. Consequently, if an unperturbed Dirichlet edge occurs at $S=3$,
\eqref{p02:eq:barrier-endpoint} supplies the recurrence endpoint, with its corresponding zero-gap exponent $1/2$.

\subsection{What is, and is not, obstructed}

For the linear energy certificate $(a-\mathcal E_Q)K$,
Theorem~\ref{p02:thm:barrier-cost} shows that a better positive polynomial
cannot reduce the \emph{rectangular localization budget} below the
stated orders. In particular, for $\sigma>3$, a proof that pays
$\exp(CB(D))$ through the elementary arithmetic height estimate cannot
reach the entropy exponent $1/(\sigma-1)<1/2$ merely by changing its
positive kernel. At $\sigma=3$, the factor $\log(1/\varepsilon)$ in the
optimal rectangular budget is also unavoidable within that optimization.

There is a weaker obstruction using actual joint support height, which
requires no rectangular comparison. If $K\ge0$, $\int K=1$, and
$\int \mathcal E_wK\le\varepsilon$, its prime-$2$ marginal must have degree
$d\gg\varepsilon^{-1/2}$ by the univariate extremum. A nonzero top
coefficient of that marginal is a pure Fourier mode $k=d e_2$ of $K$.
Thus
\begin{equation}\label{p02:eq:barrier-actual}
 H(K)\ge d\log2\gg\varepsilon^{-1/2}.
\end{equation}
Accordingly, using only a worst-case separation estimate of the form
$\exp(-C H(K))$ retains this power barrier even without replacing the
actual support by a rectangle. Equation~\eqref{p02:eq:barrier-actual} does
\emph{not} establish the stronger logarithmic factor at $\sigma=3$
for arbitrary joint supports.

None of these statements is a dynamical lower bound of exponent $1/2$
when $\sigma>3$. The true lower bound for the exact prime flow remains
the entropy bound with exponent $1/(\sigma-1)$. A better arithmetic
estimate, a more selective Fourier averaging argument, or another
recurrence mechanism could improve the upper bound. The generic
frequency theorem does precisely the first of these: it controls the
small divisors by the number of modes, rather than by their elementary
rational height.


\section{Scope and further questions}

The sharp locking theorem compares two different boundary defects of
one Euler product. Its two constructions separate the obstructions:
changing a distant prime tail controls the optimal forward horizontal
modulus, whereas changing one fixed prime controls the optimal
ordinate exponents and the reverse horizontal rate. The ordinary
zeta function inherits both obstructions through its vertical
approximation of the phase hull.

The remaining recurrence questions are arithmetic. For exact prime
frequencies and $1<\sigma<3$, the established bounds are
\[
 cY_{\sigma,\kappa}(\eps)
 \le\log\mathcal R_{w,\log p}(\eps)
 \le CX_{\sigma,\kappa}(\eps).
\]
The ratio of these two scales is $\log(1/\eps)$. Closing this gap
would sharpen the critical-point gap theorem beyond its present
double-logarithmic exponent. The almost-sure perturbation theorem
shows that the weighted geometry does not require that loss; it gives
no improved small-divisor estimate at the exact prime vector.

For $w_p\asymp p^{-\sigma}$ with $\sigma>3$, the present
unperturbed bounds only imply
\[
 \frac1{\sigma-1}\le
 \liminf_{\eps\downarrow0}
 \frac{\log\log\mathcal R_{w,\log p}(\eps)}{\log(1/\eps)}
 \le
 \limsup_{\eps\downarrow0}
 \frac{\log\log\mathcal R_{w,\log p}(\eps)}{\log(1/\eps)}
 \le\frac12.
\]
The exact kernel optimization explains why the specified arithmetic
certificate stops at this upper exponent. Determining the actual
exponent requires an improvement in the arithmetic analysis or a
different recurrence argument. At $\sigma=3$ the two power exponents
coincide, although logarithmic factors remain separated.

Finally, maximal gaps and first visits from one prescribed origin are
different quantities. The long-hole construction locates large gaps
on every deterministic orbit, but it does not force such a gap to
begin at ordinate zero. The almost-sure first-waiting result therefore
does not establish a lower waiting exponent for the untwisted zeta
function at a fixed origin. This specific problem remains open in the
present work.


\clearpage

% ===== Integrated Part 2 =====
\part{Phase-localized zeta zeros and Euler-edge laws}
\label{part:zeta-localization}The recurrence theory of the preceding part concerns the prime flow
at arbitrary heights. We now restrict those heights to ordinates of
simple critical-line zeros. The first step transports a Weil-form
counting certificate to fixed phase regions; the second makes that
transport uniform for shrinking boxes and for a growing prime set.
The resulting phase measures support both quantitative Euler-edge
companions and conditional laws for their shape. The final sections
treat inverse zero detection and the separate matrix constraints on
higher-moment amplification.

\section{Fixed phase regions and the counting constant}

Write $N(T,2T)$ for the number of nontrivial zeta zeros with
$T<\Im\rho\le2T$, counted with multiplicity. For a finite prime set
$S$, let $F_S(t)=(p^{-it})_{p\in S}$ and let $\mu_S$ be normalized
Haar measure on $\T^S$. Put
\[
 c_{\rm MT}=\frac32-\frac1{\sqrt2}\cot\frac1{\sqrt2}.
\]

\begin{theorem}[Prime-phase localization]\label{p03:thm:main-phase}
For every fixed finite $S$ and every Haar-Jordan-measurable
$B\subset\T^S$,
\[
 \liminf_{T\to\infty}
 \frac{\#\{\gamma\in(T,2T]:\zeta(\tfrac12+i\gamma)=0,
 \ \zeta'(\tfrac12+i\gamma)\ne0,\ F_S(\gamma)\in B\}}
 {N(T,2T)}\ge c_{\rm MT}\mu_S(B).
\]
No Riemann hypothesis or fixed zero-free half-plane is assumed.
The set $S$ and the region $B$ are fixed before $T$ tends to infinity.
\end{theorem}

The assertion is a lower domination statement for the simple
critical-line submeasure. It neither identifies its limiting density
nor asserts equidistribution of all simple zeros. Classical
all-zero phase distribution is discussed in
\cite{AkbaryMurty,FordZaharescu,FordMengZaharescu}; the
subsequence equidistribution results of \c{C}i\c{c}ek--Gonek
\cite{CicekGonek} use RH and a spacing hypothesis.
The rank--inertia viewpoint and the unconditioned
Montgomery--Taylor certificate are from Alp\"oge--Furman
\cite{AlpogeFurman}. Lamzouri \cite{Lamzouri} gives another
proof, and Knausg\aa rd \cite{Knausgard} reports a larger global
constant. We use the specified second-moment constant throughout;
the present localization does not assert a new global percentage.

The fixed-region result is proved first without a zero-free
half-plane. Its shrinking-region extensions are established in
Sections~\ref{p04:sec:shrinking-boxes}--\ref{p04:sec:near-logarithmic}.
The simultaneous growing-prime limit in
Section~\ref{p04:sec:growing-primes} is distinct from the iterated
height-then-tolerance limits used for the Gibbs laws. The
$\mathrm H_\vartheta$ hypothesis occurs only in the labeled
half-plane transfer of Section~\ref{p03:sec:zero-free-transfer};
the Gaussian inverse and Vinogradov--Korobov arguments do not use it.



\section{The rank--inertia counting inequality}
\label{p03:sec:rank-basic}

For a Hermitian matrix $H$, let $n_+(H)$ count its strictly positive
eigenvalues. The following elementary version will be used for
localized zero counts.

\begin{lemma}\label{p03:lem:rank-basic}Suppose $G=P+Q$, where $P\succeq0$, $\rank P\le s$,
$\tr P\le s$, $n_+(Q)\le b$, and $s+2b\le M$. Then
\[
 s\ge 4\tr G-\tr G^2-2M.
\]
Additive $o(N)$ errors in the count and trace constraints change
the conclusion by only $o(N)$.
\end{lemma}
\begin{proof}
Let $\lambda_1\ge\cdots\ge\lambda_r>0$ be the positive
eigenvalues of $G$. Since $G\preceq P+Q_+$ and
$\rank Q_+\le b$, min--max gives
$\lambda_{b+j}(G)\le\lambda_j(P)$. Thus the tail after the first
$b$ positive eigenvalues has at most $s$ terms and sum at most $s$.
For $x\le0$, $4x-x^2\le0$; for all $x$, $4x-x^2\le4$;
and $4x-x^2\le2x+1$. Consequently
\[
 4\tr G-\tr G^2\le4b+3s\le2M+s.
\]
The same calculation tracks additive errors.
\end{proof}





\section{Prime-phase weights in the Weil second moment}\label{p03:sec:phase-arithmetic}

Fix a finite set $\mathcal P$ of primes and write
\[
 \Phi_{\mathcal P}(t)=(p^{-it})_{p\in\mathcal P},\qquad
 \mu_{\mathcal P}=\text{Haar probability measure on }\mathbb T^{\mathcal P}.
\]
The arithmetic observation behind our localization result is that a
nontrivial rational twist of the prime--prime term has no bulk diagonal.
Indeed, if two prime powers have a ratio supported on $\mathcal P$, then
either their ratio is one or both prime bases belong to $\mathcal P$.
We give the details, including the passage from Fourier weights to
arbitrary continuous phase weights.

Choose a nonzero, nonnegative, even function
$f\in C_c^\infty((-1/2,1/2))$ and put
\begin{equation}\label{p03:eq:phase-window}
 \psi=f^2,\quad a=\int_{\R}\psi(u)\,du,\quad
 L=\log T-D\log\log T,\quad X=e^L,\quad \phi(u)=f(u/L),
\end{equation}
where $D>4$ is fixed. Our Fourier convention in this section is
$\widehat h(t)=\int_{\R}h(u)e^{-itu}\,du$. Let
$I=[T,2T]$, $\alpha_k=T+2\pi k/L$, and retain the integers $k$ for
which $\alpha_k\in I$. Define the real symmetric Weil compression
\begin{equation}\label{p03:eq:phase-Weil-matrix}
 G_{k\ell}=\frac1{aL^2}\int_{\R}
 \widehat\phi(t-\alpha_k)\widehat\phi(t-\alpha_\ell)
 \nu_X(t)\,dt.
\end{equation}
Here the explicit-formula density is
\begin{align*}
 \nu_X(t)&=\mu(t)+\Pi_X(t)+P_X(t),\\
 \mu(t)&=\frac1{2\pi}\Re\frac{\Gamma'}{\Gamma}
                  \left(\frac14+\frac{it}2\right)
                 -\frac{\log\pi}{2\pi},\\
 \Pi_X(t)&=\frac1\pi\Re\frac{X^{1/2+it}}{1/2+it},\qquad
 P_X(t)=-\frac1\pi\sum_{n\le X}\frac{\Lambda(n)}{\sqrt n}
                                      \cos(t\log n).
\end{align*}
The compact support in \eqref{p03:eq:phase-window} makes
\eqref{p03:eq:phase-Weil-matrix} an exact compression of the zero-side
Weil form. We use its standard explicit-formula interpretation; the
weighted estimates below are proved directly on the real side.
Set
\begin{equation}\label{p03:eq:phase-R}
 R_\psi=\frac{\displaystyle\int\psi(u)^2\,du+
       \displaystyle\iint|u-v|\psi(u)\psi(v)\,du\,dv}{a^2},
 \qquad N_I=N(2T)-N(T).
\end{equation}

\begin{proposition}[Phase-weighted first and second moments]
\label{p03:prop:phase-moments}
For every fixed continuous real function
$K\in C(\mathbb T^{\mathcal P})$, put
$k(t)=K(\Phi_{\mathcal P}(t))$ and
$\overline K=\int K\,d\mu_{\mathcal P}$. Then
\begin{equation}\label{p03:eq:phase-moments}
 \sum_k k(\alpha_k)G_{kk}=(\overline K+o(1))N_I,
 \qquad
 \sum_k k(\alpha_k)\sum_\ell G_{k\ell}^2
       =(R_\psi\overline K+o(1))N_I.
\end{equation}
In particular, if $0\le K\le1$ and
$D_K=\operatorname{diag}(\sqrt{k(\alpha_k)})$, then
\begin{equation}\label{p03:eq:phase-compressed-moments}
 \operatorname{tr}(D_KGD_K)=(\overline K+o(1))N_I,
 \qquad
 \|D_KGD_K\|_{\mathrm{HS}}^2
       \le(R_\psi\overline K+o(1))N_I.
\end{equation}
All limits are taken with $\mathcal P,K,f,D$ fixed.
\end{proposition}

\begin{proof}
We separate the boundary estimates from the arithmetic computation.
Put
\[
 u_k(t)=\frac{\widehat\phi(t-\alpha_k)}{\sqrt a L}
 \quad(k\in\Z),
\]
and let $U:\ell^2(\Z)\longrightarrow L^2(\R)$ be their synthesis
operator. Parseval's formula and the Fourier-series Bessel inequality
give
\begin{equation}\label{p03:eq:phase-frame-bound}
 \|U\|^2\ll_f L^{-1}.
\end{equation}
Moreover, the Schwartz bounds for $\widehat f$ imply, for each fixed
$A>0$,
\begin{align}
 \sum_{\alpha_k\in I}|u_k(t)|^2
   &\ll_{A,f}(1+L\operatorname{dist}(t,I))^{-A}
                      &&(t\notin I),\label{p03:eq:phase-frame-tail1}\\
 \sum_{\alpha_k\notin I}|u_k(t)|^2
   &\ll_{A,f}(1+L\operatorname{dist}(t,\R\setminus I))^{-A}
                      &&(t\in I).\label{p03:eq:phase-frame-tail2}
\end{align}
The integrals of these tails are $O_f(L^{-1})$. On $|t|\le4T$ we
have $|\nu_X(t)|\ll L+\sqrt X$; outside this range its growth is at
most logarithmic in $|t|$, apart from the bounded prime polynomial.
The remote tails are therefore harmless. If $V$ is any subfamily of
$U$ and $J$ a measurable set, the elementary operator estimate
\[
 \|V^*M_{\nu_X1_J}V\|_{\mathrm{HS}}^2
 \le \|V\|^2\int_J|\nu_X(t)|^2\sum_k|v_k(t)|^2\,dt
\]
shows that restricting the integral defining $G$ to $I$ changes $G$
by a matrix of squared Hilbert--Schmidt norm
$O_f((L^2+X)/L^2)=o(N_I)$. Likewise, for the infinite matrix
$H=U^*M_{\nu_X1_I}U$, the squared Hilbert--Schmidt mass in rows or
columns whose centers lie outside $I$ is
$O_f((L^2+X)/L^2)$. Its full squared norm is $O_f(N_I)$ independently: the exact
Poisson identity gives $\sum_k|u_k(t)|^2=1$, while the standard
Dirichlet-polynomial mean-value bound gives
$\int_I|\nu_X(t)|^2dt\ll TL^2$; applying the preceding operator
estimate and \eqref{p03:eq:phase-frame-bound} yields
$\|H\|_{\mathrm{HS}}^2\ll TL$. Cauchy--Schwarz thus makes every
resulting cross-error $o(N_I)$. The linear trace boundary error is
$O_f((L+\sqrt X)/L)=o(N_I)$ by the same tail bounds. We may consequently
compute with all centers and the real integral restricted to $I$.

First take $k(t)=q^{it}$, where $q$ is a fixed positive rational
supported on $\mathcal P$, and put $\lambda=\log q$. Poisson
summation gives the exact identity
\begin{equation}\label{p03:eq:phase-weighted-Poisson}
 \sum_{k\in\Z}e^{i\lambda\alpha_k}
   \widehat\phi(t-\alpha_k)\widehat\phi(t'-\alpha_k)
 =Le^{i\lambda(t+t')/2}\Psi_\lambda(t-t'),
\end{equation}
where
\[
 \Psi_\lambda(s)=\int_{\R}
    \phi(u-\lambda/2)\phi(u+\lambda/2)e^{-isu}\,du.
\]
There are no nonzero aliases: $\operatorname{supp}f$ is a compact
subset of $(-1/2,1/2)$ and $\lambda$ is fixed. Write
$\Phi=\Psi_0=\widehat{\phi^2}$ and
$J_\lambda=\Psi_\lambda\Phi$. The weighted row energy is exactly
\begin{equation}\label{p03:eq:phase-twisted-integral}
 \frac1{a^2L^2}\iint_{I^2}
 e^{i\lambda(t+t')/2}J_\lambda(t-t')
                  \nu_X(t)\nu_X(t')\,dt\,dt'.
\end{equation}
Uniformly for fixed $\lambda$,
\begin{equation}\label{p03:eq:phase-lag-bounds}
 \int|J_\lambda(s)|\,ds\ll L,\qquad
 \int|sJ_\lambda(s)|\,ds\ll1,\qquad
 \sup_v\left|\int J_\lambda(s)e^{ivs}\,ds\right|\ll L.
\end{equation}
These follow by scaling and the Schwartz bounds for $f$.

Let $\ell=\log T$. Uniformly on $I$,
\[
 \nu_X(t)=\frac\ell{2\pi}+P_X(t)
                         +O(1+\sqrt X/T).
\]
The standard Dirichlet-polynomial mean-value inequality gives
$\int_I|P_X(t)|^2dt\ll TL^2+XL^2$. Consequently replacing
$\nu_X$ by $\ell/(2\pi)+P_X$ in
\eqref{p03:eq:phase-twisted-integral} costs $O(TL^2)=o(TL^3)$ before
normalization. In midpoint and lag coordinates, the midpoint interval
is $I$ with $|s|/2$ removed at each end. By
\eqref{p03:eq:phase-lag-bounds}, replacing it by $I$ costs
$O(L^2+X)=o(TL^3)$; the long-lag tail is smaller by Schwartz decay.

The constant--constant term is lower order if $q\ne1$. At $q=1$
it equals
\[
 \frac{T\ell^2}{2\pi}L\int\psi(u)^2\,du
\]
before normalization. For the prime--prime term, put
$a_n=\Lambda(n)/\sqrt n$. The midpoint frequencies are
\[
 \log q+\varepsilon\log n+\eta\log m,
                 \qquad \varepsilon,\eta\in\{-1,1\}.
\]
After coincident frequencies have been combined, the two frequency
lists in each product have separation at least $c_q/X$. The
Montgomery--Vaughan Hilbert inequality \cite{Hilbert}, applied to their union and
polarized, bounds all unequal-frequency terms by
\[
 O_q\left(X\sum_{n\le X}a_n^2\right)=O_q(XL^2),
\]
uniformly in the lag $s$. The coefficient phases introduced by $s$
have modulus one. Integration against $J_\lambda$ therefore gives
$O_q(XL^3)=o(TL^3)$. The separation assertion is elementary: after
writing $q=b/c$, a nonzero difference $bn-cm$ is an integer of
absolute value at least one; the same-sign cases are no worse.

If $q=1$, the exact resonances are $n=m$ and
$\varepsilon=-\eta$. Their contribution is
\begin{equation}\label{p03:eq:phase-diagonal-primes}
 \frac T\pi\sum_{n\le X}\frac{\Lambda(n)^2}{n}
                  (\phi^2*\phi^2)(\log n).
\end{equation}
Partial summation, using
$\sum_{n\le y}\Lambda(n)^2/n=\tfrac12\log^2y+O(\log y)$,
turns the sum in \eqref{p03:eq:phase-diagonal-primes} into
\[
 L^3\int_0^1v(\psi*\psi)(v)\,dv+O_f(L^2).
\]
If $q\ne1$, write a resonant pair as $n=p^j,m=r^h$. Every prime
base must belong to $\mathcal P$: a prime outside $\mathcal P$
can cancel only when $p=r$ and the signed exponents cancel,
which would force $q=1$. Hence
\begin{equation}\label{p03:eq:phase-sparse-resonance}
 \sum_{\substack{q n^\varepsilon m^\eta=1\\
                  \Lambda(n)\Lambda(m)\ne0}}
       \frac{\Lambda(n)\Lambda(m)}{\sqrt{nm}}
 \le\left(\sum_{p\in\mathcal P}\sum_{j\ge1}
                 (\log p)p^{-j/2}\right)^2\ll_{\mathcal P}1.
\end{equation}
Together with \eqref{p03:eq:phase-lag-bounds}, this bounds all nontrivial
rational resonances by $O_{\mathcal P}(TL)$.

The constant--prime terms can resonate only when $q$ or $q^{-1}$
is a prime power supported on $\mathcal P$. Their resonant part is
$O_{\mathcal P}(TL^2)$. For the remaining terms,
$|\log q\pm\log n|\ge\frac12\log n$ once $n$ exceeds a constant
depending on $q$. Direct integration in the midpoint thus bounds
them by $O_q(L^2(1+\sqrt X))=o(TL^3)$.

It follows that \eqref{p03:eq:phase-twisted-integral} is $o(N_I)$ for
$q\ne1$. For $q=1$ it is
\[
 \frac{N_I}{a^2}\left\{
   \frac\ell L\int\psi^2+
   \frac{2L}\ell\int_0^1v(\psi*\psi)(v)\,dv\right\}+o(N_I)
       =(R_\psi+o(1))N_I.
\]
Here $L/\ell\to1$, and the evenness of $\psi$ identifies the
convolution expression with that in \eqref{p03:eq:phase-R}. This sharpens the rough unweighted Hilbert--Schmidt bound used
in the boundary reduction.

For the first moment, setting $t=t'$ in
\eqref{p03:eq:phase-weighted-Poisson} yields
\[
 \sum_kq^{i\alpha_k}H_{kk}
  =\frac{b_\lambda}{aL}\int_Iq^{it}\nu_X(t)\,dt,
 \qquad
 b_\lambda=\int\phi(u-\lambda/2)\phi(u+\lambda/2)\,du
              =aL+O_\lambda(L^{-1}).
\]
The $q=1$ main term is $N_I+o(N_I)$. For $q\ne1$ the
archimedean contribution is $O_q(L)$, any prime resonance contributes
$O_q(T)$, and the remaining prime terms contribute $O_q(\sqrt X)$.
All are $o(N_I)$. We have proved both assertions for every fixed
trigonometric polynomial on the prime torus.

Finally, approximate $K$ uniformly by such a polynomial $F$ with
$\|K-F\|_\infty\le\delta$. The error in the weighted row energy
is at most $\delta\|G\|_{\mathrm{HS}}^2=O(\delta N_I)$, and the
trace error is at most
$\delta\sqrt d\,\|G\|_{\mathrm{HS}}=O(\delta N_I)$, where
$d\sim N_I$ is the matrix dimension. First let $T\to\infty$ and
then $\delta\to0$. This proves \eqref{p03:eq:phase-moments}.
The last assertion follows directly from
\[
 \|D_KGD_K\|_{\mathrm{HS}}^2
 =\sum_{k,\ell}k(\alpha_k)k(\alpha_\ell)G_{k\ell}^2
 \le\sum_k k(\alpha_k)\sum_\ell G_{k\ell}^2.
\]
\end{proof}

\begin{remark}\label{p03:rem:phase-fixed-scope}
The proposition takes the phase set and its approximating Fourier
polynomial to be fixed before $T$ tends to infinity. It supplies no
uniform estimate for a phase ball shrinking with $T$. The reduction
$X=T/(\log T)^D$ absorbs the elementary twisted mean-value error
while leaving the limiting window functional unchanged.
\end{remark}




\section{Localization on the zero side}
\label{p03:sec:phase-zero-localization}

The arithmetic moment calculation must be combined with a localization
argument before it gives a proportion of zeros in a phase set.  A
positive weight does not weight the positive index of a Hermitian form.
We use the weight on the sampling coordinates, retain the relatively
few zeros far from the critical line, and estimate the discarded
near-line vectors in Hilbert--Schmidt norm.

Fix a finite nonempty set of primes $\mathcal P$, and write
\[
 \Phi_{\mathcal P}(t)=(p^{-it})_{p\in\mathcal P},\qquad
 \mu_{\mathcal P}=\text{normalized Haar measure on }\T^{\mathcal P}.
\]
A phase set below is an open subset $B$ of this torus whose boundary
has Haar measure zero.  Let $N_B(T,2T)$ count all nontrivial zeta zeros
with ordinate in $(T,2T]$ and phase in $B$, with multiplicity, and let
$N^s_{0,B}(T,2T)$ count the simple zeros on the critical line satisfying
the same conditions.  Put $N=N(T,2T)$.

We use two classical facts about zero ordinates:
\begin{align}
 N(t,t+1)&\ll\log(|t|+3),\label{p03:eq:phase-local-zero-count}\\
 N(\sigma,T)&\ll T^{1-(\sigma-1/2)/4}\log T,
 \qquad \tfrac12\le\sigma\le1.\label{p03:eq:phase-selberg-density}
\end{align}
Here $N(\sigma,T)$ counts zeros with $|\Im\rho|\le T$ and
$\Re\rho\ge\sigma$, with multiplicity. The second is a uniform form of Selberg's density estimate; see
\cite[Lemma~2]{FordZaharescu}.  We also use the joint form of the
Hlawka equidistribution theorem,
\begin{equation}
 N_B(T,2T)=\bigl(\mu_{\mathcal P}(B)+o(1)\bigr)N.
 \label{p03:eq:phase-zero-haar}
\end{equation}
Indeed Weyl's criterion reduces this assertion to the uniform
distribution of $\gamma\log r/(2\pi)$ for a fixed positive rational
$r\ne1$.  Unique factorization ensures that each nontrivial torus
character gives such an $r$.  The one-dimensional assertion follows
from Landau's formula and the classical density estimate; see
\cite{AkbaryMurty}.  Subtracting the asymptotic at $T$ from that at
$2T$ gives \eqref{p03:eq:phase-zero-haar}.

Here is the normalization of the finite compression needed in this
section.  Fix a real even function
$f\in C_c^\infty((-1/2,1/2))$, $f\not\equiv0$, and let
\begin{gather*}
 L=\log T-D\log\log T,\quad D>4,\quad
 \phi(u)=f(u/L),\quad a=\int_{\R}f(u)^2\,du,\\
 \alpha_k=T+\frac{2\pi k}{L}\in[T,2T].
\end{gather*}
For a zero $\rho=\beta+i\gamma$ put
$z_\rho=\gamma+i(1/2-\beta)$ and
\begin{equation}
 u_\rho(k)=\frac{\widehat\phi(z_\rho-\alpha_k)}{\sqrt a\,L}.
 \label{p03:eq:phase-normalized-vector}
\end{equation}
By the explicit formula, $G$ is the sum of
$m_\rho u_\rho u_\rho^{\mathsf T}$ over all nontrivial zeros.
The Fourier support of a product of two test transforms lies
in $(-L,L)$. On that interval the transform of the pole density
is $2\cosh(u/2)$, which accounts for both pole evaluations;
the other terms are the prime and gamma terms of Weil's formula.
Rapid decay and the unit-interval zero count make the zero sum
convergent, as in \cite{AlpogeFurman}.
Let $G'$ be the sum restricted to the extended height window
$I'=[T-\sqrt T,2T+\sqrt T]$. Then $\|G-G'\|_1=o(1)$,
where $\|\cdot\|_1$ denotes trace norm.
Indeed, repeated integration by parts bounds each
normalized coordinate by
$C_M T^{1/4}(1+L|\gamma-\alpha_k|)^{-M}L^M$ for every zero.
Sum its square over the centers and over zeros at distance at least
$\sqrt T$ from the center interval, using the unit-interval count.
The result is $O_M(NT^{1-M}\log T)=o(1)$ for large fixed $M$.
The trace norm of the sum of the rank-one bilinear terms is at most
this sum of squared vector norms. Coordinate multiplication by a
contraction preserves the estimate.

Choose $K\in C^\infty(\T^{\mathcal P})$ with
$0\le K\le1$ and $\operatorname{supp}K\subset B$, and put
\[
 D_K=\operatorname{diag}
 \bigl(K(\Phi_{\mathcal P}(\alpha_k))^{1/2}\bigr),
 \qquad G_K=D_KGD_K,\quad F_{\rm out}=D_K(G-G')D_K.
\]
Only boundedness of $D_K$ is needed on the zero side; the square root
need not be smooth.  The arithmetic input from
Proposition~\ref{p03:prop:phase-moments} is
\begin{align}
 \operatorname{tr}G_K
   &=\bigl(\mu_{\mathcal P}(K)+o(1)\bigr)N,
   \label{p03:eq:phase-localized-trace}\\
 \|G_K\|_{\mathrm{HS}}^2
   &\le\bigl(R_\psi\mu_{\mathcal P}(K)+o(1)\bigr)N.
   \label{p03:eq:phase-localized-hs}
\end{align}
For the second inequality one uses
$\|D_KGD_K\|_{\mathrm{HS}}^2
\le\operatorname{tr}(D_K^2G^2)$, since $G$ is Hermitian and
$0\le D_K^2\le I$.

\begin{lemma}[Discarding shallow zeros outside a phase neighborhood]
\label{p03:lem:phase-shallow-tail}
Fix $\eta>0$, and let $B_\eta$ be the open $\eta$-neighborhood of $B$
in the Euclidean metric inherited by the torus.  Put
\[
 d_0=16\frac{\log\log T}{\log T}.
\]
Let $E$ be the contribution to $G_K$ from zeros in $I'$ for which
$|\beta-1/2|\le d_0$ and
$\Phi_{\mathcal P}(\gamma)\notin B_\eta$.  Then
\[
 |\operatorname{tr}E|=o(N),\qquad
 \|E\|_{\mathrm{HS}}=o(\sqrt N).
\]
The total multiplicity in $I'$ of zeros with
$|\beta-1/2|>d_0$ is $o(N)$.
\end{lemma}

\begin{proof}
The last assertion follows from \eqref{p03:eq:phase-selberg-density},
reflection about the critical line, and $N\asymp T\log T$: the
relevant count is $O(N(\log T)^{-4})$.

The map $\Phi_{\mathcal P}$ is Lipschitz with constant
$C_{\mathcal P}=(\sum_{p\in\mathcal P}(\log p)^2)^{1/2}$.
Consequently every discarded zero is separated from every sampling
center with nonzero weight by at least
$r_\eta=\eta/C_{\mathcal P}>0$ in real height.
Repeated integration by parts, for any fixed integer $M$, gives
\[
 |u_\rho(k)|\ll_M
 e^{L|\beta-1/2|/2}
 (1+L|\beta-1/2|)^M
 (1+L|\gamma-\alpha_k|)^{-M}.
\]
For $|\beta-1/2|\le d_0$ and sufficiently large $T$ this implies
\begin{equation}
 |u_\rho(k)|\ll_M
 (\log T)^9(1+L|\gamma-\alpha_k|)^{-M}.
 \label{p03:eq:phase-shallow-vector-bound}
\end{equation}
The factor $(\log\log T)^M$ has been absorbed in one power of
$\log T$.  Constants may depend on the fixed $M$.

Form a rectangular matrix $U$ with columns $D_Ku_\rho$ for the
discarded zeros, repeating a column according to its multiplicity.
The grid spacing and \eqref{p03:eq:phase-local-zero-count} give
\[
 \sup_k\sum_\rho|U_{k\rho}|\ll(\log T)^{10},
 \qquad
 \sup_\rho\sum_k|U_{k\rho}|\ll(\log T)^9.
\]
Schur's test yields $\|U\|^2\ll(\log T)^{19}$.  The separation
$|\gamma-\alpha_k|\ge r_\eta$ additionally gives
\[
 \sum_k|U_{k\rho}|^2
 \ll (\log T)^{18}(Lr_\eta)^{1-2M},
 \qquad
 \|U\|_{\mathrm{HS}}^2
 \ll N(\log T)^{19-2M}.
\]
Here and below the constants can depend on $\eta$.
Since $E=UU^{\mathsf T}$, standard singular-value inequalities imply
\[
 |\operatorname{tr}E|\le\|U\|_{\mathrm{HS}}^2,
 \qquad
 \|E\|_{\mathrm{HS}}^2
 \le\|U\|^2\|U\|_{\mathrm{HS}}^2
 \ll N(\log T)^{38-2M}.
\]
Taking $M=24$ proves both assertions.  Conjugate zero pairs show that
$E$ is Hermitian, although its factorization uses a transpose rather
than an adjoint.
\end{proof}

\begin{theorem}[The localization implication]
\label{p03:thm:phase-localization}
Suppose the weighted moment estimates
\eqref{p03:eq:phase-localized-trace}--\eqref{p03:eq:phase-localized-hs} hold
for every fixed smooth $K$ with $0\le K\le1$.  Then, for every phase
set $B$ as above,
\begin{equation}
 \liminf_{T\to\infty}\frac{N^s_{0,B}(T,2T)}{N(T,2T)}
 \ge (2-R_\psi)\mu_{\mathcal P}(B).
 \label{p03:eq:phase-proportion-conclusion}
\end{equation}
In particular the same proportion constant furnished by the global
second-moment certificate is valid after restriction to any fixed
finite-prime phase set.
\end{theorem}

\begin{proof}
First fix $K$ and $\eta>0$, choosing $\eta$ so that
$\mu_{\mathcal P}(\partial B_\eta)=0$.  Such choices exist
arbitrarily close to zero, since disjoint level sets of the distance
function have positive measure for at most countably many levels.
Retain all deep zeros $|\beta-1/2|>d_0$ in $I'$, independently of
their phase, and all shallow zeros whose phase belongs to $B_\eta$.
By Lemma~\ref{p03:lem:phase-shallow-tail} and
$\|F_{\rm out}\|_1=o(1)$, the exact retained matrix
$H=G_K-F_{\rm out}-E$ satisfies
\[
 \operatorname{tr}H=\mu_{\mathcal P}(K)N+o(N),\qquad
 \|H\|_{\mathrm{HS}}^2
 \le R_\psi\mu_{\mathcal P}(K)N+o(N).
\]
Let $P_1$ be its contribution from simple on-line zeros, and put
$Q=H-P_1$.  Each such zero contributes a positive matrix of rank at
most one.  Its trace is at most one, by the full-grid sampling
identity and $\|D_K\|\le1$.  Each distinct multiple on-line zero
costs at most one positive direction, and each off-line reflected
pair likewise costs at most one.  The retained deep zeros have total
multiplicity $o(N)$.  It follows that
\begin{align*}
 \operatorname{rank}P_1
 &\le N^s_{0,B_\eta}(I'),\\
 N^s_{0,B_\eta}(I')+2n_+(Q)
 &\le N_{B_\eta}(I')+o(N).
\end{align*}
Together with $\operatorname{tr}P_1\le N^s_{0,B_\eta}(I')$,
Lemma~\ref{p03:lem:rank-basic} now gives
\begin{align*}
 N^s_{0,B_\eta}(I')
 &\ge 4\operatorname{tr}H
       -2N_{B_\eta}(I')-\|H\|_{\mathrm{HS}}^2-o(N)\\
 &\ge \bigl((4-R_\psi)\mu_{\mathcal P}(K)
             -2\mu_{\mathcal P}(B_\eta)\bigr)N-o(N).
\end{align*}
We used \eqref{p03:eq:phase-zero-haar}; the two added height intervals
contribute $O(\sqrt T\log T)=o(N)$ by
\eqref{p03:eq:phase-local-zero-count}.

To obtain the result for $B$ itself, first choose an inner open phase
set $C$ with closure contained in $B$ and Haar-null boundary.  Choose
$K$ supported in $C$, with $\mu_{\mathcal P}(K)$ arbitrarily close
to $\mu_{\mathcal P}(C)$, and let $\eta\downarrow0$ through the
admissible values, always keeping $C_\eta\subset B$.  The preceding
inequality gives a lower bound approaching
$(2-R_\psi)\mu_{\mathcal P}(C)$.  Inner regularity and the
Haar-null boundary of $B$ allow
$\mu_{\mathcal P}(C)\uparrow\mu_{\mathcal P}(B)$.
All phase sets, weights, and buffers are fixed before $T$ tends to
infinity, so no unproved uniformity for shrinking targets is used.
\end{proof}

\begin{remark}
Neither a zero-free half-plane nor an independence assumption between
phase conditions and critical-line zeros enters this implication.
The small family of deep zeros is charged by its count rather than
estimated by the norms of its evaluation vectors.  This distinction
is what permits localization on a fixed physical height scale.
\end{remark}

\begin{proof}[Proof of Theorem~\ref{p03:thm:main-phase}]
Proposition~\ref{p03:prop:phase-moments} supplies the hypotheses of
Theorem~\ref{p03:thm:phase-localization}. Smooth compactly supported
nonnegative even functions $f$ can approximate the square root of
\[
 \psi_{\rm MT}(u)=\cos(\sqrt2u)\,\mathbf1_{[-1/2,1/2]}(u).
\]
Both integrals defining $R_\psi$ converge under this approximation.
Direct integration gives
$R_{\psi_{\rm MT}}=1/2+2^{-1/2}\cot(2^{-1/2})$.
Let the approximation parameter tend to zero after $T$ tends to
infinity. For a Jordan-measurable $B$, its interior has the same
Haar measure, so the open-set assertion implies the stated result.
The empty prime set reduces to the unconditioned case.
\end{proof}




\section{Uniform Weil arithmetic at variable bandwidth}
\label{p04:sec:basic-Weil}

For shrinking phase sets we use the same compression
\eqref{p03:eq:phase-Weil-matrix}, but now allow
$\ell/2\le L\le\ell$, where $\ell=\log T$ and $X=e^L$.
The window $f$ remains fixed, nonzero, even and smooth, with support
strictly inside $(-1/2,1/2)$; again $a=\int f^2$ and $\psi=f^2$.
The exact zero-side representation and the full-grid normalization
are unchanged. A reflected off-line pair with vectors $u,\bar u$
contributes $uu^{\mathsf T}+\bar u\bar u^{\mathsf T}$; writing
$u=v+iw$ makes this $2vv^{\mathsf T}-2ww^{\mathsf T}$, with at
most one positive direction. Thus the single rank--inertia
certificate, Lemma~\ref{p03:lem:rank-basic}, applies at this bandwidth
as well. We record the variable-bandwidth estimates explicitly
because their errors must be small relative to the Haar mass.

Here are the analytic identities used in the uniform arithmetic
estimate below. For the full-grid synthesis operator with functions
$u_k(t)=\widehat\phi(t-\alpha_k)/(\sqrt a L)$, Parseval gives
$\|U\|^2\ll_f L^{-1}$ and $\sum_k|u_k(t)|^2=1$.
Schwartz decay controls the mass outside the center interval by
$(1+L\operatorname{dist}(t,I))^{-M}$, and the analogous
inside-boundary tail has integral $O_f(L^{-1})$.
The Montgomery--Vaughan mean-value inequality~\cite{Hilbert} gives
$\int_I|\nu_X(t)|^2dt\ll_f TL^2$.
The operator inequality
\[
 \|V^*M_{\nu_X1_J}V\|_{\rm HS}^2
 \le\|V\|^2\int_J|\nu_X(t)|^2\sum_k|v_k(t)|^2dt
\]
then bounds the squared matrix error caused by a boundary reduction
by $O_f(1+X/L^2)$ and the full squared norm by $O_f(TL)$.

For $\lambda=\log q$ small enough to avoid aliases, Poisson
summation gives
\[
 \sum_{k\in\mathbb Z}e^{i\lambda\alpha_k}
 \widehat\phi(t-\alpha_k)\widehat\phi(t'-\alpha_k)
 =Le^{i\lambda(t+t')/2}\Psi_\lambda(t-t'),
\]
where $\Psi_\lambda(s)=\int\phi(u-\lambda/2)
\phi(u+\lambda/2)e^{-isu}du$.
Put $J_\lambda=\Psi_\lambda\Psi_0$. After the boundary
reduction, the weighted row energy is
\[
 \frac1{a^2L^2}\int_{I^2}e^{i\lambda(t+t')/2}
        J_\lambda(t-t')\nu_X(t)\nu_X(t')\,dt\,dt'.
\]
For the zero character, its prime diagonal before division by
$a^2L^2$ is
\[
 \frac T\pi\sum_{n\le X}\frac{\Lambda(n)^2}{n}
                         (\phi^2*\phi^2)(\log n).
\]
The elementary prime-number-theorem consequence
$\sum_{n\le y}\Lambda(n)^2/n=\frac12\log^2y+O(\log y)$
and partial summation give
$L^3\int_0^1v(\psi*\psi)(v)dv+O_f(L^2)$ for its sum.
The gamma--gamma term supplies
$T\ell^2L\int\psi^2/(2\pi)$ before the same division.
Together they give the functional $R_f(L/\ell)$ below.
We also write $R_\psi=R_f(1)$. Approximating
$\psi_{\rm MT}(u)=\cos(\sqrt2u)1_{[-1/2,1/2]}(u)$ by squares
of fixed smooth windows gives
$R_{\psi_{\rm MT}}=1/2+2^{-1/2}\cot(2^{-1/2})$.

\subsection{Uniform arithmetic with a small Haar mass}


\begin{proposition}[Uniform rational arithmetic]\label{p04:prop:uniform-rational-arithmetic}
Fix a nonzero even real $f\in C_c^\infty((-1/2,1/2))$ and choose
$\sigma_f<1$ with $\operatorname{supp}f\subset[-\sigma_f/2,\sigma_f/2]$.
Let $\ell=\log T$, $\ell/2\le L\le\ell$, $X=e^L$, and form the Weil
matrix $G$ with this bandwidth as in the preceding construction.
Let $S$ be a finite set of primes, allowed to vary, and let
\[
 F(z)=\sum_{\mathbf n}c_{\mathbf n}z^{\mathbf n},\qquad
 m=c_{\mathbf0}\ge0,\qquad
 |c_{\mathbf n}|\le m,\qquad
 A=\sum_{\mathbf n}|c_{\mathbf n}|.
\]
The Fourier sum is finite. For each character write its associated
rational in lowest terms as $q=b/c$ and suppose $b,c\le H$, with
\[
 \log H\le\kappa L,\qquad 0<\kappa<1-\sigma_f.
\]
Put
\[
 B_S=\sum_{p\in S}\frac{\log p}{\sqrt p-1},\qquad
 R_f(\alpha)=\frac{\alpha^{-1}\int f^4+
 \alpha\iint|u-v|f(u)^2f(v)^2\,du\,dv}{(\int f^2)^2}.
\]
Then, uniformly in $S,F,H,L$ subject to these conditions,
\begin{align*}
 \sum_k F(F_S(\alpha_k))G_{kk}
 &=mN+O_f\!\left(mN\frac{1+B_S}{L}+AN\mathcal E\right),\\
 \sum_k F(F_S(\alpha_k))\sum_jG_{kj}^2
 &=mR_f(L/\ell)N+
 O_{f,\kappa}\!\left[mN\left(\frac{1+B_S}{L}
                       +\frac{B_S^2}{L^2}\right)+AN\mathcal E\right],
\end{align*}
where $N=N(T,2T)$ and one may take
\[
 \mathcal E=\frac{HX}{T}+\frac1{\sqrt{TL}}+
             \frac{\sqrt{X/T}}{L^{3/2}}+
             \frac{\sqrt X}{TL}.
\]
The constants do not depend on the number of primes in $S$.
The coefficient condition holds, in particular, for any Fourier
truncation of a nonnegative integrable phase weight, with $m$ its
Haar integral.
\end{proposition}

\begin{proof}
The boundary and full-grid reduction uses only $|q^{it}|=1$ for
individual characters. The squared Hilbert--Schmidt error is
$O_f(1+X/L^2)$, whereas the full squared norm is $O_f(TL)$.
Cauchy--Schwarz therefore gives, after division by $N$, the second
and third terms of $\mathcal E$. Linear trace boundary errors are
smaller. Summing over characters costs $A$.

For $\lambda=\log q$, the weighted Poisson formula is exact and has
no aliases, because $|\lambda|\le\kappa L$ and
$\kappa+\sigma_f<1$. The lag kernels $J_\lambda$ obey
\[
 \int|J_\lambda(s)|\,ds\ll_f L,\qquad
 \int|sJ_\lambda(s)|\,ds\ll_f1
\]
uniformly in this entire range of $\lambda$. Indeed, after scaling,
the two translated arguments of $f$ range over a fixed compact set
and every derivative has a uniform bound.

The rational-frequency separation used in the mean-value estimate
is uniformly $\gg(HX)^{-1}$. For opposite signs this follows from
$|bn-cm|\ge1$ whenever $bn\ne cm$. For equal signs, a small
$|\log(qnm)|$ implies $bnm\asymp c\le H$, so the same lower bound
is more than sufficient. Frequencies internal to either list are
separated by $\gg X^{-1}$.

It is essential not to replace the archimedean density by a constant
in every Fourier mode. Keep
\[
 \mu(t)=\frac1{2\pi}\Re\frac{\Gamma'}\Gamma(1/4+it/2)
               -\frac{\log\pi}{2\pi}
\]
exact in all nontrivial modes. For $|s|\le T/2$, as functions of
the midpoint $\tau\in[T,2T]$, each $\mu(\tau\pm s/2)$ has
supremum plus total variation $O(\ell)$; their product has the
corresponding bound $O(\ell^2)$. Abel summation therefore extends
the Montgomery--Vaughan Hilbert inequality to these weights.
All nonresonant prime--prime and mixed terms have midpoint integral
$O(HX\ell^2)$, uniformly in the lag. The constant--constant
nontrivial mode is $O(H\ell^2)$ by direct integration.
Integration against the lag kernel and division by $a^2L^2$
give $O(NHX/T)$ per unit Fourier $\ell^1$ mass. The large-lag
Schwartz tails and the lag-boundary correction are smaller.

On $[T,2T]$, the pole density has $L^2$ norm $O(\sqrt{X/T})$.
The remaining density has $L^2$ norm $O(\sqrt T L)$. Its insertion
in the lag integral thus contributes $O_f(\sqrt X)$ after
normalization, yielding the final term of $\mathcal E$.

For the zero character the usual diagonal computation gives
$R_f(L/\ell)N+O_f(N/L)$. Here only this character uses
$\mu=\ell/(2\pi)+O(1)$, so its error is multiplied by $m$.
The nontrivial resonances retain the same factor. If
$q n^{\varepsilon}m_1^{\eta}=1$, $q\ne1$, and both von Mangoldt
factors are nonzero, both prime bases must lie in $S$. Summing
over all nontrivial $q$ in the Fourier support therefore gives
\[
 \sum_{q\ne1}|c_q|
 \sum_{qn^{\varepsilon}m_1^{\eta}=1}
 \frac{\Lambda(n)\Lambda(m_1)}{\sqrt{nm_1}}
 \le 4mB_S^2.
\]
Similarly the total mixed-resonance amplitude is at most $2mB_S$.
Their normalized contributions are respectively
$O_f(mNB_S^2/L^2)$ and $O_f(mNB_S/L)$.

For the trace, the diagonal Poisson formula reduces the calculation
to a single midpoint integral with a prefactor bounded by one.
The zero-character integral is $N+O_f(\sqrt X+\ell)$.
Nontrivial resonances sum to $O(mTB_S)$, and nonresonant prime
integrals are $O(AHXL)$ by Cauchy--Schwarz and the frequency
separation. The archimedean nontrivial terms are $O(AH\ell)$.
These bounds give the asserted first-moment estimate.
\end{proof}

\begin{remark}[A useful choice of bandwidth]
If $\log H=o(\log T)$, choose
\[
 L=\log T-\log H-D\log\log T.
\]
Then $HX/T=(\log T)^{-D}$ and $L/\log T\to1$.
For fixed $S$ and a smooth phase bump of width
$r=(\log T)^{-a}$, $a<1$, Fourier truncation at degree
$r^{-1}(\log T)^\epsilon$, $a+\epsilon<1$, has
$\log H=O_S((\log T)^{a+\epsilon})$, uniformly bounded Fourier
$\ell^1$ norm, and coefficients bounded by its Haar mass.
The arithmetic moments consequently remain asymptotic relative to
that mass. This arithmetic observation alone does not establish
a shrinking-target theorem: the zero-side inertia budget must also
be localized with a uniform relative error.
\end{remark}




\section{Shrinking prime-phase boxes}
\label{p04:sec:shrinking-boxes}

Throughout this section the prime set $S=\{p_1,\ldots,p_d\}$ is
fixed.  We identify $\mathbb T^d$ with $(\mathbb R/\mathbb Z)^d$,
and write
\[
 \Theta_S(t)=\left(-\frac{t\log p_j}{2\pi}\pmod1\right)_{j=1}^d.
\]
For a box $B$ let $N_B(T,2T)$ count all nontrivial zeros, with
multiplicity, whose ordinates lie in $(T,2T]$ and whose phase is in
$B$.  The analogous count restricted to simple critical-line zeros
is $N^s_{0,B}(T,2T)$.

\begin{theorem}[Anisotropic shrinking targets]
\label{p04:thm:shrinking-boxes}
Fix $a_1,\ldots,a_d\ge0$ and $c_1,\ldots,c_d\in(0,1/4]$. Put
\[
 A=\sum_{j=1}^d a_j,\qquad a_* =\max_j a_j,
 \qquad r_j(T)=c_j(\log T)^{-a_j}.
\]
Suppose
\begin{equation}\label{p04:eq:shrinking-range}
 a_*<1.
\end{equation}
For any choice of centers $\theta(T)\in\mathbb T^d$, let
$B_T=\prod_{j=1}^d I_j(T)$, where $I_j(T)$ is the circular interval
of length $r_j(T)$ centered at $\theta_j(T)$. Then, uniformly in the
choice of centers,
\begin{align}
 N_{B_T}(T,2T)&=(1+o(1))\mu(B_T)N(T,2T),
                 \label{p04:eq:shrinking-all}\\
 \liminf_{T\to\infty}
 \inf_{\theta\in\mathbb T^d}
 \frac{N^s_{0,B_T}(T,2T)}{\mu(B_T)N(T,2T)}
 &\ge \frac32-\frac1{\sqrt2}\cot\frac1{\sqrt2}.
                 \label{p04:eq:shrinking-simple}
\end{align}
In particular, equal side lengths $(\log T)^{-a}$, up to fixed
positive factors, are allowed whenever $a<1$.
\end{theorem}

The strict inequality is a proved sufficient range, not a proposed
optimal threshold.  The dimension and the primes remain fixed.
The centers can depend arbitrarily on $T$; no uniformity in a growing
prime set is asserted.

\subsection{A paired quantitative Landau estimate}

Set $\ell=\log T$ and $N=N(T,2T)$.  For $n\in\mathbb Z^d$ put
$q_n=\prod_jp_j^{n_j}$ and $\lambda_n=\log q_n$.  If
$\|n\|_\infty\le H$, then the numerator and denominator of $q_n$
in lowest terms are at most $\exp(C_SH)$. In particular,
$|\lambda_n|\le C_SH$, and, for $n\ne0$,
$|\lambda_n|\ge\exp(-C_SH)$ after increasing $C_S$.

\begin{lemma}\label{p04:lem:paired-Landau}
Uniformly for $H=o(\ell)$ and $0<\|n\|_\infty\le H$, let
$x_n=\max(q_n,q_n^{-1})>1$. Then
\begin{equation}\label{p04:eq:paired-Landau}
 \sum_{T<\gamma\le2T}e^{i\lambda_n\gamma}
 =-\frac{T}{2\pi}\frac{\Lambda(x_n)}{\sqrt{x_n}}
 +O_S\left(N\frac{\lambda_n^2}{\ell^2}
               +e^{C_SH}\ell^2\right),
\end{equation}
where $\Lambda(x)=0$ unless $x$ is an integral prime power, and the
sum counts multiplicities.  The same estimate holds with the opposite
sign in the exponential.
\end{lemma}
\begin{proof}
The uniform Landau formula in \cite[Lemma~1]{FordZaharescu}, applied
at $T$ and $2T$, gives, for $x>1$,
\[
 \sum_{T<\gamma\le2T}x^\rho
 =-\frac{\Lambda(n_x)}{2\pi}
       \frac{e^{2iT\log(x/n_x)}-e^{iT\log(x/n_x)}}
            {i\log(x/n_x)}
 +O\left(x\log^2(2xT)+\frac{\log(2T)}{\log x}\right),
\]
where $n_x$ is a nearest prime power and the quotient is interpreted
as $T$ when $x=n_x$. For $x=x_n$, rational separation shows that
all nonresonant terms here are $O_S(e^{C_SH}\ell^2)$, after increasing
$C_S$: a nonzero numerator of $x_n-n_x$ is an integer, whereas all
relevant numerators and denominators are at most exponential in $H$.
Dividing by $\sqrt{x_n}$ therefore gives the asserted main term for
the sum with weight $x_n^{\beta-1/2}$.

It remains to remove that real weight.  Horizontal reflection
$\beta+i\gamma\mapsto1-\beta+i\gamma$, preserving multiplicity,
cancels its linear variation.  Writing $u_\rho=|\beta-1/2|$ and
$\lambda=\log x_n$, the absolute replacement error is bounded by
\[
 \sum_{T<\gamma\le2T}\bigl(\cosh(\lambda u_\rho)-1\bigr).
\]
The uniform Selberg estimate gives
\[
 \#\{\rho:T<\gamma\le2T,\ u_\rho>v\}
       \ll N e^{-v\ell/4}\qquad(0\le v\le1/2).
\]
By layer integration, for $\lambda\le\ell/8$ the preceding error is
at most
\[
 C N\int_0^{1/2}\lambda\sinh(\lambda v)e^{-v\ell/4}\,dv
 \ll N\lambda^2\int_0^\infty v e^{-v\ell/8}\,dv
 \ll N\frac{\lambda^2}{\ell^2}.
\]
The condition $H=o(\ell)$ ensures $\lambda\le\ell/8$ eventually.
Complex conjugation handles a negative $\lambda_n$.
\end{proof}

Pairing before taking absolute values is essential: a first-order
estimate would give $N|\lambda|/\ell$ and a smaller shrinking range.

\subsection{Smooth boxes and the all-zero budget}

Let $\tau=1/\log\ell$ and $r=\min_jr_j$.  By convolving circular
intervals with nonnegative smooth mollifiers of widths comparable to
$\tau r_j$, we obtain $K^-_T,K^+_T\in C^\infty(\mathbb T^d)$ with
\begin{equation}\label{p04:eq:box-sandwich}
 0\le K^-_T\le1_{B_T}\le K^+_T\le1,
 \qquad \int K^\pm_T\,d\mu=(1+O_d(\tau))\mu(B_T).
\end{equation}
The support of $K^-_T$ is separated from the complement of $B_T$ by
at least $c\tau r$ in the torus metric.  This can be arranged by
shrinking the initial intervals slightly more before convolution.
The construction and all following estimates are invariant under
translation of the center.

Writing either Fourier series as
$K=\sum_n c_n e^{2\pi in\cdot\theta}$, put $m=\int K\,d\mu$.
Nonnegativity gives $|c_n|\le m$.  Standard scaling estimates for
the mollifiers give
\begin{align}
 \sum_n|c_n|&\ll_d\tau^{-d},\\
 \sum_n|c_n||\lambda_n|^2
       &\ll_S\tau^{-d-2}r^{-2},\\
 \sum_{\|n\|_\infty>H}|c_n|
       &\ll_{S,M}\tau^{-d}(H\tau r)^{-M}
       \quad(H\tau r\ge1),
                 \label{p04:eq:box-Fourier-estimates}
\end{align}
for every fixed $M$, with a harmless change of $M$ in the derivative
estimate. For example, the one-dimensional coefficients are bounded
by $C_Mr_j(1+\tau r_j|n_j|)^{-M}$, and summation of the product
estimates proves these assertions.

Choose $b$ with $a_*<b<1$ and set $H=\ell^b$.
For sufficiently large fixed $M$, the last error in
\eqref{p04:eq:box-Fourier-estimates} is $o(\mu(B_T))$.  Apply
Lemma~\ref{p04:lem:paired-Landau} to the truncated series.  The prime-power
main terms satisfy the sharper bound
\[
 \sum_{n\ne0}|c_n|\frac{\Lambda(x_n)}{\sqrt{x_n}}
 \le 2m\sum_{p\in S}\sum_{k\ge1}(\log p)p^{-k/2}
 \ll_S m.
\]
Indeed an integral prime power supported on $S$ has just one nonzero
coordinate in its exponent vector.  Thus
\begin{equation}\label{p04:eq:smooth-box-zero-count}
 \sum_{T<\gamma\le2T}K(\Theta_S(\gamma))
 =mN+O_S\left(\frac{mN}{\ell}
       +\frac{N\tau^{-d-2}}{\ell^2r^2}
       +\tau^{-d}e^{C_SH}\ell^2\right)+o(\mu(B_T)N).
\end{equation}
Since $\mu(B_T)\asymp\ell^{-A}$ and $r\asymp\ell^{-a_*}$,
the stronger restriction $A+2a_*<2$ makes every error
$o(\mu(B_T)N)$. Sandwiching by \eqref{p04:eq:box-sandwich} proves
\eqref{p04:eq:shrinking-all} in that smaller range.  We next remove this
restriction.

\subsection{Analytic box bounds in a thin horizontal strip}

The loss in the elementary Fourier argument can be removed by
bounding the phase box directly at complex zero coordinates.  This
avoids replacing each Fourier weight by its value on the critical
line separately.

Put $v=\log\ell$, $\tau=v^{-1}$, and choose Gaussian widths
$s_j=\tau r_j/v=r_j/v^2$. Let $B^+$ and $B^-$ denote boxes obtained
by increasing or decreasing every side of $B_T$ by $4\tau r_j$.
For large $T$ the latter box is nonempty.  On the complexified torus
define the periodized Gaussian and its convolutions
\[
 g_s(z)=\left(\prod_{j=1}^d\frac1{\sqrt{2\pi}s_j}\right)
       \sum_{k\in\mathbb Z^d}
          \exp\left(-\sum_{j=1}^d
                         \frac{(z_j+k_j)^2}{2s_j^2}\right),
 \qquad F^\pm(z)=\int_{B^\pm}g_s(z-u)\,du.
\]
Both functions are
entire and periodic in every real coordinate, and
\[
 \int_{\mathbb T^d}g_s(z-u)\,du=1.
\]
Their Fourier coefficients are
\[
 \widehat{F^\pm}(n)=\widehat{1_{B^\pm}}(n)
                    \exp\left(-2\pi^2\sum_js_j^2n_j^2\right).
\]
In particular $|\widehat{F^\pm}(n)|\le\mu(B^\pm)$ and
\begin{equation}\label{p04:eq:Gaussian-coefficient-mass}
 \sum_n|\widehat{F^\pm}(n)|
       \ll_d\mu(B^\pm)\prod_js_j^{-1}\ll_d v^{2d}.
\end{equation}

Choose a sufficiently large fixed $B_0$ and put
$d_0=4B_0\log\ell/\ell$.  A zero with
$|\beta-1/2|\le d_0$ gives a complex phase
\[
 z_\rho=\Theta_S(\gamma)+iy_\rho,
 \qquad (y_\rho)_j=(\beta-1/2)\frac{\log p_j}{2\pi}.
\]
For all these phases,
\[
 h:=\sup_\rho\left(\sum_j\frac{|(y_\rho)_j|^2}{s_j^2}
                                                    \right)^{1/2}
       \ll_{S,B_0}\ell^{a_*-1}v^3=o(1).
\]
The absolute value of the complex Gaussian is exactly its real
Gaussian times $\exp(\|y/s\|^2/2)$.  Its negative real part can
occur only when, in at least one lattice term,
\[
 \left|\sum_j\frac{(x_j+k_j)y_j}{s_j^2}\right|\ge\pi/2.
\]
Cauchy--Schwarz forces
$\|(x+k)/s\|\ge\pi/(2\|y/s\|)$. Gaussian tail summation
therefore gives, uniformly in real $x$ and $\|y/s\|\le h$,
\begin{equation}\label{p04:eq:Gaussian-negative-tail}
 \Re g_s(x+iy)\ge
       -C_d\left(\prod_js_j^{-1}\right)e^{-c_d/h^2}.
\end{equation}
The assertion is interpreted by continuity when $y=0$.

The real distance, in each relevant normal direction, from $B_T$
to the complement of $B^+$ is at least $2\tau r_j=2vs_j$.
The identity for the full Gaussian integral and the Gaussian tail
bound imply
\[
 \Re F^+(x+iy)\ge1-O_d(e^{-c_dv^2})\quad(x\in B_T).
\]
The analogous separation of $B^-$ from the complement of $B_T$
gives
\[
 |F^-(x+iy)|\ll_d e^{-c_dv^2}\quad(x\notin B_T).
\]
Globally $\Re F^-\le e^{h^2/2}=1+o(1)$, whereas
\eqref{p04:eq:Gaussian-negative-tail} gives
$\Re F^+\ge-O_d((\prod_js_j^{-1})e^{-c_d/h^2})$.
Because $a_*<1$, all the displayed exponentially small errors are
$o(\ell^{-C})$ for every fixed $C>0$. Consequently there is
$\epsilon_T=o(\ell^{-C})$ for every fixed $C$ such that, at all
shallow zero phases,
\begin{equation}\label{p04:eq:complex-box-sandwich}
 (1+o(1))^{-1}\Re F^-(z_\rho)-\epsilon_T
 \le1_{B_T}(\Theta_S(\gamma))
 \le(1+o(1))\Re F^+(z_\rho)+\epsilon_T.
\end{equation}
The multiplicative errors are independent of the phase center.

Truncate $F^\pm$ to $\|n\|_\infty\le H=\ell^b$, with
$a_*<b<1$, and call the resulting polynomials $P^\pm$.
The relation $d_0/s_{\min}^2=o(H)$ allows the growing imaginary
factor in the Fourier tail to be absorbed by Gaussian decay.
At shallow phases the truncation error is
\[
 O_{S,d}\left(v^{2d}
         \exp(-c s_{\min}^2H^2+C_Sd_0H)\right)
       =o(\ell^{-C})
\]
for every fixed $C$.  Apply the uniform Landau formula directly to
$P^\pm(z_\rho)$, without removing the real zero weights.  Horizontal
reflection or complex conjugation gives the required sign convention
for the frequencies.  The prime-power terms are
$O_S(\mu(B_T)N/\ell)$ by the coefficient bound and the convergent
one-prime series used above. The nonresonant terms are
$O_S(v^{2d}e^{C_SH}\ell^2)=o(\mu(B_T)N)$. Thus, with all zeros
included,
\begin{equation}\label{p04:eq:analytic-box-Landau}
 \sum_{T<\gamma\le2T}P^\pm(z_\rho)
       =(\mu(B^\pm)+o(\mu(B_T)))N.
\end{equation}

The contribution of deep zeros must be removed with their complex
weights, not merely by counting them. By
\eqref{p04:eq:Gaussian-coefficient-mass},
\[
 |P^\pm(z_\rho)|\ll_d v^{2d}e^{C_SH|\beta-1/2|}.
\]
Layer integration of Selberg's estimate, using $C_SH=o(\ell)$,
therefore gives
\[
 \sum_{\substack{T<\gamma\le2T\\|\beta-1/2|>d_0}}
       |P^\pm(z_\rho)|
       \ll_{S,d}Nv^{2d}\exp(-d_0\ell/8)
       =O_{S,d}(Nv^{2d}\ell^{-B_0/2}).
\]
Choose $B_0>2A+4$.  This quantity, and the unweighted count of deep
zeros, are $o(\mu(B_T)N)$.  Combining
\eqref{p04:eq:complex-box-sandwich}--\eqref{p04:eq:analytic-box-Landau} and
$\mu(B^\pm)=(1+O_d(\tau))\mu(B_T)$ proves
\eqref{p04:eq:shrinking-all} throughout the full range $a_*<1$.

\subsection{Relative moments and zero-side localization}

Fix a nonnegative even $f\in C_c^\infty((-1/2,1/2))$, and retain
the notation $a=\int f^2$, $\psi=f^2$ and $R_\psi$ for the Weil
compression. Use the variable bandwidth
\begin{equation}\label{p04:eq:shrinking-bandwidth}
 L=\ell-C_SH-D\log\ell,\qquad X=e^L,
\end{equation}
where $C_S$ is sufficiently large and $D$ is a sufficiently large
fixed constant.  Then $L/\ell\to1$.  The support margin of the
fixed $f$ ensures that the weighted Poisson identities have no
aliases for all $\|n\|_\infty\le H$.

Proposition~\ref{p04:prop:uniform-rational-arithmetic} gives, for
$K=K^-_T$ and $D_K=\operatorname{diag}(K(\Theta_S(\alpha_k))^{1/2})$,
\begin{equation}\label{p04:eq:shrinking-relative-moments}
 \tr(D_KGD_K)=(m+o(\mu(B_T)))N,
 \qquad
 \|D_KGD_K\|_{\mathrm{HS}}^2
       \le(R_\psi m+o(\mu(B_T)))N.
\end{equation}
Here the rational height costs $\exp(O_S(H))$ are absorbed by
\eqref{p04:eq:shrinking-bandwidth}.  In checking the relative error,
the nonzero resonances must be summed using $|c_n|\le m$:
their total is bounded by $m$ times
\[
 \left(\sum_{p\in S}\sum_{k\ge1}(\log p)p^{-k/2}\right)^2,
\]
or its one-prime analogue.  A bound using only $\sum|c_n|$ would
unnecessarily lose the factor $m$.  Nonresonant archimedean factors
are kept as slowly varying functions and handled by partial
summation; replacing them by a constant with a global absolute
error would likewise be insufficient for a relative estimate.
Fourier truncation costs $o(\mu(B_T)N)$ by
\eqref{p04:eq:box-Fourier-estimates} and the unweighted trace and
Hilbert--Schmidt bounds.  These observations give the stated
uniform extension of the weighted moment proof.

We detail the additional zero-side uniformity. Choose $B_0>A+2$
and put
\[
 d_0=4B_0\frac{\log\ell}{\ell},
 \qquad I'=[T-\sqrt T,2T+\sqrt T].
\]
The total multiplicity of zeros in $I'$ with
$|\beta-1/2|>d_0$ is $O(N\ell^{-B_0})=o(\mu(B_T)N)$.
Retain all these deep zeros, regardless of their phase.  Among the
remaining zeros retain precisely those whose phase lies in $B_T$.
Let $E$ be the matrix contribution of the discarded shallow zeros.
The normalized vectors satisfy, for every fixed $M$,
\[
 |u_\rho(k)|\ll_M
 \ell^{2B_0+1}(1+L|\gamma-\alpha_k|)^{-M}.
\]
Every discarded zero is separated from every nonzero weighted
sampling center by at least
$\eta=c_S\tau r$ in real height, because $\Theta_S$ is Lipschitz.
For the matrix $U$ whose columns are the corresponding weighted
vectors, repeated according to multiplicity, the grid spacing and
unit-interval zero count give
\[
 \|U\|^2\ll\ell^{4B_0+3},\qquad
 \|U\|_{\mathrm{HS}}^2
   \ll N\ell^{4B_0+2}(L\eta)^{1-2M}.
\]
Since $E=UU^{\mathsf T}$,
\begin{align*}
 |\tr E|&\le\|U\|_{\mathrm{HS}}^2,\\
 \|E\|_{\mathrm{HS}}^2
 &\le\|U\|^2\|U\|_{\mathrm{HS}}^2
 \ll N\ell^{8B_0+6-a_*-2M(1-a_*)}
                         (\log\ell)^{2M-1}.
\end{align*}
The condition \eqref{p04:eq:shrinking-range} implies $a_*<1$.
Consequently a sufficiently large fixed $M$ makes both errors
$o(\mu(B_T)N)$, where for the Hilbert--Schmidt statement this means
$\|E\|_{\mathrm{HS}}^2=o(\mu(B_T)N)$.

Subtract also the exact contribution $F_{\rm out}$ of zeros outside
$I'$. The full-grid decay gives $\|F_{\rm out}\|_1=O(T^{-10})$
after choosing sufficiently many integrations by parts.  Thus the
exact retained matrix
\[
 H_K=D_KGD_K-F_{\rm out}-E
\]
still satisfies \eqref{p04:eq:shrinking-relative-moments}.  This exact
subtraction is necessary: a small-norm matrix need not have small
positive index.

Let $s=N^s_{0,B_T}(I')$ and decompose $H_K=P+Q$, where $P$ is the
sum over the retained simple critical-line zeros.  Full-grid
normalization gives $\rank P\le s$ and $\tr P\le s$.
Every reflected off-line pair costs at most one positive direction;
every multiple on-line zero costs at most one. Hence
\[
 s+2n_+(Q)\le N_{B_T}(I')+o(\mu(B_T)N).
\]
The two added height intervals contain only $O(\sqrt T\log T)$
zeros, uniformly in the phase condition. Applying the rank--inertia
inequality, \eqref{p04:eq:shrinking-all}, and
\eqref{p04:eq:shrinking-relative-moments}, gives
\[
 N^s_{0,B_T}(T,2T)
 \ge\bigl(2-R_\psi-o(1)\bigr)\mu(B_T)N.
\]
All constants and errors above are uniform in the box center.
Finally fix $f$ successively closer to the Montgomery--Taylor
window. For every fixed such $f$ the displayed lower bound holds;
taking the supremum over this family after the $T$ limit yields
\eqref{p04:eq:shrinking-simple}. No varying-window estimate is required.

\begin{corollary}[A covering statement at simple critical zeros]
\label{p04:cor:shrinking-covering}
Fix $S$ with $|S|=d$ and any $0<a<1$. For all sufficiently
large $T$, every phase box with side lengths $(\log T)^{-a}$ contains
the phase of a simple critical-line zero with ordinate in $(T,2T]$.
In fact each such box contains at least
$(c_{\rm MT}-o(1))(\log T)^{-ad}N(T,2T)$ such ordinates,
uniformly in its center.
\end{corollary}

This is a quantitative covering conclusion for a fixed finite prime
torus. It does not, by itself, give a rate for approximating an
infinite Euler product: such a conclusion would also require
estimates uniform as the prime set grows.




\section{Near-logarithmic phase resolution}
\label{p04:sec:near-logarithmic}

The polynomial shrinking theorem can be sharpened by choosing
Gevrey windows.  This replaces a fixed inverse power of the
sampling distance by a stretched exponential, which is needed when
the phase buffer is only a power of $\log\log T$ sampling spacings.

\begin{theorem}[Near-logarithmic boxes]
\label{p04:thm:near-logarithmic}
Fix a finite prime set $S=\{p_1,\ldots,p_d\}$, a constant $c_0>0$,
and $C>2$. Put $\ell=\log T$ and $v=\log\log T$.
Let $B_T$ be any phase box, with arbitrary center and side lengths
satisfying
\begin{equation}\label{p04:eq:nearlog-width}
 c_0\frac{v^C}{\ell}\le r_j(T)\le\frac14
       \qquad(1\le j\le d).
\end{equation}
Uniformly in the centers and all these side lengths,
\[
 N_{B_T}(T,2T)=(1+o(1))\mu(B_T)N(T,2T),
\]
and
\[
 \liminf_{T\to\infty}\inf_{B_T}
 \frac{N^s_{0,B_T}(T,2T)}{\mu(B_T)N(T,2T)}
       \ge c_{\rm MT}.
\]
Thus every such box contains many simple critical-line zero phases
in every sufficiently large dyadic height interval.
\end{theorem}

The theorem concerns a fixed prime set.  It does not reach side
length $1/\log T$, and the power $C>2$ is a sufficient condition of
this argument, not an optimality assertion.

\begin{lemma}[A complex Gevrey estimate]\label{p04:lem:complex-Gevrey}
Suppose $f\in C_c^\infty([-a,a])$ satisfies
\[
 \|f^{(m)}\|_{L^1}\le A B^m(m!)^2\qquad(m\ge0)
\]
with fixed $A,B>0$. Then there are constants $c,C>0$, depending
only on these data, such that
\begin{equation}\label{p04:eq:complex-Gevrey}
 |\widehat f(x+iy)|
 \le C\exp\left(a|y|+C\sqrt{|y|}-c\sqrt{|x|}\right)
       \qquad(x,y\in\mathbb R).
\end{equation}
There exist nonnegative even functions of this class, supported
strictly inside $(-1/2,1/2)$, whose squares approximate the
Montgomery--Taylor window in the integrals defining $R_\psi$.
\end{lemma}
\begin{proof}
Leibniz's formula, with $j$ derivatives falling on $e^{yu}$, gives
\begin{align*}
 \|(f(u)e^{yu})^{(m)}\|_1
 &\le A e^{a|y|} B^m(m!)^2
    \sum_{j=0}^m\frac{(|y|/B)^j(m-j)!}{j!m!}\\
 &\le A e^{a|y|} B^m(m!)^2
    \sum_{j=0}^\infty\frac{(|y|/B)^j}{(j!)^2}\\
 &\le A\exp\left(a|y|+2\sqrt{|y|/B}\right)B^m(m!)^2.
\end{align*}
The last inequality follows by comparing the series with
$\exp(2\sqrt{|y|/B})$, using
$\binom{2j}{j}\le4^j$. Repeated integration by parts bounds the
Fourier transform by this expression divided by $|x|^m$.
Taking $m$ to be a sufficiently small fixed multiple of $\sqrt{|x|}$
and using Stirling's bound proves \eqref{p04:eq:complex-Gevrey}; bounded
$|x|$ is absorbed into the constant.

For existence, the standard bump
$b(u)=\exp(-1/(1-u^2))$ for $|u|<1$ and $b(u)=0$ otherwise is
Gevrey of order two. One direct verification applies Cauchy's
estimate on a complex circle of radius a small constant times
$1-|u|$, on which the real part of $1/(1-z^2)$ is bounded below by
$c/(1-|u|)$. Maximizing
$(1-|u|)^{-m}\exp(-c/(1-|u|))$ gives a bound
$C^{m+1}(m!)^2$. Translates, dilates, products, and convolution
preserve this class, with constants allowed to depend on the chosen
functions. Smooth cutoffs formed from $b$ can therefore approximate
the interval indicator from within. Multiplying their even square
roots in the form $f=\chi\sqrt{\cos(\sqrt2u)}$, where
$\chi$ itself is a nonnegative even Gevrey cutoff, gives
$f^2\to\cos(\sqrt2u)1_{[-1/2,1/2]}$ in the required integrals.
The analytic square root is harmless since the cosine is strictly
positive on this interval.
\end{proof}

\begin{proof}[Proof of Theorem~\ref{p04:thm:near-logarithmic}]
Choose fixed $u,k>0$ such that
\begin{equation}\label{p04:eq:nearlog-parameters}
 u+k<C-2,
\end{equation}
and put
\[
 \eta=v^{-u},\qquad H=\ell/v^k,
 \qquad L=\ell-C_SH-Dv.
\]
Here $C_S$ and $D$ are sufficiently large fixed constants.
We use a fixed nonnegative even Gevrey window $f$ from
Lemma~\ref{p04:lem:complex-Gevrey}. All the uniform arithmetic moment
bounds apply because $H=o(\ell)$ and $L/\ell\to1$.

We first record the phase weights. Convolve slightly shrunken
interval indicators with dilates of a fixed nonnegative Gevrey-two
mollifier, with smoothing widths comparable to $\eta r_j$.
The resulting compact inner cutoff $K$ satisfies
\[
 0\le K\le1_{B_T},\quad
 m:=\int K=(1+O_d(\eta))\mu(B_T),\quad
 \operatorname{dist}(\operatorname{supp}K,\mathbb T^d\setminus B_T)
                      \gg_d\eta r_{\min}.
\]
Its Fourier coefficients $c_n$ obey $|c_n|\le m$ and
\begin{align*}
 \sum_n|c_n|&\ll_d\eta^{-d},\\
 \sum_{\|n\|_\infty>H}|c_n|
   &\ll_d\eta^{-d}\exp(-c\sqrt{H\eta r_{\min}}).
\end{align*}
Polynomial factors in the tail summation have been absorbed by
reducing $c$.  These estimates follow by multiplying the interval
coefficient bound by the Gevrey Fourier decay of the mollifier and
summing the resulting product. By \eqref{p04:eq:nearlog-width} and
\eqref{p04:eq:nearlog-parameters},
\[
 H\eta r_{\min}\gg v^{C-u-k},\qquad C-u-k>2.
\]
The Fourier tail is consequently $o(\ell^{-M})$ for every fixed
$M>0$. Since $\mu(B_T)\gg\ell^{-d}$, Fourier truncation and the
uniform rational-character arithmetic give relative first and
second moment estimates for this $K$, exactly as in
\eqref{p04:eq:shrinking-relative-moments}.

We next prove the required all-zero budget. Use the Gaussian
majorants and minorants from the preceding section, now expanding
or contracting each side by $4\eta r_j$ and taking widths
\[
 s_j=\frac{\eta r_j}{\sqrt{K_0v}},
\]
where the fixed constant $K_0$ is sufficiently large in terms of
$d$. Their Fourier coefficient mass is
$O_{d,K_0}(v^{d(u+1/2)})$. For shallow zeros with
$|\beta-1/2|\le d_0=4B_0v/\ell$,
\[
 h:=\|y/s\|\ll_{S,B_0,K_0}v^{3/2+u-C}.
\]
The complementary Gaussian tails on the target box and outside the
contracted box are $O_d(e^{-c_dK_0v})$. Choose $K_0$ so these are
$o(\ell^{-d-2})$. The global negative part of the upper analytic
majorant is at most
\[
 C_d\left(\prod_js_j^{-1}\right)
       \exp(-c/h^2)
 \ll\ell^d v^{O_d(1)}
       \exp(-c v^{2C-3-2u}).
\]
Since $C-u>2$, the exponent $2C-3-2u$ exceeds one, so this error
is $o(\ell^{-M})$ for every fixed $M$. The lower analytic bound's global
upper bound remains $e^{h^2/2}=1+o(1)$.
The Gaussian Fourier truncation at $H$ is also super-polynomially
accurate in $\ell$, since
\[
 s_{\min}H\gg v^{C-u-k-1/2},\qquad C-u-k>2,
 \qquad d_0/s_{\min}^2=o(H).
\]
The last relation ensures that the growing imaginary factor in
the tail is absorbed by Gaussian decay.

Apply Landau's formula to the truncated analytic polynomials at
the complex zero coordinates. Exact prime-power resonances have
total contribution $O_S(\mu(B_T)N/\ell)$, and the rational-height
error is $T^{o(1)}v^{O_d(1)}=o(\mu(B_T)N)$. The weighted
contribution of deep zeros is at most
\[
 Nv^{O_d(1)}\exp(-d_0\ell/8)
       =Nv^{O_d(1)}\ell^{-B_0/2}.
\]
Taking $B_0>2d+4$ makes it $o(\mu(B_T)N)$. The analytic
sandwich therefore proves the asserted all-zero asymptotic,
uniformly in the box.

It remains to localize the zero-side inertia with the smaller
buffer. Set $\delta_\rho=\beta-1/2$. Lemma~\ref{p04:lem:complex-Gevrey}
gives for the normalized sampling vector
\[
 |u_\rho(k)|\ll_f
 \exp\left(C_fL|\delta_\rho|
                    -c_f\sqrt{L|\gamma-\alpha_k|}\right).
\]
An increase in the constant covers bounded $L|\delta_\rho|$.
For a discarded shallow zero, the real separation from a weighted
sampling center is at least
$c_S\eta r_{\min}$. Thus
\[
 L|\gamma-\alpha_k|\gg v^{C-u}.
\]
The unit-interval zero count and the grid spacing give Schur bounds
for the matrix $U$ of discarded vectors of the form
$\|U\|^2\ll\ell^{C_fB_0+C_f}$. Summing the stretched-exponential
tail on the grid gives
\[
 \|U\|_{\mathrm{HS}}^2
 \ll N\ell^{C_fB_0+C_f}
                    \exp(-c_f v^{(C-u)/2}).
\]
Because $(C-u)/2>1$, both
$|\tr(UU^{\mathsf T})|$ and
$\|UU^{\mathsf T}\|_{\mathrm{HS}}^2$
are $o(N\ell^{-M})$ for every fixed $M$. They are in particular
$o(\mu(B_T)N)$. The exact outer-height tail is subtracted as
before, and the deep zeros cost $o(\mu(B_T)N)$ in the count budget.

The rank--inertia argument now gives
\[
 N^s_{0,B_T}(T,2T)
       \ge(2-R_{f^2}-o(1))\mu(B_T)N(T,2T).
\]
Finally optimize through the fixed Gevrey windows supplied by
Lemma~\ref{p04:lem:complex-Gevrey}, after taking the $T$ limit. This
recovers $c_{\rm MT}$ without requiring any estimate uniform in
the approximating window.
\end{proof}


\begin{corollary}[A quantitatively equidistributed selection]
\label{p04:cor:equidistributed-selection}
Fix a finite prime set $S$, $0<c<c_{\rm MT}$, and $C>2$.
For every sufficiently large $T$ one can select a set
$\mathcal Z_T$ of simple critical-line zero ordinates in $(T,2T]$
with $|\mathcal Z_T|=(c+o(1))N(T,2T)$ such that its normalized
phase distribution has box discrepancy
\[
 \sup_{B\text{ a phase box}}
 \left|\frac{\#\{\gamma\in\mathcal Z_T:F_S(\gamma)\in B\}}
                   {|\mathcal Z_T|}-\mu(B)\right|
 \ll_{S,c,C}\frac{(\log\log T)^C}{\log T}.
\]
This is a selected subset; equidistribution of all simple
critical-line zeros is not asserted.
\end{corollary}
\begin{proof}
Partition the torus into $m^d$ equal boxes, where
$m\asymp\log T/(\log\log T)^C$. Theorem~\ref{p04:thm:near-logarithmic}
ensures each cell contains more than $cN/m^d$ suitable ordinates.
Select exactly $\lfloor cN/m^d\rfloor$ in every cell, using
increasing height to make the selection definite. There are
$(c+o(1))N$ selected ordinates. Any phase box meets at most
$O_d(m^{d-1})$ boundary cells. Fully included cells carry exactly
their Haar fraction of the selected measure, and the boundary
cells give discrepancy $O_d(m^{-1})$.
\end{proof}




\section{From phase-localized zeros to Euler-edge triples}
\label{p03:sec:edge-triples}Let $\Omega=\prod_p\T$ with product Haar measure $\mu$. For
$z\in\Omega$ put
\[
 Z_z(s)=\prod_p(1-z_pp^{-s})^{-1},\qquad
 Q_z(s)=-\frac{Z_z'(s)}{Z_z(s)}
       =\sum_p\frac{(\log p)z_pp^{-s}}{1-z_pp^{-s}}.
\]
All these expressions converge normally in $\Re s>1$.
The frozen configuration is $z_2^*=-1$, $z_p^*=1$ for odd $p$.
The classical constants $A_*$ and $E_*$ are characterized by
\begin{align*}
 \sum_{\substack{n\ge3\\n\text{ odd}}}n^{-A_*}&=2^{-A_*},\qquad
 A_*=1.940101\ldots,\\
 \sum_{p\ge3}\frac{\log p}{p^{E_*}-1}
 &=\frac{\log2}{2^{E_*}+1},\qquad E_*=2.813014\ldots.
\end{align*}
These are the right boundaries for $Z_z(s)=1$ and $Q_z(s)=0$;
see \cite{AriasLune}. In particular no such point has real part
larger than its boundary. Their identification is not a new claim.
The frozen functions have simple roots at these boundaries:
$a_E=Q_{z^*}'(E_*)<0$ and $a_A=Z_{z^*}'(A_*)<0$.
The inequalities also follow directly by differentiating the defining
positive sums. Define
\[
 D_E(z)=\sum_p(\log p)p^{-E_*}|z_p-z_p^*|^2,\qquad
 D_A(z)=\sum_pp^{-A_*}|z_p-z_p^*|^2,
\]
and put $\beta_E=(E_*-1)^{-1}$,
$\tau=(A_*-1)/(E_*-1)$ and
$\Psi(\delta)=\delta^\tau/\log(1/\delta)$.

\begin{lemma}[Local edge estimates]\label{p03:lem:local-edges}For sufficiently small $D=D_E(z)>0$, there are simple roots
$Q_z(E_*+u_E)=0$ and $Z_z(A_*+u_A)=1$, unique in fixed
small disks about the respective boundary points, such that
\begin{align*}
 0\le-\Re u_E&\le CD,& |\Im u_E|&\le C\sqrt D,\\
 0\le-\Re u_A&\le C\Psi(D),& |\Im u_A|&\le C\sqrt D.
\end{align*}
Conversely, a critical point at $E_*-\delta$ in a twist implies
$D_E(z)\le C\delta$, and a $1$-point at $A_*-\eps$ implies
$D_A(z)\le C\eps$, when the deficits are sufficiently small.
Here the phases are evaluated at the ordinate of the point.
\end{lemma}
\begin{proof}
The converse assertions are Lemmas~\ref{p02:eul:rigidity} and
\ref{p02:eul:one-rigidity}. Critical-root realization, including its
quadratic real-part estimate, is Lemma~\ref{p02:eul:local-critical}.
For completeness, the passage from a phase defect to the other
local root uses the estimates already established in the proof of
Theorem~\ref{p02:eul:locking}. In \eqref{p02:eul:interpolation} take
$x=D^{-1/(E_*-1)}$ to obtain $D_A\ll\Psi(D)$. The inequality
$2A_*-E_*>1$, certified in Lemma~\ref{p02:lem:edge-enclosures},
gives $\|Z_z-Z_{z^*}\|_{C^1}=O(\sqrt D)$ on a fixed small disk
about $A_*$. At its center,
$\Re(Z_z(A_*)-1)=O(D_A)$, so Taylor's formula gives
\[
 u_A=-\frac{Z_z(A_*)-1}{a_A}+O(D),\qquad
 |u_A|\ll\sqrt D,\qquad
 |\Re u_A|\ll D_A+D\ll\Psi(D).
\]
The two global boundary inequalities give the signs. Choose the
fixed disks to contain only the simple frozen roots. The same
$C^1$ estimates at $E_*$ and $A_*$ imply uniform convergence to
these frozen germs as $D\to0$; Rouch\'e's theorem gives exactly
one root in each disk, and convergence of their derivatives proves
simplicity. Thus all the local estimates have the same fixed-disk
interpretation used below.
\end{proof}

\begin{proposition}[Infinite phase regions]\label{p03:prop:infinite-phases}
Let $U\subset\Omega$ be open. The lower density of simple
critical-line zero ordinates whose full phase vector belongs to
$U$ is at least $c_{\rm MT}\mu(U)$, relative to $N(T,2T)$.
For a continuity set of product Haar measure, the density of all
zero ordinates is its Haar measure.
\end{proposition}
\begin{proof}
Finite cylinder functions are uniformly dense in $C(\Omega)$.
The finite-dimensional all-zero equidistribution thus implies weak
convergence of the full all-zero phase measures to $\mu$.
Every open set is the increasing union of finite unions of open
cylinders with Haar-null boundaries, up to an arbitrarily small
loss of measure. Apply Theorem~\ref{p03:thm:main-phase} to each such
finite union and then use inner regularity. The assertion for all
zeros follows from weak convergence on continuity sets.
\end{proof}

\begin{corollary}[Euler-product patterns along simple critical zeros]
\label{p03:cor:Euler-patterns}
Let $\omega\in\Omega$, let $K$ be a compact subset of $\Re s>1$,
let $m\ge0$ be an integer and let $\eps>0$. Then a positive lower
proportion, relative to all zeros in $(T,2T]$, of simple critical-line
zero ordinates $\gamma$ satisfy
\[
 \max_{0\le j\le m}\sup_{s\in K}
 |\zeta^{(j)}(s+i\gamma)-Z_\omega^{(j)}(s)|<\eps.
\]
More precisely, every weak limit $\nu$ of the empirical subprobability
measures of simple critical-line prime phases satisfies
$c_{\rm MT}\mu\le\nu\le\mu$ on $\Omega$.
\end{corollary}
\begin{proof}
For nonnegative continuous finite-cylinder functions, the lower
integral inequality follows from Theorem~\ref{p03:thm:main-phase} by
finite step approximations on Haar-continuity rectangles. Uniform
approximation extends it to every nonnegative function in
$C(\Omega)$. The upper inequality follows because the corresponding
all-zero phase measures converge to $\mu$ and dominate the
simple-critical submeasures. This proves the assertion about weak
limits.

Normal convergence of the Euler products and their derivatives
makes the map $z\mapsto (Z_z^{(j)}|_K)_{j\le m}$ continuous in
the product topology. The displayed approximation condition defines
an open neighborhood of $\omega$, hence has positive Haar measure.
Proposition~\ref{p03:prop:infinite-phases} supplies the claimed lower
proportion.
\end{proof}

The target class in Corollary~\ref{p03:cor:Euler-patterns} consists of
unimodular Euler-product twists in the half-plane of absolute
convergence. This is distinct from approximation of arbitrary
holomorphic targets in a Voronin universality strip.

\begin{theorem}[Critical-line zeros with two edge companions]
\label{p03:thm:triples}
Put $B_\delta=\{z:D_E(z)<\delta\}$ and let $M_\delta(T)$ count
simple critical-line zeros with $T<\gamma\le2T$ and
$(p^{-i\gamma})_p\in B_\delta$. For every fixed sufficiently small
$\delta>0$,
\begin{equation}\label{p03:eq:phase-density-sandwich}
 c_{\rm MT}\mu(B_\delta)\le
 \liminf_{T\to\infty}\frac{M_\delta(T)}{N(T,2T)}
 \le\limsup_{T\to\infty}\frac{M_\delta(T)}{N(T,2T)}
 \le\mu(B_\delta).
\end{equation}
Every counted ordinate has companions $w_E,w_A$ satisfying
\[
 \zeta'(w_E)=0,\quad\zeta(w_A)=1,
\]
\[
 0<E_*-\Re w_E\le C\delta,\qquad
 0<A_*-\Re w_A\le C\Psi(\delta),\qquad
 |\Im w_E-\gamma|+|\Im w_A-\gamma|\le C\sqrt\delta.
\]
Both companion roots are simple. Moreover, for some $c,C>0$,
\begin{equation}\label{p03:eq:energy-entropy}
 \exp\!\left(-\frac{C\delta^{-\beta_E}}{\log(1/\delta)}\right)
 \le\mu(B_\delta)\le
 \exp\!\left(-\frac{c\delta^{-\beta_E}}{\log(1/\delta)}\right).
\end{equation}
\end{theorem}
\begin{proof}
The continuous energy $D_E$ is a uniformly convergent positive
series. Its distribution has no atoms: condition on every coordinate
except one, whose nonconstant cosine energy has an atomless
distribution. Thus $B_\delta$ is a continuity set, and
Proposition~\ref{p03:prop:infinite-phases} proves
\eqref{p03:eq:phase-density-sandwich}.
For $z_p=p^{-i\gamma}$ one has $Z_z(s)=\zeta(s+i\gamma)$.
Lemma~\ref{p03:lem:local-edges} gives the companions. Equality at either
boundary would force the exact frozen configuration, impossible
for an ordinary prime flow: the phases at $3$ and $5$ force
$\gamma=0$, while the phase at $2$ then has the wrong sign.

For completeness we give the volume estimate. Write
$x=K\delta^{-\beta_E}$, with a sufficiently large fixed $K$.
The tail energy beyond $x$ is at most
$4\sum_{p>x}(\log p)p^{-E_*}\ll x^{1-E_*}$.
For $p\le x$ restrict the angular distance to the frozen phase by
\[
 |\theta_p|\le c_0(p/x)^{E_*/2}
                  (\log x/\log p)^{1/2}.
\]
For small fixed $c_0$ this box has total energy less than $\delta$.
Its negative logarithmic measure is $O(x/\log x)$, since
\[
 \sum_{p\le x}\log(x/p)=\int_2^x\frac{\pi(u)}u\,du
 =O(x/\log x),\qquad \pi(x)=O(x/\log x).
\]
The logarithm-of-logarithm term only reduces this upper bound.
This proves the lower volume bound.

For the upper bound choose instead a sufficiently small fixed
multiple $x$ of $\delta^{-\beta_E}$ and keep primes in $[x,2x]$.
Their number is $m\asymp x/\log x$, each weight is comparable to
$x^{-E_*}\log x$, and their expected total energy is a sufficiently
large multiple of $\delta$. The independent variables
$1-\cos\theta_p$ have a fixed nondegenerate distribution.
Their Laplace transform at a fixed positive argument is strictly
less than one. Exponential Markov therefore bounds the probability
of energy below $\delta$ by $e^{-cm}$, proving the other inequality.
\end{proof}

Taking a diagonal sequence $\delta_j\downarrow0$ gives triples with
ordinates tending to infinity, with both ordinate displacements
tending to zero and both real parts approaching their boundaries.
Different critical-line zeros need not have different companions.
The limits in \eqref{p03:eq:phase-density-sandwich} always hold for
fixed $\delta$; no asserted error term allows substitution of an
arbitrary shrinking function $\delta(T)$.




For the growing-prime and probabilistic sections below, abbreviate
$E=E_*$ and $A=A_*$. The local maps $u_E,u_A$ always mean the
unique simple roots in the fixed disks specified in
Lemma~\ref{p03:lem:local-edges}.



\section{A simultaneous height law from a growing set of primes}
\label{p04:sec:growing-primes}

The fixed-dimensional theorem does not by itself supply an error rate
for an infinite Euler product. We now let the number of controlled
primes grow, keeping an explicit account of the Fourier cost. This
produces an actual joint limit in the height and the edge tolerance.
Write $\ell=\log T$ and $v=\log\ell$ throughout this section.

\begin{lemma}[A dimension-independent all-zero estimate]
\label{p04:lem:growing-landau}
Let $F(z)=\sum c_{\mathbf n}z^{\mathbf n}$ be a finite Fourier
polynomial on a prime torus, $m=c_{\mathbf0}$, and suppose
$|c_{\mathbf n}|\le m$. Put
\[
 A_0=\sum|c_{\mathbf n}|,\qquad
 A_2=\sum|c_{\mathbf n}|
       \left(\sum_p n_p\log p\right)^2.
\]
If all associated reduced rationals have numerator and denominator
at most $H$, and $\log H\le\ell/16$, then
\[
 \sum_{T<\gamma\le2T}F(F_S(\gamma))
 =mN+O\left(\frac{mNB_S}{\ell}
       +\frac{NA_2}{\ell^2}+A_0H^{C}\ell^2\right).
\]
Zeros are counted with multiplicity. The absolute exponent $C$ and
the implied constant are independent of $S$.
\end{lemma}
\begin{proof}
For a reduced rational $q>1$ with numerator and denominator at most
$H$, the uniform Landau formula of \cite[Lemma~1]{FordZaharescu}
gives
\[
 \sum_{T<\gamma\le2T}q^{\rho-1/2}
 =-\frac{T}{2\pi}\frac{\Lambda(q)}{\sqrt q}
   +O(H^{C}\ell^2),
\]
where $\Lambda(q)=0$ unless $q$ is an integral prime power.
Indeed, if $q$ is not that nearest prime power, rational separation
bounds the extra oscillatory term by a fixed power of $H$; also
$\log q\gg H^{-1}$. Reflection $\beta\mapsto1-\beta$ pairs the
terms with the same ordinate. If $\lambda=\log q$, their difference
from $q^{i\gamma}$ is controlled by
$\cosh(\lambda(\beta-1/2))-1$.
The uniform Selberg bound
$N(1/2+u,3T)\ll N e^{-u\ell/4}$ and integration by parts give
\[
 \sum_{T<\gamma\le2T}
  |\beta-1/2|^2 e^{|\lambda||\beta-1/2|}
 \ll N\ell^{-2},\qquad |\lambda|\le\ell/16.
\]
Thus replacing $q^{\rho-1/2}$ by $q^{i\gamma}$ costs
$O(N\lambda^2/\ell^2)$. For $q<1$ take the complex conjugate
of the corresponding formula for $q^{-1}$. Sum over characters.
The prime-power main terms have total modulus at most
$2mTB_S$; this proves the assertion.
\end{proof}

\begin{theorem}[Growing-prime phase boxes]
\label{p04:thm:growing-boxes}
Fix $\sigma>1$ and a sufficiently small $b>0$. There exists
$c_0>0$ such that the following holds. Put $x=c_0\log\log T$,
$S_x=\{p:p\le x\}$, and
\[
 r_p=b(p/x)^{\sigma/2}
           \sqrt{\frac{\log x}{\log p}}.
\]
For arbitrary centers $\omega_p\in\T$, let $B_T$ be the product
of the angular arcs $|\arg(z_p/\omega_p)|<r_p$. Uniformly in
these centers,
\[
 N_{B_T}=(1+o(1))\mu(B_T)N,\qquad
 N^s_{0,B_T}\ge(c_f-o(1))\mu(B_T)N,
\]
where $c_f=2-R_f(1)>0$ for any one fixed smooth test window
with this property. In particular a fixed positive constant is
available. Moreover
\[
 \exp(-C_2x/\log x)\le\mu(B_T)
       \le\exp(-C_1x/\log x).
\]
The constants and $c_0$ may depend on the fixed parameters and
test window. No optimality of the height scale is asserted.
\end{theorem}
\begin{proof}
For large $x$ all radii are less than a fixed small constant, and
$r_{\min}\ge x^{-C}$. The prime number theorem and partial
summation give
\[
 \sum_{p\le x}\log(x/p)\ll x/\log x,
 \qquad
 \sum_{p\le x}\log(\log x/\log p)\ll x/\log x.
\]
Together with $\pi(x)\asymp x/\log x$, these prove the stated
Haar-volume bounds (the lower bound also follows by dropping
the favorable square-root factor).

 Let $\eta=x^{-2}$. Construct smooth product cutoffs $K^-$ and
$K^+$ supported in the box and in its relative $O(\eta)$
enlargement, respectively, such that
$0\le K^-\le\mathbf1_{B_T}\le K^+$ and each cutoff changes
only across a relative boundary layer of width $\eta$. In particular,
choose $K^-=1$ on the box with radii $(1-3\eta)r_p$ and
$\supp K^-\subset\{ |\arg(z_p/\omega_p)|<(1-2\eta)r_p\}$;
this leaves a genuine buffer to the target boundary. For each
coordinate use a rescaled fixed smooth transition. Then
\[
 \int K^\pm=(1+o(1))\mu(B_T),
\]
since $\pi(x)\eta=o(1)$. Every fixed-order derivative norm of an
individual factor is bounded by a fixed power of $x$. Fourier
coefficient estimates of order four, followed by multiplication
over the coordinates, imply
\[
 \sum|\widehat K^\pm(\mathbf n)|+
 \sum|\widehat K^\pm(\mathbf n)|
              \left(\sum n_p\log p\right)^2\le e^{C x}.
\]
The constants are independent of the centers. For example, apply
the one-dimensional weighted $\ell^1$ bounds to one or two
coordinates and the unweighted bound to the other coordinates;
the factor $\pi(x)^2\log^2x$ is absorbed in $e^{Cx}$.

Truncate each coordinate at degree $K=\lceil e^{C'x}\rceil$.
Using a sufficiently high fixed derivative order and then choosing
$C'$ large makes the uniform truncation error at most $e^{-C''x}$,
with $C''$ as large as required. This is $o(\mu(B_T))$.
The retained coefficients obey
$|c_{\mathbf n}|\le\int K^\pm$, since the original cutoffs are
nonnegative. Their rational heights satisfy
\[
 \log H\le K\sum_{p\le x}\log p\le e^{C_3x}.
\]
Choose $c_0$ so small that $C_3c_0<1/2$ and all the constants
in the Fourier second-moment bounds satisfy $Cc_0<1$.
Lemma~\ref{p04:lem:growing-landau} then gives relative all-zero
asymptotics for $K^\pm$, and hence for the box and all the thin
enlargements used here: $A_2/\ell^2=o(\mu(B_T))$, while
$H^C=T^{o(1)}$ and $B_{S_x}\ll\sqrt x$.

For the Weil matrix choose
\[
 L=\ell-2\log H-C_4x-D\log\ell,
\]
with fixed large $C_4,D$. Then $L/\ell\to1$.
Proposition~\ref{p04:prop:uniform-rational-arithmetic}, applied to the
Fourier truncation of $K^-$, gives
\[
 \tr(D_{K^-}GD_{K^-})=(1+o(1))\mu(B_T)N,
 \|D_{K^-}GD_{K^-}\|_{\rm HS}^2
 \le(R_f(1)+o(1))\mu(B_T)N.
\]
The cutoff itself, rather than its Fourier approximation, is used
in the coordinate matrix. The uniform approximation error changes
these estimates by $o(\mu(B_T)N)$, using the unweighted norm
bound. The factors $e^{Cx}$ in the arithmetic errors are absorbed
by $C_4$ and $D$.

For clarity, the zero-side localization also remains uniform.
The prime-phase map has Lipschitz constant
$C_{S_x}\ll\sqrt{x\log x}$. The inner cutoff and the boundary of
$B_T$ are separated in physical height by
\[
 r_*\gg\frac{\eta r_{\min}}{C_{S_x}}\ge x^{-C_5}.
\]
Retain every zero with $|\beta-1/2|>B\log\ell/\ell$,
and every shallow zero whose phase lies in $B_T$.
The former have total multiplicity $O(N\ell^{-B/4})$, hence
$o(\mu(B_T)N)$, because $\mu(B_T)=\ell^{-o(1)}$.
For the discarded shallow zeros, repeated integration by parts
in the compactly supported test window gives a coordinate bound
\[
 |u_\rho(k)|\ll_M\ell^{B/2+1}
                (1+L|\gamma-\alpha_k|)^{-M}.
\]
The unit-interval zero count and Schur's test therefore give
arbitrarily large fixed power savings in $\ell$, by choosing
a sufficiently large fixed $M$: here $Lr_*\ge\ell x^{-C_5}$.
Their trace and squared Hilbert--Schmidt errors are
$o(\mu(B_T)N)$. The outer height tail is subtracted exactly,
as in the localization proof, and has trace norm $o(1)$.
The rank--inertia inequality, with the all-zero count of $B_T$
as budget, now gives the stated simple-zero lower bound.
\end{proof}

Write $D_E(t)=\sum_p(\log p)p^{-E}|p^{-it}-z_p^*|^2$,
where $z_2^*=-1$ and $z_p^*=1$ for odd $p$. The constants $E,A$
and the local root estimates are established in
Section~\ref{p03:sec:edge-triples}.

\begin{theorem}[Simultaneous shrinking Euler-edge triples]
\label{p04:thm:simultaneous-triples}
There are positive constants $C_0,C_1,c_1$ with the following
property. Put $\delta_T=C_0(\log\log T)^{1-E}$ and let $M(T)$
count simple critical-line zero ordinates $\gamma\in(T,2T]$
with $D_E(\gamma)\le\delta_T$. Then, for every sufficiently large $T$,
\begin{equation}\label{p04:eq:simultaneous-rarity}
 N\exp\!\left(-C_1\frac{\log\log T}{\log\log\log T}\right)
 \le M(T)\le
 N\exp\!\left(-c_1\frac{\log\log T}{\log\log\log T}\right).
\end{equation}
Each counted ordinate has simple companions $w_E,w_A$ satisfying
$\zeta'(w_E)=0$ and $\zeta(w_A)=1$, with
\begin{align*}
 0<E-\Re w_E&\ll(\log\log T)^{1-E},\\
 0<A-\Re w_A&\ll
       (\log\log T)^{1-A}/\log\log\log T,\\
 |\Im w_E-\gamma|+|\Im w_A-\gamma|
       &\ll(\log\log T)^{(1-E)/2}.
\end{align*}
Different ordinates are not asserted to have distinct companions.
This is a simultaneous limit, not an iterated fixed-tolerance limit.
\end{theorem}
\begin{proof}
Use Theorem~\ref{p04:thm:growing-boxes} with $\sigma=E$ and frozen
centers. Inside its box,
\[
 \sum_{p\le x}(\log p)p^{-E}|z_p-z_p^*|^2
 \le b^2x^{-E}\log x\,\pi(x)\ll x^{1-E}.
\]
The remaining tail is at most
$4\sum_{p>x}(\log p)p^{-E}\ll x^{1-E}$.
Taking $C_0$ large enough proves the lower count and all companion
claims by the local edge lemma.

We prove the upper count, so that the asserted rarity has the same
logarithmic order in this joint limit. Put
$y=a\delta_T^{-1/(E-1)}$, where $a>0$ is a sufficiently small
fixed number, and restrict energy to primes in $[y,2y]$.
Choose $s=a_1y^E/\log y$ with fixed $a_1>0$ and let
\[
 K(z)=\exp\left\{-s\sum_{y\le p\le2y}
                    (\log p)p^{-E}|z_p-z_p^*|^2\right\}.
\]
If $h(u)=2u-\log I_0(2u)$, then
$m=\int K=\exp(-\sum h(s(\log p)p^{-E}))$.
Each argument of $h$ lies in a fixed positive compact interval;
hence $m\le\exp(-c y/\log y)$. Choose $a$ small enough that
$s\delta_T\le(c/2)y/\log y$.

The Bessel Fourier expansion of each factor has nonnegative
coefficients before rotation by $z_p^*$. Consequently its absolute
Fourier coefficients form a probability distribution, and
\[
 A_0=1,\qquad
 A_2=2s\sum_{y\le p\le2y}(\log p)^3p^{-E}
       \ll y\log y.
\]
Truncate each coordinate at degree $\lceil y^2\rceil$.
The factorial tail bound for $I_n(u)$, uniform for the bounded
arguments here, makes the uniform error $\exp(-c'y^2\log y)$.
The rational height satisfies $\log H\ll y^3=o(\ell)$.
Lemma~\ref{p04:lem:growing-landau} therefore gives
\[
 \sum_{T<\gamma\le2T}K(F(\gamma))
 \le (1+o(1))mN+O(Ny\log y/\ell^2)+T^{o(1)}.
\]
Multiply by $e^{s\delta_T}$ and use the pointwise exponential
majorant for the event $D_E\le\delta_T$. Since
$y\asymp\log\log T$, all the error terms are smaller than
$N\exp(-c''y/\log y)$. This bounds even the count of all zeros
with the energy condition and proves the upper bound for $M(T)$.
\end{proof}

\begin{remark}
Equation~\eqref{p04:eq:simultaneous-rarity} counts zeros selected by
the stated phase-energy condition. It does not count every possible
triple with comparable displacements, and it does not assert that
the displayed approach rate is the fastest attainable one.
\end{remark}




\section{Bounded-density transfer and conditioned phase laws}
\label{p04:sec:bounded-transfer}

The preceding localization theorem does not identify a unique phase
distribution for simple critical-line zeros.  Its two-sided domination
nevertheless determines every diverging logarithmic probability scale.
It also transports concentration statements under arbitrarily rare
conditioning.  The elementary distinction between these conclusions
and convergence to a nondegenerate law is useful here.

\begin{theorem}[Bounded-density transfer]\label{p04:thm:bounded-transfer}
Let $\mu$ be a probability measure on a measurable space $\Omega$,
let $0<c<1$, and let $\nu$ be a finite measure satisfying
$c\mu\le\nu\le\mu$.  Then the following statements hold uniformly
over all such $\nu$.
\begin{enumerate}[label=\textup{(\roman*)}]
\item For every nonnegative measurable $F$,
\[
 c\int F\,d\mu\le\int F\,d\nu\le\int F\,d\mu.
\]
In particular, whenever the integrals are positive and finite,
\[
 \left|\log\int F\,d\nu-\log\int F\,d\mu\right|
 \le |\log c|.
\]
\item Let $W\ge0$ with $0<\int W\,d\mu<\infty$, and normalize
$W\mu$ and $W\nu$ to probability measures $\mu_W$ and $\nu_W$.
For a measurable event $A$, writing $p=\mu_W(A)$, one has
\begin{equation}\label{p04:eq:conditioned-sharp-bounds}
 \frac{cp}{1-(1-c)p}\le\nu_W(A)
 \le\frac{p}{c+(1-c)p}.
\end{equation}
Both bounds are sharp.  In particular,
$c\mu_W(A)\le\nu_W(A)\le c^{-1}\mu_W(A)$.
\item With total variation defined by
$\operatorname{TV}(P,Q)=\sup_A|P(A)-Q(A)|$, one has the sharp bound
\begin{equation}\label{p04:eq:sharp-TV-transfer}
 \operatorname{TV}(\nu_W,\mu_W)
 \le\frac{1-\sqrt c}{1+\sqrt c}.
\end{equation}
Moreover, putting $a_c=-c\log c/(1-c)$,
\begin{equation}\label{p04:eq:sharp-KL-transfer}
 D_{\rm KL}(\nu_W\Vert\mu_W)
 \le a_c-1-\log a_c,
\end{equation}
again with equality possible.
\end{enumerate}
\end{theorem}
\begin{proof}
The Radon--Nikodym theorem gives $d\nu=h\,d\mu$ with
$c\le h\le1$ almost everywhere, proving the first assertion.
For the second, put $a=\int_Ah\,d\mu_W$ and
$b=\int_{A^c}h\,d\mu_W$.  Then
$cp\le a\le p$, $c(1-p)\le b\le1-p$, and
$\nu_W(A)=a/(a+b)$.  Monotonicity in $a,b$ proves
\eqref{p04:eq:conditioned-sharp-bounds}, and the endpoint-valued
choices of $h$ attain the bounds.

Write $\alpha=\int h\,d\mu_W\in[c,1]$.  The secant-line
bound for the convex function $|x-\alpha|$ on $[c,1]$ gives
\[
 \operatorname{TV}(\nu_W,\mu_W)
 =\frac1{2\alpha}\int|h-\alpha|\,d\mu_W
 \le\frac{(\alpha-c)(1-\alpha)}{\alpha(1-c)}.
\]
The last expression is maximal at $\alpha=\sqrt c$ and gives
\eqref{p04:eq:sharp-TV-transfer}.  Equality holds when $h$ only takes
the values $c,1$ and has this mean.  Likewise, the secant-line
bound for $x\log x$ yields
\[
 D_{\rm KL}(\nu_W\Vert\mu_W)
 \le -\log\alpha+
       \frac{c(1-\alpha)\log c}{\alpha(1-c)}.
\]
This is maximal at $\alpha=a_c$, where its value is
$a_c-1-\log a_c$.  An endpoint-valued density with mean $a_c$
attains equality.
\end{proof}

\begin{corollary}[Logarithmic laws and concentration]
\label{p04:cor:LD-transfer}
Suppose $a_r\to\infty$.  Under the hypotheses of
Theorem~\ref{p04:thm:bounded-transfer}, any limit
\[
 \frac1{a_r}\log\int e^{V_r}\,d\mu
\]
is unchanged when $\mu$ is replaced by $\nu$.
If $-\log\mu(A_r)\sim I a_r$ for $I>0$, then
$-\log\nu(A_r)\sim I a_r$ uniformly over admissible $\nu$.
More generally any large-deviation principle for measurable maps
$X_r:\Omega\to E$ is preserved, after normalizing $\nu$ if needed.

The same conclusions hold conditionally, with $\mu_{W_r}$ and
$\nu_{W_r}$ in place of $\mu$ and $\nu$.  In particular, if
$X_r-x_r\to0$ in $\mu_{W_r}$-probability for deterministic $x_r$,
then the convergence holds in $\nu_{W_r}$-probability uniformly
over admissible $\nu$.
\end{corollary}
\begin{proof}
The logarithms in the first statements differ by at most a constant.
After normalization or conditioning, the comparison factors are
between $c$ and $c^{-1}$.  Division by $a_r$ removes these constants
in the open- and closed-set bounds defining a large-deviation
principle.  For concentration, apply the upper bound to each
deviation event.
\end{proof}

\begin{remark}[What the transfer does not determine]
Bounded-density comparison does not preserve a nondegenerate weak
limit.  For example, let $\Omega=[0,1]^2$ with Lebesgue measure,
condition on $0<x<\varepsilon$, and set $h=c$ for $y<1/2$ and
$h=1$ for $y\ge1/2$.  The conditional $y$-law remains uniform
under $\mu$, but has two different densities under $h\mu$ for
every $\varepsilon>0$.  Thus concentration at a deterministic
profile transfers; a Gaussian fluctuation law requires additional
information about the density.
\end{remark}

For the application, let $\Omega=\prod_p\mathbb T$ with Haar
measure $\mu$ and define
\[
 \eta_T=\frac1{N(T,2T)}
 \sum_{\substack{T<\gamma\le2T\\
 \zeta(1/2+i\gamma)=0,\ \zeta'(1/2+i\gamma)\ne0}}
 \delta_{(p^{-i\gamma})_p}.
\]
The fixed-box case of Theorem~\ref{p04:thm:shrinking-boxes} gives
$c_{\rm MT}\mu\le\nu\le\mu$ for every weak limit $\nu$ of
$\eta_T$. To justify this on the full product, first use disjoint
finite boxes to approximate nonnegative cylinder functions from below.
The all-zero phase measures converge to Haar by the same box theorem.
Cylinder functions are uniformly dense in $C(\Omega)$, so both
inequalities pass to every weak limit on the compact product torus.  The preceding theorem therefore applies with
$c=c_{\rm MT}$.  The sharp total-variation and relative-entropy
bounds are, respectively, $0.0988640259\ldots$ and
$0.0196336622\ldots$ nats.  These are bounds for weak limits,
including their conditional laws.  They are not bounds for the
finite atomic measures $\eta_T$ in total variation.

\begin{corollary}[Iterated rare-event limits]
\label{p04:cor:iterated-rare-transfer}
Let $B_\varepsilon$ be Haar-continuity sets with positive measure,
and suppose
$-\log\mu(B_\varepsilon)\sim J(\varepsilon)\to\infty$.
Put
\[
 d_-(\varepsilon)=\liminf_{T\to\infty}\eta_T(B_\varepsilon),
 \qquad
 d_+(\varepsilon)=\limsup_{T\to\infty}\eta_T(B_\varepsilon).
\]
Then
\[
 -\log d_-(\varepsilon)
 \sim-\log d_+(\varepsilon)\sim J(\varepsilon).
\]
If $A_\varepsilon\cap B_\varepsilon$ are also Haar-continuity
sets and $\mu(A_\varepsilon\mid B_\varepsilon)\to0$, then
\[
 \lim_{\varepsilon\downarrow0}\limsup_{T\to\infty}
 \frac{\eta_T(A_\varepsilon\cap B_\varepsilon)}
      {\eta_T(B_\varepsilon)}=0.
\]
\end{corollary}
\begin{proof}
Compactness of $\Omega$ gives subsequential weak limits of
$\eta_T$.  Their density bounds make every Haar-continuity set
a continuity set for each such limit.  Hence
$c_{\rm MT}\mu(B_\varepsilon)\le d_-(\varepsilon)
\le d_+(\varepsilon)\le\mu(B_\varepsilon)$, proving the first
assertion.  Along any weakly convergent subsequence the conditional
ratio tends to $\nu(A_\varepsilon\cap B_\varepsilon)/
\nu(B_\varepsilon)$, whose denominator is positive.  Apply
Theorem~\ref{p04:thm:bounded-transfer}(ii) and then let
$\varepsilon\downarrow0$.
\end{proof}

All these statements take $T\to\infty$ before
$\varepsilon\downarrow0$.  They do not supply estimates for a
prescribed shrinking function $\varepsilon(T)$.

Finally, the Euler-product map $\omega\mapsto Z_\omega$ is a
continuous injection from the compact space $\Omega$ into the
space of holomorphic functions on $\Re s>1$ with the compact-open
topology.  Injectivity follows by uniqueness of the absolutely
convergent Dirichlet coefficients.  It is consequently a
homeomorphism onto its compact image.  Pushing the measure
inequalities forward gives the same bounded-density transfer for
the random Euler-product law and for all compact-derivative
observables.  Robust root configurations are included by
Rouch\'e's theorem.  The target class remains the unimodular
Euler-product family in the half-plane of absolute convergence.




\section{Exact phase entropy and a conditional law for the two Euler edges}
\label{p04:sec:exact-joint-edges}

Write $V(\theta)=2(1-\cos\theta)$ and introduce the real analytic function
\[
 h(u)=-\log\left(\frac1{2\pi}\int_{-\pi}^{\pi}
                  e^{-uV(\theta)}\,d\theta\right)
     =2u-\log I_0(2u),\qquad u\in\mathbb R.
\]
For $\sigma>1$ set
\begin{equation}\label{p04:eq:entropy-Csigma}
 H_\sigma=\int_0^\infty h(y^{-\sigma})\,dy,
 \qquad C_\sigma=\frac{H_\sigma}{\sigma}
 =\int_0^\infty y^{-\sigma}h'(y^{-\sigma})\,dy.
\end{equation}
The equality follows by integration by parts. These constants are
finite and strictly positive. Indeed, $h(u)=2u+O(u^2)$ as $u\to0$,
whereas $h(u)=\frac12\log(4\pi u)+O(u^{-1})$ as $u\to+\infty$.
The method below is the classical exponential-tilting method for
small deviations of weighted independent variables; the one-variable
asymptotic is also a specialization of Aurzada~\cite{Aurzada},
using $p_n\sim n\log n$. We include the proof
to fix the integral constant and the tilt used in the joint-edge
results; no new general small-deviation theorem is claimed.

\begin{theorem}[Prime-weight small-ball specialization]
\label{p04:thm:exact-phase-entropy}
Let $a>0$, $\sigma>1$, $\kappa\in\mathbb R$, and let independent
uniform angles $\theta_p$ be indexed by the primes. Suppose $w_p>0$ and
$w_p\sim a p^{-\sigma}(\log p)^\kappa$ and put
$D_w=\sum_p w_pV(\theta_p)$. With $\beta=(\sigma-1)^{-1}$,
\begin{equation}\label{p04:eq:exact-general-entropy}
 -\log\mu(D_w\le\varepsilon)
 \sim K_{\sigma,\kappa,a}\,
 \varepsilon^{-\beta}
 [\log(1/\varepsilon)]^{(\kappa-1)\beta-1},
\end{equation}
where
\[
 K_{\sigma,\kappa,a}
   =(\sigma-1)^2 C_\sigma
     (aC_\sigma)^\beta\beta^{(\kappa-1)\beta}.
\]
Changing or deleting finitely many positive weights does not change
this constant.
\end{theorem}

\begin{proof}
Let $F(s)=-\log\mathbb E e^{-sD_w}=\sum_p h(sw_p)$, and choose
$x=x(s)\to\infty$ by $sa x^{-\sigma}(\log x)^\kappa=1$.
The prime number theorem and partial summation give
\begin{equation}\label{p04:eq:saddle-prime-sums}
 F(s)\sim H_\sigma\frac{x}{\log x},\quad
 sF'(s)\sim C_\sigma\frac{x}{\log x},\quad
 -s^2F''(s)\asymp\frac{x}{\log x}.
\end{equation}
Here is a justification of the improper-integral passage. On
$\eta x\le p\le Mx$, the ratio $sw_p$ converges uniformly to
$(p/x)^{-\sigma}$, and the prime number theorem applies to continuous
functions on $[\eta,M]$. For $p>Mx$ use $h(u)\le2u$ and partial
summation; the resulting normalized tail is
$O(M^{1-\sigma+\xi})$ for any fixed $0<\xi<\sigma-1$.
For $p<\eta x$, use $h(u)\ll1+\log(1+u)$; primes below $x^c$,
with fixed $c<1$, contribute $o(x/\log x)$, and the remaining
normalized contribution is $O(\eta^{1-\xi}(1+|\log\eta|))$.
The same argument applies to $uh'(u)$ and $-u^2h''(u)$, which are
respectively $O(\min(u,1))$ and $O(\min(u^2,1))$ for $u>0$.
The second-derivative integral is positive. Finite initial terms
are $O(\log s)$ in the first sum and bounded in the other two.

Under the tilted probability $d\mu_s=e^{-sD_w+F(s)}d\mu$, the
mean and variance of $D_w$ are $F'(s)$ and $-F''(s)$.
For the upper probability bound choose $F'(s)=\varepsilon$:
\[
 \mu(D_w\le\varepsilon)\le e^{s\varepsilon-F(s)}.
\]
For the matching logarithmic lower bound choose a nearby tilt with
mean $(1-\eta_s)\varepsilon$, where $\eta_s\to0$ and
$\eta_s\sqrt{x/\log x}\to\infty$. Chebyshev's inequality gives
\[
 \mu_s((1-2\eta_s)\varepsilon\le D_w\le\varepsilon)=1-o(1).
\]
Changing measure on this interval loses only
$2\eta_s s\varepsilon=o(x/\log x)$ in the logarithm. Thus
\[
 -\log\mu(D_w\le\varepsilon)
 \sim (\sigma-1)C_\sigma\frac{x}{\log x},\qquad
 \varepsilon\sim aC_\sigma
 x^{1-\sigma}(\log x)^{\kappa-1}.
\]
Solving the last relation gives
$x\sim(aC_\sigma)^\beta\beta^{(\kappa-1)\beta}
\varepsilon^{-\beta}[\log(1/\varepsilon)]^{(\kappa-1)\beta}$,
which proves the stated constant.
\end{proof}

We now use $E=E_*$, $A=A_*$, $\beta=(E-1)^{-1}$,
$\tau=(A-1)/(E-1)$, and
$\Psi(\delta)=\delta^\tau/\log(1/\delta)$.
The frozen phases are $z_2^*=-1$ and $z_p^*=1$ for odd primes.
Write $z_p=z_p^*e^{i\theta_p}$ and retain the definitions
\[
 D_E=\sum_p(\log p)p^{-E}V(\theta_p),\qquad
 D_A=\sum_pp^{-A}V(\theta_p),\qquad B_\delta=\{D_E\le\delta\}.
\]
Define two explicit constants
\begin{equation}\label{p04:eq:joint-edge-constants}
 \mathcal K_E=(E-1)^2 C_E^{E/(E-1)},\qquad
 J_{E,A}=\int_0^\infty y^{-A}h'(y^{-E})\,dy,\qquad
 \mathcal J_{E,A}=(E-1)C_E^{-\tau}J_{E,A}.
\end{equation}
The integral defining $J_{E,A}$ converges since $1<A<E$.

\begin{theorem}[Conditional phase-energy profile]
\label{p04:thm:conditional-energy-profile}
As $\delta\downarrow0$,
\begin{equation}\label{p04:eq:exact-E-entropy}
 -\log\mu(B_\delta)
 \sim\mathcal K_E\frac{\delta^{-\beta}}{\log(1/\delta)}.
\end{equation}
Under the conditional law $\mu(\,\cdot\mid B_\delta)$,
\begin{equation}\label{p04:eq:conditional-energy-profile}
 \frac{D_E}{\delta}\longrightarrow1,\qquad
 \frac{D_A}{\Psi(\delta)}\longrightarrow\mathcal J_{E,A}
 \quad\hbox{in probability}.
\end{equation}
For each fixed finite set of primes $F$,
\begin{equation}\label{p04:eq:fixed-primes-negligible}
 \frac1\delta\sum_{p\in F}(\log p)p^{-E}V(\theta_p)
 \longrightarrow0\quad\hbox{in conditional probability}.
\end{equation}
The same conclusions hold, uniformly, for every finite measure
$\nu$ satisfying $c\mu\le\nu\le\mu$, with any fixed $c>0$.
\end{theorem}

\begin{proof}
The entropy formula is Theorem~\ref{p04:thm:exact-phase-entropy} with
$(\sigma,\kappa,a)=(E,1,1)$. It also proves the first assertion
in \eqref{p04:eq:conditional-energy-profile}: for each $\eta>0$, the
ratio $\mu(B_{(1-\eta)\delta})/\mu(B_\delta)$ tends to zero
exponentially on the scale $\delta^{-\beta}/\log(1/\delta)$.

Choose a tilted law with $\mathbb E_s D_E=(1-\eta_s)\delta$,
$\eta_s\to0$, $\eta_s\sqrt m\to\infty$, where
\[
 s=x^E/\log x,\qquad m=x/\log x,\qquad
 \delta\sim C_Ex^{1-E}.
\]
Let $b_x=x^{1-A}/\log x$. For each real $t$ in a fixed
neighborhood of zero, prime-number-theorem summation gives
\begin{equation}\label{p04:eq:cross-energy-log-mgf}
 \frac1m\log\mathbb E_s\exp(t x^A D_A)
 \longrightarrow
 \Lambda(t):=\int_0^\infty
 \{h(y^{-E})-h(y^{-E}-ty^{-A})\}\,dy.
\end{equation}
This integral is well defined also for $t>0$. At zero in $y$,
$y^{-E}$ dominates $ty^{-A}$, while at infinity both arguments
are small and the integrand is $O(y^{-A})$. Its first two
derivatives with respect to $t$ are locally integrable uniformly
near $t=0$: their behaviors near zero are bounded by constants
times $y^{E-A}$ and $y^{2E-2A}$, and at infinity by $y^{-A}$
and $y^{-2A}$. The same bounds justify the prime-sum limits.
Consequently $\Lambda'(0)=J_{E,A}$, and exponential Markov
inequalities for positive and negative $t$ imply, for every
$\eta>0$ and some $c_\eta>0$,
\[
 \mu_s\bigl(|D_A/b_x-J_{E,A}|>\eta\bigr)
 \le e^{-c_\eta m}
\]
for all sufficiently large $x$.

The change of measure used in the preceding proof implies that
any event with $\mu_s$-probability at most $e^{-cm}$ has
$\mu(\,\cdot\mid B_\delta)$-probability tending to zero: its
conditional probability is at most $e^{o(m)}\mu_s(\text{event})$.
Since
\[
 b_x\sim(E-1)C_E^{-\tau}\Psi(\delta),
\]
this proves the second energy limit. For a fixed finite $F$ put
$D_F=\sum_{p\in F}(\log p)p^{-E}V(\theta_p)$.
For fixed $0<q<1$,
\[
 \log\mathbb E_s e^{qsD_F}
 =\sum_{p\in F}\{h(sw_p)-h((1-q)sw_p)\}=O_F(1).
\]
Thus $\mu_s(D_F>\eta\delta)\le\exp(-q\eta s\delta+O_F(1))$,
and the same change of measure proves
\eqref{p04:eq:fixed-primes-negligible}.

Finally, for every measurable event $C$,
\[
 \nu(C\mid B_\delta)
 \le c^{-1}\mu(C\mid B_\delta).
\]
This transfers every vanishing exceptional probability, uniformly
over the specified class of measures. It transfers concentration,
not arbitrary nondegenerate conditional distribution laws.
\end{proof}

\begin{theorem}[The phase profile of a rare Euler-edge configuration]
\label{p04:thm:conditional-prime-profile}
Put $x_\delta=(C_E/\delta)^{1/(E-1)}$. For $y>0$ define the
probability measure
\[
 dq_y(\theta)=\exp\{-y^{-E}V(\theta)+h(y^{-E})\}
                   \frac{d\theta}{2\pi}.
\]
For every continuous $g$ with compact support in $(0,\infty)\times
\mathbb T$, under $\mu(\,\cdot\mid B_\delta)$ one has
\begin{equation}\label{p04:eq:conditional-prime-profile}
 \frac{\log x_\delta}{x_\delta}
       \sum_p g(p/x_\delta,\theta_p)
 \longrightarrow\int_0^\infty\int_{\mathbb T}
                 g(y,\theta)\,dq_y(\theta)\,dy
\end{equation}
in probability. Any finite family of such observables converges
jointly. Both statements transfer uniformly to $c\mu\le\nu\le\mu$.

Moreover, $q_y$ is the unique entropy-minimizing measurable family
of probability measures $\rho_y$ on the circle subject to
\[
 \int_0^\infty y^{-E}\int V\,d\rho_y\,dy\le C_E.
\]
The minimum of $\int_0^\infty D(\rho_y\Vert d\theta/(2\pi))\,dy$
is $(E-1)C_E$, with uniqueness up to a set of $y$ of measure zero.
\end{theorem}

\begin{proof}
Use the undershooting tilt from the preceding proof. On the
compact interval of $y$ supporting $g$, its angular marginal at
the prime $p$ converges uniformly to $q_{p/x}$. The prime number
theorem gives the limit of the expected sum. There are $O(m)$
relevant independent terms, each bounded, so the exponential
Markov bound (or the elementary bounded-variable exponential
moment estimate) bounds any fixed deviation of their normalized
sum by $e^{-c_g m}$. The previous change-of-measure argument
transfers this to the conditioned law. Since $x/x_\delta\to1$,
uniform continuity of $g$ replaces $x$ by $x_\delta$. A finite
union bound proves joint convergence.

For the variational statement, the relative-entropy identity is
\[
 \int_0^\infty D(\rho_y\Vert d\theta/(2\pi))\,dy
 =\int_0^\infty D(\rho_y\Vert q_y)\,dy
   +H_E-\int_0^\infty y^{-E}\int V\,d\rho_y\,dy.
\]
It holds first when the integrals are finite, by substitution of
the displayed density, and then with the usual extended-integral
interpretation since the energy and $H_E$ are finite.
The constraint and nonnegativity of relative entropy give the
lower bound $H_E-C_E=(E-1)C_E$. The family $q_y$ attains it
by \eqref{p04:eq:entropy-Csigma}; equality forces the constraint to
be saturated and $D(\rho_y\Vert q_y)=0$ for almost every $y$.
\end{proof}

\begin{theorem}[Typical shape of the two edge companions]
\label{p04:thm:typical-edge-shape}
Let $u_E(z)$ and $u_A(z)$ be the local simple roots defined by
$Q_z(E+u_E)=0$ and $Z_z(A+u_A)=1$. Put
$a_E=Q_{z^*}'(E)<0$ and $a_A=Z_{z^*}'(A)<0$.
Under $\mu(\,\cdot\mid B_\delta)$, and uniformly under every
$\nu(\,\cdot\mid B_\delta)$ with $c\mu\le\nu\le\mu$,
\begin{equation}\label{p04:eq:typical-edge-shape}
 \frac{-\Re u_E}{\delta}\longrightarrow\frac1{2|a_E|},
 \qquad
 \frac{-\Re u_A}{\Psi(\delta)}
 \longrightarrow\frac{\mathcal J_{E,A}}{2|a_A|},
 \qquad
 \frac{|\Im u_E|+|\Im u_A|}{\sqrt\delta}\longrightarrow0
\end{equation}
in probability.
\end{theorem}

\begin{proof}
We give the additional estimates that replace the earlier
worst-case inequalities by exact conditional constants. Fix a
small disk about $E$, of radius $r$ with $E-2r>1$. For its
$C^1$ norm, Cauchy--Schwarz and a split at a fixed prime bound $P$
give
\[
 \|Q_z-Q_{z^*}\|_{C^1}^2
 \le C_P D_{E,\le P}+\epsilon_P D_E,
 \qquad \epsilon_P\longrightarrow0.
\]
For the derivative, for example, the tail constant is controlled by
\[
 \sum_{p>P}(\log p)^3p^{-E+2r}.
\]
Theorem~\ref{p04:thm:conditional-energy-profile}, first with $\delta$
tending to zero and then with $P$ tending to infinity, yields
\begin{equation}\label{p04:eq:conditional-C1-small}
 \|Q_z-Q_{z^*}\|_{C^1}^2=o_{\rm prob}(\delta).
\end{equation}
For odd primes, expansion into prime powers gives
\[
 \Re\left\{\frac{(\log p)z_pp^{-E}}{1-z_pp^{-E}}
       -\frac{(\log p)p^{-E}}{1-p^{-E}}\right\}
 =-\tfrac12(\log p)p^{-E}V(\theta_p)
       +O((\log p)p^{-2E}V(\theta_p)).
\]
Here $1-\cos(k\theta)\le k^2(1-\cos\theta)$ bounds the remainder.
The prime $2$ is included in the finite-prime error. Splitting the
remainder at $P$ and using \eqref{p04:eq:fixed-primes-negligible} shows
\[
 \Re Q_z(E)=-D_E/2+o_{\rm prob}(\delta).
\]
The quantitative implicit-function expansion
$u_E=-Q_z(E)/a_E+O(\|Q_z-Q_{z^*}\|_{C^1}^2)$ proves the
first limit and the improved imaginary bound in
\eqref{p04:eq:typical-edge-shape}.

For the other edge, define the normally convergent logarithm
$L_z=\log Z_z(A)-\log Z_{z^*}(A)$, where $Z_{z^*}(A)=1$.
The same prime-power calculation gives
\[
 \Re L_z=-D_A/2+O(D_{A,\le P}+P^{-A}D_A).
\]
For fixed $P$, $D_{A,\le P}\ll_P D_E\le\delta=o(\Psi(\delta))$.
Choose a fixed disk radius $r>0$ with $2A-E-2r>1$ and $A-r>1$.
The strict inequality $2A-E>1$ gives $|L_z|^2\ll D_E$, and
also a $C^1$ bound $\|Z_z-Z_{z^*}\|_{C^1}^2\ll D_E$ on this
disk. In fact the same fixed-prime/tail splitting improves this
to $o_{\rm prob}(\delta)$: the tail coefficient for the derivative
is bounded by a constant times
$\sum_{p>P}(\log p)p^{E-2A+2r}$, which tends to zero.
It follows also that $u_A=o_{\rm prob}(\sqrt\delta)$. Furthermore,
\[
 \Re(Z_z(A)-1)=-D_A/2+o_{\rm prob}(\Psi(\delta)),\qquad
 u_A=-\frac{Z_z(A)-1}{a_A}+O(\delta).
\]
The second energy limit proves the asserted constant. Uniform
transfer to $\nu$ follows from the conditional measure inequality.
\end{proof}

\begin{corollary}[Exact logarithmic rarity along simple critical zeros]
\label{p04:cor:exact-simple-rarity}
The fixed-prime phase-domination theorem proved above implies that every
weak limit $\nu$ of the simple-critical-zero phase measures obeys
$c_{\rm MT}\mu\le\nu\le\mu$. Let $\underline d(\delta)$ and
$\overline d(\delta)$ be the lower and upper densities, relative
to all zeros, of simple critical-line ordinates whose phase lies in
$B_\delta$. Then both satisfy
\[
 -\log\underline d(\delta)\sim
 -\log\overline d(\delta)\sim
 \mathcal K_E\frac{\delta^{-\beta}}{\log(1/\delta)}.
\]
Moreover, for every fixed tolerance $\eta>0$, the upper relative
proportion of exceptions to the two real-part limits in
\eqref{p04:eq:typical-edge-shape}, among those selected simple zeros,
tends to zero as $\delta\downarrow0$, with the height limit taken
first. The analogous assertion holds for
$|\Im u_E|+|\Im u_A|>\eta\sqrt\delta$.
\end{corollary}

\begin{proof}
The density sandwich
$c_{\rm MT}\mu(B_\delta)\le\underline d(\delta)
\le\overline d(\delta)\le\mu(B_\delta)$ gives the logarithmic
asymptotics. For the relative assertions use compactness of the
phase measures and their weak-limit domination. Equivalently,
bound the closed exceptional set with deviation at least $\eta$
inside $B_\delta$ by its Haar measure and divide by
$c_{\rm MT}\mu(B_\delta)$. The local root maps are continuous
on this compact set, so the upper bound follows from Portmanteau.
The conditional concentration theorem makes this bound tend to
zero. No rate connecting $\delta$ to the height is asserted.
\end{proof}




\section{A joint lower-tail rate and three active-constraint regimes}
\label{p04:sec:joint-rate}

The two energies need not be tied to the typical relation found above.
One can instead prescribe both tolerances and determine their joint
logarithmic cost. The inequality $E<A+1$, which follows already from
$E<2.82$ and $A>1.93$, is useful here.

\begin{theorem}[Joint energy rate function]
\label{p04:thm:joint-energy-rate}
Suppose $1<A<E<A+1$. Under product Haar measure let
\[
 D_E=\sum_p(\log p)p^{-E}V(\theta_p),\qquad
 D_A=\sum_pp^{-A}V(\theta_p),
\]
and put
\[
 U_x=x^{E-1}D_E,\qquad
 W_x=x^{A-1}(\log x)D_A,\qquad m_x=x/\log x.
\]
For every $e,a>0$,
\begin{equation}\label{p04:eq:joint-energy-rate}
 -\frac1{m_x}\log\mu(U_x\le e,\ W_x\le a)
 \longrightarrow \mathcal I(e,a)
 :=\max_{s,t\ge0}\{\mathcal H(s,t)-se-ta\},
\end{equation}
where
\[
 \mathcal H(s,t)=\int_0^\infty h(sy^{-E}+ty^{-A})\,dy.
\]
The maximum has a unique nonzero optimizer $(s_*,t_*)$. Write
$J_{E,A}=\int_0^\infty y^{-A}h'(y^{-E})dy$,
$J_{A,E}=\int_0^\infty y^{-E}h'(y^{-A})dy$, and
$\tau=(A-1)/(E-1)$. Its three regimes are as follows:
\begin{enumerate}
\item If $a\ge J_{E,A}(e/C_E)^\tau$, then
\[
 t_*=0,\quad s_*=(C_E/e)^{E/(E-1)},\quad
 \mathcal I(e,a)=(E-1)C_E^{E/(E-1)}e^{-1/(E-1)}.
\]
\item If $e\ge J_{A,E}(a/C_A)^{1/\tau}$, then
\[
 s_*=0,\quad t_*=(C_A/a)^{A/(A-1)},\quad
 \mathcal I(e,a)=(A-1)C_A^{A/(A-1)}a^{-1/(A-1)}.
\]
\item Otherwise $s_*,t_*>0$ and they are uniquely determined by
\[
 \partial_s\mathcal H(s_*,t_*)=e,\qquad
 \partial_t\mathcal H(s_*,t_*)=a.
\]
\end{enumerate}
The first two parameter regions are disjoint. In the third region
the joint rate is strictly larger than each separate one-constraint
rate. The same logarithmic limit holds with $\mu$ replaced by any
finite measure $\nu$ satisfying $c\mu\le\nu\le\mu$.
\end{theorem}

\begin{proof}
For fixed $s,t\ge0$, the prime number theorem gives
\begin{equation}\label{p04:eq:joint-laplace-limit}
 -\frac1{m_x}\log\mathbb E
       \exp\{-m_x(sU_x+tW_x)\}
 \longrightarrow\mathcal H(s,t).
\end{equation}
Indeed, the argument of $h$ at the prime $p$ is
$s(p/x)^{-E}(\log p/\log x)+t(p/x)^{-A}$.
The small and large prime tails are controlled exactly as in the
one-variable Laplace calculation. Exponential Markov immediately
gives the lower bound for the limiting negative logarithm in
\eqref{p04:eq:joint-energy-rate}.

For existence and uniqueness of the saddle, note that $h$ is strictly
concave on $\mathbb R$, and $h(u+v)\le h(u)+h(v)$ for $u,v\ge0$.
Thus $\mathcal H$ is strictly concave away from the origin and
\[
 \mathcal H(s,t)\le H_Es^{1/E}+H_At^{1/A}.
\]
The objective is consequently coercive on the quadrant. It increases
from zero along either axis sufficiently near the origin, so its
unique maximum is nonzero. The gradient exists at every nonzero
point of the quadrant, including the axes: at $s=0,t>0$, the
only potentially problematic integrand is
$y^{-E}h'(ty^{-A})=O(y^{A-E})$ near zero, which is integrable
by $E<A+1$. The other cases are easier.

To prove the probability lower bound, tilt the independent prime
phases by the parameter $((1+\eta)s_*,(1+\eta)t_*)$, with
$\eta>0$. Under this law the means of $U_x,W_x$ tend to
$\nabla\mathcal H((1+\eta)s_*,(1+\eta)t_*)$, and both limits
are strictly below $e,a$. Both normalized energies converge in
probability to these means. To see this without assuming a global
variance bound, first restrict to $\epsilon x\le p\le Mx$;
there each summand is $O(1/m_x)$ and the sum has variance
$O(1/m_x)$. The normalized expectations of the omitted tails
are uniformly small as $\epsilon\downarrow0$, $M\uparrow\infty$,
by the integrability of the two gradient integrands. Markov's
inequality controls these nonnegative tails.

On a small rectangle about the limiting tilted means, contained in
$\{U_x<e,W_x<a\}$, the change-of-measure formula now supplies
the matching logarithmic lower probability bound. Let the rectangle
width tend to zero and then let $\eta\downarrow0$.
The complementary-slackness identities
\[
 s_*\partial_s\mathcal H(s_*,t_*)=s_*e,\qquad
 t_*\partial_t\mathcal H(s_*,t_*)=t_*a
\]
show that the resulting exponent is exactly
$\mathcal H(s_*,t_*)-s_*e-t_*a$.

On the $t=0$ axis, differentiation of
$\mathcal H(s,0)=H_Es^{1/E}$ gives
$s_*=(C_E/e)^{E/(E-1)}$, and the unused directional derivative
is $J_{E,A}s_*^{(1-A)/E}-a$. Its nonpositivity is exactly the
first displayed regime. The same calculation on the other axis
gives the second regime. Strict concavity excludes overlap of the
two regimes and implies that in the remaining case the maximizer
lies in the interior. There the optimizer differs from either
axis optimizer, proving the strict increase of the joint cost.
Finally, comparison by the fixed factors $c$ and $1$ changes a
logarithm by at most $|\log c|=o(m_x)$.
\end{proof}

\begin{corollary}[Jointly constrained phase profile]
\label{p04:cor:joint-constrained-profile}
Condition on $U_x\le e,W_x\le a$, and let $(s_*,t_*)$ be the
optimizer above. For every continuous compactly supported
$g:(0,\infty)\times\mathbb T\to\mathbb R$,
\[
 \frac1{m_x}\sum_p g(p/x,\theta_p)
 \longrightarrow \int_0^\infty\int g(y,\theta)
 e^{-(s_*y^{-E}+t_*y^{-A})V(\theta)
       +h(s_*y^{-E}+t_*y^{-A})}\frac{d\theta}{2\pi}\,dy
\]
in conditional probability. This also holds uniformly for
$c\mu\le\nu\le\mu$.
\end{corollary}

\begin{proof}
Under the tilt at the exact optimizer, the normalized compact-band
sum has the stated limiting mean and exponentially small fixed
deviations, by independence and boundedness of its $O(m_x)$ terms.
On the conditioning event the Radon--Nikodym derivative returning
to Haar measure is at most
$\exp(m_x(s_*e+t_*a)-F_x(s_*,t_*))$.
The joint probability formula just proved equals this exponential
up to $\exp(o(m_x))$. Thus exponential concentration survives
conditioning. The inequality
$\nu(C\mid B)\le c^{-1}\mu(C\mid B)$ completes the transfer.
\end{proof}

This theorem identifies when demanding additional $1$-point phase
coherence has no leading exponential cost, when it determines the
entire cost, and when it changes the entropy-minimizing prime-phase
profile. It concerns the two phase energies; it does not by itself
give a local central limit theorem for the companion roots.


\begin{corollary}[A single normalized crossover function]
\label{p04:cor:joint-homogeneity}
Put $\beta=1/(E-1)$ and $\Phi(q)=\mathcal I(1,q)$. Then
\[
 \mathcal I(e,a)=e^{-\beta}\Phi(a/e^\tau).
\]
In the normalized variable $q=a/e^\tau$, the three regimes are
separated by two fixed constants
\[
 q_-=C_AJ_{A,E}^{-\tau}<q_+=J_{E,A}C_E^{-\tau}.
\]
The $A$-energy condition alone is active for $q\le q_-$, both
conditions are active for $q_-<q<q_+$, and the $E$-energy
condition alone is active for $q\ge q_+$.
\end{corollary}
\begin{proof}
Changing variable in the integral and then taking the dual gives
\[
 \mathcal H(\lambda^Es,\lambda^At)=\lambda\mathcal H(s,t),
 \qquad
 \mathcal I(\lambda^{1-E}e,\lambda^{1-A}a)
 =\lambda\mathcal I(e,a).
\]
Take $\lambda=e^{1/(E-1)}$ and use the two regime conditions
in Theorem~\ref{p04:thm:joint-energy-rate}. Their disjointness implies
the strict order of the two constants.
\end{proof}




\section{Asymmetric transitions between the two edge constraints}
\label{p04:sec:joint-transition}

The boundary between the second and third regimes in
Theorem~\ref{p04:thm:joint-energy-rate} has an explicit nonquadratic
local law. The relevant exponent depends on the difference of the
two classical zeta boundaries.

\begin{theorem}[Nonquadratic joint-cost transition]
\label{p04:thm:nonquadratic-transition}
Suppose $1<A<E<A+1$ and $d=E-A\in(1/2,1)$. Fix $a>0$, and set
\[
 t_0=(C_A/a)^{A/(A-1)},\qquad
 e_0=J_{A,E}t_0^{(1-E)/A},\qquad
 \mathcal I_A(a)=(A-1)C_A^{A/(A-1)}a^{-1/(A-1)}.
\]
Define
\[
 K_d=\frac{\pi}{2\sin(\pi(1/d-1))},\qquad
 L_d(a)=(1-d)\left(\frac d{K_d}\right)^{d/(1-d)}
                     t_0^{1/(1-d)}.
\]
Then $\mathcal I(e,a)=\mathcal I_A(a)$ for $e\ge e_0$, whereas
\begin{equation}\label{p04:eq:nonquadratic-transition}
 \mathcal I(e_0-\epsilon,a)-\mathcal I_A(a)
 \sim L_d(a)\,\epsilon^{1/(1-d)}
 \qquad(\epsilon\downarrow0).
\end{equation}
The corresponding multiplier satisfies
\begin{equation}\label{p04:eq:transition-saddle}
 s_*\sim t_0^{1/(1-d)}
               \left(\frac{d\epsilon}{K_d}\right)^{d/(1-d)},
 \qquad t_*=t_0+O(s_*).
\end{equation}
For the classical zeta boundaries the onset exponent is
\[
 \frac1{1-(E_*-A_*)}=7.86858\ldots.
\]
Thus the additional leading exponential cost of imposing critical
edge coherence begins with this fractional power, rather than a
quadratic law.
\end{theorem}

\begin{proof}
We first prove the uniform expansion, for $t$ in a compact
subinterval of $(0,\infty)$,
\begin{equation}\label{p04:eq:singular-axis-expansion}
 \mathcal H(s,t)=\mathcal H(0,t)
       +s\partial_s\mathcal H(0,t)
       -K_d(s/t)^{1/d}+o(s^{1/d}),\qquad s\downarrow0.
\end{equation}
Put $v=ty^{-A}$, $w=sy^{-E}$, and
$R(v,w)=h(v+w)-h(v)-wh'(v)$. Since $h$ is concave,
$R\le0$. For $v\ge1$, the bounds $h'(v)\ll v^{-1}$ and
$|h''(v)|\ll v^{-2}$ give
\[
 |R(v,w)|\ll\min(w/v,(w/v)^2).
\]
For $v<1$, boundedness of $h''$ on $[0,\infty)$ gives an
integrated contribution $O(s^2)$, uniformly for the specified $t$.

In the remaining integral substitute $y=(s/t)^{1/d}z$.
Then $w/v=z^{-d}$ and $v\to\infty$ for each fixed $z>0$.
Using $h(v)=\frac12\log(4\pi v)+O(v^{-1})$ and
$h'(v)=1/(2v)+O(v^{-2})$, its rescaled integrand tends to
\[
 \frac12\{\log(1+z^{-d})-z^{-d}\}.
\]
It is dominated by a constant times
$\min(z^{-d},z^{-2d})$, which is integrable precisely under
the restrictions $d<1$ and $2d>1$. Dominated convergence is
uniform for $t$ in compact positive intervals. Integration by
parts followed by $u=z^d$ yields
\[
 \begin{split}
 \int_0^\infty\{\log(1+z^{-d})-z^{-d}\}\,dz
 &=-d\int_0^\infty\frac{z^{-d}}{1+z^d}\,dz\\
 &=-\int_0^\infty\frac{u^{1/d-2}}{1+u}\,du
 =-\frac\pi{\sin(\pi(1/d-1))}.
 \end{split}
\]
Because $s^2=o(s^{1/d})$, this proves
\eqref{p04:eq:singular-axis-expansion} with the asserted constant.

Next maximize over $t$. At $(0,t_0)$ the derivatives
$\mathcal H_{tt}$ and $\mathcal H_{st}$ are finite, with
$\mathcal H_{tt}<0$. For example the mixed-derivative integrand
near $y=0$ is $O(y^{A-E})=O(y^{-d})$, which is integrable.
These two derivatives are continuous as $s\downarrow0$, by the
bound $|h''(u)|\ll(1+u)^{-2}$. It follows by a one-sided
monotonicity argument that the maximizing solution of
$\mathcal H_t(s,t_s)=a$ has $t_s=t_0+O(s)$.
Indeed, $\mathcal H_t(s,t_0)-a=O(s)$ and
$\mathcal H_{tt}$ stays bounded away from zero in a small
one-sided neighborhood. No extension of $\mathcal H$ to negative
$s$ is used.

The change in the objective caused by replacing $t_s$ by $t_0$
is $O(s^2)$: the first variation of
$\mathcal H(0,t)-at$ vanishes at $t_0$, and the mixed linear
term changes by $O(s|t_s-t_0|)$. Hence, with $r=1/d$,
\begin{equation}\label{p04:eq:reduced-axis-expansion}
 \max_{t\ge0}\{\mathcal H(s,t)-at\}
 =\mathcal I_A(a)+e_0s-Bs^r+o(s^r),
 \qquad B=K_dt_0^{-r}.
\end{equation}
Uniqueness and coercivity of the joint optimizer imply that as
$e\uparrow e_0$ from below the optimizer tends to $(0,t_0)$.
The remaining maximization is therefore locally that of
$\epsilon s-Bs^r+o(s^r)$. To avoid differentiating a remainder,
bound it between $\epsilon s-(B+\eta)s^r$ and
$\epsilon s-(B-\eta)s^r$, maximize these explicit functions,
and then let $\eta\downarrow0$. Their common limiting value is
\[
 (r-1)r^{-r/(r-1)}B^{-1/(r-1)}
              \epsilon^{r/(r-1)}.
\]
Since $r/(r-1)=1/(1-d)$, this coefficient equals $L_d(a)$.
The same rescaling, and uniqueness of the maximizer of
$u-Bu^r$, gives
$s_*\sim(\epsilon/(rB))^{1/(r-1)}$, which is
\eqref{p04:eq:transition-saddle}. The constant branch $e\ge e_0$
is the second regime of Theorem~\ref{p04:thm:joint-energy-rate}.
\end{proof}

The source of the fractional onset is a divergent second derivative
$\mathcal H_{ss}(0,t)$: its integrand behaves like $y^{-2d}$
near zero. The first derivative stays finite because $d<1$.
The resulting nonquadratic law is compatible with the finite
mixed derivative needed to optimize $t$; replacing the singular
second derivative by a quadratic approximation would give an
incorrect transition exponent.

\begin{theorem}[Quadratic onset at the critical-edge-dominated boundary]
\label{p04:thm:quadratic-transition}
Suppose $1<A<E$. Fix $e>0$ and put
\[
 s_0=(C_E/e)^{E/(E-1)},\qquad
 a_0=J_{E,A}s_0^{(1-A)/E},\qquad
 \mathcal I_E(e)=(E-1)C_E^{E/(E-1)}e^{-1/(E-1)}.
\]
Define the positive Gram matrix
\[
 M_{bc}(s_0)=\int_0^\infty
           y^{-b-c}[-h''(s_0y^{-E})]\,dy,
             \qquad b,c\in\{E,A\},
\]
and its positive Schur complement
\[
 \kappa_e=M_{AA}(s_0)-\frac{M_{EA}(s_0)^2}{M_{EE}(s_0)}>0.
\]
Then $\mathcal I(e,a)=\mathcal I_E(e)$ for $a\ge a_0$, and
\begin{equation}\label{p04:eq:quadratic-transition}
 \mathcal I(e,a_0-\epsilon)-\mathcal I_E(e)
 \sim\frac{\epsilon^2}{2\kappa_e},\qquad
 t_*\sim\frac\epsilon{\kappa_e}\qquad(\epsilon\downarrow0).
\end{equation}
In particular, the two transition boundaries have different leading
orders for the zeta-edge parameters: quadratic at this boundary,
and order $7.86858\ldots$ at the boundary of
Theorem~\ref{p04:thm:nonquadratic-transition}.

When $2A-E>1$, as holds for the classical zeta boundaries, the
curvature constant can also be written
\begin{equation}\label{p04:eq:explicit-transition-curvature}
 \kappa_e=\frac{s_0^{(1-2A)/E}}E
 \left\{(2A-E-1)J_{E,\,2A-E}
       -\frac{(A-1)^2J_{E,A}^2}{(E-1)C_E}\right\},
\end{equation}
where $J_{E,b}=\int_0^\infty y^{-b}h'(y^{-E})\,dy$.
\end{theorem}

\begin{proof}
Every entry of $M$ is finite. Near zero its integrand is
$O(y^{2E-b-c})$, and at infinity it is $O(y^{-b-c})$.
The exponents are integrable for $b,c\in\{E,A\}$ and
$1<A<E$. Since $-h''(u)$ is the strictly positive variance of
$V$ under a nondegenerate angular tilt, the matrix $M$ is the
Gram matrix of the linearly independent functions $y^{-E}$ and
$y^{-A}$ in a positively weighted $L^2$ space. Thus it is
positive definite and $\kappa_e>0$.

All second derivatives of $\mathcal H$ are continuous in a
one-sided neighborhood of $(s_0,0)$. Maximizing first in $s$
gives a unique solution $s_t=s_0+O(t)$ of
$\mathcal H_s(s_t,t)=e$, and ordinary second-order expansion gives
\[
 \max_{s\ge0}\{\mathcal H(s,t)-se\}
 =\mathcal I_E(e)+a_0t-\tfrac12\kappa_e t^2+o(t^2).
\]
Indeed its second derivative is
$\mathcal H_{tt}-\mathcal H_{ts}^2/\mathcal H_{ss}=-\kappa_e$.
As in the preceding proof, uniqueness and coercivity localize the
joint maximizer near the axis. Maximization of
$\epsilon t-\kappa_e t^2/2+o(t^2)$, using upper and lower
quadratic bounds, proves \eqref{p04:eq:quadratic-transition}.

Finally, scaling the integrals and differentiating the axis formulas
give
\[
 M_{EE}(s)=\frac{(E-1)C_E}{E}s^{(1-2E)/E},\qquad
 M_{EA}(s)=\frac{(A-1)J_{E,A}}E s^{(1-E-A)/E}.
\]
If $2A-E>1$, integration by parts also gives
\[
 M_{AA}(s)=\frac{2A-E-1}E
      J_{E,\,2A-E}s^{(1-2A)/E}.
\]
Substitution yields \eqref{p04:eq:explicit-transition-curvature}.
\end{proof}




\section{Zero-free half-planes and the first height of an extreme point}
\label{p03:sec:zero-free-transfer}

This section concerns the ordinary prime flow and phase patterns that
agree with a Dirichlet character outside finitely many primes. It does
not assert a first-height obstruction for arbitrary Euler-product twists.
For a primitive character $\chi$ of conductor $q$, including the
character of conductor $1$, and $1/2\le\vartheta<1$, write
\begin{equation}\label{p03:eq:H-theta}
 \mathrm H_\vartheta(\chi):\qquad
 L(\rho,\chi)=0,\quad 0<\Re\rho<1
 \quad\Longrightarrow\quad \Re\rho\le\vartheta.
\end{equation}
The numerical specialization $\vartheta=7/8$ invokes
\cite[Theorem~1.1]{OpenAIQuasiRH}, which asserts this hypothesis for
every Dirichlet character. The transfer proved below remains valid
under its explicit hypothesis independently of that input.

\begin{lemma}[A smooth prime block]\label{p03:lem:smooth-prime-block}
Fix a nonnegative nonzero $\phi\in C_c^\infty((1,2))$ and put
\[
 \widehat\phi(s)=\int_0^\infty\phi(u)u^{s-1}\,du,\qquad
 I_\phi=\widehat\phi(1)>0.
\]
Under $\mathrm H_\vartheta(\chi)$, uniformly for $X\ge2$ and real $t$,
\begin{equation}\label{p03:eq:smooth-prime-block}
 S_\chi(X,t):=\sum_p(\log p)\phi(p/X)\chi(p)p^{-it}
 =\mathbf1_{\chi=1}X^{1-it}\widehat\phi(1-it)
 +O_\phi\!\left(X^\vartheta\log(q(|t|+3))\right).
\end{equation}
Here $\chi=1$ denotes the character of conductor $1$.
\end{lemma}
\begin{proof}
Mellin inversion gives
\[
 \sum_n\Lambda(n)\chi(n)n^{-it}\phi(n/X)
 =\frac1{2\pi i}\int_{(2)}
       -\frac{L'}L(s+it,\chi)X^s\widehat\phi(s)\,ds.
\]
Moving the contour to $\Re s=-1/2$ gives the smooth explicit formula
\begin{equation}\label{p03:eq:smooth-zero-formula}
 \mathbf1_{\chi=1}X^{1-it}\widehat\phi(1-it)
 -\sum_\rho X^{\rho-it}\widehat\phi(\rho-it)
 +O_\phi\!\left(1+X^{-1/2}\log(q(|t|+3))\right).
\end{equation}
The sum runs over the nontrivial zeros with multiplicity. A possible
trivial zero at $0$ contributes the bounded term. On the new line the
functional equation gives
$L'/L(-1/2+iu,\chi)=O(\log(q(|u|+3)))$; rapid vertical decay of
$\widehat\phi$ bounds the remaining integral. Horizontal contours
avoiding zeros justify the shift; see the classical explicit formula
and zero counts in \cite[Chapters~10 and~12]{MontgomeryVaughan}.

For every fixed $N$, uniformly for $0\le\beta\le1$,
$|\widehat\phi(\beta+iv)|\le C_{\phi,N}(1+|v|)^{-N}$.
Together with the unit-interval zero count
$\#\{\rho:u\le\Im\rho\le u+1\}\ll\log(q(|u|+3))$, this makes the
zero sum absolutely convergent and gives
\[
 \sum_\rho|X^{\rho-it}\widehat\phi(\rho-it)|
 \ll_\phi X^\vartheta\log(q(|t|+3)).
\]
The prime-power terms of exponent at least two are
\[
 \ll\sum_{2\le k\le\log(2X)/\log2}
       \sum_{p\le(2X)^{1/k}}\log p
 \ll X^{1/2}+X^{1/3}\log X\ll X^{1/2}.
\]
They are absorbed since $\vartheta\ge1/2$. This proves
\eqref{p03:eq:smooth-prime-block}, without a factor $|t|$ or an extra
$\log X$.
\end{proof}

Fix a finite set $S$ and unit complex numbers $b_p$, for $p\nmid q$,
satisfying $b_p=\chi(p)$ outside $S$. For
$w_p\asymp p^{-\sigma}(\log p)^\kappa$, $\sigma>1$, define
\begin{equation}\label{p03:eq:arithmetic-coherence}
 \mathcal C_{b,w}(t;u)=
 \sum_{p\nmid q}w_p\bigl(1-\Re(ub_pp^{-it})\bigr),
 \qquad |u|=1.
\end{equation}
For later use define the inverse energy scale
\[
 X_{\sigma,\kappa}(\eps)
 =\eps^{-1/(\sigma-1)}[\log(1/\eps)]^{(\kappa-1)/(\sigma-1)}.
\]
The optional common phase includes the simultaneous rotation that
occurs for relative prime frequencies.

\begin{theorem}[A deterministic coherence barrier]\label{p03:thm:coherence-barrier}
Assume $\mathrm H_\vartheta(\chi)$ and put
$d=(\sigma-1)/(1-\vartheta)$. There are $c>0$ and $T_0$ such that,
for all $|t|\ge T_0$ and $|u|=1$,
\begin{equation}\label{p03:eq:coherence-barrier}
 \mathcal C_{b,w}(t;u)\ge
 c\bigl(\log(|t|+3)\bigr)^{-d}
  \bigl(\log\log(|t|+3)\bigr)^{\kappa-1}.
\end{equation}
Consequently, for sufficiently small $\eps$, a hit
$\mathcal C_{b,w}(t;u)\le C_1\eps$ with $|t|\ge T_0$ requires
\begin{equation}\label{p03:eq:coherence-first-height}
 \log(|t|+3)\ge cX_{\sigma,\kappa}(\eps)^{1-\vartheta}.
\end{equation}
Constants may depend on the fixed character, exceptional set and weights.
\end{theorem}
\begin{proof}
Put $L=\log(q(|t|+3))$ and $X=KL^{1/(1-\vartheta)}$, where $K$
is a large fixed constant. For large $|t|$, the block primes lie
outside $S$ and are coprime to $q$. The prime number theorem gives
\[
 M(X):=\sum_p(\log p)\phi(p/X)=I_\phi X+o(X).
\]
The principal main term of \eqref{p03:eq:smooth-prime-block} has modulus
at most $I_\phi X/8$ for large $|t|$. Its error is at most
$CK^{-(1-\vartheta)}X$, also at most $I_\phi X/8$ after choosing
$K$. Hence $M(X)-\Re(uS_\chi(X,t))\ge I_\phi X/2$, uniformly in
$u$. On this block,
\[
 w_p\ge cX^{-\sigma}(\log X)^{\kappa-1}\log p.
\]
Nonnegativity and boundedness of $\phi$ imply
$\mathcal C_{b,w}(t;u)\ge cX^{1-\sigma}(\log X)^{\kappa-1}$.
Substitution proves \eqref{p03:eq:coherence-barrier}.

For the last assertion take $X=c_2X_{\sigma,\kappa}(\eps)$ with
$c_2>0$ small. The energy assumption bounds
$M(X)-\Re(uS_\chi(X,t))$ by $I_\phi X/4$. Thus
$|S_\chi(X,t)|\gg X$, and \eqref{p03:eq:smooth-prime-block}, after
absorbing its principal term, forces
$X^{1-\vartheta}\ll\log(q(|t|+3))$.
\end{proof}

\begin{theorem}[A local inverse statement]\label{p03:thm:local-zero-forcing}
No hypothesis $\mathrm H_\vartheta$ is assumed here. Fix $\eta>0$
and $c_1>0$. There are $C_1,C_2>0$ such that, for sufficiently large
$X$ and $|t|$, if
\[
 |S_\chi(X,t)|\ge c_1X,\qquad
 X^{1-b}>C_1\log(q(|t|+3)),\qquad \tfrac12\le b<1,
\]
then $L(s,\chi)$ has a nontrivial zero $\rho=\beta+i\gamma$ with
\begin{equation}\label{p03:eq:local-zero-forcing}
 \beta>b,\qquad
 |\gamma-t|\le C_2\bigl(\log(q(|t|+3))\bigr)^\eta.
\end{equation}
The constants are uniform in $b$.
\end{theorem}
\begin{proof}
Put $L=\log(q(|t|+3))$. Formula \eqref{p03:eq:smooth-zero-formula}
is unconditional. Its principal term is $o(X)$ as $|t|\to\infty$.
The condition $L\ll X^{1-b}$ with $b\ge1/2$ makes the remaining
integral $O(1)$, while the prime powers contribute $O(X^{1/2})$.
The absolute zero sum is consequently at least $c_1X/2$.

Choose an integer $N>1+1/\eta$ and put $H=C_2L^\eta$.
If there is no zero as in \eqref{p03:eq:local-zero-forcing}, zeros with
$|\gamma-t|\le H$ contribute at most $C_NX^bL$; the others contribute
at most
\[
 C_NX\sum_{|\gamma-t|>H}(1+|\gamma-t|)^{-N}
 \ll_N XLH^{1-N}.
\]
The unit-interval count proves the last estimate; its extra $\log H$
is absorbed by $L$. This term is less than $c_1X/8$ after choosing
$C_2$ and the lower threshold, while $X^{1-b}>C_1L$ makes the near
term less than $c_1X/8$ after choosing $C_1$. This is a contradiction.
\end{proof}

\subsection{Consequences at the two ordinary-zeta edges}

Recall $z_2^*=-1$, $z_p^*=1$ for odd $p$. The phase-defect estimates
already proved above give
\[
 \zeta'(E_*-\delta+it)=0\ \Longrightarrow\
 \sum_p(\log p)p^{-E_*}|p^{-it}-z_p^*|^2\ll\delta,
\]
and
\[
 \zeta(A_*-\eps+it)=1\ \Longrightarrow\
 \sum_pp^{-A_*}|p^{-it}-z_p^*|^2\ll\eps.
\]
Apply Theorem~\ref{p03:thm:coherence-barrier} with $\chi=1$,
$b_p=\overline{z_p^*}$ and $u=1$. Under
$\mathrm H_\vartheta(\zeta)$ this yields
\begin{align}
 E_*-\Re\rho&\ge
 c\bigl(\log(|\Im\rho|+3)\bigr)^{-(E_*-1)/(1-\vartheta)},
 &&\zeta'(\rho)=0,\label{p03:eq:zeta-critical-boundary}\\
 A_*-\Re s&\ge
 \frac{c\bigl(\log(|\Im s|+3)\bigr)^{-(A_*-1)/(1-\vartheta)}}
 {\log\log(|\Im s|+3)},
 &&\zeta(s)=1.\label{p03:eq:zeta-one-boundary}
\end{align}
These statements apply to sufficiently small positive edge deficits.
Such points have unbounded ordinates: the exact extremal pattern cannot
occur on a bounded segment of the ordinary prime flow, and its
continuous defect has a positive minimum there.

Let $T_E(\delta)$ and $T_A(\eps)$ denote the first positive ordinates
in the respective edge strips. The deterministic recurrence and local
root-transfer upper bounds in Section~\ref{p02:sec:euler} combine with these lower bounds
to give, with $\beta_E=(E_*-1)^{-1}$ and $\beta_A=(A_*-1)^{-1}$,
\begin{align}
 \exp(c\delta^{-\beta_E(1-\vartheta)})
 &\le T_E(\delta)\le\exp(C\delta^{-\beta_E}),
 \label{p03:eq:critical-first-height-bracket}\\
 \exp\!\left(c[\eps\log(1/\eps)]^{-\beta_A(1-\vartheta)}\right)
 &\le T_A(\eps)\le
 \exp\!\left(C[\eps\log(1/\eps)]^{-\beta_A}\right).
 \label{p03:eq:one-first-height-bracket}
\end{align}
For $\vartheta=7/8$ the lower-height powers are
\[
 \frac1{8(E_*-1)}=0.0689459643\ldots,\qquad
 \frac1{8(A_*-1)}=0.1329643401\ldots.
\]
The boundary powers are $8(E_*-1)=14.5041121620\ldots$ and
$8(A_*-1)=7.5208134699\ldots$.

The local inverse theorem also gives an unconditional alternative.
If an ordinary-zeta critical point of deficit $\delta$ occurs at a
height satisfying
\[
 \delta^{-(1-b)/(E_*-1)}>C\log(|t|+3),
\]
then a zeta zero lies to the right of $b$ within
$O((\log(|t|+3))^\eta)$ of that same ordinate, for any fixed
$\eta>0$. For a $1$-point replace the left side by
$[\eps\log(1/\eps)]^{-(1-b)/(A_*-1)}$.

These first-height restrictions for arithmetic phase patterns neither
identify an exact first-height exponent nor improve the deterministic
uniform covering bound. In particular, a zero-free half-plane supplies
no new small-divisor bound for $\sum_pk_p\log p=\log(a/b)$ beyond
those used in the covering argument of Theorem~\ref{p02:thm:arithmetic-flow}.




\section{Gaussian inverse principles and unconditional edge repulsion}
\label{p04:sec:zi-gaussian}

This section improves the vertical localization in the inverse
coherence theorem and removes any fixed-half-plane hypothesis from a
weaker first-height obstruction.  The analytic input is the classical
explicit formula; the additional ingredient is to exploit positive
lognormal smoothing on the prime side.  Throughout, $\chi$ is a fixed
primitive Dirichlet character of conductor $q$, including $q=1$.
Write
\[
 L_t=\log(q(|t|+3)),\qquad M_t=\log L_t.
\]
All assertions involving $t$ are for sufficiently large $|t|$.

\subsection{A Gaussian prime block}

For $X\ge2$ put
\[
 G_\chi(X,t)=\sum_p(\log p)
       \exp\!\left(-\frac{\log^2(p/X)}2\right)\chi(p)p^{-it}.
\]
The series converges absolutely.  Mellin inversion, followed by a
contour shift to $\Re s=-1/2$, gives
\begin{align}
G_\chi(X,t)
={}&\mathbf1_{\chi=1}\sqrt{2\pi}\,
 X^{1-it}e^{(1-it)^2/2}
 -\sqrt{2\pi}\sum_\rho X^{\rho-it}e^{(\rho-it)^2/2}
 \notag\\
&+O\!\left(X^{1/2}+1+X^{-1/2}L_t\right).
\label{p04:eq:zi-gaussian-explicit}
\end{align}
The zero sum includes multiplicities and is absolutely convergent.
Indeed, the Mellin transform of $e^{-\log^2 u/2}$ is
$\sqrt{2\pi}e^{s^2/2}$, and
\[
 |e^{(\beta+i(\gamma-t))^2/2}|
 =e^{\beta^2/2-(\gamma-t)^2/2}.
\]
The unit-interval zero count is $O(\log(q(|\gamma|+3)))$.
The functional equation bounds the new contour by
$O(X^{-1/2}L_t)$.  A possible trivial zero at zero gives $O(1)$.
Finally, the prime-power counting function of exponents at least two
is $O(y^{1/2})$; integration by parts against the lognormal weight
bounds their contribution by $O(X^{1/2})$.  Gaussian decay on
horizontal contours justifies the shift.  Thus every term in
\eqref{p04:eq:zi-gaussian-explicit} is uniform in real $t$.

\begin{theorem}[Gaussian localization of a forced zero]
\label{p04:thm:zi-gaussian-inverse}
Fix $c_0>0$.  There are constants $C_0,C_1>0$ such that, if
\[
 |G_\chi(X,t)|\ge c_0X,\qquad
 X^{1-b}>C_0L_t,\qquad \tfrac12\le b<1,
\]
then $L(s,\chi)$ has a nontrivial zero $\beta+i\gamma$ with
\[
 \beta>b,\qquad |\gamma-t|\le C_1\sqrt{\log L_t}.
\]
The constants are uniform in $b$.
\end{theorem}
\begin{proof}
The pole term in \eqref{p04:eq:zi-gaussian-explicit} is $o(X)$ as
$|t|\to\infty$.  The threshold implies
$X^{-1/2}L_t\ll X^{1/2-b}\le1$, and the prime powers are $o(X)$.
Consequently the absolute zero sum is at least a fixed multiple of
$X$.  Suppose the asserted zero does not exist, and set
$H=C_1\sqrt{\log L_t}$.  Zeros with $|\gamma-t|\le H$ have
$\beta\le b$, so their contribution is $O(X^bL_t)$.
For the remaining zeros use $X^\beta\le X$ and split the Gaussian:
\[
 \sum_{|\gamma-t|>H}
 X^\beta e^{-(\gamma-t)^2/2}
 \ll XL_t e^{-H^2/4}.
\]
Here the residual sum with weight $e^{-(\gamma-t)^2/4}$ is
$O(L_t)$ by the unit-interval zero count.  Choosing $C_0$ and $C_1$
large makes both contributions smaller than the required multiple
of $X$, a contradiction.
\end{proof}

Let $w_p\asymp p^{-\sigma}(\log p)^\kappa$, where $\sigma>1$ and
$\kappa\in\mathbb R$.  Let $b_p$ be unit complex numbers agreeing
with $\chi(p)$ outside a fixed finite set, for $p\nmid q$, and set
\[
 \mathcal C(t;u)=\sum_{p\nmid q}w_p
          (1-\Re(ub_pp^{-it})),\qquad |u|=1.
\]
As in the preceding phase-defect applications, the common phase $u$
may depend on $t$.  Define
\[
 X_{\sigma,\kappa}(\varepsilon)
 =\varepsilon^{-1/(\sigma-1)}
   [\log(1/\varepsilon)]^{(\kappa-1)/(\sigma-1)}.
\]
Completing a square in $\log p$ shows that
\begin{equation}
 (\log p)e^{-\log^2(p/X)/2}
 \le C X^\sigma(\log X)^{1-\kappa}w_p
 \qquad(X\ge2).
\label{p04:eq:zi-fixed-domination}
\end{equation}
For completeness, after division by $w_p$ the exponential factor is
$X^\sigma e^{\sigma^2/2}
e^{-(\log p-\log X-\sigma)^2/2}$, multiplied by
$(\log p)^{1-\kappa}$.  Its supremum is
$O((\log X)^{1-\kappa})$ after the Gaussian is included; the portion
$\log p<\tfrac12\log X$ is exponentially smaller than every power
of $\log X$.  Also, the prime number theorem gives
$G_1(X,0)\sim\sqrt{2\pi e}\,X$.
Fixed exceptional primes contribute $o(X)$.  Therefore
$\mathcal C(t;u)\le\varepsilon$ implies
$|G_\chi(X,t)|\gg X$ at
$X=cX_{\sigma,\kappa}(\varepsilon)$ with a sufficiently small
fixed $c>0$.  This proves the following quantitative version.

\begin{corollary}[A nearby rightward zero from early coherence]
\label{p04:cor:zi-coherence-nearby}
There are $C_0,C_1>0$ such that a hit
$\mathcal C(t;u)\le\varepsilon$ forces a zero with
\[
 \beta>b,\qquad |\gamma-t|\le C_1\sqrt{\log L_t},
\]
whenever
$X_{\sigma,\kappa}(\varepsilon)^{1-b}>C_0L_t$ and
$1/2\le b<1$.  In particular, if
$X=cX_{\sigma,\kappa}(\varepsilon)$ and
$\log X\ge2\log(2C_0L_t)$, one may take
\[
 b=1-\frac{\log(2C_0L_t)}{\log X}.
\]
\end{corollary}

\subsection{Adaptive smoothing and joint localization}

The preceding window still grows, although very slowly.  Allowing
the Gaussian width to increase gives a complementary conclusion:
exceptionally small phase energy forces a zero close both to the
$1$-line and to the same ordinate.  For $v\ge1$ put
\[
 k_v(y)=\frac1{\sqrt{2\pi v}}
   \exp\!\left(-\frac v2-\frac{\log^2 y}{2v}\right),\qquad
 \mathcal H_\chi(v,t)=\sum_p(\log p)k_v(p)\chi(p)p^{-it}.
\]
The Mellin transform is exactly
\begin{equation}
 \int_0^\infty k_v(y)y^{s-1}\,dy
       =e^{v(s^2-1)/2}.
\label{p04:eq:zi-heat-mellin}
\end{equation}
The prime number theorem implies
\begin{equation}
 \mathcal H_1(v,0)=1+o(1)\qquad(v\to\infty).
\label{p04:eq:zi-heat-mass}
\end{equation}
One way to see this uniformly is to integrate against
$\vartheta(y)=y+o(y)$.  The continuous comparison measure has
$\log y$ normally distributed with mean $v$ and variance $v$,
so its mass on every bounded interval tends to zero.

The same contour shift as before gives, uniformly for $v\ge1$,
\begin{align}
 \mathcal H_\chi(v,t)
 ={}&\mathbf1_{\chi=1}e^{-vt^2/2-ivt}
   -\sum_\rho e^{v((\rho-it)^2-1)/2}
   +O\!\left(L_t e^{-3v/8}\right).
\label{p04:eq:zi-heat-explicit}
\end{align}
On $\Re s=-1/2$, the Gaussian has modulus
$e^{-3v/8}e^{-v(\Im s)^2/2}$, giving the stated contour bound.
The possible zero at zero contributes at most $e^{-v/2}$.
The prime-power contribution is $O(e^{-3v/8})$: integrate the
bound $\sum_{p^j\le y,\,j\ge2}\log p\ll y^{1/2}$ against
$k_v(y)$ and complete a square.  These observations also establish
the uniformity required when $v$ grows with $t$.

\begin{theorem}[A quadratic zero constraint]
\label{p04:thm:zi-quadratic}
Fix $\sigma,\kappa$, the weights, the character and the finite
exceptional set.  There are constants $C,c>0$ with the following
property.  Suppose
\[
 \mathcal C(t;u)\le\varepsilon,\qquad
 R=\log(1/\varepsilon)\ge C\log L_t.
\]
Then a nontrivial zero $\rho=\beta+i\gamma$ satisfies
\begin{equation}
 \beta^2-(\gamma-t)^2>
          1-c\frac{\log L_t}{R}.
\label{p04:eq:zi-quadratic}
\end{equation}
Whenever the right-hand side is positive, this implies
\[
 \beta>\sqrt{1-c\frac{\log L_t}{R}},\qquad
 |\gamma-t|<\sqrt{c\frac{\log L_t}{R}}.
\]
In particular, along any sequence of such hits for which
$R/\log L_t\to\infty$, the forced zeros satisfy
$\beta\to1$ and $\gamma-t\to0$.
\end{theorem}
\begin{proof}
Completing the square once more gives the pointwise estimate
\begin{equation}
 \frac{(\log p)k_v(p)}{w_p}
 \le C v^{1/2-\kappa}e^{v(\sigma^2-1)/2}.
\label{p04:eq:zi-heat-domination}
\end{equation}
Indeed, apart from a constant, the left side equals
\[
 v^{-1/2}e^{v(\sigma^2-1)/2}
 (\log p)^{1-\kappa}
 e^{-(\log p-\sigma v)^2/(2v)}.
\]
Its remaining supremum is $O(v^{1-\kappa})$ for $v\ge1$;
near $\sigma v$ this is immediate, and the two Gaussian tails
absorb all polynomial factors.  Fixed exceptional primes and
primes dividing $q$ have total weight $O(e^{-v/2})$.

Take $v=R/(\sigma^2-1)$.  Equations
\eqref{p04:eq:zi-heat-mass} and \eqref{p04:eq:zi-heat-domination} imply
\[
 \Re(u\mathcal H_\chi(v,t))
 \ge1-o(1)-C e^{-R/2}R^{1/2-\kappa},
\]
and hence $|\mathcal H_\chi(v,t)|\ge1/2$.
By increasing the constant in $R\ge C\log L_t$, the error in
\eqref{p04:eq:zi-heat-explicit} is less than $1/8$; its pole term is
also less than $1/8$.  Thus
\begin{equation}
 \sum_\rho e^{-v(1-\beta^2+(\gamma-t)^2)/2}\ge\tfrac14.
\label{p04:eq:zi-zero-mass}
\end{equation}
Since $0<\beta<1$, all quantities $1-\beta^2$ are nonnegative.
If every zero satisfied
$v(1-\beta^2+(\gamma-t)^2)>4\log(KL_t)$, splitting the
exponential into two equal factors would bound the left side by
\[
 \frac1{KL_t}\sum_\rho e^{-v(\gamma-t)^2/4}
 \ll\frac1K.
\]
The last bound is uniform for $v\ge1$ by the unit-interval zero
count.  Choosing $K$ large contradicts \eqref{p04:eq:zi-zero-mass}.
Substituting $v=R/(\sigma^2-1)$ and absorbing the fixed constant
inside $\log(KL_t)$ proves \eqref{p04:eq:zi-quadratic}.
\end{proof}

\subsection{An unconditional first-height obstruction}

The classical Vinogradov--Korobov zero-free region states that,
for some $a>0$, every sufficiently high zeta zero satisfies
\begin{equation}
 1-\beta\ge
 \frac{a}{(\log(|\gamma|+3))^{2/3}
              (\log\log(|\gamma|+3))^{1/3}}.
\label{p04:eq:zi-vk}
\end{equation}
For example, this form appears in Ford's zero-free-region paper~\cite{FordZeroFree};
explicit refinements do not affect the argument below.

\begin{corollary}[Unconditional arithmetic coherence repulsion]
\label{p04:cor:zi-vk-energy}
Take $\chi=1$ in the definition of $\mathcal C$.  Uniformly for
$|u|=1$, there is $C>0$ such that
\begin{equation}
 \mathcal C(t;u)\ge
 \exp\!\left[-C(\log(|t|+3))^{2/3}
                  (\log\log(|t|+3))^{4/3}\right].
\label{p04:eq:zi-vk-energy}
\end{equation}
Consequently a hit $\mathcal C(t;u)\le\varepsilon$, with
$|t|$ sufficiently large and $\varepsilon$ sufficiently small,
requires
\begin{equation}
 \log(|t|+3)\ge
 c\frac{[\log(1/\varepsilon)]^{3/2}}
          {[\log\log(1/\varepsilon)]^2}.
\label{p04:eq:zi-vk-first-height}
\end{equation}
\end{corollary}
\begin{proof}
Put $L=\log(|t|+3)$ and choose
\[
 b=1-\frac{a_1}{L^{2/3}(\log L)^{1/3}}
\]
with $a_1>0$ smaller than the constant in
\eqref{p04:eq:zi-vk}.  Every zero in the window
$|\gamma-t|\le C_1\sqrt{\log L}$ has $\beta\le b$.
Corollary~\ref{p04:cor:zi-coherence-nearby} therefore requires
\[
 (1-b)\log X_{\sigma,\kappa}(\varepsilon)
       \le\log(C_0L)
\]
for any sufficiently small coherence hit.  Since
\[
 \log X_{\sigma,\kappa}(\varepsilon)
 =\frac{\log(1/\varepsilon)
       +(\kappa-1)\log\log(1/\varepsilon)}{\sigma-1},
\]
this gives \eqref{p04:eq:zi-vk-energy}.  Inverting
$R\le C L^{2/3}(\log L)^{4/3}$ gives
$L\ge cR^{3/2}/(\log R)^2$, proving
\eqref{p04:eq:zi-vk-first-height}.
\end{proof}

Apply this corollary to the frozen ordinary-zeta phase pattern
$z_2^*=-1$, $z_p^*=1$ for odd $p$, using the phase-defect
necessity at the two Euler edges.  For
$\zeta'(E_*-\delta+it)=0$ use
$\sigma=E_*$, $\kappa=1$ and
$\mathcal C(t;1)\ll\delta$; for
$\zeta(A_*-\varepsilon+it)=1$ use
$\sigma=A_*$, $\kappa=0$ and
$\mathcal C(t;1)\ll\varepsilon$.
Define $T_E(\delta)$ as the infimum of positive ordinates of
critical points with $0<E_*-\Re s\le\delta$, and define
$T_A(\varepsilon)$ analogously for $1$-points with
$0<A_*-\Re s\le\varepsilon$. Existence for every positive
tolerance follows from Theorem~\ref{p04:thm:simultaneous-triples}.
These first positive edge ordinates satisfy, unconditionally,
\begin{align}
 \log T_E(\delta)&\ge
 c\frac{[\log(1/\delta)]^{3/2}}
          {[\log\log(1/\delta)]^2},\notag\\
 \log T_A(\varepsilon)&\ge
 c\frac{[\log(1/\varepsilon)]^{3/2}}
          {[\log\log(1/\varepsilon)]^2}.
\label{p04:eq:zi-vk-edges}
\end{align}
As before, sufficiently small edge deficits cannot occur in a fixed
bounded interval because the exact frozen phase pattern is absent
from the ordinary prime flow.  These inequalities are weaker than
the consequences of a fixed zero-free half-plane, but have no such
hypothesis.  The inverse statements do not establish a converse
equivalence with the Riemann hypothesis: an off-line zero alone
does not supply the almost complete prime-phase coherence assumed
here.

\paragraph{Comparison with other inverse theorems.}
Gaussian zero detection is an established method, not a new principle
claimed here.  In particular, Dong, Wang, Wang and Zhang's preprint
\emph{Large zeta sums and zeros of the Riemann zeta function}
\cite{DongInverse} gives a zero-forcing theorem from large ordinary
partial sums $\sum_{n\le x}n^{it}$, adapting the character-sum method
of Granville and Soundararajan.  The present hypotheses concern
near-complete weighted prime coherence and therefore differ from
their hypotheses.  The conclusions above are direct coherence and
Euler-edge applications with explicit vertical windows; no claim
of a general improvement over that partial-sum theorem is intended.




\section{Exact spectral surgery and a sharp fourth-moment certificate}
\label{p03:sec:fourmoments}

This section concerns finite Hermitian matrices. Its conclusions are unconditional linear algebra. The application to zeta in Corollary~\ref{p03:cor:conditional-four} explicitly assumes two additional trace asymptotics. The numerical constant $16/21$ also appears in the public conditional certificates of Jesse/SolidSF \cite{Jesse}; no priority claim for that constant is made here. Our purpose is to give a self-contained reduction, an exact spectral feasibility criterion, and a matching block extremal.

Write $n_+(H)$ for the number of strictly positive eigenvalues of a Hermitian matrix, with multiplicity.

\begin{proposition}[Exact spectral surgery]\label{p03:prop:spectral-surgery}
Let $G$ be Hermitian and let $\lambda_1\ge\cdots\ge\lambda_r>0$ be its positive eigenvalues. For integers $s,b\ge0$ and $\tau\ge0$, a matrix $P$ satisfying
\[
 P\succeq0,\qquad \operatorname{rank}P\le s,\qquad
 \operatorname{Tr}P\le\tau,\qquad n_+(G-P)\le b
\]
exists if and only if
\[
 (r-b)_+\le s,\qquad \sum_{j=b+1}^{r}\lambda_j\le\tau,
\]
where an empty sum is zero.
\end{proposition}
\begin{proof}
Put $Q=G-P$. Since $G\preceq P+Q_+$ and $\operatorname{rank}Q_+\le b$, the min--max principle gives
\[
 \lambda_{b+j}(G)\le\lambda_j(P)\qquad(1\le j\le r-b).
\]
This proves necessity. Conversely, with orthonormal eigenvectors $u_j$ for the positive spectrum, set
\[
 P=\sum_{j=b+1}^{r}\lambda_j u_j u_j^*.
\]
This has exactly the required rank and trace and leaves at most $b$ positive eigenvalues.
\end{proof}

The criterion makes no positivity assumption on $G$ or $Q$. If the additional condition $n_-(G-P)\le b$ is imposed, the further necessary and sufficient condition is $n_-(G)\le b$: necessity follows from $G\succeq G-P$, and the displayed construction leaves the negative spectrum unchanged.

\begin{lemma}[Spectral certificates]\label{p03:lem:spectral-certificates}
Suppose
\begin{equation}\label{p03:eq:spectral-count-model}G=P+Q,\quad P\succeq0,\quad
 \operatorname{rank}P\le s,\quad \operatorname{Tr}P\le s,
 \quad n_+(Q)\le b,\quad s+2b\le N.
\end{equation}
Let $A,B,C\ge0$ and let $q$ be a real function on $\mathbb R$ satisfying
\[
 q(x)\le0\quad(x\le0),\qquad
 q(x)\le\min\{A+Bx,C\}\quad(x>0).
\]
Then
\[
 \operatorname{Tr}q(G)\le(A+B)s+Cb.
\]
In particular, if $A=1$ and $C=2B$, then
\[
 s\ge\operatorname{Tr}q(G)-BN.
\]
\end{lemma}
\begin{proof}
The nonpositive spectrum contributes at most zero. Bound the largest $b$ positive eigenvalues by $C$ each. Proposition~\ref{p03:prop:spectral-surgery} bounds the number of remaining positive eigenvalues by $s$ and their sum by $s$. Their contribution is therefore at most $(A+B)s$. The case of fewer than $b$ positive eigenvalues has an empty tail. The final assertion follows from $2b\le N-s$.
\end{proof}

\begin{theorem}[A fourth-moment ratio]\label{p03:thm:four-moment-ratio}
Under \eqref{p03:eq:spectral-count-model}, define
\[
 L=4\operatorname{Tr}G-\operatorname{Tr}G^2-2N,
 \qquad D=\operatorname{Tr}\!\left[(G(2I-G))^2\right].
\]
If $D>0$, then
\begin{equation}\label{p03:eq:four-moment-ratio}
 s\ge\frac{(L_+)^2}{D}.
\end{equation}
If $D=0$, the hypotheses imply $L\le0$.
\end{theorem}
\begin{proof}
For $\theta>0$ put
\[
 q_\theta(x)=\frac{8\theta x-(4+2\theta)x^2+4x^3-x^4}{\theta^2}.
\]
Every term in its numerator is nonpositive when $x\le0$. Moreover the exact identities
\begin{align*}
 1+\frac{4x}{\theta}-q_\theta(x)
 &=\frac{(x^2-2x+\theta)^2}{\theta^2},\\
 \frac8\theta-q_\theta(x)
 &=\frac{(x-2)^2(x^2+2\theta)}{\theta^2}
\end{align*}
give the hypotheses of Lemma~\ref{p03:lem:spectral-certificates} with $A=1$, $B=4/\theta$, and $C=8/\theta$. Consequently
\[
 s\ge\operatorname{Tr}q_\theta(G)-\frac{4N}{\theta}
 =\frac{2L}{\theta}-\frac{D}{\theta^2}.
\]
For $L,D>0$, take $\theta=D/L$. When $L\le0$, use $s\ge0$. If $D=0$ and $L>0$, letting $\theta\downarrow0$ contradicts the displayed inequality.
\end{proof}

\begin{corollary}[Conditional moment transfer]\label{p03:cor:conditional-four}
Consider a sequence of decompositions satisfying \eqref{p03:eq:spectral-count-model}, with $N\to\infty$, and suppose
\[
 \frac1N\bigl(\operatorname{Tr}G,\operatorname{Tr}G^2,
 \operatorname{Tr}G^3,\operatorname{Tr}G^4\bigr)
 \longrightarrow\left(1,\frac43,2,\frac{13}{4}\right).
\]
Then $\liminf s/N\ge16/21$. The same conclusion holds if each count constraint has an additive $o(N)$ error.
\end{corollary}
\begin{proof}
We have $L/N\to2/3$ and $D/N\to7/12$, and apply Theorem~\ref{p03:thm:four-moment-ratio}. For the version with errors, use the fixed parameter $\theta=7/8$ in its proof; every error then remains $o(N)$.
\end{proof}

At $\theta=7/8$, the certificate is
\[
 q(x)=\frac{448x-368x^2+256x^3-64x^4}{49},\qquad B=\frac{32}{7}.
\]
The four prescribed moments give $N^{-1}\operatorname{Tr}q(G)\to16/3$, and $16/3-32/7=16/21$. More generally, if the normalized moments tend to $(1,m_2,m_3,m_4)$ with $m_4-4m_3+4m_2>0$, then the same argument gives
\[
 \liminf\frac{s}{N}\ge
 \frac{\bigl((2-m_2)_+\bigr)^2}{m_4-4m_3+4m_2}.
\]

\begin{proposition}[Sharp block extremal]\label{p03:prop:four-moment-sharp}
The constant $16/21$ cannot be increased using only the four moments in Corollary~\ref{p03:cor:conditional-four}, the constraints \eqref{p03:eq:spectral-count-model}, and the additional requirements that $d=N$, that $P$ be a sum of $s$ unit rank-one operators, and that $Q$ be a sum of positive multiplicity-two blocks.
\end{proposition}
\begin{proof}
Take $N=d=42$ and $a=1/(2\sqrt2)$. Let $P$ have sixteen eigenvalues $1-a$, sixteen eigenvalues $1+a$, and ten zero eigenvalues. Let $Q$ equal twice an orthogonal projection of rank five supported in $\ker P$. Then $s=32$, $b=5$, and $G=P+Q$ has five zero eigenvalues. Directly,
\[
 (\operatorname{Tr}G,\operatorname{Tr}G^2,
 \operatorname{Tr}G^3,\operatorname{Tr}G^4)
 =\left(42,56,84,\frac{273}{2}\right).
\]
Thus $L=28$, $D=49/2$, and $(L_+)^2/D=32=16N/21$.

For each pair of eigenvectors $e_+,e_-$ of $P$, set
\[
 v_\pm=\sqrt{\frac{1+a}{2}}\,e_+
       \mathbin{\pm}\sqrt{\frac{1-a}{2}}\,e_-.
\]
Both vectors have norm one, and the sum of their rank-one operators is $\operatorname{diag}(1+a,1-a)$. Sixteen such pairs realize $P$ as the required sum. The five summands $2ee^*$ realize $Q$. Orthogonal sums of the construction give arbitrarily large exact examples.
\end{proof}

\begin{theorem}[A surface of sharp moment data]\label{p03:thm:sharp-moment-surface}
Let $1\le R<2$, set $u=2-R$, and suppose $u^2\le D_0\le u$. Among asymptotic matrix sequences satisfying \eqref{p03:eq:spectral-count-model}, $d=N$, and normalized moment limits
\[
 (m_1,m_2,m_3,m_4)=(1,R,3R-2,8R-8+D_0),
\]
the optimal lower bound for $\liminf s/N$ is exactly $u^2/D_0$. Sharpness persists when $P$ is required to be a sum of unit rank-one operators and $Q$ a sum of positive multiplicity-two blocks.
\end{theorem}
\begin{proof}
Theorem~\ref{p03:thm:four-moment-ratio} gives the lower bound. Put $s_0=u^2/D_0$ and $\theta=D_0/u$, so $0<s_0\le1$ and $0<\theta\le1$. Consider the spectral probability measure with masses $(1-s_0)/2$ at each of $0,2$ and masses $s_0/2$ at each of $1-\sqrt{1-\theta},1+\sqrt{1-\theta}$. Assign the latter two masses to $P$ and the mass at $2$ to $Q$. Its first moment is one, and its next three moments are
\[
 2-s_0\theta=R,\qquad
 4-3s_0\theta=3R-2,\qquad
 8-8s_0\theta+s_0\theta^2=8R-8+D_0.
\]
For rational $s_0$, choose dimensions making all multiplicities integral. For general $s_0$, choose even integers $s_N$ with $s_N/N\to s_0$ and $N-s_N$ even, and keep $\theta$ fixed in the same construction. Then $\operatorname{rank}P=\operatorname{Tr}P=s_N$, and $Q$ has rank $(N-s_N)/2$. The two-vector construction in Proposition~\ref{p03:prop:four-moment-sharp} supplies the required unit-vector realization, and the moments converge to the prescribed values.
\end{proof}

These extremals concern the matrix count model, not realizability by actual zeta ordinates or by a band-limited evaluation frame. Such geometry is additional information. In the Weil-compression application, the first two moments are available unconditionally at the critical bandwidth. The third and fourth moments assumed in Corollary~\ref{p03:cor:conditional-four} are the unproved analytic input; this section supplies no unconditional improvement to a zeta zero proportion.




\section{Spectral stability and limitations of moment amplification}
\label{p04:sec:spectral-frontier}

Continue with the matrix count model
\eqref{p03:eq:spectral-count-model}. The results here refine the
preceding certificates and remain finite-dimensional statements;
they assert no new trace asymptotics for the zeta compression.

\begin{proposition}[A dimension obstruction]\label{p04:prop:frontier-dimension}
For every Hermitian matrix $G$ of dimension $d\le N/2$,
\eqref{p03:eq:spectral-count-model} is feasible with $s=0$.
If, in addition, $G\succeq0$, $\operatorname{tr}G=N$, and $N$ is even,
then $G$ is a sum of $N/2$ operators of the form $2vv^*$ with
$\|v\|=1$.
\end{proposition}
\begin{proof}
For the first assertion take $P=0$, $Q=G$, and $b=n_+(G)$.
For the second, set $n=N/2$. The nonnegative eigenvalues of $G/2$,
padded with zeros to length $n$, have sum $n$ and majorize
$(1,\ldots,1)$. The finite Schur--Horn theorem gives a positive
semidefinite $n$-by-$n$ Gram matrix with this spectrum and diagonal
$(1,\ldots,1)$. Factoring it gives $n$ unit vectors whose frame
operator is unitarily equivalent to $G/2$.
\end{proof}

The unit vectors in this construction need not be evaluation vectors
of a band-limited function space. Such geometry is extra information.
For a sampling grid of spacing $2\pi/L$ in $[T,2T]$, the ratio of
dimension to total zero count tends to $\eta$ when
$L/\log T\to\eta$. The elementary fourth-moment diagonal argument
for prime products of length $e^{2L}$ has the safe range
$\eta<1/2$. Proposition~\ref{p04:prop:frontier-dimension} explains why
this direct shortening of bandwidth cannot, by itself, produce a
positive simple-zero proportion in the count model.

\begin{proposition}[A third-moment obstruction]\label{p04:prop:frontier-third}
For $1\le R\le2$, there are asymptotic examples satisfying
\eqref{p03:eq:spectral-count-model}, with $d=N$, such that
\[
 \frac1N\bigl(\operatorname{tr}G,
 \operatorname{tr}G^2,\operatorname{tr}G^3\bigr)
 \longrightarrow (1,R,3R-2),
 \qquad \frac{s}{N}\longrightarrow 2-R.
\]
Here $P$ can be a sum of unit rank-one operators and $Q$ a sum of
multiplicity-two unit rank-one operators.
\end{proposition}
\begin{proof}
Put $u=2-R$. Assign spectral masses $u$ to $1$ and $(1-u)/2$ to
each of $0,2$. Put the $1$-eigenspace in $P$ and the $2$-eigenspace
in $Q$. All constraints are equalities. The moments are as stated,
and the fourth normalized moment is $7R-6$. Rational approximation
of the multiplicities treats general $u$.
\end{proof}

In particular the three moments $(1,4/3,2)$ alone cannot improve
$2/3$ within this model. There is also an elementary certificate
obstruction: a nonconstant real polynomial which is nonpositive
on $(-\infty,0]$ and bounded above on $[0,\infty)$ has even degree
and negative leading coefficient. Thus a certificate in this class
of degree at most three already has degree at most two.

\begin{theorem}[An exact stability remainder]\label{p04:thm:frontier-stability}
Assume \eqref{p03:eq:spectral-count-model}. Set
\[
 L_0=4\operatorname{tr}G-\operatorname{tr}G^2-2N,
 \qquad D_0=\operatorname{tr}\!\left(\bigl(G(2I-G)\bigr)^2\right).
\]
For $0<\theta<1$ define
\[
 \mathcal S_\theta=\{0,2,1-\sqrt{1-\theta},1+\sqrt{1-\theta}\},
 \qquad
 c_\theta=\min\left\{\frac2\theta,
                   \frac{1-\theta}{\theta^2}\right\}.
\]
Writing all eigenvalues of $G$ with multiplicity as $\lambda_j$, one has
\begin{equation}\label{p04:eq:frontier-stability}
 s\ge \frac{2L_0}{\theta}-\frac{D_0}{\theta^2}
       +c_\theta\sum_{j=1}^d
            \operatorname{dist}(\lambda_j,\mathcal S_\theta)^2.
\end{equation}
If $L_0,D_0>0$ and $D_0/L_0\in(0,1)$, the first two terms at
$\theta=D_0/L_0$ equal $L_0^2/D_0$.
\end{theorem}
\begin{proof}
Let $\lambda_1\ge\cdots\ge\lambda_r>0$ be the positive spectrum,
put $k=\min(b,r)$, $n=r-k$, and
$S=\sum_{j=k+1}^r\lambda_j$. The min--max principle applied to
$G\preceq P+Q_+$ gives $n\le s$ and $S\le s$.
Set
\[
 q_\theta(x)=
 \frac{8\theta x-(4+2\theta)x^2+4x^3-x^4}{\theta^2}.
\]
The identities
\begin{align*}
 \frac8\theta-q_\theta(x)
 &=\frac{(x-2)^2(x^2+2\theta)}{\theta^2},\\
 1+\frac{4x}\theta-q_\theta(x)
 &=\frac{(x^2-2x+\theta)^2}{\theta^2}
\end{align*}
give the exact decomposition
\begin{align}\label{p04:eq:frontier-defect-identity}
 s-\frac{2L_0}\theta+\frac{D_0}{\theta^2}
 ={}&\frac4\theta(N-s-2b)+(s-n)+\frac4\theta(s-S)
                +\frac8\theta(b-k)\notag\\
 &+\sum_{j=1}^k
       \frac{(\lambda_j-2)^2(\lambda_j^2+2\theta)}{\theta^2}
   +\sum_{j=k+1}^r
       \frac{(\lambda_j^2-2\lambda_j+\theta)^2}{\theta^2}
   -\sum_{\lambda_j<0}q_\theta(\lambda_j).
\end{align}
Every term on the right is nonnegative. The first spectral sum
dominates $(2/\theta)\sum_{j\le k}|\lambda_j-2|^2$.
For the second put $a_\pm=1\pm\sqrt{1-\theta}$. At least one of
$|x-a_-|,|x-a_+|$ is at least $\sqrt{1-\theta}$, so
\[
 (x^2-2x+\theta)^2
 \ge(1-\theta)\min\{|x-a_-|,|x-a_+|\}^2.
\]
Finally, for $y\ge0$,
\[
 -q_\theta(-y)=
 \frac{8\theta y+(4+2\theta)y^2+4y^3+y^4}{\theta^2}
 \ge\frac{4+2\theta}{\theta^2}y^2.
\]
Comparison with the distance to $\mathcal S_\theta$ proves
\eqref{p04:eq:frontier-stability}.
\end{proof}

\begin{corollary}[A geometric strengthening of the conditional target]
\label{p04:cor:frontier-geometric}
Suppose a sequence satisfying \eqref{p03:eq:spectral-count-model} has normalized
first four moments tending to $(1,4/3,2,13/4)$. If
\[
 \liminf\frac1N\sum_{j=1}^d
   \operatorname{dist}\left(\lambda_j,
           \left\{0,2,1-\frac1{2\sqrt2},1+\frac1{2\sqrt2}\right\}\right)^2
 \ge\varepsilon,
\]
then
\[
 \liminf\frac{s}{N}\ge\frac{16}{21}+\frac8{49}\varepsilon.
\]
\end{corollary}
\begin{proof}
Use $\theta=7/8$ in Theorem~\ref{p04:thm:frontier-stability}.
\end{proof}

The numerical value $16/21$ is already present in public conditional
certificates~\cite{Jesse}; the additional statement here is the
explicit spectral-distance remainder. Neither the four-moment
hypothesis nor a positive spectral-distance
lower bound is proved here for the zeta compression. The corollary
isolates the extra information that could make evaluation geometry
stronger than the unrestricted matrix model.

\begin{proposition}[Exact nonattainment for sinc evaluation vectors]
\label{p04:prop:frontier-sinc}
Let $v_x$ be unit vectors in a Hilbert space, indexed by real $x$,
with
\[
 \langle v_x,v_y\rangle=
 \frac{\sin\pi(x-y)}{\pi(x-y)}.
\]
The kernel takes its continuous value $1$ on the diagonal.
Let $x_1,\ldots,x_s,y_1,\ldots,y_b$ be distinct real numbers, with
$s,b>0$, and put
\[
 P=\sum_{j=1}^s v_{x_j}v_{x_j}^*,\qquad
 Q=2\sum_{j=1}^b v_{y_j}v_{y_j}^*,\qquad
 G=P+Q,\qquad N=s+2b.
\]
For every $0<\theta<1$, the defect on the left of
\eqref{p04:eq:frontier-defect-identity} is strictly positive.
\end{proposition}
\begin{proof}
Work on the finite-dimensional span of the displayed vectors. If
the defect were zero, its nonnegative decomposition would force
exactly $b$ top positive eigenvalues equal to $2$, all other
positive eigenvalues in $\{a_-,a_+\}$, and a positive tail of size
$s$, where $a_\pm=1\pm\sqrt{1-\theta}<2$.
Thus $G\preceq2I$ and, since $0\preceq Q\preceq G$,
$Q\preceq2I$. But $\operatorname{rank}Q\le b$ and
$\operatorname{tr}Q=2b$, so $Q$ is twice an orthogonal projection
of rank $b$. Positivity of $2I-G$ then shows that $P$ vanishes on
the range of $Q$. Consequently every $v_{x_j}$ is orthogonal to
$v_{y_1}$, and hence $x_j-y_1$ is a nonzero integer. The distinct
$x_j$ therefore differ from one another by nonzero integers, so
the $v_{x_j}$ are mutually orthogonal. Thus $P$ is a projection.
As $P$ and $Q$ have orthogonal ranges, the positive tail spectrum
of $G$ is identically $1$, contradicting $a_\pm\ne1$.
\end{proof}

This finite nonattainment statement does not supply a uniform
positive defect per zero as the dimension grows. It also concerns
real sinc evaluation vectors and a positive simple-plus-double
decomposition; off-line Weil pairs need not satisfy these
hypotheses. Converting nonattainment into a quantitative
asymptotic gap is a separate problem.




\section{Remaining analytic questions}

The fixed-phase theorem and its uniform extensions have different
quantifiers. Theorem~\ref{p04:thm:near-logarithmic} treats a fixed
finite prime set, whereas Theorem~\ref{p04:thm:simultaneous-triples}
uses a growing set at its stated tolerance. The conditional profiles
of Sections~\ref{p04:sec:bounded-transfer}--\ref{p04:sec:joint-transition}
take the height limit first. No argument above permits exchanging
these orders of limits at an arbitrary prescribed tolerance.

Further analytic questions are whether one can remove the extra
power of $\log\log T$ from fixed-prime resolution, obtain the exact
entropy constant in a joint height-and-tolerance limit, improve the
edge approach rate, or determine the limiting phase density of all
simple critical zeros. A nondegenerate joint fluctuation law for
the companion roots is not supplied by bounded-density comparison.

Higher-moment amplification has two distinct missing inputs: the
required third and fourth trace asymptotics at a usable bandwidth,
and quantitative information on the geometry of the evaluation
vectors. The exact matrix extremals show what the unrestricted
count model allows; finite sinc nonattainment does not furnish an
asymptotic spectral gap. Similarly, a zero-free half-plane bounds
smooth prime sums but does not automatically improve the rational
frequency separation used in deterministic covering.


\clearpage

% ===== Integrated Part 3 =====
\part{Rare-event laws and the recovery of hidden spectra}
\label{part:rare-events}The preceding prime-phase analysis leads to an inverse question: which
features of a weighted independent-site model survive in its rare events?
The common local probability law and the individual spectral locations
are different unknowns. A normalized continuum rate can recover the
first while erasing the second; a centered fluctuation law can recover
the information lost at the continuum scale. This part proves that
separation, calibrates unknown power weights, and then examines what
changes when the local variables are correlated vectors.

All identification assertions use exact functions or exact probability
laws, with the normalizations and model classes stated in each result.
An open patch of a rate function is infinite information. A channel
count is not a sample-complexity bound, and local derivative data are
not supplied automatically by noisy probability measurements. The
hard-conditioning statements below concern the specified independent
product model. The final sinc construction is an abstract sampling
obstruction and does not construct a configuration of zeta zeros.

\paragraph{Relation to established methods.}
Scalar weighted small deviations are treated by Aurzada
\cite{Aurzada}, and Mellin methods for small-deviation expansions by
Fill and Torcaso \cite{FillTorcaso}. Nonvanishing Mellin transforms
as an inverse and cancellation principle are central to Jacobsen,
Mikosch, Rosi\'nski, and Samorodnitsky \cite{JMRS}. Spectral
uniqueness for Gamma convolutions is classical \cite{JRY};
second-chaos cumulants and limits are studied in
\cite{NP,AEK}. Joint constraints and condensation are treated
in \cite{SEM}. General $\Upsilon$-transforms, matrix mixing, and
homogeneous-measure decompositions are established subjects
\cite{BRT,BPT,EM}; dominant-atom signed L\'evy expansions
and random-integral inversion are likewise classical
\cite{BKL,Jurek}. The results below concern particular
observation designs and the information they retain. The use of these
classical mechanisms is not itself a novelty claim.

\section{From a regularly varying spectrum to joint rare events}
\label{p05:sec:rv}

This section supplies the probabilistic meaning of the continuum transform.
Let $0<r_1\le r_2\le\cdots\to\infty$, with repetitions permitted, and put
\[
 N(x)=\#\{n:r_n\le x\}.
\]
Assume that $N$ is regularly varying with positive index $\rho$:
$N(cx)/N(x)\to c^\rho$ for every $c>0$.
Let $V_n$ be independent copies of a nonnegative, nonconstant random
variable $V$ with essential infimum zero and $\mathbb E V^2<\infty$.
Throughout this section suppose
\begin{equation}\label{p05:eq:laplace-assumptions}
 h(u)=-\log\mathbb E e^{-uV},\qquad
 h'(u)\sim\frac{\alpha}{u},\qquad
 -h''(u)\sim\frac{\alpha}{u^2}\quad(u\to\infty),\qquad \alpha>0.
\end{equation}
In particular $h(u)=O(1+\log(1+u))$, and
$g(u):=-h''(u)>0$. The function $h$ is twice continuously differentiable
at zero and smooth for $u>0$. The assumptions hold, for example, when $V$ has a
density $c v^{\alpha-1}(1+O(v^\eta))$ near zero, with $c,\eta>0$, and finite
second moment. They include both bounded phase potentials and Gamma laws.

Fix $E>A>\rho$ and define the almost surely finite sums
\[
 D_E=\sum_n r_n^{-E}V_n,\qquad D_A=\sum_n r_n^{-A}V_n,
 \qquad X_x=\frac{x^E D_E}{N(x)},\quad Y_x=\frac{x^A D_A}{N(x)}.
\]
The continuum Laplace exponent and its lower-tail dual are
\begin{align}
 H(s,t)&=\rho\int_0^\infty y^{\rho-1}
 h(sy^{-E}+ty^{-A})\,dy,\label{p05:eq:H-general}\\
 I(e,a)&=\sup_{s,t\ge0}\{H(s,t)-se-ta\},\qquad e,a>0.
 \label{p05:eq:I-general}
\end{align}

\begin{theorem}[Joint lower-tail limit]\label{p05:thm:rv-ldp}
Under these assumptions, for every $e,a>0$,
\[
 \lim_{x\to\infty}-\frac{1}{N(x)}
 \log\mathbb P(X_x\le e,\ Y_x\le a)=I(e,a).
\]
The optimizing pair in \eqref{p05:eq:I-general} exists, is unique, and is
nonzero. Furthermore,
\begin{equation}\label{p05:eq:duality-recover}
 H(s,t)=\inf_{e,a>0}\{I(e,a)+se+ta\},\qquad s,t>0.
\end{equation}
\end{theorem}

\begin{proof}
Regular variation gives vague convergence
\[
 \frac{1}{N(x)}\sum_n\delta_{r_n/x}
 \ \longrightarrow\ \rho y^{\rho-1}\,dy
 \quad\text{on }(0,\infty).
\]
This convergence extends to the integrand defining $H$.
Here are the endpoint estimates needed for that extension. Fix
$0<\delta<\min(\rho,A-\rho)$. Potter bounds control the normalized number
of points below $xy$ by $C y^{\rho-\delta}$ for $y\le1$ and by
$C y^{\rho+\delta}$ for $y\ge1$, apart from finitely many original
points. Near zero the integrand is $O(1+|\log y|)$; near infinity it
is $O(y^{-A})$. Integration by parts therefore makes both discarded
tails uniformly small. Each exceptional fixed point contributes
$O(\log x)/N(x)=o(1)$. Consequently
\begin{equation}\label{p05:eq:laplace-limit}
 \frac{1}{N(x)}\sum_n h\bigl(s(x/r_n)^E+t(x/r_n)^A\bigr)
 \longrightarrow H(s,t).
\end{equation}
The same argument applies locally uniformly for $(s,t)$ in the positive
quadrant, and to first and second derivatives there. For example, at
$y\downarrow0$ the $ss$-derivative is bounded by $C y^{\rho-1}$ when
$s>0$; the other derivatives are no worse. At infinity use boundedness
of the first two derivatives of $h$ and $A>\rho$.

The function $H$ is strictly concave on the nonnegative quadrant:
apply the strict concavity of $h$ pointwise, observing that two distinct
parameter pairs give different arguments except possibly at one value
of $y$. In the positive quadrant its negative directional second
derivative is the integral of
\[
 g(sy^{-E}+ty^{-A})(u y^{-E}+v y^{-A})^2,
\]
which is positive because $E\ne A$. Since $h$ is increasing and concave
with $h(0)=0$, it is subadditive. Scaling yields
\[
 H(s,t)\le H(s,0)+H(0,t)=C_Es^{\rho/E}+C_At^{\rho/A}
\]
for positive constants $C_E,C_A$. Thus the objective in
\eqref{p05:eq:I-general} is coercive and has a unique maximizer. The
maximizer is not the origin, since $s^{\rho/E}$ dominates $s$ as
$s\downarrow0$.

Chernoff's inequality and \eqref{p05:eq:laplace-limit} give the asserted
lower bound on the negative logarithm. For the reverse bound first
consider positive tilt parameters $(s,t)$. Change the product law by
the density proportional to
\[
 \exp\{-s x^E D_E-t x^A D_A\}.
\]
Under this law $(X_x,Y_x)$ converges in probability to
$(H_s(s,t),H_t(s,t))$: expectations converge by the derivative version
of \eqref{p05:eq:laplace-limit}, and both variances are $O(1/N(x))$.
On a small rectangle around this limit, the inverse density has a
uniform exponential lower bound. This proves the matching probability
lower bound at all gradients of positive parameters.

Boundary optimizers require one endpoint observation. At $(s,0)$ with
$s>0$, both first derivatives are finite. At $(0,t)$ with $t>0$,
$H_s(0,t)$ is finite precisely when $E-A<\rho$, because its integrand
near zero is asymptotic to $(\rho\alpha/t)y^{\rho-1-(E-A)}$.
If this derivative is infinite, $(0,t)$ cannot maximize the objective
for finite $e$. If it is finite, the tilted first moments converge;
concentration follows by restricting to $\eta\le r_n/x\le R$, where
the variance is $O(1/N(x))$, and making the omitted nonnegative tails
small in expectation. This avoids any assumption of a finite boundary
Hessian.

At any maximizing pair $(s_*,t_*)$, multiply both coordinates by
$1+\varepsilon$. Each finite tilted first moment strictly decreases.
The active constraints then have a strict margin below their target;
the inactive constraints already have a weak margin by the
one-sided optimality inequalities and also strictly decrease. Apply
the preceding tilted lower bound on a sufficiently small rectangle.
Let its width tend to zero and then let $\varepsilon\downarrow0$.
Terms belonging to inactive constraints have zero multiplier, so the
limit of the change-of-measure cost is exactly $I(e,a)$.

Finally, the definition of $I$ gives one inequality in
\eqref{p05:eq:duality-recover}. For its reverse, take
$(e,a)=\nabla H(s,t)$, whose coordinates are positive and finite.
Concavity makes $(s,t)$ the optimizer for this pair, giving equality.
\end{proof}

\begin{corollary}[Universality and its lost information]\label{p05:cor:support-universality}
For fixed $\rho,E,A$ and microscopic law $V$, all spectra with a
regularly varying counting function of index $\rho$ have the same
normalized joint rate $I$. In particular, when $\rho=1$, the positive
integers and the primes have identical continuum rates after
normalization by their respective counting functions.
\end{corollary}

The normalization is essential: $N(x)\sim x$ for integers and
$N(x)\sim x/\log x$ for primes. The corollary says that the continuum
rate forgets the individual locations of the spectrum, not that the
unnormalized probabilities are equal. The later fluctuation theorem
recovers information that this leading-order limit discards.


\section{Microscopic recovery from a joint rate}
\label{p05:sec:inverse}

Fix $E>A>\rho$ and use the assumptions and notation of
Section~\ref{p05:sec:rv}. The single-coordinate continuum rate retains only
one number from the law of $V$. The joint rate, by contrast, determines
that law completely.

\begin{theorem}[Two-energy identifiability]\label{p05:thm:two-energy-inverse}
For fixed $\rho,E,A$, knowledge of $I(e,a)$ for all $e,a>0$ uniquely
determines the distribution of $V$. More explicitly, recover $H$ by
\eqref{p05:eq:duality-recover}, and put
\[
 K(t)=-H_{st}(1,t),\qquad g(u)=-h''(u),\qquad
 (\mathcal Mf)(z)=\int_0^\infty u^{z-1}f(u)\,du.
\]
For
\[
 a(z)=1+\frac{A-\rho}{E}-\frac{A}{E}z,
 \qquad \omega(z)=1+\frac{A-\rho}{E}+\frac{E-A}{E}z,
\]
one has
\begin{equation}\label{p05:eq:inverse-mellin}
 (\mathcal MK)(z)=\frac{\rho}{E}
 \frac{\Gamma(a(z))\Gamma(z)}{\Gamma(\omega(z))}
 (\mathcal Mg)(\omega(z)),\qquad \Re z=1.
\end{equation}
The multiplier on the right is nonzero. Hence this identity gives the
Mellin transform of $g$ on the whole vertical line
$\Re w=2-\rho/E$, and uniquely determines $g$. The remaining recovery is
\begin{equation}\label{p05:eq:recover-h-law}
 h(u)=\int_0^\infty\min(u,v)g(v)\,dv,
 \qquad \mathbb E e^{-uV}=e^{-h(u)}.
\end{equation}
\end{theorem}

\begin{proof}
The hypotheses imply
$0<g(u)\le C(1+u)^{-2}$, since $g$ is continuous at zero and has the
specified asymptotic at infinity. Differentiating $H$ in the positive
quadrant gives
\[
 K(t)=\rho\int_0^\infty
 y^{\rho-E-A-1}g(y^{-E}+t y^{-A})\,dy.
\]
For real $z=1$, Tonelli's theorem permits exchanging integrals.
For complex $z$ on this line the same nonnegative integral proves
absolute convergence and hence Fubini's theorem applies. Substitute
$u=t y^{-A}$ and then $v=y^{-E}$. This gives
\[
 (\mathcal MK)(z)
 =\frac{\rho}{E}\int_0^\infty\int_0^\infty
 v^{a(z)-1}u^{z-1}g(v+u)\,du\,dv.
\]
On the specified line,
$\Re a(z)=1-\rho/E>0$ and
$\Re\omega(z)=2-\rho/E\in(1,2)$.
Changing to $r=u+v$ and $q=u/(u+v)$ evaluates the inner beta integral
and proves \eqref{p05:eq:inverse-mellin}. Its absolute convergence also
follows directly from $g(u)=O(1)$ near zero and $g(u)=O(u^{-2})$ at
infinity.

The Gamma function has neither zeros nor poles in the right
half-plane. The multiplier is therefore finite and nonzero. As
$\Im z$ ranges over $\mathbb R$, so does
$\Im\omega(z)$ because $E\ne A$. Mellin uniqueness here is simply
Fourier uniqueness applied to the integrable function
$v\mapsto e^{cv}g(e^v)$, where $c=2-\rho/E$.
Thus $g$ is determined almost everywhere and, by continuity, everywhere.
Finally $h'(u)\to0$ and $h(0)=0$. Integrating $-h''=g$ twice gives
\eqref{p05:eq:recover-h-law}; both integrals converge by the preceding
bounds. Uniqueness of the Laplace transform of a probability measure
finishes the proof.
\end{proof}

The assertion is uniqueness from exact limiting data. It is not a claim
that finitely many noisy rare-event measurements give stable recovery.
The explicit multiplier will quantify that distinction below.

\begin{proposition}[No finite list of marginal rates is sufficient]
\label{p05:prop:marginals-insufficient}
Fix any finite list $\sigma_1,\ldots,\sigma_m>\rho$.
There exist distinct laws satisfying \eqref{p05:eq:laplace-assumptions}
whose scalar continuum transforms
\[
 H_{\sigma_j}(s)=\rho\int_0^\infty
 y^{\rho-1}h(sy^{-\sigma_j})\,dy
\]
agree for every $s\ge0$ and every $j$. Their scalar lower-tail rates
therefore agree as well. The examples can have the same value of
$\alpha$ in \eqref{p05:eq:laplace-assumptions}.
\end{proposition}

\begin{proof}
Choose $m+2$ distinct positive scales $c_1,\ldots,c_{m+2}$.
The $(m+1)\times(m+2)$ real matrix with rows
\[
 (1,\ldots,1),\qquad
 (c_1^{\rho/\sigma_j},\ldots,c_{m+2}^{\rho/\sigma_j}),
 \quad 1\le j\le m,
\]
has a nonzero null vector $b$. Choose $\eta>0$ so small that
$\alpha_i^\pm=1\pm\eta b_i>0$. Let $V^\pm$ be the sum of independent
Gamma variables of shapes $\alpha_i^\pm$ and scales $c_i$.
Their Laplace exponents are
\[
 h_\pm(u)=\sum_i\alpha_i^\pm\log(1+c_i u).
\]
These laws have finite second moments, essential infimum zero, and
the same asymptotic coefficient
$\alpha=\sum_i\alpha_i^\pm=m+2$. Direct scaling and one integration
by parts give
\[
 H_{\sigma_j,\pm}(s)
 =\frac{\pi}{\sin(\pi\rho/\sigma_j)}
 s^{\rho/\sigma_j}\sum_i\alpha_i^\pm c_i^{\rho/\sigma_j}.
\]
The null-vector equations make these expressions identical.
The laws are nevertheless distinct: equality of $h_+'$ and $h_-'$
would imply
$\sum_i b_i c_i/(1+c_i u)=0$ for all $u>0$.
Distinct poles of this rational function force every $b_i=0$, a
contradiction.
\end{proof}

\subsection{Exact data on one active patch suffice}

\begin{lemma}[Analyticity in the positive quadrant]
\label{p05:lem:H-analytic}
Under \eqref{p05:eq:laplace-assumptions}, $H$ is real analytic on
$(0,\infty)^2$. Its Hessian is negative definite there.
\end{lemma}

\begin{proof}
Negative definiteness was proved in Theorem~\ref{p05:thm:rv-ldp}.
For analyticity let $L(u)=\mathbb E e^{-uV}$ and note that
$M=\sup_{u>0}u h'(u)<\infty$.
For $u>0$ and $|z-u|\le\delta u$ with $0<\delta<1$, integration along
the segment from $u$ to $z$ yields
\[
 \left|\frac{L(z)}{L(u)}-1\right|
 \le\delta u\,\frac{L((1-\delta)u)}{L(u)}h'((1-\delta)u)
 \le\frac{\delta M}{(1-\delta)^{M+1}}.
\]
Choose $\delta$ so that the last expression is less than $1/2$.
Thus $L$ is nonzero on these relative disks, and the branch of
$-\log L$ agreeing with $h$ on the positive axis is holomorphic there.
Fix $s_0,t_0>0$. For sufficiently small complex relative perturbations
of $s_0,t_0$, the number $sy^{-E}+ty^{-A}$ lies in such a disk about
$s_0y^{-E}+t_0y^{-A}$, uniformly in $y$.
Near $y=0$ the complex integrand is bounded by
$C(1+|\log y|)$. Near infinity it is $O(y^{-A})$: use
$|L(z)-1|\le |z|\mathbb EV$ for $\Re z\ge0$ and the logarithm near
$1$. These bounds are integrable against $y^{\rho-1}dy$.
Holomorphic parameter integration proves the assertion.
\end{proof}

\begin{theorem}[Local active-patch identifiability]
\label{p05:thm:local-inverse}
Let $I_1,I_2$ arise from two microscopic laws satisfying the stated
assumptions, with the same known $\rho,E,A$.
Suppose they agree on a nonempty open set $U\subset(0,\infty)^2$
on which both optimizing multipliers are positive for both laws.
Then the two microscopic laws coincide.
\end{theorem}

\begin{proof}
For a positive optimizer, the inverse function theorem applied to
$\nabla H$ gives a locally real analytic optimizer as a function of
$(e,a)$. On this region
\[
 (s,t)=-\nabla I(e,a),\qquad
 \nabla^2 I(e,a)=-[\nabla^2 H(s,t)]^{-1},\qquad
 H(s,t)=I(e,a)+se+ta.
\]
Consequently agreement of the two rates on $U$ gives agreement of
their $H$ functions on a nonempty open set in the positive quadrant.
The real analytic identity theorem on that connected quadrant extends
the agreement everywhere. Theorem~\ref{p05:thm:two-energy-inverse} applies.
\end{proof}

The active-patch condition is essential. In a region where only one
constraint is active, the rate is independent of the other coordinate
and supplies only marginal information. Thus the information gain
comes from observing the genuinely joint part of the rate surface.


\section{Exact smoothing and a regularized inverse}
\label{p05:sec:inverse-stability}

The injectivity theorem concerns exact data.  Its inverse is strongly
ill-conditioned, and the degree of ill-conditioning can be computed
explicitly.  This section concerns the linear map from the tilted variance
$g=-h''$ to the mixed derivative $K(t)=-H_{st}(1,t)$; it does not assert a
minimax lower bound over probability distributions.

Put
\[
 r=\frac{\rho}{E},\qquad b=\frac{A}{E},\qquad q=1-b,
 \qquad c=2-r,
\]
so that $0<r<b<1$.  Define logarithmic coordinates
\[
 f(v)=e^{cv}g(e^v),\qquad k(v)=e^vK(e^v),
 \qquad d(v)=qk(qv).
\]
We use the Fourier convention
$\widehat f(\xi)=\int_{\mathbb R}e^{i\xi v}f(v)\,dv$.
The bounds $g(u)=O(1)$ at zero and $g(u)=O(u^{-2})$ at infinity
imply $f\in L^1(\mathbb R)\cap L^2(\mathbb R)$.  The Mellin identity gives
\begin{equation}\label{p05:eq:inverse-multiplier}
 \widehat d(\xi)=m(\xi)\widehat f(\xi),\qquad
 m(\xi)=r B\!\left(1-r-i\frac{b}{q}\xi,
                         1+i\frac{\xi}{q}\right).
\end{equation}
In particular $d\in L^2(\mathbb R)$.

\begin{proposition}[Exact smoothing scale]\label{p05:prop:inverse-smoothing}
The multiplier in \eqref{p05:eq:inverse-multiplier} never vanishes, and
\begin{equation}\label{p05:eq:inverse-multiplier-asymptotic}
 |m(\xi)|\sim
 r\sqrt{2\pi}\,b^{1/2-r}q^{r-1}
 |\xi|^{-1/2}
 \exp\!\left(-\frac{\pi A}{E-A}|\xi|\right)
 \qquad (|\xi|\longrightarrow\infty).
\end{equation}
Consequently the inverse of this linear map is unbounded on its
$L^2$ range.
\end{proposition}

\begin{proof}
The two arguments of the beta function have positive real parts, as does
their sum.  The gamma function has no zeros and has no poles at these
arguments.  Thus $m$ does not vanish.  Apply
\[
 |\Gamma(a+it)|\sim\sqrt{2\pi}\,|t|^{a-1/2}e^{-\pi|t|/2}
\]
to the two gamma functions in the numerator and the one in the
denominator.  The imaginary parts are $-b\xi/q$, $\xi/q$, and $\xi$;
their exponential factors give
$\exp(-\pi b|\xi|/q)$.  The powers of $|\xi|$ combine to
$|\xi|^{-1/2}$, and the remaining factors give the displayed constant.
For unboundedness, take real $L^2$ functions of norm one whose Fourier
transforms are supported in $[n,n+1]\cup[-n-1,-n]$.  Their images have
norm at most $\sup_{|\xi|\ge n}|m(\xi)|$, which tends to zero.
These perturbations are ambient $L^2$ functions; they need not be tilted
variances of probability laws.
\end{proof}

\begin{proposition}[Spectral cutoff bound]\label{p05:prop:inverse-cutoff}
Let $\ell>0$ and suppose $\|f\|_{H^\ell}\le M$.  If measured data satisfy
$\|d_\delta-d\|_2\le\delta$, define
\[
 \widehat f_{\delta,R}(\xi)
 =\mathbf 1_{\{|\xi|\le R\}}\frac{\widehat d_\delta(\xi)}{m(\xi)}.
\]
For $R\ge1$,
\begin{equation}\label{p05:eq:inverse-cutoff-bound}
 \|f_{\delta,R}-f\|_2
 \le C\delta(1+R)^{1/2}
       \exp\!\left(\frac{\pi A}{E-A}R\right)
       +M(1+R^2)^{-\ell/2},
\end{equation}
where $C$ depends only on $E,A,\rho$.  For any fixed $0<\eta<1$,
the choice
\[
 R=\frac{\eta(E-A)}{\pi A}\log(1/\delta)
\]
gives, for sufficiently small $\delta$ and fixed $M>0$,
$\|f_{\delta,R}-f\|_2=O((\log(1/\delta))^{-\ell})$.
\end{proposition}

\begin{proof}
Continuity, nonvanishing, and
\eqref{p05:eq:inverse-multiplier-asymptotic} imply
\[
 \sup_{|\xi|\le R}|m(\xi)|^{-1}
 \le C(1+R)^{1/2}e^{\pi AR/(E-A)}.
\]
Plancherel's theorem bounds the error on $[-R,R]$ by this supremum
times $\delta$.  The Sobolev bound gives the second term outside that
interval.  For the stated choice of $R$, the first term is
$O(\delta^{1-\eta}\sqrt{\log(1/\delta)})$ and the second is
$O((\log(1/\delta))^{-\ell})$.
\end{proof}

\begin{remark}
If the error is instead specified as $\|k_\delta-k\|_2\le\delta$, then
$\|d_\delta-d\|_2\le\sqrt q\,\delta$, and the same conclusion holds.
The noise assumption is on a mixed derivative of the recovered free
energy, after passage to logarithmic coordinates.  A noise bound on the
original rate surface does not automatically imply this assumption:
convex duality and numerical differentiation require a separate error
analysis.  The factor $\pi A/(E-A)$ tends to infinity when the two
exponents coalesce.  At $E=A$ the transform ceases to be injective.
\end{remark}


\section{Microscopic fluctuations retain the discrete spectrum}
\label{p05:sec:microscopic-fluctuations}

The following statement requires no counting law for the weights.  It
isolates information that a leading large-deviation integral can erase.
Let $V\ge0$ have finite second moment and write
\[
 h(u)=-\log\mathbb E e^{-uV},\qquad
 g(u)=-h''(u)=\operatorname{Var}_{u}(V),\qquad u\ge0,
\]
where $\mathbb P_u$ denotes the one-coordinate exponential tilt.
Assume, for a fixed $\alpha>0$,
\begin{equation}\label{p05:eq:microscopic-one-site-hypotheses}
 uh'(u)\longrightarrow\alpha,
 \qquad u^2g(u)\longrightarrow\alpha
 \quad(u\to\infty).
\end{equation}
In particular $M:=\sup_{u\ge0}u^2g(u)<\infty$; finiteness near zero
follows from $\mathbb E V^2<\infty$.

\begin{theorem}[Square-summable ratio fluctuation theorem]
\label{p05:thm:square-summable-ratio-fluctuations}
Let $V_n$ be independent copies of $V$, and let $b_n>0$ satisfy
$\sum_n b_n<\infty$.  For a real sequence $a=(a_n)\in\ell^2$ set
\[
 D=\sum_n b_nV_n,\qquad X=\sum_n a_nb_nV_n,
 \qquad
 d\mathbb P_t=\frac{e^{-tD}}{\mathbb E e^{-tD}}\,d\mathbb P.
\]
Both series are absolutely convergent almost surely.  Put
$m_t=\mathbb E_tX=\sum_n a_nb_nh'(tb_n)$.
Then, as $t\to\infty$,
\begin{equation}\label{p05:eq:canonical-gamma-series-limit}
 t(X-m_t)\ \Longrightarrow\
 Y_a:=\sum_{n\ge1}a_n(G_n-\alpha),
 \qquad G_n\overset{\rm iid}{\sim}\operatorname{Gamma}(\alpha,1).
\end{equation}
The series defining $Y_a$ converges in $L^2$ and almost surely.
The variances also converge:
\[
 t^2\operatorname{Var}_t(X)\longrightarrow
 \alpha\sum_n a_n^2.
\]
The characteristic function and the cumulants of the limiting law are
\begin{align}
 \mathbb E e^{izY_a}
 &=\prod_n e^{-i\alpha za_n}(1-iza_n)^{-\alpha},
 \label{p05:eq:gamma-series-characteristic-function}\\
 \kappa_k(Y_a)&=\alpha(k-1)!\sum_n a_n^k,
 \qquad k\ge2.
 \label{p05:eq:gamma-series-cumulants}
\end{align}
In particular the limit is non-Gaussian whenever $a\ne0$.
\end{theorem}

\begin{proof}
Cauchy--Schwarz gives
$\sum_n|a_n|b_n\le\|a\|_2\|b\|_2<\infty$.
Together with $\mathbb EV<\infty$, this proves absolute convergence
of both original series.  Exponential tilting preserves independence,
and coordinate $n$ has parameter $u=tb_n$.
For $z\ge0$, the Laplace transform of $uV$ under $\mathbb P_u$ is
\[
 \mathbb E_u e^{-zuV}
 =\exp\{h(u)-h((1+z)u)\}
 \longrightarrow(1+z)^{-\alpha},
\]
because
$h((1+z)u)-h(u)=\int_1^{1+z}uh'(uv)\,dv
\to\alpha\log(1+z)$.
Also $\mathbb E_u(uV)=uh'(u)\to\alpha$ and
$\operatorname{Var}_u(uV)=u^2g(u)\to\alpha$.
Thus each finite vector of the centered coordinates converges to
independent copies of $G-\alpha$.

For every $R$ and every $t>0$, independence gives the uniform tail bound
\begin{equation}\label{p05:eq:gamma-series-uniform-tail}
 \mathbb E_t\left|
 \sum_{n>R}a_n\{tb_nV_n-tb_nh'(tb_n)\}\right|^2
 \le M\sum_{n>R}a_n^2.
\end{equation}
Finite-dimensional convergence followed by $R\to\infty$ proves
\eqref{p05:eq:canonical-gamma-series-limit}.  Dominated convergence in
$\sum_n a_n^2(tb_n)^2g(tb_n)$ proves variance convergence.
The limiting independent centered series has summable variances,
which also proves its almost-sure convergence.

The logarithm of each factor in
\eqref{p05:eq:gamma-series-characteristic-function} is $O_z(a_n^2)$
for all sufficiently large $n$, so the product converges.  Since an
$\ell^2$ sequence is bounded, the logarithm is analytic in a fixed
neighborhood of zero and can be expanded absolutely there.  This
gives \eqref{p05:eq:gamma-series-cumulants}.  For $a\ne0$ its fourth
cumulant is $6\alpha\sum_n a_n^4>0$, excluding a Gaussian law.
\end{proof}

\begin{corollary}[The law determines the nonzero coefficient spectrum]
\label{p05:cor:fluctuation-spectral-recovery}
For fixed known $\alpha>0$, the law of $Y_a$ determines the multiset
of nonzero real coefficients $a_n$, including multiplicities.
\end{corollary}

\begin{proof}
The finite measure $\lambda_a=\sum_n a_n^2\delta_{a_n}$ has compact
support.  All its moments are determined by the law, since
\[
 \int x^j\,d\lambda_a(x)
 =\frac{\kappa_{j+2}(Y_a)}{\alpha(j+1)!},\qquad j\ge0.
\]
Compactly supported finite measures are determined by their moments:
on a compact interval containing the supports of two candidates,
polynomials are uniformly dense in continuous functions.  For each
nonzero $x$, its multiplicity is $\lambda_a(\{x\})/x^2$.
Zero coefficients are invisible and are expressly excluded.
\end{proof}

\begin{remark}[Scope of the inversion]
The recovery corollary is exact identifiability, not a quantitative
or numerically stable recovery algorithm.  It reconstructs a multiset;
the ordering of the coefficients and the underlying weights $b_n$
cannot be recovered without additional structure.  Neither theorem
asserts that conditioning on $D\le\varepsilon$ may replace the
canonical tilt.  Such a replacement needs a separate local limit
argument.  Nor can a nondegenerate limiting law be transferred merely
from bounds $c\mu\le\nu\le\mu$.
\end{remark}


\section{Identifying the law and the spectrum together}
\label{p05:sec:two-scale}

The preceding results can be combined into an exact inverse theorem
for two different layers of a model. All parameters and normalizations
in the following statement are part of the observed model; in
particular, the two power exponents are known.

\begin{theorem}[Two-scale identifiability]\label{p05:thm:two-scale}
Fix $E>A>\rho>0$ with $2(E-A)>\rho$. For each of two candidate models,
let its positive spectrum tend to infinity with a counting function
regularly varying of index $\rho$, and let its common microscopic law
satisfy \eqref{p05:eq:laplace-assumptions} and have finite second moment.
Use the normalization of Theorem~\ref{p05:thm:rv-ldp}.
Suppose the models have the following two pieces of identical data:
\begin{enumerate}[label=(\roman*)]
\item their joint lower-tail rate surfaces agree on a nonempty open
patch where both constraints are active for both models;
\item under the tilt proportional to $e^{-tD_A}$, their limiting
centered fluctuation laws of $t(D_E-\mathbb E_tD_E)$ agree.
\end{enumerate}
Then the microscopic laws coincide and the spectral location
multisets coincide, including multiplicities.
\end{theorem}

\begin{proof}
The first datum and Theorem~\ref{p05:thm:local-inverse} identify the
microscopic law, and in particular the coefficient $\alpha$.
Let $d=E-A$. Regular variation and $2d>\rho$ imply
$\sum_n r_n^{-2d}<\infty$. Apply
Theorem~\ref{p05:thm:square-summable-ratio-fluctuations} with
$b_n=r_n^{-A}$ and $a_n=r_n^{-d}$. The second datum and
Corollary~\ref{p05:cor:fluctuation-spectral-recovery} identify the positive
coefficient multiset $(r_n^{-d})$. The injective map
$a\mapsto a^{-1/d}$ recovers the original spectral locations.
\end{proof}

The two data sources are mathematically distinct. The first is a
limiting logarithmic probability surface, normalized by the number
of modes. The second is an order-one distributional limit under a
specified tilted law. Neither is replaced here by a finite sample or
by a bounded comparison between probability measures. The theorem is
a uniqueness statement for exact asymptotic data; a statistical
estimation theorem would require additional assumptions and error
analysis.


\section{A two-scale distinction between integers and primes}
\label{p05:sec:two-scale-example}

The two inverse questions concern different objects. The continuum rate
identifies the common law of the independent variables, whereas a
non-Gaussian fluctuation law can identify the individual spectral
locations. The following example separates these statements explicitly.

Take independent unit exponential variables $V_n$, so
$h(u)=\log(1+u)$, $g(u)=(1+u)^{-2}$, and $\alpha=1$.
For a spectrum $\mathcal R$, put
\[
 D_2^{\mathcal R}=\sum_{r\in\mathcal R}r^{-2}V_r,
 \qquad D_3^{\mathcal R}=\sum_{r\in\mathcal R}r^{-3}V_r.
\]
Consider either $\mathcal R=\mathbb N$ or $\mathcal R=\mathbb P$,
the set of primes. The counting functions are respectively asymptotic
to $x$ and $x/\log x$. Both have index $\rho=1$, so their normalized
joint lower-tail rate is the same function:
\[
 H(s,t)=\int_0^\infty\log(1+s y^{-3}+t y^{-2})\,dy,
 \qquad I(e,a)=\sup_{s,t\ge0}\{H(s,t)-se-ta\}.
\]
The inverse theorem recovers the unit exponential microscopic law from
this common surface. It cannot recover whether the spectrum consists
of primes or all positive integers from this surface alone.

Now impose the canonical tilt proportional to $e^{-tD_2^{\mathcal R}}$.
Under it the variables remain independent and $V_r$ has exponential
rate $1+t/r^2$. Writing $m_{\mathcal R}(t)$ for the tilted mean of
$D_3^{\mathcal R}$, the fluctuation theorem gives
\begin{equation}\label{p05:eq:integer-prime-fluctuation}
 t\bigl(D_3^{\mathcal R}-m_{\mathcal R}(t)\bigr)
 \ \Longrightarrow\ Z_{\mathcal R}
 :=\sum_{r\in\mathcal R}\frac{G_r-1}{r},
 \qquad G_r\ \text{independent unit exponentials}.
\end{equation}
The series converges in $L^2$ and almost surely. Its characteristic
function and cumulants are
\[
 \mathbb E e^{iuZ_{\mathcal R}}
 =\prod_{r\in\mathcal R}e^{-iu/r}(1-iu/r)^{-1},
 \qquad
 \kappa_k(Z_{\mathcal R})=(k-1)!\sum_{r\in\mathcal R}r^{-k},\quad k\ge2.
\]
For $\mathcal R=\mathbb N$, the classical Weierstrass product
for the Gamma function \cite{DLMF} gives
\[
 \mathbb E e^{iuZ_{\mathbb N}}=e^{-i\gamma u}\Gamma(1-iu).
\]
Thus $Z_{\mathbb N}$ is a standard Gumbel random variable centered at
its mean $\gamma$. For primes one instead obtains
\[
 \kappa_k(Z_{\mathbb P})=(k-1)!P(k),\qquad
 P(z)=\sum_p p^{-z}\quad(\Re z>1).
\]
Already the variances differ: $\zeta(2)$ for the integer spectrum and
$P(2)$ for the prime spectrum. All cumulants together determine the
entire reciprocal spectrum, including multiplicities, by the compact
moment argument in the fluctuation section.

\begin{remark}[What the example does and does not establish]
The Gamma product, the Gumbel representation, and uniqueness of the
associated spectral measure are classical ingredients. Their role
here is to exhibit exactly which information is discarded by the
continuum rare-event limit and recovered by a second asymptotic scale.
Equation \eqref{p05:eq:integer-prime-fluctuation} is a statement under an
explicit probability tilt. It is not, by itself, a limit law along
the zeros of the Riemann zeta function, nor a prime-detection algorithm
from finitely many noisy observations.
\end{remark}

\subsection{An exactly computable transform check}
For this exponential microscopic law,
\[
 (\mathcal M g)(w)=\Gamma(w)\Gamma(2-w),\qquad 0<\Re w<2.
\]
The mixed Mellin data therefore equal
\[
 (\mathcal M K)(z)
 =\frac{\rho}{E}B(a(z),z)\Gamma(\omega(z))\Gamma(2-\omega(z)),
\]
with $a(z)$ and $\omega(z)$ defined in the inverse theorem.
In particular
\[
 \int_0^\infty K(t)\,dt
 =\frac{\pi\rho/E}{\sin(\pi\rho/E)}.
\]
This identity checks both the density factor $\rho$ and the Mellin
normalization without a numerical inverse problem.

\section{Analytic and probabilistic foundations}
\label{p06:sec:scalar-foundation}

We extend the continuum foundation to any finite number of distinct
scalar power channels and record the inversion facts needed for blind
calibration. The two-power and fluctuation proofs have already been
given above and will be used directly.
Throughout this section $V\ge0$ is nondegenerate,
$\mathbb EV^2<\infty$, and
\begin{equation}\label{p06:eq:scalar-foundation-hypotheses}
 h(u)=-\log\mathbb E e^{-uV},\qquad
 uh'(u)\longrightarrow\alpha>0,\qquad
 u^2g(u)\longrightarrow\alpha,\qquad g=-h''.
\end{equation}
Differentiation of the Laplace transform shows that $h'(u)$ is the
mean and $g(u)>0$ is the variance of $V$ under the tilt proportional
to $e^{-uV}$.  In particular
\begin{equation}\label{p06:eq:scalar-foundation-bounds}
 h(u)=O(\log(2+u)),\quad h(u)\le u\mathbb EV,\quad
 h'(u)\le\frac{C}{1+u},\quad
 0<g(u)\le\frac{C}{(1+u)^2}.
\end{equation}
Here the first bound is meant globally and follows by integrating
$h'$; continuity at zero follows from the finite second moment.

\begin{lemma}[Continuum transform and the fully active rate]
\label{p06:lem:scalar-continuum-foundation}
Let $r_n>0$ tend to infinity and suppose
$N(x)=\#\{n:r_n\le x\}$ is regularly varying of index $\rho>0$:
$N(xy)/N(x)\to y^\rho$ for every $y>0$.
Fix distinct powers $\sigma_1,\ldots,\sigma_m>\rho$ and define
\[
 D_j=\sum_n r_n^{-\sigma_j}V_n,\qquad
 Z_{x,j}=\frac{x^{\sigma_j}D_j}{N(x)},\qquad
 H(s)=\rho\int_0^\infty y^{\rho-1}
 h\left(\sum_j s_jy^{-\sigma_j}\right)dy,
\]
where the $V_n$ are independent copies of $V$.  The series $D_j$
are finite almost surely.  For $s$ in the positive orthant,
\[
 -\frac1{N(x)}\log\mathbb E e^{-N(x)s\cdot Z_x}
 =\frac1{N(x)}\sum_nh\left(\sum_j s_j(x/r_n)^{\sigma_j}\right)
 \longrightarrow H(s).
\]
The convergence also holds for derivatives through order two locally
in that orthant.  The function $H$ is real analytic and strictly
concave.  At every $z=\nabla H(s)$ with $s>0$,
\begin{equation}\label{p06:eq:scalar-active-ldp}
 -\frac1{N(x)}\log\mathbb P\{Z_{x,j}\le z_j\text{ for all }j\}
 \longrightarrow I(z):=H(s)-s\cdot z
 =\sup_{v\ge0}\{H(v)-v\cdot z\}.
\end{equation}
Near each such point,
$-\nabla I(z)=s$ and $\nabla^2I(z)=-[\nabla^2H(s)]^{-1}$.
\end{lemma}

\begin{proof}
Regular variation gives $N(x)=O(x^{\rho+\eta})$ for every fixed
$\eta>0$, so summation over dyadic intervals proves
$\sum_nr_n^{-\sigma_j}<\infty$.  This and $\mathbb EV<\infty$
prove almost-sure summability.  The measures
$N(x)^{-1}\sum_n\delta_{r_n/x}$ converge on compact subintervals of
$(0,\infty)$ to $\rho y^{\rho-1}dy$.  We detail the required tail
control so that compact convergence suffices here.
For every sufficiently small $\eta>0$, the regular-variation ratio
bound gives, whenever $x$ and $xy$ exceed a fixed constant,
\[
 \frac{N(xy)}{N(x)}\le C_\eta
 \begin{cases}y^{\rho-\eta},&0<y\le1,\\
 y^{\rho+\eta},&y\ge1.
 \end{cases}
\]
This standard bound follows by iterating the defining ratio limit
on a fixed geometric mesh and using monotonicity between mesh points.
Near $y=0$ the summand is $O(1+|\log y|)$, by
\eqref{p06:eq:scalar-foundation-bounds}; near infinity it is
$O(y^{-\min_j\sigma_j})$.  Summing over geometric annuli and choosing
$\eta<\min(\rho,\min_j\sigma_j-\rho)$ gives uniformly vanishing
tails.  The finitely many terms with bounded $r_n$ contribute at most
$O(\log x/N(x))=o(1)$.

For the first two derivatives the relevant summands are
$y^{-\sigma_i}h'(u)$ and
$y^{-\sigma_i-\sigma_j}g(u)$, with
$u=\sum_ks_ky^{-\sigma_k}$.  On compact positive parameter sets,
both are bounded near zero, since each $s_ky^{-\sigma_k}\le u$;
near infinity they are bounded by the respective integrable powers.
The same argument gives locally uniform derivative convergence.
Moreover $-\nabla^2H$ is the Gram matrix of the functions
$y^{-\sigma_j}$ in the positive measure
$\rho y^{\rho-1}g(u)dy$.  Distinct powers are linearly independent,
so the matrix is positive definite.

For analyticity write $L(u)=\mathbb E e^{-uV}$ and
$M_1=\sup_{u>0}u h'(u)<\infty$.
If $|w-u|\le\delta u$, $u>0$, and $0<\delta<1$, then
\[
 \left|\frac{L(w)}{L(u)}-1\right|
 \le\delta u\,\frac{L((1-\delta)u)}{L(u)}h'((1-\delta)u)
 \le\frac{\delta M_1}{(1-\delta)^{M_1+1}}.
\]
The first inequality follows by integrating $L'$ along the segment;
the second follows by integrating $h'(v)\le M_1/v$.
Choose $\delta$ so the final bound is less than $1/2$.
Thus $L$ has a holomorphic logarithm in every such relative disk.
A sufficiently small complex relative perturbation of all $s_j$
keeps $u$ in the corresponding disk, uniformly in $y$.
The resulting logarithm is bounded by $C(1+|\log y|)$ near zero
and $Cy^{-\min\sigma_j}$ near infinity.  For the latter bound use
$|L(w)-1|\le |w|\mathbb EV$ when $\Re w\ge0$.
Holomorphic parameter integration proves real analyticity of $H$.

Write $H_x$ for the normalized logarithmic Laplace transform.
The exponential tilt with density
$\exp\{-N(x)s\cdot Z_x+N(x)H_x(s)\}$ has
mean $\nabla H_x(s)$ and covariance
$-\nabla^2H_x(s)/N(x)$ for $Z_x$.
The established derivative convergence therefore gives concentration
at $\nabla H(s)$.  The Chernoff bound proves the lower bound on the
negative logarithm in \eqref{p06:eq:scalar-active-ldp}.
For its reverse, use the tilt with parameter $(1+\varepsilon)s$.
Every coordinate of $\nabla H((1+\varepsilon)s)$ is strictly less
than the corresponding coordinate of $z$, because all entries of
$\nabla^2H$ are negative.  A small box around this new mean lies
inside the lower orthant with corner $z$ and has tilted probability
tending to one.  Undo the tilt on that box; then let its radius tend
to zero and $\varepsilon\downarrow0$.  This gives the reverse
bound.  Strict concavity makes $s$ the unique optimizer, and the
inverse function theorem gives the stated derivative identities.
\end{proof}

\begin{lemma}[Two-power inversion from an active patch]
\label{p06:lem:scalar-mellin-foundation}
For known $E>A>\rho>0$, the values of the joint rate $I$ on any
nonempty open fully active patch uniquely determine the function
$h$ and hence the distribution of $V$.
More explicitly, once $H$ is recovered, put
$K(t)=-H_{st}(1,t)$ and
$(\mathcal Mf)(z)=\int_0^\infty u^{z-1}f(u)du$.
For $\Re z=1$, set
\[
 a(z)=1+\frac{A-\rho}{E}-\frac AEz,\qquad
 w(z)=1+\frac{A-\rho}{E}+\frac{E-A}{E}z.
\]
Then
\begin{equation}\label{p06:eq:scalar-foundation-mellin}
 (\mathcal MK)(z)=\frac\rho E
 \frac{\Gamma(a(z))\Gamma(z)}{\Gamma(w(z))}
 (\mathcal Mg)(w(z)),
 \qquad
 h(u)=\int_0^\infty\min(u,v)g(v)dv.
\end{equation}
The same inversion determines a positive constant multiple of a
rescaled $h$ if such a function is used in place of $h$.
\end{lemma}

\begin{proof}
The local-patch assertion is Theorem~\ref{p05:thm:local-inverse},
and the displayed Mellin identity and reconstruction are exactly
Theorem~\ref{p05:thm:two-energy-inverse}, with $w=\omega$.
Those proofs establish absolute convergence, the nonvanishing multiplier,
and Mellin uniqueness. All inversion steps are linear in $h$ or $g$,
and their bounds persist under positive multiplication and positive
rescaling of the argument. This proves the final assertion without
assuming that the resulting function is itself a probability
Laplace exponent.
\end{proof}

\begin{lemma}[Canonical Gamma-series fluctuation limit]
\label{p06:lem:scalar-gamma-foundation}
Let $b_n>0$ with $\sum_nb_n<\infty$, and let $a=(a_n)\in\ell^2$
be real.  For independent copies $V_n$ of $V$, put
$D=\sum_nb_nV_n$, $X=\sum_na_nb_nV_n$, and let
$\mathbb P_t$ have density $e^{-tD}/\mathbb E e^{-tD}$.
Both original series converge absolutely almost surely, and
\begin{equation}\label{p06:eq:scalar-gamma-foundation-limit}
 t(X-\mathbb E_tX)\ \Longrightarrow\
 Y_a=\sum_na_n(G_n-\alpha),\qquad
 G_n\overset{\mathrm{iid}}{\sim}\operatorname{Gamma}(\alpha,1).
\end{equation}
The limiting series converges almost surely and in $L^2$,
its variance is $\alpha\sum_na_n^2$, and its cumulants are
\begin{equation}\label{p06:eq:scalar-gamma-foundation-cumulants}
 \kappa_k(Y_a)=\alpha(k-1)!\sum_na_n^k,
 \qquad k\ge2.
\end{equation}
For known $\alpha$, its law determines the multiset of nonzero
coefficients $a_n$, including multiplicities.
\end{lemma}

\begin{proof}
The convergence, summability, variance, and cumulant assertions are
Theorem~\ref{p05:thm:square-summable-ratio-fluctuations}.
Recovery of the nonzero coefficient multiset, including multiplicity,
is Corollary~\ref{p05:cor:fluctuation-spectral-recovery}. Those results
allow arbitrary real square-summable coefficients, as required here.
\end{proof}

For $b_n=r_n^{-A}$ and $a_n=r_n^{-(E-A)}$, regular variation,
$A>\rho$, and $2(E-A)>\rho$ give precisely the summability required
in Lemma~\ref{p06:lem:scalar-gamma-foundation}.  This identifies the
continuum inversion and the microscopic fluctuation statements as
two distinct, fully justified pieces of data.

\section{Blind calibration and the exact coordinate ambiguity}
\label{p06:sec:blind-exponents}

The power exponents need not be supplied to the inverse problem.
Their ratios to the counting exponent can be recovered from the local
curvature of the joint rate.  This is a finite-dimensional consequence
of weighted Euler homogeneity, with a sign structure that rules out
singular calibration matrices.

Let $m\ge2$, let $\sigma_1,\ldots,\sigma_m>\rho>0$ be distinct,
and set $p_j=\sigma_j/\rho$.  For the one-coordinate Laplace exponent
$h(u)=-\log\mathbb E e^{-uV}$, assume $V\ge0$ is nondegenerate,
$\mathbb EV^2<\infty$, and, for some $\alpha>0$,
\[
 uh'(u)\longrightarrow\alpha,
 \qquad -u^2h''(u)\longrightarrow\alpha.
\]
Put
\[
 H(s)=\rho\int_0^\infty y^{\rho-1}
 h\left(\sum_{j=1}^m s_jy^{-\sigma_j}\right)dy
 =\int_0^\infty h\left(\sum_{j=1}^m s_jx^{-p_j}\right)dx,
 \qquad s\in(0,\infty)^m,
\]
and $I(z)=\sup_{s\ge0}\{H(s)-s\cdot z\}$.
The second representation follows from $x=y^\rho$.
The Hessian of $H$ is negative definite: with $g=-h''>0$,
its negative is the Gram matrix of the linearly independent functions
$x^{-p_j}$ in the measure
$g(\sum_j s_jx^{-p_j})dx$.
Every entry of $-\nabla^2H$ is positive.

\begin{theorem}[Calibration from a single second-order jet]
\label{p06:thm:blind-jet-calibration}
At any point $z$ where the optimizing vector is strictly positive,
define, using only the first and second derivatives of the observed rate,
\[
 s=-\nabla I(z),\qquad
 M=[\nabla^2I(z)]^{-1},\qquad
 B=\operatorname{diag}(z)-M\operatorname{diag}(s).
\]
Then $B$ is invertible, $B^{-1}$ is entrywise nonnegative, and
\begin{equation}\label{p06:eq:blind-calibration-formula}
 (p_1,\ldots,p_m)^{\mathsf T}=B^{-1}z.
\end{equation}
Thus the first and second derivatives of $I$ at a single fully active
point determine every ratio $\sigma_j/\rho$, independently of the
unknown microscopic distribution.
\end{theorem}

\begin{proof}
The change of variable $x=\lambda v$ gives
$H(\lambda^{p_1}s_1,\ldots,\lambda^{p_m}s_m)=\lambda H(s)$.
Euler differentiation gives $\sum_jp_js_jH_j=H$.  Differentiating
once more, for every $i$, yields
\[
 p_iH_i+\sum_jp_js_jH_{ij}=H_i.
\]
At the saddle point, $\nabla H=z$ and $-\nabla^2H=M$, so these
identities say $Bp=z$.  The off-diagonal entries of $B$ are negative.
Let $P=\operatorname{diag}(p)$.  The matrix $BP$ has positive row
sums $z_i$, nonpositive off-diagonal entries, and therefore positive
diagonal entries satisfying strict row diagonal dominance.  Writing
$BP=D-C$ with $D$ positive diagonal and $C\ge0$ having zero diagonal,
the row sums of $D^{-1}C$ are strictly less than one.  Hence
\[
 (BP)^{-1}=\sum_{k\ge0}(D^{-1}C)^kD^{-1}\ge0.
\]
This proves invertibility and inverse nonnegativity, and then the
formula.  The identities connecting the derivatives of $I$ and $H$
follow from the inverse function theorem and strict concavity.
\end{proof}

The matrix argument is the classical nonsingular $M$-matrix
criterion; its role here is to prove that calibration cannot fail at
an exceptional fully active point; see also~\cite{BP}. The theorem concerns exact
derivatives.  It does not replace derivative estimation from noisy data.
There is also a useful exact residual certificate.  For any trial vector
$\widehat p$, put $r=B\widehat p-z$.  Since
$P^{-1}B^{-1}\operatorname{diag}(z)$ is nonnegative and has every row
sum equal to one,
\begin{equation}\label{p06:eq:blind-calibration-certificate}
 \max_j\left|\frac{\widehat p_j-p_j}{p_j}\right|
 \le \max_i\frac{|r_i|}{z_i}.
\end{equation}
The right side is computable from the exact local data without knowing
$p$.

\begin{theorem}[Blind two-scale identification modulo a power change]
\label{p06:thm:blind-two-scale}
Consider two models with unknown parameters $E>A>\rho>0$, positive spectra
$r_n\to\infty$ whose counting functions are regularly varying of index
$\rho$, and microscopic laws satisfying the preceding assumptions.
Write $N(x)=\#\{n:r_n\le x\}$,
$D_E=\sum_n r_n^{-E}V_n$, and $D_A=\sum_n r_n^{-A}V_n$.
For each model use $N(x)$ as the large-deviation speed and
$(x^ED_E/N(x),x^AD_A/N(x))$ as the normalized threshold variables,
and assume $2(E-A)>\rho$.  The data below are only the resulting
limiting functions and laws, not the full pointwise function $N$.
Suppose that
\begin{enumerate}
\item the normalized joint rate surfaces agree on a nonempty open patch
on which both constraints are active;
\item under the tilt proportional to $e^{-tD_A}$, the limiting laws of
$t(D_E-\mathbb E_tD_E)$ agree.
\end{enumerate}
Then the microscopic laws coincide, the ratios
$p=E/\rho$ and $q=A/\rho$ coincide, and the multisets
\[
 \{r_n^\rho:n\ge1\}
\]
coincide, including multiplicities.  Equivalently, the only remaining
ambiguity is
\begin{equation}\label{p06:eq:power-coordinate-gauge}
 (\rho,E,A,r_n)\longmapsto
 (c\rho,cE,cA,r_n^{1/c}),\qquad c>0,
\end{equation}
up to a permutation of the spectrum.  Every transformation in
\eqref{p06:eq:power-coordinate-gauge} does preserve both observations.
\end{theorem}

\begin{proof}
Theorem~\ref{p06:thm:blind-jet-calibration} identifies $p,q$ from any one
point of the common patch.  Replacing $y$ by $y^\rho$ reduces the
continuum transform to the same known exponents $p,q$ and counting
exponent one.  The active patch recovers a patch of $H$ by Legendre
duality.  Lemma~\ref{p06:lem:scalar-mellin-foundation} then identifies
$h$ and the law of $V$, including $\alpha$.

By Lemma~\ref{p06:lem:scalar-gamma-foundation}, the microscopic
fluctuation limit is
\[
 Y=\sum_n a_n(G_n-\alpha),\qquad
 a_n=r_n^{-(E-A)},\quad
 G_n\overset{\mathrm{iid}}{\sim}\operatorname{Gamma}(\alpha,1).
\]
Here $\sum_n a_n^2<\infty$ by regular variation and the assumed
inequality.  Its cumulants determine
$\sum_n a_n^2\delta_{a_n}$, a compactly supported finite measure,
because
$\kappa_{k+2}(Y)/[\alpha(k+1)!]=\sum_n a_n^{k+2}$.
Moment uniqueness on a compact interval determines every positive
$a_n$ and its multiplicity.  Hence the data determine
\[
 r_n^\rho=a_n^{-1/(p-q)}.
\]
Writing $c=\rho_2/\rho_1$ gives precisely
\eqref{p06:eq:power-coordinate-gauge}.  Conversely, this transformation
leaves every coefficient $r_n^{-E}$ and $r_n^{-A}$ unchanged,
and sends $N(x)$ to $N(x^c)$, which is regularly varying with index
$c\rho$.  It also leaves $p,q,h$ and the normalized continuum
transform unchanged.  Both observations are therefore preserved.
\end{proof}

One calibrated exponent, or the value of $\rho$, removes the ambiguity.
Alternatively, one identified spectral location $r_*\ne1$ together
with its recovered canonical location $R_*=r_*^\rho$ determines
$\rho=\log R_*/\log r_*$, and therefore every absolute exponent and
spectral location.  Knowing a location equal to one supplies no such
calibration.  Since the power change leaves all weighted sums unchanged
pathwise, further observations formed from the same transformed powers
cannot distinguish the equivalent models.

\begin{proposition}[Finite-jet non-identification of the microscopic law]
\label{p06:prop:finite-jets-insufficient}
Fix distinct powers $p_1,\ldots,p_m>1$, finitely many positive parameter
points, and any finite derivative order $k\ge2$.  There exist two
distinct admissible microscopic laws having the same $\alpha$ such that
their continuum transforms have identical jets through order $k$ at
all those points.  At the corresponding common fully active threshold
points, the rate functions also have identical jets through order $k$.
On any fixed admissible spectrum, their limiting centered Gamma-series
fluctuation laws also agree, since those limits depend on the
microscopic law only through $\alpha$.
\end{proposition}

\begin{proof}
Choose positive distinct scales $c_1,\ldots,c_N$, where $N$ is larger
than the total number of prescribed real jet entries plus one.
Let
\[
 h_\pm(u)=\sum_{\ell=1}^N(1\pm\eta b_\ell)\log(1+c_\ell u).
\]
Every prescribed derivative of the corresponding $H$ is linear in
the coefficients $1\pm\eta b_\ell$.  Impose equality of all of them,
and impose $\sum_\ell b_\ell=0$.  This homogeneous real linear system
has a nonzero solution $b$.  Choose $\eta>0$ small enough that all
shapes $1\pm\eta b_\ell$ are positive.  Then $h_\pm$ are the Laplace
exponents of sums of independent Gamma variables, with the same
asymptotic coefficient $\alpha=N$, and their specified jets agree.
The laws differ: equality of their derivatives would give the rational
identity $\sum_\ell b_\ell c_\ell/(1+c_\ell u)=0$; its distinct poles
force all $b_\ell=0$.

Matching gradients give the same threshold points.  Strict concavity
makes the specified positive parameter points their unique optimizers.
The derivatives of an analytic Legendre transform through order $k$
are determined by the corresponding jet through order $k$ whenever
the Hessian is nonsingular.  This follows inductively by differentiating
$\nabla H(s(z))=z$, and proves the last assertion.
\end{proof}

The hierarchy is therefore strict: local curvature calibrates all
exponent ratios, but no finite list of finite-order derivatives determines
the microscopic law.  An exact open active patch supplies the missing
infinite information; the fluctuation law subsequently recovers the
discrete canonical spectrum.

\begin{remark}[Why the speed normalization is substantive]
If a one-coordinate law is infinitely divisible, duplicate every
spectral location $k$ times and replace that law by its $k$th
convolution root.  The complete joint law of every finite collection
of weighted sums is unchanged, since the $k$ variables at each location
sum to the original one-coordinate law.  For example,
$\operatorname{Gamma}(\alpha,1)$ becomes $k$ independent copies of
$\operatorname{Gamma}(\alpha/k,1)$ at each site.  The raw observations
cannot distinguish these descriptions.  The counting function,
however, is multiplied by $k$, the normalized threshold variables
are divided by $k$, and the continuum transform is divided by $k$.
Thus the normalization defining the limiting data in
Theorem~\ref{p06:thm:blind-two-scale} is substantive, rather than an
innocuous omission.  Without it, microscopic multiplicity and
convolution divisibility create an additional ambiguity.
\end{remark}

\begin{proposition}[The exact shape--multiplicity ambiguity]
\label{p06:prop:shape-multiplicity-ambiguity}
Let $a_n>0$, $\sum_n a_n^2<\infty$, and let
$Y=\sum_n a_n(G_n-\alpha)$ with unknown $\alpha>0$.
Its law determines the set of distinct coefficient values $x$ and
the positive numbers $w_x=\alpha m_x$, where $m_x$ is the multiplicity
of $x$.  There is a largest number $\alpha_*>0$ such that
$w_x/\alpha_*$ is an integer for every $x$.  All representations of
the given law within this class have exactly the shapes and
multiplicities
\[
 \alpha=\frac{\alpha_*}{k},\qquad
 m_x=k\frac{w_x}{\alpha_*},\qquad k=1,2,\ldots.
\]
In particular, imposing the normalization $\gcd_xm_x=1$ identifies
both $\alpha$ and the coefficient multiset.  A simple spectrum
automatically has this normalization.
\end{proposition}

\begin{proof}
Without knowing $\alpha$, the cumulants determine all moments of
the compactly supported finite measure
$\lambda=\alpha\sum_n a_n^2\delta_{a_n}$:
\[
 \int x^j\,d\lambda(x)=\frac{\kappa_{j+2}(Y)}{(j+1)!},
 \qquad j\ge0.
\]
Moment uniqueness recovers $\lambda$, and hence
$w_x=\lambda(\{x\})/x^2$ for every nonzero atom.
In any one representation put $g=\gcd_x m_x$ and
$\alpha_*=\alpha g$.  The gcd of an infinite family of positive
integers is the gcd of a finite subfamily, so B\'ezout's identity
applies to the integers $m_x/g$.  If $c>0$ satisfies
$w_x/c\in\mathbb N$ for all $x$, then
$(\alpha_*/c)(m_x/g)$ is an integer for every $x$.  A finite
B\'ezout combination proves $\alpha_*/c\in\mathbb N$.
Thus $c=\alpha_*/k$, and the stated multiplicities follow.
Conversely, splitting a Gamma variable into $k$ independent Gamma
variables of shape $1/k$ times the original shape preserves the
centered series law and gives every listed representation.
\end{proof}

Consequently a microscopic fluctuation law can itself calibrate the
counting granularity, provided primitive multiplicities are part of
the model.  For instance, suppose the measured continuum transform is
$\widetilde H=cH$ for an unknown constant $c>0$, with the Laplace
coordinates fixed.  Its derivative calibration still recovers the
power ratios, and Mellin inversion recovers $\widetilde h=ch$.
The coefficient $\widetilde\alpha=\lim_{u\to\infty}u\widetilde h'(u)$
is therefore known.  The primitive microscopic law gives $\alpha$,
and $c=\widetilde\alpha/\alpha$ identifies the missing amplitude.
The qualification about fixed Laplace coordinates matters: multiplying
a rate by a constant while keeping its threshold coordinates fixed
also rescales the dual Laplace coordinates.

\begin{corollary}[Calibration of an unknown rate amplitude]
\label{p06:cor:unknown-rate-amplitude}
Suppose that the observed rate is $J=cI$ on a nonempty open fully
active patch, with unknown $c>0$ and with threshold coordinates fixed.
Suppose also that the microscopic fluctuation law is observed as in
Theorem~\ref{p06:thm:blind-two-scale}, and restrict the model class to
spectra with primitive multiplicities.  These observations identify
$c$, $h$, all exponent ratios,
and the canonical spectrum $\{r_n^\rho\}$.
\end{corollary}

\begin{proof}
Multiplying $I$ by $c$ sends $s$ to $cs$ and $M$ to $M/c$,
so the calibration matrix $B$ is unchanged.  The exponent ratios
are therefore recovered directly from $J$.  The dual transform of
$J$ is
\[
 H_J(s)=cH(s/c)
 =\int_0^\infty\psi\left(\sum_j s_jx^{-p_j}\right)dx,
 \qquad \psi(u)=ch(u/c).
\]
Lemma~\ref{p06:lem:scalar-mellin-foundation} recovers
$\psi$; the inversion is linear in its second derivative and does
not require $e^{-\psi}$ to be a probability Laplace transform.
Its coefficient is
$\widetilde\alpha=\lim_{u\to\infty}u\psi'(u)=c\alpha$.
Proposition~\ref{p06:prop:shape-multiplicity-ambiguity} identifies
$\alpha$ from the primitive microscopic law, so that
$c=\widetilde\alpha/\alpha$ and
$h(u)=\psi(cu)/c$.  Spectral recovery then proceeds as before.
\end{proof}


\section{A finite-observation obstruction for microscopic spectra}
\label{p06:sec:finite-cumulant-obstruction}

The inverse theorem based on the full fluctuation law cannot in general
be replaced by a fixed finite list of cumulants.  The following elementary
construction makes this distinction explicit while leaving the entire
tail of the microscopic spectrum unchanged.  Its algebraic ingredients
are the classical power-sum identities; we do not claim a new general
theory of truncated moments.

\begin{proposition}[A flow invisible to finitely many cumulants]
\label{p06:prop:finite-cumulant-invisible-flow}
Fix an integer $K\ge2$, $\alpha>0$, and distinct positive numbers
$x_1>\cdots>x_K>0$.  There is a real-analytic curve
$x(\tau)=(x_1(\tau),\ldots,x_K(\tau))$, defined for sufficiently small
$|\tau|$, through $x(0)=x$, with all coordinates positive and distinct,
such that
\begin{align*}
 \sum_{i=1}^{K}x_i(\tau)^r&=\sum_{i=1}^{K}x_i^r,
       &&2\le r\le K,\\
 \sum_{i=1}^{K}x_i(\tau)^{K+1}
   &=\sum_{i=1}^{K}x_i^{K+1}+(K+1)\tau.
\end{align*}
Moreover $x_1'(0)>0$.  Append any fixed positive square-summable tail
$(c_n)$ with $\sup_n c_n<x_K$.  The centered Gamma series
\[
 Y_\tau=\sum_{i=1}^{K}x_i(\tau)(G_i-\alpha)
             +\sum_{n\ge1}c_n(\widetilde G_n-\alpha),
 \qquad G_i,\widetilde G_n\overset{\mathrm{iid}}{\sim}
             \operatorname{Gamma}(\alpha,1),
\]
have exactly the same cumulants through order $K$, although their largest
coefficients vary.  Their next cumulant satisfies
\[
 \kappa_{K+1}(Y_\tau)-\kappa_{K+1}(Y_0)
        =\alpha(K+1)!\tau.
\]
Any asymptotic counting law depending only on the common
coefficient-sequence tail is unchanged.  This assertion does not concern
the probability tails of $Y_\tau$: its right-tail exponential rate
changes when its largest coefficient changes.
\end{proposition}

\begin{proof}
For $P_\tau(z)=\prod_{j=1}^{K}(z-x_j(\tau))$, solve the analytic ODE
\[
 x_i'(\tau)=\frac{1}{x_i(\tau)P_\tau'(x_i(\tau))}.
\]
Its vector field is analytic on the open set of distinct nonzero
coordinates, so the asserted local solution exists and remains ordered
and positive.  Lagrange interpolation gives
\[
 \sum_{i=1}^{K}\frac{x_i^m}{P'(x_i)}=0\quad(0\le m\le K-2),
 \qquad
 \sum_{i=1}^{K}\frac{x_i^{K-1}}{P'(x_i)}=1.
\]
Therefore the derivative of the $r$th power sum is zero for
$2\le r\le K$, and equals $K+1$ for $r=K+1$.
Since $x_1>\cdots>x_K>0$, every factor of $P'(x_1)$ is positive,
and hence $x_1'(0)>0$.  Positivity of the separation from the appended
tail persists after shortening the parameter interval.

The series converge in $L^2$, and their cumulants are
$\alpha(r-1)!$ times the corresponding power sums for $r\ge2$.
This proves all assertions.  In fact, using the same Gamma variables
for each parameter gives the quantitative continuity estimate
\[
 W_2(\mathcal L(Y_\tau),\mathcal L(Y_0))
 \le\sqrt{\alpha}\,\|x(\tau)-x(0)\|_2=O(|\tau|).
\]
\end{proof}

\begin{remark}[An exact polynomial evolution]
The construction can equivalently be written
\[
 \partial_\tau P_\tau(z)
      =\frac{P_\tau(z)/P_\tau(0)-1}{z}.
\]
The apparent singularity at $z=0$ is removable. Partial fractions give
\[
 \sum_i\frac{1}{x_iP'(x_i)}=-\frac{1}{P(0)}.
\]
This supplies a direct finite
algebraic description of the invisible deformation.
\end{remark}

\begin{remark}[Scope]
The proposition proves noninjectivity for arbitrary fixed finite
cumulant data within the positive square-summable Gamma-series class.
It does not contradict finite-cumulant convergence criteria when the
target spectrum has fixed finite rank, nor reconstruction under a
specified finite-dimensional model.  Both are substantial classical
subjects.  It also does not say that every prescribed moment vector
has such a deformation: boundary moment vectors can be rigid.
\end{remark}

\section{The same arithmetic fluctuations under a hard energy constraint}
\label{p05:sec:hard-prime-constraint}

For this section let $V(\theta)=2(1-\cos\theta)$ with $\theta$
uniform on $[-\pi,\pi]$, and let $h$ and $g=-h''$ be as above.
The law of $V$ has density
$[\pi\sqrt{v(4-v)}]^{-1}$ on $(0,4)$, and the one-site parameter
in \eqref{p05:eq:microscopic-one-site-hypotheses} is $\alpha=1/2$.
Fix $A>1$ and write
\[
 D_A=\sum_p p^{-A}V_p,\qquad
 F(t)=\sum_p h(tp^{-A}),\qquad
 \delta_t=F'(t)=\mathbb E_tD_A.
\]

\begin{theorem}[Hard conditioning and an independent boundary slack]
\label{p05:thm:hard-prime-gamma-limit}
For every real $a=(a_p)_p\in\ell^2$ set
\[
 X_a=\sum_p a_pp^{-A}V_p,\qquad
 m_t=\mathbb E_tX_a.
\]
Under the original product law conditioned on $D_A\le\delta_t$,
\begin{equation}\label{p05:eq:hard-prime-joint-limit}
 \big(t(X_a-m_t),\ t(\delta_t-D_A)\big)
 \ \Longrightarrow\ (Y_a,R),
\end{equation}
where $Y_a=\sum_pa_p(G_p-1/2)$,
$G_p\overset{\rm iid}{\sim}\operatorname{Gamma}(1/2,1)$,
and $R\sim\operatorname{Exp}(1)$ is independent of all the $G_p$.
Consequently, for $E-A>1/2$ and
$D_E=\sum_p(\log p)p^{-E}V_p$, one may take
$a_p=(\log p)p^{A-E}$.
The centering is the exact tilted mean; replacing it by a leading
prime-number-theorem equivalent is not justified at this fluctuation
scale.
\end{theorem}

\begin{proof}
We supply the local argument that is missing from a purely logarithmic
Gibbs-conditioning calculation.  Put $x=t^{1/A}$ and $m=x/\log x$.
Under $\mathbb P_t$ let
\[
 Z_t=t(D_A-\delta_t),\qquad
 \sigma_t^2=\operatorname{Var}_t(Z_t)
 =\sum_p(tp^{-A})^2g(tp^{-A}).
\]
Prime-number-theorem summation gives
\begin{equation}\label{p05:eq:hard-prime-variance}
 \sigma_t^2\sim v_A m,
 \qquad v_A=\int_0^\infty y^{-2A}g(y^{-A})\,dy\in(0,\infty).
\end{equation}
Indeed the integrand is bounded near zero and $O(y^{-2A})$ at
infinity.  The same assertion, with half the constant, holds if the
primes are restricted to either the odd or the even indices in the
increasing sequence of primes.  Their counting functions differ from
half the prime counting function by at most one.  Deleting any fixed
finite set does not change these asymptotics.

Here is a uniform local central limit theorem for these sums.  For
$u>0$ let $W_u=uV$ under its one-coordinate tilt.  Direct integration
against the displayed arcsine density gives
\[
 \mathbb E_u|W_u-\mathbb E_uW_u|^3
 \le C\min(u^3,1).
\]
For $u\ge1$, the estimate follows by splitting $0<v<4$ at $v=2$:
the first part is bounded by a Gamma integral after $w=uv$, and the
second is exponentially small despite its integrable endpoint
singularity.  For $u\le1$, use $0\le W_u\le4u$.
The sum of these third moments is $O(m)$, so the normalized sums
satisfy the central limit theorem by the Lyapunov condition.

For primes in $[x,2x]$, the parameter $u$ lies in the compact
interval $[2^{-A},1]$.  Uniformly on this interval, the characteristic
function $\psi_u(\xi)=\mathbb E_u e^{i\xi W_u}$ satisfies
\[
 |\psi_u(\xi)|\le e^{-c\xi^2}\quad(|\xi|\le\eta),
 \qquad
 \sup_{\eta\le|\xi|\le B}|\psi_u(\xi)|<1,
 \qquad
 |\psi_u(\xi)|\le C|\xi|^{-1/2}\quad(|\xi|\ge1).
\]
The first bound follows from the uniformly positive variances, the
second from the continuous nonlattice densities and compactness,
and the last by splitting the arcsine integral at its endpoints
and one integration by parts on the intervening interval.
There are $\asymp m$ such primes in each of the two index classes.
Multiplication of their characteristic functions, followed by Fourier
inversion, upgrades the central limit theorem to
\begin{equation}\label{p05:eq:hard-prime-local-limit}
 \sup_{z\in\mathbb R}
 \left|\sigma_t f_t(\sigma_t z)
       -(2\pi)^{-1/2}e^{-z^2/2}\right|\longrightarrow0,
\end{equation}
where $f_t$ is the density of $Z_t$.
For completeness, the Fourier integral splits into
$|\xi|\le\eta$, $\eta<|\xi|\le B$, and $|\xi|>B$.
The first normalized integrand is dominated by $e^{-c'v^2}$;
the second contributes $O(\sigma_t\rho^{cm})$ with $\rho<1$;
choosing $B>C^2$ makes the integral of
$(C|\xi|^{-1/2})^{cm}$ on the last region exponentially small.
This proves uniform Fourier-integral convergence.
The same proof applies after deleting finitely many primes and to
either index class.  In particular all these reservoir densities
have supremum $O(m^{-1/2})$.

Changing measure gives
\[
 \mathbb P(D_A\le\delta_t)
 =e^{t\delta_t-F(t)}I_t,
 \qquad I_t=\int_{-\infty}^0e^z f_t(z)\,dz
 \sim(2\pi)^{-1/2}\sigma_t^{-1}.
\]
Let $I$ be a fixed finite set of primes, and center its tilted energy
as $Z_I=\sum_{p\in I}(tp^{-A}V_p-tp^{-A}h'(tp^{-A}))$.
Let $f_{I^c,t}$ be the density of the independently centered
complementary energy.  The joint conditional law of the coordinates
in $I$ and the slack $r=t(\delta_t-D_A)$ has density, relative to
$d\mathbb P_{t,I}\,dr$, equal to
\begin{equation}\label{p05:eq:hard-prime-finite-marginal}
 \mathbf1_{r\ge0}\,
 \frac{e^{-r}f_{I^c,t}(-Z_I-r)}{I_t}.
\end{equation}
By \eqref{p05:eq:hard-prime-local-limit}, the density ratio tends to
$e^{-r}$ uniformly for bounded $Z_I,r$, and is bounded by $Ce^{-r}$
everywhere.  Since the finite coordinates $tp^{-A}V_p$ converge to
independent Gamma$(1/2,1)$ variables under $\mathbb P_t$, this proves
their joint conditional convergence with an independent exponential
slack.

It remains to pass from finitely many coordinates to $X_a$.
Whenever a set of observed primes leaves one whole index class
unobserved, integrating \eqref{p05:eq:hard-prime-finite-marginal} in $r$
(the formula also holds for an infinite observed subseries) bounds
its conditional marginal density relative to its tilted marginal
by a fixed constant.  This follows from the reservoir supremum
$O(m^{-1/2})$ and $I_t\asymp m^{-1/2}$; convolution with any extra
coordinates preserves that supremum bound.
Split the primes exceeding a fixed cutoff $Q$ into their two index
classes.  Each part has the other class as a reservoir.  Therefore
independence under the tilted law and the bound
$\sup_u u^2g(u)<\infty$ give
\[
 \limsup_{t\to\infty}
 \mathbb E\left[
 \left|\sum_{p>Q}a_p\{tp^{-A}V_p-tp^{-A}h'(tp^{-A})\}\right|^2
 \mathrel{\Big|}D_A\le\delta_t\right]
 \le C\sum_{p>Q}a_p^2.
\]
The right side tends to zero.  The finite-coordinate joint limit
therefore proves \eqref{p05:eq:hard-prime-joint-limit}.
\end{proof}

\begin{remark}[A limitation for zeta-zero sampling]
The hard-conditioning theorem concerns product Haar measure.
A comparison $c\mu\le\nu\le\mu$ transfers concentration at a
deterministic limit but does not determine the distribution of
$Y_a$.  For example, weighting by a bounded nonconstant function
of the fluctuation changes its limiting law while preserving such
a comparison.  Thus this theorem alone gives no Gamma-series law
for the full set of simple zeta zeros.
\end{remark}


\section{Activation universality and the spatial response profile}
\label{p05:sec:activation}

We retain the assumptions and normalization of Section~\ref{p05:sec:rv}.
Write $d=E-A$. The change from one active constraint to two is governed
by $d/\rho$, rather than by the complete microscopic law. Its spatial
mechanism changes at $d=\rho/2$.

Fix $a>0$. Let $t_0>0$ be the unique solution of
$H_t(0,t_0)=a$, and put
\[
 I_A(a)=H(0,t_0)-t_0a.
\]
If $d<\rho$, define $e_0=H_s(0,t_0)<\infty$. Then
$I(e,a)=I_A(a)$ for $e\ge e_0$: the $E$-constraint is inactive.
For $e<e_0$ sufficiently close to $e_0$, both multipliers are positive.

\begin{theorem}[Three activation regimes]\label{p05:thm:activation}
Suppose $0<d<\rho$ and let $\varepsilon\downarrow0$.
\begin{enumerate}
\item If $d<\rho/2$, set
\[
 M_{bc}=\rho\int_0^\infty y^{\rho-1-b-c}
 g(t_0y^{-A})\,dy\quad(b,c\in\{E,A\}),\qquad
 \kappa=M_{EE}-\frac{M_{EA}^2}{M_{AA}}>0.
\]
Then
\[
 I(e_0-\varepsilon,a)-I_A(a)
 \sim\frac{\varepsilon^2}{2\kappa}.
\]
\item If $d=\rho/2$, then
\[
 I(e_0-\varepsilon,a)-I_A(a)
 \sim\frac{t_0^2}{4\alpha}
       \frac{\varepsilon^2}{\log(1/\varepsilon)}.
\]
\item If $\rho/2<d<\rho$, put
\[
 q=\frac{\rho}{d}\in(1,2),\qquad
 K=\frac{\alpha\pi}{\sin(\pi(q-1))}.
\]
Then
\begin{equation}\label{p05:eq:activation-fractional}
 I(e_0-\varepsilon,a)-I_A(a)
 \sim (q-1)q^{-q/(q-1)}
       K^{-1/(q-1)}t_0^{q/(q-1)}
       \varepsilon^{q/(q-1)}.
\end{equation}
\end{enumerate}
If $d\ge\rho$, then $H_s(0,t)=\infty$ for every $t>0$;
there is no $A$-only optimizing face at finite $(e,a)$.
\end{theorem}

\begin{proof}
The last assertion follows from
$y^{\rho-1-E}h'(ty^{-A})\sim(\alpha/t)y^{\rho-1-d}$
at zero. Assume henceforth $d<\rho$.
The mixed derivative $H_{st}(0,t)$ and $H_{tt}(0,t)$ are finite,
continuous from the positive quadrant, and $H_{tt}(0,t)<0$.
The implicit function theorem applied to $H_t(s,t_s)=a$ therefore
gives $t_s=t_0+O(s)$ for $s\downarrow0$.
For clarity, the theorem can be applied after extending $H_t$ across
$s=0$ by a $C^1$ extension; equivalently, its usual one-sided proof
uses the mean value theorem and $H_{tt}(0,t_0)\ne0$.
Define globally
\[
 F(s)=\sup_{t\ge0}\{H(s,t)-at\},\qquad
 D(s)=F(0)+e_0s-F(s).
\]
Locally, $F(s)=H(s,t_s)-at_s$.
Strict concavity shows that the maximizing pair for $e\uparrow e_0$
converges to $(0,t_0)$, and consequently
\begin{equation}\label{p05:eq:activation-dual}
 I(e_0-\varepsilon,a)-I_A(a)
 =\sup_{s\ge0}\{\varepsilon s-D(s)\}
\end{equation}
with the maximizer tending to zero.

When $2d<\rho$, the full boundary Hessian is finite. Differentiation
of $H_t(s,t_s)=a$ gives $D(s)=\kappa s^2/2+o(s^2)$.
The strict Cauchy--Schwarz inequality gives $\kappa>0$, since
$y^{-E}$ and $y^{-A}$ are not proportional. This proves the first case.

For the other two cases it suffices first to keep $t=t_0$ fixed:
profiling $t$ changes $H(s,t)-at$ by $O(s^2)$. Indeed its first
$t$-derivative at $t_0$ is $O(s)$, its second derivative remains
bounded, and $t_s-t_0=O(s)$.
Put $y_s=(s/t)^{1/d}$. The bounds
\begin{equation}\label{p05:eq:h-global-bounds}
 0<h'(u)\le\frac{C}{1+u},\qquad
 0<g(u)\le\frac{C}{(1+u)^2}
\end{equation}
follow from \eqref{p05:eq:laplace-assumptions} and continuity on compact
intervals. In the case $1<q<2$, substitution $y=y_sz$ and dominated
convergence give
\begin{align*}
 &H(s,t)-H(0,t)-sH_s(0,t)\\
 &\quad\sim \rho\alpha(s/t)^q
 \int_0^\infty z^{\rho-1}
 \{\log(1+z^{-d})-z^{-d}\}\,dz
 =-K(s/t)^q.
\end{align*}
The domination is $Cz^{\rho-1}\min(z^{-d},z^{-2d})$, integrable
at both endpoints precisely for $d<\rho<2d$.
The integral is evaluated by setting $v=z^{-d}$ and integrating by
parts:
\[
 q\int_0^\infty(v-\log(1+v))v^{-q-1}\,dv
 =\int_0^\infty\frac{v^{1-q}}{1+v}\,dv
 =\frac{\pi}{\sin(\pi(q-1))}.
\]
Thus $D(s)\sim Kt_0^{-q}s^q$.

At $2d=\rho$, split the integral defining $-H_{ss}(s,t)$ at
$y_s$ and at a small fixed number. On the middle range its integrand
is asymptotic to $(\rho\alpha/t^2)y^{-1}$, uniformly away from
the two ends. The lower range contributes $O(1)$ and the upper range
also contributes $O(1)$ after the fixed splitting point is chosen.
More explicitly, \eqref{p05:eq:h-global-bounds} bounds the lower range
by $Cs^{-2}y_s^\rho=O(1)$; the middle logarithmic coefficient follows
by making $ty^{-A}$ large and $s y^{-E}/(ty^{-A})$ small.
It follows that
\[
 -H_{ss}(s,t)\sim\frac{2\alpha}{t^2}\log(1/s),\qquad
 D(s)\sim\frac{\alpha}{t_0^2}s^2\log(1/s).
\]
Integrating the second derivative twice is justified by its
logarithmic bound. Finally optimize the three displayed asymptotics
in \eqref{p05:eq:activation-dual}. Quadratic and power functions give
the first and third constants. In the critical case the maximizing
scale is
$s\sim t_0^2\varepsilon/(2\alpha\log(1/\varepsilon))$,
giving the stated constant. These optimizations can alternatively
be obtained by upper and lower multiplicative bounds on $D$;
no differentiation of an unspecified remainder is needed.
\end{proof}

The preceding exponent has a spatial explanation. The following
statement concerns a deterministic measure of the change in tilted
mean energy; it is not a claim that a randomly chosen coordinate
has the stated distribution.
For fixed $t>0$ and $0<d<\rho$, define the positive measure
\[
 dR_{s,t}(y)=\rho y^{\rho-1-E}
 \{h'(ty^{-A})-h'(sy^{-E}+ty^{-A})\}\,dy.
\]
Its mass is the reduction of the $E$-energy mean when the multiplier
$s$ is activated at fixed $t$.

\begin{theorem}[Universal activation profile]\label{p05:thm:activation-profile}
Suppose $\rho/2<d<\rho$, and normalize $R_{s,t}$ to mass one.
Under the change of variable $W=(y/y_s)^d$, where
$y_s=(s/t)^{1/d}$, these measures converge weakly as $s\downarrow0$
to the beta-prime law
\begin{equation}\label{p05:eq:beta-front}
 \frac{\sin(\pi(q-1))}{\pi}\,
 \frac{w^{q-2}}{1+w}\,dw,\qquad w>0,
 \quad q=\rho/d.
\end{equation}
At $d=\rho/2$, the image of the normalized measure under
\[
 y\longmapsto\frac{\log(y/y_s)}{\log(1/y_s)}
\]
converges weakly on $\mathbb R$ to the uniform law on $[0,1]$.
For $d<\rho/2$, there is instead an unscaled limit:
$s^{-1}dR_{s,t}(y)$ converges in total mass to
$\rho y^{\rho-1-2E}g(ty^{-A})\,dy$.
\end{theorem}

\begin{proof}
For fixed $z>0$, after $y=y_sz$, the density of $R_{s,t}$ divided by
$(\rho\alpha/t)y_s^{\rho-d}$ tends to
\[
 \frac{z^{\rho-d-1}}{1+z^d}.
\]
The bounds \eqref{p05:eq:h-global-bounds} dominate it by
$Cz^{\rho-d-1}\min(1,z^{-d})$. Hence convergence holds in
$L^1$ when $d<\rho<2d$. Its integral is
$\pi/[d\sin(\pi(q-1))]$; setting $w=z^d$ proves
\eqref{p05:eq:beta-front}.
At $2d=\rho$, the density on $y_s\ll y\ll1$ is
asymptotic to $(\rho\alpha s/t^2)dy/y$.
Each fixed closed subinterval of the logarithmic coordinate
inside $(0,1)$ therefore receives asymptotically its length times
the total mass. The two endpoint ranges are controlled by
\eqref{p05:eq:h-global-bounds}: below $y_s$ and above a fixed positive
constant their masses are $O(s)$, while the full mass is
asymptotic to $(2\alpha/t^2)s\log(1/s)$.
Letting the closed subintervals increase to $(0,1)$ proves the
critical assertion. If $2d<\rho$, ordinary dominated convergence
applied to the difference quotient proves the last assertion.
\end{proof}

The fractional response profile also describes the leading response
along the constrained path $t=t_s$. Define the signed measure
\[
 d\widetilde R_s(y)=\rho y^{\rho-1-E}
 \{h'(t_0y^{-A})-h'(sy^{-E}+t_sy^{-A})\}\,dy.
\]
Its difference from $R_{s,t_0}$ has total variation $O(s)$,
because $|t_s-t_0|=O(s)$ and $H_{st}$ is locally bounded.
This is negligible compared with $s^{q-1}$ in the fractional regime,
or with $s\log(1/s)$ in the critical case. The signed measure
therefore has the same normalized leading limit as the positive
response measure above.

There is a further local interpretation. Let $Q_u$ be the tilt of
the law of $V$ with density $e^{-uV+h(u)}$. At a fixed rescaled
location $y=y_sz$, put $u_0=ty^{-A}$. Under
$Q_{u_0(1+z^{-d})}$, the random variable $u_0V$ tends to a Gamma
law with shape $\alpha$ and rate $1+z^{-d}$. Indeed, for $\lambda\ge0$
its Laplace transform tends to
\[
 \exp\{h(u_0(1+z^{-d}))-h(u_0(1+z^{-d}+\lambda))\}
 \longrightarrow
 \left(\frac{1+z^{-d}}{1+z^{-d}+\lambda}\right)^\alpha,
\]
by integrating $h'(u)\sim\alpha/u$.
Thus the singular regime has a universal moving response profile,
the critical regime spreads across logarithmic scales, and the
regular regime retains a finite quadratic information metric.
These conclusions use the continuum limit first and activation
second; they do not assert a uniform two-parameter limit for an
arbitrary regularly varying spectrum.


\section{One-prime and collective tails of the arithmetic fluctuation}
\label{p05:sec:prime-fluctuation-tails}

Let $1/2<d<1$ and define
\[
 a_p=(\log p)p^{-d},\qquad
 Y_d=\sum_pa_p(G_p-1/2),\qquad
 G_p\overset{\rm iid}{\sim}\operatorname{Gamma}(1/2,1).
\]
This is the hard-conditioned fluctuation in
Theorem~\ref{p05:thm:hard-prime-gamma-limit} when $d=E-A$.
Write
\[
 K_d=\frac{\pi}{2\sin(\pi(1/d-1))},\qquad
 B_d=(1-d)^2\left(\frac d{K_d}\right)^{d/(1-d)},
 \qquad a_* =\max_p a_p.
\]

\begin{theorem}[Asymmetric arithmetic tails]
\label{p05:thm:arithmetic-fluctuation-tails}
As $z\to\infty$,
\begin{equation}\label{p05:eq:prime-gamma-left-tail}
 -\log\mathbb P(Y_d\le-z)
 \sim B_d\frac{z^{1/(1-d)}}{\log z},
\end{equation}
whereas
\begin{equation}\label{p05:eq:prime-gamma-right-tail}
 -\log\mathbb P(Y_d\ge z)\sim z/a_*.
\end{equation}
If $a_*$ is attained at a unique prime $p_*$, the sharper right-tail
equivalent is
\[
 \mathbb P(Y_d\ge z)
 \sim\sqrt{\frac{a_*}{\pi}}e^{-1/2}
 \prod_{p\ne p_*}
 \left[e^{-a_p/(2a_*)}(1-a_p/a_*)^{-1/2}\right]
 z^{-1/2}e^{-z/a_*}.
\]
For $d=2.813014020252898-1.940101683743625$, the unique maximizing
prime is $p_*=3$, and the exponent in
\eqref{p05:eq:prime-gamma-left-tail} is approximately $7.86858435$.
The numerical values only illustrate the exact formulas.
\end{theorem}

\begin{proof}
For $s\ge0$, independence gives the everywhere-finite transform
\[
 L(s)=\log\mathbb E e^{-sY_d}
 =\frac12\sum_p\{sa_p-\log(1+sa_p)\}.
\]
Let $x=x(s)$ be the large solution of $s(\log x)x^{-d}=1$.
The prime number theorem and the change of variable $p=xy$ yield
\[
 L(s)\sim\frac{x}{\log x}\frac12
 \int_0^\infty\{y^{-d}-\log(1+y^{-d})\}\,dy
 =K_d\frac{x}{\log x}.
\]
The improper passage is valid because the limit superior, as
$x\to\infty$ with fixed $\eta,R$, of the normalized small-prime
tail is $O(\eta^{1-d})$ for $p<\eta x$, and that of the large-prime
tail is $O(R^{1-2d})$ for $p>Rx$. These bounds follow respectively
from $u-\log(1+u)\le u$ and
$u-\log(1+u)\le u^2/2$, followed by partial summation.
To evaluate the integral substitute $u=y^{-d}$ and integrate by
parts; the result is
$\frac12\int_0^\infty u^{1-1/d}/(1+u)\,du=K_d$.
Since
$x\sim d^{-1/d}s^{1/d}(\log s)^{1/d}$, this gives
\begin{equation}\label{p05:eq:prime-gamma-loglaplace}
 L(s)\sim K_dd^{1-1/d}s^{1/d}(\log s)^{1/d-1}.
\end{equation}

For clarity we spell out the exponential-Tauberian step.
Convexity and regular variation in
\eqref{p05:eq:prime-gamma-loglaplace} imply
$L'(s)\sim L(s)/(ds)$.  This also follows by direct differentiation
of the prime sum.  Moreover
\[
 0<L''(s)=\frac12\sum_p\frac{a_p^2}{(1+sa_p)^2}
 \le L'(s)/s.
\]
For $Z=-Y_d$, the exponential tilt by $e^{sZ-L(s)}$ has mean
$L'(s)$ and variance $L''(s)$.  At the saddle $L'(s)=z$, the
relative variance tends to zero.  Exponential Markov gives the
upper probability bound.  For the reverse bound, choose a saddle
with $L'(s)=(1+\eta)z$ and restrict $Z$ to
$[z,(1+2\eta)z]$.  Its tilted probability tends to one by the
variance estimate.  Changing measure and then letting
$\eta\downarrow0$ proves
$-\log\mathbb P(Z\ge z)\sim\sup_{s\ge0}\{sz-L(s)\}$.
Legendre optimization of \eqref{p05:eq:prime-gamma-loglaplace} gives
exactly the constant $B_d$ in
\eqref{p05:eq:prime-gamma-left-tail}.

The moment generating function of $Y_d$ is finite for
$0\le q<1/a_*$, yielding
$\liminf z^{-1}[-\log\mathbb P(Y_d\ge z)]\ge1/a_*$.
Keeping one maximizing Gamma variable and restricting the
independent remaining sum to a fixed interval of positive
probability gives the reverse inequality.
If the maximum is unique, write
$Y_d=a_*(G_{p_*}-1/2)+R$.
Since $a_p\to0$, the second largest coefficient is strictly smaller
than $a_*$.  Thus $R$ has a finite moment generating function in a
neighborhood of $1/a_*$.  Integrating the Gamma tail equivalent
$\mathbb P(G\ge u)\sim\pi^{-1/2}u^{-1/2}e^{-u}$ against $R$
gives the stated constant; the extra exponential moment controls
the complement of bounded $R/z$.
Finally, $u\mapsto(\log u)u^{-d}$ increases up to $u=e^{1/d}$
and decreases thereafter.  At the displayed value of $d$,
$3<e^{1/d}<5$, and direct comparison of $a_3$ and $a_5$ selects
$p_*=3$.
\end{proof}

The left-tail exponent is the same singular exponent that occurs
when the second constraint activates at the $A$-only edge.  Here
it has a different, microscopic interpretation: downward fluctuation
requires coordinated suppression of infinitely many centered Gamma
terms, whereas the far upward fluctuation is set by a single
coefficient.  The use of exponential tilting and exponential
Tauberian reasoning is classical; the assertion identifies their
concrete arithmetic consequences in this prime-indexed limit law.

\section{Correlated local laws: exact aliases and one-channel restoration}
\label{p06:sec:correlated-identification}

The preceding sections use a common scalar local law. We now change
the model explicitly: let $X$ be a nonnegative random vector, set
$h(u)=-\log\mathbb E e^{-u\cdot X}$, and identify a nonnegative linear
functional with its coefficient vector. Channels $\ell_i$ with powers
$F_i>\rho>0$ have response
\begin{equation}\label{p06:eq:intro-response}
 H(s)=\rho\int_0^\infty y^{\rho-1}
 h\left(\sum_i s_i y^{-F_i}\ell_i\right)\,dy.
\end{equation}
The coordinates of $X$ need not be independent. Dilation can now be
compensated by convolution, so even the entire joint response can
lose the local law. A repeated component or a positive mixed component
at a second power can remove this ambiguity under the respective
hypotheses below. These classes must be kept distinct: a gated law
with an atom at zero does not satisfy the positive Gamma-shape
asymptotic used in the scalar calibration theorem.

Fix $d\ge2$, $\rho>0$, and $E_1,\ldots,E_d>\rho$.
For a probability measure $\mu$ on $\mathbb R_+^d$ write
\[
 h_\mu(v)=-\log\int e^{-v\cdot x}\,\mu(dx),\qquad
 \delta_c v=(c^{E_1}v_1,\ldots,c^{E_d}v_d).
\]
We use the following class, which includes arbitrary continuous distributions
away from the origin as well as finite discrete examples:
\begin{equation}\label{p06:eq:origin-gap-class}
 \mathcal C_d=\left\{p_0\delta_0+(1-p_0)\nu:
 0<p_0<1,\quad\operatorname{supp}\nu\subset[a,b]^d
 \text{ for some }0<a<b<\infty\right\}.
\end{equation}
The constants $p_0,a,b$ need not be known and can differ between laws.
For $\mu\in\mathcal C_d$ its homogeneous response is
\begin{equation}\label{p06:eq:correlated-H}
 \mathcal H_\mu(s)=\rho\int_0^\infty
 y^{\rho-1}h_\mu(s_1y^{-E_1},\ldots,s_dy^{-E_d})\,dy.
\end{equation}
This integral is finite: $h_\mu\le-\log p_0$ and
$h_\mu(v)\le\mathbb E_\mu X\cdot v$.
It is concave and differentiable on the positive orthant, and
\[
 \mathcal H_\mu(\delta_c s)=c^\rho\mathcal H_\mu(s).
\]
The term ``response'' here refers to this limiting log-Laplace object.
It does not assert that $\exp(-\mathcal H_\mu)$ is a Laplace transform
of an infinitely divisible law for every $\mu$.

\begin{theorem}[Dilation-convolution aliases]\label{p06:thm:correlated-alias}
Let $X^{(1)},\ldots,X^{(n)}$ be independent copies with law
$\mu\in\mathcal C_d$, and let $c_1,\ldots,c_n>0$. Define
\[
 Y_j=\sum_{i=1}^n c_i^{E_j}X_j^{(i)},\qquad 1\le j\le d.
\]
Then the law $\mu_c$ of $Y$ belongs to $\mathcal C_d$, and
\begin{equation}\label{p06:eq:alias-scaling}
 \mathcal H_{\mu_c}(s)=
 \left(\sum_{i=1}^n c_i^\rho\right)\mathcal H_\mu(s).
\end{equation}
Consequently the full function $\mathcal H_\mu$ does not determine
$\mu$ in $\mathcal C_d$.
\end{theorem}
\begin{proof}
Independence gives
$h_{\mu_c}(v)=\sum_i h_\mu(c_i^{E_1}v_1,\ldots,c_i^{E_d}v_d)$.
In the $i$th integral in \eqref{p06:eq:correlated-H}, substitute $y=c_i r$.
This proves \eqref{p06:eq:alias-scaling}. The atom at zero has mass $p_0^n$,
and every nonzero outcome has all coordinates positive and bounded away
from zero, uniformly on the support. Thus $\mu_c\in\mathcal C_d$.
Taking $n=2$ and $c_1=c_2=2^{-1/\rho}$ preserves $\mathcal H$ but changes
$\mathbb E Y_j$ by the factor $2^{1-E_j/\rho}<1$.
\end{proof}

\begin{proposition}[Aliases survive every finite calibration]
\label{p06:prop:finite-calibration-alias}
Fix any integer $M\ge1$. There exist distinct $\mu,\widetilde\mu\in
\mathcal C_d$ with the same homogeneous response and the same mixed
moments of total degree at most $M$. They may also be required to have
identical masses at zero and identical coordinatewise upper support
bounds. In addition, any prescribed finite collection of one-component
responses
\[
 \mathcal R_{j,F}(s)=\rho\int_0^\infty y^{\rho-1}
 h_\mu(0,\ldots,sy^{-F},\ldots,0)\,dy,
 \qquad F>\rho,
\]
can be made identical between the two laws, for all $s>0$.
\end{proposition}
\begin{proof}
Let $B$ be Bernoulli with a transcendental success probability in $(0,1)$,
and start with $X=(B,\ldots,B)$. Every cumulant $\kappa_m(B)$ is a nonzero
polynomial in its success probability with integer coefficients: its
coefficient of the linear term is one. Hence every $\kappa_m(B)$ is nonzero.
For the dilation-convolution law above, the joint cumulant of multi-order
$k=(k_1,\ldots,k_d)\ne0$ is
\begin{equation}\label{p06:eq:cumulant-power-sums}
 \kappa_k(Y)=\kappa_{|k|}(B)
 \sum_i c_i^{k\cdot E}.
\end{equation}
Thus equality of power sums for the finite set
\[
 P=\{\rho\}\cup\{k\cdot E:1\le |k|\le M\}
\]
is enough. To preserve the additional scalar responses, also include
$E_j\rho/F$ for every specified pair $(j,F)$: direct substitution gives
$\mathcal R_{j,F}^{\mu_c}=(\sum_i c_i^{E_j\rho/F})\mathcal R_{j,F}^{\mu}$.
Remove repeated elements of $P$, and put $q=|P|$.

Choose $q+1$ distinct positive numbers $c_i$. The Jacobian of
$c\mapsto(\sum_i c_i^p)_{p\in P}$ has rank $q$.
Indeed, after harmless nonzero row factors, its entries are $c_i^{p-1}$;
distinct real powers form a Chebyshev system on $(0,\infty)$.
For completeness, after $c_i=e^{u_i}$ this assertion follows by induction
from Rolle's theorem: divide a nonzero linear combination of exponentials
by its smallest exponential, differentiate, and count its zeros.
The implicit function theorem therefore gives a nonconstant smooth curve
of scale vectors along which all these power sums are constant.

Choose an integer $m$ with $mE_1>\max P$. The augmented Jacobian obtained
by adding $\sum_i c_i^{mE_1}$ has full rank $q+1$, by the same argument.
Consequently this extra power sum has nonzero derivative along the local
one-dimensional level set. Two sufficiently close points on that level
set give different $m$th cumulants of $Y_1$, by
\eqref{p06:eq:cumulant-power-sums}, and thus distinct probability laws.
All required lower moments match because moments are polynomials in
cumulants. The number $q+1$ of Bernoulli summands is fixed, so the atom at
zero matches. The coordinatewise maxima are $\sum_i c_i^{E_j}$ and also
match, since $M\ge1$.
\end{proof}

\begin{remark}[A concrete alias]
For $d=2$, $\rho=1$, $(E_1,E_2)=(3,2)$, and any nondegenerate Bernoulli
$B$, the laws
\[
 (B,B)\quad\text{and}\quad
 ((B_1+B_2)/8,(B_1+B_2)/4)
\]
have identical complete two-component responses. A new channel observing
the first component with exponent $C=4$ distinguishes this particular
pair even through its marginal response: the second response is
$2^{1/4}$ times the first. The preceding proposition shows why finitely
many separate marginal measurements do not suffice in general.
\end{remark}

\subsection{One redundant channel restores the entire correlated law}

Choose $C>\rho$ with $C\ne E_1$. We observe the first component once more,
with this different exponent, and define
\begin{equation}\label{p06:eq:augmented-correlated-H}
 \mathcal K_\mu(s_1,\ldots,s_d,u)
 =\rho\int_0^\infty y^{\rho-1}
 h_\mu(s_1y^{-E_1}+uy^{-C},s_2y^{-E_2},\ldots,s_dy^{-E_d})\,dy.
\end{equation}
The new channel measures no new local random variable. Its information
comes from applying a different spatial weight to the same component.

\begin{theorem}[Identification by a repeated component]
\label{p06:thm:one-channel-restoration}
For $d\ge2$, the map $\mu\mapsto\mathcal K_\mu$ is injective on
$\mathcal C_d$. In particular, $d$ coordinate channels fail to identify
arbitrary correlated local laws even after any finite moment calibration,
whereas $d+1$ channels in \eqref{p06:eq:augmented-correlated-H} identify them.
This is a threshold for this specified observation architecture, not a
claim about all possible measurement schemes.
\end{theorem}
\begin{proof}
Let
\[
 f_\mu(v)=-\partial_{v_1}^2h_\mu(v)
 =\operatorname{Var}_{\mu,v}(X_1),
\]
where $\mu,v$ denotes exponential tilting by $e^{-v\cdot X}$.
The support and atom assumptions imply
\begin{equation}\label{p06:eq:tilted-variance-bound}
 0\le f_\mu(v)\le \frac{b^2}{p_0}
 \exp\{-a(v_1+\cdots+v_d)\},\qquad v\in\mathbb R_+^d.
\end{equation}
In particular its $d$-fold Mellin transform exists whenever all real
parts of its arguments are positive.

Set $s_d=1$ and define
\[
 G(s_1,u,s_2,\ldots,s_{d-1})
 =-\partial_{s_1}\partial_u
 \mathcal K_\mu(s_1,s_2,\ldots,s_{d-1},1,u).
\]
Differentiation under the integral is justified by
\eqref{p06:eq:tilted-variance-bound} at zero and by $E_1+C>\rho$ at infinity.
For complex $z,w,z_2,\ldots,z_{d-1}$ put
\begin{align*}
 a_1&=z+w, & a_j&=z_j\quad(2\le j<d),\\
 a_d&=\frac{E_1+C-\rho-E_1z-Cw-
                    \sum_{j=2}^{d-1}E_jz_j}{E_d}.
\end{align*}
Whenever $\Re z,\Re w,\Re z_j,\Re a_d>0$, its Mellin transform satisfies
\begin{equation}\label{p06:eq:correlated-mellin-identity}
 \mathcal M G(z,w,z_2,\ldots,z_{d-1})
 =\frac{\rho}{E_d}\,\mathrm B(z,w)
     \mathcal M f_\mu(a_1,\ldots,a_d).
\end{equation}
To verify this, substitute
$v_1=s_1y^{-E_1}$ and $v_0=uy^{-C}$ in the first two Mellin integrals,
and $v_j=s_jy^{-E_j}$ in the others. The first two integrals combine by
\[
 \int_0^\infty\!\int_0^\infty
 v_1^{z-1}v_0^{w-1}F(v_1+v_0)\,dv_1\,dv_0
 =\mathrm B(z,w)\int_0^\infty r^{z+w-1}F(r)\,dr.
\]
The remaining substitution $v_d=y^{-E_d}$ gives the factor
$1/E_d$ and the exponent $a_d$ in \eqref{p06:eq:correlated-mellin-identity}.
For real parameters Tonelli's theorem applies. For complex parameters
absolute convergence follows from the same calculation with real parts.

Choose small positive real $z_0,w_0,z_{j,0}$ so that $a_{d,0}>0$.
The affine map from $(z,w,z_2,\ldots,z_{d-1})$ to $(a_1,\ldots,a_d)$
has nonzero determinant, up to sign,
\[
 \frac{E_1-C}{E_d}.
\]
Therefore all imaginary parts of $(a_1,\ldots,a_d)$ are obtained while
their real parts stay fixed at a positive vector. The beta factor has
no zeros in the domain in question. Equation
\eqref{p06:eq:correlated-mellin-identity} determines the Mellin transform of
$f_\mu$ on an entire vertical $d$-plane. Uniqueness of the Fourier
transform, after logarithmic coordinates, determines $f_\mu$.
The bound \eqref{p06:eq:tilted-variance-bound} supplies the required
integrability. Continuity makes the identification pointwise.

It remains to recover the integration constants lost by differentiating.
As $r\to\infty$, uniformly for fixed $v'\in\mathbb R_+^{d-1}$,
\[
 h_\mu(r,v')\longrightarrow c_\mu=-\log p_0,
 \qquad \partial_1h_\mu(r,v')\longrightarrow0.
\]
Consequently
\begin{equation}\label{p06:eq:recover-h-from-variance}
 h_\mu(v_1,v')=c_\mu-
    \int_{v_1}^\infty(r-v_1)f_\mu(r,v')\,dr,
 \qquad
 c_\mu=\int_0^\infty r f_\mu(r,0,\ldots,0)\,dr.
\end{equation}
The second equality uses $h_\mu(0)=0$; the boundary values of $f_\mu$
are obtained by continuity and dominated convergence from the positive
orthant. Hence $h_\mu$, then its Laplace transform $e^{-h_\mu}$, and
finally $\mu$ are determined.
\end{proof}

\begin{corollary}[The same distinction for dual rate functions]
\label{p06:cor:rate-identification}
Define, for either response $F=\mathcal H_\mu$ or $\mathcal K_\mu$,
\[
 I_F(e)=\sup_{q\ge0}\{F(q)-q\cdot e\}.
\]
Then $I_F$ determines $F$ on the positive orthant by
\[
 F(q)=\inf_{e\ge0}\{I_F(e)+q\cdot e\}.
\]
Thus the nonidentifiability and restoration conclusions above hold
unchanged when the exact complete dual rate function is the observable.
\end{corollary}
\begin{proof}
The inequality $F(q)\le I_F(e)+q\cdot e$ is the definition of $I_F$.
For equality choose $e=\nabla F(q)$, which is coordinatewise nonnegative.
Concavity implies that $q$ maximizes $F(v)-v\cdot e$, proving the formula.
The statement is about exact functions; it does not identify them from a
finite noisy sample, nor by itself establish an array large-deviation
principle.
\end{proof}

\begin{remark}[Conditioning and stability]
The reconstruction is explicit but not uniformly stable. In the simplest
case $d=2$, along $z=z_0+iR$, $w=w_0-iR$, the beta multiplier satisfies
\[
 |\mathrm B(z,w)|\asymp
 |R|^{z_0+w_0-1}e^{-\pi|R|}.
\]
Recovering the corresponding Mellin modes therefore requires exponential
amplification. This multiplier calculation alone is not a minimax lower
bound over the nonlinear probability-law class. Such a claim would
require separate admissible perturbations and a statistical noise model.
\end{remark}

\subsection{A scale window gives a different, local reconstruction}

The information loss is also removed by retaining one finite radial
window. For $c>1$ define
\[
 A_{\mu,c}(s)=\rho\int_1^c y^{\rho-1}h_\mu(\delta_{1/y}s)\,dy,
 \qquad D=\sum_{j=1}^dE_js_j\partial_{s_j}.
\]
The following inversion only needs a nonnegative law with finite means;
it does not need the atom or support-gap assumptions.
\begin{proposition}[One-window inversion]\label{p06:prop:window-inversion}
Let $E_*:=\min_jE_j>\rho$. Then
\begin{align}
 g(s)&:=\frac1\rho(D-\rho)A_{\mu,c}(s)
       =h_\mu(s)-c^\rho h_\mu(\delta_{1/c}s),\label{p06:eq:window-defect}\\
 h_\mu(s)&=\sum_{n=0}^\infty
       c^{n\rho}g(\delta_{c^{-n}}s).\label{p06:eq:window-inverse}
\end{align}
The series converges absolutely and locally uniformly. In the norm
$\|f\|_*=\sup_{s\in\mathbb R_+^d\setminus\{0\}}|f(s)|/(s_1+\cdots+s_d)$,
its linear inverse has norm at most
$(1-c^{\rho-E_*})^{-1}$. Thus recovering $h$ from the differentiated
window data loses no further exponential factor.
\end{proposition}
\begin{proof}
Since $D h(\delta_{1/y}s)=-y\partial_yh(\delta_{1/y}s)$,
integration by parts proves \eqref{p06:eq:window-defect}.
Iterate this identity; the remainder after $n$ terms is
$c^{n\rho}h(\delta_{c^{-n}}s)$, bounded by
$c^{n(\rho-E_*)}(\max_j\mathbb EX_j)\sum_js_j$.
It tends to zero. The dilation operator
$Tf(s)=c^\rho f(\delta_{1/c}s)$ has norm at most
$c^{\rho-E_*}<1$, so the Neumann series gives the stated bound.
\end{proof}

\paragraph{Prior-art boundary.}
Homogeneous dilation, convolution homomorphisms, and inversion of
random-integral transforms are established themes in probability;
Jurek's work on random integral mappings~\cite{Jurek} is directly relevant.
The beta-integral identity and Mellin uniqueness used above are also
classical. The prospective contribution here is the exact information
comparison for these particular rare-event response maps: finite
calibrations do not eliminate genuine correlated aliases, while one
repeated component restores the entire joint local law.
No priority claim for this combination should be made without a more
complete literature review. In particular, the response
$\mathcal H_\mu$ for a general bounded non-infinitely-divisible local law
must not be identified automatically with an operator-stable Laplace
exponent.


\section{Weighted-series realization of the correlated responses}
\label{p06:sec:correlated-ldp}

The response functions above are realized by genuine rare events without
assuming independence between the coordinates at one site. The following
theorem also permits repeated local coordinates.

\begin{theorem}[Joint lower tails with arbitrary local correlations]
\label{p06:thm:correlated-array-ldp}
Let $k\ge1$, $\rho>0$, and $F_i>\rho$ for $1\le i\le k$.
Let $Y^{(n)}$ be independent identically distributed vectors with law
$p_0\delta_0+(1-p_0)\nu$, where $0<p_0<1$ and
$\operatorname{supp}\nu\subset[a,b]^k$ for some $0<a<b$.
Coordinates of each vector may be arbitrarily dependent, including
identical. Put
\[
 D_i=\sum_{n\ge1}n^{-F_i/\rho}Y_i^{(n)},\qquad
 U_{x,i}=x^{F_i-\rho}D_i,
\]
and
\[
 H(q)=\rho\int_0^\infty y^{\rho-1}
       h(q_1y^{-F_1},\ldots,q_ky^{-F_k})\,dy,
 \quad h(v)=-\log\mathbb E e^{-v\cdot Y}.
\]
Then, for every vector $e\in(0,\infty)^k$,
\begin{equation}\label{p06:eq:correlated-lower-tail-ldp}
 \lim_{x\to\infty}-x^{-\rho}
 \log\mathbb P\{U_{x,i}\le e_i\text{ for every }i\}
 =I(e):=\sup_{q\in[0,\infty)^k}\{H(q)-q\cdot e\}.
\end{equation}
In particular, the $d+1$-channel response in
\eqref{p06:eq:augmented-correlated-H} is obtained by taking
$Y=(X_1,\ldots,X_d,X_1)$ and $F=(E_1,\ldots,E_d,C)$.
\end{theorem}
\begin{proof}
All $D_i$ are almost surely finite, by finite expectation and
$F_i/\rho>1$. Set $M=x^\rho$. The exact normalized negative logarithm of
the Laplace transform is
\[
 H_x(q):=-M^{-1}\log\mathbb E e^{-M q\cdot U_x}
 =M^{-1}\sum_{n\ge1}
 h\big(q_1(n/M)^{-F_1/\rho},\ldots,
       q_k(n/M)^{-F_k/\rho}\big).
\]
For each fixed $q\ne0$, this Riemann sum converges to $H(q)$.
Indeed, $h\le-\log p_0$ controls the origin, and
$h(v)\le b\sum_i v_i$ supplies an integrable bound at infinity.
The same Riemann-sum argument gives convergence of its first and second
partial derivatives at any fixed $q\in[0,\infty)^k\setminus\{0\}$,
including one-sided derivatives at zero coordinates.
Near the origin of the Riemann-sum variable, the tilted first moments
and covariances have exponential bounds of the form
\[
 |\mathbb E_vY_i|+|\operatorname{Cov}_v(Y_i,Y_j)|
 \le K\exp\{-a\textstyle\sum_\ell v_\ell\};
\]
at infinity, the factors $t^{-F_i/\rho}$ and
$t^{-(F_i+F_j)/\rho}$ are integrable. Splitting into a compact interval
and these two tails justifies the derivative limits.

Under the tilted probability measure with density proportional to
$e^{-M q\cdot U_x}$, one has
\[
 \mathbb E_q U_x=\nabla H_x(q)\longrightarrow\nabla H(q),\qquad
 \operatorname{Cov}_q(U_x)=-M^{-1}\nabla^2H_x(q)=O(M^{-1}).
\]
Thus $U_x\to\nabla H(q)$ in probability under every such fixed tilt.

For $e>0$, the supremum defining $I(e)$ is attained. To see coercivity,
write $c_0=-\log p_0$ and use
\[
 h(v)\le\min(c_0,b\textstyle\sum_i v_i)
 \le\sum_i\min(c_0,bv_i).
\]
Integration gives $H(q)\le\sum_i K_i q_i^{\rho/F_i}$, whereas
$q\cdot e$ grows linearly. The maximizing vector $q_*$ is nonzero:
on every coordinate axis, $H(te_i)=t^{\rho/F_i}H(e_i)$, so for small
positive $t$ the objective exceeds its value zero at the origin.
Concavity and the optimality conditions give
\[
 m_i:=\partial_iH(q_*)\le e_i,
 \qquad q_{*,i}(m_i-e_i)=0.
\]
Every $m_i$ is strictly positive. No sign assumption on off-diagonal
covariances is used.

For $c>1$, put $q_{c,i}=c^{F_i}q_{*,i}$. Homogeneity yields
\[
 H(q_c)=c^\rho H(q_*),\qquad
 \partial_iH(q_c)=c^{\rho-F_i}m_i<e_i.
\]
Thus a sufficiently small fixed neighborhood of $\nabla H(q_c)$ is
contained in the lower orthant $\{u\le e\}$ and has tilted probability
tending to one. Changing measure on that neighborhood gives
\[
 \limsup_{x\to\infty}-M^{-1}\log\mathbb P\{U_x\le e\}
 \le H(q_c)-q_c\cdot\nabla H(q_c)
 =c^\rho I(e),
\]
where the neighborhood radius is sent to zero first. Letting $c\downarrow1$
proves the upper bound on the logarithmic cost. Chernoff's inequality,
valid for every $q\ge0$, gives the reverse bound. This proves
\eqref{p06:eq:correlated-lower-tail-ldp}.
\end{proof}

\begin{remark}
The anisotropic dilation in this proof is useful when coordinates are
negatively correlated: increasing all ordinary tilt coordinates need not
reduce every tilted mean. The homogeneity identity gives a simultaneous
strict decrease by the explicit factors $c^{\rho-F_i}$ instead.
\end{remark}

\begin{proposition}[Active lower tails for arbitrary bounded gated vectors]
\label{p06:prop:gated-active-ldp}
In the preceding weighted-series construction, let $Y$ be any bounded
nonnegative vector with $\mathbb P(Y=0)>0$. At each strictly positive
tilt $q$, write $e=\nabla H(q)$. After deleting identically zero
coordinates, $e>0$ and the limit in
\eqref{p06:eq:correlated-lower-tail-ldp} holds at this $e$.
In particular, this applies to bounded vectors observed through
nonnegative linear channels and a shared Bernoulli gate.
\end{proposition}
\begin{proof}
The log-Laplace Riemann sums have the same endpoint bounds as before.
For derivatives at $q>0$, put
$S=\sum_i q_i y^{-F_i}Y_i$. Each
$y^{-F_i}Y_i\le S/q_i$, and the tilted denominator is at least
$\mathbb P(Y=0)$. Since $S^j e^{-S}$ is bounded for $j=1,2$,
the first and second derivative integrands are bounded multiples of
$y^{\rho-1}$ at zero. At infinity boundedness of $Y$ gives the same
integrable power bounds as before. Consequently the derivative
Riemann sums converge and the tilted variances vanish. Concavity makes
$q$ an optimizer at $e=\nabla H(q)$. Under
$q_{c,i}=c^{F_i}q_i$, every positive mean is multiplied by
$c^{\rho-F_i}<1$. The change-of-measure proof just given then applies.
No sign condition on correlations or support gap is needed at a
strictly positive tilt.
\end{proof}


\subsection{Power-law behavior at the origin in a positive cone}
\label{p06:sec:continuous-correlated-extension}

The origin atom in Theorem~\ref{p06:thm:one-channel-restoration} can be
replaced by a power law near zero. The following hypothesis permits
arbitrary dependence between radial size and direction.

\begin{definition}
A bounded nonnegative random vector $X$ belongs to $\mathcal P_d$ if
there are a random variable $R\in(0,1]$, constants $a,b,c_-,c_+>0$,
and $\kappa>0$ such that almost surely
\[
 aR\le X_j\le bR\quad(1\le j\le d),\qquad
 c_-r^\kappa\le\mathbb P\{R\le r\}\le c_+r^\kappa
 \quad(0<r\le1).
\]
The constants and exponent may depend on the law. No independence of
$R$ and the direction of $X$ is required.
\end{definition}

\begin{proposition}[Restoration for power-law radial behavior]
\label{p06:prop:continuous-restoration}
The injectivity conclusion of
Theorem~\ref{p06:thm:one-channel-restoration} holds with $\mathcal C_d$
replaced by $\mathcal P_d$. In addition, the exact alias construction
of Theorem~\ref{p06:thm:correlated-alias} remains within $\mathcal P_d$.
Distinct aliases can preserve any prescribed finite collection of mixed
moments and one-component responses as in
Proposition~\ref{p06:prop:finite-calibration-alias}.
\end{proposition}
\begin{proof}
Write $T=\sum_jv_j$. The assumptions and integration by parts for the
distribution function of $R$ give, for each integer $m\ge0$,
\[
 \mathbb E e^{-v\cdot X}\ge c(1+T)^{-\kappa},\qquad
 \mathbb E[R^m e^{-v\cdot X}]
 \le C_m(1+T)^{-\kappa-m}.
\]
For the lower bound restrict to $R\le(1+T)^{-1}$.
For the upper bound use $v\cdot X\ge aTR$; the estimate follows by
splitting into intervals of length $(1+T)^{-1}$ and summing an
exponentially decaying series, or by Stieltjes integration by parts.
Thus every tilted moment of total order $m$ is
$O((1+T)^{-m})$. In particular,
\begin{equation}\label{p06:eq:continuous-variance-decay}
 f(v):=-h_{11}(v)=O((1+T)^{-2}),\quad
 h_j(v)=O((1+T)^{-1}),\quad
 |\partial_j f(v)|=O((1+T)^{-3}).
\end{equation}
Also $h(v)=O(1+\log(1+T))$ as $T\to\infty$, so both response
integrals remain finite. The $d$-fold Mellin transform of $f$ exists
on the nonempty tube
\[
 \Re a_j>0,\qquad \sum_j\Re a_j<2.
\]
The identity \eqref{p06:eq:correlated-mellin-identity} is still valid whenever
its transformed parameters lie in this tube. Such a vertical plane
is attainable: take all the Mellin variables small and positive except
the variable attached to the larger of $E_1$ and $C$, and choose that
variable just below
\[
 \frac{E_1+C-\rho}{\max(E_1,C)}<2.
\]
The remaining transformed parameter $a_d$ can be made positive and
arbitrarily small, while the sum of all $a_j$ stays below two.
The same invertible affine map and Mellin uniqueness therefore identify
$f$.

To check the integration constants without an atom at zero, suppose two
laws give the same $f$. Their logarithmic Laplace transforms differ by
$A(v')v_1+B(v')$. The derivative bound in
\eqref{p06:eq:continuous-variance-decay} forces $A(v')=0$.
Furthermore, for each fixed $v'\ge0$,
\[
 |h(v_1,v')-h(v_1,0)|\le
 C\|v'\|_1/(1+v_1)\longrightarrow0.
\]
It follows that $B(v')=B(0)=0$, since both transforms vanish at the
origin. Thus the complete law is identified.
If an explicit reconstruction is desired, the same bounds justify
\begin{align*}
 h(v_1,0)&=\int_0^\infty\min(v_1,r)f(r,0)\,dr,\\
 h(v_1,v')-h(v_1,0)
 &=-\int_{v_1}^\infty(r-v_1)
       [f(r,v')-f(r,0)]\,dr.
\end{align*}
The second integrand is integrable by the third estimate in
\eqref{p06:eq:continuous-variance-decay}.

Finally, a finite dilation-convolution is comparable in every coordinate
to a positive weighted sum of independent copies of $R$.
Its lower-tail distribution is bounded above and below by constants
 times $r^{n\kappa}$: use the inclusions that a sum at most $r$ requires
each summand at most $r$, and is guaranteed when each summand is at most
$r/n$, with the fixed positive coefficients absorbed into constants.
After a fixed rescaling this proves membership in $\mathcal P_d$.

For the finite-calibration assertion, replace the Bernoulli seed in
Proposition~\ref{p06:prop:finite-calibration-alias} by a common
$\operatorname{Beta}(\alpha,1)$ variable, with a transcendental
$\alpha>0$. Every cumulant is a rational function of $\alpha$, since
its $m$th moment is $\alpha/(\alpha+m)$. That rational function is
nonzero: as $\alpha\downarrow0$, its leading term is $\alpha/m$.
Consequently all cumulants are nonzero at transcendental $\alpha$.
The identical power-sum and implicit-function argument now supplies the
required distinct laws in $\mathcal P_d$.
\end{proof}

\begin{remark}
At fully active maximizing tilts, the lower-tail proof of
Theorem~\ref{p06:thm:correlated-array-ldp} also applies to $\mathcal P_k$.
For clarity, the necessary replacements are
$h=O(1+\log(1+T))$, integrable at the lower endpoint of the response
integral, and the tilted moment estimates above. Near zero in the
Riemann-sum variable, at a strictly positive tilt, the first and second
derivative integrands are bounded because the same positive tilt
contains all coordinates; hence that proof immediately covers fully
active maximizing tilts. Boundary maximizing tilts require separate
integrability checks when some exponents differ and are not asserted
here without them.
\end{remark}


\section{One additional mixed probe identifies a vector law}
\label{p06:sec:mixed-probe}

For unrestricted nonnegative vector laws, repeating one coordinate is
insufficient to extend scalar two-weight identification to a vector.
We instead add one positive
linear combination of all coordinates at a different power.
The construction is an explicit anisotropic polar decomposition,
within the classical theory of L\'evy mixing and homogeneous measures
\cite{BRT,BPT,EM}. The candidate contribution is the sharp
observation design and its applicability after one common random gate,
not a new L\'evy representation theorem.

\subsection{An explicit signed-measure inverse}

Put $C_k=[0,\infty)^k\setminus\{0\}$ and use the $\ell^1$ norm.
Fix $\rho>0$, $p,q>\rho$, $p\ne q$, and $v\in(0,\infty)^k$.
Let $\nu$ be a signed Radon measure on $C_k$ such that
\begin{equation}\label{p06:eq:mp-moment}
 \int_{C_k}(\|x\|^{\rho/p}+\|x\|^{\rho/q})|\nu|(dx)<\infty.
\end{equation}
Define
\begin{align}
 h_\nu(a)&=\int_{C_k}(1-e^{-a\cdot x})\nu(dx),\label{p06:eq:mp-h}\\
 \mathcal H_\nu(a,t)
 &=\int_0^\infty\rho y^{\rho-1}
       h_\nu(y^{-p}a+t y^{-q}v)\,dy,\qquad a>0,\ t>0.
       \label{p06:eq:mp-H}
\end{align}
For $\nu\ge0$, $h_\nu$ is the Laplace exponent of a drift-free
nonnegative infinitely divisible vector. There are $k$ coordinate
probes at power $p$ and one shared probe $v\cdot X$ at power $q$.

\begin{theorem}[Explicit mixed-probe recovery]\label{p06:thm:mp-recovery}
Exact values of $\mathcal H_\nu$ on any nonempty open subset of
$(0,\infty)^{k+1}$ determine $\nu$ uniquely.
Let $\eta$ be the unique signed Radon measure on $C_{k+1}$, with
$\int(1\wedge\|b\|)|\eta|(db)<\infty$, satisfying
\begin{equation}\label{p06:eq:mp-eta-exponent}
 \mathcal H_\nu(a,t)
       =\int(1-e^{-a\cdot z-tw})\eta(dz,dw).
\end{equation}
It is supported on $z\ne0,w>0$. Set
\begin{equation}\label{p06:eq:mp-inverse-map}
 R(z,w)=\left(\frac{w}{v\cdot z}\right)^{1/(p-q)},
 \qquad X(z,w)=R(z,w)^p z.
\end{equation}
For each continuous $f$ compactly supported in $C_k$,
\begin{equation}\label{p06:eq:mp-inverse-measure}
 \int f\,d\nu=\frac1{2^\rho-1}
       \int_{1\le R(z,w)\le2}f(X(z,w))\eta(dz,dw).
\end{equation}
\end{theorem}

\begin{proof}
For $r>\rho$ and $c\ge0$, substitution and integration by parts give
\begin{equation}\label{p06:eq:mp-gamma-identity}
 \int_0^\infty\rho y^{\rho-1}(1-e^{-cy^{-r}})\,dy
       =\Gamma(1-\rho/r)c^{\rho/r}.
\end{equation}
Together with $1-e^{-(A+B)}\le(1-e^{-A})+(1-e^{-B})$,
this proves the absolute integrability needed below.
Push $\rho y^{\rho-1}dy\,\nu(dx)$ forward by
\[
 F(y,x)=(y^{-p}x,\ y^{-q}v\cdot x).
\]
Its total variation is bounded by the same pushforward of
$\rho y^{\rho-1}dy\,|\nu|(dx)$. Equation
\eqref{p06:eq:mp-gamma-identity} and the comparability of $1-e^{-u}$
and $1\wedge u$ show that this image $\eta$ has the asserted
integrability. Absolute Fubini proves \eqref{p06:eq:mp-eta-exponent}.

The signed representation is unique by finite-measure Laplace
uniqueness. Indeed, the sum of all first partial derivatives,
evaluated at $(a,t)=(1,\ldots,1)+\lambda$, is the Laplace transform
at $\lambda\ge0$ of the finite signed measure
$\|b\|e^{-\|b\|}\eta(db)$. Its weight is positive on $C_{k+1}$.
Finite Laplace uniqueness itself follows by setting $u_i=e^{-b_i}$
and using density of polynomials on $[0,1]^{k+1}$.
The integral is holomorphic when all parameters have positive real
parts, so real-analytic continuation makes any real open patch
sufficient.

Finally $w/(v\cdot z)=y^{p-q}$. Thus $F$ is a Borel isomorphism
from $(0,\infty)\times C_k$ onto $\{(z,w):z\ne0,w>0\}$,
with inverse \eqref{p06:eq:mp-inverse-map}. Pull back $\eta$ and restrict
$y$ to $[1,2]$. Its mass in the scale variable is $2^\rho-1$,
which gives \eqref{p06:eq:mp-inverse-measure}.
\end{proof}

\begin{remark}\label{p06:rem:mp-data}
The data are a full continuum function, not $k+1$ numbers.
The number of probes counts aggregate energies whose joint limiting
Laplace exponent is supplied. Analytic continuation and Laplace
inversion are part of recovery; the explicit scale formula alone
does not establish numerical stability.
\end{remark}

\begin{corollary}[Unequal base powers]\label{p06:cor:mp-unequal}
Identification still holds with coordinate powers $p_i>\rho$
and a shared power $q>\max_i p_i$ or $\rho<q<\min_i p_i$,
assuming
\[
 \int\left(\sum_i x_i^{\rho/p_i}+(v\cdot x)^{\rho/q}\right)
       |\nu|(dx)<\infty.
\]
The inverse scale is the unique positive solution of
$w=\sum_i v_i z_i y^{p_i-q}$; then $x_i=y^{p_i}z_i$.
\end{corollary}
\begin{proof}
Use $z_i=y^{-p_i}x_i$ in the preceding proof.
The right side of the scale equation is strictly monotone, onto
$(0,\infty)$, since every nonzero summand has the same exponent sign.
\end{proof}

\begin{corollary}[Unknown counting exponent]\label{p06:cor:mp-rho}
For nonzero $\nu\ge0$, knowledge of $p,q,v$ and $\mathcal H_\nu$
also determines $\rho$ through
\[
 \mathcal H_\nu(\lambda^p a,\lambda^q t)
       =\lambda^\rho\mathcal H_\nu(a,t).
\]
\end{corollary}
\begin{proof}
Substitute $y=\lambda u$. The function is positive, so its
homogeneity degree is unique.
\end{proof}

\subsection{A sharp observation threshold}

For nonnegative linear forms $\ell_i$, powers $p_i>\rho$, and
$m$ probes, set
\begin{equation}\label{p06:eq:mp-general-design}
 \mathcal K_\nu(s)=
 \int_0^\infty\!\rho y^{\rho-1}
 \int_{C_k}\!\left(1-\exp\left\{-\sum_{i=1}^m
       s_i y^{-p_i}\ell_i(x)\right\}\right)\nu(dx)\,dy.
\end{equation}

\begin{theorem}[At most $k$ linear probes never suffice]
\label{p06:thm:mp-minimality}
For $m\le k$, $\nu\mapsto\mathcal K_\nu$ is not injective even
on finite positive measures with finite support in $(0,\infty)^k$.
The two measures may have the same total mass, as small as desired.
Thus the associated distinct compound-Poisson laws may have the
same prescribed atom at zero, arbitrarily close to one.
The minimum number for uniform identification of the stated
drift-free infinitely divisible class is therefore exactly $k+1$.
\end{theorem}
\begin{proof}
Let $M$ be the $m\times k$ matrix of the forms.
If $\operatorname{rank}M<k$, choose $x_0>0$ and nonzero
$u\in\ker M$. Small $\epsilon\ne0$ keeps $x_0+\epsilon u>0$,
and $\lambda\delta_{x_0}$ and
$\lambda\delta_{x_0+\epsilon u}$ are indistinguishable.

Otherwise $m=k$ and $M$ is invertible. For $c$ close to one put
$x_c=M^{-1}\operatorname{diag}(c^{p_i})Mx_0>0$.
Substitution $y=cu$ shows
$\mathcal K_{\delta_{x_c}}=c^\rho\mathcal K_{\delta_{x_0}}$.
Choose $c_-<1<c_+$ close to one and $a\in(0,1)$ with
$ac_-^\rho+(1-a)c_+^\rho=1$. The distinct measures
\[
 \nu_1=\lambda\delta_{x_0},\qquad
 \nu_2=\lambda\{a\delta_{x_{c_-}}+(1-a)\delta_{x_{c_+}}\}
\]
have the same observations and mass $\lambda$. Their compound-Poisson
laws have atom $e^{-\lambda}$ at zero because all jumps are positive.
The upper bound is Theorem~\ref{p06:thm:mp-recovery}.
\end{proof}

\begin{proposition}[The extra probe must see every hidden coordinate]
\label{p06:prop:mp-wrong-repeat}
The two-dimensional design
\[
 \int_0^\infty\rho y^{\rho-1}
       h(sy^{-E}+ty^{-A},uy^{-B})\,dy
\]
does not identify the joint law, even for independent coordinates
and distinct $E,A,B>\rho$.
\end{proposition}
\begin{proof}
Independence gives $h(a,b)=h_V(a)+h_W(b)$.
The second coordinate contributes only
$u^{\rho/B}\int_0^\infty\rho y^{\rho-1}h_W(y^{-B})\,dy$.
For a compound-Poisson $W$ with jump measure $\xi$, this coefficient
is $\Gamma(1-\rho/B)\int x^{\rho/B}\xi(dx)$.
Choose two distinct $\xi$ of the same mass and fractional moment,
one atom versus a mixture of atoms on either side, and hold $V$ fixed.
\end{proof}

At $p=q$ even the positive shared probe fails, since
\[
 \mathcal H_\nu(a,t)=\Gamma(1-\rho/p)
                  \int((a+tv)\cdot x)^{\rho/p}\nu(dx).
\]
The additional power breaks the otherwise invisible dilation.

\subsection{A common Bernoulli gate gives universal recovery}

\begin{corollary}[Arbitrary vector laws after gating]
\label{p06:cor:mp-gate}
Let $\mu$ be any law on $[0,\infty)^k$ satisfying
$\mathbb E\|X\|^{\rho/p}+\mathbb E\|X\|^{\rho/q}<\infty$;
every bounded nonnegative vector is allowed.
Let $B$ be independent of $X$, with known
$\mathbb P(B=0)=p_0\in(1/2,1)$ and $\mathbb P(B=1)=1-p_0$.
The mixed-probe continuum function of $BX$ uniquely determines
the original joint law $\mu$.
The same Bernoulli variable must multiply all coordinates of each
local copy; different copies use independent gates.
\end{corollary}
\begin{proof}
Put $r=(1-p_0)/p_0<1$. The classical logarithm series gives the
signed measure, with its atom at zero discarded,
\begin{equation}\label{p06:eq:mp-quasi-levy}
 \nu=\sum_{n\ge1}\frac{(-1)^{n+1}r^n}{n}\mu^{*n}
                       \quad\hbox{on }C_k,
\end{equation}
and
\[
 -\log\bigl(p_0+(1-p_0)\mathbb E e^{-a\cdot X}\bigr)
       =\int(1-e^{-a\cdot x})\nu(dx).
\]
For $0<\alpha<1$, subadditivity yields
\[
 \int\|x\|^\alpha|\nu|(dx)
 \le\sum_{n\ge1}\frac{r^n}{n}
       \mathbb E\|X_1+\cdots+X_n\|^\alpha
 \le\frac{r}{1-r}\mathbb E\|X\|^\alpha.
\]
Apply Theorem~\ref{p06:thm:mp-recovery}; the recovered $h_\nu$ gives
$\mathbb E e^{-a\cdot X}=(e^{-h_\nu(a)}-p_0)/(1-p_0)$,
which determines $\mu$.
\end{proof}

The representation \eqref{p06:eq:mp-quasi-levy} is Cuppens' classical
quasi-infinite-divisibility construction, reproduced in
\cite[Theorem 3.2]{BKL}. Gating is an additional intervention,
not an inference from the ungated data of an arbitrary law.

\begin{theorem}[Sharpness within the bounded gated class]
\label{p06:thm:mp-bounded-minimal}
For any fixed known $p_0\in(1/2,1)$, no collection of at most
$k$ nonnegative linear power probes with exponents above $\rho$
identifies every bounded nonnegative vector law from the continuum
functions of its copies gated with probability $p_0$.
\end{theorem}
\begin{proof}
Write $M$ for the probe matrix. Its kernel gives distinct positive
vectors with identical probe values if $\operatorname{rank}M<k$.
Gating their point masses proves that case.

Otherwise $M$ is invertible. Choose $x_0>0$ and write
$T_c=M^{-1}\operatorname{diag}(c^{p_i})M$ for $c$ close to one.
Fix $b\in(\sqrt{p_0},1)$, and let $Z_1,Z_2$ be independent,
each equal to zero with probability $b$ and to $x_0$ otherwise.
For sufficiently small $\epsilon>0$ put
\[
 d_\pm=(1\pm\epsilon)^{1/\rho},\qquad
 Y_1=Z_1+Z_2,\qquad Y_2=T_{d_+}Z_1+T_{d_-}Z_2.
\]
Both vectors are bounded and nonnegative and have atom $b^2>p_0$
at zero. Thus each has the form
$p_0\delta_0+(1-p_0)\mu_i$ with a bounded probability law $\mu_i$.
Independence makes logarithmic Laplace transforms additive, and
substitution $y=cu$ multiplies the contribution of $T_cZ$ by
$c^\rho$. Since $d_+^\rho+d_-^\rho=2$, the complete observed
functions of $Y_1$ and $Y_2$ agree.

The laws differ: for any nonzero probe row $\ell_j$ the ratio
of its mean under $Y_2$ to that under $Y_1$ is
\[
 \frac{d_+^{p_j}+d_-^{p_j}}2
 =\frac{(1+\epsilon)^{p_j/\rho}
                  +(1-\epsilon)^{p_j/\rho}}2>1
\]
by strict convexity. Hence $\mu_1\ne\mu_2$. Combined with
Corollary~\ref{p06:cor:mp-gate}, this proves the exact threshold.
\end{proof}

\begin{proposition}[Exact placement of the extra exponent]\label{p06:prop:exact-placement}
Fix $p_0\in(1/2,1)$ and let $\mathcal P_b$ be the class of all compactly
supported probability laws on $[0,\infty)^k$, allowing mass at the origin.
For $\mu\in\mathcal P_b$, write
\[
 P_\mu=p_0\delta_0+(1-p_0)\mu,\qquad
 h_\mu(a)=-\log\int e^{-a\cdot x}\,P_\mu(dx).
\]
Fix $p_1,\ldots,p_k,q>\rho>0$ and $v_i>0$, and define
\[
 H_\mu(a,t)=\int_0^\infty \rho y^{\rho-1}
 h_\mu\bigl((a_i y^{-p_i})_{i=1}^k+t y^{-q}v\bigr)\,dy.
\]
If $q\in[\min_i p_i,\max_i p_i]$, the map
$\mu\mapsto H_\mu$ is not injective on $\mathcal P_b$.
Consequently, together with the inversion theorem, universal identification
in this fixed-gate class holds exactly when
$q<\min_i p_i$ or $q>\max_i p_i$.
\end{proposition}

\begin{proof}
Suppose first that $\min_i p_i<q<\max_i p_i$.
Choose distinct adjacent positive integers $n,m$ and put
$c=(n/m)^{1/\rho}$. The coefficients
$v_i(c^{p_i}-c^q)$ have both signs. We can therefore choose
$x\in(0,\infty)^k$ such that
\[
 \sum_i v_i(c^{p_i}-c^q)x_i=0.
\]
Set $x'_i=c^{p_i}x_i$; then $v\cdot x'=c^q v\cdot x$.
Choose $\theta\in(0,1)$ small enough that
$(1-\theta)^{\max(n,m)}>p_0$, and let
\[
 V=xB_n,\qquad V'=x'B_m,
 \qquad B_n\sim\operatorname{Bin}(n,\theta),\quad
 B_m\sim\operatorname{Bin}(m,\theta).
\]
Their laws $P,P'$ are distinct, since their masses at zero are
$(1-\theta)^n$ and $(1-\theta)^m$. Both exceed $p_0$, so
\[
 \mu=\frac{P-p_0\delta_0}{1-p_0},\qquad
 \mu'=\frac{P'-p_0\delta_0}{1-p_0}
\]
are distinct elements of $\mathcal P_b$.
Writing $g(u)=-\log(1-\theta+\theta e^{-u})$, their Laplace exponents
are $h_\mu(a)=ng(a\cdot x)$ and
$h_{\mu'}(a)=mg(a\cdot x')$. The substitution $y=cu$ gives,
for all $a\in[0,\infty)^k$ and $t\ge0$,
\[
\begin{split}
 H_{\mu'}(a,t)
 &=m\int_0^\infty\rho y^{\rho-1}
 g\left(\sum_i a_i c^{p_i}x_i y^{-p_i}
       +t c^q(v\cdot x)y^{-q}\right)\,dy\\
 &=mc^\rho\int_0^\infty\rho u^{\rho-1}
 g\left(\sum_i a_i x_i u^{-p_i}
       +t(v\cdot x)u^{-q}\right)\,du
 =H_\mu(a,t),
\end{split}
\]
because $mc^\rho=n$. The integrals converge: $g$ is bounded at infinity
and is $O(u)$ at zero, while every probe exponent exceeds $\rho$.

If $q$ equals an endpoint of the exponent interval, choose an index $j$
with $p_j=q$, take $x=e_j$, and put $x'=c^q e_j$.
The same construction and calculation apply. This also covers the case
$p_1=\cdots=p_k=q$.
\end{proof}

\begin{remark}
The allowance of an atom at zero in $\mu$ is essential to this particular
counterexample: the fixed gate $p_0$ is a known lower contribution to the
observed zero atom, not an assertion that the observed zero atom equals
$p_0$. If one instead requires $\mu(\{0\})=0$, the above de-gating
argument does not prove the stated lower bound for that smaller class.
\end{remark}

\subsection{Local rate data and a singular power collision}

\begin{example}[An interior exponent can add no identifying power]
Take $\rho=1$, coordinate exponents $(2,4)$, and mixed exponent $3$
with $v=(1,1)$. For independent binomial counts of common success
probability $\theta$, the two vectors
\[
 (2,1)\operatorname{Bin}(2,\theta)
 \quad\hbox{and}\quad
 (8,16)\operatorname{Bin}(1,\theta)
\]
have identical three-channel responses. Indeed, $c=2$ scales the
coordinates by $2^2,2^4$, while $24=2^3\cdot3$ scales their sum;
the factor $c^\rho=2$ compensates for the number of summands.
At $\theta=1/100$, both zero atoms exceed the same known gate
probability $p_0=3/4$, so de-gating gives two different bounded input
laws with identical observed responses. Replacing exponent $3$ by $5$
gives a universally identifying design in this gated class.
\end{example}

\begin{corollary}[An active rate patch suffices]
\label{p06:cor:mp-rate-patch}
Assume $\nu\ge0$, \eqref{p06:eq:mp-moment}, and
$\operatorname{span}(\operatorname{supp}\nu)=\mathbb R^k$.
Define
$\mathcal I(e)=\sup_{\lambda\ge0}
 \{\mathcal H_\nu(\lambda)-\lambda\cdot e\}$.
Any nonempty exact open patch of $\mathcal I$ on which the
maximizer exists and has positive coordinates determines $\nu$.
\end{corollary}
\begin{proof}
The negative Hessian of $\mathcal H_\nu$ has quadratic form
$\int(b\cdot c)^2e^{-\lambda\cdot b}\eta(db)$.
If this vanishes for $c=(c_0,c_1)$, then
$y^{-p}c_0\cdot x+c_1y^{-q}v\cdot x=0$ for almost every $(y,x)$.
The distinct powers and $v\cdot x>0$ force $c_1=0$ and then
$c_0=0$ by the span assumption. The Hessian is therefore negative
definite. The inverse function theorem and strict concavity give
\[
 \lambda=-\nabla\mathcal I(e),\qquad
 \mathcal H_\nu(\lambda)
       =\mathcal I(e)-e\cdot\nabla\mathcal I(e)
\]
on an open image set. Apply Theorem~\ref{p06:thm:mp-recovery}.
\end{proof}

Although $\log R=(\log w-\log(v\cdot z))/(p-q)$ explicitly
exhibits one conditioning factor, much stronger degeneracies occur.

\begin{proposition}[Arbitrarily high-order coalescence]
\label{p06:prop:mp-coalescence}
Fix $p>\rho$, $v>0$, and $\theta\in(0,\infty)^k$.
For each integer $m\ge0$
there are distinct finite positive measures $\nu_\pm$ of equal mass,
supported on finitely many points of the ray $\{r\theta:r>0\}$,
such that
\[
 \left.\partial_q^j\mathcal H_{\nu_+}(a,t;q)\right|_{q=p}
 =\left.\partial_q^j\mathcal H_{\nu_-}(a,t;q)\right|_{q=p},
 \qquad 0\le j\le m,\quad a>0,\ t>0.
\]
Their difference is $O(|q-p|^{m+1})$ uniformly on compact
positive parameter sets, although every fixed $q\ne p$ distinguishes
the two measures.
\end{proposition}
\begin{proof}
Set $\alpha=\rho/p$. Choose $m+3$ distinct positive radii $r_i$
in a compact interval and a nonzero real vector $(c_i)$ satisfying
\[
 \sum_i c_i=0,\qquad
 \sum_i c_i r_i^\alpha(\log r_i)^j=0\quad(0\le j\le m).
\]
There are more unknowns than equations. Perturb positive weights
by $\pm\epsilon c_i$ to obtain the two positive measures.
Substitution $y=r^{1/p}z$ gives
\[
 \mathcal H_{\delta_{r\theta}}(a,t;q)=
 r^\alpha\int_0^\infty\!\rho z^{\rho-1}
 \left(1-e^{-z^{-p}a\cdot\theta
             -t r^{1-q/p}z^{-q}v\cdot\theta}\right)dz.
\]
Its $j$th $q$-derivative at $p$ is $r^\alpha$ times a polynomial
of degree at most $j$ in $\log r$. The coefficient functions do
not depend on $r$, so the annihilating moments prove the claim.
Differentiation and Taylor remainders are uniform on compact sets:
as $z\downarrow0$, $e^{-c z^{-p}}$ dominates every derivative;
as $z\to\infty$, the bound is
$C z^{\rho-1-q_*}(1+|\log z|^{m+1})$ for some $q_*>\rho$.
\end{proof}

These are uniqueness results for exact continuum information.
They provide neither a finite-sample statistical guarantee nor
a priority claim for the classical transform machinery.

\section{A periodic sampling obstruction to a spectral improvement}
\label{p05:sec:periodic-sampling-barrier}

The following construction explains why exact finite-dimensional
nonattainment of a moment certificate need not produce an asymptotic
improvement, even for genuine band-limited evaluation vectors with
uniform spacing and exceptionally regular weighted counting density.
It concerns an abstract real sampling configuration, not the actual
zeros of the zeta function.

\begin{theorem}[A regular sinc configuration saturating the quartic target]
\label{p05:thm:periodic-sinc-barrier}
Let $v_x(t)=e^{2\pi ixt}$ in $L^2([-1/2,1/2])$, so that
$\langle v_x,v_y\rangle=\operatorname{sinc}(x-y)$, where
$\operatorname{sinc}(u)=\sin(\pi u)/(\pi u)$ with its continuous
value at zero. Put
\[
 c=\frac1{2\sqrt2},\qquad
 d=\frac2\pi\arccos c\in(0,1),\qquad \theta=\frac78.
\]
For each integer $k\ge1$, take the simple locations
\[
 \{2n,2n+d:0\le n<16k\}
\]
and the double locations
\[
 \{2n:16k\le n<21k\}.
\]
Define
\[
 G_k=\sum_{0\le n<16k}
       (v_{2n}v_{2n}^*+v_{2n+d}v_{2n+d}^*)
       +2\sum_{16k\le n<21k}v_{2n}v_{2n}^*,
 \qquad N_k=42k.
\]
Then the number of simple locations is $s_k=32k$, the number of
double locations is $b_k=5k$, and
\begin{equation}\label{p05:eq:periodic-sinc-moments}
 \frac1{N_k}\bigl(\operatorname{tr}G_k,
    \operatorname{tr}G_k^2,\operatorname{tr}G_k^3,
    \operatorname{tr}G_k^4\bigr)
 =\left(1,\frac43,2,\frac{13}4\right)
       +O\left(\frac{\log(k+2)}k\right),
 \qquad \frac{s_k}{N_k}=\frac{16}{21}.
\end{equation}
The minimum distance between distinct locations is at least $d$.
If $\mathcal N_k(I)$ counts these locations with their indicated
multiplicities, then every interval $I\subset[0,N_k)$ satisfies
\begin{equation}\label{p05:eq:periodic-sinc-discrepancy}
 |\mathcal N_k(I)-|I||\le4.
\end{equation}
For
\[
 L_k=4\operatorname{tr}G_k-\operatorname{tr}G_k^2-2N_k,
 \qquad D_k=\operatorname{tr}\bigl(G_k(2I-G_k)\bigr)^2,
\]
the nonnegative quartic certificate defect obeys
\begin{equation}\label{p05:eq:periodic-sinc-defect}
 0\le s_k-\frac{2L_k}{\theta}+\frac{D_k}{\theta^2}
       =O(\log(k+2)).
\end{equation}
Consequently separation and the bounded counting discrepancy
\eqref{p05:eq:periodic-sinc-discrepancy} alone cannot force a positive
normalized improvement over $16/21$ from this quartic certificate.
\end{theorem}

\begin{proof}
We first record an elementary two-coset identity.  The doubly
infinite matrix
\[
 C=(\operatorname{sinc}(2(i-j)+d))_{i,j\in\mathbb Z}
\]
defines an operator on $\ell^2(\mathbb Z)$ satisfying
\begin{equation}\label{p05:eq:periodic-sinc-unitary}
 CC^*=C^*C=\cos^2(\pi d/2)I=c^2I.
\end{equation}
Indeed, the families $(v_{2n})$ and $(v_{2n+d})$ are individually
orthonormal.  Identify functions on $[-1/2,1/2]$ with two-component
functions on $[-1/2,0]$ by $f(t)\mapsto(f(t),f(t+1/2))$.
The respective closed spans have fiber directions
$(1,1)/\sqrt2$ and
$e^{2\pi idt}(1,e^{\pi id})/\sqrt2$.
Their fiberwise inner product has constant modulus
$|(1+e^{\pi id})/2|=c$.  Multiplication by its phase is unitary,
which proves \eqref{p05:eq:periodic-sinc-unitary} in the orthonormal
bases indexed by $n$.

Let $C_m$ be the compression of $C$ to $0\le i,j<m$.
If $\sigma_1,\ldots,\sigma_m$ are its singular values, then
$0\le\sigma_j\le c$ and
\begin{align}
 \sum_{j=1}^m(c^2-\sigma_j^2)
 &=\sum_{0\le i<m}\sum_{j\notin[0,m-1]}
             |\operatorname{sinc}(2(i-j)+d)|^2\notag\\
 &=O(\log(m+2)).\label{p05:eq:periodic-sinc-tail}
\end{align}
For the last estimate use
$|\operatorname{sinc}(2r+d)|\ll_d(1+|r|)^{-1}$;
the sum of the exterior tails for row $i$ is
$O((i+1)^{-1}+(m-i)^{-1})$.

The Gram matrix $A_m$ of the $2m$ simple vectors has block form
\[
 A_m=\begin{pmatrix}I&C_m\\ C_m^*&I\end{pmatrix}
\]
and eigenvalues $1\pm\sigma_j$.  For each fixed integer
$1\le r\le4$, the expression
$(1+\sigma)^r+(1-\sigma)^r$ is a polynomial in $\sigma^2$.
It follows from \eqref{p05:eq:periodic-sinc-tail} that
\begin{equation}\label{p05:eq:periodic-sinc-simple-moments}
 \operatorname{tr}A_m^r
 =m\bigl((1+c)^r+(1-c)^r\bigr)+O(\log(m+2)).
\end{equation}
The assertion for $r=1$ is exact.

The nonzero spectrum of $G_k$ is the spectrum of the weighted
Gram matrix $K_k$ of the vectors with factors $1$ at simple
locations and $\sqrt2$ at double locations.  With $m=16k$ and
$b=5k$, write
\[
 K_k=B_k+E_k,\qquad B_k=\operatorname{diag}(A_m,2I_b).
\]
Here $E_k$ has only off-diagonal blocks between the simple and
double groups.  The inner products between base vectors $v_{2i}$
and double vectors $v_{2j}$ vanish.  The other inner products give
\[
 \|E_k\|_{\mathrm{HS}}^2
 \ll_d\sum_{0\le i<m}\sum_{m\le j<m+b}
               (1+j-i)^{-2}
 =O(\log(k+2)).
\]
All operator norms are bounded independently of $k$: the base
vectors $v_{2n}$ are orthonormal and carry weights at most $2$,
while the offset vectors $v_{2n+d}$ are another orthonormal
family carrying weight $1$.  Thus $\|G_k\|=\|K_k\|\le3$;
also $\|B_k\|\le2$ and $\|E_k\|\le5$.
For $1\le r\le4$, expand the trace of $(B_k+E_k)^r$ into words.
A word with exactly one $E_k$ has zero trace, because $B_k$ is
block diagonal and $E_k$ is block off-diagonal.  Every word with
at least two $E_k$ has absolute trace bounded by a constant times
$\|E_k\|_{\mathrm{HS}}^2$, by the Hilbert--Schmidt trace
inequality and the uniform operator norm bounds.  Therefore
\[
 \operatorname{tr}K_k^r
 =m\bigl((1+c)^r+(1-c)^r\bigr)+b2^r
       +O(\log(k+2)).
\]
Substituting $m=16k$, $b=5k$, $c^2=1/8$, and $N_k=42k$
gives \eqref{p05:eq:periodic-sinc-moments}.  More explicitly the
simple-part normalized second, third, and fourth moments are
$6/7$, $22/21$, and $113/84$; the double-part contributions are
$10/21$, $20/21$, and $40/21$, respectively.

Every cell $[2n,2n+2)$ contains total weight exactly $2$.
Removing the at most two partial cells at the ends of an interval
proves \eqref{p05:eq:periodic-sinc-discrepancy}.  Within a simple cell
the two locations have distance $d$; all other consecutive gaps
are at least $2-d>d$.  This proves the spacing assertion.

Finally, we verify nonnegativity without invoking any zeta input.
Write $G_k=P+Q$, where $P$ is the simple part and $Q$ the double
part.  In the decreasing positive spectrum, remove the top $b_k$
eigenvalues.  Rank interlacing with $Q$, which has rank $b_k$,
shows that the remaining number is at most $s_k$ and their sum
is at most $\operatorname{tr}P=s_k$.  For
\[
 q_\theta(x)=
 \frac{8\theta x-(4+2\theta)x^2+4x^3-x^4}{\theta^2},
\]
the identities
\[
 \frac8\theta-q_\theta(x)
 =\frac{(x-2)^2(x^2+2\theta)}{\theta^2},\qquad
 1+\frac{4x}\theta-q_\theta(x)
 =\frac{(x^2-2x+\theta)^2}{\theta^2}
\]
give
$\operatorname{tr}q_\theta(G_k)
\le8b_k/\theta+s_k+4s_k/\theta
=s_k+4N_k/\theta$.
The difference on the right is precisely the defect in
\eqref{p05:eq:periodic-sinc-defect}.
The moment limits above give
$L_k=(2/3)N_k+O(\log(k+2))$ and
$D_k=(7/12)N_k+O(\log(k+2))$.
At $\theta=7/8$ their certificate main term is
$(16/21)N_k=s_k$, proving \eqref{p05:eq:periodic-sinc-defect}.
\end{proof}

The same construction works for any $0<\theta<1$ by taking
$d=(2/\pi)\arccos\sqrt{1-\theta}$ and any fixed positive
proportion of simple cells.  It produces the two simple spectral
values $1\pm\sqrt{1-\theta}$ together with the double spectral
value $2$, with only logarithmic finite-section and interface
errors.  The obstruction is phase separation in the type of
sampling cell, despite uniformly regular multiplicity-weighted
counting.  Additional information about the distribution of
multiplicities, beyond ordinary spacing or density bounds, is
therefore necessary to exclude this mechanism for zeta zeros.

\section{Boundaries of the rare-event reconstruction results}
\label{integrated:sec:rare-event-boundaries}

The results distinguish complete continuum responses, local curvature,
and entire microscopic fluctuation laws. A homogeneous response can
retain angular dependence while losing scale; an extra channel can
separate the dilation orbits. Scalar local curvature determines all
weight ratios, but finitely many jets do not determine an unrestricted
local law. A full Gamma-series fluctuation law determines its nonzero
coefficients, whereas finitely many cumulants can miss even the largest.

The noisy inverse problem remains separate from these uniqueness
statements. Proposition~\ref{p05:prop:inverse-cutoff} starts from an error
bound on mixed-derivative data. Controlling the preceding duality and
differentiation operations under a specified observation model is an
additional problem. Likewise, the mixed-probe inverse requires
quantitative control of the reconstruction of a signed measure from
its Laplace response before the exact channel threshold yields a
statistical guarantee. Analytic continuation from a patch is used
only for uniqueness.

Unknown scalar exponents are handled by
Theorem~\ref{p06:thm:blind-jet-calibration}, and the remaining coordinate
and normalization ambiguities are classified in
Theorem~\ref{p06:thm:blind-two-scale} and
Proposition~\ref{p06:prop:shape-multiplicity-ambiguity}. Genuinely
correlated vector laws lack the entrywise covariance positivity used
in that scalar argument; blind calibration for those classes requires
additional hypotheses. The repeated-component inverse also has no
unrestricted extension, as
Proposition~\ref{p06:prop:mp-wrong-repeat} demonstrates.

At $2(E-A)=\rho$, slowly varying counting factors can decide whether
the discrete coefficient ratios belong to $\ell^2$. A uniform
classification there is not furnished by the continuum activation
calculation. The hard-conditioning result uses the arcsine density
and a large prime-indexed independent reservoir; its extension needs
an appropriate local limit estimate, not merely a logarithmic
probability asymptotic. Finally, transferring a nondegenerate
fluctuation law to sampling by simple zeta zeros requires more than
a bounded comparison with product Haar measure. The sinc obstruction
shows why ordinary spacing and multiplicity-weighted counting bounds
alone do not close that gap.


\clearpage

% ===== Integrated Part 4 =====
\part{Controlled observation and blind calibration}
\label{part:controlled-observation}

\section{From observation design to hidden structure}
\label{p07:sec:introduction}The inverse problems of the preceding part motivate a change in the
observation law itself. A continuously varied Bernoulli gate first
recovers an ordered Laplace profile; arithmetic structure in that profile
then separates the original coordinate powers from the powers generated
by mixed moments. A positive interior baseline extends this mechanism to
arbitrary nonnegative summable laws, without assuming moments of the
original distribution. The finite-window normalization also permits an
unknown observation gain and a common unknown baseline to be recovered.

The reconstruction rests on three operations that remain distinct.
Compact moment inversion recovers a distribution of profile values.
Monotonicity restores their order in the observation parameter. Finally,
positivity of mixed moments identifies the original exponents as the
additively indecomposable nonzero atoms of an exponent measure. The
observation is a full function, not a single scalar value. Inactive
coordinates and uncalibrated coordinate amplitudes retain the explicit
ambiguities described below.

After the exact reconstruction results, the finite-data analysis proceeds
from noisy profile values to normalized gate derivatives and then to
ordinary gate values. These guarantees require quantitative visibility
conditions and recover dimension and powers; they do not provide uniform
recovery of the full law. Separated flat scales give a complementary
interior probe: its actual values determine every finite multivariate
measure, while all derivatives at its endpoint can fail to distinguish
smooth positive densities. Radial windows and adaptive linear queries
then give two different sharp information bounds.

The logarithmic and Stieltjes-transform relation and compact moment
uniqueness are classical
\cite{TulinoVerdu2004,Rao2007,Hausdorff1921}.
Known-power arithmetic encoding has antecedents in
\cite{Lin2016,Lenzi2026}; the distinction here is the blind recovery
of exponents and cardinality, followed by the explicitly stated
observation modifications and obstructions. Restricted-direction
determination is related to sharp Cram\'er--Wold theory
\cite{Cuesta2007}, but that theory observes projection laws rather than
one Laplace value at each point of a nonlinear curve. The later
quantitative arguments use classical Prony, approximation, convolution,
and information-width methods. Their hypotheses and precise recovery
targets are recorded with each result.

\section{Inverting a continuously varied Bernoulli gate}
\label{p07:sec:gate-tomography}

Let $X=(X_1,\ldots,X_d)$ be an almost surely finite nonnegative random
vector, and write
\[
 L_X(s)=\E\exp(-s\cdot X),\qquad s\in[0,\infty)^d.
\]
Fix $\rho>0$, powers $p_i>\rho$, and $a_i>0$. Put
$q_i=p_i/\rho$ and $\alpha_i=1/q_i$. A shared Bernoulli gate with
known success probability $\theta\in(0,1)$ replaces $X$ by $B_\theta X$,
where $B_\theta$ is independent of $X$. Its continuum response is
\begin{equation}\label{p07:eq:gate-response}
 G_a(\theta)=\int_0^\infty \rho y^{\rho-1}
 \left[-\log\left(1-\theta+
 \theta L_X(a_1y^{-p_1},\ldots,a_dy^{-p_d})\right)\right]\,dy.
\end{equation}
Only the gate probability is varied here; the amplitudes and powers remain
fixed. The substitution $r=y^\rho$ gives
\begin{equation}\label{p07:eq:gate-f}
 f_a(r)=1-L_X(a_1r^{-q_1},\ldots,a_dr^{-q_d}),\qquad
 G_a(\theta)=\int_0^\infty-\log(1-\theta f_a(r))\,dr.
\end{equation}
The usefulness of the gate is that its logarithm supplies every power of
$f_a$, rather than only one homogeneous average.

\begin{lemma}[Exact integrability condition]\label{p07:lem:gate-integrability}
The following are equivalent:
\begin{enumerate}[label=(\roman*)]
\item $G_a(\theta)<\infty$ for some $\theta\in(0,1)$;
\item $G_a(\theta)<\infty$ for every $\theta\in(0,1)$;
\item $\E X_i^{\alpha_i}<\infty$ for every $i$.
\end{enumerate}
More precisely,
\begin{equation}\label{p07:eq:gate-fractional-bounds}
 \max_i\left\{\Gamma(1-\alpha_i)a_i^{\alpha_i}
                   \E X_i^{\alpha_i}\right\}
 \leq\int_0^\infty f_a(r)\,dr
 \leq\sum_i\Gamma(1-\alpha_i)a_i^{\alpha_i}
                   \E X_i^{\alpha_i}.
\end{equation}
\end{lemma}
\begin{proof}
For nonnegative $z_i$,
\[
 \max_i(1-e^{-z_i})\leq1-e^{-\sum_i z_i}
 \leq\sum_i(1-e^{-z_i}).
\]
Tonelli's theorem and the change of variable $t=c r^{-q}$ give
\[
 \int_0^\infty(1-e^{-c r^{-q}})\,dr
 =\Gamma(1-1/q)c^{1/q},\qquad q>1,
\]
including $c=0$. This proves \eqref{p07:eq:gate-fractional-bounds}. Finally,
for $0\leq u\leq1$ and $0<\theta<1$,
\[
 \theta u\leq-\log(1-\theta u)
 \leq\frac{\theta}{1-\theta}u.
\]
\end{proof}

\begin{theorem}[Gate inversion on one anisotropic ray]
\label{p07:thm:gate-inversion}
Assume the equivalent conditions of Lemma~\ref{p07:lem:gate-integrability}.
The function $G_a$ on any nonempty open subinterval of $(0,1)$ determines
$f_a(r)$ for every $r>0$. It therefore determines the Laplace transform of
$X$ on the entire curve
\[
 \{(a_1r^{-q_1},\ldots,a_dr^{-q_d}):r>0\}.
\]
The same conclusion holds if $G_a$ is known on a set of distinct points
having an accumulation point in $[0,1)$.

An explicit inversion is as follows. Define the finite measure
\begin{equation}\label{p07:eq:gate-measure}
 \nu_a(A)=\int_0^\infty f_a(r)
                  \mathbf1_{\{f_a(r)\in A\}}\,dr,
 \qquad A\subset[0,1]\text{ Borel}.
\end{equation}
Its moments are
\begin{equation}\label{p07:eq:gate-moments}
 \int_0^1 u^n\,\nu_a(du)=\frac{G_a^{(n+1)}(0)}{n!},
 \qquad n\geq0.
\end{equation}
If $X$ is not almost surely zero, put $m=1-\Pp(X=0)$ and
\begin{equation}\label{p07:eq:gate-tail}
 R_a(t)=\int_{[t,1]}u^{-1}\,\nu_a(du),\qquad 0<t<m.
\end{equation}
Then $R_a$ is the ordinary inverse of the strictly decreasing continuous
function $f_a:(0,\infty)\to(0,m)$. In particular, $L_X(a)=1-t_a$,
where $t_a$ is the unique solution of $R_a(t_a)=1$.
\end{theorem}
\begin{proof}
For $|\theta|<1$, the series
\[
 -\log(1-\theta f_a(r))
 =\sum_{n\geq1}\frac{\theta^n f_a(r)^n}{n}
\]
converges after integration, locally uniformly in $\theta$, because
$0\leq f_a\leq1$ and $\int f_a<\infty$. Thus $G_a$ is holomorphic
in the unit disk and \eqref{p07:eq:gate-moments} follows. The identity theorem
proves the assertion concerning restricted gate data. Compactly supported
finite measures are determined by their moments: equality on polynomials
implies equality on continuous functions by uniform polynomial
approximation. Consequently the data determine $\nu_a$.

Suppose $\Pp(X\ne0)>0$. As $r$ increases, every nonzero realization of
$X$ strictly decreases $\sum_i a_iX_i r^{-q_i}$; hence $f_a$ is strictly
decreasing. Dominated convergence gives continuity and the endpoint limits
$f_a(0+)=m$, $f_a(\infty)=0$. The measure $\nu_a$ has no atoms,
since the preimage of a point under $f_a$ is at most a singleton. For
$0<t<m$, cancellation of the factor $f_a(r)$ gives
\[
 R_a(t)=\int_0^\infty\mathbf1_{\{f_a(r)\geq t\}}\,dr
       =f_a^{-1}(t).
\]
This proves every assertion. If $X=0$ almost surely then $G_a=0$; conversely,
vanishing $G_a$ implies $f_a=0$ almost everywhere and hence everywhere.
\end{proof}

\begin{corollary}[Recovery of an arbitrary correlated law]
\label{p07:cor:gate-all-a}
For fixed known $p_i,\rho$, the responses $G_a(\theta)$ for
$a\in(0,\infty)^d$ and $\theta$ in any nonempty open interval determine
the law of every nonnegative $X$ satisfying
$\E X_i^{\rho/p_i}<\infty$. No independence, infinite divisibility,
support gap, or atom at zero is required.
\end{corollary}
\begin{proof}
Theorem~\ref{p07:thm:gate-inversion} recovers $L_X(a)$ for every positive $a$.
Uniqueness of the multivariate Laplace transform concludes the proof.
\end{proof}

The moment inversion can be centered at a gate value that is actually
observed. This avoids making analytic continuation to zero part of the
formal inverse, although recovering all derivatives remains an exact-data
operation.

\begin{proposition}[Inversion from the full jet at one gate value]
\label{p07:prop:gate-local-jet}
Fix $\theta_0\in(0,1)$ and let
\[
 v_a(r)=\frac{f_a(r)}{1-\theta_0 f_a(r)},\qquad
 \nu_{a,\theta_0}(A)=\int_0^\infty v_a(r)
                         \mathbf1_{\{v_a(r)\in A\}}\,dr.
\]
This is a finite measure supported in $[0,(1-\theta_0)^{-1}]$, and
\begin{equation}\label{p07:eq:gate-centered-moments}
 \int v^n\,\nu_{a,\theta_0}(dv)
     =\frac{G_a^{(n+1)}(\theta_0)}{n!},\qquad n\geq0.
\end{equation}
The tail $\int_{[t,\infty)}v^{-1}\nu_{a,\theta_0}(dv)$ recovers
$v_a^{-1}(t)$, and $f_a=v_a/(1+\theta_0v_a)$. Thus the full derivative
sequence at one gate value determines the same ray data.
\end{proposition}
\begin{proof}
Differentiate under the integral, using
$v_a\leq f_a/(1-\theta_0)$ and a neighborhood of $\theta_0$ contained
in $(0,1)$. Then
$G_a^{(n+1)}(\theta_0)=n!\int v_a(r)^{n+1}\,dr$.
The rest repeats the compact moment and monotone inversion arguments.
\end{proof}

\begin{remark}[What is classical in the inverse]
The compact moment uniqueness and the relation between logarithmic
transforms and Stieltjes transforms are classical. For example,
\[
 \frac1zG_a'(1/z)=\int\frac{\nu_a(du)}{z-u},\qquad |z|>1.
\]
This is the Stieltjes transform of the finite measure $\nu_a$; compare the
Shannon-transform viewpoint in \cite{TulinoVerdu2004,Rao2007} and the
compact moment problem in \cite{Hausdorff1921}. The additional step here
is that the mixing variable is a strictly monotone Laplace profile.
Consequently its pushforward measure determines its order in $r$, and not
merely its unordered value distribution. Without monotonicity, arbitrary
measure-preserving rearrangements of $f_a$ leave $G_a$ unchanged.
\end{remark}

\subsection{No finite set of gate values is universally sufficient}

The continuum gate parameter cannot be replaced by any fixed finite list
on the full class of bounded scalar laws. Here ``a profile'' includes all
positive amplitudes, not just finitely many measured values.

\begin{theorem}[Finite-gate obstruction with moment calibration]
\label{p07:thm:finite-gate-obstruction}
Fix $0<b<B<\infty$, $p>\rho>0$, finitely many gate values
$\theta_j\in(0,1)$, and nonnegative integers $k_j$. Also fix an integer
$Q\geq0$. There are distinct finitely supported scalar probability laws
$\mu_+$ and $\mu_-$ supported in $[b,B]$ with the same first $Q$ ordinary
moments such that, for every $s>0$,
\begin{equation}\label{p07:eq:finite-gate-equality}
 \left.\partial_\theta^\ell G_s^{\mu_+}(\theta)\right|_{\theta=\theta_j}
 =\left.\partial_\theta^\ell G_s^{\mu_-}(\theta)\right|_{\theta=\theta_j},
 \qquad 0\leq\ell\leq k_j.
\end{equation}
Indeed, the common data have a positive-dimensional local family of
distinct laws. Thus two gate values, or any fixed finite number of gate
values and finite jets, do not identify arbitrary bounded scalar laws.
\end{theorem}
\begin{proof}
Write $\alpha=\rho/p$. Scaling in \eqref{p07:eq:gate-response} gives
\[
 G_s^\mu(\theta)=s^\alpha C_\mu(\theta),\qquad
 C_\mu(\theta)=\alpha\int_0^\infty
  -\log\{1-\theta+\theta L_\mu(u)\}\,u^{-\alpha-1}\,du.
\]
Consequently all the equalities in \eqref{p07:eq:finite-gate-equality} reduce
to $M=\sum_j(k_j+1)$ real scalar constraints.

Choose $N=M+Q+2$ distinct points $x_1,\ldots,x_N\in[b,B]$.
In the open probability simplex consider weights $w=(w_i)$ with
$\sum_iw_ix_i^q$ fixed at their uniform-weight values for $1\leq q\leq Q$.
Together with $\sum_iw_i=1$, these constraints have rank $Q+1$ by the
Vandermonde determinant. They define an open subset $U$ of an affine space
of dimension $N-Q-1=M+1$.

The map
\[
 F:U\longrightarrow\R^M,\qquad
 w\longmapsto\left(C_{\sum_iw_i\delta_{x_i}}^{(\ell)}(\theta_j)
              \right)_{j,\,0\leq\ell\leq k_j}
\]
is smooth. To justify differentiation, write any tangent perturbation as
$\sum_i h_i=0$. Then
$\sum_i h_i e^{-ux_i}=O(u)$ at zero and $O(e^{-bu})$ at infinity,
whereas the gated denominator is bounded below by $1-\theta_j>0$.
The same bounds apply after finitely many gate or weight derivatives.

Let $r$ be the maximum rank of $DF$ on $U$ and choose a point where that
rank is attained. Some $r$ by $r$ minor remains nonzero in a neighborhood,
and maximality makes the rank constant there. The constant-rank theorem
then gives local fibers of dimension $M+1-r\geq1$. Any two distinct
weights in such a fiber give the desired laws. Their finite supports are
the same, and all their weights can be kept strictly positive.
\end{proof}

\begin{remark}[Scope of the obstruction]
Embedding these scalar laws on a coordinate axis gives the same failure
for any fixed vector dimension and any prescribed coordinate powers,
provided the class allows laws supported on that axis. The theorem does
not assert that finitely many gate values fail in every more restrictive
model, nor that joint observations at additional distinct powers cannot
restore information. It concerns the stated observation architecture and
the unrestricted bounded-law class.
\end{remark}

\subsection{A scalar exponent can be calibrated from a gate curve}

\begin{proposition}[Blind scalar calibration under a finite mean]
\label{p07:prop:gate-scalar-calibration}
Let $d=1$, $a=1$, and suppose $0<\E X<\infty$. The function
$G_1(\theta)$ on any open gate interval determines both the unknown ratio
$q=p/\rho>1$ and the law of $X$. In terms of the recovered profile $f$,
\begin{equation}\label{p07:eq:gate-blind-q}
 q=-\lim_{r\to\infty}\frac{\log f(r)}{\log r},\qquad
 \E X=\lim_{r\to\infty}r^q f(r),\qquad
 L_X(s)=1-f(s^{-1/q}).
\end{equation}
\end{proposition}
\begin{proof}
The gate inverse reconstructs $f$ without using the exponent. The finite
mean assumption gives $1-L_X(s)\sim s\E X$ as $s\downarrow0$ and hence
$f(r)\sim(\E X)r^{-q}$. This proves the formulas.
\end{proof}

\begin{remark}[Why the mean assumption matters]
Without a tail restriction the exponent and the law can be confounded.
If $q'>q>1$, put $\beta=q/q'\in(0,1)$ and let $X'$ be the positive
$\beta$-stable random variable with Laplace transform
$L_{X'}(s)=e^{-s^\beta}$. The models $(X\equiv1,q)$ and $(X',q')$ have
the identical profile
$1-e^{-r^{-q}}$ and hence the identical gate response. Both satisfy the
fractional-moment condition needed above, while $X'$ has infinite mean.
The gate by itself therefore determines a Laplace profile before it
determines the exponent used to parameterize that profile.
\end{remark}

\section{The gate curve as a scalar rare-event observable}
\label{p07:sec:gate-ldp}

We give a direct probabilistic realization of the continuum data. This
also specifies exactly what information is meant by a measured gate
curve. Let $X_n=(X_{1,n},\ldots,X_{d,n})$, $n\geq1$, be independent
copies of a bounded nonnegative vector $X$ with $\Pp(X\ne0)>0$.
Fix $a_i>0$, $q_i>1$, and $\theta\in(0,1)$, and let $B_{\theta,n}$
be independent Bernoulli variables with success probability $\theta$,
independent of the vectors. For positive integers $N$ define the single
scalar statistic
\begin{equation}\label{p07:eq:gate-discrete-statistic}
 Z_N=\frac1N\sum_{n\geq1}B_{\theta,n}
                   \sum_{i=1}^d a_i(N/n)^{q_i}X_{i,n}.
\end{equation}
For each fixed $N$ this sum is bounded by a finite deterministic constant,
because $X$ is bounded and each $q_i>1$. Set
\begin{align}
 T(r)&=\sum_i a_i r^{-q_i}X_i,\notag\\
 h_\theta(\lambda,r)
   &=-\log\{1-\theta+\theta\E e^{-\lambda T(r)}\},\notag\\
 F_\theta(\lambda)&=\int_0^\infty h_\theta(\lambda,r)\,dr,
                \qquad\lambda>0.\label{p07:eq:gate-continuum-free-energy}
\end{align}
The function $F_\theta(1)$ is exactly the gate response $G_a(\theta)$
under the normalization $r=y^\rho$, $q_i=p_i/\rho$.

\begin{theorem}[Scalar lower-tail principle and dual recovery]
\label{p07:thm:gate-scalar-ldp}
Under the preceding assumptions, for every $\lambda>0$,
\begin{equation}\label{p07:eq:gate-laplace-limit}
 -\frac1N\log\E e^{-N\lambda Z_N}\longrightarrow F_\theta(\lambda).
\end{equation}
The function $F_\theta$ is twice continuously differentiable and strictly
concave on $(0,\infty)$, with
\[
 F_\theta'(0+)=\infty,\qquad F_\theta'(\infty)=0.
\]
Consequently, for each $e>0$ there is a unique $\lambda_e>0$ satisfying
$F_\theta'(\lambda_e)=e$, and
\begin{equation}\label{p07:eq:gate-rate-definition}
 I_\theta(e)=\sup_{\lambda>0}\{F_\theta(\lambda)-\lambda e\}
            =F_\theta(\lambda_e)-\lambda_e e
\end{equation}
is finite and continuous on $(0,\infty)$. For every $e>0$,
\begin{equation}\label{p07:eq:gate-lower-tail-ldp}
 \lim_{N\to\infty}\frac1N\log\Pp(Z_N\leq e)=-I_\theta(e).
\end{equation}
The full lower-tail rate function determines the gate response by
\begin{equation}\label{p07:eq:gate-rate-duality}
 G_a(\theta)=F_\theta(1)
             =\inf_{e>0}\{I_\theta(e)+e\}.
\end{equation}
Thus a family of exact scalar lower-tail rate functions, indexed by the
known gate probability, supplies the data used in the inverse theorems.
\end{theorem}

\begin{proof}
Independence and convergence of the defining sums give
\begin{equation}\label{p07:eq:gate-discrete-free-energy}
 F_{\theta,N}(\lambda):=-\frac1N\log\E e^{-N\lambda Z_N}
 =\frac1N\sum_{n\geq1}h_\theta(\lambda,n/N).
\end{equation}
Write $q_{\min}=\min_iq_i>1$. Uniformly for $\lambda$ in a compact
subinterval of $(0,\infty)$, the integrand is bounded at $0<r\leq1$
by $-\log(1-\theta)$ and is $O(r^{-q_{\min}})$ for $r\geq1$.
It is continuous in $r>0$. Splitting the integral and the Riemann sum
into $(0,\delta)$, $[\delta,R]$, and $(R,\infty)$ proves
\eqref{p07:eq:gate-laplace-limit}: the first part is $O(\delta+N^{-1})$,
the last is $O(R^{1-q_{\min}})$, and the middle part is an ordinary
Riemann sum.

At each $r$, tilt the law of $B_\theta T(r)$ by the density proportional
to $e^{-\lambda B_\theta T(r)}$. Its mean and variance satisfy
\begin{equation}\label{p07:eq:gate-integrand-derivatives}
 \partial_\lambda h_\theta(\lambda,r)
       =\E_{\lambda,r}[B_\theta T(r)],\qquad
 -\partial_\lambda^2 h_\theta(\lambda,r)
       =\operatorname{Var}_{\lambda,r}(B_\theta T(r)).
\end{equation}
For $j=1,2$, the tilted $j$th raw moment is bounded by
\[
 \frac{\theta}{1-\theta}\E[T(r)^j e^{-\lambda T(r)}].
\]
On $0<r\leq1$ this is bounded uniformly for $\lambda$ in compact
positive intervals, using $\sup_{t\geq0}t^je^{-\lambda t}<\infty$.
On $r\geq1$, boundedness of $X$ makes it $O(r^{-j q_{\min}})$.
These bounds justify differentiating the continuum integral and show,
by the same three-part Riemann-sum argument, that
\begin{equation}\label{p07:eq:gate-derivative-convergence}
 F_{\theta,N}^{(j)}(\lambda)\longrightarrow
 F_\theta^{(j)}(\lambda),\qquad j=1,2.
\end{equation}
The variance in \eqref{p07:eq:gate-integrand-derivatives} is strictly
positive for every $r$: the Bernoulli variable supplies a zero atom,
while $\Pp(T(r)>0)>0$. Hence $F_\theta''(\lambda)<0$.

As $\lambda\downarrow0$, the tilted mean increases to
$\theta\E T(r)$, whose integral over $0<r<1$ is infinite because
at least one $\E X_i$ is positive and $q_i>1$. Monotone convergence
gives $F_\theta'(0+)=\infty$. As $\lambda\to\infty$, the tilted
mean tends pointwise to zero, since the denominator stays above
$1-\theta$. It is bounded by its value at any fixed positive tilt,
an integrable function. Dominated convergence gives
$F_\theta'(\infty)=0$. These facts establish the unique maximizer in
\eqref{p07:eq:gate-rate-definition}. Continuity follows from the strict
monotonicity and continuity of $F_\theta'$ and its inverse. Also
$F_\theta(0+)=0$ by dominated convergence, so no missing endpoint
term is involved in the supremum.

For the upper bound in \eqref{p07:eq:gate-lower-tail-ldp}, Markov's
inequality gives, for every $\lambda>0$,
\[
 \Pp(Z_N\leq e)\leq
 \exp\{N\lambda e-NF_{\theta,N}(\lambda)\}.
\]
Taking the limit and optimizing proves the required upper bound.
For the lower bound, fix $0<\delta<e$ and let $\lambda_\delta$ be
the unique tilt with $F_\theta'(\lambda_\delta)=e-\delta$.
Define the tilted law by
\[
 \frac{d\mathbb Q_{N,\delta}}{d\mathbb P}
 =\exp\{-N\lambda_\delta Z_N+NF_{\theta,N}(\lambda_\delta)\}.
\]
Equations \eqref{p07:eq:gate-derivative-convergence} and elementary
differentiation of the log-Laplace transform show that
\[
 \E_{\mathbb Q_{N,\delta}}Z_N\longrightarrow e-\delta,
 \qquad
 \operatorname{Var}_{\mathbb Q_{N,\delta}}Z_N
       =-N^{-1}F_{\theta,N}''(\lambda_\delta)=O(N^{-1}).
\]
Therefore, for every $0<\varepsilon<\min\{\delta,e-\delta\}$,
the event $A_N=\{|Z_N-(e-\delta)|<\varepsilon\}$ has tilted
probability tending to one and is contained in $\{Z_N\leq e\}$.
Changing measure on $A_N$ yields
\[
 \Pp(Z_N\leq e)\geq
 \exp\{-NF_{\theta,N}(\lambda_\delta)
       +N\lambda_\delta(e-\delta-\varepsilon)\}
 \mathbb Q_{N,\delta}(A_N).
\]
Take the logarithmic limit, then let $\varepsilon\downarrow0$ and
$\delta\downarrow0$. The result is $-I_\theta(e)$ by continuity,
which proves the lower bound.

Finally, the definition of $I_\theta$ gives
$I_\theta(e)+\lambda e\geq F_\theta(\lambda)$ for every $e>0$,
with equality at $e=F_\theta'(\lambda)$. Taking $\lambda=1$ proves
\eqref{p07:eq:gate-rate-duality}.
\end{proof}

\section{A sharp arithmetic criterion for one aggregate}
\label{p07:sec:gate-resonance}

We now keep the positive coefficient vector $a=(a_1,\ldots,a_d)$ fixed.
The observed statistic is a single aggregate with coordinate-dependent
spatial powers.  Its gate profile is
\begin{equation}\label{p07:eq:res-gate}
 G_{\mu,a,p}(\theta)=\int_0^\infty \rho y^{\rho-1}
 \left\{-\log\left[1-\theta+
 \theta L_\mu(a_1y^{-p_1},\ldots,a_dy^{-p_d})\right]\right\}\,dy,
 \qquad 0<\theta<1,
\end{equation}
where $p_i>\rho$ and $L_\mu(s)=\int e^{-s\cdot x}\mu(dx)$.
The gate inversion established above says that this profile determines
\begin{equation}\label{p07:eq:res-curve}
 F_\mu(t)=L_\mu(a_1t^{p_1},\ldots,a_dt^{p_d}),\qquad t>0.
\end{equation}
It is therefore meaningful to ask when this one-dimensional restriction
determines a genuinely multivariate law. The positive coefficient-encoding
argument below is a finite-dimensional instance of the prescribed-index
method in \cite[Lemmas~3.2--3.3]{Lenzi2026}; it is reproduced to make
the sharp criterion and its converse self-contained.

\begin{theorem}[Exact resonance criterion]\label{p07:thm:resonance-criterion}
Fix $a\in(0,\infty)^d$, $p\in(\rho,\infty)^d$, and a nonempty open
interval $J\subset(0,1)$.  The following are equivalent.
\begin{enumerate}
\item The numbers $p_1,\ldots,p_d$ are linearly independent over
$\mathbb Q$.
\item The map $\mu\mapsto G_{\mu,a,p}|_J$ is injective on the class of
compactly supported probability measures on $[0,\infty)^d$.
\item The map $\mu\mapsto F_\mu$ is injective on that class.
\end{enumerate}
If rational independence fails, then, in every prescribed nonempty open
box $B\subset(0,\infty)^d$ and for every integer $M\geq0$, there are
distinct probability measures $\mu_+$ and $\mu_-$ with smooth compactly
supported densities in $B$ such that their gate profiles agree for all
$0<\theta<1$ and
\begin{equation}\label{p07:eq:res-moments}
 \int x^n\,\mu_+(dx)=\int x^n\,\mu_-(dx)
 \qquad (n\in\mathbb N_0^d,\ |n|\leq M).
\end{equation}
\end{theorem}

\begin{proof}
For a compactly supported law, the multivariate Laplace transform is
entire and, for $t>0$,
\begin{equation}\label{p07:eq:res-series}
 F_\mu(t)=\sum_{n\in\mathbb N_0^d}
 \frac{(-1)^{|n|}a^n}{n!}\left(\int x^n\mu(dx)\right)t^{p\cdot n}.
\end{equation}
This series converges absolutely.  Since $p_i>0$, the set of indices
with $p\cdot n\leq R$ is finite for every $R$.  Under rational
independence, all the exponents $p\cdot n$ are distinct.  We can
therefore list them as $0=\lambda_0<\lambda_1<\cdots\to\infty$.
Writing $F_\mu(t)=\sum_{k\geq0}c_kt^{\lambda_k}$, each coefficient is
recovered recursively by
\begin{equation}\label{p07:eq:res-extract}
 c_k=\lim_{t\downarrow0}t^{-\lambda_k}
 \left[F_\mu(t)-\sum_{j<k}c_jt^{\lambda_j}\right].
\end{equation}
To justify this limit without any separation estimate on adjacent
exponents, choose an integer $N$ with $N\min_i p_i>\lambda_k$.
The Taylor remainder containing terms of total degree at least $N$ is
$O(t^{N\min_i p_i})$ for $0<t\leq1$; only finitely many other terms
remain.  Formula \eqref{p07:eq:res-extract} follows.
Thus $F_\mu$ determines all mixed moments.  If two compactly supported
laws have the same moments, put both supports in one compact box and
apply polynomial approximation on that box to obtain equality of the
laws.  This proves the positive direction.  The previously established
gate inverse and analyticity in $\theta$ make the observations on $J$
equivalent to the full curve \eqref{p07:eq:res-curve}.

For the converse, rational dependence and positivity of the $p_i$
give distinct nonzero multi-indices $\alpha,\beta\in\mathbb N_0^d$
such that $p\cdot\alpha=p\cdot\beta$.  Define the nonzero polynomial
\begin{equation}\label{p07:eq:res-binomial}
 P(s)=a^{-\alpha}s^\alpha-a^{-\beta}s^\beta.
\end{equation}
It vanishes on $s=(a_1t^{p_1},\ldots,a_dt^{p_d})$.
Choose a nonzero nonnegative $\varphi\in C_c^\infty(B)$ and set
\begin{equation}\label{p07:eq:res-null-density}
 Q(s)=s_1^{M+1}P(s),\qquad g=Q(\partial)\varphi.
\end{equation}
Repeated integration by parts gives
\begin{equation}\label{p07:eq:res-laplace-kernel}
 \int e^{-s\cdot x}g(x)\,dx
   =Q(s)\int e^{-s\cdot x}\varphi(x)\,dx.
\end{equation}
The last Laplace transform is strictly positive for real positive $s$,
and $Q$ is a nonzero polynomial.  Consequently $g\not\equiv0$.
Every monomial of $Q$ has degree greater than $M$, so integration by
parts also gives $\int x^n g(x)\,dx=0$ for $|n|\leq M$, including
zero total mass.  Take a smooth compactly supported probability
density $f_0$ in $B$ which is strictly positive on
$\operatorname{supp}\varphi$.  Since
$\operatorname{supp}g\subset\operatorname{supp}\varphi$, the
functions $f_0\pm\varepsilon g$ are nonnegative probability densities
for all sufficiently small $\varepsilon>0$.  Their difference is
nonzero, all the moments in \eqref{p07:eq:res-moments} agree, and
\eqref{p07:eq:res-laplace-kernel} shows equality of their Laplace transforms
on the entire curve.  Their gate profiles therefore agree as well.
\end{proof}

\begin{corollary}[Moment-determinate laws]\label{p07:cor:res-moment-determinate}
The injectivity conclusion for rationally independent powers also
holds on the class of nonnegative laws having all mixed moments finite
and determined by those moments.  Compact support and a finite joint
exponential moment are sufficient, but are not necessary, hypotheses.
\end{corollary}

\begin{proof}
For a fixed integer $N$, Taylor's theorem on the nonnegative real axis
gives
\[
 \left|e^{-z}-\sum_{k=0}^N\frac{(-z)^k}{k!}\right|
 \leq\frac{z^{N+1}}{(N+1)!},\qquad z\geq0.
\]
With $z=\sum_i a_it^{p_i}X_i$ and $0<t\leq1$, taking expectations
shows that the terms with $|n|\leq N$ in
\eqref{p07:eq:res-series} have a remainder bounded by
\[
 \frac{\mathbb E(\sum_i a_iX_i)^{N+1}}{(N+1)!}
 t^{(N+1)\min_i p_i}.
\]
This finite-order expansion justifies
\eqref{p07:eq:res-extract} at every prescribed exponent, without
requiring convergence of the infinite moment series.
Thus equal gate profiles imply equal mixed moments, and the assumed
moment determinacy identifies the law.  The gate integrability follows
already from finite first moments and $p_i>\rho$.
\end{proof}

\begin{remark}[Observation architecture]\label{p07:rem:res-architecture}
The aggregate in \eqref{p07:eq:res-gate} assigns the distinct weights
$y^{-p_i}$ to the different coordinates before summing.  It is not a
channel of the form $y^{-p}\ell(X)$ with one common exponent and one
fixed linear functional $\ell$.  Moreover the data contain an exact
continuum of gate probabilities.  Thus the theorem neither recovers a
multivariate law from a single scalar number nor contradicts lower
bounds for finitely many fixed linear channels with a fixed gate.
\end{remark}

\begin{corollary}[Finite collections of resonant designs]
\label{p07:cor:res-finite-designs}
Fix finitely many designs $(a^{(j)},p^{(j)})$, $1\leq j\leq m$, each
with positive coefficients and powers greater than $\rho$.
Their combined full gate profiles identify every compactly supported
nonnegative $d$-dimensional law if and only if at least one vector
$p^{(j)}$ has rationally independent coordinates.  If all the designs
are resonant, the indistinguishable laws can have smooth densities
in any prescribed positive box and matching mixed moments through
any prescribed finite total degree.
\end{corollary}

\begin{proof}
Sufficiency follows from Theorem~\ref{p07:thm:resonance-criterion}.
If every design is resonant, construct its polynomial $P_j$ as in
\eqref{p07:eq:res-binomial} and use
$Q(s)=s_1^{M+1}\prod_{j=1}^m P_j(s)$ in
\eqref{p07:eq:res-null-density}.  The product is nonzero and vanishes
on every observed curve.  The same density perturbation proves all
the assertions.
\end{proof}

\begin{example}[An irrational design and a resonant design]
For $\rho=1$ and $d=2$, the exponent vector $(2,2\sqrt2)$ is
rationally independent, so one full gate profile with any fixed
$a_1,a_2>0$ determines every compactly supported joint law.
For $(2,4)$ the polynomial
$s_1^2/a_1^2-s_2/a_2$ generates invisible smooth perturbations.
Merely changing the coefficients a finite number of times does not
remove this obstruction.  Pairwise irrational ratios are not enough
in dimension three: $(2,2\sqrt2,2+2\sqrt2)$ has a rational relation.
\end{example}

\subsection{Instability near a resonance}

Exact injectivity has no uniform quantitative meaning as a design
approaches a rational relation.  The next estimate makes this precise
without asserting a minimax rate for a fixed design.

\begin{proposition}[Deterioration near resonant powers]
\label{p07:prop:res-stability}
Fix $a\in(0,\infty)^d$, $R>0$, $\eta>0$, and
$0<\theta_*<1$.  For every law $\mu$ supported in $[0,R]^d$ and
every $p,q\in[\rho+\eta,\infty)^d$,
\begin{equation}\label{p07:eq:res-lipschitz-design}
 \sup_{0\leq\theta\leq\theta_*}
 |G_{\mu,a,p}(\theta)-G_{\mu,a,q}(\theta)|
 \leq\sum_{i=1}^d C_i|p_i-q_i|,
 \quad
 C_i=\frac{\theta_*}{1-\theta_*}
 \left(\frac1{e\rho}+\frac{\rho a_iR}{\eta^2}\right).
\end{equation}
In particular, for every resonant $p^0\in(\rho,\infty)^d$ there
are fixed distinct smooth compactly supported laws $\mu_\pm$ and
a sequence of rationally independent vectors $p^{(n)}\to p^0$ such
that their gate data differ by $O(\|p^{(n)}-p^0\|_1)$ on
$[0,\theta_*]$, although their Wasserstein distance is a fixed
strictly positive number.
\end{proposition}

\begin{proof}
Write $A=\sum_i a_iX_i y^{-p_i}$ and
$u=1-\mathbb E e^{-A}$.  Differentiation gives
\[
 |\partial_{p_i}u|
 \leq |\log y|\,\mathbb E[a_iX_i y^{-p_i}e^{-A}].
\]
For $0<y\leq1$, the expectation is at most
$\mathbb E[Ae^{-A}]\leq e^{-1}$.
For $y\geq1$, it is at most $a_iR y^{-p_i}$.
Also $\theta/(1-\theta u)\leq\theta_* /(1-\theta_*)$.
The two needed integrals are
\[
 \int_0^1\rho y^{\rho-1}|\log y|\,dy=\rho^{-1},
 \qquad
 \int_1^\infty\rho y^{\rho-1-p_i}\log y\,dy
       =\frac{\rho}{(p_i-\rho)^2}.
\]
These bounds justify differentiation under the integral and give
\eqref{p07:eq:res-lipschitz-design} by integrating along the segment
joining $p$ and $q$.
Choose $\mu_\pm$ from Theorem~\ref{p07:thm:resonance-criterion}
for $p^0$.  Their gate data agree at $p^0$, so apply
\eqref{p07:eq:res-lipschitz-design} twice.  Rationally independent vectors
are dense, being the complement of countably many proper hyperplanes.
Finally $W_1(\mu_+,\mu_-)>0$ because the two compactly supported
probability measures are distinct.
\end{proof}

Thus no inverse modulus $W_1\leq\omega(\|\Delta G\|)$ with
$\omega(\varepsilon)\to0$ can hold uniformly over rationally
independent designs approaching a resonance.  In particular, a
H\"older estimate cannot have a uniform constant there.
The rational independence criterion is an exact identifiability
result, not a claim of numerical robustness.

\subsection{Why a tail hypothesis is essential}

\begin{proposition}[Stable-law aliasing]\label{p07:prop:res-stable-alias}
The compact-support assumption in
Theorem~\ref{p07:thm:resonance-criterion} cannot be replaced by the
fractional-moment conditions needed to make the gate profile finite.
For any $p_1>p_2>\rho$ and any $a_1,a_2>0$, there are distinct
nonnegative two-dimensional laws, one bounded and one unbounded,
with identical full gate profiles for the design $(a,p)$.
This remains true when $p_1/p_2$ is irrational.
\end{proposition}

\begin{proof}
Put $\alpha=p_2/p_1\in(0,1)$, and let $S$ be the positive
$\alpha$-stable random variable normalized by
$\mathbb E e^{-sS}=e^{-s^\alpha}$.
For completeness, its existence follows from the subordinator
L\'evy measure
$\alpha x^{-1-\alpha}\,dx/\Gamma(1-\alpha)$, whose Laplace
exponent is $s^\alpha$.
Set $c=a_2/a_1^\alpha$ and
\[
 X=(c^{1/\alpha}S,0),\qquad Y=(0,1).
\]
Then for every $t>0$,
\[
 \mathbb E e^{-a_1t^{p_1}X_1-a_2t^{p_2}X_2}
 =e^{-c(a_1t^{p_1})^\alpha}
 =e^{-a_2t^{p_2}}
 =\mathbb E e^{-a_1t^{p_1}Y_1-a_2t^{p_2}Y_2}.
\]
Hence the gate profiles coincide.
Both profiles are finite since their common curve satisfies
$1-e^{-a_2y^{-p_2}}=O(y^{-p_2})$ at infinity and $p_2>\rho$.
They also satisfy the natural coordinatewise fractional-moment
conditions: $\mathbb E S^r<\infty$ for $0<r<\alpha$,
whereas $\rho/p_1<\alpha$.
The moment assertion follows directly from
\[
 \mathbb E S^r
 =\frac{r}{\Gamma(1-r)}\int_0^\infty
       (1-e^{-s^\alpha})s^{-r-1}\,ds,
 \qquad 0<r<\alpha.
\]
Adding the same deterministic vector in $(0,\infty)^2$ to $X$
and $Y$ gives an example supported strictly inside the positive
orthant, with the same conclusion.
\end{proof}

Corollary~\ref{p07:cor:res-moment-determinate} allows tails well beyond
compact support.  The stable example isolates a different failure
mechanism: a fractional tail power can imitate another coordinate's
spatial power when the required higher moments do not exist.

\subsection{An explicit pair and its near-resonant measurements}
\label{p07:sec:explicit-collision}
The following pair gives an elementary concrete instance of the preceding
phenomenon. On $[1,2]^2$, put $t_i=x_i-1$,
\[
 \phi(t)=t^3(1-t)^3,\qquad
 g(x)=\phi''(t_1)\phi(t_2)-\phi(t_1)\phi'(t_2),
 \qquad f_\pm(x)=1\pm32g(x).
\]
Extend these densities by zero outside the box. They are probability
densities: $\int g=0$, and
$\|\phi\|_\infty\le1/64$,
$\|\phi'\|_\infty\le3/16$,
$\|\phi''\|_\infty\le3/8$ imply $|32g|\le9/32<1$.
Their Laplace-transform difference is exactly
\[
 L_+(s_1,s_2)-L_-(s_1,s_2)
 =64(s_1^2-s_2)\Phi(s_1)\Phi(s_2),\qquad
 \Phi(s)=\int_1^2 e^{-sx}\phi(x-1)\,dx.
\]
Integration by parts justifies this identity, since $\phi$ and its
first derivative vanish at both endpoints. Thus at $\rho=1$, $a=(1,1)$,
and $p=(2,4)$, their complete gate profiles agree exactly.
They nevertheless satisfy the fixed separation
\[
 W_1(\mu_+,\mu_-)\ge
 |\mathbb E_+X_2-\mathbb E_-X_2|
 =\frac{64}{19600}>0,
\]
because $x\mapsto x_2$ is $1$-Lipschitz and $\int_0^1\phi=1/140$.

For the irrational designs $p=(2,4+\delta)$ with
$\delta=\sqrt2\,10^{-k}$, exact uniqueness holds, but these fixed
laws become harder to distinguish as $k$ grows. The table gives
deterministic quadrature values; it is an illustration, not an input to
the proof or a statistical experiment.
\begin{center}
\begin{tabular}{cc}
\toprule
$\delta$ & $|G_+(1/2)-G_-(1/2)|$\\
\midrule
$\sqrt{2}\,10^{-1}$ & $1.559669\times10^{-5}$ \\
$\sqrt{2}\,10^{-2}$ & $1.631627\times10^{-6}$ \\
$\sqrt{2}\,10^{-3}$ & $1.639186\times10^{-7}$ \\
$\sqrt{2}\,10^{-4}$ & $1.639946\times10^{-8}$ \\
\bottomrule
\end{tabular}
\end{center}
The computation integrates in $z=\log y$ and uses the stable formula
$-2e^z\operatorname{arctanh}(\theta D/(1-\theta+\theta L_0))$,
where $L_0$ is the Laplace transform of the uniform law on the box
and $D=32(s_1^2-s_2)\Phi(s_1)\Phi(s_2)$.
An 80-node Gauss--Legendre rule evaluates $\Phi$;
adaptive quadrature uses $-6\le z\le24$.
The reported quadrature error estimates are below $10^{-12}$.
All equality and separation assertions above follow from the exact
formulas, independently of numerical error estimates.

\section{Blind recovery of dimension, powers, and the correlated law}
\label{p07:sec:blind-gate-arithmetic}

The arithmetic criterion has a stronger consequence when the coordinate
powers are not known in advance. Nonnegativity makes the pattern of
nonzero mixed moments hereditary under taking sub-multi-indices. This
allows the primitive spatial powers to be read from the asymptotic
exponents themselves.

For a subset $S\subset(0,\infty)$ define its set of \emph{additively
indecomposable elements} by
\begin{equation}\label{p07:eq:blind-indecomposables}
 \mathcal A(S)=\{s\in S:s\ne u+v\text{ for every }u,v\in S\}.
\end{equation}
No closure of $S$ under addition is assumed. That qualification matters
because a joint moment can vanish when the corresponding coordinates are
never simultaneously positive.

\begin{lemma}[Primitive exponents from positive moments]
\label{p07:lem:blind-primitive-exponents}
Let $q_1,\ldots,q_d>0$ be linearly independent over $\mathbb Q$.
Let $X\in[0,\infty)^d$ have finite mixed moments of every order and
$\Pp(X_i>0)>0$ for each $i$. Define
\begin{equation}\label{p07:eq:blind-support}
 S_X(q)=\{q\cdot n:n\in\mathbb N_0^d\setminus\{0\},
                           \E X^n>0\}.
\end{equation}
Then
\begin{equation}\label{p07:eq:blind-atoms}
 \mathcal A(S_X(q))=\{q_1,\ldots,q_d\}.
\end{equation}
The same conclusion holds for a finite or countably infinite index set
and a nonzero finite positive measure in place of the probability law,
provided all finite mixed moments exist and every coordinate has positive
first moment. In that form all multi-indices have finite support, and
linear independence means independence of every finite subfamily.
\end{lemma}
\begin{proof}
The first moments are strictly positive, so $q_i\in S_X(q)$.
If $q_i=q\cdot m+q\cdot n$ for nonzero $m,n\in\mathbb N_0^d$,
rational independence implies $e_i=m+n$, which is impossible.
Thus every $q_i$ is indecomposable.

Conversely, suppose $s=q\cdot n\in S_X(q)$ and $|n|\geq2$.
Choose $i$ with $n_i>0$. Since $X^n\geq0$ and $\E X^n>0$,
the event $\{X^n>0\}$ has positive probability. On this event both
$X_i$ and $X^{n-e_i}$ are positive. Their expectations are therefore
strictly positive, so both $q_i$ and $q\cdot(n-e_i)$ belong to
$S_X(q)$. They sum to $s$, which is consequently decomposable. The argument involves only the
finitely many coordinates in the support of the relevant multi-index;
it therefore proves the countable and finite-positive-measure versions
as well, with expectations replaced by integrals.
\end{proof}

\begin{theorem}[Fully blind identification from one gate curve]
\label{p07:thm:blind-gate-identification}
Consider the class of models $(d,q,\mu)$ satisfying:
\begin{enumerate}[label=(\roman*)]
\item $d$ is an unknown finite positive integer;
\item $q=(q_1,\ldots,q_d)\in(1,\infty)^d$ has coordinates linearly
      independent over $\mathbb Q$;
\item $\mu$ is a probability law on $[0,\infty)^d$ with all mixed
      moments finite, determined by those moments;
\item every coordinate is nontrivial: $\mu\{x:x_i>0\}>0$.
\end{enumerate}
For $X\sim\mu$, observe only
\begin{equation}\label{p07:eq:blind-gate-data}
 G(\theta)=\int_0^\infty
   -\log\left[1-\theta+
      \theta\E\exp\left(-\sum_{i=1}^dX_i r^{-q_i}\right)\right]\,dr
\end{equation}
on a nonempty open interval of $\theta\in(0,1)$, with the displayed
normalization of $dr$ known. Then $G$ determines $d$, the unordered
set $\{q_1,\ldots,q_d\}$, and the joint law $\mu$, up to the same
permutation of powers and coordinates. In particular, the result holds
for every bounded nonnegative vector with nontrivial coordinates.

More explicitly, the gate inverse recovers
\begin{equation}\label{p07:eq:blind-curve}
 F(t)=\E\exp\left(-\sum_{i=1}^dX_i t^{q_i}\right),\qquad t>0.
\end{equation}
Let $S(F)$ be the positive support of its generalized asymptotic
expansion at zero. Then
\begin{equation}\label{p07:eq:blind-recovery}
 \{q_1,\ldots,q_d\}=\mathcal A(S(F)),\qquad
 d=|\mathcal A(S(F))|.
\end{equation}
If $c_{q\cdot n}$ denotes the coefficient at exponent $q\cdot n$,
with coefficient zero when that exponent is absent, then
\begin{equation}\label{p07:eq:blind-moment-recovery}
 \E X^n=(-1)^{|n|}n!\,c_{q\cdot n},\qquad n\in\mathbb N_0^d.
\end{equation}
\end{theorem}
\begin{proof}
Finite first moments imply the fractional-moment condition in
Lemma~\ref{p07:lem:gate-integrability}. The gate inversion applies directly
to \eqref{p07:eq:blind-gate-data}; it does not require the values of $d$ or
$q$. It determines the strictly decreasing function
$f(r)=1-F(1/r)$, hence determines $F$.

We make the asymptotic-support assertion precise. For every integer
$N\geq0$, Taylor's theorem and nonnegativity give, for $0<t\leq1$,
\begin{equation}\label{p07:eq:blind-finite-expansion}
 F(t)=\sum_{|n|\leq N}\frac{(-1)^{|n|}}{n!}\E X^n\,t^{q\cdot n}
       +O\left(t^{(N+1)q_{\min}}\right),
 \qquad q_{\min}=\min_i q_i.
\end{equation}
Indeed, the absolute remainder is at most
\[
 \frac{\E(\sum_i X_i)^{N+1}}{(N+1)!}
             t^{(N+1)q_{\min}}.
\]
The set $\{q\cdot n:n\in\mathbb N_0^d\}$ is locally finite, because
all $q_i$ are strictly positive. Rational independence makes all these
exponents distinct. Consequently the nonzero asymptotic coefficients
of $F-1$ have precisely the support $S_X(q)$ in
\eqref{p07:eq:blind-support}; no coefficients can cancel at a common exponent.

This support and its coefficients are intrinsic, even though $q$ is
initially unknown. To see this, compare two candidate models satisfying
the hypotheses. The union of their two locally finite exponent sets is
again locally finite. If their coefficient arrays differ, there is a
smallest exponent $\lambda$ where they differ. Choose the Taylor order
in each model so that its remainder in
\eqref{p07:eq:blind-finite-expansion} has order strictly greater than
$\lambda$. Subtract the finite expansions and divide by $t^\lambda$.
As $t\downarrow0$, the difference tends to the nonzero difference of
the two coefficients. This contradicts equality of their functions $F$.
Thus both models have the same asymptotic support and the same
coefficients on it.

Lemma~\ref{p07:lem:blind-primitive-exponents} now identifies the power set
from that support and gives \eqref{p07:eq:blind-recovery}. Relabel the
coordinates so that the powers are in increasing order. Rational
independence identifies each multi-index $n$ from its exponent
$q\cdot n$, and \eqref{p07:eq:blind-moment-recovery} recovers every mixed
moment. The moment-determinacy hypothesis concludes the proof.
\end{proof}

\begin{example}[Mutually exclusive coordinates]
Suppose $d=2$, $q_1/q_2\notin\mathbb Q$, and
$X=(U,0)$ with probability $1/2$ and $X=(0,V)$ with probability
$1/2$, where $U,V$ are strictly positive bounded variables. Every mixed
moment involving both coordinates vanishes. The observed support is
\[
 S(F)=\{nq_1:n\geq1\}\cup\{nq_2:n\geq1\},
\]
which is not closed under addition. Nevertheless its only
indecomposable elements are $q_1$ and $q_2$. The missing mixed
coefficients recover the mutual exclusion, so the theorem does not
require a positive joint density or simultaneous positivity.
\end{example}

\begin{remark}[The remaining gauges and necessary qualifications]
An identically zero coordinate can always be appended without changing
the data; the nontrivial-coordinate condition removes this unavoidable
dimension ambiguity. If unknown positive amplitudes $a_i$ are inserted
in \eqref{p07:eq:blind-gate-data}, the identified variables are $a_iX_i$;
their amplitudes and marginal scales cannot be separated without an
external calibration. In the original variables $q_i=p_i/\rho$, the
curve identifies these ratios and not the common absolute scale of
$p_1,\ldots,p_d,\rho$.

The normalization of the continuum measure is also material. If the
unknown factor $C>0$ multiplies $dr$, substitution $r'=Cr$ gives
\[
 C\int_0^\infty-\log\left[1-\theta+
     \theta\E e^{-\sum_iX_i r^{-q_i}}\right]dr
 =\int_0^\infty-\log\left[1-\theta+
     \theta\E e^{-\sum_i(C^{q_i}X_i)(r')^{-q_i}}\right]dr'.
\]
Thus an unknown intensity cannot be separated from the corresponding
anisotropic rescaling of the law. Finally, the theorem uses the exact
full gate curve and arbitrarily deep asymptotics; it is not a finite
sample procedure or a stable dimension-selection algorithm.
\end{remark}

\begin{proposition}[Dimension is not uniformly stable]
\label{p07:prop:blind-dimension-instability}
Fix an admissible bounded model $(d,q,X)$ from
Theorem~\ref{p07:thm:blind-gate-identification} and $0<\theta_*<1$.
Choose $q_*>1$ such that $q_1,\ldots,q_d,q_*$ are rationally
independent. For $\varepsilon>0$, append the deterministic coordinate
$X_{d+1}=\varepsilon$ with power $q_*$. The enlarged model is admissible,
has dimension $d+1$, and its gate curve $G_\varepsilon$ satisfies
\begin{equation}\label{p07:eq:blind-dimension-bound}
 \sup_{0\leq\theta\leq\theta_*}
 |G_\varepsilon(\theta)-G(\theta)|
 \leq\frac{\theta_*}{1-\theta_*}
       \Gamma(1-1/q_*)\varepsilon^{1/q_*}.
\end{equation}
Therefore latent dimension has no uniform continuity in the full gate
data on such intervals, even when the powers are fixed and no resonance
is approached.
\end{proposition}
\begin{proof}
Writing $L(r)=\E\exp(-\sum_iX_i r^{-q_i})$, the new Laplace profile is
$L(r)e^{-\varepsilon r^{-q_*}}$. Its difference from $L(r)$ is at most
$1-e^{-\varepsilon r^{-q_*}}$. The derivative in $u\in[0,1]$ of
$-\log(1-\theta u)$ is at most $\theta_* /(1-\theta_*)$.
Integrating and applying the scalar identity in the proof of
Lemma~\ref{p07:lem:gate-integrability} proves
\eqref{p07:eq:blind-dimension-bound}. Boundedness implies moment determinacy,
and all the coordinates of the enlarged model are nontrivial.
\end{proof}

\section{Blind recovery with accumulating exponent spectra}
\label{p07:sec:countable-blind-gate}

The finite-dimensional reconstruction admits a countable-dimensional
extension. In particular, an accumulation of weighting exponents does not
destroy exact identification. The appropriate intermediate object is a
finite signed measure on the exponent axis, rather than a generalized
series that can be ordered by successive smallest exponents.

We use the signed-exponent-measure encoding already employed by
Lenzi~\cite[Lemmas~3.2--3.3]{Lenzi2026} for prescribed rationally independent
indices and the recovery of fractional power-sum profiles. The further
step here identifies the unknown indices themselves: they are precisely
the additively indecomposable nonzero atoms of that measure. Positivity
of the latent coordinates is essential to this step.

Let $J$ be a nonempty finite or countably infinite index set. Write
$\ell^1_+(J)$ for the nonnegative summable families indexed by $J$,
equipped with its norm Borel $\sigma$-algebra. Consider a family
$q=(q_j)_{j\in J}$ and an $\ell^1_+(J)$-valued random variable
$X=(X_j)_{j\in J}$ satisfying the following conditions:
\begin{enumerate}[label=(\roman*)]
\item there are finite constants $1<q_-\leq q_+<\infty$ such that
      $q_-\leq q_j\leq q_+$ for every $j$;
\item every finite rational relation between the $q_j$ is trivial;
\item $\Pp(X_j>0)>0$ for every $j$;
\item for some finite $M$, $\sum_{j\in J}X_j\leq M$ almost surely.
\end{enumerate}
The index set, exponents, constants in these assumptions, and joint law
are all allowed to be unknown. Known positive amplitudes can be absorbed
into $X_j$; if the amplitudes are themselves unknown, the theorem recovers
these scaled coordinates. No independence between the coordinates of
$X$ is assumed.

Define the normalized continuum gate profile
\begin{equation}\label{p07:eq:countable-gate-profile}
 \mathcal G_{q,X}(\theta)
 =\int_0^\infty
 -\log\left(1-\theta+
       \theta\E\exp\left\{-\sum_{j\in J}r^{-q_j}X_j\right\}\right)
 \,dr,\qquad 0<\theta<1.
\end{equation}
The exponents in \eqref{p07:eq:countable-gate-profile} are the normalized
ratios $q_j=p_j/\rho$ when the original radial measure is
$\rho y^{\rho-1}\,dy$ and $r=y^\rho$.

\begin{theorem}[Blind countable identification]
\label{p07:thm:countable-blind-gate}
Within the class just specified, the exact function
$\mathcal G_{q,X}$ on any nonempty open subinterval of $(0,1)$ determines
the index cardinality, the exponent set $\{q_j:j\in J\}$, and the joint
law of $X$, up to the simultaneous relabeling of coordinates and
exponents.

More explicitly, if two such models $(J,q,X)$ and $(\widetilde J,
\widetilde q,\widetilde X)$ have the same profile on such an interval,
then there is a bijection $\pi:J\to\widetilde J$ such that
\[
 \begin{gathered}
 q_j=\widetilde q_{\pi(j)}\quad(j\in J),\\
 (X_j)_{j\in J}\ \stackrel{d}{=}
 (\widetilde X_{\pi(j)})_{j\in J}
 \quad\hbox{as random elements of }\ell^1_+(J).
 \end{gathered}
\]
The exponents may have finite accumulation points. The theorem is an
identification statement for exact function data and does not assert a
finite-sample reconstruction rate.
\end{theorem}

\begin{proof}
We give the reconstruction in four steps.

\emph{Step 1: recover the monotone radial profile.}
Put
\[
 f(r)=1-\E\exp\left\{-\sum_{j\in J}r^{-q_j}X_j\right\}.
\]
For each $r>0$ the series in the exponent is finite almost surely, since
\[
 \sum_jr^{-q_j}X_j
 \leq M\max\{r^{-q_-},r^{-q_+}\}.
\]
The same estimate gives continuity in $r$ on compact subintervals of
$(0,\infty)$. In addition,
\begin{equation}\label{p07:eq:countable-gate-tail-bound}
 0\leq f(r)\leq1\quad(0<r\leq1),\qquad
 0\leq f(r)\leq M r^{-q_-}\quad(r\geq1).
\end{equation}
Thus $\int_0^\infty f(r)\,dr\leq1+M/(q_--1)<\infty$.
For $r_1<r_2$ the series at $r_1$ is strictly larger than that at $r_2$
on the event $\sum_jX_j>0$, which has positive probability. Therefore
$f$ is strictly decreasing. By dominated convergence its endpoint limits
are $m=\Pp(\sum_jX_j>0)$ at zero and zero at infinity.

The hypotheses used in the proof of
Theorem~\ref{p07:thm:gate-inversion} are precisely continuity, strict
monotonicity, and integrability of $f$, all just verified. That proof
does not use finite dimension. Applying its weighted pushforward and
tail inverse recovers $f$, and therefore
\begin{equation}\label{p07:eq:countable-blind-g}
 g(t)=1-f(1/t)
      =\E\exp\left\{-\sum_jt^{q_j}X_j\right\},\qquad t>0.
\end{equation}

\emph{Step 2: recover a finite signed measure of exponents.}
Let $\mathcal I_J$ be the countable set of nonnegative integer
multiindices $n=(n_j)_{j\in J}$ with finite support. Set
\[
 |n|=\sum_jn_j,\qquad n!=\prod_jn_j!,\qquad
 X^n=\prod_jX_j^{n_j},\qquad n\cdot q=\sum_jn_jq_j,
\]
using the empty-product convention. Define
\begin{equation}\label{p07:eq:blind-exponent-measure}
 c_n=\frac{(-1)^{|n|}}{n!}\E X^n,
 \qquad
 \eta=\sum_{n\in\mathcal I_J}c_n\delta_{n\cdot q}.
\end{equation}
Tonelli's theorem and increasing finite subsets of $J$ give
\begin{equation}\label{p07:eq:blind-exponent-variation}
 \sum_{n\in\mathcal I_J}|c_n|
   =\E\exp\left(\sum_jX_j\right)\leq e^M.
\end{equation}
Consequently \eqref{p07:eq:blind-exponent-measure} defines a finite signed
Borel measure on $[0,\infty)$. For $s>0$, absolute convergence allows
expansion of the exponential in \eqref{p07:eq:countable-blind-g}, yielding
\begin{equation}\label{p07:eq:blind-exponent-laplace}
 g(e^{-s})
   =\sum_{n\in\mathcal I_J}c_ne^{-s(n\cdot q)}
   =\int_{[0,\infty)}e^{-s\lambda}\,\eta(d\lambda).
\end{equation}
Thus $g$ determines $\eta$ by uniqueness of the Laplace transform of
finite signed measures. For completeness, if a finite signed measure
$\zeta$ has zero Laplace transform on $(0,\infty)$, push
$e^{-\lambda}\zeta(d\lambda)$ forward by $u=e^{-\lambda}$ to a
measure on $[0,1]$. Every nonnegative integer moment of the resulting
measure is zero. Polynomial density makes that measure zero, and then
restriction to $[0,R]$ and multiplication by $e^\lambda$ show
$\zeta=0$. This proof is valid in the presence of accumulating atoms.

\emph{Step 3: identify the unknown coordinates by indecomposability.}
Let
\begin{equation}\label{p07:eq:blind-atom-set}
 A=\{\lambda>0:\eta(\{\lambda\})\ne0\},
 \qquad
 Q=A\setminus(A+A),
 \qquad A+A=\{a+b:a,b\in A\}.
\end{equation}
Here $A$ is the set of nonzero atoms, not the closed topological support
of $\eta$. Rational independence makes the map $n\mapsto n\cdot q$
injective. Since $X$ is nonnegative, there can be no cancellation within
an atom, and
\begin{equation}\label{p07:eq:blind-atom-positivity}
 A=\{n\cdot q:n\in\mathcal I_J\setminus\{0\},\ \E X^n>0\}.
\end{equation}
The countable form of
Lemma~\ref{p07:lem:blind-primitive-exponents} now applies directly to
\eqref{p07:eq:blind-atom-positivity}. It proves the intrinsic formula
\begin{equation}\label{p07:eq:blind-intrinsic-indices}
 Q=\{q_j:j\in J\}.
\end{equation}
In particular no discreteness or local finiteness of $A$ is needed.
The argument uses hereditary positivity under taking submulti-indices,
not closure of $A$ under addition.

\emph{Step 4: recover the joint law.}
The recovered set $Q$ determines the unknown cardinality. Choose any
enumeration of $Q$ to label the coordinates. Rational independence and
\eqref{p07:eq:blind-exponent-measure} give, for every finite-support
multiindex,
\begin{equation}\label{p07:eq:blind-cylinder-moments}
 \E X^n=(-1)^{|n|}n!\,
          \eta\left(\left\{\sum_jn_jq_j\right\}\right).
\end{equation}
This formula also recovers zero moments. Each finite subvector is
bounded, so its mixed moments determine its law. When comparing two
models with different unknown bounds, use the larger of the two bounds
for this moment-determinacy argument. Agreement of every finite marginal
then determines the law on the coordinate $\sigma$-algebra.

Finally, on $\ell^1(J)$ the coordinate $\sigma$-algebra equals the norm
Borel $\sigma$-algebra. One inclusion follows from continuity of the
coordinate maps. For the reverse inclusion, for each fixed
$x\in\ell^1(J)$ the function $z\mapsto\|z-x\|_1$ is the increasing
limit of measurable finite sums of coordinate distances. Norm balls
are therefore coordinate measurable; separability of $\ell^1(J)$ gives
a countable base of such balls. Thus the recovered law is the law as an
$\ell^1$-valued random element, and \eqref{p07:eq:blind-intrinsic-indices}
establishes the asserted simultaneous relabeling.
\end{proof}

\begin{remark}[What is being recovered]
The theorem identifies a model satisfying the structural assumptions;
it does not infer those assumptions from a finite experiment. A
coordinate that is identically zero is invisible and is excluded.
Likewise, the normalized ratios $q_j$ do not separate an unknown
counting exponent $\rho$ from the physical powers $p_j$. The data consist
of a full gate-response function, so this result does not contradict a
lower bound proved for finitely many fixed-gate channels.
\end{remark}

\section{Correlated calibration: a local obstruction and its repair}
\label{p07:sec:correlated-calibration}

Let $W=(W_1,\ldots,W_m)$ be a bounded nonnegative random vector,
assume $\mathbb P(W=0)>0$ and $\mathbb P(W_j>0)>0$ for every $j$,
and let $p_j>1$.  Write
\[
 h(u)=-\log\mathbb E e^{-u\cdot W},\qquad
 H(s)=\int_0^\infty h(s_1x^{-p_1},\ldots,s_mx^{-p_m})\,dx.
\]
The counting exponent has been normalized to one; thus $p_j$ denotes
the physical exponent divided by the counting exponent.  The function
$H$ is finite, increasing, concave, and real analytic on the positive
orthant.  It satisfies
\begin{equation}\label{p07:eq:vector-weighted-euler}
 \sum_{j=1}^m p_j s_j H_j(s)=H(s).
\end{equation}

For completeness, analyticity does not require independence or a
positive covariance assumption.  If $s$ is positive and
$|v_j|\le\delta s_j$, the perturbation of the Laplace transform at
$(s_jx^{-p_j})_j$ is bounded by
$\mathbb E[e^{-T}(e^{\delta T}-1)]$, where
$T=\sum_j s_jx^{-p_j}W_j$.  This bound tends to zero uniformly in
$x$ as $\delta\downarrow0$, while the unperturbed transform is bounded
below by $\mathbb P(W=0)$.  A fixed relative complex neighborhood
therefore has a zero-free Laplace transform and a consistent logarithm.
Its logarithm is uniformly bounded for $x\le1$, and is
$O(x^{-\min_jp_j})$ for $x\ge1$.  Holomorphic integration proves the
assertion.  Homogeneity follows by a change of variables.

At a point $s$, put $z=\nabla H(s)$, $M=-\nabla^2H(s)$, and
\[
 B=\operatorname{diag}(z)-M\operatorname{diag}(s).
\]
Differentiation of \eqref{p07:eq:vector-weighted-euler} gives $Bp=z$.
For common-scalar or nonnegative infinitely divisible models, all
entries of $M$ are nonnegative, and the usual nonsingular $M$-matrix
argument proves that $B$ is invertible.  Correlation can remove this
sign property.  The following is an exact failure of identification,
rather than merely a failure of that proof.

\begin{theorem}[A full local curvature jet can conceal the exponents]
\label{p07:thm:correlated-jet-counterexample}
Fix any gate probability $\pi_0\in(1/2,1)$.  There exist $P>1$,
$\epsilon_0>0$, and a smooth family of bounded random vectors
$W^{(\epsilon)}\in[0,\infty)^2$, $|\epsilon|<\epsilon_0$, such that:
\begin{enumerate}
\item $\mathbb P(W^{(\epsilon)}=0)=\pi_0$ for every $\epsilon$;
conditionally on being nonzero, each law has two strictly positive,
linearly independent support points;
\item the associated exponents are
$(p_1,p_2)=(P+\epsilon,P-\epsilon)$;
\item the values, gradients, and Hessians of all associated transforms
$H_\epsilon$ at $s=(1,1)$ are identical.
\end{enumerate}
All these Hessians are negative definite.  Consequently the local
Legendre rates $I_\epsilon(z)=\sup_{s\ge0}\{H_\epsilon(s)-s\cdot z\}$
have the same full second-order jet at a common fully active point.
In particular, even a known shared Bernoulli gate with arbitrarily
large zero probability does not make one-point curvature calibration
valid for all correlated laws.
\end{theorem}

\begin{proof}
Set $c=(1-\pi_0)/\pi_0>0$.  For $0<a<1$, define
\[
 Q(t)=\frac{ce^{-t}}{1+ce^{-t}},\quad
 J_1(a)=\int_0^\infty t^{-a}Q(t)\,dt,\quad
 J_2(a)=\int_0^\infty t^{1-a}Q(t)\,dt.
\]
As $a\downarrow0$,
\[
 J_1(a)\longrightarrow\log(1+c),\qquad
 J_2(a)\longrightarrow\int_0^c\frac{\log(1+u)}{u}\,du.
\]
The latter limit is strictly larger than the former: their difference
vanishes at $c=0$ and has derivative
$\log(1+c)/c-1/(1+c)>0$.  Choose $a$ small enough that $J_2(a)>J_1(a)$,
put $P=1/a$, and set $b_0=\sqrt{J_1(a)/J_2(a)}\in(0,1)$.

Consider laws depending on $(L,\theta,b)$ near $(1,0,b_0)$.
Give the origin probability $\pi_0$, and give the points
\[
 x_+=\frac L2(1+b,1-b),\qquad
 x_-=\frac L2(1-b,1+b)
\]
probabilities $(1-\pi_0)(1+\theta)/2$ and
$(1-\pi_0)(1-\theta)/2$, respectively.  At equal exponents $(P,P)$
and $s=(1,1)$, direct differentiation and the substitution
$t=Lx^{-P}$ give
\begin{align*}
 H&=L^aJ_1,\\
 D:=H_1-H_2&=aL^a\theta bJ_1,\\
 R:=H_{11}+H_{22}-2H_{12}
 &=-aL^ab^2\{J_2-\theta^2J_3\},
 \qquad J_3=\int_0^\infty t^{1-a}Q(t)^2\,dt.
\end{align*}
For the first identity one integrates by parts in
$a\int_0^\infty t^{-a-1}\log[(1+c)/(1+ce^{-t})]dt$.
The Jacobian of $(H,D,R)$ with respect to $(L,\theta,b)$ at
$(1,0,b_0)$ is nonsingular: its three successive diagonal entries
are $aJ_1$, $ab_0J_1$, and $-2ab_0J_2$.

At this base point, $z_1=z_2=z=aJ_1/2$ and
$H_{11}-H_{12}=R/2=-z$.  Thus $B(1,-1)^{\mathsf T}=0$.
In addition to the true Euler identities with weight vector $(P,P)$,
the entire base second-order jet therefore satisfies the Euler
identities with every weight vector $(P+\epsilon,P-\epsilon)$.

Apply the implicit function theorem to the three observations
$(H,D,R)$ as the exponents vary from $(P,P)$ to
$(P+\epsilon,P-\epsilon)$.  It provides smooth functions
$L(\epsilon),\theta(\epsilon),b(\epsilon)$ keeping these observations
fixed.  At any positive exponent pair, the three Euler equations
\[
 p_1H_1+p_2H_2=H,\quad
 p_1H_{11}+p_2H_{12}=(1-p_1)H_1,\quad
 p_1H_{12}+p_2H_{22}=(1-p_2)H_2
\]
together with $D$ and $R$ uniquely determine all remaining entries
of the second-order jet.  Indeed the first equation and $D$ determine
the gradient; the final two equations and $R$ determine the three
Hessian entries.  The coefficient matrix in the latter step has
nonzero determinant of absolute value $(p_1+p_2)^2$.  Since the base
jet also satisfies these same equations for every $\epsilon$, all
jets agree exactly.

The three support points $0,x_+,x_-$ are affinely independent and
have positive probabilities.  Every tilted covariance matrix is
therefore positive definite, and integration preserves strict
negative definiteness of the Hessian of $H_\epsilon$.  The common
positive saddle $s=(1,1)$ has a common gradient $z$, and Legendre
inversion gives
$I_\epsilon(z)=H_\epsilon(1,1)-2z$,
$\nabla I_\epsilon(z)=-(1,1)$, and
$\nabla^2I_\epsilon(z)=[-\nabla^2H_\epsilon(1,1)]^{-1}$.
This proves the rate assertion.
\end{proof}

For example, $c=1/2$ and $P=10$ give numerically
$J_1\simeq0.4271863541$, $J_2\simeq0.4280990318$, and
$b_0\simeq0.9989334657$.  These numbers illustrate the construction;
existence in the theorem follows from the preceding analytic argument.
Taking two distinct small positive values of $\epsilon$ yields
distinct, ordered exponent pairs, so the phenomenon is not merely an
interchange of coordinate labels.

\begin{theorem}[Generic multipoint calibration and eventual finite jets]
\label{p07:thm:correlated-multipoint-calibration}
Under the bounded-vector assumptions at the beginning of this section:
\begin{enumerate}
\item the exponent vector $p$ is uniquely determined by $H$ on any
nonempty open subset of the positive orthant;
\item in every such subset, a generic $m$-tuple of points
$s^{(1)},\ldots,s^{(m)}$ determines $p$ from only the values and first
derivatives of $H$, by the linear system
\begin{equation}\label{p07:eq:multipoint-calibration}
 \sum_j s_j^{(i)}H_j(s^{(i)})p_j=H(s^{(i)}),
 \qquad 1\le i\le m;
\end{equation}
\item at every fixed positive point, there is some finite derivative
order, possibly depending on the law and the point, whose jet
uniquely determines $p$ by differentiated Euler identities.
\end{enumerate}
Here ``generic'' means outside the zero set of a nonzero real-analytic
determinant; in particular, the exceptional set has empty interior
and Lebesgue measure zero.  No uniform bound on the derivative order
in the last assertion is claimed.
\end{theorem}

\begin{proof}
The transform extends continuously to the nonnegative orthant and
has coordinate restrictions
\begin{equation}\label{p07:eq:axis-calibration}
 H(te_j)=C_jt^{1/p_j},\qquad
 C_j=\int_0^\infty-\log\mathbb E e^{-x^{-p_j}W_j}\,dx>0.
\end{equation}
Continuity follows by dominated convergence, using the atom at zero
near $x=0$ and boundedness of $W$ near $x=\infty$.
Analytic continuation of a known open patch and
\eqref{p07:eq:axis-calibration} already prove the first assertion.

For the constructive statement, put
$A(s)=(s_1H_1(s),\ldots,s_mH_m(s))$.
These vectors span $\mathbb R^m$ as $s$ ranges over any nonempty open
set $U$.  Otherwise some nonzero $q$ satisfies
$\sum_jq_js_jH_j(s)=0$ on $U$, hence throughout the positive orthant
by analyticity.  Integration along the flow gives
$H(e^{tq_1}s_1,\ldots,e^{tq_m}s_m)=H(s)$ for every real $t$.
Passing to a coordinate axis and using \eqref{p07:eq:axis-calibration}
forces $e^{tq_j/p_j}=1$ for every $j$, hence $q=0$, a contradiction.
Thus there is an invertible matrix with $m$ rows $A(s^{(i)})$ in
$U$.  Its determinant is analytic and not identically zero on the
connected positive orthant product, which proves the generic claim.

For the final statement, differentiate Euler's identity.  For each
multi-index $\beta$ one obtains the linear equation
\begin{equation}\label{p07:eq:all-jet-euler}
 \sum_j\left(s_j\partial^{\beta+e_j}H
          +\beta_j\partial^\beta H\right)p_j
       =\partial^\beta H.
\end{equation}
If the rows of these equations did not span $\mathbb R^m$ at the
chosen point, a nonzero $q$ would annihilate all of them.  Every
Taylor coefficient of the analytic function
$\sum_jq_js_jH_j(s)$ would then vanish there.  The same dilation-flow
argument would imply $q=0$.  The rows therefore span a finite-dimensional
space of full dimension, so some finite collection of them already
has full rank.  They involve a finite jet and recover $p$ uniquely.
\end{proof}

When an active rate $I$ is observed instead, one sets
$s=-\nabla I(z)$ and $H(s)=I(z)+s\cdot z$.  Thus the multipoint
system can be written directly as
\[
 \sum_j s_j^{(i)}z_j^{(i)}p_j
       =I(z^{(i)})+s^{(i)}\cdot z^{(i)}.
\]
The argument distinguishes exact identifiability from numerical
conditioning: the error amplification of the finite system is
controlled by its smallest singular value, which need not be bounded
away from zero uniformly over the model class.  Differentiated Euler
identities and symmetry-based rank tests are classical tools; the
specific contribution here is the explicit gated correlated family
with coincident curvature data, together with a calibration repair
that does not assume signs for cross-covariances.



\section{Blind recovery by an interior power probe}
\label{p08:sec:tempered-probe}

Moving a Laplace probe from the boundary of the positive orthant to
its interior changes the relevant moment problem. The moments to be
decoded are then exponentially tilted moments, which exist and determine
the tilted law without any moment assumption on the original law.
For known powers, rationally independent moment encoding is the method
of \cite[Lemmas~3.2--3.3]{Lenzi2026}; the contribution considered here is
its combination with blind recovery of the powers, the dimension, a
common unknown baseline, and an unknown observation gain.

Let $I$ be an unknown nonempty finite or countably infinite index set.
Write $\ell^1_+(I)$ for the nonnegative summable sequences, and let
$X=(X_i)_{i\in I}$ have a Borel probability law $P$ on this space. Set
$S(x)=\sum_i x_i$. The probe parameters are an unknown baseline $a>0$,
an unknown gain $C>0$, and unknown positive powers $q_i$ satisfying
\begin{equation}\label{p08:eq:temp-assumptions}
 \inf_{i\in I}q_i>0,\qquad
 \{q_i:i\in I\}\text{ is linearly independent over }\mathbb Q,
 \qquad P(X_i>0)>0\quad(i\in I).
\end{equation}
Linear independence in the countable case means independence of every
finite subfamily. No upper bound on the powers is imposed. Define
\begin{align}
 F(t)&=\mathbb E_P\exp\left[-aS(X)-\sum_{i\in I}t^{q_i}X_i\right],
                  &&0\leq t\leq1,\label{p08:eq:temp-F}\\
 \mathcal G(\theta)&=C\int_0^1
       -\log\bigl[1-\theta+\theta F(t)\bigr]\,dt,
                  &&0<\theta<1.\label{p08:eq:temp-G}
\end{align}
The parameter $t$ has the displayed known window and Lebesgue clock;
the gain $C$ is not assumed known. Since $S(X)<\infty$ almost surely,
these quantities exist for every probability law under consideration.
In particular, no positive ordinary moment need exist.

\begin{theorem}[Blind dimension, powers, baseline, gain, and law]
\label{p08:thm:temp-fully-blind}
Knowledge of $\mathcal G$ on any nonempty open interval contained in
$(0,1)$ determines $C$, $a$, the cardinality of $I$, the unordered set
$\{q_i:i\in I\}$, and the joint law $P$, up to the same permutation of
the powers and coordinates, throughout the class
\eqref{p08:eq:temp-assumptions}--\eqref{p08:eq:temp-G}.

More precisely, the observation first determines the curve $F$ and
then the finite measure
\begin{equation}\label{p08:eq:temp-nu}
             \nu(dx)=e^{-aS(x)}P(dx).
\end{equation}
The baseline is the unique positive solution of
\begin{equation}\label{p08:eq:temp-normalize}
                  \int e^{bS(x)}\nu(dx)=1,
\end{equation}
and $P(dx)=e^{aS(x)}\nu(dx)$.
\end{theorem}

\begin{proof}
\emph{Step 1: gate inversion and the gain.}
Put $f(t)=1-F(t)$. Dominated convergence, first for the summable series
and then for the bounded exponential, shows that $F$ is continuous on
$[0,1]$. It is strictly decreasing: if $0\leq t<u\leq1$, then
$\sum_i(u^{q_i}-t^{q_i})X_i>0$ on the event $\{S(X)>0\}$, which has
positive probability. Also
\[
             0<f(0)\leq f(t)\leq f(1)<1.
\]
The first strict inequality follows from $a>0$ and
$P(S(X)>0)>0$; the last follows from $S(X)<\infty$ almost surely.
The function $\mathcal G$ extends holomorphically to $|\theta|<1$,
so the prescribed data determine all its derivatives at zero.
Define the finite positive measure on $[0,1]$
\[
             \kappa(A)=C\int_0^1 f(t)\mathbf1_A(f(t))\,dt.
\]
For every $n\geq0$,
\begin{equation}\label{p08:eq:temp-gate-moments}
       \int u^n\kappa(du)=\frac{\mathcal G^{(n+1)}(0)}{n!}.
\end{equation}
Compact moment uniqueness determines $\kappa$. Consequently it
determines
\[
 \sigma(du)=\mathbf1_{(0,1]}(u)u^{-1}\kappa(du),\qquad
 C=\sigma([0,1]).
\]
The probability measure $\sigma/C$ is the distribution of $f(U)$ for
$U$ uniform on $[0,1]$. Its quantile function is the continuous strictly
increasing function $f$, including the endpoints by limits. Thus $C$
and the whole curve $F$ are recovered.

\emph{Step 2: a finite signed measure of exponents.}
For a finitely supported multi-index $n\in\mathbb N_0^{(I)}$, put
$|n|=\sum_i n_i$, $n!=\prod_i n_i!$, $n\cdot q=\sum_i n_iq_i$, and
\[
                        m_n=\int x^n\nu(dx).
\]
All these moments are finite, because
$x^n e^{-aS(x)}\leq S(x)^{|n|}e^{-aS(x)}$ is bounded.
Let $q_- =\inf_i q_i>0$ and choose $0<\tau<1$ so small that
$\tau^{q_-}<a$. The signed atomic measure
\begin{equation}\label{p08:eq:temp-eta}
 \eta_\tau=\sum_{n\in\mathbb N_0^{(I)}}
       \frac{(-1)^{|n|}\tau^{n\cdot q}m_n}{n!}\,
                         \delta_{n\cdot q}
\end{equation}
has finite total variation. Indeed, the countable multinomial theorem
and monotone convergence give
\begin{align}
 \sum_n\frac{\tau^{n\cdot q}m_n}{n!}
  &=\int\exp\left(\sum_i\tau^{q_i}x_i\right)\nu(dx)\notag\\
  &\leq\mathbb E_P e^{-(a-\tau^{q_-})S(X)}\leq1.
                                                    \label{p08:eq:temp-TV}
\end{align}
The absolutely convergent expansion therefore yields
\begin{equation}\label{p08:eq:temp-Laplace}
       F(\tau e^{-s})=\int_{[0,\infty)}e^{-s\lambda}\eta_\tau(d\lambda),
                    \qquad s\geq0.
\end{equation}
Uniqueness of the Laplace transform of finite signed measures recovers
$\eta_\tau$. One elementary proof pushes the measure forward under
$\lambda\mapsto e^{-\lambda}$ and uses compact moment uniqueness at
the nonnegative integer values of $s$.

Since $a$ and $q_-$ are unknown, the choice of $\tau$ is not asserted
to be a data-effective stopping rule. For the identifiability claim,
compare two admissible models with the same data and choose a common
$\tau$ small enough for both. Step 1 gives the same $F$, and
\eqref{p08:eq:temp-Laplace} then gives the same signed measure
$\eta_\tau$. This comparison is sufficient for all subsequent steps.

\emph{Step 3: the primitive powers.}
Use the set of actual nonzero positive atoms
\[
          A_\tau=\{\lambda>0:\eta_\tau(\{\lambda\})\ne0\},
\]
not its topological closure. Rational independence prevents collisions
between distinct multi-indices in \eqref{p08:eq:temp-eta}, so
\[
 A_\tau=\{n\cdot q:n\ne0,\ m_n>0\}.
\]
The finite-positive-measure version of
Lemma~\ref{p07:lem:blind-primitive-exponents}, applied to $\nu$, shows that
the additively indecomposable elements of $A_\tau$ are exactly the powers
$q_i$. Indeed $m_{e_i}>0$, and a positive mixed moment remains positive
upon deleting a coordinate factor. The lemma uses actual atoms rather
than topological support, so it remains valid when the powers and their
finite sums accumulate. It identifies whether the dimension is finite
or countably infinite as well as its cardinality.
Once the powers have been matched, all moments are
recovered from
\begin{equation}\label{p08:eq:temp-coefficients}
       m_n=(-1)^{|n|}n!\tau^{-n\cdot q}
                         \eta_\tau(\{n\cdot q\}).
\end{equation}
In particular the zeroth moment is $\nu(\ell^1_+)=F(0)>0$.

\emph{Step 4: the tilted law, its normalization, and the baseline.}
For every finite $J\subset I$ and $0<b\leq a$,
\[
       \int\exp\left(b\sum_{i\in J}x_i\right)\nu(dx)\leq1.
\]
Each finite-coordinate marginal of $\nu$ thus has a moment generating
function in a neighborhood of zero and is determined by its moments.
The finite-coordinate marginals determine the measure on the countable
product, and therefore on $\ell^1_+(I)$: its Borel sigma field equals
the trace of the coordinate sigma field. Thus the recovered moments
determine $\nu$.

Finally $H(b)=\int e^{bS(x)}\nu(dx)$ is strictly increasing wherever
finite, because $\nu(S>0)>0$. It satisfies $H(a)=1$, whereas
$H(0)=F(0)<1$. It can therefore equal one for only one positive $b$.
This identifies $a$, and then \eqref{p08:eq:temp-nu} identifies $P$.
\end{proof}

\begin{remark}[What is and is not assumed]
The original law can have infinite mean, no finite positive ordinary
moment, or moment-indeterminate finite moments. In finite dimension,
every nonnegative finite-valued vector satisfies the summability
condition. In countable dimension the condition is $S(X)<\infty$
almost surely, with no deterministic bound and no integrability
condition on $S$. The lower bound on the powers is a hypothesis of
the countable theorem. An identically zero coordinate remains
unobservable and has explicitly been excluded. Unknown coefficients
multiplying $t^{q_i}X_i$ introduce the usual coordinate scale ambiguity;
those coefficients are normalized to one here.

This is an exact-data identification theorem. An open interval of the
gate curve contains infinitely much data; neither finite-sample
recovery nor a stable procedure for resolving infinitely many powers
is asserted. The construction of the finite signed exponent measure
is an adaptation of the known-power method in
\cite[Lemmas~3.2--3.3]{Lenzi2026}; exponential tilting and compact
moment uniqueness are classical. The extra blind parameters and the
precise ambiguity below are the candidate contribution, without a
claim that its priority has been exhaustively established.
\end{remark}

\begin{example}[How the interior probe breaks a stable-law alias]
Let $0<\beta<1$ and let $Z_\beta$ have Laplace transform
$\mathbb E e^{-sZ_\beta}=e^{-s^\beta}$. At zero baseline, the
deterministic variable $X=1$ with power $q$ and the variable
$Z_\beta$ with power $q/\beta$ give the identical curve $e^{-t^q}$.
For positive baselines $a$ and $b$, their interior curves instead are
\[
 e^{-a-t^q},\qquad e^{-(b+t^{q/\beta})^\beta}.
\]
Their first nonconstant exponents at $t=0$ are respectively $q$ and
$q/\beta$, with nonzero coefficients. Thus they cannot agree, even if
their baselines are independently unknown. The gate and gain inverse
ensures that an unknown positive observation gain cannot conceal this
difference either. This explicitly removes the heavy-tail alias rather
than imposing a moment condition that excludes it.
\end{example}

\subsection{The exact ambiguity for an unknown vector baseline}

For finite $d$, replace the scalar $a$ in \eqref{p08:eq:temp-F} by an
unknown vector $b=(b_1,\ldots,b_d)\in(0,\infty)^d$:
\[
 F(t)=\int\exp\left[-b\cdot x-\sum_{i=1}^dt^{q_i}x_i\right]P(dx).
\]
The preceding proof still identifies the gain, the powers, and the
finite measure $\nu(dx)=e^{-b\cdot x}P(dx)$.

\begin{proposition}[Complete vector-baseline fiber]
\label{p08:prop:temp-vector-fiber}
After powers and coordinates have been matched, all models producing
this same data are exactly
\begin{equation}\label{p08:eq:temp-fiber}
 \mathcal B_\nu=\left\{c\in(0,\infty)^d:
                      \int e^{c\cdot x}\nu(dx)=1\right\},
       \qquad P_c(dx)=e^{c\cdot x}\nu(dx).
\end{equation}
If $\nu$ has compact support and every coordinate is nontrivial,
$\mathcal B_\nu$ is locally a smooth hypersurface of dimension $d-1$
near every one of its points. Restricting the baseline to a known
positive ray $c=av$, $v_i>0$, leaves at most one admissible $a$, and
exactly one whenever a model on that ray is assumed to exist.
\end{proposition}
\begin{proof}
The equality in \eqref{p08:eq:temp-fiber} is precisely what makes $P_c$
a probability measure. Substitution shows that its probe is
$\int e^{-\sum_i t^{q_i}x_i}\nu(dx)$, independent of $c$.
Conversely, recovery of $\nu$ in the preceding proof shows that every
admissible model has this form. In the compactly supported case the
function $H(c)=\int e^{c\cdot x}\nu(dx)$ is smooth, and every partial
derivative $\partial_iH=\int x_i e^{c\cdot x}\nu(dx)$ is strictly
positive. The implicit function theorem gives the hypersurface
statement. Along the prescribed ray, $a\mapsto H(av)$ is strictly
increasing wherever finite, proving uniqueness there.
\end{proof}

\subsection{A sharp fixed-power criterion}

The positive baseline removes tail assumptions, but it does not remove
arithmetic resonance. This remains a genuine observational obstruction.

\begin{proposition}[Necessity and sufficiency of rational independence]
\label{p08:prop:temp-sharp-criterion}
Fix a finite positive baseline vector $b$, a positive gain $C$, and
positive powers $q_1,\ldots,q_d$. The map taking an arbitrary probability
law $P$ on $[0,\infty)^d$ to its finite-window gate curve is injective
if and only if the powers are linearly independent over $\mathbb Q$.
If the powers are dependent, any prescribed positive open box contains
two distinct smooth compactly supported probability densities with
identical full gate curves. They can additionally be chosen to have
the same ordinary mixed moments through any prescribed finite order.
\end{proposition}
\begin{proof}
The zero curve corresponds exactly to $P=\delta_0$. For every other
law, sufficiency follows from the tilted-moment part of the preceding
proof, without requiring every coordinate to be nontrivial when the
dimension and powers are fixed. For necessity, choose distinct
$\alpha,\beta\in\mathbb N_0^d$ with
$\alpha\cdot q=\beta\cdot q$. The nonzero polynomial
\[
            R(s)=(s-b)^\alpha-(s-b)^\beta
\]
vanishes along $s=b+(t^{q_1},\ldots,t^{q_d})$. If moments through
order $N$ are to be preserved, put $Q(s)=s_1^{N+1}R(s)$.
Choose a nonzero nonnegative smooth bump $\varphi$ supported inside
the prescribed box, and set $g=Q(\partial)\varphi$.
Repeated integration by parts gives
\[
          \int e^{-s\cdot x}g(x)\,dx
                 =Q(s)\int e^{-s\cdot x}\varphi(x)\,dx.
\]
Thus the Laplace transform of $g$ vanishes on the probe curve, and
all its derivatives through order $N$ at zero vanish. In particular
$\int g=0$, and its ordinary moments through that order vanish.
The function $g$ is nonzero: the displayed transform cannot vanish
identically, since $Q$ is nonzero and the Laplace transform of
$\varphi$ is strictly positive at every real positive $s$.
Choose a smooth probability density $p_0$ supported in the same box
and bounded below by a positive constant on the support of $g$.
For sufficiently small $\varepsilon>0$, the functions
$p_0\pm\varepsilon g$ are nonnegative probability densities. Their
Laplace curves, and consequently their full gate curves, agree, while
the densities are distinct and their prescribed moments agree.
\end{proof}

\subsection{Universal identification does not give uniform tail stability}

\begin{proposition}[An exponential tail barrier]
\label{p08:prop:temp-tail-barrier}
Even with known dimension $d=1$, baseline $a>0$, power $q>0$, and
gain $C>0$, the inverse observation map is not uniformly continuous
over all nonnegative probability laws in total variation or in the
Wasserstein distance $W_1$, even if the entire gate interval including
$\theta=1$ is observed. More precisely, fix $0<\varepsilon<1$ and,
for $R>0$, put
\[
 P_R=(1-\varepsilon)\delta_0+\varepsilon\delta_R,
 \qquad Q_R=(1-\varepsilon)\delta_0+\varepsilon\delta_{R+1}.
\]
With total variation normalized to lie in $[0,1]$, one has
\[
     \operatorname{TV}(P_R,Q_R)=W_1(P_R,Q_R)=\varepsilon,
\]
whereas
\begin{equation}\label{p08:eq:temp-tail-bound}
 \sup_{0\leq\theta\leq1}
       |\mathcal G_{P_R}(\theta)-\mathcal G_{Q_R}(\theta)|
       \leq\frac{C\varepsilon}{1-\varepsilon}e^{-aR}.
\end{equation}
\end{proposition}
\begin{proof}
The total variation is immediate. Transporting the mass $\varepsilon$
from $R$ to $R+1$ gives $W_1\leq\varepsilon$, and the difference
of means gives the opposite inequality. Their Laplace profiles differ
by at most $\varepsilon e^{-aR}$ uniformly on $0\leq t\leq1$, and
both profiles are at least $1-\varepsilon$. For every
$0\leq\theta\leq1$, the function
$u\mapsto-\log(1-\theta+\theta u)$ has derivative of absolute value
at most $(1-\varepsilon)^{-1}$ on $[1-\varepsilon,1]$.
Integrating proves \eqref{p08:eq:temp-tail-bound}. Letting $R\to\infty$
gives the failure of every uniform inverse modulus tending to zero.
\end{proof}

The proposition does not rule out stability on a fixed tight class
with additional quantitative restrictions. It separates the universal
exact identification theorem from a uniform noisy-data recovery claim;
the common atom at zero shows that this instability persists even when
the endpoint $\theta=1$ is available.


\begin{example}[Countably many coordinates with no positive moments]
\label{p08:ex:no-positive-moments}
Let $Y$ have density $(1+y)^{-2}$ on $[0,\infty)$, put $T=e^Y-1$,
and set $X_i=2^{-i}T$ for $i\geq1$. Then
$\sum_iX_i=T<\infty$ almost surely, but
\[
             \E X_i^r=\infty\qquad(i\geq1,\ r>0).
\]
Indeed the exponential growth of $(e^y-1)^r$ dominates the polynomial
decay of the density. Choose $q_i=\sqrt{p_i}$, where $p_i$ is the
$i$th prime. The square roots of distinct primes are rationally
independent: in each finite multiquadratic field, the automorphism
reversing one square root and fixing the others isolates its coefficient
in a putative rational relation. These powers have a positive lower
bound and no upper bound. For every $a,C>0$, the response is well
defined and Theorem~\ref{p08:thm:temp-fully-blind} identifies the model
within its stated class. This example simultaneously uses countably
many coordinates, unbounded powers, arbitrary dependence, and the
absence of every positive ordinary moment.
\end{example}

\section{An interior probe invisible to every ordinary jet}
\label{p08:sec:flat-probe}

The origin is the wrong expansion point if the unknown distribution has
heavy tails.  At a strictly positive Laplace parameter every finite
measure has an analytic transform, regardless of its ordinary moments.
We first give an explicit uniqueness curve for analytic germs and then
apply this elementary analytic fact to probability measures.  Analytic
uniqueness sequences are an established subject
\cite{BoivinGauthierManolaki2018}; multidimensional Laplace uniqueness
sequences are studied in \cite{Kostic2025}.  The results below give a
self-contained construction, without a claim of priority for the
general uniqueness-set principle.

\begin{definition}[Separated scales]
Let $x_j:(0,\epsilon)\to(0,1)$, $1\le j\le d$, tend to zero as
$t\downarrow0$.  They are separated if
\[
 \frac{\log(1/x_j(t))}{\log(1/x_{j-1}(t))}\longrightarrow\infty,
 \qquad 2\le j\le d.
\]
\end{definition}

\begin{theorem}[A leading monomial on separated scales]
\label{p08:thm:flat-leading}
For separated scales and every nonzero real or complex analytic germ
$A$ at $0\in\mathbb R^d$, there exist a unique
$\alpha\in\mathbb N_0^d$ and $c\ne0$ such that
\begin{equation}
 A(x_1(t),\ldots,x_d(t))
   =c\prod_{j=1}^d x_j(t)^{\alpha_j}(1+o(1)).
 \label{p08:eq:flat-leading}
\end{equation}
Here $\alpha$ is the smallest index with a nonzero Taylor coefficient
when indices are compared first in coordinate $d$, then $d-1$, and so
on.  In particular, the composition has no zero for all sufficiently
small positive $t$.
\end{theorem}

\begin{proof}
For $d=1$, factor the ordinary order of vanishing from the convergent
power series.  Suppose the assertion holds in dimension $d-1$.
The nonempty set of powers of $z_d$ occurring in $A$ has a least
element $m$.  Convergence of the series gives the analytic factorization
\[
 A(z',z_d)=z_d^m\bigl(A_0(z')+z_d B(z',z_d)\bigr),
 \qquad A_0\not\equiv0,
\]
where $B$ is bounded on a sufficiently small closed polydisc.
By induction, $A_0(x_1(t),\ldots,x_{d-1}(t))$
is asymptotic to $c\prod_{j<d}x_j(t)^{\alpha_j}$.
Writing $\ell_j(t)=\log(1/x_j(t))$, separation implies
$\sum_{j<d}\alpha_j\ell_j(t)=o(\ell_d(t))$.  Consequently
\[
 \frac{x_d(t)}{\prod_{j<d}x_j(t)^{\alpha_j}}
 =\exp\left(-\ell_d(t)+\sum_{j<d}\alpha_j\ell_j(t)\right)
 \longrightarrow0.
\]
The term containing $B$ is negligible.  This proves the assertion with
$\alpha_d=m$, including the description of its coefficient.
Two distinct monomials on these scales have a ratio tending to zero
or infinity, so the leading index and coefficient are unique.
\end{proof}

\begin{corollary}[A universal discrete interior probe]
\label{p08:cor:flat-laplace}
Fix $a\in(0,\infty)^d$ and any sequence $t_n>0$ tending to zero.
For $d\ge2$, put
\begin{equation}
 c(t)=\bigl(t,e^{-1/t},e^{-1/t^2},\ldots,e^{-1/t^{d-1}}\bigr),
 \qquad s_n=a+c(t_n).
 \label{p08:eq:flat-curve}
\end{equation}
For $d=1$, put $c(t)=t$.
If finite Borel measures $\mu,\nu$ on $[0,\infty)^d$ satisfy
\[
 \int e^{-\langle s_n,x\rangle}\,\mu(dx)
 =\int e^{-\langle s_n,x\rangle}\,\nu(dx)
 \qquad\text{for every }n,
\]
then $\mu=\nu$.  No moment, tail, support, independence, or
moment-determinacy assumption is required.  The same statement holds
for finite signed or complex measures.
\end{corollary}

\begin{proof}
The difference of the transforms is holomorphic on
$\{s\in\mathbb C^d:\Re s_j>0\text{ for all }j\}$.
Indeed, differentiation of every order is justified on a neighborhood
of $a$ by
\[
 \sup_{x\ge0}x^\beta
       \exp\left(-\frac12\sum_j a_jx_j\right)<\infty,
\]
with the product convention for $x^\beta$.
The scales in \eqref{p08:eq:flat-curve} are separated.
If the transform difference had a nonzero germ at $a$,
Theorem~\ref{p08:thm:flat-leading} would make its values nonzero at all
sufficiently late $s_n$, a contradiction.  Thus the transforms agree
on a neighborhood of $a$, and the identity theorem extends this
equality to the connected positive tube.
For a fully explicit final uniqueness step, fix $b\in(0,\infty)^d$.
The finite signed or complex measure
$e^{-\langle b,x\rangle}(\mu-\nu)(dx)$ has Fourier transform
equal to the transform difference at $b+i\xi$, hence zero everywhere.
Fourier uniqueness and positivity of the exponential weight give
$\mu=\nu$.
\end{proof}

\begin{corollary}[Exponential approximation and weak convergence]
\label{p08:cor:flat-density}
For the sample parameters $s_n$ of Corollary~\ref{p08:cor:flat-laplace},
the linear span of
\[
 x\longmapsto e^{-\langle s_n,x\rangle},\qquad n\ge1,
\]
is dense in $C_0([0,\infty)^d)$ in the uniform norm and in
$L^p([0,\infty)^d)$ for every $1\le p<\infty$.
Moreover, for probability measures $\mu_k,\mu$ on $[0,\infty)^d$,
\[
 \mu_k\Rightarrow\mu
 \quad\Longleftrightarrow\quad
 L_{\mu_k}(s_n)\longrightarrow L_\mu(s_n)
 \text{ for every }n.
\]
In particular,
\[
 d_c(\mu,\nu)=\sum_{n\ge1}2^{-n}
           |L_\mu(s_n)-L_\nu(s_n)|
\]
is a metric inducing the weak topology on probability measures.
\end{corollary}

\begin{proof}
If the span were not dense in $C_0$, Hahn--Banach and
Riesz--Markov would give a nonzero finite signed or complex measure
annihilating it, contrary to Corollary~\ref{p08:cor:flat-laplace}.
For $L^p$, an annihilator is a function $h\in L^{p'}$.
Fix $b\in(0,\infty)^d$ with $b_j<a_j$.
H\"older's inequality shows that
$e^{-\langle b,x\rangle}h(x)\,dx$ is a finite complex measure.
Its Laplace transform vanishes at $s_n-b$, so the same corollary
forces $h=0$.

Weak convergence implies convergence of these bounded continuous
test functions.  Conversely, uniform density and the common total
mass bound imply convergence of integrals against every function
in $C_0$.  This is vague convergence, and because the limit $\mu$
has total mass one it implies weak convergence.  Finally, the
summable weights and $0\le L_\mu(s_n)\le1$ show that convergence
in $d_c$ is equivalent to coordinatewise convergence of the probe
values.
\end{proof}

\begin{remark}[The anchoring step is not exclusive to flat probes]
\label{p08:rem:flat-power-comparison}
If $q_1,\ldots,q_d>0$ are linearly independent over $\mathbb Q$,
the shifted power curve $a+(t^{q_1},\ldots,t^{q_d})$ also determines
every finite measure.  To see this, substitute into the convergent
Taylor series at $a$.  The exponents $\alpha\cdot q$ are distinct
and form a locally finite set.  A nonzero series has a least
exponent and a nonzero leading coefficient.  The same argument as
above then applies.  Thus removal of tail assumptions comes from
interior anchoring.  Separated flat scales provide a different,
explicit design with no arithmetic independence requirement.
\end{remark}

\begin{proposition}[The entire endpoint jet is insufficient]
\label{p08:prop:flat-jet-blind}
For $d\ge2$, extend $c(0)=0$ in \eqref{p08:eq:flat-curve}.
For any analytic germ $A$ at $a$, the function
$A(a+c(t))$ is smooth from the right at zero and
\begin{equation}
 \left.\frac{d^k}{dt^k}A(a+c(t))\right|_{t=0+}
   =\partial_1^k A(a),\qquad k\ge0.
 \label{p08:eq:flat-jet}
\end{equation}
There exist two distinct compactly supported probability measures
whose Laplace observations along the curve have identical derivatives
of every order at zero, although their observations differ for every
sufficiently small positive $t$.
\end{proposition}

\begin{proof}
Every component of $c$ after the first is flat at zero.
By writing the difference
$A(a+c(t))-A(a+(t,0,\ldots,0))$ as the integral of the derivatives
of $A$ along the line segment between these two arguments,
one sees that the difference is a smooth flat function.  This proves
\eqref{p08:eq:flat-jet}.

For an explicit pair, take $a=(1,\ldots,1)$,
$p=1/(e+1)$, and
\[
 \mu=\delta_{(0,1,0,\ldots,0)},\qquad
 \nu=p\delta_0+(1-p)\delta_{(0,2,0,\ldots,0)}.
\]
Their transform difference is
$D(s)=e^{-s_2}-p-(1-p)e^{-2s_2}$.
It satisfies $D(a)=0$ and is independent of $s_1$, so all the
derivatives in \eqref{p08:eq:flat-jet} vanish.  On the other hand,
\[
 D(a+c(t))
 =\frac{1-e}{e(e+1)}e^{-1/t}+O(e^{-2/t}),
\]
which is nonzero for small positive $t$.
The measures have disjoint supports and total variation distance one.
\end{proof}

\begin{remark}[What is measured]
The data in Corollary~\ref{p08:cor:flat-laplace} are countably many exact
real numbers with an accumulation point, not one scalar measurement.
The corollary is a uniqueness result, not a noise-stable recovery
algorithm.  Proposition~\ref{p08:prop:flat-jet-blind} also shows that
arbitrarily accurate ordinary Taylor models can omit the information
responsible for uniqueness.
\end{remark}

\subsection{Smooth laws with identical infinite jets}
\label{p08:sec:flat-smooth-barriers}

The atomic endpoint-jet example above admits a stronger smooth-density
version with a quantified observation gap. The observations remain direct
Laplace values along the same interior flat curve.

For a probability measure $\mu$ on $[0,1]^d$, write
\[
 L_\mu(s)=\int_{[0,1]^d}e^{-s\cdot x}\,d\mu(x).
\]
Fix $a\in(0,\infty)^d$ and $d\geq2$, and consider the particular curve
\begin{equation}\label{p08:eq:barrier-flat-curve}
 \Gamma_1(t)=a_1+t,\qquad
 \Gamma_j(t)=a_j+\exp(-t^{-(j-1)})\quad(2\leq j\leq d),\qquad t>0.
\end{equation}
Set $\Gamma(0)=a$.  Every component is $C^\infty$ on $[0,\infty)$,
although the components with index at least two are not analytic at zero.

\begin{proposition}[Identical infinite jets with separated laws]
\label{p08:prop:flat-infinite-jets}
For every positive integer $M$, there exist two distinct probability measures
$\mu_{M,+},\mu_{M,-}$ with smooth densities on $[0,1]^d$, bounded between
$1/2$ and $3/2$, and constants $c_M,C_M>0$, such that
\begin{align}
 W_1(\mu_{M,+},\mu_{M,-})&\geq c_M,\label{p08:eq:flat-w1}\\
 \left.\frac{d^n}{dt^n}L_{\mu_{M,+}}(\Gamma(t))\right|_{t=0}
 &=\left.\frac{d^n}{dt^n}L_{\mu_{M,-}}(\Gamma(t))\right|_{t=0}
 &&(n\geq0),\label{p08:eq:flat-jets}\\
 \sup_{0\leq t\leq\epsilon}
 \big|L_{\mu_{M,+}}(\Gamma(t))-L_{\mu_{M,-}}(\Gamma(t))\big|
 &\leq C_M\exp(-M\epsilon^{-(d-1)})
 &&(\epsilon>0).\label{p08:eq:flat-small-signal}
\end{align}
The measures and constants in this statement do not depend on $\epsilon$.
Their one-dimensional marginals in coordinates $1,\ldots,d-1$ are identical.
\end{proposition}

\begin{proof}
In $L^2[0,1]$, let $V_M$ be the finite-dimensional span of
\[
 1,\quad e^{-a_d x},\quad x e^{-a_d x},\quad\ldots,\quad
 x^{M-1}e^{-a_d x},
\]
and let $h=x-\operatorname{proj}_{V_M}x$.  This is a smooth real-valued
function, and $h\ne0$: an identity
$x=A+P(x)e^{-a_d x}$ on $[0,1]$ would extend to all real $x$, which is
impossible as $x\to+\infty$.  Orthogonality gives
\begin{equation}\label{p08:eq:flat-orthogonality}
 \int_0^1h=0,\qquad
 \int_0^1x^k e^{-a_d x}h(x)\,dx=0\quad(0\leq k<M),\qquad
 \int_0^1x h(x)\,dx=\|h\|_2^2>0.
\end{equation}
Set $b=(2\|h\|_\infty)^{-1}$ and define densities
$f_\pm(x_1,\ldots,x_d)=1\pm b h(x_d)$.  They are probability densities
in the stated range.  The test function $x\mapsto x_d$ is $1$-Lipschitz,
so Kantorovich duality gives
$W_1(\mu_+,\mu_-)\geq2b\|h\|_2^2=:c_M$.

Put
$\Psi(s)=\int_0^1 e^{-sx}h(x)\,dx$ and
$\Phi(s)=\int_0^1e^{-sx}\,dx$.  Then
\[
 L_{\mu_+}(s)-L_{\mu_-}(s)
 =2b\Psi(s_d)\prod_{j<d}\Phi(s_j).
\]
By \eqref{p08:eq:flat-orthogonality}, $\Psi$ has a zero of order at least $M$
at $a_d$.  Taylor's formula, using $s\geq a_d>0$, yields
\[
 |\Psi(a_d+u)|\leq\frac{u^M}{M!}\int_0^1x^M|h(x)|\,dx
 \quad(u\geq0).
\]
Since $0<\Phi(s)\leq1$ for $s\geq0$, this proves
\eqref{p08:eq:flat-small-signal} with
$C_M=2b(M!)^{-1}\int_0^1x^M|h(x)|\,dx$.
The function $\Psi(a_d+u)$ factors as $u^M$ times an analytic function.
Substituting $u=\exp(-t^{-(d-1)})$ therefore gives a smooth function flat
at zero.  The other factors are smooth there, proving
\eqref{p08:eq:flat-jets}.
\end{proof}

\begin{corollary}[A horizon-dependent noise obstruction]
\label{p08:cor:flat-horizon}
Suppose the entire function $L_\mu(\Gamma(t))$, $0\leq t\leq\epsilon$,
is observed with adversarial error at most $\delta$ in supremum norm.
If
\begin{equation}\label{p08:eq:flat-noise-threshold}
 2\delta\geq C_M\exp(-M\epsilon^{-(d-1)}),
\end{equation}
then the worst-case $W_1$ reconstruction error over smooth densities
$1/2\leq f\leq3/2$ is at least $c_M/2$.
\end{corollary}

\begin{proof}
The midpoint of the two observation curves lies within $\delta$ of each
curve.  Any single reconstruction from this common noisy observation has
distance at least half the separation from one of the two laws.
\end{proof}

In particular, even complete knowledge of the Taylor jet at the base point
need not identify a law that the full curve identifies.  The constants
$c_M$ may shrink with $M$; no uniform lower bound in $M$ is asserted.
Also, this result varies the accessible horizon $\epsilon$.  It does not
contradict continuity of the inverse on a fixed compact class when a fixed
identifying curve segment is observed exactly.

\section{A finite resolution theorem for the hidden dimension}
\label{p08:sec:finite-blind}

The exact gate inverse and finite noisy sampling are different observation
models.  In this section the available data are noisy values of the
\emph{Laplace profile}
\[
 F(t)=\mathbb E\exp\left(-\sum_{i=1}^d X_i t^{q_i}\right),
 \qquad 0<t\leq1.
\]
No stability assertion for recovering these values from noisy gate data
is implicit.  That additional step is supplied by
Sections~\ref{p08:sec:gate-jet-stability}--\ref{p08:sec:gate-value-stability}.  The result below is an application of classical approximate Prony
and matrix-pencil ideas; compare \cite{PottsTasche2010,PottsTasche2013}.
Its purpose is to give a complete finite resolution guarantee for the
present nonlinear random-vector model, including its approximation error.

\begin{theorem}[Finite noisy recovery of dimension and powers]
\label{p08:thm:finite-blind}
Fix an integer $D\geq1$, numbers $0<a\leq b<2a$, $m>0$, $S>0$,
and $0<\Delta\leq b$.  Consider all nonnegative random vectors with
unknown dimension $1\leq d\leq D$ such that
\begin{equation}\label{p08:eq:finite-blind-class}
 a\leq q_i\leq b,\qquad |q_i-q_j|\geq\Delta\ (i\ne j),\qquad
 \mathbb E X_i\geq m,\qquad
 \mathbb E\left(\sum_iX_i\right)^2\leq S.
\end{equation}
The separation condition is vacuous when $d=1$.  In particular, no
rational independence of the powers is assumed.  Put
\begin{equation}\label{p08:eq:finite-blind-constants}
 \zeta=\frac{\Delta}{eb},\qquad
 \kappa=\frac1D\left(\frac{\zeta}{2}\right)^{2D-2}.
\end{equation}
Choose $0<\tau\leq1$ and suppose that the $2D$ real observations satisfy
\[
 t_j=\tau e^{-j/b},\qquad
 |\widetilde F_j-F(t_j)|\leq\delta,
 \qquad j=0,\ldots,2D-1.
\]
Define
\begin{equation}\label{p08:eq:finite-blind-error}
 \epsilon=\delta+\frac S2\tau^{2a},\qquad
 E=D\epsilon,\qquad s=m\kappa\tau^b.
\end{equation}
If $E<s/2$, then the number of eigenvalues greater than $s/2$ of
\[
 \widetilde H_D=(1-\widetilde F_{j+k})_{j,k=0}^{D-1}
\]
is exactly $d$.  Thus the first $2D-1$ observations already determine
the hidden dimension throughout this class.

After recovering $d$, form the $d$ by $d$ matrices
\[
 \widetilde H_0=(1-\widetilde F_{j+k})_{j,k=0}^{d-1},\qquad
 \widetilde H_1=(1-\widetilde F_{j+k+1})_{j,k=0}^{d-1}.
\]
The first matrix is positive definite.  There is a matching of the
generalized eigenvalues $\widetilde z_i$ of
$(\widetilde H_1,\widetilde H_0)$ to $z_i=e^{-q_i/b}$ for which
\begin{equation}\label{p08:eq:finite-blind-node-error}
 |\widetilde z_i-z_i|\leq
 \eta:=\frac{2E}{s-E}.
\end{equation}
If also $\eta<1/(2e)$, the estimates
$\widetilde q_i=-b\log\widetilde z_i$ are real and satisfy
\begin{equation}\label{p08:eq:finite-blind-power-error}
 \min_{\pi}\max_i|\widetilde q_i-q_{\pi(i)}|
 \leq 2eb\eta.
\end{equation}
Here the minimum ranges over permutations.
\end{theorem}

\begin{proof}
For $u\geq0$ one has
$0\leq u-(1-e^{-u})\leq u^2/2$.  Since $t\leq1$ and $q_i\geq a$,
\begin{equation}\label{p08:eq:finite-blind-linearization}
 1-F(t)=\sum_i\mu_i t^{q_i}+R(t),\qquad
 |R(t)|\leq \frac S2t^{2a},\qquad \mu_i=\mathbb E X_i.
\end{equation}
Set $w_i=\mu_i\tau^{q_i}$ and $z_i=e^{-q_i/b}$.  The principal
part of the sampled sequence is therefore
\[
 y_j=\sum_i w_i z_i^j,\qquad w_i\geq m\tau^b.
\]
The total entrywise error in $1-\widetilde F_j$ is at most $\epsilon$.
The nodes lie in $[e^{-1},e^{-a/b}]$, and the mean value theorem gives
$|z_i-z_j|\geq\zeta$.

We record an elementary Vandermonde bound.  Let
$V_d=(z_i^j)_{0\leq j\leq d-1,\,1\leq i\leq d}$.
The coefficients of each Lagrange polynomial
\[
 \ell_i(z)=\prod_{k\ne i}\frac{z-z_k}{z_i-z_k}
\]
have total absolute value at most $(2/\zeta)^{d-1}$.
These coefficient vectors are the rows of $V_d^{-1}$.
Consequently
\[
 \|V_d^{-1}\|_\infty\leq(2/\zeta)^{d-1},\qquad
 \|V_d^{-1}\|_1\leq d(2/\zeta)^{d-1},
\]
and hence
\begin{equation}\label{p08:eq:finite-blind-vandermonde}
 \sigma_{\min}(V_d)^2\geq
 \frac1d(\zeta/2)^{2d-2}\geq\kappa.
\end{equation}
Adding rows cannot decrease the smallest singular value of a matrix
with full column rank, so the same bound applies to
$V_D=(z_i^j)_{0\leq j\leq D-1,\,1\leq i\leq d}$.

The ideal matrix
$H_D=(y_{j+k})_{j,k=0}^{D-1}=V_D\operatorname{diag}(w_i)V_D^T$
is positive semidefinite of rank $d$, and each nonzero eigenvalue is
at least $s$.  The perturbation has operator norm at most $D\epsilon=E$.
Weyl's inequality now puts exactly $d$ eigenvalues of
$\widetilde H_D$ above $s/2$.

For the recovered size $d$, the ideal matrices satisfy
\[
 H_0=V_d\operatorname{diag}(w_i)V_d^T,\qquad
 H_1=V_d\operatorname{diag}(w_i z_i)V_d^T.
\]
Thus $H_0\succeq sI$, their generalized eigenvalues are the $z_i$,
and each perturbation has norm at most $d\epsilon\leq E$.
For a nonzero real vector $v$, write
$r(v)=v^TH_1v/(v^TH_0v)$.  This quotient lies in $[0,1]$.
Its perturbed version obeys
\[
 |\widetilde r(v)-r(v)|
 \leq\frac{\|\widetilde H_1-H_1\|+
              |r(v)|\|\widetilde H_0-H_0\|}{s-E}
 \leq\frac{2E}{s-E}.
\]
The generalized Courant--Fischer principle gives
\eqref{p08:eq:finite-blind-node-error} in sorted order.  When
$\eta<1/(2e)$, the true and estimated nodes are at least $1/(2e)$.
The derivative of $-b\log z$ on this interval is bounded by $2eb$,
which proves \eqref{p08:eq:finite-blind-power-error}.
\end{proof}

\begin{corollary}[An explicit noise threshold]
\label{p08:cor:finite-blind-noise}
Under the assumptions of Theorem~\ref{p08:thm:finite-blind}, suppose
$0<\delta\leq S$ and choose
\[
 \tau=(\delta/S)^{1/(2a)},\qquad
 \alpha=1-\frac b{2a}>0.
\]
A sufficient condition for exact dimension recovery is
\begin{equation}\label{p08:eq:finite-blind-threshold}
 \delta^\alpha<\frac{m\kappa}{4D S^{b/(2a)}}.
\end{equation}
For all sufficiently small $\delta$, the same $2D$ samples give
\[
 \min_\pi\max_i|\widetilde q_i-q_{\pi(i)}|
 \leq \frac{16ebD S^{b/(2a)}}{m\kappa}\,
                 \delta^{\,1-b/(2a)}.
\]
The last assertion is an upper bound for this procedure, not a claim of
minimax optimality.
\end{corollary}
\begin{proof}
For the displayed choice of $\tau$, one has
$\epsilon=3\delta/2\leq2\delta$, and therefore
\[
 \frac Es\leq
 \frac{2D S^{b/(2a)}}{m\kappa}\delta^\alpha.
\]
Condition \eqref{p08:eq:finite-blind-threshold} makes this ratio smaller
than $1/2$.  Then $\eta=2E/(s-E)\leq4E/s$.
As $\delta\downarrow0$, the additional logarithm condition is
eventually satisfied, and the claimed estimate follows.
\end{proof}

\begin{corollary}[Positive tilting supplies the moment bound]
\label{p08:cor:finite-blind-tilt}
Let $P$ be any probability law of a finite nonnegative vector, without
assuming that any positive moment of $P$ is finite.  Write
$T=\sum_iX_i$.  Fix a known baseline $c\geq c_0>0$ and suppose
\[
 \nu(dx)=e^{-c\sum_i x_i}P(dx),\qquad Z=\nu(\mathbb R_+^d)\geq z_0>0.
\]
Consider the normalized tilted profile
\[
 F_c(t)=\frac1Z\int
       \exp\left(-\sum_i t^{q_i}x_i\right)\nu(dx).
\]
If $1\leq d\leq D$, the powers satisfy the band and separation assumptions of
Theorem~\ref{p08:thm:finite-blind}, and each tilted coordinate mean satisfies
$\int x_i\,\nu(dx)/Z\geq m$, that theorem applies with the explicit choice
\begin{equation}\label{p08:eq:finite-blind-tilted-moment}
 S=\frac{4}{e^2c_0^2z_0}.
\end{equation}
In particular, the finite dimension and power guarantees can hold even
when every positive moment of the original law $P$ is infinite.

If noisy unnormalized profile values have error at most $\delta_F$,
and a mass estimate satisfies $|\widetilde Z-Z|\leq\delta_Z\leq z_0/2$,
then dividing those observations by $\widetilde Z$ gives the normalized
error bound
\[
 \delta=\frac{2(\delta_F+\delta_Z)}{z_0}.
\]
These are errors at the tilted-profile and mass stages.
Section~\ref{p08:sec:end-to-end-jet} connects them to errors in gate measurements.
\end{corollary}
\begin{proof}
The normalized tilted law satisfies
\[
 \frac1Z\mathbb E_P[T^2e^{-cT}]
 \leq\frac1Z\sup_{u\geq0}u^2e^{-cu}
 =\frac4{e^2c^2Z}\leq\frac4{e^2c_0^2z_0}.
\]
It is therefore in the class of Theorem~\ref{p08:thm:finite-blind}.
For the error claim, write the exact unnormalized profile value as $A$,
so that $0\leq A\leq Z$.  Since $\widetilde Z\geq z_0/2$,
\[
 \left|\frac{\widetilde A}{\widetilde Z}-\frac A Z\right|
 \leq\frac{\delta_F}{\widetilde Z}
       +\frac{A\delta_Z}{Z\widetilde Z}
 \leq\frac{2(\delta_F+\delta_Z)}{z_0}.
\]
\end{proof}

\begin{remark}[Scope and the role of the gap]
The class includes arbitrarily dependent coordinates and unbounded laws
with the indicated second-moment bound.  Compact support with
$\sum_iX_i\leq M$ implies the hypothesis with $S=M^2$.
The condition $b<2a$ separates the original powers from all quadratic
and higher powers in the expansion of $F$.  It allows recovery of
dimension and powers even when the full joint law is not identifiable
because of a higher-order rational resonance.  For example, the powers
$2$ and $3$ lie below the first quadratic threshold $4$, although
$3\cdot2=2\cdot3$ produces a higher-order resonance.
The condition is sufficient for this simple method; its failure is not
an impossibility theorem.  The constants are deliberately conservative.
Neither the Vandermonde mechanism nor the matrix pencil is asserted
to be new.
\end{remark}

\begin{proposition}[Two obstructions despite a positive mean floor]
\label{p08:prop:finite-blind-obstructions}
Fix $m>0$.  The observation in this proposition is the direct profile
$F(t)$ on $0<t\leq1$.
\begin{enumerate}
\item If no uniform tail bound is imposed, a positive lower bound on
every coordinate mean does not make dimension uniformly stable, even
with fixed separated rationally independent powers.  More precisely,
append to any fixed admissible vector an independent coordinate which
equals $m/p$ with probability $p$ and zero otherwise, using a fixed new
power $q_*>0$.  Its mean is $m$, whereas the profiles satisfy
\[
 \sup_{0<t\leq1}|F_p(t)-F(t)|\leq p\longrightarrow0.
\]
\item If power separation is removed, dimension is not uniformly
stable even for deterministic bounded vectors with the same positive
mean floor and a uniform second-moment bound.  Fix $q>0$.  Compare
the one-coordinate model $(X,q)=(2m,q)$ with the two-coordinate model
$(X_1,X_2)=(m,m)$ and powers $(q,q+u)$, where $u\downarrow0$ and
$u/q$ is irrational.  Their profiles satisfy
\[
 \sup_{0<t\leq1}|F_u(t)-F_0(t)|\leq\frac{mu}{eq}.
\]
\end{enumerate}
\end{proposition}
\begin{proof}
For the first assertion, independence gives
\[
 F_p(t)=F(t)\bigl(1-p+p e^{-(m/p)t^{q_*}}\bigr).
\]
Since $0\leq F(t)\leq1$, the difference is at most $p$.
The new coordinate has second moment $m^2/p$,
so this example does not satisfy the uniform second-moment hypothesis
of Theorem~\ref{p08:thm:finite-blind}.

For the second assertion, on $0<t\leq1$ the difference between
the two exponents inside the Laplace transform is at most
$mu\,t^q|\log t|\leq mu/(eq)$.
The map $v\mapsto e^{-v}$ is $1$-Lipschitz on $[0,\infty)$,
which proves the profile bound.
Both models have total mass $2m$ deterministically,
and therefore the common second-moment bound is $S=4m^2$.
\end{proof}

\section{Finite gate jets control the entire ordered profile}
\label{p08:sec:gate-jet-stability}

This section provides a quantitative bridge from finitely many gate
derivatives to the function encoded by the gate experiment.  It is a
deterministic pairwise stability statement.  Derivatives are measured in
the factorial normalization below; noise in an unnormalized high derivative
is a different observation model.

Fix constants $0<\beta<1$, $0<C_-\leq C_+<\infty$, and $L>0$.
An admissible pair $(C,f)$ consists of $C\in[C_-,C_+]$ and a
nondecreasing $L$-Lipschitz function $f:[0,1]\to[\beta,1]$.  Strict
increase, though available in the applications, is not needed here.  Define
\begin{equation}\label{p08:eq:finite-gate-J}
 G_{C,f}(\theta)=C\int_0^1-\log(1-\theta f(t))\,dt,\qquad
 J_k(C,f)=\frac{G_{C,f}^{(k)}(0)}{(k-1)!}
 =C\int_0^1 f(t)^k\,dt\quad(k\geq1).
\end{equation}

\begin{theorem}[A finite-jet modulus for gate inversion]
\label{p08:thm:finite-gate-modulus}
There is an absolute constant $K>0$ with the following property;
$K=1024$ is an admissible choice.
Let $(C,f)$ and $(\widetilde C,\widetilde f)$ be admissible, and suppose
for some integer $N\geq2$ and $\eta\geq0$ that
\begin{equation}\label{p08:eq:gate-jet-noise}
 |J_k(C,f)-J_k(\widetilde C,\widetilde f)|\leq\eta
 \qquad(1\leq k\leq N).
\end{equation}
Set
\begin{align}
 D_N&=2C_+(1-\beta)^N+(2^N-1)\eta,\label{p08:eq:gate-DN}\\
 E_N&=\frac{1}{C_-}\left\{
 \frac{2C_+K}{N}+5\,6^N\eta+2C_+(1-\beta)^N\right\}.
 \label{p08:eq:gate-EN}
\end{align}
Then
\begin{equation}\label{p08:eq:gate-profile-modulus}
 |C-\widetilde C|\leq D_N,\qquad
 \|f-\widetilde f\|_\infty
 \leq 2E_N+\sqrt{2LE_N}.
\end{equation}
In particular, if $0<\eta\leq6^{-4}$ and
\begin{equation}\label{p08:eq:gate-Nchoice}
 N=\left\lfloor\frac{\log(1/\eta)}{2\log6}\right\rfloor,
\end{equation}
then, as $\eta\downarrow0$ with the model constants fixed,
\begin{equation}\label{p08:eq:gate-log-rate}
 |C-\widetilde C|+\|f-\widetilde f\|_\infty
 =O\big((\log(1/\eta))^{-1/2}\big).
\end{equation}
No optimality of this modulus is asserted.
\end{theorem}

\begin{lemma}[Polynomial approximation with controlled coefficients]
\label{p08:lem:gate-polynomial-approx}
The absolute constant $K$ can be chosen so that for every $1$-Lipschitz
$h:[0,1]\to\mathbb R$ with $h(0)=0$ and every integer $N\geq2$, there
is a polynomial $P(u)=\sum_{k=1}^N b_k u^k$ such that
\begin{equation}\label{p08:eq:gate-polynomial-bounds}
 \|h-P\|_\infty\leq K/N,\qquad
 \sum_{k=1}^N|b_k|\leq4\,6^N.
\end{equation}
\end{lemma}

\begin{proof}
We include a standard Jackson-kernel construction to specify the coefficient
cost needed here.  The even $2\pi$-periodic function
$H(v)=h((1+\cos v)/2)$ has supremum norm at most one and Lipschitz constant
at most $1/2$ with respect to circular distance.  For an integer $n\geq2$,
normalize
\[
 \left(\frac{\sin(nv/2)}{\sin(v/2)}\right)^4
\]
to a nonnegative trigonometric polynomial $K_n$ with
$(2\pi)^{-1}\int_{-\pi}^{\pi}K_n(v)\,dv=1$.
Its degree is $2(n-1)$.  The estimates
\[
 \left|\frac{\sin(nv/2)}{\sin(v/2)}\right|
 \leq\min\{n,\pi/|v|\},\qquad
 \left|\frac{\sin(nv/2)}{\sin(v/2)}\right|
 \geq2n/\pi\quad(|v|\leq1/n)
\]
show that the unnormalized integral is at least $32n^3/\pi^4$,
while its first absolute moment is at most $(1+\pi^4)n^2$.
Consequently
\[
 \frac1{2\pi}\int_{-\pi}^{\pi}|v|K_n(v)\,dv\leq A/n
\]
with $A=\pi^4(1+\pi^4)/32<300$. Convolution on the circle gives an even trigonometric
polynomial $Q_n=H*K_n$ with
$\|Q_n-H\|_\infty\leq A/(2n)$ and $\|Q_n\|_\infty\leq1$.

Take $n=\lfloor N/2\rfloor+1$.  Write
$Q_n(v)=a_0+\sum_{k=1}^{2(n-1)}a_k\cos(kv)$, where
$|a_0|\leq1$ and $|a_k|\leq2$.  Substituting
$\cos(kv)=T_k(2u-1)$ produces a polynomial $P_0$ of degree at most $N$
with $\|P_0-h\|_\infty\leq A/N$.  Here $T_k$ is the Chebyshev
polynomial.  All its shifted roots lie in $(0,1)$, so the coefficients of
$T_k(2u-1)$ alternate in sign.  Thus their sum of absolute values is
\[
 |T_k(-3)|=T_k(3)
 =\tfrac12\big((3+2\sqrt2)^k+(3-2\sqrt2)^k\big)
 \leq(3+2\sqrt2)^k.
\]
It follows that the sum of absolute values of the monomial coefficients of
$P_0$ is at most
$1+2\sum_{k=1}^N(3+2\sqrt2)^k\leq4\,6^N$.
Finally set $P=P_0-P_0(0)$.  Its constant coefficient vanishes, its
nonconstant coefficients are unchanged, and
$\|P-h\|_\infty\leq2A/N$. This proves the lemma with $K=2A$,
and hence with the convenient choice $K=1024$.
\end{proof}

\begin{proof}[Proof of Theorem~\ref{p08:thm:finite-gate-modulus}]
Let $\nu=f_\#dt$ and $\widetilde\nu=\widetilde f_\#dt$ be probability
measures on $[\beta,1]$, and set $\sigma=C\nu$,
$\widetilde\sigma=\widetilde C\widetilde\nu$.  The observed values
$J_k$ are the nonzero-order moments of these finite measures.

First, the polynomial $p_N(u)=1-(1-u)^N$ has no constant coefficient and
the sum of absolute values of its coefficients is $2^N-1$.  Moreover,
\[
 0\leq C-\int p_N\,d\sigma
 =\int(1-u)^N\,d\sigma\leq C_+(1-\beta)^N.
\]
The analogous estimate for $\widetilde\sigma$, followed by
\eqref{p08:eq:gate-jet-noise}, proves $|C-\widetilde C|\leq D_N$.

Next, for any $1$-Lipschitz $h$ with $h(0)=0$, apply
Lemma~\ref{p08:lem:gate-polynomial-approx} to obtain
\[
 \left|\int h\,d\sigma-\int h\,d\widetilde\sigma\right|
 \leq\frac{2C_+K}{N}+4\,6^N\eta.
\]
Since $\|h\|_\infty\leq1$, division by $C$ and comparison with
$\widetilde C$ gives
\[
 \left|\int h\,d\nu-\int h\,d\widetilde\nu\right|
 \leq\frac1{C_-}\left\{\frac{2C_+K}{N}+4\,6^N\eta+D_N\right\}
 \leq E_N.
\]
Kantorovich duality on $[0,1]$ therefore gives
$W_1(\nu,\widetilde\nu)\leq E_N$.  Notice that approximation was
performed before normalizing the measures.  Normalizing their moments
first would unnecessarily amplify the mass error.

Because $f$ and $\widetilde f$ are continuous and nondecreasing, they are
quantile representatives of their pushforward laws; hence
\begin{equation}\label{p08:eq:gate-quantile-distance}
 W_1(\nu,\widetilde\nu)=\int_0^1|f(t)-\widetilde f(t)|\,dt.
\end{equation}
The difference $r=f-\widetilde f$ is $L$-Lipschitz, rather than merely
$2L$-Lipschitz: for $s<t$, both increments belong to $[0,L(t-s)]$.
If $b=\|r\|_\infty$ is attained at $t_0$, then
$|r(t)|\geq(b-L|t-t_0|)_+$.  Integrating this lower bound over $[0,1]$
is minimized when $t_0$ is an endpoint.  It follows that
\[
 W_1(\nu,\widetilde\nu)\geq
 \begin{cases}
 b^2/(2L),&b\leq L,\\
 b-L/2\geq b/2,&b>L.
 \end{cases}
\]
Thus $b\leq\sqrt{2LE_N}+2E_N$, proving
\eqref{p08:eq:gate-profile-modulus}.

Under \eqref{p08:eq:gate-Nchoice}, $6^N\eta\leq\eta^{1/2}$ and
$(1-\beta)^N=O(\eta^{a_\beta})$, where
$a_\beta=-\log(1-\beta)/(2\log6)>0$.
Therefore $E_N=O(1/\log(1/\eta))$, while $D_N$ decays as a positive
power of $\eta$.  This proves \eqref{p08:eq:gate-log-rate}.
\end{proof}

\begin{remark}[Interpretation of noisy observations]
If measured normalized jets $\widehat J_k$ have error at most $\eta_0$,
any two admissible models fitting the data to that tolerance satisfy
\eqref{p08:eq:gate-jet-noise} with $\eta=2\eta_0$.  The theorem bounds the
diameter of this feasible set.  It does not estimate derivatives from
ordinary noisy values of $G$; that step is supplied in
Section~\ref{p08:sec:gate-value-stability}. It does not assert that a large gate jet
can be acquired at a cost independent of its order.  With exact jets,
$\eta=0$, the same proof gives profile diameter $O(N^{-1/2})$ as
$N\to\infty$.
\end{remark}

\begin{remark}[A tail-free source of the Lipschitz and positivity bounds]
Suppose $X_i\geq0$, $Z=\sum_iX_i<\infty$ almost surely,
$a_i\geq a_0>0$, and $1\leq q_i\leq v<\infty$.  For
\[
 f(t)=1-\mathbb E\exp\left(-\sum_i(a_i+t^{q_i})X_i\right),
 \qquad 0\leq t\leq1,
\]
finite or countably many coordinates are allowed, provided the defining
exponent is meaningful.  The estimate
\[
 0\leq f'(t)
 \leq v\,\mathbb E[Z e^{-a_0Z}]
 \leq\frac{v}{e a_0}
\]
gives a common Lipschitz constant without any ordinary moment assumption.
One may justify the inequality first for finite sums and then pass to the
limit by monotone convergence of the Laplace integrands; the Lipschitz
bound is preserved.  If $Z\geq z_0>0$ almost surely, then
$f(t)\geq1-e^{-a_0z_0}$, providing an admissible $\beta$.
If $Z$ has no positive deterministic lower bound, a separately specified
uniform lower bound on $f(0)$ suffices instead.
\end{remark}

\section{From finite gate jets to stable hidden structure}
\label{p08:sec:end-to-end-jet}

The preceding finite-profile theorem can be combined with the gate-jet
estimate to obtain a guarantee at the gate observation stage itself.
The observations here are normalized derivatives of the gate curve at
zero, not finitely many values of that curve.  These are derivatives in
the gate parameter $\theta$, rather than the
endpoint derivatives in the profile parameter $t$ considered in the
infinite-jet blindness example. Section~\ref{p08:sec:gate-value-stability}
then controls these gate derivatives from ordinary noisy gate values.

\begin{corollary}[Finite gate jets resolve dimension and powers]
\label{p08:cor:gate-finite-structure}
Fix an integer $D\geq1$, $1\leq u\leq v<2u$, $0<\Delta\leq v$, $m>0$,
$c_* >0$, $0<\beta<1$, $0<z_*\leq1-\beta$, and $0<C_-\leq C_+<\infty$.
Consider models consisting of a finite nonnegative random vector $X$,
powers $q_i$, a baseline $c$, and a gain $C$, satisfying
\begin{gather}
 1\leq d\leq D,\qquad u\leq q_i\leq v,\qquad
 |q_i-q_j|\geq\Delta\quad(i\ne j),\qquad
 c\geq c_*,\qquad C_-\leq C\leq C_+,\label{p08:eq:jet-structure-class}\\
 T=\sum_iX_i,\qquad Z=\mathbb E e^{-cT}\in[z_*,1-\beta],\qquad
 \frac{\mathbb E[X_i e^{-cT}]}{Z}\geq m.\label{p08:eq:jet-tilted-class}
\end{gather}
No positive moment of the original law is assumed finite.  Define
\[
 A(t)=\mathbb E\exp\left(-cT-\sum_iX_it^{q_i}\right),\qquad
 f(t)=1-A(t),\qquad
 G(\theta)=C\int_0^1-\log(1-\theta f(t))\,dt.
\]
Let $N\geq2$ be an integer and let $\eta\geq0$. Suppose two models in this class have
normalized gate jets satisfying
\begin{equation}\label{p08:eq:jet-pair-data}
 \left|\frac{G_1^{(k)}(0)-G_2^{(k)}(0)}{(k-1)!}\right|\leq\eta,
 \qquad 1\leq k\leq N.
\end{equation}
Subscripts $1,2$ distinguish the models.  Set
\[
 L_* =\frac{v}{ec_*},\qquad S_* =\frac{4}{e^2c_*^2z_*},\qquad
 \kappa=\frac1D\left(\frac{\Delta}{2ev}\right)^{2D-2},\qquad
 \alpha=1-\frac{v}{2u}>0.
\]
Let $E_N$ be the bound in \eqref{p08:eq:gate-EN} of
Theorem~\ref{p08:thm:finite-gate-modulus}, with intensity bounds $C_-,C_+$,
floor $\beta$, and moment error $\eta$.
Put
\begin{equation}\label{p08:eq:jet-profile-error}
 h_N=2E_N+\sqrt{2L_*E_N},\qquad \delta_N=\frac{2h_N}{z_*}.
\end{equation}
If
\begin{equation}\label{p08:eq:jet-structure-threshold}
 0<\delta_N\leq S_*,\qquad
 \delta_N^\alpha<\frac{m\kappa}{16eD S_*^{v/(2u)}},
\end{equation}
then the dimensions of the two models are equal, and their powers obey
\begin{equation}\label{p08:eq:jet-power-resolution}
 \min_\pi\max_i|q_{1,i}-q_{2,\pi(i)}|
 \leq\frac{32evD S_*^{v/(2u)}}{m\kappa}\,\delta_N^\alpha.
\end{equation}
For fixed class parameters and
$N=\lfloor\log(1/\eta)/(2\log6)\rfloor$, all sufficiently small
$\eta>0$ satisfy these conditions, and
\begin{equation}\label{p08:eq:jet-log-rate}
 \min_\pi\max_i|q_{1,i}-q_{2,\pi(i)}|
 =O\left((\log(1/\eta))^{-\alpha/2}\right).
\end{equation}
Thus sufficiently accurate finite gate jets determine the hidden dimension
throughout this quantitative class and constrain its powers with the
stated modulus.
\end{corollary}

\begin{proof}
For $0\leq t\leq1$, the function $f$ is nondecreasing and takes values
in $[\beta,1]$.  For $0<t\leq1$, differentiation under the expectation
is justified by the bounded envelope $vT e^{-cT}$ and gives
\[
 0\leq f'(t)
 =\mathbb E\left[e^{-cT-\sum_iX_it^{q_i}}
                   \sum_iq_iX_it^{q_i-1}\right]
 \leq v\mathbb E[T e^{-cT}]\leq\frac{v}{ec_*}=L_*.
\]
Continuity at zero extends this Lipschitz bound to the closed interval.
Theorem~\ref{p08:thm:finite-gate-modulus} therefore gives
$\|f_1-f_2\|_\infty\leq h_N$.
Consequently $\|A_1-A_2\|_\infty\leq h_N$ and
$|Z_1-Z_2|=|A_1(0)-A_2(0)|\leq h_N$.
Since $0\leq A_2(t)\leq Z_2$ and $Z_1,Z_2\geq z_*$,
\[
 \left\|\frac{A_1}{Z_1}-\frac{A_2}{Z_2}\right\|_\infty
 \leq\frac{2h_N}{z_*}=\delta_N.
\]

Both normalized profiles are Laplace profiles of the exponentially tilted
laws.  These laws have coordinate means at least $m$, and the elementary
bound on $T^2e^{-cT}$ gives their common second-moment bound $S_*$.
Apply Theorem~\ref{p08:thm:finite-blind} and
Corollary~\ref{p08:cor:finite-blind-noise}, with
$(a,b,S,\delta)=(u,v,S_*,\delta_N)$ and
$\tau=(\delta_N/S_*)^{1/(2u)}$.
Choose as common sample data the exact values of the first normalized
profile at the prescribed $2D$ sample sites.  Their error is zero for the
first model and at most $\delta_N$ for the second.
The common Hankel rank estimate must therefore equal the dimension of
both models.

For completeness, the ratio of Hankel perturbation size to its signal
lower bound in this application is at most
\[
 \frac{2D S_*^{v/(2u)}}{m\kappa}\delta_N^\alpha<\frac1{8e}.
\]
It is below $1/2$, and the node-error bound is less than $1/(2e)$.
Thus the power estimates are real, and each of the two true power lists
is within
$16evD S_*^{v/(2u)}\delta_N^\alpha/(m\kappa)$ of the same estimated
list, after sorting.  The triangle inequality proves
\eqref{p08:eq:jet-power-resolution}.

With the stated choice of $N$, the gate-jet estimate is
$E_N=O((\log(1/\eta))^{-1})$.  Therefore
$\delta_N=O((\log(1/\eta))^{-1/2})$, which proves the asymptotic
claim and eventually satisfies the two strict small-error conditions.
\end{proof}

The class assumptions are quantitative visibility assumptions.  In
particular, the gain floor, the tilted mass floor, the tilted mean floor,
and the separation of powers exclude distinct ways in which hidden
components can become arbitrarily difficult to see.  The corollary gives
stability of dimension and powers; it does not assert recovery of the
whole law when the powers have rational relations.  It is a consistency
statement: if each of two models fits the same observed gate jets to
accuracy $\eta/2$, their mutual jet discrepancy is at most $\eta$, so the
conclusion applies.  No minimax optimality of the logarithmic exponent
is asserted.

\section{Ordinary gate values give a finite-data guarantee}
\label{p08:sec:gate-value-stability}

The gate derivatives in the preceding section can themselves be controlled
from finitely many ordinary gate values.  A deliberately conservative
forward-difference construction suffices.  Its sampling locations depend on
the noise tolerance and cluster near zero; no numerical efficiency or
optimality is asserted.

\begin{theorem}[A finite-value modulus for gate inversion]
\label{p08:thm:finite-gate-values}
Retain the admissible class of Theorem~\ref{p08:thm:finite-gate-modulus}:
$C\in[C_-,C_+]$ and $f:[0,1]\to[\beta,1]$ is nondecreasing and
$L$-Lipschitz. Let $N\geq2$ be an integer, $0<h\leq N^{-2}$, and $\delta\geq0$.
For two admissible gate curves $G_1,G_2$, suppose
\begin{equation}\label{p08:eq:gate-value-data}
 |G_1(jh)-G_2(jh)|\leq\delta\qquad(1\leq j\leq N).
\end{equation}
The zeroth value $G_1(0)=G_2(0)=0$ is known exactly.  Define
\begin{equation}\label{p08:eq:gate-value-jet-error}
 \eta_{N,h,\delta}
 =16C_+N^2h+2^N\delta h^{-N}.
\end{equation}
Then the normalized gate jets satisfy
\begin{equation}\label{p08:eq:gate-value-to-jet}
 \left|\frac{G_1^{(k)}(0)-G_2^{(k)}(0)}{(k-1)!}\right|
 \leq\eta_{N,h,\delta}\qquad(1\leq k\leq N).
\end{equation}
Consequently, with $D_N,E_N$ defined by
\eqref{p08:eq:gate-DN}--\eqref{p08:eq:gate-EN} using
$\eta=\eta_{N,h,\delta}$,
\begin{equation}\label{p08:eq:gate-value-profile-bound}
 |C_1-C_2|\leq D_N,\qquad
 \|f_1-f_2\|_\infty\leq2E_N+\sqrt{2LE_N}.
\end{equation}

In particular, if $0<\delta\leq e^{-64}$, choose
\begin{equation}\label{p08:eq:gate-value-design}
 N=\left\lfloor\frac{\sqrt{\log(1/\delta)}}4\right\rfloor,
 \qquad h=e^{-4N}.
\end{equation}
For fixed class parameters, these $N$ positive gate settings give
\begin{equation}\label{p08:eq:gate-value-log-rate}
 |C_1-C_2|+\|f_1-f_2\|_\infty
 =O\left((\log(1/\delta))^{-1/4}\right)
 \qquad(\delta\downarrow0).
\end{equation}
\end{theorem}

\begin{proof}
Write
\[
 \Delta_h^k G(0)
 =\sum_{j=0}^k(-1)^{k-j}\binom kj G(jh).
\]
Repeated use of the fundamental theorem of calculus gives
\begin{equation}\label{p08:eq:gate-difference-integral}
 \frac{\Delta_h^kG(0)}{h^k}
 =\int_{[0,1]^k}G^{(k)}\bigl(h(u_1+\cdots+u_k)\bigr)
          \,du_1\cdots du_k.
\end{equation}
Every gate curve in the class is smooth on $[0,Nh]$, since $Nh\leq1/2$.
Differentiation of its defining integral gives
\[
 \frac{G^{(k+1)}(\theta)}{(k-1)!}
 =Ck\int_0^1\frac{f(t)^{k+1}}{(1-\theta f(t))^{k+1}}\,dt
 \leq C_+k(1-Nh)^{-(k+1)}
 \quad(0\leq\theta\leq Nh).
\]
The mean value theorem in \eqref{p08:eq:gate-difference-integral} therefore
bounds the forward-difference bias by
\begin{equation}\label{p08:eq:gate-difference-bias}
 \left|\frac{\Delta_h^kG(0)}{h^k(k-1)!}
              -\frac{G^{(k)}(0)}{(k-1)!}\right|
 \leq C_+k^2h(1-Nh)^{-(k+1)}\leq8C_+N^2h.
\end{equation}
For the last inequality, use
$(1-Nh)^{-(k+1)}\leq(1-1/N)^{-(N+1)}\leq8$ for $N\geq2$.
For example, the logarithm of the middle expression, viewed as a
function of real $N\geq2$, is decreasing: its derivative is
$\log(1+1/(N-1))-(N+1)/(N(N-1))<0$.

The discrepancy between the two normalized forward differences is at most
\[
 \frac{2^k\delta}{h^k(k-1)!}\leq2^N\delta h^{-N}.
\]
Combining this with the two bias bounds proves
\eqref{p08:eq:gate-value-to-jet}.  Theorem~\ref{p08:thm:finite-gate-modulus}
then proves \eqref{p08:eq:gate-value-profile-bound}.

For \eqref{p08:eq:gate-value-design}, one has $N\geq2$ and $h\leq N^{-2}$.
Also $\log(1/\delta)\geq16N^2$, and hence
\begin{align*}
 6^N\eta_{N,h,\delta}
 &=16C_+N^2e^{-(4-\log6)N}
       +12^N\delta e^{4N^2}\\
 &\leq16C_+N^2e^{-(4-\log6)N}
       +e^{N\log12-12N^2}=o(N^{-1}).
\end{align*}
Since $(1-\beta)^N=o(N^{-1})$, it follows that $E_N=O(N^{-1})$;
the gain bound $D_N$ is smaller.  Thus the profile error is $O(N^{-1/2})$.
Finally $N\asymp\sqrt{\log(1/\delta)}$, proving
\eqref{p08:eq:gate-value-log-rate}.
\end{proof}

\begin{corollary}[Stable hidden structure from ordinary gate values]
\label{p08:cor:gate-values-finite-structure}
Consider two models in the class
\eqref{p08:eq:jet-structure-class}--\eqref{p08:eq:jet-tilted-class}, with
$\alpha=1-v/(2u)>0$.  Suppose their ordinary gate values satisfy
\eqref{p08:eq:gate-value-data}.  In
Corollary~\ref{p08:cor:gate-finite-structure}, use
$\eta=\eta_{N,h,\delta}$ from \eqref{p08:eq:gate-value-jet-error} to form
$E_N$, $h_N$, and $\delta_N$.  Whenever
\eqref{p08:eq:jet-structure-threshold} holds, the dimensions agree and
\eqref{p08:eq:jet-power-resolution} bounds their powers.

In particular, under \eqref{p08:eq:gate-value-design}, all sufficiently small
$\delta>0$ force equality of the dimensions and give
\begin{equation}\label{p08:eq:gate-value-power-rate}
 \min_\pi\max_i|q_{1,i}-q_{2,\pi(i)}|
 =O\left((\log(1/\delta))^{-\alpha/4}\right).
\end{equation}
No ordinary moment assumption on the original laws is needed, and no
rational independence of their powers is required.
\end{corollary}

\begin{proof}
Theorem~\ref{p08:thm:finite-gate-values} supplies precisely the finite-jet
hypothesis of Corollary~\ref{p08:cor:gate-finite-structure}.
For the stated design, $E_N=O(N^{-1})$ and
$\delta_N=O(N^{-1/2})=O((\log(1/\delta))^{-1/4})$.
The strict threshold in that corollary is therefore satisfied eventually,
and its power bound is $O(\delta_N^\alpha)$, as claimed.
\end{proof}

If measured gate values have individual error at most $\delta_0$, any two
models fitting those same data to that tolerance satisfy
\eqref{p08:eq:gate-value-data} with $\delta=2\delta_0$.  The conclusions
therefore bound the diameter of the compatible model class from finitely
many ordinary noisy observations.  The proof does not require derivative
observations.  It does require the specified quantitative class bounds and
noise-dependent sampling design.  The powers may have rational relations,
so stable dimension and power recovery still does not imply recovery of
the full joint law.  The logarithmic rate is a sufficient bound for this
construction, not a minimax or computational-efficiency claim.

\section{Sharp noise cost of radial window design}
\label{p07:sec:window-noise}

The finite-window identity of
Proposition~\ref{p06:prop:window-inversion} is an exact-data inverse.
We now quantify how a radial window changes the conditioning of recovery
of the local Laplace exponent.  We specify the target and the noise norm
carefully: the result below concerns the function $h=-\log L$, not the
subsequent inversion of $L$ to a probability distribution.  It is a
deterministic noise theorem, not a statement about the sample size needed
to estimate a rare-event rate.

Fix $p>\rho>0$, $L>0$, $p_0\in(1/2,1)$, and $M>0$.  Let
\[
 \mathcal P=\{P\text{ on }[0,\infty):P\{0\}\ge p_0,
                       \ \int x\,P(dx)\le M\},
 \qquad h_P(s)=-\log\int e^{-sx}\,P(dx).
\]
Choose $0<\gamma<p$, with $\gamma\ne\rho$, and put
\begin{equation}\label{p07:eq:window-target}
 f_P(t)=e^{-\gamma t}h_P(e^{pt}),\qquad
 g_P(t)=e^{-\gamma t}\rho\int_0^L
                  e^{\rho u}h_P(e^{p(t-u)})\,du.
\end{equation}
Thus $g_P=Kf_P$, where
\begin{equation}\label{p07:eq:window-convolution}
 (Kf)(t)=\int k(u)f(t-u)\,du,
 \qquad k(u)=\rho e^{(\rho-\gamma)u}\mathbf1_{[0,L]}(u).
\end{equation}
All norms in this section are $L^2(\mathbb R)$ norms unless indicated
otherwise.  The observed function is $g_P+e$, where $\|e\|_2\le\delta$.
For an observation operator $T$, write
\begin{equation}\label{p07:eq:window-minimax}
 \mathcal E_T(\delta)=
 \inf_{\mathcal A}\sup_{P\in\mathcal P}
 \sup_{\|e\|_2\le\delta}
       \|\mathcal A(Tf_P+e)-f_P\|_2.
\end{equation}
For several channels the observation and noise belong to the Hilbert
direct sum of their $L^2$ spaces.

\begin{theorem}[Optimal radial-window noise rate]
\label{p07:thm:window-optimal}
There are constants $0<c<C<\infty$, depending only on the fixed
parameters above, such that, for all sufficiently small $\delta>0$,
\begin{equation}\label{p07:eq:window-rate}
 c\delta\log(1/\delta)\le\mathcal E_K(\delta)
                   \le C\delta\log(1/\delta).
\end{equation}
The lower bound continues to hold if the exact atom at zero and the
entire sequence of ordinary moments of $P$ are supplied to the
reconstruction procedure.  More precisely, there is one fixed atom and
one fixed moment sequence such that the lower bound already holds on
the subclass of $\mathcal P$ with those data.

Moreover, let $k_1,\ldots,k_J\in L^1(\mathbb R)\cap BV(\mathbb R)$,
extended by zero wherever a compactly supported window vanishes, and
suppose
\begin{equation}\label{p07:eq:window-budget}
 \left(\sum_{j=1}^J
       (\|k_j\|_1+\operatorname{TV}(k_j))^2\right)^{1/2}\le B.
\end{equation}
For $Tf=(k_1*f,\ldots,k_J*f)$, the lower bound
\begin{equation}\label{p07:eq:all-window-lower}
 \mathcal E_T(\delta)\ge c_B\delta\log(1/\delta)
\end{equation}
holds uniformly over all such designs.  The constants do not depend
on $J$ or on the chosen kernels.  Hence one tilted boxcar attains the
best order possible under this window-variation budget.
\end{theorem}

The lower bound allows a design to depend on $\delta$, provided its
budget remains bounded.  The statement does not include atomic kernels:
an ideal Dirac observation has no integration loss and is outside
$L^1\cap BV$.

\subsection{The analytic property of genuine Laplace exponents}

\begin{lemma}\label{p07:lem:window-analytic}
Put $b_*=\pi/(2p)$. There is $R<\infty$, depending only on
$p_0,M,p,\gamma$, such that every $f_P$ extends holomorphically to
$|\operatorname{Im}t|<b_*$ and satisfies
\begin{equation}\label{p07:eq:window-analytic-bound}
 \int_{\mathbb R}e^{2b_*|\xi|}|\widehat f_P(\xi)|^2\,d\xi
 \le 2R^2.
\end{equation}
Here the Fourier transform is unitary.  In particular this is a
uniform property of the probability class, not an additional Sobolev
or analyticity assumption on the unknown target.
\end{lemma}
\begin{proof}
For $\operatorname{Re}s\ge0$, the Laplace transform $L_P(s)$ belongs
to the disk with center $q=P\{0\}$ and radius $1-q$.
Consequently $\operatorname{Re}L_P(s)\ge2p_0-1>0$, and the principal
logarithm defines $h_P$ without zeros or branch ambiguity.  The
logarithmic series about $q$ gives
\[
 |h_P(s)|\le C_0:=-\log(2p_0-1).
\]
Integrating the derivative of the logarithm on the segment joining
$1$ to $L_P(s)$, and using
$|1-e^{-sx}|\le |s|x$, gives the second bound
\[
 |h_P(s)|\le C_1|s|,\qquad C_1=M/(2p_0-1).
\]
If $|v|\le b_*$, then $\operatorname{Re}e^{p(t+iv)}\ge0$.
It follows that
\[
 |f_P(t+iv)|\le e^{-\gamma t}
                         \min(C_0,C_1e^{pt}).
\]
The right side belongs to both $L^1$ and $L^2$, uniformly in $v$.
For example one may take
\[
 R^2=\frac{C_1^2}{2(p-\gamma)}+
                         \frac{C_0^2}{2\gamma}.
\]
For $|v|<b_*$, contour shifting, justified by the same exponential
bounds at the two ends of a horizontal strip, gives
$\widehat{f_P(\cdot+iv)}(\xi)=e^{-v\xi}\widehat f_P(\xi)$.
Apply Plancherel at $v$ and $-v$, let $v\uparrow b_*$ using
Fatou's lemma, and add the two boundary inequalities.
\end{proof}

\subsection{An estimator and its uniform error bound}

Set $a=\rho-\gamma\ne0$.  The Fourier multiplier of $K$ is
\begin{equation}\label{p07:eq:window-symbol}
 m(\xi)=\rho\frac{e^{(a-i\xi)L}-1}{a-i\xi}.
\end{equation}
In particular,
\begin{equation}\label{p07:eq:window-symbol-bounds}
 \frac{\rho|e^{aL}-1|}{\sqrt{a^2+\xi^2}}
 \le |m(\xi)|\le
 \frac{\rho(e^{aL}+1)}{\sqrt{a^2+\xi^2}}.
\end{equation}
The off-critical choice $\gamma\ne\rho$ removes the zeros of the
ordinary boxcar multiplier.  For observed $g^\delta$, define the
spectral-cutoff estimator by
\[
 \widehat f_T^\delta(\xi)=
 \mathbf1_{\{|\xi|\le T\}}\frac{\widehat g^\delta(\xi)}{m(\xi)}.
\]
Equations \eqref{p07:eq:window-analytic-bound} and
\eqref{p07:eq:window-symbol-bounds} imply the explicit estimate
\begin{equation}\label{p07:eq:window-cutoff-error}
 \|f_T^\delta-f_P\|_2\le
 \frac{\sqrt{a^2+T^2}}{\rho|e^{aL}-1|}\delta
                        +\sqrt2R e^{-b_*T}.
\end{equation}
Taking $T=b_*^{-1}\log(\sqrt2R/\delta)$ when $\delta$ is sufficiently
small proves the upper bound in Theorem~\ref{p07:thm:window-optimal}.
This estimator needs no numerical derivative of the noisy data.

\subsection{A lower bound inside an actual probability class}

The following construction uses the classical oscillatory
lognormal Stieltjes class \cite{Stoyanov2004,Kleiber2013}; its moment
cancellation is not a new
moment-indeterminacy result.  Its role here is to make the noise
lower bound hold for genuine local laws, even after every ordinary
moment has been specified.

Choose
\[
 0<\lambda<\min\{-\log p_0,Me^{-1/2}\},\qquad
 0<\varepsilon\le1/2,\qquad \omega_n=2\pi n.
\]
Let $P_{n,\pm,\varepsilon}$ be the compound-Poisson laws with Levy
measures obtained by pushing forward
\begin{equation}\label{p07:eq:window-stieltjes}
 \lambda\frac{e^{-u^2/2}}{\sqrt{2\pi}}
                    (1\pm\varepsilon\sin(\omega_n u))\,du
\end{equation}
under $u\mapsto e^u$.  These are positive finite measures.  For every
integer $j\ge0$,
\begin{equation}\label{p07:eq:window-moment-cancel}
 \int_{\mathbb R}e^{ju}\frac{e^{-u^2/2}}{\sqrt{2\pi}}
                     \sin(2\pi n u)\,du
 =e^{(j^2-4\pi^2n^2)/2}\sin(2\pi nj)=0.
\end{equation}
All measures in \eqref{p07:eq:window-stieltjes} therefore have mass
$\lambda$ and the same moments $\lambda e^{j^2/2}$.  Their
compound-Poisson laws have atom $e^{-\lambda}$ at zero and mean
$\lambda e^{1/2}$, and lie in $\mathcal P$.  They also have one common
full sequence of ordinary moments: the $j$th cumulant is
$\lambda e^{j^2/2}$, and each moment is the corresponding complete
Bell polynomial in the cumulants.  In particular these side data
do not distinguish any of the alternatives.

Write $\Delta f_{n,\varepsilon}=f_{P_{n,+,\varepsilon}}-
f_{P_{n,-,\varepsilon}}$, and set $b=\gamma/p\in(0,1)$.
The Mellin integral
\[
 \int_0^\infty x^{-z-1}(1-e^{-x})\,dx=-\Gamma(-z),
                      \qquad 0<\operatorname{Re}z<1,
\]
and Fubini give, with $z=b+i\xi/p$,
\begin{equation}\label{p07:eq:window-exact-packet}
 \widehat{\Delta f}_{n,\varepsilon}(\xi)=
 \frac{\lambda\varepsilon}{ip\sqrt{2\pi}}
 [-\Gamma(-z)]
 \left(e^{(z+i\omega_n)^2/2}-e^{(z-i\omega_n)^2/2}\right).
\end{equation}
Fubini is justified by the integrability of $x^{-b-1}(1-e^{-x})$
and the finite $b$th moment of the lognormal jump measure.

\begin{lemma}[Frequency-packet estimates]\label{p07:lem:window-packets}
Put $A_n=\omega_n^{-b-1/2}e^{-\pi\omega_n/2}$.
There are positive constants $c,C,C_r$, independent of $n$ and
$\varepsilon$, such that, for all sufficiently large $n$,
\begin{align}
 c\lambda\varepsilon A_n
 &\le \|\Delta f_{n,\varepsilon}\|_2
                   \le C\lambda\varepsilon A_n,
                    \label{p07:eq:window-packet-size}\\
 \int_{\mathbb R}(1+\xi^2)^{-r}
       |\widehat{\Delta f}_{n,\varepsilon}(\xi)|^2\,d\xi
 &\le C_r(\lambda\varepsilon)^2\omega_n^{-2r}A_n^2,
                  \qquad r\ge0.\label{p07:eq:window-packet-weight}
\end{align}
\end{lemma}
\begin{proof}
For $b\in(0,1)$, Stirling's estimate, together with continuity and
the absence of poles on the line $-b+i\mathbb R$, gives two-sided
bounds
\[
 |\Gamma(-b-iv)|\asymp
                   (1+|v|)^{-b-1/2}e^{-\pi|v|/2}.
\]
The two Gaussian factors in \eqref{p07:eq:window-exact-packet} have
absolute values
$\exp\{[b^2-(\xi/p\pm\omega_n)^2]/2\}$.
On $|\xi/p-\omega_n|\le1$, one Gaussian is bounded below by a
positive constant and the other is exponentially small in
$\omega_n^2$.  This proves the lower bound in
\eqref{p07:eq:window-packet-size}.

For the upper bounds use $|A-B|^2\le2|A|^2+2|B|^2$ and symmetry.
Writing $v=\xi/p$ reduces the relevant integral, up to constants,
to
\[
 \int_{\mathbb R}(1+|v|)^{-2b-1-2r}
          e^{-\pi|v|}e^{-(v-\omega_n)^2}\,dv.
\]
On $|v-\omega_n|\le\omega_n/2$, divide by
$\omega_n^{-2b-1-2r}e^{-\pi\omega_n}$ and use the integrable
majorant $C e^{-w^2-\pi w}$, $w=v-\omega_n$.
On the complement the Gaussian is at most exponentially small in
$\omega_n^2$, after its product with $e^{-\pi|v|}$ is integrated;
this is negligible compared with the displayed divisor.
These estimates establish \eqref{p07:eq:window-packet-size} and
\eqref{p07:eq:window-packet-weight}.
\end{proof}

For the general design in \eqref{p07:eq:window-budget}, integration by
parts against the finite signed derivative of a BV function yields
\[
 |\widehat k_j^{\,\mathrm{mult}}(\xi)|
 \le\min\{\|k_j\|_1,\operatorname{TV}(k_j)/|\xi|\},
\]
where $\widehat k_j^{\,\mathrm{mult}}(\xi)=\int k_j(u)e^{-i\xi u}\,du$
is the convolution multiplier.  Hence
\[
 \sum_j|\widehat k_j^{\,\mathrm{mult}}(\xi)|^2
                         \le \frac{2B^2}{1+\xi^2}.
\]
Lemma~\ref{p07:lem:window-packets} implies
\begin{equation}\label{p07:eq:window-two-point}
 \|T\Delta f_{n,\varepsilon}\|_2
       \le C_B\lambda\varepsilon A_n/\omega_n,
 \qquad
 \|\Delta f_{n,\varepsilon}\|_2
       \ge c\lambda\varepsilon A_n.
\end{equation}
Fix $\varepsilon_0=1/2$. For small $\delta$, take the largest $n$
for which $C_B\lambda\varepsilon_0 A_n/\omega_n\ge\delta$, and
then choose $\varepsilon\le\varepsilon_0$ so the right side of the
first inequality in \eqref{p07:eq:window-two-point} is $\delta$.
Because $\omega_n=2\pi n$,
\[
 n\sim\pi^{-2}\log(1/\delta),\qquad
 \|\Delta f_{n,\varepsilon}\|_2
                          \ge c'_B\delta\log(1/\delta).
\]
The midpoint of the two noiseless observations can be explained
by either law with noise norm at most $\delta/2$.
Every reconstruction makes an error of at least half their target
separation for one of the laws. This proves both lower bounds in
Theorem~\ref{p07:thm:window-optimal}, including the assertion about
full moment side information.

\begin{corollary}[Cost of smoothing the edges of a window]
\label{p07:cor:window-spline}
For any fixed integer $r\ge1$, observe $K^rf_P$ instead of $Kf_P$.
Then
\[
 \mathcal E_{K^r}(\delta)\asymp
                         \delta[\log(1/\delta)]^r.
\]
The log-kernel of $K^r$ is
$\rho^r e^{(\rho-\gamma)u}\mathbf1_{[0,L]}^{*r}(u)$,
so this corresponds to a nonnegative B-spline radial mask,
supported on $[0,rL]$ in logarithmic radius.  Multiplication by a
fixed normalizing constant does not alter the asserted order.
\end{corollary}
\begin{proof}
The multiplier is $m(\xi)^r$ and thus comparable to
$(1+|\xi|)^{-r}$. The cutoff argument gives error
$O(\delta T^r+e^{-b_*T})$. For the lower bound use
\eqref{p07:eq:window-packet-weight} with exponent $r$ and repeat the
two-point argument with $A_n/\omega_n^r$.
\end{proof}

\begin{proposition}[No uniform total-variation inversion]
\label{p07:prop:window-tv-barrier}
On the same class $\mathcal P$, the local distribution cannot be
recovered uniformly consistently in total variation from the noisy
window data of Theorem~\ref{p07:thm:window-optimal}. This remains true if
the exact atom at zero and all ordinary moments are known. The claim
also holds for any design satisfying a fixed budget
\eqref{p07:eq:window-budget}.
\end{proposition}
\begin{proof}
Keep $\varepsilon\in(0,1/2]$ fixed in
\eqref{p07:eq:window-stieltjes}. Write $\nu_\pm$ for the two finite
Levy measures, each of mass $\lambda$, and write $\|\cdot\|_{\rm var}$
for the total variation norm of a signed measure. The compound-Poisson
expansion and the telescoping convolution identity imply
\[
 P_+-P_-=e^{-\lambda}(\nu_+-\nu_-)+R,
 \qquad
 \|\nu_+^{*j}-\nu_-^{*j}\|_{\rm var}
       \le j\lambda^{j-1}\|\nu_+-\nu_-\|_{\rm var}.
\]
Consequently
\[
 \|R\|_{\rm var}
 \le e^{-\lambda}\sum_{j\ge2}
       \frac{j\lambda^{j-1}}{j!}\|\nu_+-\nu_-\|_{\rm var}
 =(1-e^{-\lambda})\|\nu_+-\nu_-\|_{\rm var}.
\]
Since $e^{-\lambda}>p_0>1/2$, this yields
\[
 d_{\rm TV}(P_+,P_-)
 \ge (2e^{-\lambda}-1)\lambda\varepsilon
                  \mathbb E|\sin(2\pi nU)|,
 \qquad U\sim N(0,1),
\]
where $d_{\rm TV}(P,Q)=\frac12\|P-Q\|_{\rm var}$.
The last expectation tends to $2/\pi$, by the Fourier series of
$|\sin|$ and the Gaussian characteristic function. Thus the
distributional separation stays bounded away from zero, while
\eqref{p07:eq:window-two-point} makes the observation separation tend
to zero. The midpoint argument proves the asserted impossibility.
The common moment sequence and common atom were already verified
in \eqref{p07:eq:window-moment-cancel}.
\end{proof}

\begin{remark}[Scope and relation to classical inverse problems]
Mellin convolution, spectral cutoff, ordinary-smooth
deconvolution, and lognormal moment-indeterminacy are classical;
see, in particular, the boxcar channel-design theory of
\cite{PenskySapatinas2011} and the Stieltjes classes in
\cite{Stoyanov2004,Kleiber2013}.
The role of the present result is to connect an implementable
radial design to a sharp noise modulus on the actual class of
local log-Laplace exponents.  The uniform BV-budget lower bound
also rules out improvement by merely adding finitely many
ordinary radial windows.  Neither the moment construction nor
the Fourier cutoff method is claimed as a new general technique.
The sharp lower bounds here use unbounded compound-Poisson laws;
they do not assert the same lower bound under a fixed compact
support restriction.  Recovery of the distribution from its
real Laplace transform remains a separate inverse problem.
\end{remark}

\section{The dimension cost of adaptive linear readings}
\label{p08:sec:information-barriers}

An adaptive linear experiment with $m$ readings sequentially chooses bounded
Borel functions $\varphi_1,\ldots,\varphi_m$ on $Q=[0,1]^d$, with each
choice allowed to depend arbitrarily on the previous exact answers, and
receives $\int_Q\varphi_i\,d\mu$.  Neither continuity of the adaptive rule
nor a bound on the magnitude of the query functions is imposed.  Laplace
values, compactly supported moment data, and Laplace derivatives are examples.
The result below is a direct information-width argument, included to prevent
an interpretation of curve uniqueness as finite-data compression.  Its
method is part of the established theory of adaptive information; see
\cite{KriegNovakUllrich2025}.

\begin{theorem}[Sharp order for finite exact linear information]
\label{p08:thm:adaptive-linear-w1}
For every fixed $d\geq1$, there is a constant $b_d>0$ with the following
property.  For any deterministic adaptive experiment using $m\geq1$ linear
readings, there are two smooth probability densities $f_+,f_-$ on $Q$, with
$1/2\leq f_\pm\leq3/2$, producing exactly the same complete transcript and
satisfying
\begin{equation}\label{p08:eq:adaptive-w1-lower}
 W_1(f_+\,dx,f_-\,dx)\geq b_d m^{-1/d}.
\end{equation}
Consequently its worst-case reconstruction error is at least
$\tfrac12b_dm^{-1/d}$.  If $q$ additional prescribed bounded linear
calibrations are supplied, the same assertion holds with $m$ replaced by
$m+q$.  Conversely, arbitrary linear information admits a nonadaptive
reconstruction of every probability law on $Q$ with $m$ readings and error
at most $B_d m^{-1/d}$, for a constant $B_d$ depending only on $d$.
\end{theorem}

\begin{proof}
Choose smooth functions $\psi,\phi$ supported in the interior of $Q$ such
that
\[
 \int_Q\psi=0,\qquad \|\psi\|_\infty\leq1,\qquad
 \operatorname{Lip}(\phi)\leq1,\qquad
 A:=\int_Q\phi\psi>0.
\]
For example, take a nonnegative smooth interior bump $\eta$, let $\psi$ be
a positive normalization of $-\partial_1\eta$, and let $\phi$ agree with
a sufficiently small positive multiple of $x_1$ on its support and vanish
near the boundary of $Q$.

Let $k=\lceil(2m)^{1/d}\rceil$, $N=k^d$, and $h=1/k$.  Subdivide $Q$
into $N$ cubes of side length $h$, with lower corners $z_j$, and set
\[
 \psi_j(x)=\psi((x-z_j)/h),\qquad
 \phi_j(x)=h\phi((x-z_j)/h)
\]
inside the corresponding cube and zero outside.  Run the proposed algorithm
on the uniform density.  This fixes a particular list of query functions
$\varphi_1,\ldots,\varphi_m$.  Let $V\subset\mathbb R^N$ be the nullspace
of the matrix with entries $\int_Q\varphi_i\psi_j$.  Then
$\dim V\geq N-m$.

At an extreme point $v$ of the nontrivial compact polytope
$V\cap[-1,1]^N$, at least $\dim V$ coordinates have absolute value one.
Indeed, if fewer coordinate constraints are active, a nonzero direction in
$V$ annihilating the active coordinates allows small displacements of both
signs within the polytope, contradicting extremality.  Thus we can choose
$v$ with $|v_j|\leq1$ and $\sum_j|v_j|\geq N-m\geq N/2$.

The densities
\[
 f_\pm=1\pm\tfrac12\sum_{j=1}^N v_j\psi_j
\]
are smooth probability densities in $[1/2,3/2]$.  They give the same answer
as the uniform density to every query in the fixed list.  Induction over
the query number therefore shows that the adaptive algorithm follows this
same list and produces exactly the same transcript for both inputs.

The function $\Phi_v=\sum_j\operatorname{sgn}(v_j)\phi_j$ is
$1$-Lipschitz on $Q$: on each cube it is $1$-Lipschitz and vanishes near the
boundary, and the same bound follows across cubes by subdividing a straight
line at its intersections with their boundaries.  Hence
\begin{align*}
 W_1(f_+\,dx,f_-\,dx)
 &\geq\int_Q\Phi_v(f_+-f_-)\,dx
 =A h^{d+1}\sum_j|v_j|\\
 &\geq\frac{A}{2}h\geq\frac{A}{4\,2^{1/d}}m^{-1/d}.
\end{align*}
This proves the lower bound.  Appending the $q$ calibration functions to
the nullspace constraints gives the calibrated version, by the same proof.
The densities then share the calibrations of the uniform law; no assertion
is made for every possible prescribed calibration value.

For the upper bound, choose
$k=\lfloor(m+1)^{1/d}\rfloor\geq1$ and query the masses of all but one
of the $k^d$ grid cubes.  The remaining mass follows from total mass one.
Move each cube's mass to its center.  This costs at most
$\sqrt d/(2k)\leq\sqrt d\,m^{-1/d}$ in $W_1$.
\end{proof}

The upper bound permits arbitrary bounded linear query functions, including
indicators of cubes.  It does not claim a matching upper bound for Laplace
queries alone.  The theorem is deterministic; no randomized complexity
claim is made.  Its lower bound also permits derivatives of unbounded order
to be queried exactly, provided each derivative value counts as a separate
linear reading.  Smoothness here is qualitative: no uniform derivative
bound on the densities is imposed.

\section{Quantitative questions for controlled observations}
\label{p07:sec:questions}The finite-value guarantees resolve one part of the quantitative problem
left by exact blind identification. The sharp modulus from ordinary gate
values, the best placement of finitely many settings, and recovery on a
fixed interior gate interval remain open in this framework. In particular,
the bound in Corollary~\ref{p08:cor:gate-values-finite-structure} is a
sufficient rate, and the linear-information lower bound does not establish
its optimality for the nonlinear gate functional.

Recovering the entire joint law requires additional control of tails,
regularity, or the topology of error, as well as quantitative separation
of higher mixed-moment exponents. Pairwise separation of the original
powers only addresses the first exponent band. A classification of the
models representing a fixed resonant gate curve would further distinguish
power collisions from vanishing mixed moments. The smooth null densities
above rule out an unrestricted universal positive statement.

Empirical estimation of the rare-event response and propagation of that
error through the inverse form a separate sampling problem. The
deterministic window rate is not a sample-complexity result. In countable
dimension, a quantitative bound on $\sum_{j>m}X_j$ may support an effective
dimension notion, but a detection threshold must be specified because
arbitrarily small added coordinates destroy uniform cardinality recovery.

The finite-window clock is known throughout the interior-probe theorem.
An unknown clock changes the parameterization while preserving the
weighted distribution of profile values, so additional clock calibration
would be needed. Likewise, the proof with countably many powers uses
$\inf_iq_i>0$ to choose a common convergence radius for the signed
exponent measure; exponents tending to zero require a different argument.

\clearpage

% ===== Integrated Part 5 =====
\part{Infinite mixtures, moment coincidences, and statistical recovery}
\label{part:mixtures}

The preceding parts concern what arithmetic and analytic observations
reveal about hidden structure. Here the unknown object is a probability
measure. We study three questions separately: whether an exact collection
of observed laws determines that measure; whether replication removes
otherwise invisible directions; and whether a finite sample permits
uniform recovery. These are different properties, even for strictly
positive observation kernels that extend to entire functions.

For a known component family $P_\theta$ and an unknown mixing
probability $\mu$, the law of $q$ conditionally independent observations
sharing one latent parameter is
\[
 A_q\mu=\int P_\theta^{\otimes q}\,d\mu(\theta).
\]
The replication threshold is the least positive integer $q$ for which
$A_q$ is injective on all Borel mixing probabilities, with threshold
$\infty$ if no such integer exists. This counts observations within
one group, not the number $N$ of independent groups available for
estimation. Separately redrawing the latent parameter for every
observation produces the product of the one-observation law and
does not repair its ambiguities.

The results develop in three stages. First, a bounded-holomorphic
Fourier--Laplace transform gives the sharp reciprocal-sum criterion
for Gaussian mixtures observed through partial-sum marginals, and
classifies fixed-shape stable families. Second, a missing Legendre
mode separates finite parameter-derivative tests from global mixture
identification. Finite-dimensional ambiguity becomes atomic after two
observations, whereas infinite-dimensional ambiguity supports arbitrary
replication thresholds on both the circle and an ordinary interval.
Finally, multiplicative smoothing gives localized Gevrey annihilators
and exact moment-coincidence sets; their interaction with compactness
distinguishes analytic identification from uniform statistical recovery.
The factorial Legendre example then supplies fixed-replication rates,
growing-replication budget bounds, and a contrasting construction with
arbitrarily slow recovery.

\paragraph{Literature and scope.}
Lin's normal-mixture question is used in the form stated in
\cite[Section~2.6]{Lin2015}. The differential-barrier question
concerns precisely the finite operator classes in
\cite[Definitions~1 and~4]{LeBarilettoRinaldoHo2026}; no absence of
nonlocal or infinite-order obstructions is asserted.
Vandermeulen and Scott~\cite{VandermeulenScott2015} obtain the sharp
$2m-1$ bound within the corresponding class of mixtures with $m$
components, and Vandermeulen and Saitenmacher~\cite{VandermeulenSaitenmacher2022}
refine finite-component bounds using independence assumptions.
Here a fixed component family is allowed arbitrary Borel mixing laws,
so finite-support identification and unrestricted identification
must be distinguished.

The analytic ingredients are classical. Fourier descriptions of
characteristic kernels appear in Fukumizu et
al.~\cite{Fukumizu2008}; arithmetic-progression criteria for strict
positive definiteness and Skolem--Mahler--Lech arguments appear in
Pinkus~\cite[Section~3]{Pinkus2004} and
Bell--Lagarias~\cite[Theorem~1.3]{BellLagarias2015}.
Identification by the full tensor-moment hierarchy is an existing
principle~\cite{Bonnier2023}. The M\"untz and Clarkson--Erd\H{o}s
theorems are used in the forms of Borwein and
Erd\'elyi~\cite[Theorems~3.1--3.2]{BorweinErdelyi1995}; sparse
moments have also been used in mixture identification by Stoyanov and
Lin~\cite{StoyanovLin2011}.
The localized smooth-annihilator theorem below proves the mathematical
assertion in B\"urgin--Epperlein--Wirth's
Conjecture~17~\cite{BurginEpperleinWirth2022}, with stronger
smoothness, support, and exact-coincidence conclusions.
These source-specific claims do not assert historical priority or
external acceptance. In particular, the rate
$\log\log N/\log N$ is already known in other mixture models,
including Poisson mixtures~\cite{MiaoEtAl2024}.



\section{Normal mixtures from partial-sum marginals}
\label{p09:sec:normal-resolution}

\subsection{The question being answered}

Lin's preprint \cite[Section~2.6]{Lin2015} asks whether arbitrary
mixtures of iid normal sequences are determined by the marginal laws
of their partial sums. Its Proposition~2.6.4 proves this under an
exponential-moment assumption on the random variance. We remove that
assumption. The observation remains exactly the one in the question:
the individual law of each partial sum, with no joint observations of
different partial sums.

Let $(A,V)$ have an arbitrary Borel probability law on
$\mathbb R\times[0,\infty)$, and let $Z_1,Z_2,\ldots$ be independent
standard normal random variables independent of $(A,V)$. Define
\begin{equation}\label{p09:eq:normal-mixture-model}
       X_j=A+\sqrt V Z_j,\qquad S_n=\sum_{j=1}^nX_j.
\end{equation}
The parameters $A$ and $V$ may be dependent, and neither is subject to
any moment condition. Every mixture of iid normal sequences, including
degenerate normal components, has this representation in distribution.

\begin{theorem}[Normal mixtures from partial-sum marginals]
\label{p09:thm:normal-no-moments}
The collection $(\mathcal L(S_n))_{n\geq1}$ determines the entire
joint law of $(A,V)$ in \eqref{p09:eq:normal-mixture-model}, and hence the
law of the exchangeable sequence $(X_j)_{j\geq1}$.

More generally, it is enough to know the marginal laws of $S_n$ for
$n$ in any infinite set $\mathcal N\subset\mathbb N$ satisfying
\begin{equation}\label{p09:eq:normal-sparse-condition}
                  \sum_{n\in\mathcal N}\frac1n=\infty.
\end{equation}
\end{theorem}

The key is not analytic continuation across the boundary point zero.
The relevant transform is bounded and holomorphic on the open right
half-plane. Its prescribed zeros approach zero too densely to be the
zeros of a nonzero bounded holomorphic function.

\subsection{Two transform lemmas}

We include the required half-plane uniqueness argument. It is the
classical necessary condition on the zeros of a bounded holomorphic
function, often expressed through Blaschke products.

\begin{lemma}[A boundary uniqueness criterion]
\label{p09:lem:half-plane-zeros}
Let $H$ be a bounded holomorphic function on $\{z:\Re z>0\}$.
Suppose $r_k>0$ are distinct, $H(r_k)=0$, and
\begin{equation}\label{p09:eq:half-plane-divergence}
                   \sum_k\frac{r_k}{1+r_k^2}=\infty.
\end{equation}
Then $H$ is identically zero.
\end{lemma}
\begin{proof}
Put $h(w)=H((1+w)/(1-w))$ for $|w|<1$. If $h$ is not identically
zero, factor out its finite-order zero at zero, if present, to obtain
a bounded holomorphic function $g$ with $g(0)\ne0$. Boundedness of
the quotient follows from the maximum principle on disks and passage
of their radii to one. Jensen's formula implies that its nonzero
zeros $w_j$, counted with multiplicity, satisfy
\[
 \sum_{|w_j|<r}\log\frac r{|w_j|}
 \leq\log\|g\|_\infty-\log|g(0)|
 \qquad(0<r<1),
\]
with radii meeting zeros omitted if needed. Monotone convergence as
$r\uparrow1$ gives $\sum_j(1-|w_j|)<\infty$, since
$1-u\leq-\log u$ for $0<u\leq1$.

Our zeros correspond to $w_k=(r_k-1)/(r_k+1)$. Omitting the possible
single zero $r_k=1$, we have
\[
 1-|w_k|
 =\frac{1-|w_k|^2}{1+|w_k|}
 \geq\frac{2r_k}{(1+r_k)^2}
 \geq\frac{r_k}{1+r_k^2}.
\]
This contradicts \eqref{p09:eq:half-plane-divergence}.
\end{proof}

\begin{lemma}[Fourier--Laplace uniqueness]
\label{p09:lem:fourier-laplace}
A finite complex Borel measure $\rho$ on
$\mathbb R\times[0,\infty)$ is zero if
\begin{equation}\label{p09:eq:FL-zero}
       \int e^{ita-sb}\rho(da,db)=0
              \qquad(t\in\mathbb R,\ s>0).
\end{equation}
\end{lemma}
\begin{proof}
For a fixed $s>0$, the finite measure
$D\mapsto\int_{D\times[0,\infty)}e^{-sb}\rho(da,db)$ on
$\mathbb R$ has zero Fourier transform and is therefore zero.
Consequently, for every Borel $D\subset\mathbb R$, the finite measure
$\rho_D(E)=\rho(D\times E)$ on $[0,\infty)$ has zero Laplace
transform at every $s>0$. Dominated convergence gives its vanishing
total mass as well. Push it forward by $b\mapsto e^{-b}$ and regard
the resulting measure as a measure on $[0,1]$. All its nonnegative
integer moments vanish. Density of polynomials in $C([0,1])$ and
finite-measure uniqueness show that $\rho_D=0$. Thus $\rho$ vanishes
on every Borel rectangle and hence is zero.
\end{proof}

\begin{proof}[Proof of Theorem~\ref{p09:thm:normal-no-moments}]
Let $\mu_1,\mu_2$ be two candidate laws for $(A,V)$ that yield the
same stated observations, and set $\rho=\mu_1-\mu_2$.
Conditionally on $(A,V)$, the average $S_n/n$ is normal with mean
$A$ and variance $V/n$. Equality of its characteristic functions
therefore gives
\begin{equation}\label{p09:eq:normal-data-zero}
       \int e^{ita-t^2v/(2n)}\rho(da,dv)=0
                   \qquad(t\in\mathbb R,\ n\in\mathcal N).
\end{equation}
Fix $t\ne0$ and define
\[
                H_t(z)=\int e^{ita-zv}\rho(da,dv),\qquad\Re z>0.
\]
Its absolute value is at most $|\rho|(\mathbb R\times[0,\infty))$.
It is holomorphic: on $\Re z\geq\varepsilon>0$, differentiation
under the integral is justified by boundedness of
$v^m e^{-\varepsilon v}$ for every fixed $m$.
Its distinct positive zeros include $r_n=t^2/(2n)$ for
$n\in\mathcal N$. Condition \eqref{p09:eq:normal-sparse-condition}
implies
\[
           \sum_{n\in\mathcal N}\frac{r_n}{1+r_n^2}=\infty.
\]
Lemma~\ref{p09:lem:half-plane-zeros} yields $H_t(s)=0$ for every $s>0$.
Continuity in $t$ also gives this equality at $t=0$.
Lemma~\ref{p09:lem:fourier-laplace} now gives $\rho=0$.
The common mixing law determines every finite-dimensional distribution
of the sequence by conditional independence, completing the proof.
\end{proof}

\begin{remark}[What changes in the earlier argument]
In the cited preprint the variance exponential moment lets the
transform extend across zero, after which the usual identity theorem
applies to $t^2/(2n)\to0$. The present argument instead uses
bounded holomorphy on the existing half-plane. No extension across zero
and no variance exponential moment are needed. This also explains why
an arbitrary sequence of boundary-approaching observations cannot
simply be substituted: the divergence condition is the step that
forces uniqueness.
\end{remark}





\section{Sharp sparse identifiability for Gaussian location--scale mixtures}

Let $\mathcal N\subset\{1,2,\ldots\}$ and, conditionally on a random
pair $(A,B)\in\mathbb R\times[0,\infty)$, let
$X_j=A+\sqrt B Z_j$, where the $Z_j$ are independent standard normal
variables, independent of $(A,B)$.  Write $S_n=X_1+\cdots+X_n$.
The preceding theorem proves the positive direction. We now establish
the converse with one joint density realizing all the required transform
zeros simultaneously.

\begin{theorem}[Sharp sampling criterion]
\label{p09:thm:sharp-gaussian}
The map $\mu\mapsto(\mathcal L_\mu(S_n))_{n\in\mathcal N}$ is injective
on the Borel probability measures on
$\mathbb R\times[0,\infty)$ if and only if
\[
   \sum_{n\in\mathcal N}\frac1n=\infty.
\]
If the sum converges, there are two distinct, absolutely continuous
mixing laws with all polynomial moments finite and with identical
selected sum laws.  In particular the obstruction persists even if
$B>0$ almost surely.
\end{theorem}

\begin{proof}[Sufficiency]
This is Theorem~\ref{p09:thm:normal-no-moments}, whose proof applies to
every set $\mathcal N$ with divergent reciprocal sum and imposes no
moment condition.
\end{proof}

\begin{proof}[Necessity, including all regularity details]
Put $\lambda_n=(2n)^{-1}$ and suppose
$\Lambda:=\sum_{n\in\mathcal N}\lambda_n<\infty$.  The finite or
infinite product
\[
 P(z)=\prod_{n\in\mathcal N}\frac{z-\lambda_n}{z+\lambda_n}
       \qquad(\Re z>0)
\]
converges locally uniformly, is not identically zero, has zeros at all
$\lambda_n$, and satisfies $|P(z)|\le1$.  These assertions follow from
$\sum_n\lambda_n<\infty$ and
$(z-\lambda)/(z+\lambda)=1-2\lambda/(z+\lambda)$; no additional phase
normalization is needed.  Its boundary values satisfy $|P(i\xi)|=1$
for $\xi\ne0$.

Use the principal square root and define
\[
 F(z)=P(z)\exp\{-\sqrt z-1/\sqrt z\},\qquad \Re z>0.
\]
This is a bounded, nonzero holomorphic function with $F(\lambda_n)=0$
and $F(\overline z)=\overline{F(z)}$.  We verify that its boundary
function $H(\xi)=F(i\xi)$, extended by $H(0)=0$, is Schwartz.
On the imaginary axis away from zero, logarithmic differentiation gives
\[
 Q(z):=\frac{P'(z)}{P(z)}
 =\sum_n\left(\frac1{z-\lambda_n}-\frac1{z+\lambda_n}\right).
\]
For each integer $r\ge0$,
\[
 |Q^{(r)}(i\xi)|\le
       2(r+1)!\Lambda |\xi|^{-r-2}.
\]
Indeed, express the difference of the two $(r+1)$st inverse powers
as an integral over the real interval $[-\lambda_n,\lambda_n]$;
throughout that interval $|i\xi-x|\ge|\xi|$.  The usual recursion
$P^{(j+1)}/P=(P^{(j)}/P)'+Q(P^{(j)}/P)$ therefore gives, for each
$j\ge1$,
\[
 |P^{(j)}(i\xi)|\le C_j|\xi|^{-2j}\quad(0<|\xi|\le1),
 \qquad
 |P^{(j)}(i\xi)|\le C_j\quad(|\xi|\ge1).
\]
Here the product and all derivatives converge uniformly on compact
subsets of the imaginary axis avoiding zero.  Since
\[
 \Re\sqrt{i\xi}=\sqrt{|\xi|/2},\qquad
 \Re\frac1{\sqrt{i\xi}}=\frac1{\sqrt{2|\xi|}},
\]
the two exponential factors dominate every polynomial derivative
bound both at zero and at infinity.  Thus $H\in\mathcal S(\mathbb R)$,
with all derivatives zero at zero.

Set
\[
 f(b)=\frac1{2\pi}\int_{\mathbb R}e^{i\xi b}H(\xi)\,d\xi.
\]
Then $f\in\mathcal S(\mathbb R)$, it is real, and it is nonzero.
An elementary contour shift proves that $f(b)=0$ for $b<0$.
For fixed $R>0$ shift the integration line from $i\mathbb R$ to
$R+i\mathbb R$ in the integral of $e^{bz}F(z)$.  The horizontal
integrals tend to zero, since for $0\le x\le R$,
$\Re\sqrt{x+iT}\ge\sqrt{|T|/2}$.  The boundary value at zero causes
no difficulty because $F$ extends continuously there by zero.
It follows that
\[
 |f(b)|\le \frac{e^{bR}}{2\pi}
           \int_{\mathbb R}|F(R+i\xi)|\,d\xi
 \le \frac{e^{bR}}{2\pi}
           \int_{\mathbb R}e^{-\sqrt{|\xi|/2}}\,d\xi.
\]
The last integral is independent of $R$, and $R\to\infty$ proves the
support assertion for $b<0$.  More formally one may first shift from
$\varepsilon+i\mathbb R$ and let $\varepsilon\downarrow0$ by dominated
convergence, using the same integrable bound.

For $s>0$, Fubini and the Cauchy integral formula on a right half-disk
give
\[
 \int_0^\infty e^{-sb}f(b)\,db
 =\frac1{2\pi}\int_{\mathbb R}\frac{H(\xi)}{s-i\xi}\,d\xi
 =F(s).
\]
The large semicircle contributes zero because
$|F(z)|\le e^{-\sqrt{|z|/2}}$ there.  A small semicircle at zero can
also be deleted, or treated by its continuous zero boundary value.

Choose any nonzero, real, even function
$h\in C_c^\infty(\{t:1<|t|<2\})$ and define
\[
 g(a,b)=\frac1{2\pi}\int_{\mathbb R}
       e^{-ita}h(t)t^2f(bt^2)\,dt.
\]
This is real.  The function $h(t)t^2f(bt^2)$ is smooth with compact
$t$-support bounded away from zero; all of its $t$-derivatives decay
faster than every power of $|b|$, uniformly in $t$.  Repeated
integration by parts in $t$ proves, for all $M,N\ge0$,
\[
 |g(a,b)|\le C_{M,N}(1+|a|)^{-N}(1+|b|)^{-M}.
\]
The same estimates hold for all $(a,b)$-derivatives.  Hence
$g\in\mathcal S(\mathbb R^2)$, and $g=0$ when $b<0$.  Fourier
inversion in $a$ and the Laplace identity for $f$ give
\[
 \int_{\mathbb R}\int_0^\infty e^{ita-sb}g(a,b)\,db\,da
      =h(t)F(s/t^2) \quad(t\ne0,\ s>0).
\]
The left side is zero when $t=0$, because the partial Fourier
transform in $a$ is zero there.  In particular $\int g=0$.
It is nonzero because $h$ and $f$ are nonzero.  Consequently, if
$g=g_+-g_-$ and $m=\int g_+=\int g_->0$, the two measures
\[
 \mu_+(da,db)=m^{-1}g_+(a,b)\,da\,db,
 \qquad
 \mu_-(da,db)=m^{-1}g_-(a,b)\,da\,db
\]
are distinct probabilities on $\mathbb R\times[0,\infty)$, with all
polynomial moments finite.  They give no mass to $b=0$.

Finally, for $n\in\mathcal N$ and a characteristic-function argument
$v\ne0$, set $t=nv$ and $s=nv^2/2$.  Then $s/t^2=1/(2n)=\lambda_n$,
so the difference of the two characteristic functions of $S_n$ is
\[
 m^{-1}h(nv)F(\lambda_n)=0.
\]
It is also zero at $v=0$.  Thus all selected sum laws agree, completing
the necessity proof.
\end{proof}

\begin{corollary}[Equal marginals and all moments do not repair sparse sampling]
\label{p09:cor:same-moments}
If $\sum_{n\in\mathcal N}1/n<\infty$, the two counterexample laws can
be chosen to have the same $A$ marginal, the same $B$ marginal, and
identical mixed polynomial moments of every order, all finite.
They may additionally satisfy $B\ge b_0$ almost surely for any
prescribed $b_0>0$.
\end{corollary}
\begin{proof}
Use the signed Schwartz density $g$ constructed above. Its partial
Fourier transform is
\[
 G(t,b):=\int_{\mathbb R}e^{ita}g(a,b)\,da=h(t)t^2f(bt^2).
\]
It vanishes on a neighborhood of $t=0$. Consequently, for each $b$,
\[
 \int_{\mathbb R}a^j g(a,b)\,da=i^{-j}\partial_t^jG(0,b)=0
 \qquad(j\ge0).
\]
The case $j=0$ gives a zero signed $B$ marginal. Multiplying by
$b^k$ and integrating gives zero for every mixed moment, with all
interchanges justified by Schwartz decay.

Also $\int_0^\infty f(b)\,db=\lim_{s\downarrow0}F(s)=0$, since
$|P(s)|\le1$ and the regularizing factor tends to zero.
The Fourier transform of the signed $A$ marginal is therefore
\[
 \int_0^\infty G(t,b)\,db
 =h(t)\int_0^\infty f(u)\,du=0
\]
for $t\ne0$, and is zero at $t=0$ as well. Fourier uniqueness proves
that this marginal vanishes. The normalized positive and negative
parts of $g$ thus have the asserted equal marginals and moments.

Finally replace $B$ by $B+b_0$ in both laws. This multiplies their
signed Fourier--Laplace transform by $e^{-b_0s}$, preserving every
required zero. It preserves equality of the coordinate marginals
and, by the binomial formula, equality of all mixed moments.
\end{proof}

\begin{remark}
All polynomial moments in the counterexamples are finite, but no
finite exponential moment is asserted. Thus agreement of all mixed
moments here does not contradict any moment-determinacy theorem.
\end{remark}





\subsection{Sampling examples and random drift--variance mixtures}

For normal mixtures, Theorem~\ref{p09:thm:sharp-gaussian} gives a complete
classification of the sets of observed integer lengths:
\begin{center}
\begin{tabular}{@{}lll@{}}
\toprule
Observed lengths & Reciprocal sum & Every mixing law identified?\\
\midrule
All positive integers & Divergent & Yes\\
An infinite arithmetic progression & Divergent & Yes\\
Squares $1,4,9,\ldots$ & Convergent & No\\
Powers of two $1,2,4,\ldots$ & Convergent & No\\
Any finite set & Finite & No\\
\bottomrule
\end{tabular}
\end{center}
A negative entry means that distinct mixing laws can give identical
observations. It does not mean that every mixing law has a competitor.
By Corollary~\ref{p09:cor:same-moments}, neither knowing both parameter
marginals nor knowing every mixed polynomial moment removes these
counterexamples. A common strictly positive variance lower bound
also does not remove them.

The Gaussian theorem applies as well to mixtures of Brownian motions:
if $Y_t=At+\sqrt V W_t$, with $W$ independent of $(A,V)$, then
$Y_n$ and $S_n$ have the same law. Thus integer-time marginal laws
identify the random drift and variance under exactly the same
sampling condition. No joint observation at different times is needed.


\section{A classification of fixed-shape stable families}
\label{p09:sec:stable-classification}

The same boundary uniqueness mechanism extends beyond the Gaussian family.
In this section the stability index and skewness are fixed; only the location
and scale are mixed.  The distinction is essential: the theorem does not
permit the shape parameters themselves to vary inside the mixing measure.

Fix $0<\alpha\leq 2$ and $-1\leq\beta\leq1$.  For $t\ne0$ put
\begin{equation}\label{p09:eq:stable-exponent}
 \psi_{\alpha,\beta}(t)=
 \begin{cases}
 |t|^\alpha\bigl(1-i\beta\tan(\pi\alpha/2)\operatorname{sgn}t\bigr),
       &\alpha\ne1,\\
 |t|\bigl(1+i(2\beta/\pi)\operatorname{sgn}t\log|t|\bigr),
       &\alpha=1,
 \end{cases}
\end{equation}
and put $\psi_{\alpha,\beta}(0)=0$.
Let $K_{a,b}^{\alpha,\beta}$ be the probability measure with characteristic
function
\begin{equation}\label{p09:eq:stable-kernel}
 \widehat K_{a,b}^{\alpha,\beta}(t)
       =\exp\{ita-b\psi_{\alpha,\beta}(t)\},
 \qquad a\in\mathbb R,\quad b\in[0,\infty).
\end{equation}
This is the standard $1$-parameterization of stable laws: in the usual
scale notation $\gamma$, $b=\gamma^\alpha$ for $\alpha\ne1$, and
$b=\gamma$ for $\alpha=1$; see \cite[Definition~1.8]{Nolan2018}.
At $b=0$ the measure is $\delta_a$.  At $\alpha=2$ the choice of $\beta$
has no effect, and the variance is $2b$.
The parameter $a$ is a location parameter, not in general an expectation.

We first record the analytic criterion in a form which covers both real
and nonreal sampling points.

\begin{lemma}[A Fourier--Laplace sampling criterion]
\label{p09:lem:stable-sampling}
Let $\mu,\nu$ be probability measures on $\mathbb R\times[0,\infty)$.
Suppose that for every $t\in\mathbb R\setminus\{0\}$ there are distinct
points $z_j(t)$ in the open right half-plane such that
\begin{equation}\label{p09:eq:stable-zero-data}
 \int e^{ita-z_j(t)b}\,d\mu(a,b)
 =\int e^{ita-z_j(t)b}\,d\nu(a,b)
\end{equation}
and
\begin{equation}\label{p09:eq:stable-zero-divergence}
 \sum_j\frac{\operatorname{Re}z_j(t)}{1+|z_j(t)|^2}=\infty.
\end{equation}
Then $\mu=\nu$.
\end{lemma}

\begin{proof}
Write $\rho=\mu-\nu$ and, for fixed real $t$, define
\[
 H_t(z)=\int e^{ita-zb}\,d\rho(a,b),\qquad \operatorname{Re}z>0.
\]
This is bounded by $2$ and is holomorphic.  Indeed, on a compact subset
of the right half-plane there is $\varepsilon>0$ with
$\operatorname{Re}z\geq\varepsilon$, and each integrand derivative is
bounded in absolute value by $b^k e^{-\varepsilon b}$, a bounded function
of $b\geq0$.  Thus differentiation under the finite signed measure is
justified without a moment assumption.

A nonzero bounded holomorphic function on the right half-plane has zeros
satisfying the Blaschke condition
$\sum_j\operatorname{Re}z_j/(1+|z_j|^2)<\infty$.
For completeness, map the half-plane to the disk by
$w=(z-1)/(z+1)$.  Jensen's formula for a nonzero bounded holomorphic
function on the disk gives $\sum_j(1-|w_j|)<\infty$, after removing a
possible zero at the origin.  Moreover,
\[
 1-|w|^2=\frac{4\operatorname{Re}z}{|z+1|^2},\qquad
 1+|z|^2\leq |z+1|^2\leq2(1+|z|^2),
\]
so the disk and half-plane conditions are equivalent.
Equations \eqref{p09:eq:stable-zero-data}--\eqref{p09:eq:stable-zero-divergence}
therefore imply $H_t\equiv0$ for every $t\ne0$.
Dominated convergence in $t$ gives $H_0\equiv0$ as well.

The finite signed measure $d\rho_1(a,b)=e^{-b}\,d\rho(a,b)$ now has
Fourier transform
\[
 \int e^{ita+iub}\,d\rho_1(a,b)=H_t(1-iu)=0
 \qquad(t,u\in\mathbb R).
\]
Fourier uniqueness yields $\rho_1=0$.
Multiplication by $e^b$ on each strip $0\leq b\leq m$, followed by
$m\to\infty$, yields $\rho=0$.
\end{proof}

\begin{theorem}[Stable-family classification]
\label{p09:thm:stable-classification}
Fix $0<\alpha\leq2$ and $-1\leq\beta\leq1$.
Consider exchangeable sequences obtained by first drawing an arbitrary
random pair $(A,B)\in\mathbb R\times[0,\infty)$ and, conditionally on this
pair, drawing independent variables with law $K_{A,B}^{\alpha,\beta}$.
The one-dimensional laws of all partial sums determine the law of $(A,B)$,
and hence the entire exchangeable law, if and only if
\[
                         (\alpha,\beta)\ne(1,0).
\]
No independence between $A$ and $B$, and no moment or exponential-moment
condition on their joint law, is required.

More generally, let $\mathcal I\subset\mathbb N$ be a set of observed
partial-sum lengths.  The following are sufficient for determination:
\begin{align}
 &\sum_{n\in\mathcal I}n^{-|1-\alpha|}=\infty,
       &&\alpha\ne1;\label{p09:eq:stable-sparse-power}\\
 &\sum_{n\in\mathcal I}\frac{1}{1+(\log n)^2}=\infty,
       &&\alpha=1,\quad\beta\ne0.\label{p09:eq:stable-sparse-log}
\end{align}
The sparse sampling conditions are asserted only as sufficient conditions.
\end{theorem}

\begin{proof}
Let $\mu$ and $\nu$ be two mixing laws with identical observed
partial-sum distributions, and write $S_n$ for a partial sum.
Conditional independence gives
\begin{equation}\label{p09:eq:stable-normalized-data}
 \mathbb E e^{itS_n/n}
   =\int\exp\{ita-nb\psi_{\alpha,\beta}(t/n)\}\,d\mu(a,b).
\end{equation}
The same formula holds with $\nu$, and the two quantities agree for
$n\in\mathcal I$.

Suppose first that $\alpha\ne1$.
For fixed $t\ne0$, set
\[
 c_t=|t|^\alpha
       \bigl(1-i\beta\tan(\pi\alpha/2)\operatorname{sgn}t\bigr).
\]
The sampling points in \eqref{p09:eq:stable-normalized-data} are
$z_n(t)=c_t n^{1-\alpha}$.
They are distinct and have positive real part.  For fixed $t$,
\[
 \frac{\operatorname{Re}z_n(t)}{1+|z_n(t)|^2}
 =\frac{(\operatorname{Re}c_t)n^{1-\alpha}}
         {1+|c_t|^2 n^{2(1-\alpha)}}
 \asymp n^{-|1-\alpha|}\qquad(n\to\infty).
\]
Consequently \eqref{p09:eq:stable-sparse-power} allows application of
Lemma~\ref{p09:lem:stable-sampling} and gives $\mu=\nu$.
For the full set $\mathcal I=\mathbb N$, this condition holds because
$0<|1-\alpha|\leq1$.

Now suppose that $\alpha=1$ and $\beta\ne0$.
Put $k=2\beta/\pi$.  Directly from \eqref{p09:eq:stable-exponent},
\begin{align*}
 z_n(t)=n\psi_{1,\beta}(t/n)
 &=|t|+ikt\log|t|-ikt\log n.
\end{align*}
These are distinct points on the vertical line
$\operatorname{Re}z=|t|>0$.  The logarithmic location shift must be
retained; omitting it would incorrectly identify this case with the
symmetric Cauchy case.  Since $kt\ne0$,
\[
 \frac{\operatorname{Re}z_n(t)}{1+|z_n(t)|^2}
  =\frac{|t|}{1+t^2+k^2t^2(\log|t|-\log n)^2}
  \asymp\frac{1}{1+(\log n)^2}.
\]
Condition \eqref{p09:eq:stable-sparse-log} and the same lemma give $\mu=\nu$.
In particular the condition holds on $\mathbb N$.

It remains to exhibit failure for $\alpha=1$, $\beta=0$.
Fix $b,c>0$ and let $Z_0,Z_1,Z_2,\ldots$ and
$Z'_1,Z'_2,\ldots$ be independent standard symmetric Cauchy variables.
Define
\[
 X_j=cZ_0+bZ_j,\qquad Y_j=(b+c)Z'_j.
\]
The first sequence has random location $A=cZ_0$ and fixed scale $B=b$;
the second has fixed parameters $(A,B)=(0,b+c)$.
For every $n\geq1$, both partial sums are centered Cauchy variables
with scale $n(b+c)$.  Nevertheless $X_1-X_2$ has Cauchy scale $2b$,
whereas $Y_1-Y_2$ has scale $2(b+c)$, so the exchangeable laws differ.
All conditional distributions in this counterexample are nondegenerate.
\end{proof}

\begin{remark}[What the classification says]
For each fixed $(\alpha,\beta)$, the parameters $(a,b)$ specify distinct
laws in \eqref{p09:eq:stable-kernel}: the modulus of the characteristic
function determines $b$, after which its phase determines $a$.
The parameterization is continuous into the space of probability measures
with its weak topology.  Its injectivity and the Lusin--Souslin theorem
give a Borel inverse on its Borel image.  Thus the preceding theorem is
also a statement about the goodness of the corresponding measurable
class of directing distributions.
It does not address mixtures over different values of $\alpha$ or $\beta$.
The analytic ingredients and the stable scaling formulas are classical;
the conclusion here is their application to identification from marginal
partial-sum distributions.
\end{remark}




\section{Non-identifiability without finite differential barriers}
\label{p09:sec:no-differential-barrier}

The remark following Proposition~6 and the first question in the
discussion of \cite{LeBarilettoRinaldoHo2026} ask whether absence of
parameter differential structure suffices for identifiability of an
infinite mixture.  We give a negative answer for both classes explicitly
defined there: finite-order parameter differential operators and finite
shift, finite-order difference--differential parameter operators.  The
obstruction in the example is a missing global integral mode.  No claim
is made about exclusion of arbitrary infinite-order or nonlocal
operators.

Let $P_n$ be the ordinary Legendre polynomial with $P_n(1)=1$, and set
\[
 L_n(\theta)=P_n(2\theta-1),\qquad 0\le\theta\le1.
\]
We use the standard elementary identities
\begin{equation}
 L_0=1,\qquad L_1=2\theta-1,\qquad |L_n(\theta)|\le1,
 \qquad
 \int_0^1 L_n(\theta)L_m(\theta)\,d\theta
       =\frac{\mathbf1_{\{n=m\}}}{2n+1}.
 \label{p09:eq:legendre-orthogonality}
\end{equation}
The sequence $(L_0,\ldots,L_N)$ is a basis for polynomials of degree at
most $N$.

\begin{lemma}[A missing integral mode has no finite jet representation]
\label{p09:lem:no-finite-jet}
Let $y_1,\ldots,y_J$ be distinct points of $[0,1]$, let
$r_1,\ldots,r_J\in\mathbb N_0$, and define, on polynomials,
\[
 T(p)=\sum_{j=1}^J\sum_{k=0}^{r_j} b_{j,k}p^{(k)}(y_j).
\]
If $T(1)=0$ and $T(L_n)=0$ for every $n\ge2$, then every coefficient
$b_{j,k}$ is zero.
\end{lemma}

\begin{proof}
The Legendre expansion of a polynomial and
\eqref{p09:eq:legendre-orthogonality} imply
\begin{equation}
 T(p)=3T(L_1)\int_0^1 p(\theta)L_1(\theta)\,d\theta
 \quad\text{for every polynomial }p.
 \label{p09:eq:finite-jet-integral}
\end{equation}
Put $Q(\theta)=\prod_{j=1}^J(\theta-y_j)^{r_j+1}$.
All jets entering $T$ vanish on $Q^2L_1$, hence
\[
 0=T(Q^2L_1)
  =3T(L_1)\int_0^1 Q(\theta)^2L_1(\theta)^2\,d\theta.
\]
The integral is strictly positive, so $T(L_1)=0$.  Consequently $T$
annihilates every polynomial.  Hermite interpolation at the distinct
points $y_j$ provides, for each pair $(j,k)$, a polynomial whose listed
jets are zero except for $p^{(k)}(y_j)=1$.  Applying $T$ to it gives
$b_{j,k}=0$.
\end{proof}

\begin{theorem}[A strictly positive polynomial kernel with an invisible mode]
\label{p09:thm:missing-mode-discrete}
Let $\Theta=[0,1]$, let
$\mathcal X=\{\sigma n:\sigma\in\{-1,1\},\ n\ge2\}\subset\mathbb R$
carry counting measure, and put $w_n=2^{1-n}$.  Define
\begin{equation}
 f(\sigma n\mid\theta)
    =\frac{w_n}{2}\left(1+\frac{\sigma}{2}L_n(\theta)\right).
 \label{p09:eq:discrete-missing-mode-kernel}
\end{equation}
Then all the following hold.
\begin{enumerate}
\item For every $\theta$, $f(\cdot\mid\theta)$ is a strictly positive
probability mass function.  Each parameter section is a polynomial,
and $\theta\mapsto f(\cdot\mid\theta)$ is injective.
\item For the mixture operator $A\nu=\int f(\cdot\mid\theta)\nu(d\theta)$
on finite signed measures,
\begin{equation}
 \ker A=\{bL_1(\theta)\,d\theta:b\in\mathbb R\}.
 \label{p09:eq:exact-missing-mode-nullspace}
\end{equation}
In particular, for every $0<c<1$ the distinct probabilities
\begin{equation}
 G_\pm(d\theta)=(1\pm cL_1(\theta))\,d\theta
 \label{p09:eq:indistinguishable-smooth-mixings}
\end{equation}
have strictly positive analytic densities and satisfy
$AG_+=AG_-= (w_n/2)_{\sigma n\in\mathcal X}$.
\item No nonzero finite combination of parameter evaluations and
derivatives at finitely many distinct points annihilates every section
$f(x\mid\cdot)$.  Consequently the kernel satisfies no nontrivial
parameter PDE of finite order on any nonempty open subset of
$(0,1)$, and no nontrivial finite-shift, finite-order
difference--differential parameter relation on any allowed open set.
The latter statements hold even with arbitrary coefficient functions,
provided the asserted identities hold pointwise.
\end{enumerate}
\end{theorem}

\begin{proof}
Since $\sum_{n\ge2}w_n=1$ and $|L_n|\le1$, the values in
\eqref{p09:eq:discrete-missing-mode-kernel} lie in
$[w_n/4,3w_n/4]$ and sum to one after summing the two signs and then $n$.
Measurability is immediate.

Let $\nu$ be a finite signed measure.  Adding the two values
$(A\nu)(n)$ and $(A\nu)(-n)$ shows that $A\nu=0$ implies
$\nu([0,1])=0$.  Subtracting them then gives
$\int L_n\,d\nu=0$ for every $n\ge2$.
For any polynomial $p$, expansion in the Legendre basis therefore yields
\[
 \int_0^1p\,d\nu
 =3\left(\int_0^1L_1\,d\nu\right)
       \int_0^1p(\theta)L_1(\theta)\,d\theta.
\]
Polynomials are uniformly dense in $C([0,1])$, so finite-measure
uniqueness gives
$\nu=bL_1(\theta)d\theta$, with
$b=3\int L_1\,d\nu$.  Conversely every such measure is annihilated by
$A$, by orthogonality.  This proves
\eqref{p09:eq:exact-missing-mode-nullspace} and the assertions about
$G_\pm$.  If $f(\cdot\mid\theta)=f(\cdot\mid\theta')$, then
$\delta_\theta-\delta_{\theta'}$ is in this null space.  A measure that
is both purely atomic and absolutely continuous must be zero, so
$\theta=\theta'$.

Now let $T$ be a finite jet functional as in
Lemma~\ref{p09:lem:no-finite-jet}, and suppose
$T(f(\sigma n\mid\cdot))=0$ for every sign and every $n\ge2$.
Adding the two identities gives $T(1)=0$, and subtracting them gives
$T(L_n)=0$.  The lemma forces $T=0$ coefficientwise.

For a parameter differential operator
$\mathcal L_\theta=\sum_{k=0}^r c_k(\theta)\partial_\theta^k$,
fix any parameter value where the annihilation identity holds and apply
the preceding conclusion with one node $y_1=\theta$.  All coefficients
vanish there.  For a difference--differential operator as defined in
\cite[Definition~4]{LeBarilettoRinaldoHo2026}, fix $\theta$ and use the
distinct nodes $y_j=\theta+v_j$ and coefficients $b_{j,k}=c_{j,k}(\theta)$.
Every coefficient again vanishes.  With counting measure the
almost-everywhere statement in the sample variable excludes no sample
points, so these pointwise arguments apply exactly to the operator
definitions in that paper.
\end{proof}

\begin{corollary}[Singular mixing measures remain identifiable]
\label{p09:cor:singular-identification}
For the kernel in Theorem~\ref{p09:thm:missing-mode-discrete}, every mixing
probability singular with respect to Lebesgue measure is identifiable
among all Borel mixing probabilities.  This includes every finite or
countably discrete mixing probability.
\end{corollary}

\begin{proof}
If $AG=AH$ and $G$ is singular, then
$H=G+bL_1(\theta)d\theta$ by
\eqref{p09:eq:exact-missing-mode-nullspace}.  The absolutely continuous part
of the nonnegative measure $H$ is $bL_1(\theta)d\theta$.
Because $L_1$ has both signs on sets of positive Lebesgue measure, this
can be nonnegative only when $b=0$.
\end{proof}

\begin{corollary}[Two observations sharing a latent parameter restore identification]
\label{p09:cor:two-replicate-repair}
The mixture model with kernel
$f(x\mid\theta)f(y\mid\theta)$ on $\mathcal X^2$ is identifiable
over all Borel probability measures on $[0,1]$.
The two observations are conditionally independent given the same
latent parameter; two independent draws with separately redrawn
latent parameters do not have this property.
\end{corollary}

\begin{proof}
Let a finite signed measure $\nu$ annihilate the two-observation kernel.
Marginalizing over $y$ shows that $A\nu=0$, so
$\nu=bL_1(\theta)d\theta$.  The marginalization is legitimate because
the sums of absolute values are bounded by $\|\nu\|_{\rm TV}$.
Take the signed sum over the four pairs $(x,y)=(\sigma2,\tau3)$:
\[
 \sum_{\sigma,\tau\in\{-1,1\}}\sigma\tau
    f(\sigma2\mid\theta)f(\tau3\mid\theta)
       =\frac{w_2w_3}{4}L_2(\theta)L_3(\theta).
\]
It follows that $b\int_0^1L_1L_2L_3=0$.
The Legendre recurrence gives
$L_1L_2=(3/5)L_3+(2/5)L_1$, so
\[
 \int_0^1 L_1L_2L_3\,d\theta=\frac{3}{35}>0.
\]
Hence $b=0$ and $\nu=0$.
\end{proof}

\begin{remark}[Exact separation and the scope of the counterexample]
For \eqref{p09:eq:indistinguishable-smooth-mixings}, direct integration of
their distribution functions gives
$W_1(G_+,G_-)=c/3$, while their induced observation laws coincide.
Thus the failure is exact, despite distinct component distributions,
positive smooth mixing densities, and absence of both finite operator
classes.  It answers the general sufficiency question negatively for
the two precisely specified differential structures in
\cite{LeBarilettoRinaldoHo2026}.  It does not contradict their theorems,
which assert sufficiency of such structures for non-identifiability.
An arbitrary nonlocal obstruction is present: integration against
$L_1(\theta)d\theta$ annihilates the kernel by construction.
\end{remark}

\begin{corollary}[A jointly entire continuous-density counterexample]
\label{p09:cor:entire-continuous-counterexample}
Let both the sample space and the parameter space be $[0,1]$, with
Lebesgue measure as the dominating sample measure, and define
\begin{equation}
 k(x\mid\theta)
   =1+\sum_{n=2}^{\infty}\frac{L_n(x)L_n(\theta)}{n!}.
 \label{p09:eq:continuous-missing-mode-kernel}\end{equation}
This is a probability density in $x$ for every $\theta$, with
$k(x\mid\theta)\ge3-e>0$ on $[0,1]^2$, and extends to an entire
function of the two complex variables $(x,\theta)$.
Its mixture operator has exactly the null space
$\{bL_1(\theta)d\theta:b\in\mathbb R\}$.
Its component distributions are distinct, every finite family of
parameter jets at distinct parameter points is linearly independent
as a family of functions of $x$, and it admits neither of the finite
differential operator classes excluded in
Theorem~\ref{p09:thm:missing-mode-discrete}.
The probabilities $G_\pm$ in
\eqref{p09:eq:indistinguishable-smooth-mixings} both give the uniform sample
density.  Every singular mixing probability remains identifiable, and
two conditionally independent observations sharing one latent
parameter restore identifiability over all mixing probabilities.
\end{corollary}

\begin{proof}
On the real unit square, absolute uniform convergence and
$\sum_{n\ge2}1/n!=e-2$ give the stated lower bound.  Orthogonality
allows termwise integration in $x$ and gives integral one.
For completeness, the elementary polynomial formula
\[
 P_n(z)=2^{-n}\sum_{j=0}^n\binom nj^2
             (z-1)^{n-j}(z+1)^j
\]
shows that on every complex disc $|z|\le R$ one has
$|P_n(z)|\le[2(R+1)]^n$, using
$\sum_j\binom nj^2=\binom{2n}{n}\le4^n$.
After replacing $z$ by $2z-1$, this gives an exponential-in-$n$
bound for $L_n$ on every compact complex set.  The factorial denominator
therefore makes \eqref{p09:eq:continuous-missing-mode-kernel} locally
uniformly convergent on $\mathbb C^2$.  It defines an entire function,
and the series can be differentiated termwise to every fixed order,
locally uniformly.  These observations justify all subsequent
coefficient extraction and parameter differentiation.

If $B\nu(x)=\int k(x\mid\theta)\nu(d\theta)$, then
\begin{equation}
 \int_0^1 B\nu(x)\,dx=\nu([0,1]),\qquad
 \int_0^1 B\nu(x)L_m(x)\,dx
   =\frac{\int_0^1L_m(\theta)\nu(d\theta)}{m!(2m+1)}
 \quad(m\ge2).
 \label{p09:eq:continuous-legendre-coefficients}
\end{equation}
Thus $B\nu=0$ if and only if $\nu$ has zero mass and zero Legendre
moments of degrees $m\ge2$.  The proof of
\eqref{p09:eq:exact-missing-mode-nullspace} applies without change, and
also proves distinctness of component distributions, the two uniform
mixtures, and identifiability of singular mixing probabilities.

Let $T$ be any finite jet functional as in
Lemma~\ref{p09:lem:no-finite-jet}.  If $T(k(x\mid\cdot))=0$ for almost
every $x$, integrate against $1$ and against $L_m(x)$ in
\eqref{p09:eq:continuous-legendre-coefficients}.  Parameter differentiation
commutes with these integrals, so $T(1)=T(L_m)=0$ for all $m\ge2$.
The lemma forces every coefficient of $T$ to vanish.  In particular
every finite family of parameter jets is linearly independent.
Applying this argument at each fixed parameter point excludes both
finite differential and finite difference--differential operators,
exactly as in the discrete example.

Finally suppose a finite signed measure $\nu$ annihilates the
two-observation kernel $k(x\mid\theta)k(y\mid\theta)$.
Marginalization gives $\nu=bL_1(\theta)d\theta$.
Integrating the joint mixture against $L_2(x)L_3(y)$ and using
\eqref{p09:eq:continuous-legendre-coefficients} gives
\[
 0=\frac{1}{2!\,5\,3!\,7}
          \int_0^1L_2(\theta)L_3(\theta)\nu(d\theta)
   =\frac{b}{2!\,5\,3!\,7}\frac{3}{35}.
\]
Therefore $b=0$ and $\nu=0$.
\end{proof}




\section{Finite observational defects and replicated observations}

The one-dimensional invisible space above is a special case of a general
finite-dimensional phenomenon. Two replicated observations force every
remaining invisible measure onto one fixed finite set.
\label{p10:sec:finite-defect}

Let $\Theta$ and $\mathcal X$ be compact metric spaces, and let
$\theta\mapsto P_\theta$ be a weakly continuous family of Borel probability
measures on $\mathcal X$.  Write $\mathcal M(\Theta)$ for the real vector space
of finite signed Borel measures on $\Theta$.  The observation operator with
$m$ conditionally independent observations sharing the same latent parameter is
\[
 A_m\mu=\int_\Theta P_\theta^{\otimes m}\,d\mu(\theta),
 \qquad K_m=\ker A_m.
\]
Thus $A_1$ describes one observation from each independently selected latent
parameter, whereas $A_m$ describes a group of $m$ observations from one common
parameter.  These are different statistical experiments.

For $f\in C(\mathcal X;\mathbb R)$, put
\[
 Tf(\theta)=\int f\,dP_\theta,
 \qquad
 V=\overline{T(C(\mathcal X;\mathbb R))}^{\|\cdot\|_\infty}
       \subset C(\Theta;\mathbb R).
\]
The function $1$ belongs to $V$.  The family $P_\theta$ is injectively
parameterized if and only if $V$ separates the points of $\Theta$.
Products of functions of the separate observation coordinates are dense in
$C(\mathcal X^m)$, by the real Stone--Weierstrass theorem.  Consequently,
\begin{equation}\label{p10:eq:finite-defect-annihilator}
 K_m=\left\{\nu:\int v_1\cdots v_m\,d\nu=0
                    \text{ for all }v_1,\ldots,v_m\in V\right\}.
\end{equation}
In particular, $K_1=V^\perp$.  If $d=\dim K_1<\infty$, then $V$ has
codimension $d$ in $C(\Theta;\mathbb R)$: the dual of $C(\Theta)/V$ is
$V^\perp$, and a normed space whose dual is finite dimensional is itself
finite dimensional, with the same dimension.

\begin{theorem}[Two observations reduce every finite defect to atoms]
\label{p10:thm:finite-defect-atoms}
Suppose $d=\dim K_1<\infty$.  Every $\nu\in K_2$ is supported on at most
$2d$ points.  More precisely, its positive and negative parts are each
supported on at most $d$ points.  There is one fixed set $E\subset\Theta$
of at most $2d$ points supporting every measure in $K_2$.  Therefore, if $K_1$ contains no nonzero
finitely supported signed measure, then $K_2=\{0\}$.

In particular, two observations identify every mixing probability measure
whenever all measures in $K_1$ are absolutely continuous with respect to a
fixed atomless measure.  The same conclusion follows if every finite
subfamily of $\{P_\theta:\theta\in\Theta\}$ is linearly independent.
\end{theorem}

\begin{proof}
For $\nu\in K_2$, consider the multiplication operator
\[
 M_\nu:C(\Theta;\mathbb R)\longrightarrow\mathcal M(\Theta),
 \qquad M_\nu(g)=g\nu.
\]
Equation~\eqref{p10:eq:finite-defect-annihilator} shows that
$M_\nu(V)\subset K_1$.  Since $V$ has codimension $d$, a vector space
complement $F$ of dimension $d$ gives
\[
 \operatorname{rank}M_\nu
 \leq \dim M_\nu(V)+\dim M_\nu(F)\leq 2d.
\]
If the support of $|\nu|$ contained $r$ distinct points, choose pairwise
disjoint neighborhoods of those points and continuous functions $g_j$
supported in those neighborhoods, each nonzero on a set of positive
$|\nu|$-measure.  The measures $g_j\nu$ would be nonzero and have disjoint
supports, hence be linearly independent.  Thus a multiplication operator of
rank at most $2d$ forces $\nu$ to have at most $2d$ atoms.

For the more precise assertion, let
$\nu=\sum_{i=1}^r c_i\delta_{\theta_i}$ with all $c_i\neq0$, and let $p$
and $q$ be the numbers of positive and negative coefficients.  The restriction
space
\[
 W=\{(v(\theta_1),\ldots,v(\theta_r)):v\in V\}\subset\mathbb R^r
\]
has codimension at most $d$: its annihilator identifies with the subspace of
$K_1$ supported on these $r$ points.  Moreover,
\[
 \sum_{i=1}^r c_i w_i z_i=0\qquad(w,z\in W).
\]
A totally isotropic subspace for a nondegenerate real form of signature
$(p,q)$ has dimension at most $\min(p,q)$.  This follows, for example,
because its projections onto the positive and negative coordinate spaces
are both injective.  Hence
\[
 p+q-d\leq\dim W\leq\min(p,q),
 \qquad \max(p,q)\leq d.
\]
Since $1\in V$, we have $K_2\subset K_1$.  To obtain a fixed set, choose a
basis of the finite-dimensional space $K_2$ and let $E$ be the union of its
members' finite supports.  A linear
combination of the basis vectors whose coefficients avoid finitely many
proper hyperplanes has nonzero mass at every point of $E$.  Its support is
therefore $E$, so the already established bound gives $|E|\leq2d$.
Finally, if $K_1$ contains no nonzero
finitely supported measure, the first part forces $K_2=\{0\}$.  The two
stated sufficient conditions each exclude such a measure.
\end{proof}

The theorem concerns exact identifiability, not uniform statistical stability.
It also states more than the special case $K_2=0$: for an arbitrary
finite-dimensional defect, every non-atomic part of the mixing measure is
determined by the two-observation law, and any remaining difference is finite
atomic.

\begin{theorem}[A finite bound on replication]
\label{p10:thm:finite-defect-replication}
Suppose the component map $\theta\mapsto P_\theta$ is injective and
$d=\dim K_1<\infty$.  Then $K_{d+1}=\{0\}$.  More precisely, until the
nullspace becomes zero, each additional shared-parameter observation strictly
decreases its dimension.  The bound $d+1$ cannot be decreased in general.
\end{theorem}

\begin{proof}
Let $U_m$ be the uniform closure of the span of products of $m$ elements of
$V$.  Since $1\in V$, the spaces $U_m$ increase, their annihilators $K_m$
decrease, and $\dim K_m\leq d$.  If $K_m=K_{m+1}$, the Hahn--Banach theorem
gives $U_m=U_{m+1}$.  For $v\in V$ and $u\in U_m$, approximation by finite
linear combinations of products shows that $vu\in U_{m+1}=U_m$.  Thus $U_m$
contains the unital algebra generated by $V$.  This algebra separates points,
so Stone--Weierstrass gives $U_m=C(\Theta)$ and $K_m=\{0\}$.  Every nonzero
step therefore decreases a nonnegative integer dimension, giving
$K_{d+1}=\{0\}$.

For sharpness, take $d+2$ distinct parameters
$0<t_1<\cdots<t_{d+2}<1$ and $P_{t_i}=\operatorname{Bernoulli}(t_i)$.
The one-observation operator has rank $2$ and nullity $d$.  The
$m$-observation law determines exactly the moments of the mixing measure
through degree $m$.  The corresponding Vandermonde matrix has rank
$\min(m+1,d+2)$.  Hence $K_d\neq0$ when $d\geq1$, whereas $K_{d+1}=0$.
Every nonzero null measure has total mass zero, so its normalized Jordan
parts are distinct mixing probability measures with the same observation
law.  This proves sharpness for probability mixtures as well as for signed
measures.  When $d=0$, one observation already identifies the mixture.
\end{proof}



\section{Analytic kernels with an arbitrary finite polynomial defect}
\label{p10:sec:legendre-finite-defect}

Continue to use the shifted Legendre polynomials and normalization
in \eqref{p09:eq:legendre-orthogonality}. We now prescribe an arbitrary
finite collection of missing modes, rather than only $L_1$.
Fix any nonempty finite set $S\subset\{1,2,\ldots\}$ and put
$d=|S|$, $M=\max S$, and $a_n=2^{-n^2-2}$.  Consider the probability density
\begin{equation}\label{p10:eq:legendre-finite-kernel}
 k_S(x,\theta)=1+\sum_{\substack{n\geq1\\n\notin S}}
                   a_nL_n(x)L_n(\theta),
 \qquad 0\leq x,\theta\leq1.
\end{equation}
The series converges locally uniformly on $\mathbb C^2$, since Legendre
polynomials grow at most exponentially in $n$ on each compact set, while
$a_n$ decays quadratically exponentially.  It therefore defines an entire
function of the two variables.  On the real unit square,
\[
 \frac34<k_S(x,\theta)<\frac54,
 \qquad\int_0^1 k_S(x,\theta)\,dx=1.
\]
We use Lebesgue measure on $[0,1]$ as the observation dominating measure.

\begin{theorem}[An exactly prescribed global ambiguity]
\label{p10:thm:legendre-exact-defect}
For the kernel~\eqref{p10:eq:legendre-finite-kernel},
\begin{equation}\label{p10:eq:legendre-nullspace}
 K_1=\left\{h(\theta)\,d\theta:
                   h\in\operatorname{span}\{L_s:s\in S\}\right\}.
\end{equation}
Consequently:
\begin{enumerate}
\item The defect has exactly dimension $d$; positive analytic mixing
densities can be nonidentifiable from one observation.
\item The singular part, relative to Lebesgue measure, of an arbitrary
mixing measure is determined by its one-observation law.  In particular,
the operator is injective on all singular mixing measures, including all
discrete and all singular continuous mixing measures.
\item Every finite set of component densities is linearly independent.
\item Two shared-parameter observations identify every mixing probability
measure.
\end{enumerate}
\end{theorem}

\begin{proof}
For $\mu\in\mathcal M([0,1])$, termwise integration and orthogonality give
\[
 \int A_1\mu(x)\,dx=\mu([0,1]),\qquad
 \int L_n(x)A_1\mu(x)\,dx
       =\frac{a_n}{2n+1}\int L_n(\theta)\,d\mu(\theta)
       \quad(n\geq1,\ n\notin S).
\]
Thus $A_1\mu=0$ implies that the mass and all moments indexed outside $S$
vanish.  Set
\[
 h(\theta)=\sum_{s\in S}(2s+1)
                  \left(\int L_s\,d\mu\right)L_s(\theta).
\]
The measure $\mu-h(\theta)d\theta$ annihilates every Legendre polynomial,
hence every polynomial, and hence every continuous function by Weierstrass.
It is zero.  Conversely, every measure on the right side of
\eqref{p10:eq:legendre-nullspace} is annihilated by the kernel.

Small nonzero polynomial perturbations $1+h$ and $1-h$ of the uniform
mixing density establish the first conclusion.  Equality of two
one-observation laws forces the difference of the mixing measures to be
absolutely continuous, giving the second conclusion by uniqueness of the
Lebesgue decomposition.  No nonzero finitely supported measure is
absolutely continuous with respect to Lebesgue measure, giving the third
conclusion.  The fourth follows from
Theorem~\ref{p10:thm:finite-defect-atoms}.
\end{proof}

\begin{theorem}[Every finite parameter-jet test is independent]
\label{p10:thm:legendre-jet-independence}
Let $\theta_1,\ldots,\theta_q$ be distinct points of $[0,1]$ and let
$r_1,\ldots,r_q$ be nonnegative integers.  The functions of $x$
\[
 \left\{\partial_\theta^j k_S(x,\theta_i):
                   1\leq i\leq q,\ 0\leq j\leq r_i\right\}
\]
are linearly independent.  In particular, no nonzero finite-order
parameter differential operator, or finite sum of shifted finite-order
parameter differential operators, annihilates this kernel wherever its
evaluation points are admissible.
\end{theorem}

\begin{proof}
It suffices to treat real coefficients, since the kernel is real.
Suppose a finite jet functional
\[
 T(p)=\sum_{i=1}^q\sum_{j=0}^{r_i} c_{ij}p^{(j)}(\theta_i)
\]
annihilates $k_S(x,\cdot)$ for every $x$.  Orthogonality in $x$, with
termwise differentiation justified by local uniform holomorphic
convergence, gives $T(1)=0$ and $T(L_n)=0$ for $n\notin S$.  Hence for
every polynomial $p$,
\[
 T(p)=\int_0^1 p(t)h(t)\,dt,
 \qquad h(t)=\sum_{s\in S}(2s+1)T(L_s)L_s(t).
\]
Set $Q(t)=\prod_{i=1}^q(t-\theta_i)^{r_i+1}$.  Every prescribed jet of
$Q^2h$ vanishes, so
\[
 0=T(Q^2h)=\int_0^1 Q(t)^2h(t)^2\,dt.
\]
Thus $h=0$, and $T$ annihilates all polynomials.  Hermite interpolation at
the distinct nodes $\theta_i$ then gives $c_{ij}=0$ for all $i,j$.
For complex coefficients, apply the real argument separately to their
real and imaginary parts.  For a shifted differential operator, fix the
base parameter, collect coincident evaluation nodes, and apply the same
finite-jet assertion; only the zero functional can remain.
\end{proof}

The assertion concerns finite differential and finite shift-differential
operators.  The kernel does have the nonlocal integral annihilators
$h(t)dt$ in~\eqref{p10:eq:legendre-nullspace}; no absence of nonlocal or
infinite-order obstructions is claimed.

\subsection{Exactly $d$ scalar correlations repair the defect}

The two-observation conclusion admits an explicit finite reconstruction.
Let
\[
 u_r=\int_0^1 L_r(\theta)\,d\mu(\theta).
\]
The one-observation law gives $u_0=1$ and every $u_r$ for $r\notin S$.
Fix any integer $N>M$.  For $s\in S$, the joint law of two observations
$(X_1,X_2)$ sharing the same latent parameter gives
\begin{equation}\label{p10:eq:legendre-repair-correlation}
 b_s=\frac{(2N+1)(2N+2s+1)}{a_Na_{N+s}}
             \mathbb E[L_N(X_1)L_{N+s}(X_2)]
     =\int_0^1 L_N(\theta)L_{N+s}(\theta)\,d\mu(\theta).
\end{equation}
The scaling follows from
$\mathbb E[L_n(X_i)\mid\theta]=a_nL_n(\theta)/(2n+1)$ for $n\notin S$.

\begin{proposition}[A minimal number of scalar repairs]
\label{p10:prop:legendre-minimal-repair}
The one-observation law together with the $d$ correlations in
\eqref{p10:eq:legendre-repair-correlation} determines $\mu$ uniquely.
Among procedures that add expectations of a fixed finite list of bounded
two-observation statistics to the one-observation law, fewer than $d$
scalar expectations cannot identify every mixing probability measure.
\end{proposition}

\begin{proof}
Write the finite polynomial identity
\[
 L_N L_{N+s}=\sum_{r=0}^{2N+s}c_{sr}L_r,
 \qquad
 c_{sr}=(2r+1)\int_0^1 L_N(t)L_{N+s}(t)L_r(t)\,dt.
\]
If $r<s$, the polynomial $L_NL_r$ has degree less than $N+s$, so
$c_{sr}=0$.  Let $\kappa_j=\binom{2j}{j}$ be the leading coefficient of
$L_j$.  Comparing leading coefficients in $L_NL_s$ and using
orthogonality yields
\[
 c_{ss}=\frac{(2s+1)\kappa_N\kappa_s}
                    {(2N+2s+1)\kappa_{N+s}}>0.
\]
Subtracting the already known terms $c_{sr}u_r$ with $r\notin S$ from
$b_s$ leaves a triangular linear system for the $d$ missing moments $u_s$.
It is solved by descending through $S$, since its diagonal entries are
nonzero.  All Legendre moments are then known, and polynomial density on
$[0,1]$ determines $\mu$.

For minimality, consider the open family of mixing densities
\[
 p_c(\theta)=1+\sum_{s\in S}c_sL_s(\theta),
 \qquad \sum_{s\in S}|c_s|<\frac12.
\]
All these densities have the same one-observation law.  The expectation
of any fixed bounded statistic of the two observations is an affine
function of $c\in\mathbb R^d$, by linearity of mixing.  Fewer than $d$
such expectations define an affine map with nontrivial linear kernel,
which is not injective on this open parameter set.  This proves the
claimed lower bound for scalar expectation repairs.
\end{proof}

The reconstruction is exact but may be ill-conditioned: the factors
$1/(a_Na_{N+s})$ can be very large.  No finite-sample accuracy guarantee is
implied by this identifiability statement.




\section{Arbitrary replication thresholds on the analytic circle}

When the one-observation nullspace is infinite dimensional, the sharp
finite-defect bound does not apply. The following construction realizes
every threshold while keeping all finite parameter jets independent.
Its Fourier mechanism differs from the interval construction in
Section~\ref{p11:sec:interval-core}: here the indistinguishable mixing
densities themselves can be analytic.
\label{p10:sec:replicate-hierarchy}

Let $\mathbb T=\mathbb R/(2\pi\mathbb Z)$ and write
$m(d\theta)=d\theta/(2\pi)$ for its Haar probability measure.
The parameter space in this section is the analytic circle itself: the
representatives $0$ and $2\pi$ denote the same parameter, not two distinct
points of a closed interval.  Both the observations and the parameter take
values in $\mathbb T$.

For a conditional density $k(x\mid\theta)$ and a Borel probability measure
$\mu$ on $\mathbb T$, let
\[
 M_q^\mu(dx_1\cdots dx_q)
 =\left[\int_{\mathbb T}\prod_{j=1}^q
          k(x_j\mid\theta)\,\mu(d\theta)\right]
       m(dx_1)\cdots m(dx_q).
\]
Thus the $q$ observations share one latent parameter and are conditionally
independent.  Identifiability at order $q$ means injectivity of
$\mu\mapsto M_q^\mu$ on all Borel probability measures.

\begin{theorem}[Every finite replication threshold, and an infinite one]
\label{p10:thm:every-replication-threshold}
For every $r\in\{1,2,3,\ldots\}\cup\{\infty\}$ there is a strictly positive
conditional density $k_r$ on $\mathbb T\times\mathbb T$ with the following
properties.
\begin{enumerate}
\item It extends to an entire, $2\pi$-periodic function of each complex
variable jointly, and its real values are bounded below by a positive
absolute constant.
\item For any distinct $\theta_1,\ldots,\theta_J\in\mathbb T$ and any
nonnegative integers $d_1,\ldots,d_J$, the functions
\[
 \left\{\partial_\theta^\ell k_r(\,\cdot\mid\theta_j):
       1\le j\le J,\ 0\le\ell\le d_j\right\}
\]
are linearly independent over $\mathbb C$.
\item If $r<\infty$, identifiability at order $q$ holds exactly when
$q\ge r$.  If $r=\infty$, it fails for every finite $q$, although the
collection $(M_q^\mu)_{q\ge1}$ identifies $\mu$.
\item Every finitely supported mixing measure is identifiable among
all Borel probability measures from one observation.  At every order
where unrestricted identifiability fails, the failure can be witnessed
by two distinct, strictly positive analytic mixing densities.
For the assertion about unrestricted competitors, see
Corollary~\ref{p10:cor:finite-atomic-all-competitors}.
\end{enumerate}
In particular, absence of finite-order differential or finite-shift
differential annihilators supplies no finite universal bound on the
number of shared-latent observations required for unrestricted mixture
identification.
\end{theorem}

We prove the theorem using a Fourier criterion and an arithmetic
independence lemma.  For a symmetric subset $S\subseteq\mathbb Z$
containing zero, put
\begin{equation}\label{p10:eq:spectral-kernel}
 k_S(x\mid\theta)
 =1+\sum_{n\in S\cap\mathbb N}e^{-n^2}\cos(n(x-\theta)).
\end{equation}
Here $\mathbb N=\{1,2,\ldots\}$.  The series and all its derivatives
converge uniformly on compact subsets of $\mathbb C^2$, since on every
such compact set the summands are bounded by $e^{-n^2+Cn}$ times a
polynomial in $n$.  Its integral in $x$ is one, and
\[
 k_S(x\mid\theta)\ge 1-\sum_{n\ge1}e^{-n^2}
 \ge 1-\frac1{e-1}>0.
\]

For $q\ge1$ let
$qS=\{s_1+\cdots+s_q:s_1,\ldots,s_q\in S\}$.
The inclusion $0\in S$ makes these sets increasing in $q$.

\begin{lemma}[Replication is addition of spectral sets]
\label{p10:lem:replication-sumset}
For the kernel \eqref{p10:eq:spectral-kernel}, $M_q^\mu=M_q^\nu$ if and
only if $\widehat\mu(n)=\widehat\nu(n)$ for every $n\in qS$, where
$\widehat\mu(n)=\int e^{-in\theta}\,\mu(d\theta)$.
Consequently, identifiability at order $q$ holds if and only if
$qS=\mathbb Z$.
\end{lemma}

\begin{proof}
The Fourier coefficient of $k_S(\cdot\mid\theta)$ at $n$ is
$a_n e^{-in\theta}$, where $a_0=1$,
$a_n=e^{-n^2}/2$ for $n\in S\setminus\{0\}$, and $a_n=0$ otherwise.
Thus the Fourier coefficient of $M_q^\mu$ at
$(n_1,\ldots,n_q)\in\mathbb Z^q$ is
\[
 \left(\prod_{j=1}^q a_{n_j}\right)
 \widehat\mu(n_1+\cdots+n_q).
\]
Uniqueness of Fourier coefficients for finite measures proves the first
assertion and the sufficient condition.  If $b\notin qS$, then $b\ne0$
and symmetry also gives $-b\notin qS$.  For any $0<c<1$ the densities
$1+c\cos(b\theta)$ and $1-c\cos(b\theta)$ relative to $m$ are distinct
and strictly positive, but their Fourier coefficients agree on $qS$.
\end{proof}

\begin{lemma}[Arithmetic richness implies independence of all finite jets]
\label{p10:lem:AP-jets}
Suppose that $S\cap\mathbb N$ meets every residue class modulo every
positive integer infinitely often.  Then $k_S$ has the independence
property in part~\textup{(2)} of
Theorem~\ref{p10:thm:every-replication-threshold}.
\end{lemma}

\begin{proof}
Suppose a linear combination of the displayed functions vanishes.
Taking its Fourier coefficients at $n\in S\cap\mathbb N$ gives
\[
 E(n)=\sum_{j=1}^J P_j(n)\lambda_j^n=0,
 \qquad
 \lambda_j=e^{-i\theta_j},\quad
 P_j(n)=\sum_{\ell=0}^{d_j}c_{j\ell}(-in)^\ell.
\]
The numbers $\lambda_j$ are distinct and nonzero.  The sequence $E(n)$
satisfies a constant-coefficient linear recurrence over $\mathbb C$.
By the Skolem--Mahler--Lech theorem, its zero set in the nonnegative
integers is a finite union of arithmetic progressions together with
a finite set; see \cite[Theorem~1.3]{BellLagarias2015}.
In particular its zero set is eventually periodic, with some period
$Q\ge1$.  If any residue class modulo $Q$ failed eventually to consist
of zeros, its infinitely many intersections with $S\cap\mathbb N$
would give a contradiction.  Hence $E(n)=0$ for all sufficiently large
$n$.

For completeness, an exponential polynomial with distinct nonzero bases
cannot be eventually zero unless every $P_j$ is zero.  If $P_j\ne0$,
apply to the sequence the shift operator
$\prod_{h\ne j}(\mathcal E-\lambda_h)^{d_h+1}$, where
$(\mathcal E u)(n)=u(n+1)$.  It kills every term except the $j$th and
leaves $\lambda_j^n Q_j(n)$ for a polynomial $Q_j$ of the same degree
as $P_j$, with nonzero leading coefficient.  Such a sequence is not
eventually zero.  This contradiction proves $P_j=0$ for every $j$,
and hence all the original coefficients vanish.
\end{proof}

This lemma also excludes operators with finitely many shifted parameter
evaluations and derivatives, even if their coefficients vary arbitrarily
with the base point: at each base point, combine coincident nodes and
apply the lemma to the resulting distinct nodes.  Only the zero
functional remains.  It does not exclude infinite-order or genuinely
nonlocal annihilators.

\subsection{A finite threshold prescribed in advance}

Fix $r\ge2$ and choose
\begin{equation}\label{p10:eq:Bohr-parameters}
 \alpha_r=\frac12-\frac{\sqrt2}{10r},\qquad
 \delta_r=\frac{2r-1}{4r(r-1)},\qquad
 S_r=\{n\in\mathbb Z:\|n\alpha_r\|_{\mathbb R/\mathbb Z}<\delta_r\}.
\end{equation}
These choices satisfy
\begin{equation}\label{p10:eq:Bohr-inequalities}
 r\delta_r>\frac12,
 \qquad
 (r-1)\delta_r=\frac12-\frac1{4r}<\alpha_r<\frac12.
\end{equation}

\begin{lemma}\label{p10:lem:Bohr-order}
The set $S_r$ meets every residue class modulo every positive integer
infinitely often, and
\[
 rS_r=\mathbb Z,\qquad 1\notin(r-1)S_r.
\]
\end{lemma}

\begin{proof}
For a fixed residue class $a$ modulo $Q$, the numbers
$(a+Qj)\alpha_r$ are dense modulo one as $j\to\infty$, because
$Q\alpha_r$ is irrational.  They therefore enter the open interval
$(-\delta_r,\delta_r)$ modulo one infinitely often.
The triangle inequality for distance modulo one and
\eqref{p10:eq:Bohr-inequalities} give
$1\notin(r-1)S_r$.

Now fix $b\in\mathbb Z$, and let $y\in[-1/2,1/2]$ represent
$b\alpha_r$ modulo one.  We have $|y/r|<\delta_r$.
Choose integers $s_1,\ldots,s_{r-1}$ such that their respective
representatives $y_1,\ldots,y_{r-1}$ modulo one are sufficiently close
to $y/r$.  Irrationality supplies these integers, and the choices can
be made so that all $|y_j|<\delta_r$ and
$|y-\sum_{j<r}y_j|<\delta_r$.
Set $s_r=b-\sum_{j<r}s_j$.  Its phase modulo one is represented by
$y-\sum_{j<r}y_j$, so $s_r\in S_r$ as well.  This represents
$b\in rS_r$.
\end{proof}

Taking $k_r=k_{S_r}$ proves all the assertions for finite $r\ge2$.
The same pair of mixing densities $1\pm c\cos\theta$ witnesses
nonidentifiability at every order $q<r$.
For $r=1$, use $S_1=\mathbb Z$.

\subsection{An analytic family with no finite replication threshold}

Enumerate the pairs $(Q,a)$ with $Q\ge1$ and $0\le a<Q$ in a sequence
$(Q_j,a_j)_{j\ge2}$ in which each pair appears infinitely often.
Set $n_1=1$ and recursively choose positive integers $n_j$ such that
\begin{equation}\label{p10:eq:lacunary-AP-construction}
 n_j\equiv a_j\pmod{Q_j},\qquad
 n_j>2^j\sum_{i<j}i n_i\quad(j\ge2).
\end{equation}
Each choice is possible in the prescribed infinite arithmetic
progression.  Set
\[
 S_\infty=\{0\}\cup
 \bigcup_{j\ge1}\{\pm n_j,\pm2n_j,\ldots,\pm jn_j\}.
\]
The positive blocks are disjoint, by the growth condition.
Since each $n_j$ is present, this set meets every residue class
infinitely often.  The finite blocks will also ensure determination of
finite mixing laws among unrestricted competitors.

We claim that $qS_\infty\ne\mathbb Z$ for each fixed $q\ge1$.
Choose $K\ge2$ large enough that $2^j>4q$ for every $j\ge K$ and
$n_K>4$.  Put $b=\lfloor n_K/2\rfloor$.  In any representation
\[
 z=s_1+\cdots+s_q=\sum_j c_j n_j,
 \qquad s_h\in S_\infty,
\]
combine all terms from the same block and cancel their integer
coefficients.  The resulting coefficients satisfy $|c_j|\le qj$.
If every nonzero
coefficient has index below $K$, then
$|z|\le q\sum_{j<K}j n_j<n_K/4<b$.
Otherwise let $J\ge K$ be the largest index with nonzero coefficient.
Then
\[
 |z|\ge n_J-q\sum_{j<J}j n_j>\frac34n_J
 \ge\frac34n_K>b.
\]
Thus $b\notin qS_\infty$.  Lemma~\ref{p10:lem:replication-sumset} gives
nonidentifiability at every finite order, and its cosine perturbations
give the asserted positive analytic witnesses (which may depend on $q$).

On the other hand, $1\in S_\infty$, so every integer belongs to
$qS_\infty$ for some finite $q$.  Equality of $M_q^\mu$ and $M_q^\nu$
for every $q$ therefore gives equality of all Fourier coefficients of
$\mu$ and $\nu$, and hence $\mu=\nu$.
This is also identification from the law of the entire conditionally
i.i.d. sequence.

Finally, for any of the constructed kernels, equality of two finite
mixtures is a finite linear dependence among the functions
$k_S(\cdot\mid\theta_j)$.  Lemma~\ref{p10:lem:AP-jets} with all derivative
orders zero shows that their signed atomic difference vanishes.
Corollary~\ref{p10:cor:finite-atomic-all-competitors} strengthens this to
identification among unrestricted competitors and completes the proof
of Theorem~\ref{p10:thm:every-replication-threshold}.

\begin{remark}[What the threshold measures]
The theorem concerns exact probability laws of grouped observations,
not the number of independently sampled groups needed for statistical
estimation.  In the infinite-threshold construction no single pair of
distinct mixing measures is indistinguishable at every order: each
finite order has a nontrivial kernel, while the intersection of all
these kernels is zero.  All finite-node derivative tests nevertheless
remain independent already at order one.
\end{remark}

\begin{corollary}[Finite support remains determined without a support bound]
\label{p10:cor:finite-atomic-all-competitors}
For each $r\in\mathbb N\cup\{\infty\}$, let $k_r$ be the kernel
constructed above.  If $\mu$
has finite support and $\nu$ is any Borel probability measure on
$\mathbb T$, then $M_1^\mu=M_1^\nu$ implies $\mu=\nu$.
\end{corollary}

\begin{proof}
The assertion is immediate for $r=1$.  Otherwise, write
$\operatorname{supp}\mu=\{\theta_1,\ldots,\theta_J\}$.
For finite $r\ge2$, choose a positive integer $n$ with
$\|n\alpha_r\|<\delta_r/J$; density of the irrational rotation
supplies one.  The triangle inequality gives
$\{-Jn,-(J-1)n,\ldots,Jn\}\subseteq S_r$.
For $r=\infty$, choose any block index $h\ge J$ and set $n=n_h$;
the same progression then lies in $S_\infty$ by construction.
The nonnegative trigonometric polynomial
\[
 p(\theta)=\prod_{j=1}^J
        |e^{in\theta}-e^{in\theta_j}|^2
\]
has its Fourier spectrum in that finite subset of $S_r$.
Consequently equality of the one-view laws gives
$\int p\,d\nu=\int p\,d\mu=0$.
Since $p\ge0$, the measure $\nu$ is supported on the zero set of $p$,
which consists of at most $nJ$ points of the circle.  The already proved
independence of the component densities then identifies the two finite
mixing measures.
\end{proof}

\begin{remark}[Uniform analytic bounds across the entire hierarchy]
The freedom in the threshold does not come from deteriorating smoothness
or positivity constants.  Every kernel above has the same positive
lower bound $1-1/(e-1)$.  On any complex set with
$|\operatorname{Im}x|+|\operatorname{Im}\theta|\le C$, each mixed
derivative of total order $d$ is bounded by the common finite majorant
\[
 \mathbf1_{\{d=0\}}+\sum_{n\ge1} n^d e^{-n^2+Cn},
\]
independently of the chosen finite or infinite threshold.
\end{remark}





\section{Localized invisible perturbations and exact moment coincidences}

To construct unrestricted thresholds on an ordinary interval, we need
smooth signed measures annihilating prescribed sparse powers. Ordinary
additive smoothing need not preserve those moments. Multiplicative
smoothing does preserve them, because moments factor under multiplication.
This also resolves the specified prescribed-moment conjecture and permits
an exact classification of coincidence sets.
\label{p11:sec:smooth-rigidity}

For $\sigma\geq1$ and an interval $J$, a smooth function belongs to
$G^\sigma(J)$ if, for some constants $C,B>0$,
\[
 \sup_{\theta\in J}|f^{(k)}(\theta)|
       \leq C B^k(k!)^\sigma\qquad(k\geq0).
\]
On a closed interval we require smoothness up to its endpoints.
For $J=\mathbb R$, the displayed bounds are global.
In particular $G^1([1,2])$ consists of the restrictions of functions real
analytic on a neighborhood of the interval. The constants may depend on the
function. The distinction between this union of classes with varying
constants and one fixed bounded class matters for statistical conclusions.

\begin{theorem}[Localized smooth M\"untz annihilators]\label{p11:thm:smooth-witness}
Let $0<a<b$, and let $T\subset\mathbb N_0$ contain $0$. The following
conditions are equivalent:
\begin{enumerate}
\item $\displaystyle\sum_{s\in T\setminus\{0\}}s^{-1}<\infty$.
\item For every nonempty open interval $I$ with
$\overline I\subset(a,b)$ and every $\sigma>1$, there exists a nonzero
real $h\in G^\sigma(\mathbb R)$ supported in $I$ such that
\begin{equation}\label{p11:eq:gevrey-invisible-moments}
 \int_a^b x^s h(x)\,dx=0\qquad(s\in T).
\end{equation}
\item There is a nonzero continuous real function $h$ supported in
$[a,b]$ satisfying \eqref{p11:eq:gevrey-invisible-moments}.
\end{enumerate}
For each $I$ and $\sigma$ in the second assertion, the space of such
annihilators is infinite dimensional. Consequently every strictly positive
$G^\sigma$ probability density $g_0$ on $[a,b]$ has distinct strictly
positive $G^\sigma$ probability perturbations $g_0\pm\varepsilon h$ with
the same $T$-moments, agreeing with $g_0$ outside $I$.
\end{theorem}

\begin{proof}
Assume the reciprocal sum converges, fix $I$, and choose a nondegenerate
closed interval $[c,d]\subset I$. The classical M\"untz theorem on an
interval away from zero says that
\[
 V=\operatorname{span}\{t^s:s\in T\}
\]
is not uniformly dense in $C([c,d];\mathbb R)$ (this is immediate if $T$
is finite); see, for example,
\cite[Theorem~3.2]{BorweinErdelyi1995}. Hahn--Banach and the Riesz
representation theorem therefore supply a nonzero finite real signed
measure $\nu$, supported in $[c,d]$, such that
\[
 \int t^s\,d\nu(t)=0\qquad(s\in T).
\]
This includes $\nu([c,d])=0$, because $0\in T$.

Choose $\delta\in(0,1)$ so small that
$[c(1-\delta),d(1+\delta)]\subset I$. By
Lemma~\ref{p11:lem:gevrey-bump}, there exists a nonnegative, nonzero
$G^\sigma$ function $\varphi$ supported in
$[1-\delta,1+\delta]$. Define its multiplicative smoothing of $\nu$ by
\begin{equation}\label{p11:eq:multiplicative-smoothing}
 h(x)=\int_{[c,d]}t^{-1}\varphi(x/t)\,d\nu(t),\qquad x\in\mathbb R.
\end{equation}
It is supported in $[c(1-\delta),d(1+\delta)]$, hence in $I$.
Differentiation under the finite measure gives
\[
 \|h^{(k)}\|_\infty
 \leq \|\nu\|_{\rm TV}c^{-k-1}\|\varphi^{(k)}\|_\infty
 \leq C' (B/c)^k(k!)^\sigma,
\]
so $h\in G^\sigma(\mathbb R)$.
For every real $\lambda\geq0$, Fubini and $x=tu$ yield the factorization
\begin{equation}\label{p11:eq:mellin-smoothing-factorization}
 \int_0^\infty x^\lambda h(x)\,dx
 =\left(\int_{[c,d]}t^\lambda\,d\nu(t)\right)
  \left(\int_0^\infty u^\lambda\varphi(u)\,du\right).
\end{equation}
Thus all moments in $T$ vanish.

It remains essential to check that smoothing has not killed $\nu$.
If every nonnegative integer moment of $\nu$ were zero, the Weierstrass
approximation theorem would imply $\nu=0$. Hence there is an integer
$m\geq0$ with $\int t^m\,d\nu(t)\ne0$. The second factor in
\eqref{p11:eq:mellin-smoothing-factorization} at $\lambda=m$ is strictly
positive, since $\varphi\geq0$ is nonzero and supported on positive real
numbers. Therefore $\int x^m h(x)\,dx\ne0$, and $h$ is nonzero.
This proves the second assertion.

For necessity, if the reciprocal sum diverges, the same M\"untz theorem
makes $V$ dense in $C[a,b]$. A continuous $h$ annihilating $V$ then
annihilates every continuous function; taking that function to be $h$
gives $\int_a^b h^2=0$, so $h=0$. The second assertion plainly implies
the third.

To obtain infinite dimensionality, apply the construction on infinitely
many pairwise disjoint nonempty open subintervals of $I$ with closures
inside $I$. The resulting nonzero functions have disjoint supports and
are linearly independent. Finally $0\in T$ gives $\int h=0$, and
boundedness of $h$ and the positive minimum of $g_0$ show that
$g_0\pm\varepsilon h$ have all the asserted properties for sufficiently
small $\varepsilon>0$.
\end{proof}

\begin{corollary}[The prescribed-moment conjecture]\label{p11:cor:bew-conjecture}
Let $0<a<b$ and let $(n_k)$ be a strictly increasing sequence of positive
integers with $\sum_k n_k^{-1}<\infty$. There exists a nonzero signed
measure with a smooth density supported strictly inside $(a,b)$ whose
$n_k$th moments vanish for every $k$. Its total mass can also be required
to vanish, its support can be localized in any prescribed interior
interval, and its density can be chosen in any prescribed Gevrey class
$G^\sigma$, $\sigma>1$.
\end{corollary}
\begin{proof}
Apply Theorem~\ref{p11:thm:smooth-witness} to $T=\{0,n_1,n_2,\ldots\}$.
\end{proof}

\begin{corollary}[Strictly log-concave moment twins]\label{p11:cor:logconcave-twins}
Under the summability assumption in Theorem~\ref{p11:thm:smooth-witness},
for every $\sigma>1$ there are distinct probability densities $g_+$ and
$g_-$ in $G^\sigma(\mathbb R)$, each supported on $[a,b]$, positive on
$(a,b)$ and strictly log-concave there, with
\[
 \int x^s g_+(x)\,dx=\int x^s g_-(x)\,dx\qquad(s\in T).
\]
Their difference can be supported in any prescribed interior interval.
Any finite list of additional nonnegative integer moments can also be
required to agree.
\end{corollary}
\begin{proof}
Put $p=(\sigma-1)^{-1}$ and normalize
\[
 g_0(x)=C\exp\bigl(-(x-a)^{-p}-(b-x)^{-p}\bigr),\quad a<x<b,
 \qquad g_0(x)=0\quad\text{otherwise}.
\]
The proof of Lemma~\ref{p11:lem:gevrey-bump} shows that
$g_0\in G^\sigma(\mathbb R)$. On $(a,b)$,
\[
 (\log g_0)''(x)
 =-p(p+1)\bigl((x-a)^{-p-2}+(b-x)^{-p-2}\bigr)<0.
\]
Choose the nonzero $h$ of Theorem~\ref{p11:thm:smooth-witness} with compact
support $K$ in the prescribed interval. On $K$, $g_0$ has a positive
minimum and $g_0g_0''-(g_0')^2$ has a strictly negative maximum.
Consequently, for sufficiently small $\varepsilon>0$,
$g_\pm=g_0\pm\varepsilon h$ remain positive on $K$ and satisfy
$g_\pm g_\pm''-(g_\pm')^2<0$ there. Outside $K$ they equal $g_0$.
Thus they are strictly log-concave on $(a,b)$, remain probability
densities, and have the required common moments. Adding finitely many
exponents to $T$ preserves the convergent reciprocal sum.
\end{proof}

This proves the mathematical assertion of Conjecture~17 in
B\"urgin--Epperlein--Wirth~\cite{BurginEpperleinWirth2022}, with stronger
smoothness and localization conclusions. The proof combines classical
M\"untz non-density with the elementary moment factorization
\eqref{p11:eq:mellin-smoothing-factorization}. A claim about historical
priority requires checking subsequent literature independently.

\begin{lemma}[Compact Gevrey functions]\label{p11:lem:gevrey-bump}
For every $\sigma>1$ and $\delta>0$ there is a nonnegative, nonzero
$\varphi\in G^\sigma(\mathbb R)$ supported in $[-\delta,\delta]$.
\end{lemma}
\begin{proof}
Put $p=(\sigma-1)^{-1}$ and
$f(t)=e^{-t^{-p}}$ for $t>0$, $f(t)=0$ for $t\leq0$.
Choose $\epsilon_p>0$ small enough that, on the complex disk
$|z-t|\leq\epsilon_p t$ about any positive real $t$, the principal branch
satisfies $\operatorname{Re}(z^{-p})\geq c_p t^{-p}$ for a constant
$c_p>0$. Cauchy's estimate yields
\[
 |f^{(k)}(t)|\leq k!(\epsilon_p t)^{-k}e^{-c_pt^{-p}}.
\]
For $k\geq1$, maximizing $t^{-k}e^{-c_pt^{-p}}$ and using
$(k/e)^k\leq k!$ shows
\[
 \sup_{t>0}|f^{(k)}(t)|\leq B_p^k(k!)^{1+1/p}.
\]
For each fixed $k$ the preceding pointwise estimate tends to zero as
$t\downarrow0$, so extension by zero is smooth. Now take
$\varphi(t)=f(t+\delta)f(\delta-t)$. It is positive on
$(-\delta,\delta)$, supported in the required interval, and Leibniz's
rule gives the claimed Gevrey estimates.
\end{proof}

\begin{theorem}[Endpoint rigidity]\label{p11:thm:analytic-rigidity}
Let $E\subset[0,\infty)$ be unbounded. If $h\in C^\infty([1,2])$
satisfies
\[
 \int_1^2\theta^s h(\theta)\,d\theta=0\qquad(s\in E),
\]
then $h^{(k)}(2)=0$ for every $k\geq0$. In particular, if $h$ is real
analytic on a neighborhood of $[1,2]$, it is identically zero.
The same conclusion holds in any linear smooth class on a neighborhood
of $[1,2]$ for which the Taylor jet at $2$ determines the function.
\end{theorem}

\begin{proof}
If the endpoint jet is not zero, let $m$ be its least nonzero order.
Taylor's formula and the substitution $t=s(1-\theta/2)$ give
\begin{equation}\label{p11:eq:endpoint-moment-asymptotic}
 \int_1^2\theta^s h(\theta)\,d\theta
  \sim (-1)^m h^{(m)}(2)\frac{2^{s+m+1}}{s^{m+1}}
       \qquad(s\to+\infty).
\end{equation}
Here is a direct justification. On the whole interval,
$|h(\theta)|\leq C(2-\theta)^m$, since all lower endpoint derivatives
vanish. After multiplying the integral by $s^{m+1}2^{-s-m-1}$,
the transformed integrand converges pointwise to
$(-1)^m h^{(m)}(2)t^m e^{-t}/m!$ and is dominated by $Ct^me^{-t}$;
we used $(1-t/s)^s\leq e^{-t}$ on the transformed domain $0\leq t\leq s/2$.
Dominated convergence proves \eqref{p11:eq:endpoint-moment-asymptotic}.
Its nonzero leading term excludes zeros on an unbounded set of positive
exponents. Thus the endpoint jet vanishes. Analyticity, or the stated
jet uniqueness property, now implies $h=0$.
\end{proof}

\begin{corollary}[A sharp regularity distinction]\label{p11:cor:regularity-threshold}
Suppose a grouped-observation model determines exactly the moments with
exponents in $qS$ at replication order $q$, where $S\subset\mathbb N_0$
is unbounded and contains $0$. If $\sum_{s\in qS\setminus\{0\}}s^{-1}<\infty$, then its $q$-observation law is not
injective even on strictly positive $G^\sigma$ densities, for any
$\sigma>1$, and the ambiguity can be localized in any interior interval.
Nevertheless its one-observation law is injective on the class of
real-analytic mixing densities.
\end{corollary}
\begin{proof}
Apply Theorem~\ref{p11:thm:smooth-witness} with $T=qS$ for the smooth
ambiguity. For two real-analytic densities with equal one-observation
laws, their difference has vanishing moments on the unbounded set $S$;
Theorem~\ref{p11:thm:analytic-rigidity} forces that difference to vanish.
\end{proof}

The last injectivity statement compares two analytic densities. It does
not say that an analytic mixing density is identifiable among arbitrary
Borel competitors: even the uniform density has the invisible smooth
perturbations supplied by Theorem~\ref{p11:thm:smooth-witness}. Likewise the
endpoint rigidity statement permits nonzero smooth annihilators, but
forces every such annihilator to have zero Taylor jet at the right
endpoint. These statements concern the regularity of the unknown mixing
law, whereas the observation kernels in the following constructions are
entire throughout.


\subsection{Exact prescribed coincidence sets}
\label{p11:sec:exact-moment-pattern}

The preceding existence theorem can be strengthened: one can exclude all
unrequested coincidences. This refinement uses the classical
Clarkson--Erd\H{o}s description of the closure of a non-dense integer
M\"untz space, rather than only its non-density.

\begin{theorem}[Exact moment-zero patterns]\label{p11:thm:exact-moment-pattern}
Let $T\subset\mathbb N_0$ contain $0$ and satisfy
$\sum_{s\in T\setminus\{0\}}s^{-1}<\infty$. Fix $0<a<b$, an open
interval $I$ with $\overline I\subset(a,b)$, and $\sigma>1$.
There is a real function $h\in G^\sigma(\mathbb R)$ supported in $I$
such that
\[
 \int_a^b x^m h(x)\,dx=0\quad\Longleftrightarrow\quad m\in T
       \qquad(m\in\mathbb N_0).
\]
Consequently the moment coincidence set of two distinct strictly
log-concave, compactly supported $G^\sigma$ probability densities can
be prescribed to be exactly any such set $T$.
\end{theorem}

\begin{proof}
Choose a nondegenerate interval $[c,d]\subset I$ and let
$V=\overline{\operatorname{span}\{x^s:s\in T\}}^{\,C[c,d]}$.
For every $m\notin T$, one has $x^m\notin V$. For finite $T$ this follows
from independence of distinct monomials. For infinite $T$, the
Clarkson--Erd\H{o}s theorem, in the form of
\cite[Theorem~3.1]{BorweinErdelyi1995}, states that every element of $V$
has a representation
\[
 f(x)=\sum_{s\in T}a_sx^s\qquad(c\leq x<d).
\]
The positive gap condition required by that theorem holds because the
exponents are distinct integers. Convergence at each $x<d$ implies that
this is a power series of radius at least $d$. If $f(x)=x^m$ on
$[c,d]$, the identity theorem and uniqueness of power-series coefficients
would give the coefficient $1$ at degree $m$, impossible when $m\notin T$.

Enumerate $\mathbb N_0\setminus T=\{m_1,m_2,\ldots\}$.
Hahn--Banach and Riesz now give real signed measures $\nu_j$, supported
on $[c,d]$, of total variation at most $1$, satisfying
\[
 \int x^s\,d\nu_j(x)=0\quad(s\in T),
 \qquad \int x^{m_j}\,d\nu_j(x)\ne0.
\]
Fix one nonnegative, nonzero Gevrey function $\varphi$ supported in a
sufficiently small interval about $1$, and apply
\eqref{p11:eq:multiplicative-smoothing} to each $\nu_j$. The resulting
functions $h_j$ have supports in one common compact subset of $I$ and
satisfy common estimates
\[
 \|h_j^{(k)}\|_\infty\leq C B^k(k!)^\sigma
 \quad\text{for all }j,k.
\]
They annihilate every moment in $T$, while the $m_j$th moment of $h_j$
is nonzero by \eqref{p11:eq:mellin-smoothing-factorization}.

For $0<t<1$, set
\[
 h_t(x)=\sum_{j=1}^\infty 2^{-j}t^j h_j(x).
\]
All derivative series converge uniformly, with common Gevrey bounds,
and the support stays in $I$. For each $m\notin T$, the scalar function
\[
 H_m(t)=\int_a^b x^m h_t(x)\,dx
\]
is analytic on $|t|<2$ and is not identically zero: if $m=m_j$, its
coefficient of $t^j$ is the nonzero number
$2^{-j}\int x^{m_j}h_j(x)\,dx$. Its real zeros in $(0,1)$ therefore
form a countable set. Avoiding the union of these zero sets over the
countably many $m\notin T$ gives a parameter $t_*$ for which $h_{t_*}$
has exactly the asserted vanishing moments. In particular it is nonzero.
The final statement follows by using this $h_{t_*}$ in the compact
strictly log-concave baseline perturbation of
Corollary~\ref{p11:cor:logconcave-twins}; multiplication by a nonzero constant
does not change its moment-zero set.
\end{proof}

\begin{corollary}[Classification of exact coincidence sets]
\label{p11:cor:coincidence-classification}
For a set $T\subset\mathbb N_0$, the following are equivalent:
\begin{enumerate}
\item $0\in T$ and $\sum_{s\in T\setminus\{0\}}s^{-1}<\infty$.
\item For every $0<a<b$ and every $\sigma>1$, there exist distinct
probability densities $g_+,g_-\in G^\sigma(\mathbb R)$, each with
support $[a,b]$ and strictly log-concave on $(a,b)$, whose integer moment
coincidence set is exactly $T$.
\item There exist distinct probability measures, both supported in
some compact interval contained in $(0,\infty)$, whose integer moment
coincidence set is exactly $T$.
\end{enumerate}
In the second assertion the difference of the two densities can be
supported in any prescribed interior interval.
\end{corollary}
\begin{proof}
Theorem~\ref{p11:thm:exact-moment-pattern} gives the first implication, and
the second is immediate. For the converse, probability normalization
forces $0\in T$. If the reciprocal sum diverged, the M\"untz theorem
would force any two compactly supported positive-interval measures
agreeing on those moments to coincide. This contradicts distinctness.
\end{proof}



\subsection{Examples and the constructive scope of the classification}

The exact coincidence classification is stronger than the existence
of two distributions sharing some infinite collection of moments.
For example, there are two distinct strictly log-concave smooth
densities on any chosen positive interval that agree at precisely
the orders $0,1,4,9,16,\ldots$ and disagree at every other integer order.
In contrast, distinct compactly supported laws cannot agree at all
prime orders, since the sum of reciprocal primes diverges.
The first example follows by taking $T=\{0\}\cup\{n^2:n\geq1\}$;
the second is the classical completeness direction.
The proof of the first is an existence construction, using
Hahn--Banach and avoidance of countably many analytic zero sets;
it is not a numerical algorithm for evaluating the densities.


\section{Arbitrary replication thresholds on an ordinary interval}
\label{p11:sec:interval-core}

Throughout this section the observation space is $[0,1]$ with Lebesgue
measure and the parameter space is $[1,2]$. The observation polynomials
$L_s$ retain the normalization in \eqref{p09:eq:legendre-orthogonality};
the parameter functions will now be monomials. The endpoints of the
parameter interval are distinct, and no periodic identification is used.
For an infinite set $S\subseteq\mathbb Z_{\geq0}$ containing $0$, define
\begin{equation}\label{p11:eq:interval-kernel}
 a_s=2^{-s^2-s-2},\qquad
 k_S(x\mid\theta)=1+\sum_{s\in S\setminus\{0\}}
               a_sL_s(x)\theta^s.
\end{equation}
For a Borel probability measure $\mu$ on $[1,2]$, let
\[
 A_{q,S}\mu(dx_1\cdots dx_q)
 =\left\{\int_{1}^{2}\prod_{j=1}^q k_S(x_j\mid\theta)
                      \,d\mu(\theta)\right\}dx_1\cdots dx_q.
\]
The same latent parameter is used for the $q$ observations in a group.
Set $qS=\{s_1+\cdots+s_q:s_i\in S\}$.  The inclusion $0\in S$ means that
moments accessible with fewer observations remain accessible with more.

\begin{lemma}[Regularity and the exact observable moments]
\label{p11:lem:interval-observability}
The function in \eqref{p11:eq:interval-kernel} is jointly entire in
$(x,\theta)\in\mathbb C^2$, integrates to one in $x$, and satisfies
\[
 \frac34<k_S(x\mid\theta)<\frac54
 \quad (x\in[0,1],\ \theta\in[1,2]).
\]
On every fixed complex compact set, every fixed mixed derivative is
bounded uniformly over the choice of $S$.  For any two finite signed
Borel measures $\mu,\nu$ on $[1,2]$,
\begin{equation}\label{p11:eq:observable-sumset}
 A_{q,S}\mu=A_{q,S}\nu
 \quad\Longleftrightarrow\quad
 \int\theta^t\,d\mu(\theta)=\int\theta^t\,d\nu(\theta)
 \quad\hbox{for every }t\in qS.
\end{equation}
\end{lemma}

\begin{proof}
On the real rectangle, the sum of the absolute values of the
nonconstant terms is bounded by
\[
 \sum_{s\geq1}2^{-s^2-s-2}2^s
 =\frac14\sum_{s\geq1}2^{-s^2}
 <\frac14\sum_{s\geq1}2^{-s}=\frac14.
\]
Orthogonality to $L_0$ gives normalization.  On each compact subset of
the complex $x$-plane, $|L_s(x)|\leq C^s$ for some $C\geq1$: for example,
this follows immediately by induction from the three-term Legendre
recurrence.  The Gaussian decrease in $s$ in $a_s$ therefore gives
normal convergence on every complex compact rectangle.  The same
majorant for all $S$, followed by Cauchy estimates on a slightly larger
rectangle, bounds every mixed derivative uniformly in $S$.

For $n\in S\setminus\{0\}$, the conditional mean of $L_n(X)$ is
$a_n\theta^n/(2n+1)$; for $n=0$ it is one.  Consequently, for
$n_1,\ldots,n_q\in S$,
\begin{equation}\label{p11:eq:interval-product-statistic}
 \int\prod_{j=1}^q L_{n_j}(x_j)\,d(A_{q,S}\mu)
 =\left\{\prod_{j:n_j>0}\frac{a_{n_j}}{2n_j+1}\right\}
   \int\theta^{n_1+\cdots+n_q}\,d\mu(\theta).
\end{equation}
This proves the forward implication in \eqref{p11:eq:observable-sumset}.
For the reverse implication, expand the product of the $q$ kernel
series.  It converges absolutely and uniformly on the real rectangle,
and every resulting coefficient depends on a moment whose exponent
belongs to $qS$.  Termwise integration against a finite signed measure
is therefore valid and gives equality of the two densities.
\end{proof}

\begin{proposition}[Exact moment criterion for replication]
\label{p11:prop:interval-muntz}
For an arbitrary infinite $S\subseteq\mathbb Z_{\geq0}$ containing
$0$, the operator $A_{q,S}$ is injective on all Borel probabilities
on $[1,2]$ if and only if
\begin{equation}\label{p11:eq:replication-muntz}
 \sum_{t\in qS\setminus\{0\}}\frac1t=\infty.
\end{equation}
If the sum converges, noninjectivity holds even for strictly positive
Gevrey densities of any prescribed order $\sigma>1$, with their
difference supported in an arbitrarily small interior interval.
\end{proposition}
\begin{proof}
When the sum diverges, the M\"untz theorem on a positive interval says
that the real span of $\{\theta^t:t\in qS\}$ is uniformly dense in
$C[1,2]$.  Thus \eqref{p11:eq:observable-sumset} implies equality of the
measures.  This use of M\"untz's theorem is classical; the formulation
needed here is also stated as
\cite[Theorem~2]{BurginEpperleinWirth2022}.
When the sum converges, Theorem~\ref{p11:thm:smooth-witness} supplies the
stated positive density perturbations, and
\eqref{p11:eq:observable-sumset} makes their grouped laws identical.
\end{proof}

\begin{lemma}[All finite parameter jets are independent]
\label{p11:lem:interval-jets}
For any distinct $\theta_1,\ldots,\theta_J\in[1,2]$ and any nonnegative
integers $m_1,\ldots,m_J$, the functions
\[
 \left\{\partial_\theta^\ell k_S(\,\cdot\mid\theta_j):
        1\leq j\leq J,\ 0\leq\ell\leq m_j\right\}
\]
are linearly independent over $\mathbb C$.
\end{lemma}

\begin{proof}
Suppose their linear combination with coefficients $c_{j\ell}$ is
zero.  Taking its $L_s$ coefficient gives, for every positive $s\in S$,
\[
 0=\sum_{j=1}^J\theta_j^s B_j(s),\qquad
 B_j(s)=\sum_{\ell=0}^{m_j}c_{j\ell}\theta_j^{-\ell}(s)_\ell,
\]
where $(s)_\ell=s(s-1)\cdots(s-\ell+1)$ and $(s)_0=1$.
Order the nodes so that $\theta_1<\cdots<\theta_J$.  If $B_J$ is a
nonzero polynomial, division by $\theta_J^s$ would give
\[
 |B_J(s)|\leq C(1+s)^M\rho^s
 \quad (s\in S,\ s\longrightarrow\infty)
\]
for some $\rho<1$, whereas a nonzero polynomial is bounded below by a
positive constant times a nonnegative power of $s$ for all large
positive $s$.  This is impossible.  Repeating the argument gives
$B_j=0$ for all $j$.  The falling factorial polynomials are linearly
independent, so all $c_{j\ell}=0$.
\end{proof}

\begin{lemma}[Sparse moments determine every finite atomic measure]
\label{p11:lem:interval-atomic}
If $\mu$ is a finitely supported probability measure on $[1,2]$ and
$\nu$ is any Borel probability measure on $[1,2]$, then
\[
 A_{1,S}\mu=A_{1,S}\nu\quad\Longrightarrow\quad\mu=\nu.
\]
\end{lemma}

\begin{proof}
We first recall an elementary zero-counting fact, and include its
proof.  A nonzero real linear combination of $m+1$ distinct real
powers of $t$ has at most $m$ zeros on $(0,\infty)$, counted with
multiplicity.  Divide by the smallest power.  Differentiation leaves
at most $m$ powers, and induction together with Rolle's theorem proves
the assertion.  More explicitly, any finite collection of zeros of
total multiplicity $N$ gives at least $N-1$ zeros of the derivative,
counted with multiplicity; division by a nonvanishing power does not
alter the original zeros.

Let the support of $\mu$ be $\{\theta_1,\ldots,\theta_J\}$.  Choose
$2J+1$ exponents $0=s_0<s_1<\cdots<s_{2J}$ from $S$.  The $2J$
homogeneous real linear equations
\[
 p(\theta_j)=p'(\theta_j)=0\quad(1\leq j\leq J),\qquad
 p(t)=\sum_{i=0}^{2J}b_i t^{s_i},
\]
have a nonzero solution.  The zero-counting fact implies that these
are all the positive zeros of $p$ and that each has multiplicity
exactly two.  Thus, after multiplying $p$ by $-1$ if necessary,
$p\geq0$ on $(0,\infty)$ and its zero set there is precisely
$\{\theta_1,\ldots,\theta_J\}$.  This argument also applies when a
node is an endpoint of $[1,2]$, since both endpoints are interior
points of $(0,\infty)$.

Equality of the one-observation laws gives equality of all the
$S$-moments.  Hence $\int p\,d\nu=\int p\,d\mu=0$.  Positivity of
$\nu$ forces its support to lie in the same finite set of nodes.
Lemma~\ref{p11:lem:interval-jets}, with derivative order zero, now gives
equality of the weights, and hence $\mu=\nu$.
\end{proof}

\begin{theorem}[Every replication threshold on an interval]
\label{p11:thm:interval-threshold}
For every $r\in\{1,2,\ldots\}\cup\{\infty\}$ there is an infinite
set $S_r\subseteq\mathbb Z_{\geq0}$ containing $0$ and $1$ such that
the kernel \eqref{p11:eq:interval-kernel} has the following properties.
\begin{enumerate}
\item It is strictly positive and jointly entire, with the uniform
      bounds in Lemma~\ref{p11:lem:interval-observability}.
\item Every finite family of parameter jets is linearly independent,
      and every finitely supported mixing probability is identified
      by one observation even among arbitrary Borel competitors.
\item If $r<\infty$, then $A_{q,S_r}$ is injective on all Borel
      probabilities exactly when $q\geq r$.
\item If $r=\infty$, no finite $q$ gives injectivity, but the
      collection $(A_{q,S_r}\mu)_{q\geq1}$ determines $\mu$.
\item Whenever injectivity fails, it already fails for distinct
      strictly positive $C^\infty$ mixing densities.  For finite
      $r\geq2$, one pair can witness failure simultaneously for
      every $q<r$.
\end{enumerate}
The smooth densities in the last assertion may, more strongly, be
chosen in any prescribed Gevrey class of order $\sigma>1$.
\end{theorem}

\begin{proof}
For $r=1$, take $S_1=\mathbb Z_{\geq0}$.  All monomial moments are
observed, so polynomial approximation on $[1,2]$ proves injectivity.

Fix $2\leq r<\infty$.  For $j=0,\ldots,r-1$, let
\[
 D_{r,j}=\left\{\sum_{k\geq0}\varepsilon_k2^{rk+j}:
       \varepsilon_k\in\{0,1\},\ \varepsilon_k=0
       \text{ for all but finitely many }k\right\},\qquad
 S_r=\bigcup_{j=0}^{r-1}D_{r,j}.
\]
Separating the binary digits of an arbitrary nonnegative integer
according to their positions modulo $r$ proves
\begin{equation}\label{p11:eq:digit-basis}
 rS_r=\mathbb Z_{\geq0}.
\end{equation}
Since $0\in S_r$, the same is true with $r$ replaced by any $q\geq r$.
Lemma~\ref{p11:lem:interval-observability} and polynomial approximation
therefore prove injectivity for $q\geq r$.

For $t\geq1$, the number of free binary digits in each $D_{r,j}$
below $t$ is at most $1+(\log_2t)/r$.  Consequently
\[
 \#(S_r\cap[0,t])\leq 2r\,t^{1/r},\qquad
 \#(qS_r\cap[0,t])\leq(2r)^q t^{q/r}.
\]
The second inequality holds because all summands are nonnegative,
so that a sum not exceeding $t$ uses only elements not exceeding
$t$.  If $q<r$, the exponent $q/r$ is strictly below one.  More generally,
if $T\subseteq\mathbb Z_{\geq0}$ satisfies
$\#(T\cap[0,t])\leq C t^\beta$ for $t\geq1$ and $\beta<1$, then
\[
 \sum_{t\in T\setminus\{0\}}\frac1t
 \leq\sum_{k\geq0}2^{-k}\#(T\cap[0,2^{k+1}])
 \leq C2^\beta\sum_{k\geq0}2^{-k(1-\beta)}<\infty.
\]
Apply
Theorem~\ref{p11:thm:smooth-witness} with $T=(r-1)S_r$: it provides
distinct strictly positive smooth densities with the same moments
at every exponent in $T$.  Since $qS_r\subseteq T$ for every $q<r$,
the same pair has identical $q$-observation laws for all such $q$.

For the infinite threshold take
\[
 S_\infty=\{0,1\}\cup\{2^{2^j}:j=0,1,2,\ldots\}.
\]
For each fixed $q$, the counting function of $qS_\infty$ is bounded
by a constant times $(1+\log\log(t+3))^q$.  In particular it is
$O_q(t^{1/2})$ as $t\longrightarrow\infty$.  Applying
Theorem~\ref{p11:thm:smooth-witness} to $T=qS_\infty$ proves failure of
injectivity for every fixed $q$.  The witnessing pair may depend
on $q$.  On the other hand, $1\in S_\infty$ implies
$\bigcup_{q\geq1}qS_\infty=\mathbb Z_{\geq0}$; equality of all grouped
laws thus gives equality of all moments and hence of the measures.
The regularity, finite-jet and finite-atomic assertions follow from
the preceding lemmas.
\end{proof}

\begin{corollary}[Localized infinite-dimensional ambiguity]
\label{p11:cor:interval-local-ambiguity}
Fix one of the models in Theorem~\ref{p11:thm:interval-threshold}, a
subcritical replication count $q<r$ (any finite $q$ when $r=\infty$),
and a nonempty open interval $I$ with $\overline I\subset(1,2)$.
For every $\sigma>1$ the nullspace of $A_{q,S_r}$ contains an
infinite-dimensional real vector space of signed densities in
$G^\sigma(\mathbb R)$ supported in $I$.
If $g_0$ is any strictly positive continuous probability density on
$[1,2]$, then, for every finite collection of such null densities
$h_1,\ldots,h_m$, all densities
\[
 g_0+\sum_{j=1}^m c_jh_j
\]
with $\sum_j|c_j|$ sufficiently small are strictly positive probability
densities with the same $q$-observation law as $g_0$.  If $g_0$ belongs
to $G^\sigma([1,2])$, then so do these perturbations.
\end{corollary}
\begin{proof}
The reciprocal sum over $qS_r\setminus\{0\}$ converges by the preceding
proof.  Choose countably many disjoint open subintervals whose closures
lie in $I$.  Apply Theorem~\ref{p11:thm:smooth-witness} separately in each
subinterval.  The resulting nonzero compactly supported functions are
linearly independent because their supports are disjoint, and
\eqref{p11:eq:observable-sumset} puts them in the nullspace.  Each has
integral zero.  For a finite collection, boundedness of the $h_j$ and
the positive minimum of $g_0$ on $[1,2]$ prove the remaining assertions.
\end{proof}

\begin{remark}[What produces the threshold]
The mechanism is an additive covering problem for the degrees of
observable monomials.  Splitting binary digits makes $r$ copies cover
every degree, whereas fewer copies have a counting exponent strictly
below one.  Positivity and analyticity of the observation kernel,
and independence of every finite collection of parameter jets, are
unchanged as this global threshold varies.  In particular, this
construction takes place on an ordinary compact interval and uses
neither an identification of endpoints nor a periodic parameter.
\end{remark}




\section{Analytic identifiability and a sharp boundary for uniform learning}
\label{p11:sec:analytic-minimax}

Identifiability within a regularity class need not imply uniform
statistical recovery on that class.  Here this distinction survives
every prescribed Gevrey bound of order strictly greater than one,
even if each density under consideration extends to an entire function.
The corresponding fixed analytic bounds behave differently.

We first state the compactness principle used to separate these issues.
For a compact metric parameter space $\Theta$, let $W_1$ denote the
Wasserstein distance for its metric.  Suppose $X$ is also a compact
metric space and $\theta\mapsto P_\theta\in\mathcal P(X)$ is continuous
in total variation.  Fix a positive integer $q$ and put
\[
 A_q\mu=\int P_\theta^{\otimes q}\,d\mu(\theta).
\]
For a nonempty class $\mathcal D\subseteq\mathcal P(\Theta)$, define
\begin{equation}\label{p11:eq:class-minimax-risk}
 R_{N,q}(\mathcal D)=\inf_{\widehat\mu}
       \sup_{\mu\in\mathcal D}
       \mathbb E_{(A_q\mu)^{\otimes N}}
                    W_1(\widehat\mu,\mu).
\end{equation}
There are $N$ independent groups, each containing $q$ observations;
the estimator may take values in all of $\mathcal P(\Theta)$.
We use $\operatorname{TV}(P,Q)=\sup_B|P(B)-Q(B)|$.

\begin{proposition}[Uniform learning depends on the closure]
\label{p11:prop:closure-learning}
Let $\overline{\mathcal D}$ denote weak closure in
$\mathcal P(\Theta)$.  Then
\begin{equation}\label{p11:eq:closure-criterion}
 R_{N,q}(\mathcal D)\longrightarrow0
 \quad\Longleftrightarrow\quad
 A_q\text{ is injective on }\overline{\mathcal D}.
\end{equation}
More precisely, for every $N\geq1$,
\begin{equation}\label{p11:eq:closure-fiber-lower-bound}
 R_{N,q}(\mathcal D)\geq\frac12
 \sup\left\{W_1(\mu,\nu):
     \mu,\nu\in\overline{\mathcal D},\ A_q\mu=A_q\nu\right\}.
\end{equation}
\end{proposition}

\begin{proof}
First, $A_q$ is continuous from the weak topology on mixing measures
to total variation on observation laws.  Indeed,
$\operatorname{TV}(P^{\otimes q},Q^{\otimes q})\leq
q\operatorname{TV}(P,Q)$, so the map
$\theta\mapsto P_\theta^{\otimes q}$ is norm continuous as a map into
the Banach space of finite signed measures.  A finite partition of
unity on $\Theta$ approximates this map uniformly by maps of the form
$\sum_j f_j(\theta)Q_j$, with each $f_j$ continuous and each $Q_j$ fixed.
Integration of each approximant is weak-to-norm continuous; uniform
approximation proves the assertion.

For arbitrary parameters $\mu,\nu$, arbitrary data laws $P,Q$, and
an arbitrary estimator, the triangle inequality gives
\begin{equation}\label{p11:eq:overlap-risk}
 \max\{\mathbb E_P W_1(\widehat\mu,\mu),
         \mathbb E_Q W_1(\widehat\mu,\nu)\}
 \geq\frac{W_1(\mu,\nu)}{2}
                         (1-\operatorname{TV}(P,Q)).
\end{equation}
To verify the factor $1/2$, dominate $P,Q$ by $P+Q$ and integrate
the pointwise inequality
$W_1(\widehat\mu,\mu)+W_1(\widehat\mu,\nu)\geq W_1(\mu,\nu)$
against the smaller of their two densities; this common measure has
mass $1-\operatorname{TV}(P,Q)$.

Now let $\mu,\nu\in\overline{\mathcal D}$ satisfy $A_q\mu=A_q\nu$.
Choose $\mu_j,\nu_j\in\mathcal D$ converging weakly to them.
Then $W_1(\mu_j,\nu_j)\to W_1(\mu,\nu)$ and, for each fixed $N$,
\[
 \operatorname{TV}\big((A_q\mu_j)^{\otimes N},
                       (A_q\nu_j)^{\otimes N}\big)
 \leq N\operatorname{TV}(A_q\mu_j,A_q\nu_j)\longrightarrow0.
\]
Apply \eqref{p11:eq:overlap-risk} and let $j\to\infty$.
The result holds for every estimator, proving
\eqref{p11:eq:closure-fiber-lower-bound}.

For the converse put $K=\overline{\mathcal D}$, which is compact.
Let $d_{\rm BL}$ be the bounded Lipschitz metric on
$\mathcal P(X^q)$.  If $A_q$ is injective on $K$, its continuous
inverse on the compact image is uniformly continuous.  Equivalently,
\[
 \omega(t)=\sup\{W_1(\mu,\nu):\mu,\nu\in K,
                   d_{\rm BL}(A_q\mu,A_q\nu)\leq t\}
 \longrightarrow0\qquad(t\downarrow0).
\]
Choose finite $1/N$-nets $F_N\subseteq\mathcal D$ for $K$ in $W_1$.
If $\widehat Q_N$ is the empirical law of the $N$ groups, choose
$\widehat\mu_N\in F_N$ minimizing
$d_{\rm BL}(A_q\widehat\mu_N,\widehat Q_N)$, breaking ties by a fixed
ordering.  This is measurable.  Uniform continuity of $A_q$ gives
numbers $b_N\to0$ such that, for every $\mu\in\mathcal D$,
\[
 d_{\rm BL}(A_q\widehat\mu_N,A_q\mu)
 \leq 2d_{\rm BL}(\widehat Q_N,A_q\mu)+b_N.
\]
On a compact sample space,
\begin{equation}\label{p11:eq:uniform-empirical-bl}
 \sup_{Q\in\mathcal P(X^q)}
       \mathbb E_Q d_{\rm BL}(\widehat Q_N,Q)\longrightarrow0.
\end{equation}
Indeed, the bounded Lipschitz unit ball has a finite uniform
$\eta$-net by Arzel\`a--Ascoli.  For fixed $\eta$, the expected maximum
of the finitely many empirical integral errors tends to zero, while
replacing an arbitrary function by a net element costs at most
$2\eta$.  Let $N\to\infty$ and then $\eta\downarrow0$.
Uniform inverse continuity, \eqref{p11:eq:uniform-empirical-bl}, and the
bounded diameter of $\Theta$ imply
$\sup_{\mu\in\mathcal D}\mathbb E W_1(\widehat\mu_N,\mu)\to0$.
\end{proof}

The lower-bound mechanism is a standard consequence of indistinguishable
limits and two-point decision theory.  In particular, the use of
nonidentifiable limits to prove constant minimax risk over finite
discrete mixing laws appears in Proposition~3 of
Le, Bariletto, Rinaldo, and Ho~\cite{LeBarilettoRinaldoHo2026}.  The purpose of the next theorem is the
regularity-restricted consequence and its exact replication threshold,
not a claim of novelty for the compactness principle itself.

For $\sigma\geq1$, $M>1$, and $B>0$, let
$\mathcal E_{\sigma,M,B}$ consist of the probability measures
$g(\theta)\,d\theta$ on $[1,2]$ such that $g$ extends to an entire
function of one complex variable, is real on $[1,2]$, and satisfies
\begin{equation}\label{p11:eq:entire-gevrey-ball}
 \frac12\leq g(\theta)\leq\frac32,\qquad
 \|g^{(k)}\|_{L^\infty([1,2])}
        \leq MB^k(k!)^\sigma\quad(k=0,1,\ldots).
\end{equation}
Thus the derivative constants are fixed for the whole class, and
every one of its members is entire, whatever the value of $\sigma$.
The uniform density belongs to every such class.

\begin{theorem}[A sharp boundary inside classes of entire densities]
\label{p11:thm:entire-gevrey-learning}
Consider the interval kernel $k_{S_r}$ of
Theorem~\ref{p11:thm:interval-threshold}, where
$r\in\{1,2,\ldots\}\cup\{\infty\}$.
For every $\sigma\geq1$, $M>1$, and $B>0$,
$A_{1,S_r}$ is injective on $\mathcal E_{\sigma,M,B}$.
Nevertheless uniform estimation has the following sharp distinctions.
\begin{enumerate}
\item For $\sigma=1$, every fixed $q\geq1$ satisfies
\[
 R_{N,q}(\mathcal E_{1,M,B})\longrightarrow0.
\]
\item For $\sigma>1$ and every finite $q<r$, there exists
$c=c(r,q,\sigma,M,B)>0$ such that
\[
 R_{N,q}(\mathcal E_{\sigma,M,B})\geq c
 \qquad\text{for every }N\geq1.
\]
If $2\leq r<\infty$, the same constant $c>0$ may be chosen for all
$q<r$ simultaneously.  If $r=\infty$, the constant may depend on $q$.
\item If $r<\infty$ and $q\geq r$, then
\[
 R_{N,q}(\mathcal P([1,2]))\longrightarrow0,
\]
and in particular uniform consistency holds on every
$\mathcal E_{\sigma,M,B}$.
\end{enumerate}
The assertion for $\sigma>1$ holds for every prescribed $M>1$ and
$B>0$, not merely for some sufficiently large derivative constants.
\end{theorem}

\begin{proof}
Any two elements of $\mathcal E_{\sigma,M,B}$ have analytic density
difference.  Their one-observation laws determine its moments on the
unbounded set $S_r$.  Theorem~\ref{p11:thm:analytic-rigidity} therefore
proves the stated injectivity within the class.

For the first assertion, the weak closure of $\mathcal E_{1,M,B}$
consists of measures having densities with the same derivative bounds.
To see this, any sequence of these densities has, by the bounds on
successive derivatives and diagonal Arzel\`a--Ascoli, a subsequence
converging in $C^k([1,2])$ for every fixed $k$.  The resulting smooth
density $g$ satisfies $\|g^{(k)}\|_\infty\leq MB^k k!$.
Taylor's remainder estimate shows that $g$ is real analytic, with
local radius at least $1/B$; the one-sided Taylor series at the two
endpoints extend it analytically to a neighborhood of $[1,2]$.
Although such a limit need not be entire, analytic moment rigidity
still applies.  Thus $A_{1,S_r}$ is injective on this weak closure,
and so is every $A_{q,S_r}$.  Proposition~\ref{p11:prop:closure-learning}
proves uniform consistency.

For the second assertion fix $\sigma>1$, and choose
$1<\tau<\sigma$.  If $2\leq r<\infty$, set $T=(r-1)S_r$;
if $r=\infty$ and $q$ is fixed, set $T=qS_\infty$.
Theorem~\ref{p11:thm:smooth-witness} supplies a nonzero real
$h\in G^\tau(\mathbb R)$ compactly supported inside $(1,2)$ with
\[
 \int_1^2\theta^s h(\theta)\,d\theta=0\quad(s\in T),
 \qquad \|h^{(k)}\|_\infty\leq A C^k(k!)^\tau
\]
for some $A,C>0$.  In particular $\int h=0$.
For $t>0$ let
\[
 \varphi_t(z)=\frac{1}{\sqrt{2\pi t}}e^{-z^2/(2t)},\qquad
 h_t(z)=\int_{\mathbb R}\varphi_t(z-y)h(y)\,dy,\qquad
 c_t=\int_1^2h_t(u)\,du.
\]
Compact support of $h$ makes $h_t$ entire.  On the real line,
integration by parts gives $h_t^{(k)}=\varphi_t*h^{(k)}$, whence
\[
 \|h_t^{(k)}\|_\infty\leq A C^k(k!)^\tau
 \quad(k\geq0),\qquad
 \|h_t-c_t\|_{L^\infty([1,2])}\leq2A.
\]
Furthermore $h_t\to h$ uniformly and $c_t\to0$ as $t\downarrow0$.
Because $\sigma-\tau>0$,
\[
 D=\max\left\{2A,
       A\sup_{k\geq1}\frac{(C/B)^k}{(k!)^{\sigma-\tau}}\right\}<\infty.
\]
Choose a fixed $\varepsilon>0$ satisfying
$\varepsilon D\leq\min\{1/2,M-1\}$ and put
\[
 g_t^\pm=1\pm\varepsilon(h_t-c_t),\qquad
 g^\pm=1\pm\varepsilon h.
\]
The interval has length one, so $\int g_t^\pm=1$ exactly.
Both $g_t^+$ and $g_t^-$ are entire, between $1/2$ and $3/2$,
and satisfy \eqref{p11:eq:entire-gevrey-ball}.  Indeed the bound for
$k=0$ is at most $1+\varepsilon D\leq M$, and for $k\geq1$
the derivative bound is at most
$\varepsilon D B^k(k!)^\sigma\leq MB^k(k!)^\sigma$.
Thus $g_t^\pm d\theta\in\mathcal E_{\sigma,M,B}$ for every $t>0$.
Their distinct limits $\mu^\pm=g^\pm d\theta$ belong to its weak
closure and have identical $q$-observation laws by
\eqref{p11:eq:observable-sumset}.  For finite $r$, this equality holds
simultaneously for every $q<r$.  The lower bound
\eqref{p11:eq:closure-fiber-lower-bound} gives the positive constant
\begin{equation}\label{p11:eq:explicit-regularity-floor}
 c=\frac12W_1(\mu^+,\mu^-)
   =\varepsilon\int_1^2
          \left|\int_1^x h(u)\,du\right|dx>0.
\end{equation}
Positivity follows because a nonzero continuous $h$ cannot have an
identically zero indefinite integral.

For the final assertion, $A_{q,S_r}$ is injective on all of the
compact space $\mathcal P([1,2])$ when $q\geq r$.
Proposition~\ref{p11:prop:closure-learning} gives uniform consistency.
\end{proof}

\begin{remark}[Interpretation of the two boundaries]
For a fixed bound with $\sigma=1$, the closure retains analytic
rigidity.  When $\sigma>1$, even if every member of the parameter
class is entire, its closure contains nonanalytic invisible
perturbations.  No finite number of independent groups overcomes
that uniform obstruction below the replication threshold.
The identities between analytic densities are still unique: the
lower bound concerns pairs of analytic densities whose data laws
become arbitrarily close, not an equality between two distinct
analytic densities.  Above the threshold, compactness restores
uniform consistency, without asserting any particular rate.
\end{remark}

\begin{corollary}[Strong log-concavity does not remove the obstruction]
\label{p11:cor:logconcave-analytic-learning}
Fix $\sigma>1$, $M>1$, $B>0$, and a finite $q<r$ in
Theorem~\ref{p11:thm:entire-gevrey-learning}.  There exist $\kappa,c>0$
such that the smaller class
\[
 \mathcal E_{\sigma,M,B}^{\kappa}
 =\left\{g\,d\theta\in\mathcal E_{\sigma,M,B}:
                (\log g)''\leq-\kappa\text{ on }[1,2]\right\}
\]
still satisfies $R_{N,q}(\mathcal E_{\sigma,M,B}^{\kappa})\geq c$
for every $N\geq1$.  For finite $r$, one choice of $\kappa,c>0$
works simultaneously for every $q<r$.
\end{corollary}

\begin{proof}
Choose $1<M_0<M$.  For a sufficiently small $a>0$, the entire
probability density
\[
 g_0(z)=\frac{e^{-a(z-3/2)^2}}
                   {\int_1^2e^{-a(u-3/2)^2}\,du}
\]
is between $3/4$ and $5/4$ on $[1,2]$, and its modulus is at most
$M_0$ on the closed complex $2/B$-neighborhood of that interval.
These assertions follow from uniform convergence to $1$ on this
complex compact set as $a\downarrow0$.  Cauchy's estimate gives
\[
 \|g_0^{(k)}\|_\infty\leq M_0(B/2)^k k!
        \leq M_0B^k(k!)^\sigma\qquad(k\geq0).
\]
Use the same invisible $h$ and entire centered approximants
$b_t=h_t-c_t$ as in the preceding proof, and put
$g_t^\pm=g_0\pm\varepsilon b_t$.
Their derivatives fit the prescribed bounds if
$\varepsilon D\leq M-M_0$.  Their integrals are one exactly.
The functions $b_t$ have uniformly bounded $C^2([1,2])$ norms.
Because $(\log g_0)''=-2a$, taking $\varepsilon>0$ still smaller
ensures simultaneously that all $g_t^\pm$ lie between $1/2$ and
$3/2$ and satisfy $(\log g_t^\pm)''\leq-a$.
This follows also directly from continuity of
$g\mapsto g''/g-(g'/g)^2$ in the $C^2$ norm on densities bounded
away from zero.  The distinct limits
$(g_0\pm\varepsilon h)d\theta$ have equal $q$-observation laws.
Proposition~\ref{p11:prop:closure-learning} supplies the positive
all-$N$ lower bound, with $\kappa=a$.
\end{proof}



\subsection{Identification and consistency across mixing classes}

For a fixed finite threshold $r\geq2$, the different consequences
of an exact observation law and of a finite statistical sample can
be displayed together. In the table, ``least group size'' refers
to identification within the stated class or to uniform $W_1$
consistency as the number of independent groups tends to infinity.
\begin{center}
\small
\begin{tabularx}{\textwidth}{@{}Xcc@{}}
\toprule
Mixing class & Exact identification & Uniform consistency\\
\midrule
All Borel probability measures & $r$ & $r$\\
All finite atomic probabilities, with no common bound on atom count
 & $1$ & $r$\\
Entire densities with fixed bounds $\mathcal E_{1,M,B}$
 & $1$ & $1$\\
Entire densities with fixed bounds $\mathcal E_{\sigma,M,B}$,
 $\sigma>1$ & $1$ & $r$\\
\bottomrule
\end{tabularx}
\end{center}
The finite atomic row follows because finite atomic probabilities
are weakly dense in all probability measures, together with
Lemma~\ref{p11:lem:interval-atomic} and
Proposition~\ref{p11:prop:closure-learning}. If a common finite upper
bound on atom count is imposed, that subclass is compact and
one-observation uniform consistency instead holds. No rate for
that compact finite-atom problem is claimed here.

The distinction between local and global obstructions is also sharp.
Lemma~\ref{p11:lem:interval-jets} excludes every nonzero finite point-derivative
functional annihilating the component kernel. In particular there is
no nontrivial finite-order differential relation at an interior
parameter point. Nevertheless the global observation operator has
localized, infinite-dimensional smooth nullspaces below the prescribed
threshold. Smoothness of an invisible perturbation does not imply
that it arises from a finite parameter differential identity.


\section{Sharp quantitative recovery for the factorial kernel}
\label{p10:sec:quantitative-factorial}

The exact recovery supplied by replication need not be numerically
well conditioned. We return to the factorial Legendre kernel
\eqref{p09:eq:continuous-missing-mode-kernel}, now denoted
$k_{\mathrm{fac}}$. Its qualitative identification properties were proved
in Corollary~\ref{p09:cor:entire-continuous-counterexample}.
For an integer $r\ge1$ and a mixing probability $G$, write
\begin{equation}\label{p10:eq:quant-r-view}
 q_{r,G}(x_1,\ldots,x_r)
   =\int_0^1\prod_{i=1}^r k_{\mathrm{fac}}(x_i\mid\theta)\,G(d\theta),
 \qquad Q_r(G)=q_{r,G}\,dx_1\cdots dx_r.
\end{equation}
The $r$ observations in a group share the same latent parameter.
Different groups in the statistical experiment have independent latent
draws. The density is normalized by Legendre orthogonality, is
bounded below by $3-e>0$, and is jointly entire by locally uniform
convergence on complex compact sets.  We use
$\operatorname{TV}(P,Q)=\tfrac12\int|p-q|$ and the ordinary
Euclidean Wasserstein distance $W_1$ on $[0,1]$.

\begin{theorem}[Sharp modulus and statistical rate]
\label{p10:thm:quant-factorial-rate}
For every fixed integer $r\ge2$, define
\[
 \omega_r(\varepsilon)=
 \sup\{W_1(G,H):\operatorname{TV}(Q_r(G),Q_r(H))\le\varepsilon\},
\]
where the supremum ranges over all Borel probabilities on $[0,1]$.
There are constants $0<c_r\le C_r<\infty$ such that, for all sufficiently
small $\varepsilon>0$,
\begin{equation}\label{p10:eq:quant-inverse-modulus}
 c_r\frac{\log\log(1/\varepsilon)}{\log(1/\varepsilon)}
 \le\omega_r(\varepsilon)
 \le C_r\frac{\log\log(1/\varepsilon)}{\log(1/\varepsilon)}.
\end{equation}
If $N$ independent groups with law $Q_r(G)$ are observed, then
\begin{equation}\label{p10:eq:quant-minimax}
 \inf_{\widehat G}\sup_G
   \mathbb E_G W_1(\widehat G,G)
 \asymp_r\frac{\log\log N}{\log N}
 \qquad(N\longrightarrow\infty).
\end{equation}
The infimum is over measurable probability-valued estimators.  For
$r=1$, the minimax risk is at least $1/12$ for every $N$, and
$\omega_1(0)\ge1/6$.
\end{theorem}

The rate $\log\log N/\log N$ is already familiar in nonparametric
mixture estimation, including compactly supported Poisson mixtures
\cite{MiaoEtAl2024}.  Our assertion is a quantitative description of
the present replication transition, not a new universal mixture rate.
In particular, any fixed number of further replicates leaves the rate
order unchanged after identifiability has been restored.  The theorem fixes the replicate count. Growing replication and an
equal-group observation budget are treated in
Section~\ref{p11:sec:replication-budget}; those bounds leave a logarithmic
gap.

\subsection{Moment recovery and approximation}

Let $\nu=G-H$ and $\delta_j=\int L_j\,d\nu$.  Orthogonality gives
\begin{align}
 \int L_j(x)\,d\{Q_2(G)-Q_2(H)\}(x,y)
   &=\frac{\delta_j}{j!(2j+1)},\qquad j\ge2,
   \label{p10:eq:quant-visible-moment}\\
 \int L_2(x)L_3(y)\,d\{Q_2(G)-Q_2(H)\}(x,y)
   &=\frac1{420}\int L_2L_3\,d\nu.
   \label{p10:eq:quant-missing-moment}
\end{align}
The elementary product identity
\begin{equation}\label{p10:eq:quant-product-identity}
 L_2L_3=\frac{10}{21}L_5+\frac4{15}L_3+\frac9{35}L_1
\end{equation}
therefore recovers the missing first moment from finitely many observable
expectations.  For $M\ge5$, consider the $M$ bounded test functions
\[
 \mathcal F_M=
 \{(x,y)\mapsto L_j(x):2\le j\le M\}
 \cup\{(x,y)\mapsto L_2(x)L_3(y)\}.
\]
If the absolute differences of their expectations under $Q_2(G)$ and
$Q_2(H)$ are bounded by $d$, then
\begin{equation}\label{p10:eq:quant-moment-bounds}
 |\delta_j|\le j!(2j+1)d\quad(2\le j\le M),
 \qquad |\delta_1|\le4200d.
\end{equation}
Indeed, the constant for the last inequality is bounded by
\[
 \frac{35}{9}
 \left(420+\frac{10}{21}\,1320+\frac4{15}\,42\right)<4200.
\]

\begin{lemma}[Polynomial approximation and moment comparison]
\label{p10:lem:quant-jackson}
There is an absolute constant $C$ such that, for every integer $M\ge5$
and every $1$-Lipschitz function $f:[0,1]\to\mathbb R$ with $f(0)=0$,
there is a polynomial $p$ of degree at most $M$ satisfying
\[
 \|p\|_\infty\le1,\qquad \|f-p\|_\infty\le C/M.
\]
Consequently, under the expectation discrepancy assumption above,
\begin{equation}\label{p10:eq:quant-finite-comparison}
 W_1(G,H)\le\frac C M+C M^3M!d.
\end{equation}
\end{lemma}

\begin{proof}
We include the classical Jackson-kernel argument.  Define the even
$2\pi$-periodic function $F(u)=f((1+\cos u)/2)$.  It is
$1/2$-Lipschitz for circular distance, and $\|F\|_\infty\le1$.
Put $\ell=\lfloor M/2\rfloor+1$ and normalize
\[
 J_\ell(u)=c_\ell
   \left(\frac{\sin(\ell u/2)}{\sin(u/2)}\right)^4,
 \qquad \int_{-\pi}^{\pi}J_\ell(u)\,du=1.
\]
The quotient at zero is understood by continuity.  This is a
nonnegative trigonometric polynomial of degree $2(\ell-1)\le M$.
On $[-\pi,\pi]$ the modulus of the quotient is bounded by
$\min(\ell,\pi/|u|)$.  On $|u|\le1/\ell$ it is at least
$2\ell/\pi$.  These inequalities show that its fourth-power integral
is comparable to $\ell^3$, whereas its fourth-power integral weighted
by $|u|$ is $O(\ell^2)$.  Hence
\[
 \int_{-\pi}^{\pi}|u|J_\ell(u)\,du\le C/\ell.
\]
The convolution $F*J_\ell$ is even, has sup norm at most one, and
differs uniformly from $F$ by at most $C/\ell$.  Every even
trigonometric polynomial of degree at most $M$ is a polynomial of
degree at most $M$ in $\cos u$.  Substituting $\cos u=2x-1$ gives
the required $p$.

Expand $p=\sum_{j=0}^M b_jL_j$.  Orthogonality and $\|p\|_\infty\le1$
give $|b_j|\le2j+1$.  Since $\nu([0,1])=0$, the
Kantorovich--Rubinstein formula and the approximation just proved yield
\[
 W_1(G,H)\le\frac C M+
       \sum_{j=1}^M(2j+1)|\delta_j|.
\]
Use \eqref{p10:eq:quant-moment-bounds} and the sufficient coarse bound
$\sum_{j=2}^M(2j+1)^2j!\le M(2M+1)^2M!$.
\end{proof}

\subsection{Explicit indistinguishable approximations}

\begin{lemma}[Moment matching with an exact transport gap]
\label{p10:lem:quant-binomial}
For every integer $m\ge1$, set $n=m+1$ and define
\begin{align*}
 G_m&=2^{1-n}\sum_{\substack{0\le j\le n\\j\ {\mathrm{even}}}}
          \binom nj\delta_{j/n},\qquad
 H_m=2^{1-n}\sum_{\substack{0\le j\le n\\j\ {\mathrm{odd}}}}
          \binom nj\delta_{j/n}.
\end{align*}
These are probabilities with equal polynomial moments through degree
$m$, and
\begin{equation}\label{p10:eq:quant-exact-transport}
 W_1(G_m,H_m)=\frac1{m+1}.
\end{equation}
For every integer $r\ge1$,
\begin{equation}\label{p10:eq:quant-factorial-tail}
 \operatorname{TV}(Q_r(G_m),Q_r(H_m))
 \le T_{m,r}:=\sum_{s=m+1}^{\infty}\frac{r^s}{s!}.
\end{equation}
If $m+2\ge2r$, then $T_{m,r}\le2r^{m+1}/(m+1)!$.
\end{lemma}

\begin{proof}
The even and odd binomial sums both equal $2^{n-1}$, so the measures
are probabilities.  Their difference is a constant multiple of an
$n$th finite difference and therefore annihilates every polynomial of
degree at most $n-1=m$.  On the interval $[k/n,(k+1)/n)$, their
cumulative distribution functions differ by
\[
 2^{1-n}\sum_{j=0}^k(-1)^j\binom nj
   =2^{1-n}(-1)^k\binom{n-1}{k}.
\]
Integrating the absolute difference proves
\eqref{p10:eq:quant-exact-transport}.

For the kernel expansion, allow indices in $\{0,2,3,\ldots\}$,
with $L_0=1$ and coefficient $1/0!=1$.  A term in the $r$-fold
product indexed by $(j_1,\ldots,j_r)$ is a polynomial in $\theta$
of degree at most $s=j_1+\cdots+j_r$.  Its integral against
$G_m-H_m$ vanishes for $s\le m$ and otherwise has absolute value
at most two.  Absolute uniform convergence justifies all
interchanges.  Enlarging the index sets to all nonnegative integers
and using the multinomial identity gives
\[
 \|q_{r,G_m}-q_{r,H_m}\|_\infty
 \le2\sum_{s=m+1}^{\infty}
       \sum_{j_1+\cdots+j_r=s}\frac1{j_1!\cdots j_r!}
 =2\sum_{s=m+1}^{\infty}\frac{r^s}{s!}.
\]
The cube has volume one, proving \eqref{p10:eq:quant-factorial-tail}.
When $m+2\ge2r$, successive tail terms have ratio at most $1/2$,
which proves the final bound.
\end{proof}

\subsection{Proof of the sharp modulus and minimax theorem}

\begin{proof}[Proof of Theorem~\ref{p10:thm:quant-factorial-rate}]
If $\operatorname{TV}(Q_2(G),Q_2(H))\le\varepsilon$, all tests in
$\mathcal F_M$ have expectation discrepancy at most $2\varepsilon$.
Write $L=\log(1/\varepsilon)$ and choose
$M=\lfloor L/(2\log L)\rfloor$.  For sufficiently large $L$ this
is at least five, and
\[
 \log(\varepsilon M^3M!)
 \le-L+M\log M+3\log M\le-L/3.
\]
Lemma~\ref{p10:lem:quant-jackson} gives the upper bound in
\eqref{p10:eq:quant-inverse-modulus}.  For $r>2$, marginalization to
the first two coordinates can only decrease total variation, so
the same upper bound applies.  For the lower bound, choose
$m=\lceil2L/\log L\rceil$.  For every fixed $r$, the elementary
factorial bound $(m+1)!\ge((m+1)/e)^{m+1}$ shows that
$T_{m,r}\le\varepsilon$ for all sufficiently large $L$.
Lemma~\ref{p10:lem:quant-binomial} then gives a transport gap
$1/(m+1)\ge c_r\log L/L$.

For the statistical upper bound, retain the first two coordinates of
each group.  Let $\widehat a$ be the vector of empirical means of the
$M$ tests in $\mathcal F_M$, and let $a(G)$ denote their population
expectations.  Choose $\widehat G$ with
\[
 \|a(\widehat G)-\widehat a\|_\infty
 \le\inf_{H}\|a(H)-\widehat a\|_\infty+N^{-1}.
\]
For completeness, a measurable choice is obtained by fixing a
countable dense family of finitely supported probabilities with
rational atoms and rational weights, taking the infimum over this
family, and selecting its first qualifying element.  Continuity
of all feature expectations in the weak topology ensures that
this infimum equals the displayed one.
If $Z=\|\widehat a-a(G)\|_\infty$, then
$\|a(\widehat G)-a(G)\|_\infty\le2Z+N^{-1}$.
Hoeffding's inequality for bounded tests, a union bound, and
integration yield
\[
 \mathbb E_G Z\le C\sqrt{\frac{\log(2M)}{N}},
\]
uniformly in $G$.  Consequently
\[
 \sup_G\mathbb E_G W_1(\widehat G,G)
 \le\frac C M+CM^3M!
       \left(\sqrt{\frac{\log(2M)}{N}}+\frac1N\right).
\]
Take $M=\lfloor\log N/(4\log\log N)\rfloor$.
Since $M!\le M^M\le N^{1/3}$ eventually, the second term is
$o(1/M)$, proving the upper rate.

For the lower bound, set $m=\lceil2\log N/\log\log N\rceil$.
For fixed $r$, $NT_{m,r}\to0$.  The total variation distance
between the two $N$-group sample laws is at most $NT_{m,r}$,
and hence at most $1/2$ eventually.  Integrating the triangle
inequality for $W_1$ against the overlap of these sample laws
shows that the larger of the two estimation risks at $G_m,H_m$
is at least $1/[4(m+1)]$.  This proves the lower rate.

Finally, the mixing densities $1+\tfrac12L_1$ and
$1-\tfrac12L_1$ both give the uniform one-view observation law.
Their signed density difference is $2\theta-1$, so their
cumulative distribution functions differ by $\theta^2-\theta$.
Their $W_1$ distance is $\int_0^1(\theta-\theta^2)d\theta=1/6$.
The identical sample laws and the triangle inequality force
the larger estimation risk to be at least $1/12$ for every $N$.
\end{proof}



\section{Growing replication under an equal-group budget}
\label{p11:sec:replication-budget}

We continue with the factorial kernel
$k_*=k_{\mathrm{fac}}$ from \eqref{p09:eq:continuous-missing-mode-kernel}.
The preceding theorem determines the rate for each fixed group size;
here the number of replicates may grow with the number of groups.
Let $Q_q(G)=\int P_\theta^{\otimes q}\,dG(\theta)$, and write
\[
 R_{N,q}=\inf_{\widehat G}\sup_{G\in\mathcal P([0,1])}
       \mathbb E_{Q_q(G)^{\otimes N}}W_1(\widehat G,G).
\]
Thus $N$ independent groups each contain $q$ observations sharing one
latent parameter.  All constants below are independent of $N,q$.

\begin{theorem}\label{p11:thm:budget-bounds}
There are constants $c,C>0$ such that, for $N\ge1$ and $q\ge2$,
\begin{align}
 R_{N,q}&\ge c\max\left\{\frac1{\sqrt N},
     \min\left(\frac1{\log(eN)},\frac1{\sqrt{q\log(eN)}}\right)\right\},
      \label{p11:eq:budget-lower}\\
 R_{N,q}&\le\min\left\{1,\frac1{2\sqrt N}
                 +\frac{2100}{\sqrt{\lfloor q/2\rfloor}}\right\}.
      \label{p11:eq:budget-upper}
\end{align}
Consequently, for sufficiently large $B$,
\begin{equation}\label{p11:eq:budget-total}
 c\{B\log(eB)\}^{-1/4}
 \le\inf_{\substack{N,q\in\mathbb N,\ q\ge2\\Nq\le B}}R_{N,q}
 \le C B^{-1/4}.
\end{equation}
In particular $q\ge N$ gives $R_{N,q}\asymp N^{-1/2}$, whereas
$R_{N,q}=O(N^{-1/2})$ along a sequence $N\to\infty$ requires
$q=\Omega(N/\log(eN))$.
\end{theorem}

These inequalities determine the equal-group budget exponent up to a
logarithmic factor; they do not resolve the complete growing-$q$
minimax curve.  The infimum in \eqref{p11:eq:budget-total} concerns equal,
nonadaptively chosen group sizes, not arbitrary allocations.
Related logarithmic gains and the limitations of empirical latent
parameter estimates are classical in binomial-mixture estimation
\cite{Vinayak2019}; no novelty of those rate phenomena is asserted.

\begin{proof}[Proof of the upper bound]
Put $A_j(x)=j!(2j+1)L_j(x)$, so
$\mathbb E_\theta A_j(X)=L_j(\theta)$ for $j\ge2$.
The product identity \eqref{p10:eq:quant-product-identity} shows that the statistic
\[
 Y(x,y)=\frac12+\frac{35}{18}
 \left\{A_2(x)A_3(y)-\frac{10}{21}A_5(x)-\frac4{15}A_3(x)\right\}
\]
has conditional expectation $\theta$ when $X,Y$ are independent
with common law $P_\theta$.  Moreover
\[
 |Y(x,y)|\le\frac12+\frac{35}{18}
       \left(420+\frac{10}{21}\,1320+\frac4{15}\,42\right)<2100.
\]
Within each group, average this statistic over
$M=\lfloor q/2\rfloor$ disjoint pairs and clip the result to $[0,1]$,
obtaining $\widehat\theta_i$.  Conditional independence and clipping give
$\mathbb E_\theta|\widehat\theta_i-\theta|\le2100/\sqrt M$.
For $\widehat G=N^{-1}\sum_i\delta_{\widehat\theta_i}$, couple with
the latent empirical measure $G_N=N^{-1}\sum_i\delta_{\theta_i}$.
If $F_N,F$ are the distribution functions of $G_N,G$, respectively,
\[
 \mathbb E W_1(\widehat G,G)
 \le\frac{2100}{\sqrt M}+\int_0^1\mathbb E|F_N(t)-F(t)|\,dt
 \le\frac{2100}{\sqrt M}+\frac1{2\sqrt N}.
\]
This proves \eqref{p11:eq:budget-upper}.
\end{proof}

\begin{proof}[Proof of the lower bound]
An oracle observing the $N$ latent parameters is more informative than
the grouped experiment.  Consider
$G_\pm=(\frac12\pm a)\delta_1+(\frac12\mp a)\delta_0$ with
$a=1/(8\sqrt N)$.  Their $W_1$ distance is $2a$ and their
$N$-sample relative entropy is at most $1/6$.
Pinsker's inequality and the two-point bound therefore give
$R_{N,q}\ge1/(16\sqrt N)$.

For the other bound, put $\theta_0=1/2$ and
$\mathcal H=L^2(P_{\theta_0};\mathbb C)$.
The normalized likelihood
\[
 f_z(x)=\frac{k_*(x\mid\theta_0+z)}{k_*(x\mid\theta_0)}
\]
is $\mathcal H$-holomorphic near $z=0$ and satisfies
$\int f_z\,dP_{\theta_0}=1$.
Since the denominator is bounded below and the numerator and its
complex derivative are bounded on compact sets, there exist
$0<R\le1/4$ and $C_0\ge1$ such that
$\|f_z-1\|_{\mathcal H}\le\sqrt{C_0}|z|$ for $|z|\le R$.
Normalization cancels the linear term in the squared norm, giving
\[
 \|f_z\|_{\mathcal H}^2=1+\|f_z-1\|_{\mathcal H}^2
      \le1+C_0|z|^2.
\]
The $q$-fold likelihood $F_z=f_z^{\otimes q}$ in
$\mathcal H_q=L^2(P_{\theta_0}^{\otimes q};\mathbb C)$ consequently
satisfies $\|F_z\|_{\mathcal H_q}\le\exp(C_0q|z|^2/2)$.
Writing $F_z=\sum_{j\ge0}T_jz^j$, the Hilbert-space Cauchy estimate
on $|z|=2h$, for $0<h\le R/2$, gives
\begin{equation}\label{p11:eq:budget-cauchy}
 \|T_j\|_{\mathcal H_q}\le
           e^{2C_0qh^2}(2h)^{-j}.
\end{equation}

Let $m\ge1$, $n=m+1$, and $\theta_j=\theta_0-h+2hj/n$.
The probabilities
\[
 G=2^{1-n}\sum_{\substack{0\le j\le n\\j\ \mathrm{even}}}
            \binom nj\delta_{\theta_j},\qquad
 H=2^{1-n}\sum_{\substack{0\le j\le n\\j\ \mathrm{odd}}}
            \binom nj\delta_{\theta_j}
\]
have matching moments through degree $m$, by the formula for an
$n$th finite difference.  Their transport distance is
$W_1(G,H)=2h/(m+1)$: on the $j$th grid interval the difference of their
cumulative distribution functions is
$2^{1-n}(-1)^j\binom{n-1}{j}$, whose absolute values sum to one.
Since $|\int(\theta-\theta_0)^j\,d(G-H)|\le2h^j$,
\eqref{p11:eq:budget-cauchy} yields
\[
 \left\|\frac{dQ_q(G)-dQ_q(H)}{dP_{\theta_0}^{\otimes q}}
       \right\|_{\mathcal H_q}
 \le2e^{2C_0qh^2}\sum_{j>m}2^{-j}
 =2e^{2C_0qh^2}2^{-m}.
\]
In particular their one-group total variation distance is at most
$e^{2C_0qh^2}2^{-m}$.
Choose
\[
 m=\left\lceil\frac4{3\log2}\log(4N)\right\rceil,
 \qquad h=\min\left\{\frac R2,
             \sqrt{\frac{m\log2}{8C_0q}}\right\}.
\]
Then the one-group distance is at most $1/(4N)$, so the
$N$-group distance is at most $1/4$.
The two-point inequality gives
\[
 R_{N,q}\ge\frac{3h}{4(m+1)}
 \ge c\min\left\{\frac1{\log(eN)},
                    \frac1{\sqrt{q\log(eN)}}\right\}.
\]
Together with the oracle argument this proves \eqref{p11:eq:budget-lower}.
Notice that this local lower-bound proof applies to any normalized
analytic likelihood satisfying the displayed $L^2$ bound; it does
not depend on the particular missing Legendre coefficient.
\end{proof}

\begin{proof}[Proof of the budget and replication conclusions]
Choosing $q=\lfloor\sqrt B\rfloor$ and $N=\lfloor B/q\rfloor$
proves the upper bound in \eqref{p11:eq:budget-total}.
Write $L=\log(eN)$.  If $q\ge L$, the lower bound is at least a
constant times
\[
 \max\{N^{-1/2},(qL)^{-1/2}\}
 \ge(NqL)^{-1/4}\ge\{B\log(eB)\}^{-1/4}.
\]
If $q<L$, it is at least $c/L\ge c/\log(eB)$, which is larger
than the asserted bound for all sufficiently large $B$.
Finally, $q\ge N$ makes the upper bound $O(N^{-1/2})$.
Conversely, $R_{N,q}=O(N^{-1/2})$ first forces $q\ge L$
eventually, since $1/L\gg N^{-1/2}$, and then forces
$qL=\Omega(N)$ by \eqref{p11:eq:budget-lower}.
\end{proof}




\section{A fixed entire kernel with arbitrarily slow recovery}
\label{p10:sec:quant-arbitrary-slow}

The preceding rate is a feature of the factorial coefficients, not a
consequence of entire regularity or of the finite-jet independence
property.  The following statement makes this separation quantitative.
It concerns worst-case estimation of the mixing probability in $W_1$;
it does not concern estimation of the observation density.

\begin{theorem}[No rate follows from the qualitative properties]
\label{p10:thm:quant-arbitrary-slow}
Let $(\rho_N)_{N\ge1}$ be any positive sequence tending to zero.
There is a single strictly positive, jointly entire probability kernel
$k_\rho$ on $[0,1]^2$ with the following properties.  Its one-view
mixture operator has precisely the signed-measure null space
$\{bL_1(\theta)d\theta:b\in\mathbb R\}$; its two-view mixture
operator is injective; and every finite family of parameter jets at
distinct parameter points is linearly independent.  If
\[
 R_{N,r}(k_\rho)=\inf_{\widehat G}\sup_G
       \mathbb E_{G,k_\rho}W_1(\widehat G,G)
\]
denotes the risk from $N$ independent groups of $r$ observations,
then, for every fixed integer $r\ge2$,
\begin{equation}\label{p10:eq:quant-arbitrarily-slow}
 R_{N,r}(k_\rho)\longrightarrow0,
 \qquad
 \lim_{N\to\infty}\frac{R_{N,r}(k_\rho)}{\rho_N}=\infty.
\end{equation}
The same kernel works for all fixed $r$, and the rate comparison holds
for all sufficiently large sample sizes, not only along a subsequence.
\end{theorem}

\begin{proof}
Choose strictly increasing positive integers $N_j$ such that
$N_j\ge j$ and $\sup_{N\ge N_j}\rho_N\le(j+1)^{-2}$.
Such a choice is possible even when $(\rho_N)$ is not monotone.  Set
\[
 a_j=\frac{2^{-j}}{j!N_j^2}\quad(j\ge2),\qquad
 k_\rho(x\mid\theta)=1+\sum_{j=2}^{\infty}
                   a_jL_j(x)L_j(\theta).
\]
All coefficients are positive, and $a_j\le1/j!$.  The Legendre
polynomials grow at most exponentially in degree on every fixed
compact subset of $\mathbb C$, so the series converges locally
uniformly on $\mathbb C^2$.  It is jointly entire and is bounded
below by $3-e>0$ on the real unit square.  Orthogonality gives
normalization in $x$.

For clarity, the qualitative properties can be checked directly.
Coefficient extraction gives
\[
 \int_0^1L_j(x)\int_0^1 k_\rho(x\mid\theta)\nu(d\theta)\,dx
   =\frac{a_j}{2j+1}\int_0^1L_j\,d\nu\quad(j\ge2).
\]
Thus a null measure has zero mass and zero Legendre moments except
possibly its first moment.  Polynomial density identifies it as
$bL_1(\theta)d\theta$.  Two views additionally give its $L_2L_3$ moment; the calculation in
Corollary~\ref{p09:cor:two-replicate-repair} gives
$\int_0^1L_1L_2L_3=3/35\ne0$, so the two-view null measure is zero.
For finite-jet independence, coefficient extraction gives $T(1)=0$
and $T(L_j)=0$ for $j\ge2$ for any purported finite-jet annihilator
$T$. Lemma~\ref{p09:lem:no-finite-jet} then forces every coefficient
of $T$ to vanish.

Fix $r\ge2$.  For every sufficiently large sample size $N$, take
the unique integer $j\ge2$ with $N_{j-1}\le N<N_j$ and put
$m=r(j-1)$.  Apply the measures in
Lemma~\ref{p10:lem:quant-binomial} with this $m$.  In an $r$-fold
kernel product, every term of total polynomial degree exceeding
$r(j-1)$ has at least one index at least $j$.  Write
$A=1+\sum_{k\ge2}a_k<\infty$.  The same absolute-convergence
argument as in that lemma, followed by a union bound on the
large index, gives
\[
 \operatorname{TV}(Q_r(G_m),Q_r(H_m))
 \le rA^{r-1}\sum_{k\ge j}a_k
 \le\frac{rA^{r-1}}{N_j^2}.
\]
The last inequality uses monotonicity of $N_k$ and
$\sum_{k\ge j}2^{-k}/k!\le1$.  For the $N$-group sample laws,
the total variation distance is at most
$rA^{r-1}N/N_j^2\le rA^{r-1}/N_j$, which tends to zero
as $N\to\infty$.  The two-point risk bound and the exact
transport gap consequently imply
\[
 R_{N,r}(k_\rho)\ge\frac1{4\{r(j-1)+1\}}
 \quad\text{for all sufficiently large }N.
\]
Our choice of $N_{j-1}$ ensures that $\rho_N\le j^{-2}$.
Hence the risk ratio is bounded below by
$j^2/[4\{r(j-1)+1\}]$, which tends to infinity.

Finally, consistency is not lost.  For each fixed approximation
degree $M\ge5$, all coefficients $a_2,\ldots,a_M$ are nonzero.
The moment-recovery and minimum-discrepancy arguments used above
therefore give, with a finite constant $C_M$ depending on this
kernel and on $M$,
\[
 R_{N,r}(k_\rho)\le C/M+C_M/\sqrt N+C_M/N.
\]
First let $N\to\infty$ and then $M\to\infty$.
\end{proof}

This construction is consistent with the general distinction between
identifiability and stability in inverse problems.  Its role here is
to show that strict positivity, joint entire regularity, independence
of all finite parameter jets, and recovery from two replicates can
coexist with an arbitrarily slow statistical inverse.






\section{Scope and remaining quantitative questions}
\label{p05:sec:mixture-scope}

Finite point and derivative independence, exact recovery from grouped
laws, and statistical conditioning govern different questions.
The circle examples retain analytic mixing-density ambiguity, whereas
the interval examples distinguish analytic rigidity from localized
smooth ambiguity. A finite-dimensional one-observation nullspace
reduces to a fixed finite atomic set after two observations; arbitrary
thresholds exceeding two therefore require an infinite-dimensional
defect when all finite component families are independent.
In the infinite-threshold constructions, the witnessing pair may vary
with the finite group size: no two distinct mixing laws agree at every
order, because the complete hierarchy identifies the law.

The stable-family classification fixes the index and skewness within
each model. Its sparse conditions outside the Gaussian case are only
sufficient, and mixtures over the shape parameters are not covered.
The exact scalar repairs for finite Legendre defects may be severely
ill-conditioned. The fixed analytic-bound learning theorem supplies
consistency without a sharp rate. Section~\ref{p11:sec:replication-budget}
leaves a logarithmic gap and treats only equal, nonadaptive group sizes;
it does not settle unequal or adaptive replication or the full
growing-replication minimax curve. The exact moment-coincidence
construction does not control the size of the first unequal moment or
give an efficient numerical representation of the densities.
The threshold kernels are constructed examples; analogous conclusions
for particular standard statistical families need additional arguments.

\paragraph{Source versions and verification status.}
The normal-mixture question is used as formulated in
\texttt{arXiv:1507.00678v1}, Section~2.6. The finite differential and
difference--differential classes are those in the author manuscript
of Le et al. dated 8 August 2026, where the sufficiency question follows
Proposition~6 and appears in Section~6. The localized moment theorem
proves the assertion of Conjecture~17 in the cited
B\"urgin--Epperlein--Wirth source. Historical priority remains
unverified; Lin's final publisher PDF was not inspected. The proofs
and the exact Legendre identity received AI-assisted checks only.

\clearpage
\section{Source correspondence and consolidation}
\label{sec:source-correspondence}
This integrated manuscript incorporates the eleven English research
drafts listed below, all authored under the name KEIO and prepared with
AI assistance on 13--14 October 2026. The list records provenance rather
than eleven independent external references. Earlier inspirational
papers that were not generated in this sequence are cited in the
bibliography where used.

\begingroup
\small
\setlength{\tabcolsep}{5pt}
\renewcommand{\arraystretch}{1.15}
\begin{longtable}{@{}p{0.06\textwidth}p{0.51\textwidth}p{0.37\textwidth}@{}}
\toprule
No. & Source manuscript & Location and retained material \\
\midrule
\endfirsthead
\toprule
No. & Source manuscript & Location and retained material \\
\midrule
\endhead
\bottomrule
\endfoot
\hypertarget{source:p01}{1} & \emph{Sharp Recurrence Exponents for Prime Flows and Finite-Height Zeros at a Frozen Boundary} & Part I. Power-weight recurrence, finite-height boundary zeros, and the explicit Fourier certificate. \\
\hypertarget{source:p02}{2} & \emph{Entropy-Scale Recurrence and Sharp Quantitative Rigidity at Zeta Edges} & Part I. Logarithmic weights, Euler-edge locking, random perturbations, and positive-kernel bounds. \\
\hypertarget{source:p03}{3} & \emph{Prime-Phase Localization of Simple Zeta Zeros and Quantitative Euler-Edge Repulsion} & Part II. Fixed prime-phase localization, edge triples, inverse principles, and spectral certificates. \\
\hypertarget{source:p04}{4} & \emph{Shrinking Prime-Phase Targets, Gibbs Profiles, and Simultaneous Zeta-Edge Recurrence} & Part II. Shrinking and growing phase targets, simultaneous heights, Gibbs laws, and joint edge transitions. \\
\hypertarget{source:p05}{5} & \emph{Two-Scale Identification from Joint Rare Events: Microscopic Laws, Spectral Fluctuations, and Singular Constraint Activation} & Part III. Microscopic-law inversion, discrete spectral fluctuations, hard constraints, and activation laws. \\
\hypertarget{source:p06}{6} & \emph{Minimal Observation and Blind Calibration of Rare-Event Laws} & Part III. Minimal channels, blind power calibration, correlated laws, and finite-information obstructions. \\
\hypertarget{source:p07}{7} & \emph{Arithmetic Observability through Controlled Thinning: Blind Reconstruction, Resonance, and Sharp Information Barriers} & Part IV. Controlled thinning, primitive exponent recovery, resonances, and optimal window noise. \\
\hypertarget{source:p08}{8} & \emph{Blind Identification beyond Moment Assumptions: Interior Probes, Hidden Coordinates, and Finite Information} & Part IV. Interior and flat probes, countable coordinates, finite gate values, and information bounds. \\
\hypertarget{source:p09}{9} & \emph{Mixture Identification beyond Moments and Differential Barriers} & Part V. Sparse partial sums, stable-family classification, and differential obstructions. \\
\hypertarget{source:p10}{10} & \emph{Replication Thresholds and Finite-Dimensional Ambiguity in Infinite Mixtures} & Part V. Finite defects, analytic-circle replication thresholds, factorial rates, and slow recovery. \\
\hypertarget{source:p11}{11} & \emph{Prescribed Moment Coincidences, Smooth Densities, and Replication Thresholds} & Part V. Prescribed moments, ordinary-interval thresholds, Gevrey boundaries, and growing replication. \\
\end{longtable}
\endgroup

\paragraph{Editorial treatment.}
Repeated foundations and identities are proved once and referenced
internally. Distinct results, constructions, and proof arguments are
retained. The later recurrence, shrinking-target, and interval-mixture
theorems supersede the corresponding earlier restrictions or open
questions. Conditional hypotheses and remaining quantitative gaps are
unchanged; verification status is discussed in the introduction.
\clearpage

\addcontentsline{toc}{section}{References}
\begin{thebibliography}{99}
\bibitem{Baroni}
L.~Baroni,
\emph{Kronecker Flow on the Infinite Torus},
arXiv:2604.02213 (2026), version 1.
\url{https://arxiv.org/abs/2604.02213}.


\bibitem{BCS}
I.~Belovas, R.~\v{C}epaityt\.{e}, and M.~Sabaliauskas,
\emph{On the zero-free region and the distribution of zeros of the
prime zeta function},
An. \c{S}tiin\c{t}. Univ. Ovidius Constan\c{t}a Ser. Mat.
\textbf{33} (2025), no.~2, 27--44.
\href{https://doi.org/10.2478/auom-2025-0017}{doi:10.2478/auom-2025-0017}.


\bibitem{Farag}
H.~M.~Farag,
\emph{Dirichlet truncations of the Riemann zeta function in the
critical strip possess zeros near every vertical line},
Int. J. Number Theory \textbf{4} (2008), no.~4, 653--662.
\href{https://doi.org/10.1142/S1793042108001596}{doi:10.1142/S1793042108001596}.


\bibitem{Keio004}
KEIO,
\emph{Prime-Zeta Divisor Tomography: Local Zero Processes,
Frozen-Edge Rigidity, and Derivative Caustics},
unpublished research manuscript, supplied in the author's
\texttt{[004]} source collection (2026).
Used for motivation; its non-elementary assertions are not inputs
to the proofs of this paper.


\bibitem{KR}
M.~A.~Korolev and I.~S.~Rezvyakova,
\emph{On simultaneous approximations to the logarithms of primes},
Chebyshevskii Sb. \textbf{23} (2022), no.~5, 87--100.
\href{https://doi.org/10.22405/2226-8383-2022-23-5-87-100}
{doi:10.22405/2226-8383-2022-23-5-87-100}.


\bibitem{MontgomeryVaughan}
H. L. Montgomery and R. C. Vaughan,
\emph{Multiplicative Number Theory I: Classical Theory},
Cambridge Studies in Advanced Mathematics 97,
Cambridge University Press, 2007.


\bibitem{SV}
J.~M.~Sepulcre and T.~Vidal,
\emph{On the real projections of zeros of analytic almost periodic
functions},
Carpathian J. Math. \textbf{38} (2022), no.~2, 489--501.
\href{https://doi.org/10.37193/CJM.2022.02.18}{doi:10.37193/CJM.2022.02.18}.
Preprint: \url{https://arxiv.org/abs/1805.02041}.


\bibitem{AriasLune}
J. Arias de Reyna and J. van de Lune,
\emph{Some bounds and limits in the theory of Riemann's zeta function},
J. Math. Anal. Appl. \textbf{396} (2012), 199--214.
\href{https://doi.org/10.1016/j.jmaa.2012.06.017}{doi:10.1016/j.jmaa.2012.06.017}.
\href{https://arxiv.org/abs/1107.5134}{arXiv:1107.5134}.


\bibitem{Bounemoura}
A. Bounemoura,
\emph{Ergodization time for linear flows on tori via geometry of numbers},
preprint (2015).
\href{https://arxiv.org/abs/1506.06469}{arXiv:1506.06469}.


\bibitem{DumasFischler}
H. S. Dumas and S. Fischler,
\emph{Filling times for linear flow on the torus with truncated
Diophantine conditions: a brief review and new proof}, preprint (2021).
\href{https://arxiv.org/abs/2112.11786}{arXiv:2112.11786}.


\bibitem{GalatoloPeterlongo}
S. Galatolo and P. Peterlongo,
\emph{Long hitting time, slow decay of correlations and arithmetical
properties}, Discrete Contin. Dyn. Syst. \textbf{27} (2010), 185--204.
\href{https://doi.org/10.3934/dcds.2010.27.185}{doi:10.3934/dcds.2010.27.185}.


\bibitem{FejerMulti}
L. J. Patton and M. Putinar,
\emph{On multivariable Fej\'er inequalities},
author manuscript, 20 September 2005.
\url{https://web.math.ucsb.edu/~mputinar/fejermultivarrev.pdf}.


\bibitem{Bourgain}
J. Bourgain,
\emph{On invariant tori of full dimension for 1D periodic NLS},
J. Funct. Anal. \textbf{229} (2005), 62--94.
\href{https://doi.org/10.1016/j.jfa.2004.10.019}{doi:10.1016/j.jfa.2004.10.019}.


\bibitem{AlpogeFurman}
L. Alp\"oge and R. Furman,
\emph{More than two thirds of the zeros of the Riemann zeta function
are simple and on the critical line}, preprint (2026),
\href{https://arxiv.org/abs/2608.13637}{arXiv:2608.13637}.


\bibitem{Lamzouri}
Y. Lamzouri,
\emph{A new proof that more than $2/3$ of the zeros of the Riemann
zeta function are simple and on the critical line}, preprint (2026),
\href{https://arxiv.org/abs/2609.02882}{arXiv:2609.02882}.


\bibitem{Knausgard}
K. M. Knausg\aa rd,
\emph{More than 83.9\% of the zeros of the Riemann zeta function
are distinct and more than 67.35\% are simple and on the critical line},
preprint (2026),
\href{https://arxiv.org/abs/2610.08965}{arXiv:2610.08965}.


\bibitem{OpenAIQuasiRH}
OpenAI,
\emph{The Quasi-Riemann Hypothesis: A Zero-Free Half-Plane
$\Re(s)>7/8$}, research manuscript, 30 September 2026,
Theorem 1.1.
\href{https://github.com/openai/math/blob/main/preprints/The-Quasi-Riemann-Hypothesis-September-30-2026/paper.pdf}{Dated manuscript in the official repository}.
The complete external proof is not audited here. Applications that
invoke its $7/8$ conclusion are explicitly conditional on that input;
the general zero-free-region transfers are stated with their hypotheses. The shrinking-target results do not require that conclusion.


\bibitem{FordZaharescu}
K. Ford and A. Zaharescu,
\emph{On the distribution of imaginary parts of zeros of the
Riemann zeta function}, Journal f\"ur die reine und angewandte
Mathematik \textbf{579} (2005), 145--158.
\href{https://ford126.web.illinois.edu/wwwpapers/fz.pdf}{Author manuscript}.


\bibitem{AkbaryMurty}
A. Akbary and M. R. Murty,
\emph{Uniform distribution of zeros of Dirichlet series},
author-hosted manuscript,
\url{https://www.cs.uleth.ca/~akbary/uniform.pdf}.


\bibitem{CicekGonek}
F. \c{C}i\c{c}ek and S. M. Gonek,
\emph{The uniform distribution modulo one of certain subsequences
of ordinates of zeros of the zeta function}, Mathematical Proceedings
of the Cambridge Philosophical Society \textbf{176} (2024), 593--608,
\href{https://doi.org/10.1017/S0305004124000045}{doi:10.1017/S0305004124000045}.


\bibitem{Hilbert}
H. L. Montgomery and R. C. Vaughan,
\emph{Hilbert's inequality}, Journal of the London Mathematical
Society (2) \textbf{8} (1974), 73--82.


\bibitem{Jesse}
Jesse/SolidSF,
\emph{Autonomous Millennium Prize Lemma Bank}, public online
conditional spectral certificates,
\url{https://jesse.my/lemmas}, accessed 13 October 2026.
The $16/21$ constant and a quartic certificate already appear there.
No external numerical or formalization claim from this resource is
used as a premise of our finite-dimensional proofs.


\bibitem{Aurzada}
F. Aurzada, \emph{On the lower tail probabilities of some random
sequences in $\ell_p$}, Journal of Theoretical Probability
\textbf{20} (2007), 843--858.
\href{https://www2.mathematik.tu-darmstadt.de/~aurzada/independent.pdf}{Author manuscript}.


\bibitem{FordMengZaharescu}
K. Ford, X. Meng and A. Zaharescu, \emph{Simultaneous distribution
of the fractional parts of Riemann zeta zeros}, Bulletin of the
London Mathematical Society \textbf{49} (2017), 1--9.
\href{https://arxiv.org/abs/1511.06814}{arXiv:1511.06814}.


\bibitem{FordZeroFree}
K. Ford, \emph{Zero-free regions for the Riemann zeta function},
in \emph{Number Theory for the Millennium, II}, A K Peters,
2002, 25--56.
\href{https://ford126.web.illinois.edu/wwwpapers/zeros.pdf}{Author manuscript}.


\bibitem{DongInverse}
Z. Dong, R. Wang, W. Wang and H. Zhang,
\emph{Large zeta sums and zeros of the Riemann zeta function},
preprint, 2026.
\href{https://arxiv.org/abs/2608.31060}{arXiv:2608.31060}.


\bibitem{FillTorcaso}
J. A. Fill and F. Torcaso, \emph{Asymptotic analysis via Mellin
transforms for small deviations in $L^2$-norm of integrated Brownian
sheets}, Probability Theory and Related Fields \textbf{130} (2004),
259--288. \href{https://arxiv.org/abs/math/0309391}{arXiv:math/0309391}.


\bibitem{JMRS}
M. Jacobsen, T. Mikosch, J. Rosi\'nski, and G. Samorodnitsky,
\emph{Inverse problems for regular variation of linear filters, a
cancellation property for $\sigma$-finite measures and identification
of stable laws}, Annals of Applied Probability \textbf{19} (2009),
210--242. \href{https://arxiv.org/abs/0712.0576}{arXiv:0712.0576}.


\bibitem{JRY}
L. F. James, B. Roynette, and M. Yor, \emph{Generalized Gamma
Convolutions, Dirichlet means, Thorin measures, with explicit examples},
Probability Surveys \textbf{5} (2008), 346--415.
\href{https://arxiv.org/abs/0708.3932}{arXiv:0708.3932}.


\bibitem{NP}
I. Nourdin and G. Poly, \emph{Convergence in law in the second
Wiener/Wigner chaos}, Electronic Communications in Probability
\textbf{17} (2012), paper 36, 1--12. Corrected version:
\href{https://arxiv.org/abs/1205.2684v3}{arXiv:1205.2684v3}.


\bibitem{AEK}
E. Azmoodeh, P. Eichelsbacher, and L. Knichel, \emph{Optimal Gamma
Approximation on Wiener Space}, ALEA, Latin American Journal of
Probability and Mathematical Statistics \textbf{17} (2020), 101--132.
\url{https://alea.impa.br/articles/v17/17-05.pdf}.


\bibitem{SEM}
J. Szavits-Nossan, M. R. Evans, and S. N. Majumdar,
\emph{Condensation transition in joint large deviations of linear
statistics}, Journal of Physics A: Mathematical and Theoretical
\textbf{47} (2014), 455004.
\href{https://arxiv.org/abs/1406.3573}{arXiv:1406.3573}.


\bibitem{DLMF}
NIST Digital Library of Mathematical Functions, Chapter 5,
Section 5.8, Infinite products for the Gamma function.
\url{https://dlmf.nist.gov/5.8}.


\bibitem{BRT}
O. E. Barndorff-Nielsen, J. Rosi\'nski and S. Thorbj\o rnsen,
\emph{General $\Upsilon$-transformations}, ALEA \textbf{4} (2008), 131--165.
\url{https://alea.impa.br/articles/v4/04-07.pdf}.


\bibitem{BPT}
O. E. Barndorff-Nielsen, V. P\'erez-Abreu and S. Thorbj\o rnsen,
\emph{L\'evy mixing}, ALEA \textbf{10}(2) (2013), 1013--1062.
\url{https://alea.math.cnrs.fr/articles/v10/10-41.pdf}.


\bibitem{EM}
S. N. Evans and I. Molchanov,
\emph{Polar decomposition of scale-homogeneous measures with application
to L\'evy measures of strictly stable laws}, arXiv:1509.09261 (2015).
\url{https://arxiv.org/abs/1509.09261}.


\bibitem{BKL}
D. Berger, M. Kutlu and A. Lindner,
\emph{On multivariate quasi-infinitely divisible distributions},
arXiv:2101.02544 (2021).
\url{https://arxiv.org/abs/2101.02544}.


\bibitem{Jurek}
Z. J. Jurek,
\emph{Calculus on random integral mappings $I^{h,r}_{(a,b]}$ and their domains},
arXiv:1210.5990 (2012).
\url{https://arxiv.org/abs/1210.5990}.


\bibitem{BP}
A. Berman and R. J. Plemmons,
\emph{Nonnegative Matrices in the Mathematical Sciences},
Classics in Applied Mathematics 9, SIAM, Philadelphia, 1994.
\url{https://doi.org/10.1137/1.9781611971262}.


\bibitem{Cuesta2007}
J. A. Cuesta-Albertos, R. Fraiman, and T. Ransford,
A sharp form of the Cram\'er--Wold theorem,
\emph{Journal of Theoretical Probability} \textbf{20} (2007), 201--209.
\href{https://personales.unican.es/cuestaj/CramerWold.pdf}{Author version}.


\bibitem{Hausdorff1921}
F. Hausdorff,
Summationsmethoden und Momentfolgen. I, II,
\emph{Mathematische Zeitschrift} \textbf{9} (1921), 74--109, 280--299.


\bibitem{Kleiber2013}
C. Kleiber,
The generalized lognormal distribution and the Stieltjes moment problem,
\emph{Journal of Theoretical Probability} \textbf{27} (2014), 1167--1177.
\href{https://arxiv.org/abs/1301.1277}{arXiv:1301.1277}.


\bibitem{Lenzi2026}
J. Lenzi,
\emph{One marked-sum law determines random mass partitions and
$\Xi$-coalescents}, preprint, version 1, 11 September 2026.
\href{https://arxiv.org/abs/2609.12790v1}{arXiv:2609.12790v1}.
Related version: \href{https://doi.org/10.5281/zenodo.22708550}{doi:10.5281/zenodo.22708550}.


\bibitem{Lin2016}
J. Lin,
A negative answer to a problem of Aldous on determination of
exchangeable sequences,
\emph{Electronic Journal of Probability} \textbf{21} (2016), paper 30,
1--26.
\href{https://arxiv.org/abs/1507.00678}{arXiv:1507.00678}.


\bibitem{PenskySapatinas2011}
M. Pensky and T. Sapatinas,
Multichannel boxcar deconvolution with growing number of channels,
\emph{Electronic Journal of Statistics} \textbf{5} (2011), 53--82.
\href{https://doi.org/10.1214/11-EJS597}{doi:10.1214/11-EJS597}.


\bibitem{Rao2007}
N. Raj Rao,
The analytic computability of the Shannon transform for a large class
of random matrix channels,
preprint, 2007.
\href{https://arxiv.org/abs/0712.0305}{arXiv:0712.0305}.


\bibitem{Stoyanov2004}
J. Stoyanov,
Stieltjes classes for moment-indeterminate probability distributions,
\emph{Journal of Applied Probability} \textbf{41A} (2004), 281--294.


\bibitem{TulinoVerdu2004}
A. M. Tulino and S. Verd\'u,
Random matrix theory and wireless communications,
\emph{Foundations and Trends in Communications and Information Theory}
\textbf{1} (2004), no. 1, 1--182.


\bibitem{BoivinGauthierManolaki2018}
A. Boivin, P. M. Gauthier, and M. Manolaki,
\emph{On zero sets of harmonic and real analytic functions},
Annales math\'ematiques du Qu\'ebec \textbf{42} (2018), 159--167.
\href{https://doi.org/10.1007/s40316-018-0098-1}{doi:10.1007/s40316-018-0098-1}.


\bibitem{Kostic2025}
M. Kosti\'c,
\emph{Uniqueness sequences for multidimensional vector-valued Laplace transform},
preprint (2025),
\href{https://arxiv.org/abs/2510.17485}{arXiv:2510.17485}.


\bibitem{KriegNovakUllrich2025}
D. Krieg, E. Novak, and M. Ullrich,
\emph{On the power of adaption and randomization},
Forum of Mathematics, Sigma \textbf{13} (2025), e152.
\href{https://doi.org/10.1017/fms.2025.10101}{doi:10.1017/fms.2025.10101}.


\bibitem{PottsTasche2010}
D. Potts and M. Tasche,
\emph{Parameter estimation for exponential sums by approximate Prony method},
Signal Processing \textbf{90} (2010), 1631--1642.
\href{https://doi.org/10.1016/j.sigpro.2009.11.012}{doi:10.1016/j.sigpro.2009.11.012}.


\bibitem{PottsTasche2013}
D. Potts and M. Tasche,
\emph{Parameter estimation for nonincreasing exponential sums by
Prony-like methods},
Linear Algebra and its Applications \textbf{439} (2013), 1024--1039.
\href{https://doi.org/10.1016/j.laa.2012.10.036}{doi:10.1016/j.laa.2012.10.036}.


\bibitem{Lin2015}
J. Lin,
\emph{A negative answer to a problem of Aldous on determination of
exchangeable sequences},
\href{https://arxiv.org/abs/1507.00678v1}{arXiv:1507.00678v1} (2015).
Section and proposition numbers here refer to this preprint.
Published in Electronic Journal of Probability \textbf{21} (2016),
paper 30, 1--26,
\href{https://doi.org/10.1214/16-EJP4403}{doi:10.1214/16-EJP4403}.


\bibitem{LeBarilettoRinaldoHo2026}
D. Le, N. Bariletto, A. Rinaldo, and N. Ho,
\emph{Partial Differential Equation Barriers to Identifiability in
Infinite Mixture Models},
\href{https://arxiv.org/abs/2608.08597}{arXiv:2608.08597} (2026).
Definitions and question locations refer to the
\href{https://nhatptnk8912.github.io/Identifiability_Infinite_Mixture_final.pdf}{author manuscript}
dated 8 August 2026.


\bibitem{Nolan2018}
J. P. Nolan,
\emph{Stable Distributions: Models for Heavy Tailed Data},
Chapter 1, author manuscript processed 30 January 2018,
Definition 1.8 and equation (1.6).
\href{https://edspace.american.edu/jpnolan/wp-content/uploads/sites/1720/2020/09/Chap1.pdf}{Author's chapter PDF}.


\bibitem{BellLagarias2015}
J. P. Bell and J. C. Lagarias,
\emph{A Skolem--Mahler--Lech theorem for iterated automorphisms of
$K$-algebras}, Canadian Journal of Mathematics \textbf{67}(2)
(2015), 286--314.
\href{https://doi.org/10.4153/CJM-2013-048-3}{doi:10.4153/CJM-2013-048-3}.
\href{https://arxiv.org/abs/1312.5305}{arXiv:1312.5305}.


\bibitem{Bonnier2023}
P. Bonnier, H. Oberhauser, and Z. Szab\'o,
\emph{Kernelized cumulants: Beyond kernel mean embeddings},
Advances in Neural Information Processing Systems \textbf{36} (2023).
\href{https://arxiv.org/abs/2301.12466}{arXiv:2301.12466}.


\bibitem{Fukumizu2008}
K. Fukumizu, B. Sriperumbudur, A. Gretton, and B. Sch\"olkopf,
\emph{Characteristic kernels on groups and semigroups},
Advances in Neural Information Processing Systems \textbf{21} (2008),
473--480.
\href{https://papers.neurips.cc/paper/3396-characteristic-kernels-on-groups-and-semigroups.pdf}{Conference paper}.


\bibitem{MiaoEtAl2024}
Z. Miao, W. Kong, R. K. Vinayak, W. Sun, and F. Han,
\emph{Fisher--Pitman permutation tests based on nonparametric Poisson
mixtures with application to single cell genomics},
Journal of the American Statistical Association \textbf{119}(545)
(2024), 394--406.
\href{https://doi.org/10.1080/01621459.2022.2120401}{doi:10.1080/01621459.2022.2120401}.
Rate comparison refers to Theorems~6.1--6.2 of the
\href{https://ramyakv.github.io/testing-poisson-npmle.pdf}{author manuscript}
dated 8 June 2021.


\bibitem{Pinkus2004}
A. Pinkus,
\emph{Strictly Hermitian positive definite functions},
Journal d'Analyse Math\'ematique \textbf{94} (2004), 293--318.
\href{https://arxiv.org/abs/math/0404013}{arXiv:math/0404013}.


\bibitem{VandermeulenScott2015}
R. A. Vandermeulen and C. D. Scott,
\emph{On the identifiability of mixture models from grouped samples},
\href{https://arxiv.org/abs/1502.06644}{arXiv:1502.06644}
(2015; revised 2022). The preprint reports that the work was subsumed
and expanded in their Annals of Statistics article,
\emph{An operator theoretic approach to nonparametric mixture models}
(2019), \href{https://doi.org/10.1214/18-AOS1762}{doi:10.1214/18-AOS1762}.


\bibitem{VandermeulenSaitenmacher2022}
R. A. Vandermeulen and R. Saitenmacher,
\emph{Generalized identifiability bounds for mixture models with grouped
samples} (2022).
\href{https://arxiv.org/abs/2207.11164}{arXiv:2207.11164}.


\bibitem{BorweinErdelyi1995}
P.~Borwein and T.~Erd\'elyi,
M\"untz spaces and Remez inequalities,
\emph{Bulletin of the American Mathematical Society (N.S.)}
\textbf{32} (1995), 38--42.
\href{https://doi.org/10.1090/S0273-0979-1995-00553-7}{doi:10.1090/S0273-0979-1995-00553-7}.
\href{https://arxiv.org/abs/math/9501225}{arXiv:math/9501225}.


\bibitem{BurginEpperleinWirth2022}
V.~B\"urgin, J.~Epperlein, and F.~Wirth,
Remarks on the tail order on moment sequences,
\emph{Journal of Mathematical Analysis and Applications}
\textbf{512} (2022), no.~1, 126135.
\href{https://doi.org/10.1016/j.jmaa.2022.126135}{doi:10.1016/j.jmaa.2022.126135}.
\href{https://arxiv.org/abs/2104.10572v3}{arXiv:2104.10572v3}
(revised 21 March 2022); see Conjecture~17.


\bibitem{StoyanovLin2011}
J.~Stoyanov and G.~D.~Lin,
Mixtures of power series distributions: identifiability via uniqueness
in problems of moments,
\emph{Annals of the Institute of Statistical Mathematics}
\textbf{63} (2011), 291--303.
\href{https://doi.org/10.1007/s10463-009-0221-9}{doi:10.1007/s10463-009-0221-9}.


\bibitem{Vinayak2019}
R.~K.~Vinayak, W.~Kong, G.~Valiant, and S.~Kakade,
Maximum likelihood estimation for learning populations of parameters,
\emph{Proceedings of Machine Learning Research}
\textbf{97} (2019), 6448--6457.
\href{https://proceedings.mlr.press/v97/vinayak19a.html}{Proceedings and full paper}.

\end{thebibliography}
\end{document}
